\documentclass[10pt]{amsart}
\usepackage{amsmath,amssymb,amsthm,mathtools,mathrsfs}
\usepackage[a4paper,hmargin=0.76in,vmargin=0.60in]{geometry}
\usepackage{booktabs}
\usepackage{array}
\usepackage{longtable}
\usepackage{tabularx}
\usepackage[hypertexnames=false]{hyperref}
\hypersetup{hidelinks}
\usepackage[final,expansion=false]{microtype}
\usepackage{url}
\emergencystretch=2.5em

\newtheorem{theorem}{Theorem}[section]
\newtheorem{proposition}[theorem]{Proposition}
\newtheorem{lemma}[theorem]{Lemma}
\newtheorem{corollary}[theorem]{Corollary}
\theoremstyle{definition}
\newtheorem{definition}[theorem]{Definition}
\newtheorem{problem}[theorem]{Problem}
\theoremstyle{remark}
\newtheorem{remark}[theorem]{Remark}

\newcommand{\C}{\mathbb C}
\newcommand{\R}{\mathbb R}
\newcommand{\N}{\mathbb N}
\newcommand{\Z}{\mathbb Z}
\newcommand{\Q}{\mathbb Q}
\newcommand{\Rel}{\operatorname{Re}}
\newcommand{\Iml}{\operatorname{Im}}
\newcommand{\supp}{\operatorname{supp}}
\newcommand{\TV}{\operatorname{TV}}
\newcommand{\Kp}{\mathcal K^{\mathrm{prim}}_{2T}}
\newcommand{\Kw}{\mathcal K_{2T}}
\newcommand{\PT}{\mathcal P_T}
\newcommand{\IT}{I_T}
\newcommand{\lamP}{\lambda^{\mathrm{prim}}_T}
\newcommand{\lamW}{\lambda^{\mathrm W}_T}
\newcommand{\lamA}{\lambda_{\mathrm A}(T)}
\newcommand{\sT}{\sigma_T}
\newcommand{\gh}[1]{\widehat g_{#1}}
\newcommand{\nuh}{\widehat\nu}
\newcommand{\wh}{\widehat w}
\newcommand{\Fh}{\widehat F}
\newcommand{\cF}{c_F}
\newcommand{\Lev}{\mathfrak E}
\newcommand{\Wcl}{\mathcal W_T}
\newcommand{\NT}{\mathcal N_T}
\newcommand{\Aconst}{\mathsf A_T}
\newcommand{\Ainf}{\mathcal A_\infty}
\newcommand{\PiT}{\Pi_T}
\newcommand{\mT}{\mathfrak m_T}
\newcommand{\laminf}{\lambda^{\infty}_{\min}}
\newcommand{\laminfP}{\lambda^{\infty,\mathrm P}_{\min}}
\newcolumntype{Y}{>{\raggedright\arraybackslash}X}

\providecommand{\anCP}{C_P}
\providecommand{\anPa}{\Phi_a}
\providecommand{\anPb}{\Phi_b}
\providecommand{\anPl}{\Phi}
\providecommand{\ancl}{\mathrm{cl}}
\providecommand{\aneu}{\mathrm{eu}_d}
\providecommand{\anrd}{\mathrm{rd}}
\providecommand{\auAm}{\mathcal A_m}
\providecommand{\auAz}{\mathcal A_0}
\providecommand{\auLat}{\mathcal L}
\providecommand{\auNm}{\operatorname{Nm}}
\providecommand{\auO}{\mathcal O}
\providecommand{\auP}{\mathcal P}
\providecommand{\auPa}{\Phi_a}
\providecommand{\auPl}{\Phi}
\providecommand{\auS}{\mathcal S}
\providecommand{\auSram}{\mathcal S^{\mathrm{ram}}}
\providecommand{\auT}{\mathcal T}
\providecommand{\auU}{\mathrm U}
\providecommand{\aubar}[1]{\overline{#1}}
\providecommand{\aucB}{c_B}
\providecommand{\auchi}{\chi_K}
\providecommand{\audT}{d_{\auT}}
\providecommand{\auerr}{\operatorname{err}}
\providecommand{\augain}{\operatorname{gain}}
\providecommand{\auhdir}{h_{\mathrm{dir}}}
\providecommand{\auhmain}{h_{\mathrm{main}}}
\providecommand{\aupp}{\mathfrak p}
\providecommand{\aurd}{\mathrm{rd}_E}
\providecommand{\aurem}{\operatorname{rem}}
\providecommand{\ausfin}{s_{\mathrm{fin}}}
\providecommand{\ausinf}{s_\infty}
\providecommand{\ausmain}{s_{\mathrm{main}}}
\providecommand{\pbdeg}{\operatorname{deg}}
\providecommand{\pbeq}{\mathsf q}
\providecommand{\pcMat}{\mathrm M}
\providecommand{\pcR}{R}
\providecommand{\pcW}{\mathbb W}
\providecommand{\pcm}{\mathfrak m}
\providecommand{\pdsK}{\mathbf K}
\providecommand{\pdsm}{m}
\providecommand{\pfqm}{\mathfrak m_\vartheta}
\providecommand{\pguK}{\mathsf K}
\providecommand{\pgupt}{\mathsf u}
\providecommand{\phA}{\mathfrak A}
\providecommand{\phB}{\mathfrak B}
\providecommand{\phcd}{c_{\mathrm{div}}}
\providecommand{\phdef}{\operatorname{def}}
\providecommand{\phk}{\boldsymbol k}
\providecommand{\pmP}{\mathsf P}
\providecommand{\ptvtr}{\operatorname{tr}}
\newenvironment{quotedthm}[2]{\par\medskip\noindent\begingroup
  \textbf{#1}#2.\ \itshape\ignorespaces}{\endgroup\par\medskip}
\newenvironment{quoteddef}[2]{\par\medskip\noindent\begingroup
  \textbf{#1}#2.\ \ignorespaces}{\endgroup\par\medskip}
\usepackage{xurl}
\begin{document}
\addtolength{\abovedisplayskip}{-3.4pt}
\addtolength{\belowdisplayskip}{-3.4pt}
\addtolength{\belowdisplayshortskip}{-2pt}
\title[The Short Forms Theorem]{The Short Forms Theorem: theta lifts, complex multiplication and divisibility}
\author{David Alejandro Trejo Pizzo}
\date{}
\begin{abstract}
The Short Forms Theorem is the first of three theorems on integral algebraic modular forms from
which the main paper of this series deduces the Riemann hypothesis for Dirichlet $L$-functions.
It concerns a Poincar\'e map from forms on a definite unitary group over a CM field to forms on an
indefinite one, followed by a chain of theta lifts, and says that suitable multiples of the images
are integral up to a bounded index and short at every conjugate. We prove the statements from which
it is deduced. A seesaw identity replaces the chain by lifts between definite groups, whose kernels
are derivatives at the origin of products of theta functions of one elliptic curve with complex
multiplication. We bound all conjugates of these derivatives, with the Chowla--Selberg period as
exact constant, and we determine their divisibility at every prime, in particular at the primes
that do not split in the field of multiplication and at the prime of the level. The paper also
contains the facts on the Weil representation and on discrete series of real unitary groups that
are used throughout the series.
\end{abstract}
\maketitle
\setcounter{tocdepth}{1}\tableofcontents

\section{Introduction}\label{A1:intro}

The main paper~\cite{P46} of this series deduces the Riemann hypothesis for Dirichlet
$L$-functions from three theorems on integral algebraic modular forms: the Short Forms Theorem,
the Spanning Theorem and the Covolume Theorem. All three concern one lattice of automorphic forms
on a unitary group over a CM field, defined by the integrality of the images under a chain of
theta lifts. The first two theorems produce many short vectors in this lattice, the third bounds
its covolume from below, and the two bounds differ by the remainder of a sum over primes, which
therefore cannot be large and positive. In \cite{P46} each of the three theorems is deduced from a
list of propositions, which are proved in the other papers of the series: those for the Spanning
Theorem in \cite{P51} and~\cite{P52}, those for the Covolume Theorem in \cite{P53}, \cite{P54},
\cite{P55} and~\cite{P56}, and those for the Short Forms Theorem in this paper, with one
lemma on Hecke operators in~\cite{P50}.

We recall the objects; the definitions are in {\cite[\S\S14.1--14.3]{P46}}. Let
$K=\Q(\sqrt{-d})$ be an imaginary quadratic field and $E=K(\zeta_{q_1})$, a CM field with a set
$\auPl$ of $f=[E:K]$ infinite places. For $0\le j\le m$ let $X_j$ be the standard positive
definite hermitian space of rank $n_j=k+j(t+1)$ over $E$, and let $Y_j=X_j\perp S_j$, where $S_j$
is a sum of hyperbolic planes and at most one line, of rank $m-1-j$; for $j=m$ the roles are
exchanged and $Y_m$ is a hyperplane of $X_m$. Thus $\auU(Y_0)$ has signature $(k+\varrho,
\varrho+1)$ at every infinite place, where $m=2\varrho+2$, the group $\auU(Y_m)$ is definite, and
consecutive groups of the \emph{upper chain} $\auU(X_0),\dots,\auU(X_m)$, and of the \emph{lower
chain} $\auU(Y_0),\dots,\auU(Y_m)$, form dual reductive pairs. The \emph{source forms} are the
algebraic modular forms, in the sense of Gross~\cite{GrossAMF}, on the definite group
$\auU(X_0)$, with values in the tensor product of the duals of the spaces of forms of degree $a$
in $k$ variables at a set $\auPa$ of infinite places, and with a level that involves $f$ primes of
$E$ above a prime $N$ which splits completely in $E$, the divisors of an element $\pi_N$ of norm
$N$ to $E^+$; $\auS$ denotes their space and $\auSram$ the subspace of the forms that are ramified
at all these primes. The \emph{Poincar\'e map} $\auP$
sends a source form to the projection of a Poincar\'e series of Bessel type to the space $\auAz$
of cusp forms on $\auU(Y_0)$, of a fixed level, that generate a fixed discrete series at the
infinite places. For every twist $\vartheta$, an automorphic character of finite order of the unitary group
of rank one, there is a composite $\iota_\vartheta$ of $m$ theta lifts along the lower chain; it
maps $\auAz$ to algebraic modular forms on $\auU(Y_m)$, and a form in $\auAz$ is called integral
at a prime if its images under the maps $\iota_\vartheta$, for the twists of a certain class that
are unramified above this prime, are integral there in the integral structure
of~\cite{GrossAMF}.

The Short Forms Theorem ({\cite[Theorem~14.13]{P46}}) is a statement about the map
$\aucB\,c_{\mathrm{ss}}^{-1}\auP$ on $\auSram$, projected to a direct summand $\auT$ of $\auAz$;
here $c_{\mathrm{ss}}$ is a non-zero complex number and $\aucB$ a positive rational number, a
product of powers of $N$ and of the primes $p<x_1$ that do not split in $K$. It says that
$\auSram$ is spanned by integral forms $\varphi$ whose images $y$, and the translates of these
images by a finite group acting on $\auT$, satisfy for every automorphism $\sigma$ of $\C$
\[
\|\iota_1(y)\|_\sigma^2\le\exp\bigl(\auhmain-m\cdot\aurem+f\,a\,(\log b)^4\bigr),
\]
and that the images of integral forms are integral up to an explicitly bounded index. Here
$\auhmain$ is a main term, which reappears in the Covolume Theorem, $\aurem$ is the remainder of a
sum over the powers of the primes that do not split in $K$, and the last term is an admissible
error. The deduction of the theorem in {\cite[\S14.7]{P46}} uses the following
statements, which we now describe together with their proofs.

\subsection*{The seesaw identity}
The lifts of the lower chain are theta lifts between indefinite groups, of which little can be
computed directly. On the image of $\auP$, however, they can be replaced by lifts along the upper
chain. For $0\le j<m$ and a Schwartz--Bruhat function $\phi$ let $\Theta^X_j(\phi)$ be the theta
lift from algebraic modular forms on $\auU(X_j)$ to those on $\auU(X_{j+1})$. Both groups are
compact at the infinite places, so the kernel of this lift is a matrix whose entries are
derivatives at the origin, along polynomials with integral coefficients, of theta series of positive definite hermitian lattices over the
integers of $E$. A \emph{composite lift} $B_{\vartheta,\Psi}$ is a finite linear combination, with
coefficients given by a function $\Psi$, of the composites $\Theta^X_{m-1}(\phi_{m-1})\circ\dots
\circ\Theta^X_0(\phi_0)$ in which each $\phi_j$ is a translate, under the Heisenberg group, of the
\emph{vacuum}, the normalised characteristic function of the lattice of integral homomorphisms;
it is followed by the restriction of forms from $\auU(X_m)$ to $\auU(Y_m)$, a projection of the
weights, and a division by powers of the conjugates of $\aubar{\pi_N}$. The seesaw identity ({\cite[Proposition~14.19]{P46}}) states that there
is a constant $c_{\mathrm{ss}}\in\C^\times$ such that for every twist $\vartheta$ there is such a
function $\Psi$, a \emph{chain datum}, with
\begin{equation}\label{A1:eq-seesaw}
\iota_\vartheta(\auP\varphi)=c_{\mathrm{ss}}\cdot B_{\vartheta,\Psi}\,\varphi\qquad\text{for all }
\varphi\in\auSram .
\end{equation}
The constant does not depend on $\vartheta$, the values of $\Psi$ are algebraic numbers that are
integral at all primes above which $\vartheta$ is unramified, and for the trivial twist the sum of
their absolute values is bounded at every conjugate.

The proof (\S\ref{ps:ss-seesaw} and \S\ref{pk:sec}) removes one lift at a time. Pairing the
left side of \eqref{A1:eq-seesaw} with a form on $\auU(Y_m)$ and moving the theta kernels to that
form, one has to compute a Bessel period of an $m$-fold theta lift from $\auU(Y_m)$ down to
$\auU(Y_0)$. The summand $S_j$ defines a subgroup $R_j$ of $\auU(Y_j)$, with reductive part $\auU(X_j)$, and
periods over $R_j$ which are of Bessel type if the rank of
$S_j$ is odd and of Fourier--Jacobi type if it is even. Proposition~\ref{ps:onestep} shows that
the period over $R_j$ of the theta lift to $\auU(Y_j)$ of a rapidly decreasing automorphic
function $z$ on $\auU(Y_{j+1})$,
against a function $\mathcal W$ on $R_j$, is the period of $z$ itself over $R_{j+1}$ against a
function $\mathcal W'$, which is obtained from $\mathcal W$ by an integral over $R_j$ against a
coefficient of the theta kernel. To prove it we expand the theta kernel in a mixed model of the
Weil representation along an isotropic space of homomorphisms from $Y_j$ to $Y_{j+1}$, show that
after the integration over $R_j$ only one orbit of rational points contributes, describe this
orbit by Witt's theorem, and unfold. This is a global form, for the nilpotent orbits that occur here, of the
transfer of generalised Whittaker models under the theta correspondence established locally by
Gomez and Zhu~\cite{GomezZhu}; it combines a seesaw in the sense of Kudla~\cite{KudlaSeesaw} with
an unfolding. After $m$ steps the summand $S_j$ is exhausted, the reductive parts $\auU(X_j)$ have run through
the upper chain, and one is left with the upper chain applied to $\varphi$ and evaluated on
$Y_m\subset X_m$. The integrals over the real groups give the constant $c_{\mathrm{ss}}$, and
the sums over cosets at the finite places give the chain datum. The delicate points are the
convergence of the unfolded integrals, the independence of the constant of $\vartheta$
(Lemmas~\ref{ps:double} and~\ref{ps:fin}), and the primes above $N$, where consecutive lifts have
to be glued in pairs (\S\ref{ps:ss-glue}) and where the hypothesis that $\varphi$ is ramified is
used.

For \eqref{A1:eq-seesaw} to have a meaning the lower chain must be defined: the theta integrals
over the non-compact quotients have to converge, each lift has to be a cusp form of the prescribed
type at the infinite places, and the last lift has to be an algebraic modular form. We define the
forms of the type of the chain on $\auU(Y_j)$ as the bounded automorphic functions reproduced by
convolution with the coefficient of the lowest $K$-type of a discrete series with an explicit
Harish-Chandra parameter, multiplied by the formal degree, and we prove in
Proposition~\ref{pz:types} that these forms are cuspidal and rapidly decreasing and that the
lifts of the chain and their transposes preserve the types. This rests on three properties of
discrete series of the groups $\auU(p,q)$ whose parameter is far from the walls: the integrability
of the coefficient and of its derivatives against powers of the norm (Theorem~\ref{par:A1}), the
vanishing of its integrals over unipotent radicals (Corollary~\ref{par:A2}), and the equality, on
a polynomial of the Fock model that is harmonic for both groups of a dual pair, of the projections
to the types of the two groups (Theorem~\ref{par:A3}). The parameters of the chain have gaps of size at least $\tfrac12r^8$
between entries attached to coordinates of different signs, where $r$ is the rank of $Y_0$, and
for such parameters we prove the three properties from an explicit integral formula for the
coefficient of the lowest $K$-type (Theorem~\ref{pco:FJ}; this is Flensted-Jensen's description
of that coefficient~\cite{FJ}, for which we give a proof that does not rely on~\cite{FJ}), from
Schmid's theorem on the $K$-types of a discrete
series~\cite{Schmid75}, the Dirac inequality in the form of~\cite{HuangPandzic} and a result of
Vogan and Zuckerman~\cite{VoganZuckerman}, and from the determination of the character by which
the Weil representation of a unitary dual pair differs from the metaplectic representation
(Theorem~\ref{pco:V}).

\subsection*{The size of the kernels}
By \eqref{A1:eq-seesaw} the Short Forms Theorem becomes a statement on the kernels of the upper
chain, at the infinite places and at every prime. Let $\theta_j(\phi;g,g')$ be the kernel of
$\Theta^X_j(\phi)$, a linear map between the spaces of values of the forms on $\auU(X_j)$ and on
$\auU(X_{j+1})$. {\cite[Proposition~14.23]{P46}} states that for every translate $\phi$ of the
vacuum, every automorphism $\sigma$ of $\C$ and all $g$, $g'$ and $\xi$,
\[
\|\sigma(\theta_j(\phi;g,g')\xi)\|\le\exp\Bigl(x^{50}+\tfrac12\sum_{\tau\in\auPl}\
\sum_{\text{rows }i}\ausinf(l_{\tau,i})\Bigr)\,\|\sigma(\xi)\|,\qquad
\ausinf(l)=\log l!+l\log\frac M\Omega,
\]
where the $l_{\tau,i}$ are the lengths of the rows of the weight of $X_{j+1}$, $M$ is the
auxiliary level of~\cite{P46} and
\[
\Omega=\exp\Bigl(\gamma-\Rel\frac{L'}{L}(1,\auchi)\Bigr)\cdot\frac{2\pi}{\aurd},
\]
with $\auchi$ the quadratic character of $K$ and $\aurd$ the root discriminant of $E$. The
constant $\Omega$ has to be exact: the lengths of the rows are of the size of the weights, about
$b$, while the error admitted in the exponent of the Short Forms Theorem is smaller than the main
term by a power of $x$, and the same constant occurs in the evaluation of the sum over primes with
which the bound is compared. The proof is in \S\ref{ps:ss-kernel}. Since $E$ is generated over $K$
by a root of unity, the theta series in question are finite combinations, with roots of unity as
coefficients, of products of translates by torsion points of the theta function $\theta$ of the
single elliptic curve $\C/\Omega_0\auO_K$, evaluated at linear forms given by the powers of that
root of unity (a \emph{resolution of the identity}, \S\ref{pa:ss-prod}). We normalise $\theta$ at a point of
order two, so that it does not vanish at the origin, and collect what we need about it in
Proposition~\ref{pk:curve}: its derivatives at the origin are algebraic numbers, integral at every prime
other than $d$; an automorphism $\sigma$ of $\C$ over $K$ maps the Taylor series of a translate of
$\theta$ by a torsion point to that of a translate of the theta function of the conjugate curve,
the new torsion point being given by an idele of $K$ (Theorem~\ref{pf:recip}, a reciprocity law
proved in \S\ref{pf:ss-recip}); and there is an algebraic number $\mathsf u$, a unit at every
prime, such that the area of the period lattice of every conjugate curve, multiplied by
$|\sigma\mathsf u|^2$, has the same value $2\pi^2e^{\gamma-\Rel(L'/L)(1,\auchi)}/\sqrt d$. The
last statement is the Chowla--Selberg formula~\cite{ChowlaSelberg}. Together they express every
conjugate of a derivative of a theta series at the origin as a Gaussian integral of a polynomial against
a bounded function, at the scale $\Omega^{-1/2}$ for all conjugates at once, and the bound
follows. Two further statements complete the archimedean estimate: the restriction from
$\auU(X_m)$ to $\auU(Y_m)$, with its division by a power of $N$, is controlled by an integral
form of the branching law ({\cite[Proposition~14.24]{P46}}, proved in \S\ref{ps:ss-restrict} with
Proposition~\ref{pk:branching}), and the passage to the summand $\auT$ by the equivariance of the
lower chain under the finite group that acts on $\auAz$ ({\cite[Proposition~14.27]{P46}}). The short
integral forms that span $\auSram$ are obtained from bases of the lattices in question with small
entries at all conjugates ({\cite[Lemma~14.21]{P46}}, proved in \S\ref{ps:ss-short}).

\subsection*{Divisibility}
The primes enter through the valuations of the same derivatives. Let $p$ be a prime that
does not split in $K$, so that the curve has supersingular reduction at $p$, let $v$ be the
$p$-adic valuation with $v(p)=1$, let $\pi_p$ be $p$ if $p$ is inert and $\sqrt{-d}$ if $p=d$,
and let $I_p$ be the set of the odd positive integers if $p$ is inert and of all positive integers
if $p=d$. The set $\auPl$ is a principal homogeneous space under $\operatorname{Gal}(E/K)=
(\Z/q_1)^\times$; for $e\in I_p$ the \emph{partner} $\tau'$ of $\tau\in\auPl$ at $p^e$ is its
translate by the class of $-p^e$. For inert $p$ this is $\tau$ composed with complex conjugation
and the $e$-th power of the Frobenius at $p$. We prove (Theorem~\ref{pa:onecurve} and
Proposition~\ref{pa:product}) that the theta products above, as power series $F$ in variables
$z_{\tau,i}$ indexed by the places $\tau\in\auPl$ and a finite set of indices $i$, have the
following normal form: for a number $\varpi$ of valuation $v(\pi_p)/(\auNm\pi_p-1)$,
\begin{equation}\label{A1:eq-normal}
F(\varpi z)=W(z)\cdot\exp\Bigl(\frac{\varpi^{p+1}}p\sum_{e}\sum_{i\ge0}\lambda_{e,i}\,
\Gamma_e^{\,p^i}\Bigr),\qquad\Gamma_e=\sum_{\tau,i}z_{\tau,i}^{\,p^e}\,z_{\tau',i},
\end{equation}
where $W$ and the $\lambda_{e,i}$ are integral and $e$ runs over $I_p$, and over $I_p\cup\{0\}$ if
$p=d$.
The factor $\varpi$ reflects the fact that the formal group of the curve has height two, and the
couplings $\Gamma_e$ come from the Frobenius: for one curve the exponent is determined, modulo
$\pi_p$, by a congruence between $\log\theta(z)$ and $\log\theta(\pi_pz)$
(Proposition~\ref{pa:frob}), and for a product the resolution of the identity turns the powers
$z^{p^e}\cdot z$ of one variable into the couplings of partners. From \eqref{A1:eq-normal} we
deduce a lower bound for the valuation of the derivatives $Q(\partial)F(0)$ of $F$ at the origin
along the entries $Q$ of a Schur matrix (Proposition~\ref{pa:count}, which is {\cite[Proposition~14.10]{P46}}): if the weight has
$a_\tau$ columns of height one at $\tau$ and $a'_\tau$ other columns, with $l_\tau$ boxes in these
and $B$ boxes in all, and all heights are less than $n'$, the valuation is at least
\[
v(\pi_p)\sum_{e\in I_p}\ \sum_{\tau\in\auPl}\Bigl(\bigl(\tfrac{a_\tau}{p^e}-a_{\tau'}-a'_{\tau'}\bigr)^+
+\bigl(\tfrac{l_\tau}{p^e}-l_{\tau'}\bigr)^+-\tfrac{a'_\tau}{p^e}\Bigr)-f\,n'^3\,(1+\log_pB),
\]
up to one further explicit term when $p=d$. The positive parts are the origin of the sum over
primes in the Short Forms Theorem. By the choice of $\varpi$, a derivative of total order $B$ of
$W(z/\varpi)$ at the origin has valuation at least $v(\pi_p)\sum_{e\in I_p}B/p^e$, up to a
logarithmic term. For an inert prime and $e=1$ the mechanism is visible in two variables
({\cite[Proposition~14.11]{P46}}): each factor $\Gamma_1$ taken from the exponential uses $p$
derivatives at a place $\tau$ and one at its partner, and costs one power of $p$, so that of the
amount $y/p$ which derivatives of order $y$ at $\tau$ contribute, at most $\min(y/p,y')$ is lost
if the order at $\tau'$ is $y'$, and $(y/p-y')^+$ remains. The couplings $\Gamma_e$ with $e>1$
account for the terms with $p^e$ in place of $p$. For $E=K$ and one
variable, divisibilities of this kind for the Taylor coefficients of elliptic theta functions at
supersingular primes are due to Katz~\cite{KatzDivisibilities}, Chellali~\cite{Chellali} and
Bannai and Kobayashi~\cite{BannaiKobayashi}. Here many places are coupled by the partner rule,
the constants have to be exact in the weights, and the bounds have to be uniform in the rank; the
passage from the entries of a Schur matrix to products of minors uses standard monomial
theory~\cite{DEP}.

Applied to the $m$ lifts of the upper chain this gives the divisibility of the composite lift
(Theorem~\ref{pa:divis}, which is {\cite[Proposition~14.20]{P46}}). Let $\ell$ be a prime and fix an
embedding of the algebraic numbers into $\C_\ell$. If $\ell\ne N$ and $\ell$ splits in $K$ or
$\ell\ge x_1$, every composite lift $B_{\vartheta,\Psi}$, for a twist $\vartheta$ that is
unramified above $\ell$, maps integral forms to integral forms. If $\ell=p<x_1$ does not split in $K$, then
with the exponents
\[
\augain_p=m\Bigl(\sum_{\tau\in\auPa}s_p(a,\tau)+k\sum_{\sigma\in\auPl}s_p(\kappa(\sigma),\sigma)
\Bigr),\qquad s_p(y,\sigma)=v(\pi_p)\sum_{e\in I_p}\Bigl(\frac y{p^e}-\kappa(\sigma')\Bigr)^+,
\]
where $\kappa(\sigma)$ is the weight at $\sigma$ and $\sigma'$ its partner at $p^e$, the image
$B_{\vartheta,\Psi}\varphi$ of an integral form $\varphi$ is divisible by $p^{\augain_p-\auerr_p}$,
in the sense that $p^{\auerr_p-\augain_p}B_{\vartheta,\Psi}\varphi$ is integral, with a loss
$\auerr_p=2f(\log b)^2a/p$, or $m$ times this if $p\equiv\pm1\pmod m$. This holds for all integral
$\varphi$ unless $p$ is \emph{exceptional}, and for an exceptional $p$ it holds for $\varphi$ in a
submodule of the integral forms whose colength divides $p^{\,i}$ with $i\le\tfrac12\auerr_p\dim
\auS$. If $\ell=N$, the image of an integral ramified form $\varphi$ has denominator at most
$N^{2mfb}$, again for $\varphi$ in a submodule, of colength dividing $N^{\,i}$ with $i\le mfb\dim
\auS$. By the identity~{\cite[(42)]{P46}}, the sum of $2\augain_p\log p$ over $p$ accounts
exactly for the sums over primes in the exponent of the Short Forms Theorem, and the losses,
multiplied by $\log\ell$, add up to less than the admitted error ({\cite[Lemma~14.18]{P46}}).

The places of $\auPa$ are those of class $1$ modulo a divisor $q_0$ of $q_1$, and a prime $p$ is
exceptional if it is inert in $K$, $p\not\equiv\pm1\pmod m$, and the order of $-p$ modulo $q_0$ is
odd and at most $\lceil\log_pm\rceil$. For such a prime the partner of a place of $\auPa$ at $p^e$
is again in $\auPa$ for some odd $e$ with $p^{e-1}<m$. At such a collision the count gives nothing
for the derivatives of order $a$, and the deficit, up to $2mfa/p$ for the $m$ lifts together,
exceeds the loss that can be afforded. We recover it on the image of the $m$-th
power of a product $Y_\Lambda$ of scaled Hecke operators $Y_{\mathfrak l}=p^{A_{\mathfrak l}}
T_{\mathfrak l}$ at primes $\mathfrak l$ of $E$ above $p$ (Lemma~\ref{pa:hecke}, proved in
\S\ref{pl:ss-hecke}); here $T_{\mathfrak l}$ is the Hecke operator of a diagonal matrix with
$\lceil\log_pm\rceil$ entries $p$, and $A_{\mathfrak l}$ is $a$ times the number of places of
$\auPa$ that induce $\mathfrak l$. The Hecke operators at $\mathfrak l$ on the two sides of a theta
lift with unramified data are related by an explicit identity
(Proposition~\ref{pl:onelift}), so that one factor $Y_\Lambda$ can be moved to each of the $m$
lifts, where it acts on the Schwartz function, and where a variant of the count
(Proposition~\ref{pl:charges}) gives the missing divisibility. That the image of $Y_\Lambda^{\,m}$, enlarged by a multiple of the lattice of
integral forms, has small colength is a statement on the elementary divisors of
$Y_{\mathfrak l}$ (Lemma~\ref{pa:eldiv}); it is the only input of the Short Forms Theorem not
proved in this paper, and is the subject of~\cite{P50}. At the prime $N$ nothing is gained, and
the task is to bound the denominators introduced by the level. Proposition~\ref{pl:direct}
rewrites the glued lifts at $N$, on ramified forms, as lifts whose data at $N$ are the
characteristic functions of the homomorphisms of lattices with a marked vector, applied to
$U_N^{\,m/2}\varphi$ for an explicit operator $U_N$ on the source forms, up to a constant of
absolute value at most $N^{mfb}$; the $m$ lifts together then lose at most $mfb$ further powers of
$N$, and the forms $\varphi$ for which $U_N^{\,m/2}\varphi$ is integral have the colength stated
above (Lemma~\ref{pl:level}). This error grows quadratically with $m$, because $\log N$ is itself
proportional to $m$; it is one of the two reasons for the choice $m\approx\sqrt x$
in~\cite{P46}. Finally, the integrality of the kernels at the primes dividing $2q_1M$ and at $N$,
to which the argument for the other primes does not apply, is
proved in Theorem~\ref{pf:std} and Corollary~\ref{pf:K1}; it uses that the model
$y^2+a_1xy=x^3+x$ of the curve defined by the chosen point of order two has good reduction at
every prime (Lemma~\ref{pf:two}).

\subsection*{The Weil representation}
The three theorems of~\cite{P46} use the Weil representation of unitary groups over the adeles of
$E^+$, split by characters as in Kudla~\cite{KudlaSplitting} and Harris, Kudla and
Sweet~\cite{HKS}, through a short list of properties (Lemma~\ref{ps:weil}). Some of them are not
in the literature in the form needed, and Section~\ref{pwf:sec} proves them from the definition
and from results of Weil~\cite{WeilActa} and of~\cite{KudlaSplitting}: the invariance of the theta
functional under the rational points for this splitting (Theorem~\ref{pwf:inv}), the mixed model
attached to an isotropic subspace that is not Lagrangian (Propositions~\ref{pwf:W2}
and~\ref{pwf:W3}), and the model at the places of $E^+$ that split in $E$
(Proposition~\ref{pwf:W5}). These results, and those of Section~\ref{A1:g2}
on discrete series, are also used in the other papers of the series.

The notation, the data and the statements of the propositions are those of~\cite{P46}. This
paper and~\cite{P50}, \cite{P51} quote one another, and we say precisely how. Lemma~\ref{pa:eldiv}
is stated in \S\ref{pa:ss-divis} and proved in~\cite{P50}, from a theorem on local rings proved there ({\cite[Theorem~2.2]{P50}}); we use it only for the exceptional primes, in part (iv) of Theorem~\ref{pa:divis}. For that proof \cite{P50} takes from us the normalisation of the weight modules
and of the Hecke operators, weak approximation, the existence of a full mark and the congruence between the places of $E$ under
the Frobenius (\S\ref{pl:ss-norm} with Remark~\ref{pl:N}, Proposition~\ref{pl:wa} and Corollary~\ref{pl:wacor}, Lemmas~\ref{pa:mark} and~\ref{pa:partner}), and for the theorem on local rings one lemma of linear algebra (Lemma~\ref{pa:elim}); none of
these depends on Lemma~\ref{pa:eldiv}. With~\cite{P51} the order is the following. That paper contains two elementary lemmas on the function
that represents the coefficient of the lowest $K$-type, an estimate for this function, and the convergence of an integral over a real group
({\cite[Lemmas~5.4 and~5.5]{P51}}, the proof of {\cite[Theorem~5.6]{P51}}, {\cite[Lemma~4.7(c)]{P51}}); none of these uses the integral
formula. With them we prove the decay of that function and then the formula (Lemma~\ref{pco:decay}, Theorem~\ref{pco:FJ}). From the formula, \cite{P51}
deduces the decay of the coefficient, the other properties of the kernel at a real place and the formal degree ({\cite[Theorem~5.6]{P51}}, parts (a) and (b) of {\cite[Lemma~4.7]{P51}}, {\cite[Lemma~4.8]{P51}},
{\cite[Corollary~4.43]{P51}}). These enter our proofs of the three archimedean properties, of Proposition~\ref{pwf:haar} and
of Proposition~\ref{pz:types}, and the latter is used in~\cite{P51} for the finiteness of the
dimension of the target space. Thus each statement is proved from statements that precede it,
although the two papers alternate. From~\cite{P53} and~\cite{P54} we quote a lemma of reduction
theory, a lemma on the Fock space, the description of the lattices and of the data modulo an ideal
at the primes above $q_1$, and the existence of the characters of the lifts; their proofs do not
use the results of this paper in whose proofs they occur. All statements quoted
from the other papers of the series are reproduced in Appendix~\ref{A1:quoted}. The main published
results we rely on have been named above; in addition we use the standard theory of elliptic
curves with complex multiplication and of their reduction~\cite{SilvermanAEC}, \cite{SilvermanAT},
\cite{LangEF}, and the theory of the Fock model for compact dual pairs~\cite{HoweTheta},
\cite{KashiwaraVergne}. A list of all statements on which the Short Forms Theorem depends, with
the place of each proof, is given in \S\ref{pz:ss-audit}.

The paper is organised as follows. Section~\ref{pwf:sec} contains the foundations of the Weil
representation. Section~\ref{A1:g2} treats the discrete series of $\auU(p,q)$: the three
archimedean properties (\S\ref{par:sec}), the formula for the coefficient of the lowest $K$-type
and the character of the Weil representation on the vacuum (\S\ref{pco:sec}).
Section~\ref{pa:sec} proves the count, the normal form of the theta products and the divisibility
of the composite lift. Section~\ref{A1:g4} is devoted to the upper bound: \S\ref{ps:sec} makes
the objects of the two chains precise and proves the equivariance, the restriction, the bound for
one kernel and the seesaw identity, and \S\ref{pk:sec} supplies the proofs postponed there, namely
the definite lifts, the unfolding, the integrality of the branching map and the properties of the
elliptic curve. Section~\ref{pl:sec} moves the Hecke operators through the lifts and treats the
prime $N$. Section~\ref{pf:sec} proves the reciprocity law, the integrality of the kernels at
every prime and a bound for numbers of classes. Section~\ref{pz:sec} defines the archimedean
types, proves that the lower chain is well defined, makes the constants explicit and lists the
inputs of the Short Forms Theorem.


\providecommand{\auA}{\mathbb A}
\providecommand{\auU}{\mathrm U}
\providecommand{\auGL}{\mathrm{GL}}
\providecommand{\auHom}{\operatorname{Hom}}
\providecommand{\auS}{\mathcal S}
\providecommand{\aubar}[1]{\overline{#1}}
\providecommand{\psHS}{\mathcal H}
\providecommand{\psrho}{\varrho}
\providecommand{\psZ}{Z}
\providecommand{\pwfA}{\mathsf A}
\providecommand{\pwfL}{\mathsf L}
\providecommand{\pwfM}{\mathsf M}
\providecommand{\pwfev}{\operatorname{ev}_0}

\section{Foundations of the Weil representation}\label{pwf:sec}

The properties of the Weil representation on which the Short Forms, Spanning and
Covolume Theorems rest are collected in Lemma~\ref{ps:weil}. For several of them the works cited there do not contain the statement: the invariance of
$\theta$ under the rational points for the splitting by characters (the cited works on the
splitting are local), the mixed model attached to a non-Lagrangian isotropic subspace, and the model
at places of $E^+$ split in $E$. This section proves these statements from the definition of
$\omega$ and from statements that are contained, with the locators given, in Weil~\cite{WeilActa}
and Kudla~\cite{KudlaSplitting}. It also replaces the appeal to Harish-Chandra's formula for the
formal degree in {\cite[Proposition~5.2]{P51}} by a result proved in the present series of papers, and proves the
archimedean normalisation {\cite[(30)]{P54}} by a computation on spaces of rank one.

\emph{Summary.} Theorem~\ref{pwf:inv} (the theta functional is invariant under $\auU(W)(E^+)$) is
proved in full from Definition~\ref{pwf:def} and the published statements (P1)--(P5) below.
Propositions~\ref{pwf:W2} and~\ref{pwf:W3} prove (W2), (W2$'$) and the identity of (W3) in full
(the bound of moderate growth in (W3) is Lemma~\ref{pz:absgrowth}(b), proved from (W0), (W1), (W2),
(W4), not from (W3)). Proposition~\ref{pwf:W5} proves (W5) and
\eqref{pl:eq-w5}, for the normalisation at split places fixed in Definition~\ref{pwf:def}(D3).
Proposition~\ref{pwf:haar} proves {\cite[Proposition~5.2]{P51}} for the parameters to which it is
applied. Lemma~\ref{pwf:delta-red} and Proposition~\ref{pwf:eplus} prove the archimedean
normalisation {\cite[(30)]{P54}}.

\subsection{Conventions, published inputs, and the definition of $\omega$}\label{pwf:ss-def}

\emph{Notation.} $k=E^+$; for a place $v$ of $k$ put $F=k_v$ and $E_v=E\otimes_kF$, which is a
quadratic field extension of $F$ or $F\times F$. All hermitian forms are conjugate-linear in the
first variable. For a non-degenerate hermitian space $W$ over $E$ (or $E_v$) put
\[
\pwfA(x,y)=\operatorname{Tr}_{E/k}\bigl(\jmath\langle y,x\rangle\bigr),\qquad
\beta(x,y)=\psi\bigl(\tfrac12\jmath\langle y,x\rangle\bigr)=\psi_k\bigl(\tfrac12\pwfA(x,y)\bigr),\qquad
[x,y]=\beta(x,y)\beta(y,x)^{-1}=\psi_k\bigl(\pwfA(x,y)\bigr),
\]
where $\psi=\psi_k\circ\operatorname{Tr}_{E/k}$ and $\psi_k=\psi_\Q\circ\operatorname{Tr}_{k/\Q}$
(\S\ref{ps:ss-weil}); the second equality holds because $\operatorname{Tr}(\jmath\langle x,y\rangle)
=\operatorname{Tr}(\aubar{\jmath\langle x,y\rangle})=-\operatorname{Tr}(\jmath\langle y,x\rangle)$.
Thus $\pwfA$ is a non-degenerate alternating $k$-bilinear form, $\beta$ is bimultiplicative, and
$\beta(x,y)=1$ whenever $\pwfA(x,y)=0$. A subspace $X$ is \emph{$\pwfA$-isotropic} if $\pwfA$
vanishes on $X\times X$; then $\beta=1$ on $X\times X$. Every $E$-subspace isotropic for
$\langle\,,\rangle$ is $\pwfA$-isotropic. $\pwfA$ is unchanged if $\langle\,,\rangle$ is replaced by
$\jmath\langle\,,\rangle$ and $\jmath$ by $1$; so $W$ with the skew-hermitian form
$\jmath\langle\,,\rangle$ is the space $\mathbb W=V\otimes_EW$ of~\cite{KudlaSplitting} for the
hermitian line $V=E$ (case $3_+$, $m=\dim V=1$), and $\auU(W)$ is the group $G=\auU(W)$ there.

A \emph{Heisenberg representation} of $W_v$ (resp.\ $W(\auA)$) is a pair $(\psHS,\psrho)$ of a
Hilbert space and a strongly continuous map $\psrho$ to unitary operators with
$\psrho(x)\psrho(y)=\beta(x,y)\psrho(x+y)$, irreducible; $\psHS^\infty$ denotes its smooth vectors.
An \emph{intertwiner} between two of them (possibly for different spaces, along an isometry
$\sigma$) is a unitary $T$ with $T\psrho(x)=\psrho'(\sigma x)T$.

\emph{Published inputs.} We use exactly the following.
\begin{itemize}
\item[(P1)] (Stone--von Neumann, and \cite[Th\'eor\`eme~1 and no.~34]{WeilActa}.) Up to
isomorphism there is one Heisenberg representation of $W_v$, and one of $W(\auA)$; an
intertwiner, if it exists, is unique up to a scalar of modulus $1$. For $g\in\mathrm{Sp}(W_v,\pwfA)$
there is a unitary $T$, unique up to a scalar, with $T\psrho(x)T^{-1}=\psrho(gx)$; it preserves
$\psHS^\infty$ and is continuous on it \cite[no.~35]{WeilActa}. A continuous linear operator on
$\psHS^\infty$ commuting with all $\psrho(x)$ is a scalar.
\item[(P2)] (\emph{Schr\"odinger model}.) Let $W=X\oplus Y$ with $X,Y$ $\pwfA$-isotropic. On
$\auS(X)$ (Schwartz--Bruhat functions) put $\psrho(x_0)\varphi(x)=\varphi(x+x_0)$ ($x_0\in X$),
$\psrho(y_0)\varphi(x)=[x,y_0]\varphi(x)$ ($y_0\in Y$) and $\psrho(x_0+y_0)=\beta(x_0,y_0)^{-1}
\psrho(x_0)\psrho(y_0)$. This is a Heisenberg representation, with $\psHS^\infty=\auS(X)$. (That
$\psrho(z)\psrho(z')=\beta(z,z')\psrho(z+z')$ is the computation $\psrho(y)\psrho(x')=[y,x']
\psrho(x')\psrho(y)$, $\beta(x+x',y+y')\beta(x,y)^{-1}\beta(x',y')^{-1}[y,x']=\beta(x,y')\beta(y,x')
=\beta(z,z')$, using $\beta=1$ on $X\times X$ and $Y\times Y$; irreducibility and the smooth
vectors are the classical facts of the Schr\"odinger model.) The same holds over $\auA$ with
$\auS(X(\auA))$.
\item[(P3)] (\emph{Kudla's splitting for split spaces}, \cite[Theorem~3.1, p.~378; (1.11)--(1.13),
p.~370; (1.14), p.~371; \S5, pp.~399--400]{KudlaSplitting}.) Let $F\ne\C$ be a local field of characteristic $0$, $E_v$ a field,
$W$ split, $Y\subset W$ a Lagrangian $E_v$-subspace and $X$ a complementary Lagrangian. Fix a
character $\xi$ of $E_v^\times$ with $\xi|_{F^\times}=\epsilon_{E_v/F}$. Kudla defines unitary
operators $\omega_{Y,\xi}(g)=\beta_{Y}(g)\,r_Y(g)$ on $\auS(X)$, where $r_Y(\sigma)\varphi(x)=
\int_{Y_\sigma}f_\sigma(x+y)\varphi(x\alpha+y\gamma)\,d\mu_\sigma(y)$ is Rao's operator for
$\sigma=\bigl(\begin{smallmatrix}\alpha&\beta\\\gamma&\delta\end{smallmatrix}\bigr)$ (with a
positive Haar measure $\mu_\sigma$ on $Y_\sigma=Y/\ker\gamma$, normalised so that $r_Y(\sigma)$ is
unitary), and $\beta_Y(g)=\xi(x(g))\gamma_F(\eta\circ RV)^{-j(g)}$ (Kudla's $\beta_V(g)$, case $3_+$,
with $\eta(t)=\psi_K(\frac12t)$ for Kudla's additive character $\psi_K$ of $F$); here $g\in P_Y\tau_jP_Y$ (Bruhat
decomposition (1.11)) and, for $g=p_1\tau_Sp_2$, $x(g)$ is the reduced norm, here the determinant, of
$p_1p_2|_Y$ (1.13)--(1.14). The map
$g\mapsto\omega_{Y,\xi}(g)$ is a representation of $\auU(W)$ on $\auS(X)$ \cite[p.~400]{KudlaSplitting}.
Kudla's groups act on the right (\cite[p.~362]{KudlaSplitting}: rows times matrices), and Rao's
operators normalise the Heisenberg operators $\psrho(x_0)\varphi(x)=\varphi(x+x_0)$,
$\psrho(y_0)\varphi(x)=\psi_K(\langle\!\langle x,y_0\rangle\!\rangle)\varphi(x)$ by
$\omega_{Y,\xi}(g)\psrho(w)\omega_{Y,\xi}(g)^{-1}=\psrho(w\cdot g^{-1})$ (for the generators of the
Siegel parabolic this is the direct computation from the formula for $r_Y$; with the opposite sign in
$\psrho(y_0)$ the unipotent and Levi elements would conjugate in opposite directions). For $W$ with
the form $\jmath\langle\,,\rangle$ and the hermitian line $V=E$, $(v_1,v_2)=\aubar v_1v_2$, Kudla's
symplectic form is $\langle\!\langle\,,\rangle\!\rangle=\frac12\pwfA$ \cite[(0.8), (0.9), p.~364]{KudlaSplitting},
so these Heisenberg operators are those of (P2) exactly when $\psi_K(t)=\psi_k(2t)$; this is the
character with which (P3) is read.
\item[(P4)] (\emph{Kudla's splitting for non-split spaces}, \cite[(4.1)--(4.5), p.~389, and
Proposition~4.1, p.~390]{KudlaSplitting}.) For $W$ not split, $W\oplus W^-$ is split ($W^-$ is $W$ with
the form negated), $\auU(W)$ embeds by $g\mapsto(g,1)$, and the splitting of $\auU(W)$ is the
restriction of the splitting (P3) of $\auU(W\oplus W^-)$, read in the decomposition
$\mathrm{Mp}(W)\times\mathrm{Mp}(W^-)\to\mathrm{Mp}(W\oplus W^-)$ attached to a polarisation
$Y'=Y'_1+Y'_2$ with $Y'_1\subset W$, $Y'_2\subset W^-$, along which the pull-back of Rao's cocycle
of $W\oplus W^-$ is the product of those of $W$ and $W^-$ \cite[p.~389]{KudlaSplitting}.
\item[(P5)] (\emph{Characters and Fourier analysis}.) $\psi$ is trivial on $E$ and non-trivial on
$\auA_E$; $\upsilon$ and $|\cdot|_{\auA_E}$ are trivial on $E^\times$; for a $k$-vector space $X$
the pairing $\psi_k(\pwfA(\cdot,\cdot))$ identifies $Y(k)$ with the character group of
$X(\auA)/X(k)$ when $X\oplus Y$ is a Lagrangian decomposition (self-duality of $\auA_k/k$); Fourier
series of smooth functions on a real torus converge in the $C^\infty$ topology; a distribution $u$
on $\R^N$ with $x_iu=0$ for all $i$ is a multiple of $\delta_0$.
\end{itemize}

\emph{Left actions.} Kudla's groups act on the right (P3); the paper's act on the left. We pass from
one to the other by letting $g$ act on the left by $g\cdot w:=w\cdot g^{-1}$, which is the unique
way to keep the group law and the identity $\omega(g)\psrho(w)\omega(g)^{-1}=\psrho(g\cdot w)$ of
(W1). Under this dictionary, for $p$ in the stabiliser $P_Y$ of $Y$, the endomorphism $p_1p_2|_Y$
of (1.13) with $p_1=p$, $p_2=1$ is the inverse of the left action of $p$ on $Y$. Hence, in the
paper's notation,
\begin{equation}\label{pwf:eq-xP}
\beta_Y(p)=\xi\bigl(\det(p|_Y)\bigr)^{-1}\qquad(p\in P_Y),
\end{equation}
since $j(p)=0$. (The value of $\xi$ on $\det(p|_Y)$ only matters modulo norms, and $\xi(\aubar
z^{-1})=\xi(z)\xi(z\aubar z)^{-1}=\xi(z)$; so the same value is obtained if $\det(p|_Y)$ is replaced
by the determinant of $p$ on a complement of $Y$ dual to $Y$.)

\begin{definition}[The Weil representation]\label{pwf:def}
Let $W$ be a non-degenerate hermitian space over $E$. For each place $v$ of $k$ a homomorphism
$\omega_v=\omega_{W,v}$ from $\auU(W)(k_v)$ to the unitary operators of a Heisenberg representation
$\psHS_v(W)$, with $\omega_v(g)\psrho(w)\omega_v(g)^{-1}=\psrho(gw)$, is fixed as follows.
\begin{itemize}
\item[(D1)] $E_v$ a field and $W_v$ split: $\omega_v=\omega_{Y,\upsilon_v}$ of (P3) for a Lagrangian
$Y\subset W_v$ (by Lemma~\ref{pwf:indep} it does not depend on $Y$).
\item[(D2)] $E_v$ a field and $W_v$ not split: the splitting (P4) with $\xi=\upsilon_v$, characterised
intrinsically as follows. The space $W_v\oplus W_v^-$ is split, and $\omega_{W\oplus W^-,v}(g,1)$ is
defined by (D1), independently of the Lagrangian (Lemma~\ref{pwf:indep}). Since $(g,1)$ fixes $W^-$
pointwise, this operator commutes with $1\otimes\psrho(w')$ for $w'\in W^-_v$ on
$\psHS_v(W\oplus W^-)=\psHS_v(W)\widehat\otimes\psHS_v(W^-)$; as $\psrho|_{W^-}$ is irreducible (P1),
the commutant of $1\otimes\psrho(W^-_v)$ is the set of operators $A\otimes1$, so
$\omega_{W\oplus W^-,v}(g,1)=A_g\otimes1$ for a unique bounded $A_g$, and we set $\omega_v(g):=A_g$.
As $A\mapsto A\otimes1$ is an injective, multiplicative, $*$-preserving map, $g\mapsto A_g$ is a
unitary representation, and $A_g\psrho(w)A_g^{-1}=\psrho(gw)$ follows from the same identity for
$(g,1)$ on $w\in W_v$. This is the operator of (P4) \cite[(4.1)--(4.5), Proposition~4.1]{KudlaSplitting},
and it involves no auxiliary choice.
\item[(D3)] $v=w\aubar w$ split in $E$: $W_v=W_w\oplus W_{\aubar w}$ with both summands
$\pwfA$-isotropic, $\auU(W)(k_v)=\auGL(W_w)$ (restriction to $W_w$), and on the model (P2)
$\auS(W_w)$,
\[
(\omega_v(g)\Psi)(x)=\upsilon_w(\det g_w)^{-1}|\det g_w|_w^{-1/2}\,\Psi(g_w^{-1}x).
\]
\end{itemize}
The global representation is $\omega=\bigotimes_v\omega_v$ on $\psHS(W)=\bigotimes'_v\psHS_v(W)$, as
in \S\ref{ps:ss-weil} and (W4).
\end{definition}

In (D3) the formula defines a representation with the required property: $g_w$ preserves $W_w$
and $W_{\aubar w}$ (as $g$ is $E_v$-linear), $\Psi\mapsto|\det g_w|^{-1/2}\Psi(g_w^{-1}\cdot)$ is
unitary, conjugates translation by $x_0$ into translation by $g_wx_0$, and conjugates multiplication
by $[\cdot,y_0]$ into multiplication by $[g_w^{-1}\cdot,y_0]=[\cdot,gy_0]$ ($g$ is an isometry);
the factor $\upsilon_w(\det g_w)^{-1}$ is a character. The cited works treat only the case of a
field $E_v$; (D3) is the normalisation used in the paper at split places (it is the one for which
\eqref{pl:eq-w5} holds, Proposition~\ref{pwf:W5}, and the one compatible with (D1), (D2) on the
rational points, Lemma~\ref{pwf:siegel}). In (D1), (D2) the identity $\omega_v(g)\psrho(w)
\omega_v(g)^{-1}=\psrho(gw)$ is (P3) read with left actions.

\begin{lemma}[Naturality; independence of the Lagrangian]\label{pwf:indep}
In case (D1), let $\sigma\colon W\to W'$ be an isometry of split spaces and $T\colon\psHS_v(W)\to
\psHS_v(W')$ an intertwiner along $\sigma$. Then $T\omega_{Y,\xi}(g)T^{-1}=\omega_{\sigma Y,\xi}
(\sigma g\sigma^{-1})$. Moreover $\omega_{Y,\xi}=\omega_{Y',\xi}$ for any two Lagrangians $Y,Y'$.
\end{lemma}
\begin{proof}
All ingredients of (P3) ($\psi$, $\xi$, the Bruhat cells and the invariants $x$, $j$, Rao's
operators and their measures, the Weil index of the form $RV$) are defined
(here $x(g)$ is the reduced norm of the $E_v$-linear endomorphism $p_1p_2|_Y$, which $\sigma$
preserves; the Weil index $\gamma_F(\eta\circ RV)$ depends only on $V$, the same on both sides; and
Rao's measures are fixed by unitarity of his operators, so they are carried over as well) from $(W,Y)$ by
formulas that are carried to those for $(W',\sigma Y)$ by $\sigma$ and by the transport of functions
on $X$ to functions on $\sigma X$; the transport of functions is an intertwiner along $\sigma$, and
by (P1) any other intertwiner $T$ differs from it by a scalar, which cancels in $T(\cdot)T^{-1}$.
This is the first claim. For the second, $Y'=\alpha Y$ for some $\alpha\in\auU(W)$ (Witt's theorem).
Apply the first claim with $W'=W$, $\sigma=\alpha$ and $T=\omega_{Y,\xi}(\alpha)$, which is an
intertwiner along $\alpha$: $\omega_{Y',\xi}(\alpha g\alpha^{-1})=\omega_{Y,\xi}(\alpha)\omega_{Y,\xi}(g)
\omega_{Y,\xi}(\alpha)^{-1}=\omega_{Y,\xi}(\alpha g\alpha^{-1})$ for all $g$.
\end{proof}

\subsection{Local lemmas: invariant functionals, the Siegel parabolic, orthogonal sums}\label{pwf:ss-local}

\begin{lemma}[Distributions supported at $0$]\label{pwf:dist}
Let $v$ be a place, $W_v=X\oplus Y\oplus W'$ with $X,Y$ $\pwfA$-isotropic, $X\oplus Y$ orthogonal
to $W'$ for $\pwfA$, and $(\psHS',\psrho')$ a Heisenberg representation of $W'$ (or $\psHS'=\C$ if
$W'=0$). Realise $\psHS_v(W)$ as $\auS(X)\widehat\otimes\psHS'$, with smooth vectors
$\auS(X;\psHS'^\infty)$, where $X\oplus Y$ acts by \textup{(P2)} on the first factor and $W'$ by
$\psrho'$ on the second. If $T\colon\auS(X;\psHS'^\infty)\to\psHS'^\infty$ is continuous linear,
$T\circ\psrho(y)=T$ for $y\in Y$ and $T\circ\psrho(w')=\psrho'(w')\circ T$ for $w'\in W'$, then
$T=c\,\pwfev$ for a constant $c$, where $\pwfev F=F(0)$.
\end{lemma}
\begin{proof}
(This is a realisation of $\psHS_v(W)$: $\beta(z,w')=1$ for $z\in X\oplus Y$, $w'\in W'$, so the
operators commute and $\psrho(z+w')=\psrho(z)\psrho(w')$; irreducibility is (P1).) Let $F\in
\auS(X;\psHS'^\infty)$ with $F(0)=0$; we show $TF=0$. If $v$ is finite, $F$ is locally constant
with compact support, so $F=F\cdot(1-\mathbf 1_U)$ for an open compact subgroup $U\subset X$; let
$M\subset Y$ be the open compact subgroup $\{y:[U',y]=1\}$ for $U'\subset U$ small enough, with
$\int_M[x,y]\,dy=\mathbf 1_{M^\perp}(x)$ ($dy$ of mass $1$) and $M^\perp\subset U$. Then
$F=F-\mathbf 1_{M^\perp}F=\int_M(F-\psrho(y)F)\,dy$, a finite average (the integrand is locally
constant in $y$), so $TF=\int_M(TF-T\psrho(y)F)dy=0$. If $v$ is real or complex, differentiate
$T(\psrho(ty)F)=T(F)$ at $t=0$: $T(\pwfA(x,y)F)=0$ for all $y$, i.e.\ $T(x_iF)=0$ for linear
coordinates $x_i$ on $X$. For a continuous functional $\lambda$ on $\psHS'^\infty$ and $h\in
\psHS'^\infty$ the distribution $u(\varphi)=\lambda(T(\varphi h))$ on $X$ satisfies $x_iu=0$, so
$u=c(\lambda,h)\delta_0$ (P5); hence $T(\varphi h)=\varphi(0)T(\phi_0h)$ for any $\phi_0$ with
$\phi_0(0)=1$, and by density of $\auS(X)\otimes\psHS'^\infty$ and continuity, $TF=T_0(F(0))$ with
$T_0(h)=T(\phi_0h)$. In both cases $T=T_0\circ\pwfev$ with $T_0$ continuous on $\psHS'^\infty$, and
$T_0\psrho'(w')=\psrho'(w')T_0$; by (P1), $T_0$ is a scalar.
\end{proof}

In the rest of this paragraph $W$ is split over $E_v$, $W=I\oplus I^*$ with $I,I^*$ Lagrangian
$E_v$-subspaces, and $\psHS_v(W)=\auS(I^*)$ by (P2) ($X=I^*$, $Y=I$). For $a\in\auGL(I)$ let $m(a)$
act by $a$ on $I$ and by the inverse adjoint $(a^*)^{-1}$ on $I^*$ (adjoint for the perfect pairing
$\langle\,,\rangle\colon I\times I^*\to E_v$), and for $b\in\auHom_{E_v}(I^*,I)$ with
$\langle bx,y\rangle+\langle x,by\rangle=0$ let $n(b)$ be the identity on $I$ and $x\mapsto x+bx$ on
$I^*$. These are elements of $\auU(W)$, $n(b)n(b')=n(b+b')$, and $m(a)n(b)m(a)^{-1}=n(m(a)\,b\,
m(a)^{-1})$. Put
\[
(\pwfL_a\varphi)(x)=|\det a|_{E_v}^{1/2}\varphi(m(a)^{-1}x),\qquad(\pwfM_b\varphi)(x)=\beta(x,bx)^{-1}
\varphi(x),
\]
where $\det a=\det_{E_v}(a|_I)$ and $|\cdot|_{E_v}$ is the module of $E_v$.

\begin{lemma}[The Siegel parabolic]\label{pwf:siegel}
At every place $v$, $\omega_v(m(a))=\upsilon_v(\det a)^{-1}\pwfL_a$ and $\omega_v(n(b))=\pwfM_b$.
\end{lemma}
\begin{proof}
\emph{The operators.} $m(a)$ acts on $I^*$ with module $|\det a|_{E_v}^{-1}$ (its
$E_v$-determinant in dual bases is $\aubar{\det a}^{-1}$), so $\pwfL_a$ is unitary; it conjugates
translation by $x_0$ into translation by $m(a)x_0$ and multiplication by $[\cdot,y_0]$ into
multiplication by $[m(a)^{-1}\cdot,y_0]=[\cdot,m(a)y_0]$, so it is an intertwiner along $m(a)$. For
$\pwfM_b$: the condition on $b$ gives $\pwfA(x,by)=\pwfA(y,bx)$, hence $\beta(x,by)=\beta(y,bx)$; with
$q(x)=\beta(x,bx)$ one gets $q(x+x_0)=q(x)q(x_0)[x,bx_0]$ and $[x_0,bx_0]=q(x_0)^2$. Therefore
$\pwfM_b\psrho(x_0)\pwfM_b^{-1}\varphi(x)=q(x_0)[x,bx_0]\varphi(x+x_0)=\beta(x_0,bx_0)^{-1}[x+x_0,bx_0]
\varphi(x+x_0)=\psrho(x_0+bx_0)\varphi(x)$, while $\pwfM_b$ commutes with $\psrho(y_0)$, $y_0\in I$;
as $n(b)x_0=x_0+bx_0$ and $n(b)y_0=y_0$, $\pwfM_b$ is an intertwiner along $n(b)$. Moreover
$\pwfM_b\pwfM_{b'}=\pwfM_{b+b'}$ and $\pwfL_a\pwfM_b\pwfL_a^{-1}=\pwfM_{m(a)bm(a)^{-1}}$ (as
$\beta(m^{-1}x,bm^{-1}x)=\beta(x,mbm^{-1}x)$).

\emph{The unipotent elements.} By (P1), $\omega_v(n(b))=c(b)\pwfM_b$ and $\omega_v(m(a))=\kappa(a)
\pwfL_a$ with scalars of modulus $1$. Since $\omega_v$ is a homomorphism, $c$ is a character of the
vector group of the $b$'s, and conjugating by $\omega_v(m(a))$ gives $c(m(a)bm(a)^{-1})=c(b)$. For
$a=2\cdot\mathrm{id}_I$, $m(a)bm(a)^{-1}=4b$; so $c(3b)=1$ for all $b$, and $c=1$.

\emph{The Levi elements.} $\kappa$ is a character and it remains to find it. Case (D1): $W_v$ is
split, $\omega_v=\omega_{I,\upsilon_v}$ (Lemma~\ref{pwf:indep}). For $\sigma=m(a)$ the lower left
block $\gamma$ vanishes, $Y_\sigma=0$, $f_\sigma=1$ (the upper right block is $0$), so $r_I(m(a))$ is a
positive multiple of $\varphi\mapsto\varphi(m(a)^{-1}\cdot)$ (read with left actions), i.e.\ of
$\pwfL_a$; and $\beta_I(m(a))=\upsilon_v(\det a)^{-1}$ by \eqref{pwf:eq-xP}. So $\kappa(a)=
\upsilon_v(\det a)^{-1}$. Case (D2) does not occur ($W_v$ is split). Case (D3): $v=w\aubar w$. The
idempotents $e_w,e_{\aubar w}$ of $E_v$ satisfy $\aubar{e_w}=e_{\aubar w}$, so
$\langle e_wx,e_wy\rangle=e_{\aubar w}e_w\langle x,y\rangle=0$: $\langle\,,\rangle$ pairs $W_w$ with
$W_{\aubar w}$, and $I^*_w$ with $I_{\aubar w}$. Let $\Phi\colon\auS(W_w)\to\auS(I^*)$ be an
intertwiner (P1) between the models (P2) for $W_w\oplus W_{\aubar w}$ and $I^*\oplus I$, and
$\pwfL'_a\Psi=|\det(m(a)|_{W_w})|^{-1/2}\Psi(m(a)^{-1}\cdot)$. Both $\Phi\pwfL'_a\Phi^{-1}$ and
$\pwfL_a$ are intertwiners along $m(a)$, so $\Phi\pwfL'_a\Phi^{-1}=\kappa'(a)\pwfL_a$ with
$|\kappa'(a)|=1$. To see $\kappa'(a)=1$, consider the functionals $\lambda_1(\Psi)=\int_{I_w}\Psi$ on
$\auS(W_w)$ and $\lambda_2=\pwfev$ on $\auS(I^*)$. Both are invariant under $\psrho(z)$, $z\in I$:
for $z\in I_w\subset W_w$, $\psrho(z)$ is a translation along $I_w$; for $z\in I_{\aubar w}$ it is
multiplication by $[x,z]$, which is $1$ on $I_w$ since $I$ is isotropic; for $\lambda_2$ by (P2). By
Lemma~\ref{pwf:dist} (with $W'=0$, in the model $\auS(I^*)$), $\lambda_1\circ\Phi^{-1}=c\lambda_2$,
and $c\ne0$. Now $\lambda_1\circ\pwfL'_a$ and $\lambda_2\circ\pwfL_a$ are positive multiples of
$\lambda_1$, $\lambda_2$ (change of variables), and $\lambda_1\circ\pwfL'_a=c\lambda_2\circ\Phi\pwfL'_a=
c\,\kappa'(a)\lambda_2\circ\pwfL_a\circ\Phi$; evaluating on $\Psi$ with $\lambda_1(\Psi)\ne0$ gives
$\kappa'(a)>0$, so $\kappa'(a)=1$. By (D3), $\omega_v(m(a))$ is $\upsilon_w(\det(m(a)|_{W_w}))^{-1}
\pwfL'_a$ in the first model. On $W_w=I_w\oplus I^*_w$, $m(a)$ acts by $a_w$ and by the inverse
transpose of $a_{\aubar w}$, so $\det(m(a)|_{W_w})=\det a_w/\det a_{\aubar w}$. As $\upsilon$ is
trivial on $k_v^\times$ ($\epsilon_{E/k}$ is trivial at split places), $\upsilon_{\aubar w}=
\upsilon_w^{-1}$ on $k_v^\times$, and $\upsilon_w(\det a_w/\det a_{\aubar w})^{-1}=\upsilon_w(\det
a_w)^{-1}\upsilon_{\aubar w}(\det a_{\aubar w})^{-1}=\upsilon_v(\det a)^{-1}$. Transporting by
$\Phi$: $\kappa(a)=\upsilon_v(\det a)^{-1}$.
\end{proof}

\begin{lemma}[Orthogonal sums and doubling]\label{pwf:sum}
Let $W=W_1\oplus W_2$ be an orthogonal sum of non-degenerate hermitian spaces over $E_v$, and
identify $\psHS_v(W)=\psHS_v(W_1)\widehat\otimes\psHS_v(W_2)$ by an intertwiner (P1). Then
$\omega_{W,v}(g_1,g_2)=\omega_{W_1,v}(g_1)\otimes\omega_{W_2,v}(g_2)$ for $g_i\in\auU(W_i)(k_v)$. In
particular $\omega_{W\oplus W^-,v}(g,1)=\omega_{W,v}(g)\otimes1$.
\end{lemma}
\begin{proof}
The identity does not depend on the intertwiner chosen (P1). \emph{Case (D3).} $W_w=W_{1,w}\oplus
W_{2,w}$, $\auS(W_w)=\auS(W_{1,w})\widehat\otimes\auS(W_{2,w})$ is an intertwiner,
$\det((g_1,g_2)_w)=\det g_{1,w}\det g_{2,w}$, and the formula of (D3) factors.

\emph{Case $E_v$ a field, $W_1,W_2$ split.} Choose Lagrangians $Y_i\subset W_i$ with complements
$X_i$; then $Y=Y_1\oplus Y_2$ is a Lagrangian of $W$, and by Lemma~\ref{pwf:indep} we may compute
$\omega_{W,v}$ with $Y$ on $\auS(X_1\oplus X_2)=\auS(X_1)\widehat\otimes\auS(X_2)$, which is an
intertwiner. If $g_i=p_i\tau_{S_i}p'_i$ are Bruhat decompositions in $\auU(W_i)$ (for bases adapted
to $W_i=X_i\oplus Y_i$), then $(g_1,g_2)=(p_1,p_2)\tau_{S_1\sqcup S_2}(p'_1,p'_2)$ is one in
$\auU(W)$; so $j(g_1,g_2)=j(g_1)+j(g_2)$ and $x(g_1,g_2)=x(g_1)x(g_2)$ by (1.14), and, the space $V$
and the Weil index being the same, $\beta_Y(g_1,g_2)=\beta_{Y_1}(g_1)\beta_{Y_2}(g_2)$. For
$\sigma=(\sigma_1,\sigma_2)$ block diagonal, $Y_\sigma=Y_{1,\sigma_1}\times Y_{2,\sigma_2}$ and
$f_\sigma=f_{\sigma_1}\otimes f_{\sigma_2}$, so $r_Y(\sigma)$ is a positive multiple of
$r_{Y_1}(\sigma_1)\otimes r_{Y_2}(\sigma_2)$ (the measures being positive); both are unitary, so the
multiple is $1$. Hence the identity.

\emph{Doubling, $E_v$ a field.} If $W$ is split, so is $W\oplus W^-$, and the case just
treated gives $\omega_{W\oplus W^-,v}(g,1)=\omega_{W,v}(g)\otimes1$ when $W$ and $W^-$ are split.
If $W$ is not split, this identity holds by the definition (D2): $\omega_{W\oplus W^-,v}(g,1)$,
which does not depend on the Lagrangian of the split space $W\oplus W^-$ (Lemma~\ref{pwf:indep}),
commutes with $1\otimes\psrho(W^-_v)$ and hence equals $A_g\otimes1$ with $A_g=\omega_{W,v}(g)$.
So in all cases
\begin{equation}\label{pwf:eq-double}
\omega_{W\oplus W^-,v}(g,1)=\omega_{W,v}(g)\otimes1\qquad(g\in\auU(W)(k_v)).
\end{equation}

\emph{General case, $E_v$ a field.} Let $D_i=W_i\oplus W_i^-$ (split), and $\sigma\colon
W\oplus W^-\to D_1\oplus D_2$ the isometry $(w_1,w_2,w'_1,w'_2)\mapsto(w_1,w'_1,w_2,w'_2)$; it
carries $((g_1,g_2),1)$ to $((g_1,1),(g_2,1))$. By \eqref{pwf:eq-double} for $W$, by
Lemma~\ref{pwf:indep} for $\sigma$, by the split case for $D_1\oplus D_2$, and by
\eqref{pwf:eq-double} for $W_1$, $W_2$:
\[
\omega_{W,v}(g_1,g_2)\otimes1=\omega_{W\oplus W^-,v}((g_1,g_2),1)\ \hat=\ \omega_{D_1,v}(g_1,1)\otimes
\omega_{D_2,v}(g_2,1)=\bigl(\omega_{W_1,v}(g_1)\otimes1\bigr)\otimes\bigl(\omega_{W_2,v}(g_2)\otimes1
\bigr),
\]
where $\hat=$ means equality after the rearrangement of tensor factors induced by $\sigma$ (an
intertwiner along $\sigma$). Comparing the factors on $\psHS_v(W)$ gives the claim; the last
assertion is \eqref{pwf:eq-double}.
\end{proof}

\subsection{Uniqueness of the theta functional}\label{pwf:ss-unique}

Globally, $\psHS(W)^\infty$ denotes the smooth vectors (at the infinite places for the action of
the Heisenberg group, at the finite places the algebraic restricted tensor product), with its
natural topology; in a Schr\"odinger model (P2) over $\auA$ it is $\auS(X(\auA))$. A \emph{theta
functional} is a non-zero continuous linear functional $\theta$ on $\psHS(W)^\infty$ with
$\theta\circ\psrho(w)=\theta$ for $w\in W(k)$; this is the functional of (W0), (W1).

\begin{lemma}[Uniqueness]\label{pwf:unique}
Let $W(k)=X(k)\oplus Y(k)$ with $X,Y$ $\pwfA$-isotropic $k$-subspaces (such a decomposition always
exists, $\pwfA$ being alternating). In the model $\auS(X(\auA))$ every continuous linear functional
$T$ with $T\circ\psrho(w)=T$ for $w\in W(k)$ is $T(\varphi)=c\sum_{\xi\in X(k)}\varphi(\xi)$. In
particular a theta functional exists and is unique up to a scalar, in every model.
\end{lemma}
\begin{proof}
The functional $\vartheta(\varphi)=\sum_{\xi\in X(k)}\varphi(\xi)$ is continuous and invariant:
translations by $X(k)$ permute the terms; multiplication by $[\cdot,y]$, $y\in Y(k)$, does not
change them, since $[\xi,y]=\psi_k(\pwfA(\xi,y))=1$ (P5); and $\beta(x_0,y_0)=1$ for rational $x_0,
y_0$. Conversely let $T$ be invariant. Let $C=X(\auA)/X(k)$, compact. Fix an open compact subgroup
$K_f\subset X(\auA_f)$ and $\Lambda=X(k)\cap K_f$, a lattice in $X_\infty$, with $X(\auA)=X(k)+(X_\infty
\times K_f)$; choose $u_\infty\in C_c^\infty(X_\infty)$ with $\sum_{\lambda\in\Lambda}u_\infty(x+
\lambda)=1$ and put $u=u_\infty\otimes\mathbf 1_{K_f}$. Then $\sum_{\xi\in X(k)}u(x+\xi)=1$ for all
$x\in X(\auA)$.

(i) \emph{Multiplication.} For a $K_f$-invariant smooth function $f$ on $C$ (a smooth function on the
real torus $X_\infty/\Lambda$, the $K_f$-invariant part of $C$) let $\tilde f$ be its lift to
$X(\auA)$. The characters of $C$ are the $[\cdot,y]$, $y\in Y(k)$ (P5), and $T([\cdot,y]\varphi)=
T(\varphi)$. Hence $T(\tilde f\varphi)=f(0)T(\varphi)$ for trigonometric polynomials $f$, and by
$C^\infty$ convergence of Fourier series (P5) and continuity of $f\mapsto\tilde f\varphi$ into
$\auS(X(\auA))$, for all such $f$.

(ii) \emph{Periodisation.} For $\varphi\in\auS(X(\auA))$, $\varphi=\sum_{\xi}\varphi\cdot u(\cdot-\xi)$
with convergence in $\auS(X(\auA))$ (only finitely many classes of $\xi$ modulo $\Lambda$ meet the
compact support of the finite part, and the sum over $\Lambda$ converges in $\auS(X_\infty)$ because
$u_\infty$ has compact support). By invariance under translation by $\xi$,
$T(\varphi\cdot u(\cdot-\xi))=T(\varphi(\cdot+\xi)u)$, so $T(\varphi)=T\bigl((P\varphi)u\bigr)$ with
$P\varphi(x)=\sum_\xi\varphi(x+\xi)$, a smooth function on $C$, $K_f'$-invariant for some open
$K'_f$ (apply (i) with $K'_f\subset K_f$). By (i), $T(\varphi)=P\varphi(0)\,T(u)=T(u)\,\vartheta(\varphi)$.
The last assertion follows because any two Heisenberg representations are intertwined (P1), and an
intertwiner is a topological isomorphism of the smooth vectors.
\end{proof}

\subsection{The theta functional is invariant under the rational points}\label{pwf:ss-M4}

\begin{lemma}[Generation]\label{pwf:gen}
Let $W=I\oplus I^*$ be split over $E$ with $I,I^*$ Lagrangian. Then $\auU(W)(k)$ is generated by the
$m(a)$, $n(b)$ of \S\ref{pwf:ss-local} and the elements $n^-(c)$ that are the identity on $I^*$ and
map $y\in I$ to $y+cy$, $c\in\auHom(I,I^*)$ with $\langle cy,y'\rangle+\langle y,cy'\rangle=0$. With
$\varepsilon\in E^\times$, $\aubar\varepsilon=-\varepsilon$, the element $w_\varepsilon$ that maps
$y\in I$ to $-\varepsilon^{-1}\iota(y)$ and $x\in I^*$ to $\varepsilon\iota^{-1}(x)$, where
$\iota\colon I\to I^*$ is the isomorphism attached to dual bases, lies in $\auU(W)(k)$, satisfies
$w_\varepsilon=n(b)n^-(c)n(b)$ for suitable $b,c$, and $w_\varepsilon n(b')w_\varepsilon^{-1}$ runs
over all $n^-(c)$ as $b'$ runs over the admissible $b'$.
\end{lemma}
\begin{proof}
Choose dual bases of $I$, $I^*$; the Gram matrix is $J=\bigl(\begin{smallmatrix}0&1\\1&0
\end{smallmatrix}\bigr)$ and $g=\bigl(\begin{smallmatrix}a&b\\c&d\end{smallmatrix}\bigr)\in\auU(W)$
iff $g^*Jg=J$, equivalently $gJg^*=J$; in particular $cd^*+dc^*=0$. Then $m(a)=\mathrm{diag}(a,
a^{*-1})$, $n(b)=\bigl(\begin{smallmatrix}1&b\\0&1\end{smallmatrix}\bigr)$ and $n^-(c)=
\bigl(\begin{smallmatrix}1&0\\c&1\end{smallmatrix}\bigr)$ with $b,c$ skew-hermitian. The matrix
identities
\[
\begin{pmatrix}1&\varepsilon\\0&1\end{pmatrix}\begin{pmatrix}1&0\\-\varepsilon^{-1}&1\end{pmatrix}
\begin{pmatrix}1&\varepsilon\\0&1\end{pmatrix}=\begin{pmatrix}0&\varepsilon\\-\varepsilon^{-1}&0
\end{pmatrix}=w_\varepsilon,\qquad w_\varepsilon n(b')w_\varepsilon^{-1}=n^-(-\varepsilon^{-2}b')
\]
(scalars times the identity of the blocks; $\pm\varepsilon^{\pm1}$ are skew-hermitian and
$\varepsilon^2\in k$) prove the assertions on $w_\varepsilon$.

Let $g\in\auU(W)(k)$. \emph{Step 1: there is a skew-hermitian $y$ with $d+cy$ invertible.} Replacing
$g$ by $m(\alpha)gm(\beta)$ replaces $(c,d)$ by $(\alpha^{*-1}c\beta,\alpha^{*-1}d\beta^{*-1})$ and
$d+cy$ by $\alpha^{*-1}(d+c\,\beta y\beta^*)\beta^{*-1}$, with $\beta y\beta^*$ skew-hermitian; so we may
assume $c=\mathrm{diag}(1_s,0)$. Writing $d$ in blocks, $cd^*+dc^*=0$ gives $d_{21}=0$ and
$d_{11}+d_{11}^*=0$; the rows of $(c\ d)$ are independent, so $d_{22}$ is invertible. Take
$y=\mathrm{diag}(\varepsilon1_s-d_{11},0)$, skew-hermitian; then $d+cy=\bigl(\begin{smallmatrix}
\varepsilon1_s&d_{12}\\0&d_{22}\end{smallmatrix}\bigr)$ is invertible.
\emph{Step 2.} Put $g'=gn(y)$, whose lower right block $d'=d+cy$ is invertible. From $c'd'^*+d'c'^*
=0$, $d'^{-1}c'$ is skew-hermitian, so $n^-(-d'^{-1}c')\in\auU(W)(k)$, and $g'n^-(-d'^{-1}c')=
\bigl(\begin{smallmatrix}*&*\\0&d'\end{smallmatrix}\bigr)$ lies in the stabiliser $P$ of $I$. An
element $p=\bigl(\begin{smallmatrix}a&b\\0&d\end{smallmatrix}\bigr)\in P$ has $d=a^{*-1}$ and
$p\,m(a)^{-1}=\bigl(\begin{smallmatrix}1&ba^*\\0&1\end{smallmatrix}\bigr)$, which is unitary, hence an
$n(b')$. So $g\in\langle m,n,n^-\rangle$.
\end{proof}

\begin{theorem}[Invariance of $\theta$; (W1)]\label{pwf:inv}
Let $W$ be a non-degenerate hermitian space over $E$, $\omega$ as in Definition~\ref{pwf:def} and
$\theta$ a theta functional. Then $\theta\circ\omega(\gamma)=\theta$ on $\psHS(W)^\infty$ for every
$\gamma\in\auU(W)(k)$.
\end{theorem}
\begin{proof}
For $\gamma\in\auU(W)(k)$, $\theta\circ\omega(\gamma)$ is again a theta functional:
$\theta(\omega(\gamma)\psrho(w)F)=\theta(\psrho(\gamma w)\omega(\gamma)F)=\theta(\omega(\gamma)F)$ for
$w\in W(k)$, and $\omega(\gamma)$ is continuous on $\psHS(W)^\infty$ (P1). By Lemma~\ref{pwf:unique},
\begin{equation}\label{pwf:eq-char}
\theta\circ\omega(\gamma)=\kappa(\gamma)\,\theta,\qquad\kappa(\gamma)\in\C^\times,
\end{equation}
and $\kappa$ is a character of $\auU(W)(k)$ (as $\omega$ is a homomorphism and $\theta\ne0$).

\emph{Step 1: $W$ split.} Write $W=I\oplus I^*$ with rational Lagrangians and use the model
$\auS(I^*(\auA))=\bigotimes'_v\auS(I^*_v)$ (P2), in which $\theta=c\sum_{\xi\in I^*(k)}\varphi(\xi)$
by Lemma~\ref{pwf:unique} ($X=I^*$, $Y=I$). For $a\in\auGL(I)(E)$, by Lemma~\ref{pwf:siegel} at all
places, $\omega(m(a))\varphi=\prod_v\upsilon_v(\det a)^{-1}|\det a|_{E_v}^{1/2}\,\varphi(m(a)^{-1}\cdot)
=\varphi(m(a)^{-1}\cdot)$ by (P5) (almost all factors are $1$), and $m(a)^{-1}$ permutes $I^*(k)$:
$\kappa(m(a))=1$. For rational $b$, $\omega(n(b))$ is multiplication by $\beta(x,bx)^{-1}$, which is
$1$ at $x\in I^*(k)$ (P5): $\kappa(n(b))=1$. For $n^-(c)$ write $n^-(c)=w_\varepsilon n(b')
w_\varepsilon^{-1}$ (Lemma~\ref{pwf:gen}); $\kappa$ being a character, $\kappa(n^-(c))=\kappa(n(b'))=1$.
By Lemma~\ref{pwf:gen}, $\kappa=1$.

\emph{Step 2: general $W$.} The space $W\oplus W^-$ is split: $I_\Delta=\{(x,x)\}$ and
$I_\nabla=\{(x,-x)\}$ are complementary Lagrangians. Choose rational $\pwfA$-isotropic
decompositions $W=X\oplus Y$, $W^-=X^-\oplus Y^-$ (Lemma~\ref{pwf:unique}); then $(X\oplus X^-)\oplus
(Y\oplus Y^-)$ is one of $W\oplus W^-$, and in the Schr\"odinger models
$\auS((X\oplus X^-)(\auA))\supset\auS(X(\auA))\otimes\auS(X^-(\auA))$, which realise
$\psHS(W\oplus W^-)=\psHS(W)\widehat\otimes\psHS(W^-)$. Let $\theta^-$, $\Theta$ be theta functionals
for $W^-$ and $W\oplus W^-$. By Lemma~\ref{pwf:unique} in these models, $\Theta(F\otimes F')=
c\,\theta(F)\theta^-(F')$ for $F,F'$ smooth, with $c\ne0$. By \eqref{pwf:eq-double} at every place,
$\omega_{W\oplus W^-}(\gamma,1)=\omega_W(\gamma)\otimes1$. By Step 1 applied to $W\oplus W^-$ and
$(\gamma,1)\in\auU(W\oplus W^-)(k)$:
\[
c\,\theta(\omega_W(\gamma)F)\theta^-(F')=\Theta\bigl(\omega_{W\oplus W^-}(\gamma,1)(F\otimes F')\bigr)
=\Theta(F\otimes F')=c\,\theta(F)\theta^-(F').
\]
Choosing $F'$ with $\theta^-(F')\ne0$ gives $\theta\circ\omega_W(\gamma)=\theta$.
\end{proof}

\begin{remark}\label{pwf:rem-inv}
The proof uses neither Weil's theorem on $\Theta$ nor the product formula for Weil indices: only
the values of the local splittings on the Levi subgroup of a Siegel parabolic
(Lemma~\ref{pwf:siegel}), whose product over the places is the Hecke character
$\upsilon(\det)^{-1}|\det|^{1/2}$, trivial on $\auGL(I)(E)$. Together with Lemma~\ref{pwf:unique}
it proves all of (W1): the commutation relation is part of Definition~\ref{pwf:def}, and existence,
uniqueness and invariance of $\theta$ are Lemma~\ref{pwf:unique} and Theorem~\ref{pwf:inv}.
\end{remark}

\begin{corollary}[Agreement with Weil's lift of rational points]\label{pwf:rk}
Realise $\psHS(W)$ as $\auS(X(\auA))$ for a rational $\pwfA$-isotropic decomposition $W=X\oplus Y$,
and let $r_k\colon\mathrm{Sp}(W)(k)\to\mathrm{Mp}(W)_\auA$ be Weil's lift of the rational points
\cite[no.~40]{WeilActa} (defined there on $Ps(X)_k$, which is identified with $\mathrm{Sp}(W)(k)$ when the
characteristic is not $2$, \cite[p.~181]{WeilActa}). Then $\omega(\gamma)=r_k(\gamma)$ as operators on $\auS(X(\auA))$ for
$\gamma\in\auU(W)(k)$.
\end{corollary}
\begin{proof}
Both operators induce $\gamma$ on the Heisenberg group, so by (P1) $\omega(\gamma)=\lambda(\gamma)
r_k(\gamma)$ with $|\lambda(\gamma)|=1$. The functional $\vartheta(\varphi)=\sum_{\xi\in X(k)}\varphi(\xi)$
is invariant under $r_k(\gamma)$ by \cite[Th\'eor\`eme~6]{WeilActa} (evaluate the function $\Theta$
there at the identity, for $S=r_k(\gamma)$ and every $\varphi$), and under $\omega(\gamma)$ by
Theorem~\ref{pwf:inv} and Lemma~\ref{pwf:unique}. Since $\vartheta\ne0$, $\lambda(\gamma)=1$.
\end{proof}

In particular, Weil's Siegel formula, stated in~\cite[Th\'eor\`eme~5]{WeilSiegel} for the lift $r_k$, applies
verbatim to $\omega$ restricted to the rational points, as the paper uses it.

\subsection{The mixed model and the zero term: (W2), (W2$'$), (W2$_\infty$), (W3)}\label{pwf:ss-mixed}

Let $I\subset W$ be a non-zero isotropic $E$-subspace (over $k$, or over $k_v$ in the local
statements). Choose an isotropic $I^*$ with $\langle\,,\rangle\colon I\times I^*\to E$ perfect and
put $W_I=(I\oplus I^*)^\perp$; then $I^\perp=I\oplus W_I$, and $W_I\to I^\perp/I$ is an isometry, by
which we identify $W_I$ with $I^\perp/I$. Let $P_I$ be the stabiliser of $I$; $p\in P_I$ preserves
$I^\perp$ and induces $\bar p\in\auU(W_I)$. Let $M_I\subset P_I$ be the stabiliser of $I$ and $I^*$
(it preserves $W_I$, and $M_I=\auGL(I)\times\auU(W_I)$, $(a,h)$ acting by $m(a)$ on $I\oplus I^*$ and
by $h$ on $W_I$), and $N_I$ the subgroup of $P_I$ acting trivially on $I$, on $I^\perp/I$ and on
$W/I^\perp$.

\begin{lemma}[Levi decomposition]\label{pwf:levi}
$P_I=N_IM_I$.
\end{lemma}
\begin{proof}
Let $p\in P_I$ and $J=pI^*$: isotropic, of the dimension of $I^*$, with $J\cap I^\perp=p(I^*\cap
I^\perp)=0$. As $W=I^\perp\oplus I^*$, $J$ is the graph $\{x+\alpha x+\beta x\}$ of a map $(\alpha,
\beta)\colon I^*\to I\oplus W_I$, and isotropy of $J$ says $\langle x,\alpha y\rangle+\langle\alpha x,
y\rangle+\langle\beta x,\beta y\rangle=0$. Define $u$ by $u(i)=i$ on $I$, $u(x)=x+\alpha x+\beta x$ on
$I^*$, $u(w)=w+\gamma w$ on $W_I$, where $\gamma\colon W_I\to I$ is defined by $\langle\gamma w,x\rangle
=-\langle w,\beta x\rangle$. The isotropy relation and the definition of $\gamma$ give $\langle uz,uz'
\rangle=\langle z,z'\rangle$ on all pairs of basis vectors, so $u\in\auU(W)$, and $u\in N_I$. Then
$u^{-1}pI^*=I^*$, so $u^{-1}p\in M_I$.
\end{proof}

\emph{The mixed model.} By Lemma~\ref{pwf:sum} for $W=(I\oplus I^*)\oplus W_I$ and the model (P2)
$\auS(I^*)$ of $\psHS(I\oplus I^*)$, $\psHS(W)^\infty=\auS(I^*;\psHS(W_I)^\infty)$, locally and over
$\auA$, where $I\oplus I^*$ acts by (P2) on the variable and $W_I$ by $\psrho_{W_I}$ on the values,
and $\omega_W(m)=\omega_{I\oplus I^*}(m(a))\otimes\omega_{W_I}(h)$ for $m=(a,h)\in M_I$. Over $\auA$
the realisation is unique up to a scalar (P1); we fix it by requiring
\begin{equation}\label{pwf:eq-thetamixed}
\theta_W(F)=\sum_{\xi\in I^*(k)}\theta_{W_I}\bigl(F(\xi)\bigr)\qquad(F\in\auS(I^*(\auA);\psHS(W_I)^\infty)),
\end{equation}
for the theta functionals $\theta_W$, $\theta_{W_I}$ fixed in \S\ref{ps:ss-weil}. This is possible:
the right side converges absolutely ($\theta_{W_I}\circ F\in\auS(I^*(\auA))$ by continuity of
$\theta_{W_I}$), is continuous in $F$, non-zero, and invariant under $\psrho(w)$, $w\in W(k)$ (for
$w\in I(k)$ it multiplies $F(\xi)$ by $[\xi,w]=1$; for $w\in I^*(k)$ it translates the variable by
$w$; for $w\in W_I(k)$ it applies $\psrho_{W_I}(w)$ to the values, and $\theta_{W_I}$ is invariant;
the factors $\beta$ between rational points are $1$); so it is a multiple of $\theta_W$ by
Lemma~\ref{pwf:unique}, and the scalar is absorbed in the realisation.

\begin{definition}[Zero term]\label{pwf:defZ}
$\psZ_I\colon\psHS(W)^\infty\to\psHS(W_I)^\infty$ is $\psZ_IF=F(0)$ in the realisation just fixed;
locally, $\psZ_{I,v}F=F(0)$ in the local mixed model (defined up to a scalar).
\end{definition}

\begin{lemma}[Characterisation of the zero term]\label{pwf:Zunique}
Let $T\colon\psHS(W)^\infty\to\psHS(W_I)^\infty$ be continuous linear with $T\psrho(w)=
\psrho(\bar w)T$ for all $w\in I^\perp$ (over $\auA$, resp.\ over $k_v$), $\bar w$ the image in
$W_I=I^\perp/I$. Then $T=c\,\psZ_I$ (resp.\ $c\,\psZ_{I,v}$).
\end{lemma}
\begin{proof}
Locally this is Lemma~\ref{pwf:dist} with $X=I^*$, $Y=I$, $W'=W_I$: the hypothesis for $w\in I$ is
$T\psrho(w)=T$, and for $w\in W_I$ it is $T\psrho(w)=\psrho_{W_I}(w)T$. Over $\auA$: the local
argument in the proof of Lemma~\ref{pwf:dist} uses only the invariance under $\psrho(I_v)$, so for
each $v$ and each fixed vector $F^v$ away from $v$, $T(\varphi_v\otimes F^v)=\varphi_v(0)\,
T(\varphi_v^0\otimes F^v)$, where $\varphi^0_v$ is the standard vector (with $\varphi^0_v(0)=1$) at
the finite places outside a finite set and any fixed vector with value $1$ at $0$ elsewhere. As every
element of $\auS(I^*(\auA))\otimes\psHS(W_I)^\infty$ is a finite sum of tensors that are standard
outside a finite set of places, $TF=T_0(F(0))$ on these, with $T_0(h)=T(\otimes_v\varphi^0_v\otimes h)$,
and by continuity everywhere. $T_0$ commutes with $\psrho_{W_I}(W_I(\auA))$, so it is a scalar (P1).
\end{proof}

\begin{proposition}[(W2), (W2$'$), (W2$_\infty$)]\label{pwf:W2}
With $\theta_W,\theta_{W_I}$ as in \S\ref{ps:ss-weil} and $di$ the Haar measure of mass $1$ on
$[I]=I(k)\backslash I(\auA)$:
\begin{itemize}
\item[(a)] $\theta_{W_I}(\psZ_IF)=\int_{[I]}\theta_W(\psrho(i)F)\,di$; this property together with
\textup{(b)} characterises $\psZ_I$.
\item[(b)] $\psZ_I\psrho(w)=\psrho(\bar w)\psZ_I$ for $w\in I^\perp(\auA)$.
\item[(c)] $\psZ_I\,\omega_W(p)=\chi_I(p)\,\omega_{W_I}(\bar p)\,\psZ_I$ for $p\in P_I(\auA)$, with
$\chi_I(p)=\upsilon(\det p|_I)^{-1}|\det p|_I|^{1/2}$ ($\det p|_I\in\auA_E^\times$, $|\cdot|$ the
idele norm of $E$). Locally, $\psZ_{I,v}\omega_{W,v}(p)=\chi_{I,v}(p)\omega_{W_I,v}(\bar p)
\psZ_{I,v}$ with $\chi_{I,v}(p)=\upsilon_v(\det p|_I)^{-1}|\det p|_I|_{E_v}^{1/2}$; at a real place
and for a line $I$ on which $p$ acts by $a\in\C^\times$, $\chi_{I,\tau}(p)=\upsilon_\tau(a)^{-1}|a|$
\textup{(W2$_\infty$)}. In particular $\chi_I=1$ on $N_I(\auA)$ and on $P_I(k)$.
\item[(d)] For isotropic $I\subset I'$, $\psZ_{I'}=\psZ_{I'/I}\circ\psZ_I$, and $\psZ_I=\psZ_{I,\infty}
\otimes\psZ_{I,f}$ (for suitable normalisations of the local maps).
\end{itemize}
\end{proposition}
\begin{proof}
(a) By (P2), $(\psrho(i)F)(x)=[x,i]F(x)$ for $i\in I(\auA)$, so by \eqref{pwf:eq-thetamixed}
$\theta_W(\psrho(i)F)=\sum_\xi[\xi,i]\,\theta_{W_I}(F(\xi))$, with $\sum_\xi|\theta_{W_I}(F(\xi))|<\infty$
independent of $i$. For $\xi\in I^*(k)$, $i\mapsto[\xi,i]$ is a character of $[I]$, non-trivial if
$\xi\ne0$ (P5). Integrating term by term leaves $\theta_{W_I}(F(0))$. If $T$ satisfies (a) and (b),
then $T=c\psZ_I$ by Lemma~\ref{pwf:Zunique}, and (a) for both gives $c\,\theta_{W_I}\circ\psZ_I=
\theta_{W_I}\circ\psZ_I\ne0$, so $c=1$.

(b) For $w=i+w'$, $i\in I$, $w'\in W_I$: $\beta(i,w')=1$, $\psrho(w)=\psrho(i)\psrho_{W_I}(w')$, and
$(\psrho(w)F)(0)=[0,i]\psrho_{W_I}(w')F(0)=\psrho(\bar w)F(0)$.

(c) \emph{Local statement.} For $p\in P_I(k_v)$ the operator $T=\omega_{W_I}(\bar p)^{-1}\psZ_{I,v}
\omega_W(p)$ satisfies the hypothesis of Lemma~\ref{pwf:Zunique}: for $w\in I^\perp$, $T\psrho(w)=
\omega(\bar p)^{-1}\psZ_{I,v}\psrho(pw)\omega(p)=\omega(\bar p)^{-1}\psrho(\bar p\bar w)\psZ_{I,v}
\omega(p)=\psrho(\bar w)T$, using (b) locally and $\overline{pw}=\bar p\bar w$. So $\psZ_{I,v}\omega(p)=
\chi(p)\omega(\bar p)\psZ_{I,v}$, and $\chi$ is a continuous character of $P_I(k_v)$ ($\psZ_{I,v}$ is
onto). It suffices to compute $\chi$ on $M_I$ and $N_I$ (Lemma~\ref{pwf:levi}). For $m=(a,h)\in
M_I$, Lemma~\ref{pwf:sum} and Lemma~\ref{pwf:siegel} give $\omega_W(m)F=\upsilon_v(\det a)^{-1}
\pwfL_a\otimes\omega_{W_I}(h)F$, and $(\pwfL_a\varphi)(0)=|\det a|_{E_v}^{1/2}\varphi(0)$; as
$\bar m=h$ and $\det m|_I=\det a$, $\chi(m)=\chi_{I,v}(m)$. On $N_I$, $\chi$ is a character with
$\chi(qnq^{-1})=\chi(n)$ for $q\in P_I$. Write $n\in N_I$ as in the proof of
Lemma~\ref{pwf:levi}: $n(x)=x+\alpha x+\beta x$ on $I^*$; then $n\mapsto\beta\in\auHom(I^*,W_I)$ is a
homomorphism (as $(nn')(x)=x+(\alpha+\alpha'+\gamma\beta')x+(\beta+\beta')x$) whose kernel $Z$ is the
abelian group of the $n$ with $\beta=0$, parametrised additively by the skew maps $\alpha$. Let
$t=m(2\cdot\mathrm{id}_I)\in M_I$ ($2$ on $I$, $\frac12$ on $I^*$, $1$ on $W_I$). Then $tnt^{-1}$
has parameters $(4\alpha,2\beta)$ ($tnt^{-1}(x)=t\,n(2x)=x+4\alpha x+2\beta x$). On $Z$:
$\chi(\alpha)=\chi(4\alpha)=\chi(\alpha)^4$, so $\chi^3=1$ on the vector group $Z$, and $\chi(\alpha)=
\chi(\alpha/3)^3=1$. Hence $\chi$ factors through the abelian group $N_I/Z$, on which
$\chi(\beta)=\chi(2\beta)=\chi(\beta)^2$, so $\chi=1$. Since $\bar n=1$ and $\det n|_I=1$,
$\chi=\chi_{I,v}$ on $N_I$, hence on $P_I(k_v)$. At a real place $\tau$ of $k$ one has $E_\tau\cong\C$
with module $|a|_{E_\tau}=|a|^2$, so $|a|_{E_\tau}^{1/2}=|a|$, which is (W2$_\infty$).

\emph{Global statement.} By (P1) the global realisation is the tensor product of local ones up to a
scalar, so $\psZ_I=c\bigotimes_v\psZ_{I,v}$, and $\omega_W=\bigotimes_v\omega_{W,v}$,
$\omega_{W_I}=\bigotimes_v\omega_{W_I,v}$; so (c) holds with $\chi_I=\prod_v\chi_{I,v}$, which is the
stated formula (almost all factors are $1$ on a given $p$). It is $1$ on $N_I(\auA)$ by the local
case and on $P_I(k)$ by (P5).

(d) The second statement is the construction ($\psZ_I=c\bigotimes_v\psZ_{I,v}$). For the first,
$I'/I$ is an isotropic subspace of $W_I$ with $W_{I'}=(I'/I)^\perp/(I'/I)$. The map
$T=\psZ_{I'/I}\psZ_I$ satisfies the hypothesis of Lemma~\ref{pwf:Zunique} for $I'$: for $w\in
I'^\perp\subset I^\perp$, $\psZ_I\psrho(w)=\psrho(\bar w)\psZ_I$ with $\bar w\in(I'/I)^\perp$, and then
(b) for $I'/I$. So $T=c\psZ_{I'}$. By (a) three times and Fubini (choose a complement $J$ of $I$ in
$I'$; $\psrho(i)\psrho(j)=\psrho(i+j)$ as $\beta=1$ on the isotropic $I'$, and $[I']$ fibres over
$[J]=[I'/I]$ with fibre $[I]$, all measures of mass $1$),
\[
\theta_{W_{I'}}(TF)=\int_{[J]}\theta_{W_I}\bigl(\psZ_I\psrho(j)F\bigr)dj=\int_{[J]}\int_{[I]}
\theta_W(\psrho(i)\psrho(j)F)\,di\,dj=\int_{[I']}\theta_W(\psrho(i')F)\,di'=\theta_{W_{I'}}(\psZ_{I'}F),
\]
using (b) for $\psZ_I\psrho(j)=\psrho(\bar j)\psZ_I$. So $c=1$.
\end{proof}

\begin{proposition}[The identity of (W3)]\label{pwf:W3}
For $F\in\psHS(W)^\infty$, $\theta_W(F)=\sum_{w\in W(k)/I^\perp(k)}\theta_{W_I}(\psZ_I\psrho(w)F)$; the
terms do not depend on the representatives, and the series converges absolutely.
\end{proposition}
\begin{proof}
$I^*(k)$ is a set of representatives of $W(k)/I^\perp(k)$, and $\psZ_I\psrho(x)F=F(x)$ for $x\in
I^*(k)$ by (P2); so the series is \eqref{pwf:eq-thetamixed}, absolutely convergent. For $u\in
I^\perp(k)$, $\psrho(u)\psrho(w)=\beta(u,w)\psrho(u+w)$ with $\beta(u,w)=1$ for rational $u,w$ (P5); so $\theta_{W_I}(\psZ_I\psrho(w+u)F)=
\theta_{W_I}(\psZ_I\psrho(u)\psrho(w)F)=\theta_{W_I}(\psrho(\bar u)\psZ_I\psrho(w)F)=\theta_{W_I}(\psZ_I
\psrho(w)F)$, by (b) and the invariance of $\theta_{W_I}$ under $\psrho(\bar u)$, $\bar u\in W_I(k)$.
\end{proof}

\begin{remark}\label{pwf:rem-W3}
The moderate growth of $(g,h)\mapsto\sum_w|\theta_{W_I}(\psZ_I\psrho(w)\omega(g,h)F)|$ asserted in
(W3) is not proved in this section: it is Lemma~\ref{pz:absgrowth}(b) (for $F=\Psi\otimes f$, with
$A=4s$ and the constant $C(2D)^{4s}\|\Psi\|_{2s}$ there), whose proof uses (W0), (W1), (W2), (W4) but
not (W3); its part (a) gives again the absolute convergence of Proposition~\ref{pwf:W3}. Proposition~\ref{pwf:W2}(c) and Lemma~\ref{pwf:Zunique} also give the uniqueness
of the zero term used in \S\ref{pf:ss-flip} and the identity $\psZ_{L_0}\circ\omega(\gamma)=
\psZ_{\gamma^{-1}L_0}$ for rational $\gamma$ and a Lagrangian $L_0$: both sides
satisfy (a) and (b) for $\gamma^{-1}L_0$, by Theorem~\ref{pwf:inv} (this is the step in the proof of
{\cite[Lemma~2.14]{P53}}).
\end{remark}

\subsection{The model at split places: (W5) and \eqref{pl:eq-w5}}\label{pwf:ss-split}

For hermitian spaces $V,V'$ of ranks $n,n'$ and $W=\auHom(V,V')$, $\omega_{\auHom(V,V')}(g,h)$ is
$\omega_W$ of the image of $(g,h)$ under $\auU(V)\times\auU(V')\to\auU(W)$, $(g,h)w=hwg^{-1}$, and
$\omega_\mu$ is its determinant twist (\S\ref{ps:ss-weil}).

\begin{proposition}[(W5)]\label{pwf:W5}
Let $v=w\aubar w$ be split in $E$. Then $\psHS_v(\auHom(V,V'))$ is realised on
$\auS(\auHom(V_w,V'_w))$, with $\psrho(a_w,0)\Psi(x)=\Psi(x+a_w)$ and $\psrho(0,a_{\aubar w})\Psi(x)=
[x,a_{\aubar w}]\Psi(x)$, and
\[
(\omega_{\auHom(V,V')}(g,h)\Psi)(x)=\upsilon_w(\det h_w)^{-n}\upsilon_w(\det g_w)^{n'}\,
|\det h_w|_w^{-n/2}|\det g_w|_w^{n'/2}\,\Psi(h_w^{-1}xg_w).
\]
If $n,n'$ are even, then $(\omega_\mu(g,h)\Psi)(x)=\chi_w(\det h_w)\chi_w(\det g_w)^{-1}|\det
h_w|_w^{-n/2}|\det g_w|_w^{n'/2}\Psi(h_w^{-1}xg_w)$, where $\chi=\mu\circ(z\mapsto z/\aubar z)$; this
is \eqref{pl:eq-w5}.
\end{proposition}
\begin{proof}
$W_v=W_w\oplus W_{\aubar w}$ with $W_w=\auHom(V_w,V'_w)$, both summands isotropic
(\S\ref{pwf:ss-local}, proof of Lemma~\ref{pwf:siegel}); the realisation and the Heisenberg operators
are (P2) with $X=W_w$, $Y=W_{\aubar w}$. The element $(g,h)$ acts on $W_w$ by $x\mapsto h_wxg_w^{-1}$,
whose determinant is $\det(h_w)^n\det(g_w)^{-n'}$. So (D3) gives the first formula. For the second,
$\auU(1)(k_v)=\{(t,t^{-1})\}\cong k_v^\times$ (component at $w$), and $\upsilon(t,t^{-1})=\upsilon_w(t)
\upsilon_{\aubar w}(t)^{-1}=\upsilon_w(t)^2$ because $\upsilon_{\aubar w}=\upsilon_w^{-1}$ on $k_v^\times$
($\upsilon|_{\auA_k^\times}=\epsilon_{E/k}$ is trivial at $v$). Hence, for even $n$,
$\upsilon_n(\det h)=(\upsilon|_{\auU(1)})^{n/2}(\det h)=\upsilon_w(\det h_w)^n$, and similarly for
$g$; the determinant twists $\upsilon_n(\det h)\upsilon_{n'}(\det g)^{-1}$ cancel the powers of
$\upsilon_w$ in the first formula. Finally $\mu(\det h)=\mu(\det h_w,(\det h_w)^{-1})$, and for
$z\in E_w^\times\subset E_v^\times$ one has $z/\aubar z=(z,z^{-1})$, so $\chi_w(z)=\mu(z,z^{-1})$:
$\mu(\det h)\mu(\det g)^{-1}=\chi_w(\det h_w)\chi_w(\det g_w)^{-1}$.
\end{proof}

\begin{remark}\label{pwf:rem-W5}
The normalisation (D3) is the only one of the form $\nu(\det g_w)|\det g_w|^{-1/2}$ for which the
Siegel Levi acts by the same Hecke character $\upsilon(\det)^{-1}|\det|^{1/2}$ at all places
(Lemma~\ref{pwf:siegel}), which is what Theorem~\ref{pwf:inv} needs. With it, the powers of
$\upsilon$ cancel in \eqref{pl:eq-w5}, as stated in (N5).
\end{remark}

\subsection{The formal degree in \texorpdfstring{{\cite[Proposition~5.2]{P51}}}{[Moment, Proposition~5.2]}}\label{pwf:ss-haar}

\begin{proposition}\label{pwf:haar}
{\cite[Proposition~5.2]{P51}} holds, with the same statement, for every parameter to which it is
applied in the paper. More precisely: for $(\lambda,s)$ satisfying the gap condition
{\cite[(49)]{P51}} and $dg$ the Haar measure compatible with the Haar measure of mass $1$ on
$\auU(N)$, $\int_G|c|^2dg=\bigl(\prod_{i<j}(\lambda_i-\lambda_j)/(j-i)\bigr)^{-1}$ for the coefficient
$c$ of a unit vector of the lowest $K$-type; and the parameters of {\cite[\S5]{P51}} satisfy
{\cite[(49)]{P51}}.
\end{proposition}
\begin{proof}
The first assertion is {\cite[Corollary~4.43]{P51}}: the formal degree $d$ of $\pi(\lambda,s)$ for
$dg$ is defined there by $\int_G|c|^2dg=d^{-1}$ with $c$ the coefficient of a unit highest weight
vector of the lowest $K$-type, and it is proved that $d=\prod_{i<j}(\lambda_i-\lambda_j)/(j-i)$; the
proof does not use Harish-Chandra's formula. The coefficient $c_\tau$ of {\cite[\S5]{P51}} is of this
kind ($|c_\tau|^2$ does not depend on the convention for inner products), with $N=r$ and $\lambda=
\varsigma_\tau$. Entries of $\varsigma_\tau$ attached to coordinates of opposite signs differ by at
least $\Delta=2r^8$ ({\cite[\S3.2.2]{P51}}: the first $k-1$ entries are $>-k/2$, the others are
$-\kappa(\tau)-j\Delta$, $1\le j\le m$, with $\kappa(\tau)>0$, and consecutive $j$ have opposite
signs; an entry $>-k/2$ and an entry $-\kappa(\tau)-j\Delta$ with $j\ge2$ differ by more than
$2\Delta-k/2>\Delta$), and $2r^8\ge500r^3+10r^5$ because $r\ge100$ ({\cite[Lemma~13.4]{P46}},
(S7)(b)): indeed $500r^3+10r^5\le510r^5\le2r^8$ as soon as $r^3\ge255$. So {\cite[(49)]{P51}} holds
with $n=r\ge2$. (The other standing hypothesis of {\cite[\S4.2]{P51}}, $\lambda\in(\frac{n-1}2+\Z)^n$ with
distinct entries, holds for $\lambda=\varsigma_\tau$ because it is the Harish-Chandra parameter of a
discrete series of $\auU(k+\varrho,\varrho+1)$.)
\end{proof}

\subsection{The archimedean scalar on the maximal compact subgroup}\label{pwf:ss-delta}

Let $\tau$ be a real place, $W$ of signature $(P,Q)$ at $\tau$, $W_\tau=W^+\oplus W^-$ the
decomposition into a positive and a negative definite part, and $K=\auU(W^+)\times\auU(W^-)$. Put
$s=1/\sqrt d$, so that $\tau(\jmath)=is$, and let $\|u\|^2=\langle u^+,u^+\rangle-\langle u^-,u^-\rangle$
be the majorant ($u=u^++u^-$); it is $K$-invariant.

\emph{The vacuum.} \S\ref{ps:ss-weil} defines the Fock model, and its vacuum $1$, for positive
definite spaces. For any non-degenerate $W_\tau$ we call \emph{vacuum} a unit vector $v\in\psHS_\tau(W)$
with
\begin{equation}\label{pwf:eq-vaccoef}
\langle\psrho(u)v,v\rangle=e^{-\pi s\|u\|^2}\qquad(u\in W_\tau).
\end{equation}
(a) \emph{Uniqueness.} A unit vector $v$ of an irreducible unitary representation is determined up
to a scalar of modulus $1$ by $u\mapsto\langle\psrho(u)v,v\rangle$: if $v,v'$ have the same coefficient,
then $\psrho(u)v\mapsto\psrho(u)v'$ extends (by $\|\sum c_i\psrho(u_i)v\|^2=\sum c_i\aubar c_j\beta(\cdot)
\langle\psrho(u_i-u_j)v,v\rangle$, the same for $v'$) to an isometry of $\psHS_\tau(W)$ onto itself
commuting with all $\psrho(u)$, a scalar by (P1).
(b) \emph{Existence.} For positive definite $W$ the Fock vacuum $1$ satisfies \eqref{pwf:eq-vaccoef}:
by \eqref{pk:eq-fock}, $\langle\psrho(y)1,1\rangle=e^{-|y|^2/2}\int e^{z\cdot y}e^{-|z|^2}d\mu(z)=e^{-|y|^2/2}$
(mean value property), and $|y|^2/2=\pi s\|u\|^2$ for $y=c_\psi u$, $c_\psi^2=2\pi s$. For negative
definite $W$, $W^-$ is positive definite and $\beta_W=\aubar{\beta_{W^-}}$ (as $\psi(-t)=\aubar{\psi(t)}$);
so the complex conjugate Hilbert space $\aubar{\psHS_\tau(W^-)}$ with the same operators $\psrho(u)$ is a
Heisenberg representation of $W$, and the vector $1$ in it has coefficient $\aubar{e^{-\pi s\|u\|^2}}=
e^{-\pi s\|u\|^2}$ (the majorants of $W$ and $W^-$ coincide). In this realisation the Fock vectors of
$W$ are the \emph{antiholomorphic} functions of the Fock variable of $W^-$, i.e.\ the holomorphic
functions of the coordinate $w$; this is the convention of {\cite[\S3.3.2]{P54}} and of (A$_\infty$) in
{\cite[\S3.3.1]{P54}} (complex structure $-J_\tau$ for $W^-$). For general $W_\tau=W^+\oplus W^-$, $\psHS_\tau(W)=\psHS_\tau(W^+)\widehat
\otimes\psHS_\tau(W^-)$ (an intertwiner, $\beta(W^+,W^-)=1$), and $v^+\otimes v^-$ satisfies
\eqref{pwf:eq-vaccoef}, the majorant being additive. Thus the vacuum exists, is unique up to a
scalar, and is compatible with orthogonal sums of definite spaces of either sign.
(c) \emph{$K$ preserves the vacuum line.} For $\gamma\in K$, $\omega_\tau(\gamma)v$ has coefficient
$\langle\psrho(u)\omega(\gamma)v,\omega(\gamma)v\rangle=\langle\psrho(\gamma^{-1}u)v,v\rangle=
e^{-\pi s\|\gamma^{-1}u\|^2}=e^{-\pi s\|u\|^2}$, so by (a) $\omega_\tau(\gamma)v=\Delta(\gamma)v$ with
$|\Delta(\gamma)|=1$, and $\Delta=\Delta_W$ is a continuous character of $K$. The operator $F\mapsto
F\circ\gamma^{-1}$ on Fock vectors ({\cite[\S3.3.2]{P54}}) also conjugates $\psrho(u)$ into $\psrho(\gamma u)$
and fixes the vacuum, so it equals $\Delta(\gamma)^{-1}\omega_\tau(\gamma)$ by (P1). Hence
{\cite[(30)]{P54}} is the assertion $\Delta(\gamma)=\aubar{\det\gamma^-}$.
(d) \emph{Naturality.} An isometry $\sigma\colon W_\tau\to W'_\tau$ carries $\omega_{W,\tau}$ to
$\omega_{W',\tau}$ (Lemma~\ref{pwf:indep} for split spaces; for non-split spaces, Lemma~\ref{pwf:indep} applied to the
isometry $\sigma\oplus\sigma^-$ of the split doubles, together with the characterisation (D2)), and
an intertwiner along $\sigma$ carries vacuum to vacuum (by (a), the majorant being transported); so
$\Delta_{W'}(\sigma\gamma\sigma^{-1})=\Delta_W(\gamma)$.

\begin{lemma}[Reduction to lines]\label{pwf:delta-red}
There are integers $e_+,e_-$, independent of $W$, with $\Delta(\gamma^+,\gamma^-)=\det(\gamma^+)^{e_+}
\det(\gamma^-)^{e_-}$ for all $W$; and $e_++e_-=-1$. Hence {\cite[(30)]{P54}} holds for all $W$
if and only if $e_+=0$, i.e.\ if and only if, for a positive definite line $\ell$, $\omega_{\ell,\tau}(z)$
fixes the vacuum for $z\in\auU(1)(\R)$.
\end{lemma}
\begin{proof}
By Lemma~\ref{pwf:sum} for $W_\tau=W^+\oplus W^-$ and for orthogonal decompositions of $W^\pm$ into
lines, and by (b) above (the vacuum of an orthogonal sum is the tensor product of the vacua),
$\Delta$ restricted to the diagonal torus of $K$ is the product of the scalars $\Delta_\pm(z)=
z^{e_\pm}$ of the positive and negative lines (continuous characters of $\auU(1)$; by (d) they do not
depend on the line, all positive, resp.\ negative, lines being isometric); a character of $\auU(P)\times\auU(Q)$ is determined by its
restriction to the torus, which gives the first assertion. For the relation, take $W=\ell\oplus\ell^-$
with $\ell$ a positive line; it is split, with Lagrangians $I_\Delta=\{(x,x)\}$, $I_\nabla=
\{(x,-x)\}$. The element $(z,z)\in K$ equals $m(z)$ in the notation of Lemma~\ref{pwf:siegel}
($I=I_\Delta$; on $I_\nabla$ it acts by $z=\aubar z^{-1}$). By Lemma~\ref{pwf:siegel},
$\omega_\tau(m(z))=\upsilon_\tau(z)^{-1}\pwfL_z=z^{-1}\pwfL_z$ in the model $\auS(I_\nabla)$, where
$(\pwfL_z\varphi)(x)=\varphi(z^{-1}x)$ ($|z|=1$). In this model the vacuum is a multiple of the radial
Gaussian $\varphi_0$ of \eqref{pwf:eq-phi0} below (that identification uses only (a) and a Gaussian
integral, not this lemma). As $\pwfL_z\varphi_0=\varphi_0$, $\omega_\tau(m(z))\varphi_0=z^{-1}\varphi_0$.
Hence $\Delta(z,z)=z^{-1}$, i.e.\ $e_++e_-=-1$.
\end{proof}

\begin{proposition}[$e_+=0$; {\cite[(30)]{P54}}]\label{pwf:eplus}
Let $\ell$ be a positive definite hermitian line at the real place $\tau$ and $\ell^-$ the negative
line. For $(z_1,z_2)\in\auU(\ell)(\R)\times\auU(\ell^-)(\R)$ the operator $\omega_{\ell\oplus\ell^-,\tau}
(z_1,z_2)$ multiplies the Fock vacuum by $z_2^{-1}$. Hence $e_+=0$, $e_-=-1$ in
Lemma~\ref{pwf:delta-red} and~{\cite[(30)]{P54}} holds: $\Delta(\gamma)=\aubar{\det\gamma^-}$.
\end{proposition}
\begin{proof}
By Lemma~\ref{pwf:sum}, $\omega_{\ell\oplus\ell^-,\tau}(z_1,z_2)=\omega_{\ell,\tau}(z_1)\otimes
\omega_{\ell^-,\tau}(z_2)$, and the vacuum of $\ell\oplus\ell^-$ is the tensor product of the vacua; so the
scalar in question is $\Delta_{\ell}(z_1)\Delta_{\ell^-}(z_2)=z_1^{e_+}z_2^{e_-}$; so it suffices to
compute it for $z_1\ne z_2$, and then $e_+=0$, $e_-=-1$, and the last assertion is
Lemma~\ref{pwf:delta-red}. Write $s=1/\sqrt d$, so that $\tau(\jmath)=is$
(\S\ref{ps:ss-weil}), and $\psi_{k,\tau}(t)=e^{-2\pi it}$.

\emph{Kudla's coordinates.} Let $W=\ell\oplus\ell^-=\C^2$ with $\langle u,w\rangle=\aubar u_1w_1-\aubar
u_2w_2$, and Kudla's skew-hermitian form $\langle u,w\rangle_K=\jmath\langle w,u\rangle$ (linear in $u$).
The vectors $e=(2\jmath)^{-1}(1,-1)$, $f=(1,1)$ span $X=I_\nabla$, $Y=I_\Delta$ and satisfy
$\langle e,e\rangle_K=\langle f,f\rangle_K=0$, $\langle e,f\rangle_K=1$, so in the coordinates
$xe+yf\leftrightarrow(x,y)$ the form is Kudla's (0.6) with $\epsilon=1$ \cite[p.~364]{KudlaSplitting}.
With $V=E$, $(v_1,v_2)=\aubar v_1v_2$ (so $m=1$, signature $(p,q)=(1,0)$) and $\kappa=\frac12$,
$\langle\!\langle(x,0),(0,y)\rangle\!\rangle=\operatorname{Re}(\aubar xy)$, and (P3) is read with
$\psi_K(t)=\psi_{k,\tau}(2t)=e^{-4\pi it}$, so $\eta(t)=e^{-2\pi it}$, and with $\xi=\upsilon_\tau$,
$\xi(z)=z/|z|$. The paper's element $g=(z_1,z_2)$ is Kudla's element acting on rows by the map
$g^{-1}$ (\S\ref{pwf:ss-def}, left actions). With $a_i=\aubar z_i=z_i^{-1}$, $g^{-1}f=(a_1,a_2)$ and
$g^{-1}e=(2\jmath)^{-1}(a_1,-a_2)$, whence Kudla's matrix
\[
\begin{pmatrix}\alpha&\beta\\\gamma&\delta\end{pmatrix}=\begin{pmatrix}\frac{a_1+a_2}2&\frac{a_1-a_2}{4\jmath}\\
\jmath(a_1-a_2)&\frac{a_1+a_2}2\end{pmatrix}.
\]
As $\gamma\ne0$, $\ker\gamma=0$ and $g$ lies in the cell $j=1$ (1.11). Writing
$g=p_1\tau p_2$ with $p_i=\bigl(\begin{smallmatrix}\alpha_i&\beta_i\\0&\delta_i\end{smallmatrix}\bigr)\in
P_Y$ and $\tau=\bigl(\begin{smallmatrix}0&-1\\1&0\end{smallmatrix}\bigr)$ (1.12), one finds
$\gamma=\delta_1\alpha_2$, and $\alpha_2\aubar\delta_2=1$ (unitarity of $p_2$); so $x(g)=\delta_1\delta_2=
\gamma/|\alpha_2|^2$ and $\xi(x(g))=\gamma/|\gamma|$.

\emph{The Weil index.} For Kudla's character $e(t)=e^{2\pi it}$, comparison of Theorem~3.1 with
Corollary~3.2 \cite[pp.~378--379]{KudlaSplitting} gives $\gamma_\R(\eta_0\circ RV)=i^{p-q}=i$ for
$\eta_0(t)=e(t/2)$. Our $\eta(t)=\aubar{\eta_0(2t)}$. Kudla's $\gamma_F(\eta\circ L)$ is Weil's
$\gamma(f)$ for the character of second degree $f(x)=\eta((x,x)_L)$ \cite[p.~364]{KudlaSplitting}, and
Weil's $\gamma(f)$ is the number of modulus $1$ such that $\int(\int\Phi(x-y)f(y)\,dy)\,dx=\gamma(f)
|\varrho_f|^{-1/2}\int\Phi(x)\,dx$ for all $\Phi\in\auS(G)$, $|\varrho_f|>0$ the module of the
symmetric isomorphism attached to $f$ \cite[Th\'eor\`eme~2, p.~161, and Corollaire~2, p.~162]{WeilActa}.
Two consequences, for $G=\R^N$: (i) if $f_\lambda(x)=f(\lambda^{1/2}x)$ with $\lambda>0$ (this is the
character $\eta(\lambda\,\cdot)\circ L$), the substitutions $y\mapsto\lambda^{-1/2}y$, $x\mapsto
\lambda^{-1/2}x$ turn the left side for $(f_\lambda,\Phi)$ into $\lambda^{-N}$ times the left side for
$(f,\Phi(\lambda^{-1/2}\cdot))$, so $\gamma(f_\lambda)|\varrho_{f_\lambda}|^{-1/2}=\lambda^{-N/2}\gamma(f)
|\varrho_f|^{-1/2}$; comparing arguments (take $\int\Phi\ne0$), $\gamma(f_\lambda)=\gamma(f)$. (ii)
Conjugating the identity for $(f,\aubar\Phi)$ gives it for $(\aubar f,\Phi)$ with the factor
$\aubar{\gamma(f)}|\varrho_f|^{-1/2}$, and $\aubar f$ is attached to $-\varrho_f$, of the same module;
so $\gamma(\aubar f)=\aubar{\gamma(f)}$. Hence $\gamma_\R(\eta\circ RV)=\aubar{\gamma_\R(\eta_0\circ RV)}=
-i$ and $\beta_Y(g)=\xi(x(g))(-i)^{-1}=i\,\gamma/|\gamma|$.

\emph{The vacuum in this model.} Let
\begin{equation}\label{pwf:eq-phi0}
\varphi_0(x)=e^{-\pi|x|^2/s}\qquad(x\in X=\C,\ \text{the point }xe).
\end{equation}
With
$\psrho(xe)\varphi(u)=\varphi(u+x)$, $\psrho(yf)\varphi(u)=e^{-4\pi i\operatorname{Re}(\aubar uy)}\varphi(u)$
and $\psrho(xe+yf)=\beta(xe,yf)^{-1}\psrho(xe)\psrho(yf)$ (P2), a Gaussian integral gives
$\langle\psrho(xe+yf)\varphi_0,\varphi_0\rangle/\|\varphi_0\|^2=\exp(-\pi|x|^2/(2s)-2\pi s|y|^2)$
(the phase $[\frac x2,y]$ produced by centring the integral is $\beta(xe,yf)$ and cancels). In the
paper's coordinates $xe+yf=(\frac x{2\jmath}+y,-\frac x{2\jmath}+y)$, and $|x|^2/(2s)+2s|y|^2=s(|u_1|^2+
|u_2|^2)=s\|u\|^2$ for $u=xe+yf$. By (a), (b) of \S\ref{pwf:ss-delta} (the vacuum of $\ell\oplus\ell^-$
is $v_0\otimes v_0^-$, with coefficient $e^{-\pi s\|u\|^2}$, $v_0^-$ the vacuum of the negative line
defined there), $\varphi_0/\|\varphi_0\|$ is the vacuum in the model $\auS(X)$, and by (c) there
$\omega(g)\varphi_0=c\,\varphi_0$ with $|c|=1$.

\emph{The scalar.} By Rao's formula at $x=0$ \cite[p.~400]{KudlaSplitting},
$c=(\omega(g)\varphi_0)(0)=\beta_Y(g)\int_\C\psi_K\bigl(\tfrac12\langle\!\langle y\gamma,y\delta\rangle\!\rangle
\bigr)\varphi_0(y\gamma)\,d\mu(y)$ with $\mu$ a positive Haar measure. Put $w=z_1\aubar z_2$. Then
$\langle\!\langle y\gamma,y\delta\rangle\!\rangle=|y|^2\operatorname{Re}(\aubar\gamma\delta)$ and
$\aubar\gamma\delta=-\frac\jmath2(z_1-z_2)(\aubar z_1+\aubar z_2)=s\operatorname{Im}w$, while
$|\gamma|^2=s^2|z_1-z_2|^2=s^2|1-w|^2$. So the integrand is $\exp(-\pi s|y|^2(|1-w|^2+2i\operatorname{Im}w))
=\exp(-2\pi s|y|^2(1-\aubar w))$, and the integral is a positive multiple of $(1-\aubar w)^{-1}$,
whose phase is $(1-w)/|1-w|$. On the other hand $\gamma/|\gamma|=i(\aubar z_1-\aubar z_2)/|z_1-z_2|=
-i\aubar z_2(1-\aubar w)/|1-w|$, so $\beta_Y(g)=\aubar z_2(1-\aubar w)/|1-w|$. Therefore $c$ is a
positive multiple of $\aubar z_2(1-\aubar w)(1-w)/|1-w|^2=\aubar z_2$, and $|c|=1$ gives $c=z_2^{-1}$.
\end{proof}

\begin{proof}[Second proof, without the Weil index and Rao's formula]
This proof uses only Lemma~\ref{pwf:siegel}, the homomorphism property of $\omega$ and (P1). Keep
$W=\ell\oplus\ell^-$, $I=I_\Delta=\C f$, $I^*=I_\nabla$, but write the points of $I^*$ as $\xi e_0$ with
$e_0=(1,-1)=2\jmath e$, so that $\varphi_0(\xi e_0)=e^{-4\pi s|\xi|^2}$ and, from $\langle yf,\xi e_0\rangle=
2\aubar y\xi$, $[\xi e_0,yf]=e^{8\pi is\operatorname{Im}(\aubar y\xi)}$. For real $t$ let $b_t\colon\xi e_0
\mapsto it\xi f$ (admissible: $it$ is imaginary). Then $\beta(\xi e_0,b_t\xi e_0)=e^{-4\pi ist|\xi|^2}$,
so $\omega(n(b_t))$ is multiplication by $e^{4\pi ist|\xi|^2}$ (Lemma~\ref{pwf:siegel}). The element
$w_0=(1,-1)\in\auU(W)$ exchanges $e_0$ and $f$, and $n^-_t:=w_0n(b_t)w_0^{-1}$ is the identity on $I^*$
and maps $yf\mapsto yf+ity\,e_0$. The operator $(T\varphi)(\xi)=\int_\C\varphi(\xi')[\xi'e_0,w_0^{-1}\xi
e_0]\,d\xi'=\int_\C\varphi(\xi')e^{8\pi is\operatorname{Im}(\aubar\xi\xi')}\,d\xi'$ is, up to a positive
factor, unitary and an intertwiner along $w_0$ (translation by $\xi_0$ becomes multiplication by
$[\xi_0e_0,w_0^{-1}\xi e_0]^{-1}=[\xi e_0,w_0\xi_0e_0]$, and multiplication by $[\cdot,y_0f]$ becomes
translation by $w_0y_0f$, both by bimultiplicativity of $[\,,]$ and $\pwfA(w_0x,x')=\pwfA(x,w_0^{-1}x')$).
By (P1), $\omega(w_0)=\lambda T$, so $\omega(n^-_t)=T\,\omega(n(b_t))\,T^{-1}$ exactly. With
$\int_\C e^{-4\pi a|\xi'|^2+8\pi is\operatorname{Im}(\aubar\xi\xi')}d\xi'=(4a)^{-1}e^{-4\pi s^2|\xi|^2/a}$ for
$\operatorname{Re}a>0$ (a two-dimensional Gaussian integral, holomorphic in $a$), the same with
$-8\pi is$ in place of $8\pi is$, and Fourier inversion ($T^{-1}$ is $16s^2$ times the operator with
the conjugate kernel):
\[
T^{-1}\varphi_0=4s\,\varphi_0,\qquad \omega(n^-_t)\varphi_0=T\bigl(4s\,e^{-4\pi s(1-it)|\xi|^2}\bigr)
=(1-it)^{-1}\exp\bigl(-4\pi s|\xi|^2/(1-it)\bigr).
\]
Now let $g=(z_1,z_2)$ with $z_1\ne-z_2$, $D=\frac{z_1+z_2}2$, $B=\frac{z_1-z_2}2$, $it=B/D$ (imaginary,
as $|z_1|=|z_2|=1$) and $a=\aubar D^{-1}$. In the basis $(f,e_0)$, $g=\bigl(\begin{smallmatrix}D&B\\B&D
\end{smallmatrix}\bigr)$, and, using $D^2-B^2=z_1z_2=D/\aubar D$,
\[
n(b_t)\,m(a)\,n^-_t=\begin{pmatrix}1&it\\0&1\end{pmatrix}\begin{pmatrix}\aubar D^{-1}&0\\0&D\end{pmatrix}
\begin{pmatrix}1&0\\it&1\end{pmatrix}=\begin{pmatrix}\aubar D^{-1}+B^2/D&B\\B&D\end{pmatrix}=g .
\]
By Lemma~\ref{pwf:siegel}, $\omega(m(a))\varphi(\xi)=\upsilon_\tau(a)^{-1}|a|\varphi(\aubar a\xi)$, and
$\upsilon_\tau(a)^{-1}|a|=\aubar a=D^{-1}$. Hence
\[
\omega(g)\varphi_0(\xi)=D^{-1}(1-it)^{-1}\exp\Bigl(-4\pi s|\xi|^2\Bigl(\frac{|a|^2}{1-it}-it\Bigr)\Bigr).
\]
Here $1-it=(D-B)/D=z_2/D$, so $D^{-1}(1-it)^{-1}=z_2^{-1}$; and with $w=z_1\aubar z_2$,
$|a|^2/(1-it)=1/(\aubar Dz_2)=2/(1+\aubar w)=2w/(1+w)$ and $it=(w-1)/(w+1)$, so the bracket is $1$.
Thus $\omega(g)\varphi_0=z_2^{-1}\varphi_0$, as in the first proof.
\end{proof}

\begin{remark}\label{pwf:delta-open}
The result depends on the signs fixed in \S\ref{ps:ss-weil}: with $\psi_{\Q,\infty}(t)=e^{+2\pi it}$
(or with $\tau(\jmath)=-i/\sqrt d$) the same computation gives $c=z_1^{-1}$, i.e.\ $e_+=-1$, $e_-=0$.
For the conventions of the paper, $\Delta(\gamma)=\aubar{\det\gamma^-}$ as used in {\cite[\S3.3.2]{P54}}
and in \S\ref{ps:ss-weil} (where a positive definite space carries no scalar on the compact group).
\end{remark}



\section{The archimedean theory of the lifts}\label{A1:g2}
\providecommand{\auU}{\mathrm U}
\providecommand{\auPl}{\Phi}
\providecommand{\auHom}{\operatorname{Hom}}
\providecommand{\psrho}{\varrho}
\providecommand{\pzF}{\mathfrak F}
\providecommand{\parsgn}{\operatorname{sgn}}

\subsection{The archimedean statements (A1)--(A3)}\label{par:sec}

This section treats the three statements (A1)--(A3) of \S\ref{pz:ss-types}. \emph{Dependencies.}
(A2) is proved from the decay bound \eqref{par:eq-decay} below (Corollary~\ref{par:A2}). (A1) is
proved from the same bound (Theorem~\ref{par:A1}); the bound rests on the explicit formula
{\cite[(59)]{P51}}, which is proved in Theorem~\ref{pco:FJ} (Remark~\ref{par:FJ}). (A3) is proved
(Theorem~\ref{par:A3}) modulo the published theorems (D1)--(D3) of \S\ref{par:ss-rigid} and the
normalisation (V) of the vacuum character (Lemma~\ref{par:vac}), which is proved in
Theorem~\ref{pco:V} from the local properties (W1$_\infty$), (W2$_\infty$), (M), (S) of the
archimedean Weil representation.

\emph{Conventions.} We fix $\tau\in\auPl$ and $0\le j\le m$. $G=G_{j,\tau}$ is identified with
$\auU(p,q)=\{g:g^*Jg=J\}$, $J=\operatorname{diag}(\varepsilon_1,\dots,\varepsilon_n)$, $n=N_j$, and
$K=G\cap\auU(n)$; $T\subset K$ is the diagonal torus and $\mathfrak g_\C=\mathfrak{gl}_n(\C)$ with
matrix units $E_{ab}$. For a parameter $\lambda\in(\tfrac{n-1}2+\Z)^n$ with distinct entries,
the positive roots are the $e_a-e_b$ with $\lambda_a>\lambda_b$; $\rho$, $\rho_{\mathrm c}$,
$\rho_{\mathrm n}$ are as in {\cite[\S5.1]{P51}}, $\Lambda^K(\lambda)=\lambda+\rho_{\mathrm n}-
\rho_{\mathrm c}$, and $\pi_\lambda$ is the discrete series with Harish-Chandra parameter $\lambda$.
We write $\Omega=\sum_{a,b}E_{ab}E_{ba}$; it acts on a representation with infinitesimal character
$\lambda$ by $\kappa_\lambda=|\lambda|^2-|\rho|^2$. The \emph{gap} of $\lambda$ is
$\min\{|\lambda_a-\lambda_b|:\varepsilon_a\ne\varepsilon_b\}$. By \eqref{pz:eq-gap} and
{\cite[\S14.1]{P46}} ($r\ge x^5$, $N_j\le40x^6$) the parameters $\Lambda^{(j)}_\tau$ have gap at
least $\tfrac12r^8\ge\tfrac12x^{40}$, which exceeds $500(N_j+N_{j+1})^3+10N_j^5+x^{25}$.

\emph{The coefficient.} Inner products are conjugate-linear in the first variable, as in
\S\ref{ps:ss-weil}. We read the coefficient of \S\ref{pz:ss-types} as
\begin{equation}\label{par:eq-coef}
c_j(x)=\langle v_j,\pi_j(x)v_j\rangle,\qquad e_j=d_j\,\overline{c_j}.
\end{equation}
With this reading, Schur's orthogonality relations give, for every unitary representation
$\sigma$, that $\sigma(e_j)=\int e_j(x)\sigma(x)\,dx$ is the orthogonal projector onto the vectors
of type $\pi_j$ that correspond to $v_j$. This is the property used in \S\ref{pz:ss-types}. (With
the other reading, $\sigma(e_j)$ projects onto the type $\overline{\pi_j}$, and the rule
\eqref{pz:eq-param} would no longer agree with Lemma~\ref{par:vac}; see Remark~\ref{par:rem-match}.)

\subsubsection{The decay bound}\label{par:ss-decay}
We say that $\lambda$ \emph{satisfies} \textup{(Dec$_c$)} if the coefficient
$c_\lambda(g)=\langle v,\pi_\lambda(g)v\rangle$, with $v$ a unit highest weight vector of the
lowest $K$-type, satisfies
\begin{equation}\label{par:eq-decay}
|c_\lambda(g)|\le(2/\|g\|)^{c}\qquad(g\in G).
\end{equation}

\begin{remark}\label{par:FJ}
{\cite[Theorem~5.6]{P51}} proves \eqref{par:eq-decay} for every integer $c$ up to the gap, for any
pair $(\lambda,\varepsilon)$, \emph{assuming} the integral formula {\cite[(59)]{P51}} for $c_\lambda$.
The proof of {\cite[Theorem~5.6]{P51}} ({\cite[Lemmas~5.4 and~5.5]{P51}}) uses only
the formula and the gap, not the signature of $Y_0$: {\cite[Lemma~5.4]{P51}} uses only $g^*Jg=J$ and that the maximal compact subgroup commutes with $J$; {\cite[Lemma~5.5]{P51}} uses only Fischer's inequality, the formula and $c\le h$; and the displayed estimate in the proof of {\cite[Theorem~5.6]{P51}} uses only {\cite[(59)]{P51}} and $\det a(u^{-1}zu)=\det a(z)$ for $u$ in the maximal compact subgroup. Hence {\cite[Theorem~5.6]{P51}} and its proof hold verbatim with $N$ replaced by $n=N_j$ and $J$ by the diagonal matrix of signs of $Y_{j,\tau}$ (of any signature), for every regular $\lambda\in(\frac{n-1}2+\Z)^n$ with gap $h$ for which {\cite[(59)]{P51}} holds.
Formula {\cite[(59)]{P51}} is Theorem~\ref{pco:FJ}, proved for every regular $\lambda\in
(\frac{n-1}2+\Z)^n$ with gap $h\ge n^2+1$, in particular for every $\Lambda^{(j)}_\tau$: steps (i)
(Lemma~\ref{pco:K}, Proposition~\ref{pco:casimir}), (ii) (Lemma~\ref{pco:decay}) and (iii)
(Proposition~\ref{pco:rigid} and the proof of Theorem~\ref{pco:FJ}) are carried out there. Hence every
$\Lambda^{(j)}_\tau$ satisfies (Dec$_c$) for every integer $c$ up to the gap.
\end{remark}

\subsubsection{(A1) and (A2)}\label{par:ss-A12}

\begin{theorem}[(A1)]\label{par:A1}
Assume that every $\Lambda^{(j)}_\tau$ satisfies \textup{(Dec$_c$)} with $c=B+N_j^2+1$, $B=x^{25}$.
Then for $0\le j\le m$ the function $c_j=\bigotimes_\tau c_{j,\tau}$ is smooth, left and right
$K_{j,\infty}$-finite, an eigenfunction of the centre of the enveloping algebra, and
$\int_{G_{j,\infty}}|\mathsf Xc_j(x)|\,\|x\|^B\,dx<\infty$ for every left- or right-invariant
differential operator $\mathsf X$.
\end{theorem}
\begin{proof}
Smoothness, $K$-finiteness and the eigenfunction property hold for every $K$-finite coefficient
of an irreducible unitary representation: $K$-finite vectors are smooth (indeed analytic)
\cite[Proposition~8.5, p.~208]{Knapp}, and the centre acts on the smooth vectors of an irreducible representation
by scalars. Fix $\tau$ and write $c=c_{j,\tau}$, $e=d\bar c$. By {\cite[Lemma~4.7(b)]{P51}}, whose proof
uses only that $c$ is a $K$-finite coefficient of an irreducible unitary representation, there
is $\alpha\in C_c^\infty(G)$ with $e=\alpha*e=e*\alpha^*$; so $\mathsf Xe=(\mathsf X\alpha)*e$ for
right-invariant and $\mathsf Xe=e*(\mathsf X\alpha^*)$ for left-invariant $\mathsf X$. With
$\|gh\|\le\|g\|\|h\|$ and \eqref{par:eq-decay},
$|\mathsf Xe(g)|\le C_{\mathsf X}\|g\|^{-c}\int\|h\|^c|\mathsf X\alpha(h)|dh$. By
{\cite[Lemma~4.7(c)]{P51}}, $\int_G\|g\|^{B-c}dg<\infty$ since $c-B>N_j^2$. The norm $\|x\|$ of
$G_{j,\infty}$ is the product of the local norms and $c_j$ is a product, so Fubini gives the
statement on $G_{j,\infty}$ (an invariant operator on the product is a sum of products of local
ones).
\end{proof}

\begin{corollary}[(A2)]\label{par:A2}
Assume \textup{(Dec$_c$)} with $c=5N_j^3$ for every $\Lambda^{(j)}_\tau$. Let $\mathsf P=\prod_\tau
\mathsf P_\tau$ be a parabolic subgroup of $G_{j,\infty}$ with $\mathsf P_\tau$ proper for at least
one $\tau$, and $\mathsf N$ its unipotent radical. Then $\int_{\mathsf N}|c_j(anb)|dn<\infty$ and
$\int_{\mathsf N}c_j(anb)\,dn=0$ for all $a,b$.
\end{corollary}
\begin{proof}
The proof of {\cite[Lemma~4.8]{P51}} uses, for the function $e_\tau$, only: (a) the bound
$|e_\tau(g)|\le C\|g\|^{-5n^3}$; (b) the derivative bounds of {\cite[Lemma~4.7(b)]{P51}}; (c) that
$e_\tau$ is an eigenfunction of the centre. Here (a) is \eqref{par:eq-decay}, (b) follows from it as
in the proof of Theorem~\ref{par:A1}, and (c) holds as there. The signature and the gap condition
{\cite[(41)]{P51}} enter {\cite[Lemma~4.8]{P51}} only through (a). So for each $\tau$ with $\mathsf
P_\tau$ proper, $\int_{\mathsf N_\tau}|c_{j,\tau}(a_\tau nb_\tau)|dn<\infty$ and the integral vanishes;
for $\mathsf P_\tau=G_{j,\tau}$, $\mathsf N_\tau=1$. Since $c_j$ is a product and $\mathsf
N=\prod_\tau\mathsf N_\tau$, Fubini gives absolute convergence, and the integral is the product of
the local integrals, one of which is $0$.
\end{proof}
This is stronger than (A2), which assumes every $\mathsf P_\tau$ proper. The hypotheses of
Theorem~\ref{par:A1} and Corollary~\ref{par:A2} hold as soon as {\cite[(59)]{P51}} holds, since the gap
exceeds both exponents.

\subsubsection{Rigidity of the lowest $K$-type}\label{par:ss-rigid}
We use three published results, for $G=\auU(p,q)$ and a regular parameter $\lambda$ as above.
\begin{enumerate}
\item[(D1)] \emph{(Dirac inequality; Parthasarathy's formula for the square of the Dirac operator \cite[Proposition~3.1]{Parthasarathy72}, there in the geometric setting of spinors on $G/K$, in the algebraic form of \cite[Proposition~1.1 and \S2]{HuangPandzic}.)} Let $X$ be an
irreducible unitary $(\mathfrak g,K)$-module on which $\Omega$ acts by $\omega_X$, $S$ the spin module
of $\mathfrak p$, $\mu$ a $K$-type of $X$ and $\tau_K$ a $K$-type of $\mu\otimes S$. Then
$\omega_X+|\rho|^2\le|\tau_K+\rho_{\mathrm c}|^2$; if equality holds, the Dirac operator kills the
$\tau_K$-isotypic part of $X\otimes S$, so the Dirac cohomology of $X$ contains $\tau_K$. Moreover, for
$\mu=\Lambda^K(\lambda)$, the $K$-type $\tau_K=\mu-\rho_{\mathrm n}=\lambda-\rho_{\mathrm c}$ occurs in
$\mu\otimes S$ (same reference).
\item[(D2)] \emph{(Huang--Pand\v zi\'c \cite[Theorem~2.3]{HuangPandzic}.)} If
the Dirac cohomology of an irreducible $(\mathfrak g,K)$-module with infinitesimal character
$\chi$ contains $\tau_K$, then $\chi$ is $W$-conjugate to $\tau_K+\rho_{\mathrm c}$.
\item[(D3)] \emph{(Lowest $K$-type and cohomology.)} Let $\Psi=\{\alpha:\langle\alpha,\lambda\rangle>0\}$; $K$-types are labelled by highest weights for $\Psi_{\mathrm c}=\Psi\cap\Delta(\mathfrak k)$.
(a) \emph{(Schmid \cite[Theorem~(1.3)]{Schmid75}.)} Up to infinitesimal equivalence, $\pi_\lambda$ is the only irreducible representation that contains the $K$-type $\Lambda^K(\lambda)$ and contains no $K$-type of highest weight $\Lambda^K(\lambda)-A$, $A$ a non-empty sum of elements of $\Psi$.
(b) \emph{(Schmid \cite[Corollary~(3.9), Corollary~(4.19) and the proof of Theorem~(1.3) on p.~560]{Schmid75}.)} $\Lambda^K(\lambda)$ occurs in $\pi_\lambda$ with multiplicity one.
(c) \emph{(Vogan--Zuckerman \cite[Proposition~5.4, Theorems~5.3 and~5.6]{VoganZuckerman}, for a connected semisimple group $G_1$ with finite centre, maximal compact $K_1$, $\mathfrak g_1=\mathfrak k_1+\mathfrak p_1$.)} Let $X$ be the Harish-Chandra module of an irreducible unitary representation of $G_1$ and $F$ an irreducible finite-dimensional representation of $G_1$ such that the Casimir operator acts by the same scalar on $X$ and on $F$. Then $H^*(\mathfrak g_1,K_1;X\otimes F)\cong\operatorname{Hom}_{K_1}(\wedge^*\mathfrak p_1,X\otimes F)$; if this is non-zero, there are a $\theta$-stable parabolic subalgebra $\mathfrak q=\mathfrak l+\mathfrak u$, defined by an element $x_{\mathfrak q}\in i\mathfrak t_{1}$ ($\mathfrak t_1$ a Cartan subalgebra of $\mathfrak k_1$; $\mathfrak u$, $\mathfrak l$ the sums of the positive, zero eigenspaces of $\operatorname{ad}x_{\mathfrak q}$), and a weight $\lambda_{\mathfrak q}$, the differential of a unitary character of the centraliser $L$ of $x_{\mathfrak q}$, with $\langle\alpha,\lambda_{\mathfrak q}\rangle\ge0$ for $\alpha\in\Delta(\mathfrak u)$, such that, for any positive system $\Delta^+(\mathfrak l\cap\mathfrak k_1)$, with $\Delta^+(\mathfrak k_1)=\Delta^+(\mathfrak l\cap\mathfrak k_1)\cup\Delta(\mathfrak u\cap\mathfrak k_1)$ used to label $K_1$-types: every $K_1$-type of $X$ has highest weight $\lambda_{\mathfrak q}+2\rho(\mathfrak u\cap\mathfrak p_1)+\sum_{\beta\in\Delta(\mathfrak u\cap\mathfrak p_1)}n_\beta\beta$, $n_\beta\in\Z_{\ge0}$; and the centre of $U(\mathfrak g_1)$ acts on $X$ by the character with Harish-Chandra parameter $\lambda_{\mathfrak q}+\rho'$, $\rho'$ the half-sum of $\Delta^+(\mathfrak l)\cup\Delta(\mathfrak u)$ for any positive system $\Delta^+(\mathfrak l)\supseteq\Delta^+(\mathfrak l\cap\mathfrak k_1)$ (here $\operatorname{rk}\mathfrak k_1=\operatorname{rk}\mathfrak g_1$, so $\mathfrak t_1$ is a Cartan subalgebra of $\mathfrak g_1$).
\end{enumerate}
We also use that $G$ is of type I, so every unitary representation on a separable Hilbert space
is a direct integral of irreducible ones \cite{Dixmier}, and that irreducible unitary
representations are admissible with smooth $K$-finite vectors \cite[Theorem~8.1 and Proposition~8.5]{Knapp}. (D1), (D2)
are stated for reductive groups with these normalisations in \cite[Proposition~1.1 and \S2]{HuangPandzic}; the centre of
$\auU(p,q)$ acts by a character and does not affect them.

\emph{On (D3).} Schmid's paper is written for a connected semisimple linear group with $\operatorname{rk}K=\operatorname{rk}G$, and covers \emph{all} discrete series, not only the holomorphic ones; it applies to $\auU(p,q)=Z\cdot\mathrm{SU}(p,q)$ ($Z$ the compact centre, acting by a character) through $\mathrm{SU}(p,q)$. The discrete series used in this paper are in general neither holomorphic nor antiholomorphic: in \eqref{pz:eq-param} the signs of the entries $-\kappa(\tau)-i\Delta$ alternate, so there are positive noncompact roots of both kinds $e_+-e_-$ and $e_--e_+$; hence no elementary (generalised Verma module) argument replaces Schmid's theorem here. What \cite{Schmid75} proves is: (i) $\pi_\lambda$ contains $\Lambda^K(\lambda)$ exactly once, and every $K$-type of $\pi_\lambda$ has highest weight $\Lambda^K(\lambda)+\sum n_\beta\beta$, $n_\beta\in\Z_{\ge0}$, over the noncompact roots $\beta$ with $\langle\beta,\lambda\rangle>0$ (Lemmas~(2.22), (3.8), Corollary~(3.9) and, for the identification of $\pi_\lambda$ with the irreducible quotient of the universal module, Corollary~(4.19)); (ii) is (D3)(a). Schmid's characterisation is by $K$-types, not by the infinitesimal character; the form used below (an irreducible unitary module with infinitesimal character $\lambda$ containing $\Lambda^K(\lambda)$ is $\pi_\lambda$) is Lemma~\ref{par:D3}, deduced from (D3)(a) and (D3)(c). In \cite{VoganZuckerman}, $G$ is connected semisimple with finite centre; Theorems~5.3 and~5.6 and Proposition~5.4 are in \S5, where the authors indicate the proofs as extensions of those given in full for trivial coefficients (Theorems~2.5 and~4.1). The positive system $\Delta^+(\mathfrak g)=\Delta^+(\mathfrak l)\cup\Delta(\mathfrak u)$ entering the infinitesimal character is the one fixed before (2.6) there, and the Harish-Chandra parametrisation of infinitesimal characters is (2.9) there: the characters with parameters $\nu,\nu'$ agree if and only if $\nu'\in W\nu$.

\begin{lemma}\label{par:D3}
Let $X$ be the Harish-Chandra module of an irreducible unitary representation of $G=\auU(p,q)$ with infinitesimal character $\lambda$ \textup{(}regular, $\lambda\in(\frac{n-1}2+\Z)^n$\textup{)} that contains the $K$-type $\Lambda^K(\lambda)$. Then every $K$-type of $X$ has highest weight $\Lambda^K(\lambda)+\sum n_\beta\beta$ with $n_\beta\in\Z_{\ge0}$ and $\beta$ running over noncompact roots in $\Psi$; in particular $X$ contains no $K$-type $\Lambda^K(\lambda)-A$ as in \textup{(D3)(a)}, and $X$ is the Harish-Chandra module of $\pi_\lambda$.
\end{lemma}
\begin{proof}
Write $\Lambda=\Lambda^K(\lambda)$. If $p=0$ or $q=0$, $G$ is compact, there are no noncompact roots, $\Lambda=\lambda-\rho$, and $X$ is the irreducible representation with highest weight $\Lambda$, which is $\pi_\lambda$. Let $p,q\ge1$.

\emph{Reduction to $G_1=\mathrm{SU}(p,q)$.} $G=Z\cdot G_1$ with $Z$ the centre, which acts on $X$ by a character (Schur's lemma); so $X$ is the Harish-Chandra module of an irreducible unitary representation of $G_1$, for $K_1=K\cap G_1$, and $K$-types of $G$ restrict to irreducible $K_1$-types. We work with $\mathfrak t_1=\mathfrak t\cap\mathfrak g_1$ and restrict all weights to it; roots, $\Psi$, $\rho$, $\rho_{\mathrm c}$, $\rho_{\mathrm n}$ live on $\mathfrak t_1$ already, and the conclusion for $G$ follows from the one for $G_1$, since a $K$-type is determined by its restriction to $K_1$ and the central character, which is common to all $K$-types of $X$.

\emph{Non-vanishing of cohomology.} Since $\lambda\in(\frac{n-1}2+\Z)^n$ and $\rho\in(\frac{n-1}2+\Z)^n$, $\lambda-\rho\in\Z^n$; for $\alpha\in\Psi$ simple, $\langle\lambda-\rho,\alpha^\vee\rangle=\langle\lambda,\alpha^\vee\rangle-1\ge0$, because $\langle\lambda,\alpha^\vee\rangle$ is a positive integer. Let $F^*$ be the irreducible finite-dimensional representation of $G$ (restricted from $\mathrm{GL}_n(\C)$) with $\Psi$-highest weight $\lambda-\rho$, and $F$ its dual. As $\mathfrak g_1$-modules ($\mathfrak g_\C=\mathfrak g_{1,\C}\oplus\C\cdot1$, and $Z(U(\mathfrak g_\C))=Z(U(\mathfrak g_{1,\C}))\otimes\C[1]$), $X$ and $F^*$ both have infinitesimal character with parameter $\lambda|_{\mathfrak t_1}$, so the Casimir operator of $\mathfrak g_1$ acts on them by the same scalar; it acts on $F$ as on $F^*$ (its scalar on the module with highest weight $\nu$ is $|\nu+\rho|^2-|\rho|^2$, and $F$ has highest weight $-w_0(\lambda-\rho)$, $w_0$ the longest element, with $|-w_0(\lambda-\rho)+\rho|=|\lambda-\rho+\rho|$ as $-w_0\rho=\rho$). So the Casimir operator acts on $F$ as on $X$. Let $\mathfrak u_{\mathrm n}=\bigoplus_{\beta\in\Psi_{\mathrm n}}\mathfrak g_\beta\subset\mathfrak p$, $R=\dim\mathfrak u_{\mathrm n}$. Since $\Psi$ is closed under addition of roots, $[\mathfrak g_\alpha,\mathfrak u_{\mathrm n}]\subset\mathfrak u_{\mathrm n}$ for $\alpha\in\Psi_{\mathrm c}$, so the line $\wedge^R\mathfrak u_{\mathrm n}\subset\wedge^R\mathfrak p$ has weight $2\rho_{\mathrm n}$ and is killed by every $\mathfrak g_\alpha$, $\alpha\in\Psi_{\mathrm c}$. With $v$ a highest weight vector of $F^*$, $e=(\wedge^R\mathfrak u_{\mathrm n})\otimes v$ is therefore a $\Psi_{\mathrm c}$-highest weight vector of weight $2\rho_{\mathrm n}+\lambda-\rho=\Lambda$ in $\wedge^R\mathfrak p\otimes F^*$, and generates a $K$-type $\Lambda$. As $X$ contains $\Lambda$, $\operatorname{Hom}_{K_1}(\wedge^R\mathfrak p\otimes F^*,X)\cong\operatorname{Hom}_{K_1}(\wedge^R\mathfrak p,X\otimes F)\ne0$. By (D3)(c), $H^R(\mathfrak g_1,K_1;X\otimes F)\ne0$, and we obtain $\mathfrak q=\mathfrak l+\mathfrak u$, $x_{\mathfrak q}$ and $\lambda_{\mathfrak q}$ as there.

\emph{Normalising $\mathfrak q$.} All Cartan subalgebras of $\mathfrak k_1$ are $K_1$-conjugate, and conjugating $(\mathfrak q,x_{\mathfrak q},\lambda_{\mathfrak q},\Delta^+(\mathfrak l\cap\mathfrak k_1))$ by an element $k\in K_1$ transforms the conclusions of (D3)(c) into the same conclusions for the conjugated data (a $K_1$-type with highest weight $\delta$ for a positive system $P$ has highest weight $w\delta$ for $wP$, $w$ the image of $k$ in the Weyl group when $k$ normalises $\mathfrak t_1$). So we may assume $x_{\mathfrak q}\in i\mathfrak t_1$ and, choosing $\Delta^+(\mathfrak l\cap\mathfrak k_1)$ and then conjugating by $N_{K_1}(\mathfrak t_1)$, that $\Delta^+(\mathfrak l\cap\mathfrak k_1)\cup\Delta(\mathfrak u\cap\mathfrak k_1)=\Psi_{\mathrm c}$. Then $\alpha(x_{\mathfrak q})>0$ for $\alpha\in\Delta(\mathfrak u)$, $\alpha(x_{\mathfrak q})=0$ for $\alpha\in\Delta(\mathfrak l)$, and $\Psi_{\mathrm c}\cap\Delta(\mathfrak u)=\Delta(\mathfrak u\cap\mathfrak k_1)$, $\Psi_{\mathrm c}\cap\Delta(\mathfrak l)=\Delta^+(\mathfrak l\cap\mathfrak k_1)$.

\emph{The comparison.} Fix $\Delta^+(\mathfrak l)\supseteq\Delta^+(\mathfrak l\cap\mathfrak k_1)$, put $P'=\Delta^+(\mathfrak l)\cup\Delta(\mathfrak u)$, $\rho'$ its half-sum, and $\lambda'=\lambda_{\mathfrak q}+\rho'$. For $\alpha$ simple in $P'$, $\langle\lambda_{\mathfrak q},\alpha\rangle\ge0$ (it is $0$ for $\alpha\in\Delta(\mathfrak l)$, since $\lambda_{\mathfrak q}$ is the differential of a character of $L$, and $\ge0$ for $\alpha\in\Delta(\mathfrak u)$), and $\langle\rho',\alpha^\vee\rangle=1$; so $\lambda'$ is $P'$-dominant, and $x_{\mathfrak q}$ is $P'$-dominant ($\alpha(x_{\mathfrak q})\ge0$ on $P'$). By (D3)(c) and (2.9) of \cite{VoganZuckerman}, $\lambda=w\lambda'$ for some $w$ in the Weyl group $W$. For dominant $\lambda'$, $\lambda'-w\lambda'$ is a sum of elements of $P'$ with non-negative integer coefficients (induction on the length: if $w=s_\alpha w_1$, $\alpha$ simple, $\ell(w)=\ell(w_1)+1$, then $w_1^{-1}\alpha\in P'$ and $\lambda'-w\lambda'=(\lambda'-w_1\lambda')+\langle\lambda',w_1^{-1}\alpha^\vee\rangle\alpha$). Hence
\begin{equation}\label{par:eq-D3a}
\lambda(x_{\mathfrak q})\le\lambda'(x_{\mathfrak q})=\lambda_{\mathfrak q}(x_{\mathfrak q})+\rho(\mathfrak u)(x_{\mathfrak q}),
\end{equation}
as $\rho'-\rho(\mathfrak u)$ is a combination of roots of $\mathfrak l$, which vanish at $x_{\mathfrak q}$. On the other hand $\Lambda$ is a $K_1$-type of $X$, so by (D3)(c), $\Lambda=\lambda_{\mathfrak q}+2\rho(\mathfrak u\cap\mathfrak p_1)+B_{\mathfrak q}$ with $B=\sum_{\beta\in\Delta(\mathfrak u\cap\mathfrak p_1)}n_\beta\beta$, $n_\beta\ge0$. Now $\lambda=\Lambda+\rho_{\mathrm c}-\rho_{\mathrm n}$, and: $2\rho_{\mathrm c}(x_{\mathfrak q})=\sum_{\alpha\in\Psi_{\mathrm c}}\alpha(x_{\mathfrak q})=2\rho(\mathfrak u\cap\mathfrak k_1)(x_{\mathfrak q})$ by the normalisation; for each noncompact $\beta$ with $\beta(x_{\mathfrak q})>0$ exactly one of $\pm\beta$ lies in $\Psi_{\mathrm n}$, and noncompact roots with $\beta(x_{\mathfrak q})=0$ contribute nothing, so $2\rho_{\mathrm n}(x_{\mathfrak q})\le\sum_{\beta\in\Delta(\mathfrak u\cap\mathfrak p_1)}\beta(x_{\mathfrak q})=2\rho(\mathfrak u\cap\mathfrak p_1)(x_{\mathfrak q})$, with equality if and only if $\Delta(\mathfrak u\cap\mathfrak p_1)\subseteq\Psi$. Therefore
\[
\lambda(x_{\mathfrak q})\ge\lambda_{\mathfrak q}(x_{\mathfrak q})+2\rho(\mathfrak u\cap\mathfrak p_1)(x_{\mathfrak q})+B_{\mathfrak q}(x_{\mathfrak q})+\rho(\mathfrak u\cap\mathfrak k_1)(x_{\mathfrak q})-\rho(\mathfrak u\cap\mathfrak p_1)(x_{\mathfrak q})=\lambda_{\mathfrak q}(x_{\mathfrak q})+\rho(\mathfrak u)(x_{\mathfrak q})+B_{\mathfrak q}(x_{\mathfrak q}).
\]
With \eqref{par:eq-D3a}, $B_{\mathfrak q}(x_{\mathfrak q})\le0$; since $\beta(x_{\mathfrak q})>0$ for $\beta\in\Delta(\mathfrak u\cap\mathfrak p_1)$, all $n_\beta=0$, so $B_{\mathfrak q}=0$, and equality holds throughout, so $\Delta(\mathfrak u\cap\mathfrak p_1)\subseteq\Psi$. Thus $\Lambda=\lambda_{\mathfrak q}+2\rho(\mathfrak u\cap\mathfrak p_1)$, and by (D3)(c) every $K_1$-type of $X$ has highest weight $\Lambda+\sum n_\beta\beta$ with $\beta\in\Delta(\mathfrak u\cap\mathfrak p_1)\subseteq\Psi_{\mathrm n}$, $n_\beta\ge0$. This is the first assertion.

If $X$ contained a $K$-type $\Lambda-A$ with $A$ a non-empty sum of elements of $\Psi$, then $-A=\sum n_\beta\beta$ with $\beta\in\Psi$, $n_\beta\ge0$; pairing with $\lambda$, which is positive on $\Psi$, gives $\langle\lambda,-A\rangle<0\le\langle\lambda,\sum n_\beta\beta\rangle$, a contradiction. By (D3)(a) (an irreducible unitary representation is quasi-simple, and infinitesimal equivalence means isomorphic Harish-Chandra modules), $X$ is the Harish-Chandra module of $\pi_\lambda$.
\end{proof}

\begin{lemma}[Rigidity]\label{par:rigid}
Let $\sigma$ be a unitary representation of $G$ on a separable Hilbert space and $y$ a smooth
vector with: $\sigma(t)y=t^{\Lambda^K(\lambda)}y$ for $t\in T$; $d\sigma(E_{ab})y=0$ for
$\varepsilon_a=\varepsilon_b$, $\lambda_a>\lambda_b$; $d\sigma(\Omega)y=\kappa_\lambda y$. Put
$e_\lambda=d_\lambda\overline{c_\lambda}$. Then $\langle y,\sigma(g)y\rangle=\|y\|^2c_\lambda(g)$ and
$\sigma(e_\lambda)y=y$, provided $e_\lambda\in L^1(G)$.
\end{lemma}
\begin{proof}
Let $H_y$ be the closed span of $\sigma(G)y$ and $H_y=\int^\oplus H_\zeta\,d\mu(\zeta)$ a
decomposition into irreducibles, $y=\int y_\zeta$. The three conditions on $y$ are equivalent to
countably many identities $\sigma(f)y=\sigma(f')y$ with $f,f'\in C_c^\infty(G)$ (for $\Omega$: $\sigma(f)
d\sigma(\Omega)y=\sigma(\Omega f)y$, with $\Omega$ acting on $f$ as a bi-invariant operator; for the
others, the projector onto $T$-weight $\Lambda^K$ and integration by parts against $f$). These
identities decompose, so for almost all $\zeta$ they hold for $y_\zeta$. Such a $y_\zeta$ lies in a
$K$-isotypic subspace, which is finite-dimensional, so $y_\zeta$ is smooth and satisfies the three
conditions in the derived action. If $y_\zeta\ne0$, then $H_\zeta$ contains the $K$-type
$\Lambda^K(\lambda)$ and $\Omega$ acts by $\kappa_\lambda$; with $\tau_K=\lambda-\rho_{\mathrm c}$,
$|\tau_K+\rho_{\mathrm c}|^2=|\lambda|^2=\kappa_\lambda+|\rho|^2$, so equality holds in (D1), and (D2)
gives the infinitesimal character $\lambda$. By Lemma~\ref{par:D3}, $H_\zeta\cong\pi_\lambda$, and $y_\zeta$, a
highest weight vector of the $K$-type of multiplicity one ((D3)(b)), is $a(\zeta)U_\zeta v$ for a unitary
$U_\zeta:\pi_\lambda\to H_\zeta$ and a scalar $a(\zeta)$. Hence $\langle y,\sigma(g)y\rangle=
\int|a|^2d\mu\cdot c_\lambda(g)=\|y\|^2c_\lambda(g)$. By Schur's orthogonality,
$\int e_\lambda(h)c_\lambda(g^{-1}h)\,dh=c_\lambda(g^{-1})$, so $\langle\sigma(g)y,\sigma(e_\lambda)y-
y\rangle=0$ for all $g$; since $\sigma(e_\lambda)y\in H_y$, $\sigma(e_\lambda)y=y$.
\end{proof}


\subsubsection{The Fock model at one place}\label{par:ss-fock}
Let $0\le j<m$. At $\tau$ let $R$ (resp.\ $C$) index the orthogonal basis of
$Y_j\otimes_{E,\tau}\C$ (resp.\ $Y_{j+1}$), with signs $\varepsilon_r$ (resp.\ $\varepsilon'_c$),
$n=|R|$, $n'=|C|=n+t$. The polynomial vectors of $\pzF_\tau$ are the polynomials in variables
$z_{rc}$; the variable $z_{rc}$ is the conjugate of the coordinate $x_{rc}$ if
$\varepsilon_r\varepsilon'_c=1$ and the coordinate itself otherwise (the positive complex structure
of \S\ref{pz:ss-types}). Put $\eta_{rc}=-\varepsilon_r\varepsilon'_c$ and define operators
$a^+_{rc}=z_{rc}$, $a^-_{rc}=\partial_{rc}$ if $\eta_{rc}=1$, and $a^+_{rc}=-\partial_{rc}$,
$a^-_{rc}=z_{rc}$ if $\eta_{rc}=-1$; then $[a^-_i,a^+_k]=\delta_{ik}$ and the $a^+$ (resp.\ $a^-$)
commute. For integers or half-integers $\gamma,\gamma'$ put, with $s_{ab}=1$ if
$\varepsilon_a=\varepsilon_b$ and $-1$ otherwise,
\begin{equation}\label{par:eq-osc}
E_{ab}=s_{ab}\sum_{c}a^+_{ac}a^-_{bc}+\gamma\delta_{ab},\qquad
E'_{cd}=-s'_{cd}\sum_ra^-_{rc}a^+_{rd}+\gamma'\delta_{cd}.
\end{equation}
These are commuting actions of $\mathfrak{gl}_n$ and $\mathfrak{gl}_{n'}$ (direct computation from the
commutation relations). For $\varepsilon_a\ne\varepsilon_b$ let $\Delta_{ab}=\sum_{c:\eta_{ac}=-1}
\partial_{ac}\partial_{bc}$ (the degree-lowering part of $E_{ab}$), and $\Delta'_{cd}$ likewise for
columns. On polynomials we use the Fischer inner product, for which $z_i$ is adjoint to $\partial_i$;
it is the inner product of $\pzF$.

\begin{lemma}[The derived action]\label{par:vac}
After conjugation by an element of $T\times T'$, the derived actions of $\omega_\infty(\cdot,1)$ and
$\omega_\infty(1,\cdot)$ at $\tau$ on polynomial vectors are \eqref{par:eq-osc}, with
\[
\gamma=\tau_0+\tfrac{n'}2,\quad\gamma'=\tau'_0+\tfrac n2,\qquad
\tau_0=\kappa_{j+1}(\tau)+\tfrac12[j\text{ odd}],\quad\tau'_0=-\kappa_{j+1}(\tau)-\tfrac12[j\text{ even}],
\]
provided the following holds.
\begin{enumerate}
\item[(V)] At $\tau$, the operator $\omega_{\auHom(Y_j,Y_{j+1})}(g,h)$ of \S\ref{ps:ss-weil} is the
genuine oscillator action of the metaplectic covers (Weyl-ordered quadratic operators) times
$\det(g)^{n'/2}\det(h)^{-n/2}$, where $\upsilon_\tau$ restricted to $\auU(1)$ is $z\mapsto z$.
\end{enumerate}
\end{lemma}
\begin{proof}
$K\times K'$ acts by linear substitutions, so the variable $z_{rc}$ has weight $\eta_{rc}e_r$ for
$T$ and $-\eta_{rc}e'_c$ for $T'$ (compare $F(h^tz\bar g)$ in the definite case of
\S\ref{pk:ss-definite}); hence the compact part is \eqref{par:eq-osc} up to the constants
$\gamma,\gamma'$. For all of $\mathfrak{gl}_n$ and $\mathfrak{gl}_{n'}$, compact and noncompact parts alike, we compute the operator $d\omega^{\mathrm{Mp}}$ of \S\ref{pco:ss-V} directly from its characterisation there: it is the unique element of the span $\mathcal V$ of the symmetrised quadratic operators in the $a^\pm_{rc}$ whose commutator with $d\varrho_\tau(w)$ is $d\varrho_\tau$ of the image of $w$. Let $x_{rc}$ be the coordinates of $w\in W_\tau$ (so that $we_r=\sum_cx_{rc}e'_c$ in the orthogonal bases, scaled as in \S\ref{pk:ss-conv}); the hermitian form is $\sum\varepsilon_r\varepsilon'_c|x_{rc}|^2$, the coordinates for the positive complex structure are $y_{rc}=x_{rc}$ if $\eta_{rc}=-1$ and $y_{rc}=\overline{x_{rc}}$ if $\eta_{rc}=1$, and by \eqref{pk:eq-fock} $d\varrho_\tau(w)=\sum_{r,c}(y_{rc}z_{rc}-\overline{y_{rc}}\partial_{rc})=\sum_{r,c}(\overline{x_{rc}}\,a^+_{rc}-\eta_{rc}x_{rc}\,a^-_{rc})$. Hence
\[
[a^-_{rc},d\varrho_\tau(w)]=\overline{x_{rc}},\qquad[a^+_{rc},d\varrho_\tau(w)]=\eta_{rc}\,x_{rc}.
\]
Let $X\in\mathfrak u(V)$; it acts on $W_\tau$ by $w\mapsto-wX$, that is, $x_{rc}\mapsto-\sum_{r'}X_{r'r}x_{r'c}$, and $\overline{X_{r'r}}=-s_{rr'}X_{rr'}$. So
\[
d\varrho_\tau(-wX)=\sum_{r,r',c}\bigl(s_{rr'}X_{rr'}\,\overline{x_{r'c}}\,a^+_{rc}+\eta_{rc}X_{r'r}\,x_{r'c}\,a^-_{rc}\bigr)
=\sum_{r,r'}s_{rr'}X_{rr'}\Bigl[\sum_ca^+_{rc}a^-_{r'c},\,d\varrho_\tau(w)\Bigr],
\]
where we inserted the two commutators, used $\eta_{rc}\eta_{r'c}=s_{rr'}$, and exchanged $r,r'$ in the second sum. The operator $\sum_{r,r'}s_{rr'}X_{rr'}\sum_c\frac12(a^+_{rc}a^-_{r'c}+a^-_{r'c}a^+_{rc})$ lies in $\mathcal V$ and differs from the one in the bracket by a scalar; therefore it is $d\omega^{\mathrm{Mp}}(X)$, and by complex linearity $d\omega^{\mathrm{Mp}}(E_{ab})=s_{ab}\sum_ca^+_{ac}a^-_{bc}+\frac{n'}2\delta_{ab}$ for all $a,b$. In the same way, $Y\in\mathfrak u(V')$ acts by $x_{rc}\mapsto\sum_{c'}Y_{cc'}x_{rc'}$ with $\overline{Y_{cc'}}=-s'_{cc'}Y_{c'c}$, and one finds $d\varrho_\tau(Yw)=-\sum_{c,c'}s'_{cc'}Y_{c'c}[\sum_ra^+_{rc}a^-_{rc'},d\varrho_\tau(w)]$, so $d\omega^{\mathrm{Mp}}(E'_{cd})=-s'_{cd}\sum_r\frac12(a^-_{rc}a^+_{rd}+a^+_{rd}a^-_{rc})$. Thus no conjugation is needed: the element of $T\times T'$ in the statement may be taken to be $1$. (For the compact parts this agrees with the weights found above.) These are the quadratic parts of \eqref{par:eq-osc}, with the constants $\tfrac{n'}2\delta_{ab}$ and $\tfrac n2\delta_{cd}$: for $\varepsilon_a=\varepsilon_b$ and $\varepsilon'_c=\varepsilon'_d$ the Weyl-ordered operators are $\sum_ca^+_{ac}a^-_{bc}+\tfrac{n'}2\delta_{ab}$ and
$-\sum_ra^-_{rc}a^+_{rd}+\tfrac n2\delta_{cd}$. The twists of $\omega_{\mu_{j+1}}$ contribute
$\det(g)^{\kappa_{j+1}}\det(g)^{-\lfloor n'/2\rfloor}\det(h)^{-\kappa_{j+1}}\det(h)^{\lfloor n/2\rfloor}$
($\mu_{j+1}$ has weight $-\kappa_{j+1}$, $\upsilon_k=\upsilon^{\lfloor k/2\rfloor}$), and (V) contributes
$\det(g)^{n'/2}\det(h)^{-n/2}$. So $\tau_0=\kappa_{j+1}+n'/2-\lfloor n'/2\rfloor$ and
$\tau'_0=-\kappa_{j+1}+\lfloor n/2\rfloor-n/2$; and $N_j=r+jt$ is odd exactly when $j$ is even.
\end{proof}
(V) is proved in Theorem~\ref{pco:V} (in the derived form (V$'$), which is what the proof above
uses), for all pairs, from (W1$_\infty$), (W2$_\infty$) and (M) of \S\ref{pco:ss-V}, and for the
definite pairs (only $j=m-1$, Remark~\ref{pco:rem-V}) also from the compatibility (S) with
orthogonal sums.

\emph{Parameters.} The row representation is $\overline{\pi_j}\cong\pi_\lambda$ with
$\lambda=-\Lambda^{(j)}_\tau$, and the column representation is $\pi_{\lambda'}$ with
$\lambda'=\Lambda^{(j+1)}_\tau$. Put $\xi_r=\lambda_r-\gamma+n'/2$ and $\xi'_c=\lambda'_c-\gamma'+n/2$.
We say that $(\lambda,\lambda')$ \emph{match} if there are a subset $B\subset C$ and a bijection
$r\mapsto c(r)$ of $R$ onto $C\smallsetminus B$ with
\begin{equation}\label{par:eq-match}
\xi'_{c(r)}=-\xi_r,\quad|\xi_r|>\tfrac t2,\quad\eta_{r\,c(r)}=\parsgn\xi_r;\qquad
\{\xi'_c:c\in B\}=\{\tfrac{t-1}2-b:0\le b<t\},\ \varepsilon'\text{ constant on }B.
\end{equation}

\begin{lemma}\label{par:match}
Under Lemma~\ref{par:vac}, $(-\Lambda^{(j)}_\tau,\Lambda^{(j+1)}_\tau)$ match.
\end{lemma}
\begin{proof}
Here $\xi_r=-\Lambda_r-\tau_0$, $\xi'_c=\Lambda'_c-\tau'_0$, and $\xi'_{c(r)}=-\xi_r$ means
$\Lambda'_{c(r)}=\Lambda_r+\tau_0+\tau'_0$, a shift by $-\tfrac12$ for $j$ even and $+\tfrac12$ for $j$
odd. By \eqref{pz:eq-param} the passage $j\to j+1$ shifts the old numbers by exactly this amount and
adds the string $-\kappa_{j+1}+\tfrac{t-1}2-b$ lowered by $\tfrac12[j\text{ even}]$, that is, centred
at $\tau'_0$: this is $B$, with signs $+$. The condition $\eta_{rc(r)}=\parsgn\xi_r$ reads
$\varepsilon'_{c(r)}=\varepsilon_r\parsgn(\Lambda'_{c(r)}-\tau'_0)$: the sign is kept above the centre
and changed below it, which is the rule of \eqref{pz:eq-param}; one checks directly from
\eqref{pz:eq-param} that the numbers $-\kappa-i\Delta$ with $i\ge j+2$ are exactly those below the
centre, and that the resulting signs are those of $\Lambda^{(j+1)}$. Finally the old numbers are at
distance at least $\tfrac12(\Delta-t)>t/2$ from the centre.
\end{proof}

\begin{remark}\label{par:rem-match}
With the other reading of $c_j$, the row parameter would be $\Lambda^{(j)}$ and the column parameter
$-\Lambda^{(j+1)}$; the first condition of \eqref{par:eq-match} would put the new string at $-\tau'_0$,
near $+\kappa_{j+1}$, contradicting \eqref{pz:eq-param}. This is why \eqref{par:eq-coef} is the
reading that makes \S\ref{pz:ss-types} consistent.
\end{remark}

\subsubsection{The joint harmonic}\label{par:ss-harm}
In this section $(\lambda,\lambda')$ is any matching pair of parameters with distinct entries.
Write $\Lambda=\Lambda^K(\lambda)$, $\Lambda'=\Lambda^K(\lambda')$. The vacuum $1$ has $E_{aa}$-weight
$v_a=\gamma-\#\{c:\eta_{ac}=-1\}$ and $E'_{cc}$-weight $v'_c=\gamma'-\#\{r:\eta_{rc}=1\}$. Put
$w_r=\Lambda_r-v_r$ and $w'_c=\Lambda'_c-v'_c$.

\begin{definition}\label{par:def-harm}
A \emph{joint harmonic} is a polynomial $\Phi\ne0$ with $E_{aa}\Phi=\Lambda_a\Phi$,
$E'_{cc}\Phi=\Lambda'_c\Phi$, $E_{ab}\Phi=0$ ($\varepsilon_a=\varepsilon_b$, $\lambda_a>\lambda_b$),
$E'_{cd}\Phi=0$ ($\varepsilon'_c=\varepsilon'_d$, $\lambda'_c>\lambda'_d$), $\Delta_{ab}\Phi=0$
($\varepsilon_a\ne\varepsilon_b$) and $\Delta'_{cd}\Phi=0$ ($\varepsilon'_c\ne\varepsilon'_d$). A
\emph{pre-joint harmonic} satisfies the same conditions except $E'_{cc}\Phi=\Lambda'_c\Phi$.
\end{definition}
In the notation of \S\ref{pz:ss-types}, $\Phi_j=\bigotimes_\tau\Phi_\tau$ with $\Phi_\tau$ a joint
harmonic at $\tau$ for $(-\Lambda^{(j)}_\tau,\Lambda^{(j+1)}_\tau)$. Being highest for the order of
$-\Lambda^{(j)}$ is being lowest for the order of $\Lambda^{(j)}$: $\Phi_j$ has the $K_j$-weight
$-\Lambda^K(\Lambda^{(j)})$ of the vector $\overline{v_j}$ of $\overline{\pi_j}$.

\begin{lemma}\label{par:relw}
\textup{(i)} $w_r\parsgn\xi_r\ge0$ for all $r$. \textup{(ii)} $w_r=-w'_{c(r)}$, and $w'_c=0$ for $c\in B$.
\textup{(iii)} If $\varepsilon_a=\varepsilon_b$ and $\lambda_a>\lambda_b$ then $\Lambda_a\ge\Lambda_b$;
likewise for $\Lambda'$.
\end{lemma}
\begin{proof}
From $2(\rho_{\mathrm n}-\rho_{\mathrm c})_r=\sum_k\varepsilon_r\varepsilon_k\parsgn(\lambda_k-\lambda_r)=:A_r$
and $2\#\{c:\eta_{rc}=-1\}=n'-\sum_c\eta_{rc}$ we get $2w_r=2\xi_r+A_r-\beta_r$, $\beta_r=\sum_c\eta_{rc}$.
By \eqref{par:eq-match}, $\varepsilon'_{c(k)}=-\varepsilon_k\parsgn\xi_k$, so
$\beta_r=\varepsilon_r\sum_k\varepsilon_k\parsgn\xi_k-\varepsilon_r\sum_{c\in B}\varepsilon'_c$. Let $\xi_r>0$
(the case $\xi_r<0$ is symmetric). Comparing $\parsgn(\xi_k-\xi_r)$ with $\parsgn\xi_k$ term by term,
\[
2w_r=2\xi_r-1-2\!\!\sum_{0<\xi_k<\xi_r}\!\!\varepsilon_r\varepsilon_k+\varepsilon_r\sum_{c\in B}\varepsilon'_c .
\]
All $\xi$ and $\xi'$ lie in $\tfrac{t-1}2+\Z$, so $|\xi_k|>t/2$ gives $|\xi_k|\ge\tfrac{t+1}2$, and at most
$\xi_r-\tfrac{t+1}2$ indices $k$ have $0<\xi_k<\xi_r$. Hence $2w_r\ge2\xi_r-1-(2\xi_r-t-1)-t=0$: this is
(i). For (ii), the same computation for the column $c=c(r)$ (with $\xi'_c=-\xi_r$, signs
$-\eta_{rc}$ for the column action, and $|\xi'_d|<t/2<\xi_r$ for $d\in B$) gives
$2w'_c=-2\xi_r+1+2\sum_{0<\xi_k<\xi_r}\varepsilon_r\varepsilon_k-\varepsilon_r\sum_{B}\varepsilon'_d=-2w_r$. For
$c\in B$ the matched columns contribute $\varepsilon'_c\sum_k\varepsilon_k$, which cancels the column
analogue of $\beta$, and what remains is $2\xi'_c+\sum_{d\in B}\parsgn(\xi'_d-\xi'_c)=0$ because
$\varepsilon'$ is constant on $B$ and $B$ is a string centred at $0$. (iii): $\Lambda_a-\Lambda_b=
\lambda_a-\lambda_b-1-\sum\varepsilon_a\varepsilon_k$ over the $k$ with $\lambda_b<\lambda_k<\lambda_a$, of
which there are at most $\lambda_a-\lambda_b-1$.
\end{proof}

\begin{proposition}[Existence]\label{par:exist}
For $r\in R$ and an integer $0\le\ell<|w_r|$, call $(r,c(r))$ a pair of \emph{level} $\ell$. For each
$\ell$ and each pair of signs $(\epsilon,\epsilon')$ let $L_{\ell,\epsilon,\epsilon'}$ be the set of
rows $r$ having a pair of level $\ell$ with $\varepsilon_r=\epsilon$, $\varepsilon'_{c(r)}=\epsilon'$, and
$M_{\ell,\epsilon,\epsilon'}=\det(z_{r\,c(r')})_{r,r'\in L_{\ell,\epsilon,\epsilon'}}$. Then
$\Phi=\prod_{\ell,\epsilon,\epsilon'}M_{\ell,\epsilon,\epsilon'}$ is a joint harmonic. Moreover no monomial
of $\Phi$ contains two variables of one column from rows of opposite signs, or of one row from columns
of opposite signs.
\end{proposition}
\begin{proof}
$\Phi\ne0$: each minor is a generic determinant in distinct variables. A variable $z_{rc}$ occurs only
if $c=c(r')$ with $r,r'$ in a common layer, so $\varepsilon_r=\varepsilon_{r'}$ and the column sign is
$\varepsilon'_{c(r)}$; this gives the last assertion, hence $\Delta\Phi=0=\Delta'\Phi$. All variables of
row $r$ in $\Phi$ have sign $\eta_{rc(r)}=\parsgn\xi_r$ (\eqref{par:eq-match}).

For $\varepsilon_a=\varepsilon_b$, \eqref{par:eq-osc} gives
$E_{ab}=D^+_{a\leftarrow b}-D^-_{b\leftarrow a}+\delta_{ab}v_a$ with
$D^s_{a\leftarrow b}=\sum_{c:\eta_{bc}=s}z_{ac}\partial_{bc}$, a derivation. \emph{Weight.} $D^s_{a\leftarrow a}$ is
the degree in the sign-$s$ variables of row $a$, which is the number $|w_a|$ of layers containing
$a$ if $s=\parsgn\xi_a$ and $0$ otherwise; by Lemma~\ref{par:relw}(i) $E_{aa}\Phi=(w_a+v_a)\Phi=
\Lambda_a\Phi$. \emph{Raising operators.} Let $\lambda_a>\lambda_b$. If $D^+_{a\leftarrow b}$ acts
non-trivially on a minor, $b$ lies in that layer with $\parsgn\xi_b=+1$; then $\xi_a>\xi_b>0$, $a$ is
matched with a column of the same sign as $c(b)$ (by $\eta_{ac(a)}=\parsgn\xi_a$), and
$w_a\ge w_b$ by Lemma~\ref{par:relw}(iii) ($v_a=v_b$); so $a$ is in the same layer, and the derivation
produces a determinant with two equal rows. The same argument applies to $D^-_{b\leftarrow a}$ with
the roles of $a$, $b$ and of the signs exchanged. So $E_{ab}\Phi=0$. \emph{Columns.} The data are
symmetric under exchanging rows and columns (which replaces $\eta$ by $-\eta$); by
Lemma~\ref{par:relw}(ii) the levels of $(r,c(r))$ seen from the column are the same, columns in $B$
carry no variable and have $w'_c=0$, so the same argument gives the column conditions.
\end{proof}

\begin{lemma}[Polarisations]\label{par:polar}
Let $P$ be killed by the row raising operators and by the $\Delta'_{cd}$. Let $\varepsilon_a=\varepsilon_b$.
\textup{(a)} $\langle D^{-s}_{b\leftarrow a}P,D^{s}_{a\leftarrow b}P\rangle=0$ for $s=\pm1$.
\textup{(b)} If $s\lambda_a>s\lambda_b$ then $D^s_{a\leftarrow b}P=0$.
\textup{(c)} If $E_{aa}P=\Lambda_aP$ and $k=s\,w_a\ge0$, then $D^s_{a\leftarrow a}P=kP$ and
$\partial_{ac}P=0$ whenever $\eta_{ac}=-s$.
\end{lemma}
\begin{proof}
Put $u_{acd}=\partial_{ac}\partial_{ad}P$ and $N_s(a,b)=\sum_{\eta_{ad}=-s,\ \eta_{ac}=s}\langle u_{acd},u_{bcd}
\rangle$. Adjointness gives $\langle D^{-s}_{b\leftarrow a}P,D^s_{a\leftarrow b}P\rangle=N_s(a,b)$ (no
commutator term, as $c\ne d$), and $N_{-s}=N_s$. If $\lambda_a>\lambda_b$, $E_{ab}P=0$ gives
$D^+_{a\leftarrow b}P=D^-_{b\leftarrow a}P$, so $N_1(a,b)\ge0$; the cases $a=b$ and $\lambda_a<\lambda_b$
are similar, so $N_1(a,b)\ge0$ for all $b$ with $\varepsilon_b=\varepsilon_a$. For fixed $c,d$ with
$\eta_{ac}=1=-\eta_{ad}$, $\sum_{b:\varepsilon_b=\varepsilon_a}\partial_{bc}\partial_{bd}P=\pm\Delta'_{cd}P=0$, so
$\sum_bN_1(a,b)=0$, hence all $N_s(a,b)=0$: (a). (b) follows from (a) and the equality of the two
polarisations. (c): $E_{aa}P=\Lambda_aP$ gives $D^s_{a\leftarrow a}P=D^{-s}_{a\leftarrow a}P+kP$; by (a)
with $b=a$, $\|D^{-s}_{a\leftarrow a}P\|^2+k\sum_{\eta_{ac}=-s}\|\partial_{ac}P\|^2=0$.
\end{proof}

\begin{proposition}[Uniqueness]\label{par:unique}
Two pre-joint harmonics are proportional. In particular the joint harmonic is unique up to a scalar.
\end{proposition}
\begin{proof}
Let $x_0$ be the matrix with $z_{rc}=1$ if $c=c(r)$ and $0$ otherwise. We show that a pre-joint
harmonic $F$ has $F(x_0)\ne0$; then for pre-joint harmonics $P,Q$, $P-\tfrac{P(x_0)}{Q(x_0)}Q$
vanishes at $x_0$ and is, if non-zero, a pre-joint harmonic, which is absurd. Suppose $F(x_0)=0$.
Expanding $F$ at $x_0$ and taking a monomial of least positive degree, we find a variable $i$ with
$\partial_iF\ne0$ such that $F$ is an eigenvector of no vector field $\sum_kz_{a_k}\partial_{b_k}$ whose
only term not vanishing at $x_0$ is a multiple of $\partial_i$ (compare the coefficients of that
monomial). Let $i=(r,c)$, $s=\eta_{rc}$. Since $|\xi_r|>t/2$, either $-s\xi_r>t/2$, and then
Lemma~\ref{par:polar}(c) with $-s$ (and Lemma~\ref{par:relw}(i)) gives $\partial_{rc}F=0$; or
$s\xi_r>t/2$, so $\eta_{rc(r)}=s$ and $\varepsilon'_{c(r)}=\varepsilon'_c$. Three cases. If $c=c(r)$, $F$ is an
eigenvector of $D^s_{r\leftarrow r}$ by Lemma~\ref{par:polar}(c), whose only term non-zero at $x_0$
is $z_{rc}\partial_{rc}$. If $s\lambda'_{c(r)}<s\lambda'_c$, the column polarisation
$\sum_{a:\eta_{ac}=s}z_{a\,c(r)}\partial_{ac}$ kills $F$ (Lemma~\ref{par:polar}(b) for columns), and its only
term non-zero at $x_0$ is $a=r$. If $s\lambda'_{c(r)}>s\lambda'_c$, then $|\xi'_c|>t/2$, so $c=c(a)$ with
$s\lambda_a>s\lambda_r$, $\varepsilon_a=\varepsilon_r$, and $D^s_{a\leftarrow r}F=0$ by Lemma~\ref{par:polar}(b),
with only term $d=c$ non-zero at $x_0$. Each case contradicts the choice of $i$.
\end{proof}

\subsubsection{The Casimir comparison and the proof of (A3)}\label{par:ss-A3}

\begin{lemma}\label{par:casimir}
Put $\Omega=\sum E_{ab}E_{ba}$, $\Omega'=\sum E'_{cd}E'_{dc}$, $\mathcal N=\sum_{a,c}a^+_{ac}a^-_{ac}=
\sum_aE_{aa}-n\gamma$ and $s_0=\gamma+\gamma'-\tfrac{n+n'}2$. On polynomials,
$\Omega'=\Omega-2s_0\mathcal N+n^2n'-2\gamma'nn'+n'\gamma'^2-n\gamma^2$. If $(\lambda,\lambda')$ match,
then on every polynomial $\Phi$ with $\sum_aE_{aa}\Phi=(\sum_a\lambda_a)\Phi$ the scalars satisfy
$\kappa_{\lambda'}=\kappa_\lambda-2s_0(\sum_a\lambda_a-n\gamma)+n^2n'-2\gamma'nn'+n'\gamma'^2-n\gamma^2$.
\end{lemma}
\begin{proof}
The signs $s_{ab}$ drop out since $s_{ab}s_{ba}=1$. Write $A=a^+$, $B=a^-$ and
$Q=\sum A_{ac}A_{bd}B_{bc}B_{ad}$. Normal ordering gives $\sum_{a,b}(\sum_cA_{ac}B_{bc})(\sum_dA_{bd}B_{ad})
=Q+n\mathcal N$, and with $F_{cd}=\sum_rB_{rc}A_{rd}$, $\sum_{c,d}F_{cd}F_{dc}=Q+(2n+n')\mathcal N+n^2n'$,
while $\sum_cF_{cc}=\mathcal N+nn'$. For the second count, moving the $A$'s to the left with $[B_i,A_k]=\delta_{ik}$ gives
\[
B_{rc}A_{rd}B_{r'd}A_{r'c}=A_{r'c}A_{rd}B_{rc}B_{r'd}+\delta_{rr'}A_{rd}B_{rd}+\delta_{cd}A_{r'c}B_{r'd}+\delta_{cd}A_{rd}B_{rc}+\delta_{cd},
\]
and summing over $r,r',c,d$ the four correction terms give $n'\mathcal N$, $n\mathcal N$, $n\mathcal N$ and $n^2n'$. Adding the shifts gives the first identity. For the second,
$\lambda'=\{-\lambda_r+s_0\}\cup\{\xi_0+\tfrac{t-1}2-b\}$ with $\xi_0=\gamma'-\tfrac n2$, by
\eqref{par:eq-match}; so $|\lambda'|^2=|\lambda|^2-2s_0\sum\lambda+ns_0^2+t\xi_0^2+\tfrac{t(t^2-1)}{12}$, and
$|\rho'|^2-|\rho|^2=\tfrac{n'(n'^2-1)-n(n^2-1)}{12}$. Substituting $\gamma=s_0-\xi_0+n'/2$,
$\gamma'=\xi_0+n/2$, both sides reduce to $ns_0^2+t\xi_0^2-nn't/4$.
\end{proof}
(Here $\sum\Lambda_a=\sum\lambda_a$, since $\rho_{\mathrm n}-\rho_{\mathrm c}$ is a combination of roots.)

\begin{theorem}[(A3)]\label{par:A3}
Assume Lemma~\ref{par:vac} (that is, (V)), the conclusion of Theorem~\ref{par:A1}, and
\textup{(D1)--(D3)}. Let $0\le j<m$. A joint harmonic $\Phi_j=\bigotimes_\tau\Phi_\tau$ exists and is
unique up to a scalar, and
\[
\int_{G_{j,\infty}}e_j(x)\,\omega_\infty(x^{-1},1)\Phi_j\,dx=\int_{G_{j+1,\infty}}e_{j+1}(x)\,\omega_\infty(1,x)
\Phi_j\,dx .
\]
\end{theorem}
\begin{proof}
Existence and uniqueness at each $\tau$ are Lemma~\ref{par:match} and
Propositions~\ref{par:exist}, \ref{par:unique}. For the product: the conditions of
Definition~\ref{par:def-harm} are conditions on each tensor factor, and on each finite-dimensional
space of polynomials of bounded degree the common kernel of operators acting on separate factors is
the tensor product of the kernels. Both sides are tensor products of the local integrals (Fubini;
$e_j\in L^1$, $\omega_\infty$ unitary), so it suffices to prove the local identity
$P\Phi=P'\Phi$, $P=\int e_{j,\tau}(x)\omega(x^{-1},1)dx$, $P'=\int e_{j+1,\tau}(x)\omega(1,x)dx$, which
are bounded commuting operators on $\pzF_\tau$.

Since $e_{j,\tau}(x^{-1})=\overline{e_{j,\tau}(x)}$ and $\overline{c_{j,\tau}}=c_{-\Lambda^{(j)}}$ (the
conjugate of $\pi_j$ is $\pi_{-\Lambda^{(j)}}$, with $\overline{v_j}$ a highest weight vector for the
order of $-\Lambda^{(j)}$ of its lowest $K$-type), $P=\sigma(e_\lambda)$ for $\sigma=\omega(\cdot,1)$ and
$\lambda=-\Lambda^{(j)}_\tau$; likewise $P'=\sigma'(e_{\lambda'})$, $\lambda'=\Lambda^{(j+1)}_\tau$.

Let $u=P\Phi$. It is smooth for $\sigma'$, and $d\sigma'(Y)u=Pd\sigma'(Y)\Phi$, because $P$ is
bounded and commutes with $\sigma'$ and $\Phi$ is a smooth vector. So $u$ has $T'$-weight $\Lambda'$
and is killed by the compact raising operators of the columns. By Lemma~\ref{par:casimir},
$d\sigma'(\Omega')u=P\Omega'\Phi=P\Omega\Phi+c\,P\Phi$ with $c=-2s_0(\sum\lambda_a-n\gamma)+\dots$, and
$P\,d\sigma(\Omega)\Phi=\kappa_\lambda P\Phi$: integrate by parts twice, using that $e_\lambda$ is an
eigenfunction of the bi-invariant operator $\Omega$ with eigenvalue $\kappa_\lambda$ and that its
derivatives are integrable (Theorem~\ref{par:A1}). So $d\sigma'(\Omega')u=\kappa_{\lambda'}u$, and
Lemma~\ref{par:rigid} for $\sigma'$ gives $P'u=u$. The symmetric argument for $y=P'\Phi$ (with
Lemma~\ref{par:casimir} read backwards) gives $Py=y$. Hence $P\Phi=P'P\Phi=PP'\Phi=P'\Phi$.
\end{proof}

\begin{remark}[The type on $Y_m$]\label{par:Ym}
In Definition~\ref{pz:def-type}, $G_{m,\infty}$ acts on $W_{\lambda^\circ}$ through $g\mapsto\rho(\bar g)$
(\S\ref{ps:ss-weil}). For $y_l(g)=l(\rho(\bar g_\infty)^{-1}\varphi(g_f))$ one has $y_l(gx)=y_{\sigma(x)l}(g)$
with $\sigma(x)l=l\circ\rho(\bar x)^{-1}$. For $t$ in the torus, $\bar t^{-1}=t$, so the dual basis vector
$e_T^*$ has $\sigma$-weight equal to the weight of $e_T$. Thus $\sigma\cong W_{\lambda^\circ}$, the
highest weight vector being $e^*_{T}$ for the highest $T$, and with \eqref{par:eq-coef}
$R(e_m)y_l=y_l$ exactly when $l$ is proportional to that vector. So the type on $Y_m$ is
$W_{\lambda^\circ}$ and not its dual, as Definition~\ref{pz:def-type} states. This agrees with the
definite case \eqref{pk:eq-schuract}, where $h$ acts on covectors $A$ by $A\mapsto A\rho(\bar h)^{-1}$.
\end{remark}



\providecommand{\auU}{\mathrm U}
\providecommand{\auHom}{\operatorname{Hom}}
\providecommand{\pcom}{\mathsf m}
\providecommand{\pcoD}{\mathsf D}
\providecommand{\pcoH}{\mathsf H}
\providecommand{\pcoB}{\mathsf B}
\providecommand{\pcotr}{\operatorname{tr}}

\subsection{The coefficient of the lowest $K$-type and the vacuum character}\label{pco:sec}

This section proves the two statements needed from \S\ref{par:sec}: the integral formula
{\cite[(59)]{P51}} for the coefficient of the lowest $K$-type (Theorem~\ref{pco:FJ}), and the
normalisation (V) of Lemma~\ref{par:vac} (Theorem~\ref{pco:V}). \emph{Dependencies.}
Theorem~\ref{pco:FJ} is proved in full from the published results (D1)--(D3) of
\S\ref{par:ss-rigid}, {\cite[Lemmas~5.4 and~5.5]{P51}} and the argument of {\cite[Lemma~4.7(c)]{P51}},
and standard facts recalled where they are used (highest weight theory for compact groups,
uniqueness of invariant measures on a quotient by a compact subgroup, admissibility). Theorem~\ref{pco:V}
is proved modulo four exactly stated properties of the archimedean Weil representation:
the equivariance (W1$_\infty$), the archimedean form (W2$_\infty$) of (W2), the formula (M) for the metaplectic representation
in a mixed model, and, for definite pairs only, the compatibility (S) with orthogonal sums;
see \S\ref{pco:ss-V}.

\subsubsection{Notation and statement}\label{pco:ss-stat}
We use the conventions of \S\ref{par:sec}: $G=\auU(p,q)=\{g:g^*Jg=J\}$ with
$J=\operatorname{diag}(\varepsilon_1,\dots,\varepsilon_n)$, $K=G\cap\auU(n)$, $T\subset K$ the
diagonal torus, $\lambda\in(\frac{n-1}2+\Z)^n$ with distinct entries, the positive roots
$e_a-e_b$ with $\lambda_a>\lambda_b$, $\rho,\rho_{\mathrm c},\rho_{\mathrm n}$,
$\Lambda=\Lambda^K(\lambda)=\lambda+\rho_{\mathrm n}-\rho_{\mathrm c}$,
\[
\mu=\lambda+\rho=\Lambda+2\rho_{\mathrm c},\qquad\kappa_\lambda=|\lambda|^2-|\rho|^2,
\]
and $h$ is the gap of $\lambda$. Since $\rho_a=\frac12(\#\{b:\lambda_b<\lambda_a\}-\#\{b:\lambda_b>
\lambda_a\})\in\frac{n-1}2+\Z$, $\mu\in\Z^n$, and $\Lambda\in\Z^n$. For $t\in T$ (or any
diagonal matrix $t$) and $w\in\Z^n$ we write $t^w=\prod_at_{aa}^{w_a}$. The coefficient is
$c_\lambda(g)=\langle v,\pi_\lambda(g)v\rangle$, $v$ a unit highest weight vector of the lowest
$K$-type $\Lambda$ of $\pi_\lambda$ (\S\ref{par:ss-decay}); it does not depend on the choice of
$v$, which is unique up to a scalar of modulus one by (D3)(b).

For $x\in M_n(\C)$ and $S\subseteq\{1,\dots,n\}$, $\Delta_S(x)$ is the principal minor on $S$
($\Delta_\emptyset=1$), and for $w\in\Z^n$ and invertible $x$, as in {\cite[(57)]{P51}},
\[
\pcom_w(x)=\prod_{t\in\Z}\Delta_{S_t(w)}(x)\,(\det x)^{-[t\le0]},\qquad S_t(w)=\{a:w_a\ge t\}.
\]
Every $g\in G$ is uniquely $g=kz$ with $k\in K$ and $z\in G$ positive definite
({\cite[\S5.2]{P51}}). Let $du$ be the Haar measure of mass $1$ on $K$, and
$d_\Lambda=(\int_K|\pcom_\Lambda|^2du)^{-1}$ (note $\pcom_\Lambda(1)=1$). Define
\begin{equation}\label{pco:eq-F}
F(kz)=d_\Lambda\int_K\frac{\pcom_\Lambda(ku)\,\overline{\pcom_\Lambda(u)}}{\pcom_{-\mu}(u^{-1}zu)}\,du
\qquad(k\in K,\ z\in G\text{ positive definite}),
\end{equation}
the right side of {\cite[(59)]{P51}}; the denominator is positive by {\cite[Lemma~5.5]{P51}}.

\begin{theorem}\label{pco:FJ}
Assume $h\ge n^2+1$. Then $F=c_\lambda$ on $G$; that is, {\cite[(59)]{P51}} holds.
\end{theorem}

The hypothesis holds for every $\Lambda^{(j)}_\tau$ (the gap is at least $\frac12x^{40}$ and
$n=N_j\le40x^6$, \S\ref{par:sec}), and for the parameters of {\cite[\S5.1]{P51}}
($h\ge2r^8$, $N=r$). The proof occupies \S\S\ref{pco:ss-minors}--\ref{pco:ss-ident}. It follows
the outline of Remark~\ref{par:FJ}: $F$ is real analytic, right $K$-highest of weight $\Lambda$,
an eigenfunction of $\Omega$ with eigenvalue $\kappa_\lambda$ (\S\ref{pco:ss-casimir}),
integrable and square integrable (by the proof of {\cite[Theorem~5.6]{P51}}); a weak form of
Lemma~\ref{par:rigid} and the left $K$-behaviour of $F$ then give $F=c_\lambda$
(\S\ref{pco:ss-ident}). The analytic difficulty is that \eqref{pco:eq-F} uses the polar
decomposition, which is not holomorphic; it is removed by showing that the polar factor $k$ may be
replaced by any block diagonal $a$ for which $a^{-1}g$ is accretive (Proposition~\ref{pco:indep}).

\subsubsection{Products of minors}\label{pco:ss-minors}
Let $\pcoD\subset M_n(\C)$ be the algebra of block diagonal matrices (those commuting with $J$),
$\pcoH=\pcoD\cap\mathrm{GL}_n(\C)\cong\mathrm{GL}_p(\C)\times\mathrm{GL}_q(\C)$, so that $K=\pcoH\cap
\auU(n)$. Let $\pcoB\subset\pcoH$ be the subgroup of $b\in\pcoH$ with $b_{ab}=0$ whenever
$\lambda_a<\lambda_b$ (upper triangular in the order of decreasing $\lambda$), and $\pcoB^-=
\{b^*:b\in\pcoB\}$. A matrix $Z$ is \emph{accretive} if $Z+Z^*$ is positive definite.

\begin{lemma}\label{pco:minors}
Let $w\in\Z^n$ and let $x$ be invertible.
\begin{enumerate}
\item[(a)] $\pcom_{-w}(x)=\pcom_w(x^{-1})$, and $\overline{\pcom_w(x)}=\pcom_w(x^*)$.
\item[(b)] Let $b$ be invertible, with $\Delta_S(b)=\prod_{a\in S}b_{aa}$ for every $S$. If
$b_{ab}=0$ whenever $w_a<w_b$, then $\pcom_w(xb)=b^w\pcom_w(x)$. If $b_{ab}=0$ whenever
$w_b<w_a$, then $\pcom_w(bx)=b^w\pcom_w(x)$. The condition on the minors holds for diagonal $b$,
for $b=1+sE_{ab}$ \textup($a\ne b$, $s\in\C$\textup), and for $b\in\pcoB\cup\pcoB^-$.
\item[(c)] If $\varepsilon_a=\varepsilon_b$ and $\lambda_a>\lambda_b$, then $\Lambda_a\ge\Lambda_b$.
For all $a,b$: $\mu_a>\mu_b$ if and only if $\lambda_a>\lambda_b$.
\item[(d)] If $Z$ is accretive and $y$ invertible, then $y^*Zy$ and $Z^{-1}$ are accretive, and
every principal minor of $Z$ is nonzero; so $\pcom_w(Z)$ is defined and nonzero.
\end{enumerate}
\end{lemma}
\begin{proof}
(a) Jacobi's identity $\Delta_{S^c}(x)=\det x\cdot\Delta_S(x^{-1})$ and
$\{a:-w_a\ge t\}=S_{1-t}(w)^c$ give, with $s=1-t$, the factor $\Delta_{S_s(w)}(x^{-1})
(\det x)^{1-[s\ge1]}=\Delta_{S_s(w)}(x^{-1})(\det x^{-1})^{-[s\le0]}$ of $\pcom_w(x^{-1})$. The
second identity holds factor by factor, since $\overline{\Delta_S(x)}=\Delta_S(\bar x^{\,t})$.

(b) If $b_{cd}=0$ for $c\notin S$, $d\in S$, then $(xb)_{SS}=x_{SS}b_{SS}$, so $\Delta_S(xb)=
\Delta_S(x)\Delta_S(b)$. For $S=S_t(w)$, $c\notin S$, $d\in S$ means $w_c<t\le w_d$, so the
hypothesis applies. With $\Delta_S(b)=\prod_{a\in S}b_{aa}$ (also for $S$ the full set, so
$\det b=\prod b_{aa}$), the factors of $\pcom_w(xb)/\pcom_w(x)$ multiply to $\prod_ab_{aa}^{e_a}$
with $e_a=\#\{t\in(-B,B]:t\le w_a\}-\#\{t\in(-B,B]:t\le0\}=w_a$ for $B$ large. The left version
is the same with rows and columns exchanged. Diagonal and elementary matrices have the stated
minors; a matrix triangular for a total order has every principal submatrix triangular.

(c) Let $m_{\mathrm c}$ (resp.\ $m_{\mathrm n}$) be the number of $c$ with $\varepsilon_c=
\varepsilon_a$ (resp.\ $\ne$) and $\lambda_b<\lambda_c<\lambda_a$. Counting from the definitions,
$\rho_{\mathrm c,a}-\rho_{\mathrm c,b}=1+m_{\mathrm c}$ and $\rho_{\mathrm n,a}-
\rho_{\mathrm n,b}=m_{\mathrm n}$. The entries of $\lambda$ are distinct and congruent modulo $\Z$,
so $m_{\mathrm c}+m_{\mathrm n}\le\lambda_a-\lambda_b-1$, and $\Lambda_a-\Lambda_b=
(\lambda_a-\lambda_b-1-m_{\mathrm c}-m_{\mathrm n})+2m_{\mathrm n}\ge0$. For the second claim,
$\rho_a-\rho_b$ has the sign of $\lambda_a-\lambda_b$.

(d) $v^*(y^*Zy+y^*Z^*y)v>0$ for $v\ne0$; $Z^{-1}+Z^{-*}=Z^{-1}(Z^*+Z)Z^{-*}$. A principal submatrix
$Z_{SS}$ is accretive, and $Z_{SS}v=0$ forces $v^*(Z_{SS}+Z_{SS}^*)v=0$, so $v=0$.
\end{proof}

\begin{corollary}\label{pco:equiv}
Let $Y$ be invertible.
\begin{enumerate}
\item[(a)] For $b\in\pcoB$: $\pcom_\Lambda(Yb)=b^\Lambda\pcom_\Lambda(Y)$, $\pcom_\mu(Yb)=
b^\mu\pcom_\mu(Y)$, $\pcom_\Lambda(b^*Y)=\overline{b^\Lambda}\pcom_\Lambda(Y)$ and $\pcom_\mu(b^*Y)
=\overline{b^\mu}\pcom_\mu(Y)$.
\item[(b)] If $\varepsilon_a=\varepsilon_b$ and $\lambda_a>\lambda_b$, then $\pcom_\Lambda(Y(1+sE_{ab}))
=\pcom_\Lambda(Y)=\pcom_\Lambda((1+sE_{ba})Y)$ for $s\in\C$. For $t$ diagonal,
$\pcom_\Lambda(Yt)=t^\Lambda\pcom_\Lambda(Y)=\pcom_\Lambda(tY)$.
\item[(c)] Put $\psi(Y)=\pcom_{-\mu}(Y)^{-1}$ where $\pcom_{-\mu}(Y)\ne0$. If $\lambda_a<\lambda_b$
\textup(any signs\textup), $\psi(Y(1+sE_{ab}))=\psi(Y)$; for $t$ diagonal, $\psi(Yt)=t^\mu\psi(Y)$.
\end{enumerate}
\end{corollary}
\begin{proof}
Apply Lemma~\ref{pco:minors}(b). (a) An off-diagonal entry $b_{cd}\ne0$ of $b\in\pcoB$ has
$\varepsilon_c=\varepsilon_d$ and $\lambda_c>\lambda_d$, hence $\Lambda_c\ge\Lambda_d$ and $\mu_c>
\mu_d$ by (c); so $b_{cd}=0$ whenever $\Lambda_c<\Lambda_d$ or $\mu_c<\mu_d$. The matrix $b^*$
satisfies the transposed condition, and $(b^*)^w=\overline{b^w}$. (b) The entry $(a,b)$ of
$1+sE_{ab}$ is allowed since $\Lambda_a\ge\Lambda_b$, and likewise the entry $(b,a)$ of $1+sE_{ba}$
on the left. (c) For $w=-\mu$ the entry $(a,b)$ is allowed on the right unless $-\mu_a<-\mu_b$,
i.e.\ unless $\lambda_a>\lambda_b$; and $t^{-(-\mu)}=t^\mu$.
\end{proof}

\subsubsection{Densities on the flag variety of $\pcoH$}\label{pco:ss-density}
Let $T_{\pcoH}=\pcoB\cap K$; a unitary triangular matrix is diagonal, so $T_{\pcoH}=T$. Put
\[
\delta(b)=\bigl|b^{2\rho_{\mathrm c}}\bigr|^2=\prod_{\varepsilon_a=\varepsilon_b,\ \lambda_a>\lambda_b}
|b_{aa}/b_{bb}|^2\qquad(b\in\pcoB).
\]

\begin{lemma}\label{pco:density}
Let $f\colon\pcoH\to\C$ be continuous with $f(yb)=\delta(b)^{-1}f(y)$ for $y\in\pcoH$, $b\in\pcoB$.
Then $\int_Kf(xu)\,du=\int_Kf(u)\,du$ for every $x\in\pcoH$.
\end{lemma}
\begin{proof}
\emph{Step 1 (Iwasawa).} Let $\pcoB_+\subset\pcoB$ be the elements with positive diagonal. The map
$K\times\pcoB_+\to\pcoH$, $(u,b)\mapsto ub$, is a homeomorphism: in each of the two blocks, list
the indices in the order of decreasing $\lambda$; then $\pcoB$ is the upper triangular group, and
the inverse map is the Gram--Schmidt process applied to the columns. In particular $\pcoH=K\pcoB$,
and $ub=u'b'$ ($u,u'\in K$, $b,b'\in\pcoB$) forces $b'b^{-1}=u'^{-1}u\in T$.

\emph{Step 2 (Haar measures of $\pcoB$).} Write $b=tn$ with $t$ diagonal invertible and $n$
unipotent in $\pcoB$; let $d^\times t$ and $dn$ be Haar measures of the two groups (both
unimodular). Left multiplication by $t_0n_0$ sends $(t,n)$ to $(t_0t,(t^{-1}n_0t)n)$, so
$d_lb=d^\times t\,dn$ is a left Haar measure. Right multiplication by $t_0$ sends $(t,n)$ to
$(tt_0,t_0^{-1}nt_0)$, and $n\mapsto t_0^{-1}nt_0$ multiplies $dn$ by $|\det\operatorname{Ad}(t_0)|
_{\mathfrak n}|^{-1}=\delta(t_0)^{-1}$ (real determinant on the Lie algebra $\mathfrak n$ of the
unipotent group, on which $\operatorname{Ad}(t_0)E_{ab}=(t_{0,a}/t_{0,b})E_{ab}$); right
multiplication by $n_0$ is $(t,n)\mapsto(t,nn_0)$. Hence $d_rb=\delta(b)\,d_lb$ is a right Haar
measure, where $\delta(tn)=\delta(t)$.

\emph{Step 3 (an integral formula).} Let $dh$ be a Haar measure of $\pcoH$, which is unimodular. For
$\varphi\in C_c(\pcoH)$ put $\Lambda'(\varphi)=\int_K\int_{\pcoB}\varphi(ub)\,d_rb\,du$; the inner
integrand has compact support $\pcoB\cap u^{-1}\operatorname{supp}\varphi$ and the result is
continuous in $u$, so $\Lambda'$ is a nonzero positive Radon functional. It is invariant under
$\varphi\mapsto\varphi(k^{-1}\cdot\,b_0)$ for $k\in K$, $b_0\in\pcoB$. The group $K\times\pcoB$ acts
on $\pcoH$ by $(k,b)\cdot y=kyb^{-1}$, transitively by Step 1, with stabiliser $\{(t,t):t\in T\}$ of
$1$, which is compact; by Step 1 the orbit map identifies $(K\times\pcoB)/T$ with $\pcoH$
homeomorphically. Both $dh$ and $\Lambda'$ are invariant measures on this homogeneous space, and
an invariant Radon measure on the quotient of a locally compact group $L$ by a compact subgroup $M$ is
unique up to a positive scalar. (If $\mu$ is one, $f\mapsto\int_{L/M}\int_Mf(xm)\,dm\,d\mu(xM)$ is a
left Haar measure on $L$; the map $f\mapsto\int_Mf(\cdot\,m)\,dm$ from $C_c(L)$ onto $C_c(L/M)$ is
surjective, since $\varphi=\int_M(\chi\cdot\varphi\circ p)(\cdot\,m)\,dm$ for $\chi\in C_c(L)$, $\chi\ge0$, with
$\int_M\chi(xm)dm=1$ on $p^{-1}(\operatorname{supp}\varphi)$, $p$ the projection; so uniqueness of Haar measure
gives that of $\mu$. See also \cite[\S2.6, Proposition~(2.48) and Theorem~(2.49), pp.~56--57]{FollandAHA}.) Normalising,
$\int_{\pcoH}\varphi\,dh=\int_K\int_{\pcoB}\varphi(ub)\,d_rb\,du$.

\emph{Step 4.} Choose $\psi\in C_c(\pcoB)$, $\psi\ge0$, $\psi\ne0$, with $\psi(tb)=\psi(b)$ for $t\in T$
(for instance a function of $|b_{aa}|$ and $|b_{ab}|$ only), and put $\Psi(ub)=\psi(b)$; by Step 1
this is well defined and continuous, with compact support $K\cdot T\operatorname{supp}\psi$. Let
$c_\psi=\int\psi\,d_lb>0$. For $x\in\pcoH$ let $A(x)=\int_{\pcoH}f(xh)\Psi(h)\,dh$. By Step 3 and
$f(xub)\psi(b)\,d_rb=f(xu)\psi(b)\,d_lb$,
\[
A(x)=\int_K\int_{\pcoB}f(xub)\psi(b)\,d_rb\,du=c_\psi\int_Kf(xu)\,du .
\]
By left invariance, $A(x)=\int_{\pcoH}f(h)\Psi(x^{-1}h)\,dh=\int_Kf(u)\bigl[\int_{\pcoB}\Psi(x^{-1}ub)
\delta(b)^{-1}d_rb\bigr]du$. Write $x^{-1}u=u'b'$ (Step 1); then $\Psi(x^{-1}ub)=\psi(b'b)$ and the
bracket is $\int\psi(b'b)\,d_lb=c_\psi$. Comparing the two expressions gives the claim.
\end{proof}

\subsubsection{Independence of the block diagonal factor}\label{pco:ss-indep}
For $x\in\mathrm{GL}_n(\C)$ let $A_x=\{a\in\pcoD:x^{-1}a\text{ is accretive}\}$; its elements are
invertible, $A_x$ is open and convex in $\pcoD$ (the accretive matrices form an open convex cone),
and $a\in A_x$ if and only if $a^{-1}x$ is accretive (Lemma~\ref{pco:minors}(d)). For $a\in A_x$ put
\begin{equation}\label{pco:eq-I}
I(x;a)=\int_K\frac{\pcom_\Lambda(au)\,\overline{\pcom_\Lambda(u)}}{\pcom_{-\mu}(u^{-1}a^{-1}xu)}\,du ;
\end{equation}
the denominator is nonzero by Lemma~\ref{pco:minors}(d), and continuous in $u$. If $x=kz\in G$ is
the polar decomposition, then $k\in A_x$ ($k^{-1}x=z$ is positive definite) and
$F(x)=d_\Lambda I(x;k)$.

\begin{proposition}\label{pco:indep}
For every $x\in\mathrm{GL}_n(\C)$, $a\mapsto I(x;a)$ is constant on $A_x$. Moreover the function
$\mathcal J(p,r)$ defined in the proof is locally constant on the open set $U_x$ defined there; in
particular $\mathcal J(p,r)=\mathcal J(1,1)$ whenever $(p,r)$ is joined to $(1,1)$ by a path in $U_x$.
\end{proposition}
\begin{proof}
By Lemma~\ref{pco:minors}(a), $\overline{\pcom_\Lambda(u)}=\pcom_\Lambda(u^*)$ and
$\pcom_{-\mu}(u^{-1}a^{-1}xu)=\pcom_\mu(u^*x^{-1}au)$. Let $U_x$ be the open set of pairs
$(p,r)\in\pcoH\times\pcoH$ with $rx^{-1}p$ accretive, and for $(p,r)\in U_x$ and $y\in\pcoH$ put
\[
f_{p,r}(y)=\frac{\pcom_\Lambda(py)\,\pcom_\Lambda(y^*r)}{\pcom_\mu(y^*rx^{-1}py)},\qquad
\mathcal J(p,r)=\int_Kf_{p,r}(u)\,du .
\]
By Lemma~\ref{pco:minors}(d), $f_{p,r}$ is defined and continuous on all of $\pcoH$, and
$I(x;a)=\mathcal J(a,1)$ for $a\in A_x$.

(i) \emph{$f_{p,r}$ is a density.} For $b\in\pcoB$, Corollary~\ref{pco:equiv}(a) gives
$f_{p,r}(yb)=b^\Lambda\,\overline{b^\Lambda}\,(\overline{b^\mu}b^\mu)^{-1}f_{p,r}(y)=|b^{\Lambda-\mu}|^2
f_{p,r}(y)=\delta(b)^{-1}f_{p,r}(y)$, since $\mu-\Lambda=2\rho_{\mathrm c}$.

(ii) \emph{Invariance.} For $y\in\pcoH$, $(py,y^*r)\in U_x$ and $f_{py,y^*r}(u)=f_{p,r}(yu)$, so
Lemma~\ref{pco:density} gives $\mathcal J(py,y^*r)=\mathcal J(p,r)$.

(iii) \emph{Holomorphy.} The integrand is jointly continuous on $U_x\times K$ and holomorphic in
$(p,r)$ (the $\pcom_w$ are rational functions whose denominators do not vanish there); since $K$ is
compact, $\mathcal J$ is holomorphic on $U_x$ and may be differentiated under the integral sign
(Cauchy's formula on small polydiscs and Fubini).

(iv) \emph{Polarisation.} Fix $(p,r)\in U_x$ and $X\in\pcoD$. For $s\in\C$ near $0$, (ii) with
$y=e^{sX}$ gives $\mathcal J(pe^{sX},e^{\bar sX^*}r)=\mathcal J(p,r)$. Write $\partial_1,\partial_2$
for the holomorphic differentials in $p$ and $r$. Differentiating along $s=\sigma\in\R$ and along
$s=i\sigma$ at $\sigma=0$,
\[
\partial_1\mathcal J[pX]+\partial_2\mathcal J[X^*r]=0,\qquad i\,\partial_1\mathcal J[pX]-
i\,\partial_2\mathcal J[X^*r]=0,
\]
so $\partial_1\mathcal J(p,r)[pX]=0$. As $X$ runs over $\pcoD$, so does $pX$; hence $\partial_1
\mathcal J=0$ on $U_x$. The function $a\mapsto\mathcal J(a,1)$ on the convex open set $A_x$ is
therefore holomorphic with vanishing differential, hence constant. Dividing the second identity by
$i$ and comparing with the first gives also $\partial_2\mathcal J(p,r)[X^*r]=0$; as $X$ runs over
$\pcoD$, so does $X^*r$ ($\pcoD$ is stable under $X\mapsto X^*$ and $r$ is invertible in $\pcoD$), so
$\partial_2\mathcal J=0$. Both holomorphic differentials of $\mathcal J$ vanish on $U_x$, so
$\mathcal J$ is locally constant there, hence constant along every path in $U_x$.
\end{proof}

\subsubsection{Local holomorphic form, $K$-behaviour and the Casimir eigenvalue}\label{pco:ss-casimir}
For $Z\in\mathfrak{gl}_n(\C)=\mathfrak g\otimes_\R\C$ and a $C^1$ function $\phi$ on $G$ let $R(Z)\phi$ be
the right derivative, $(R(X)\phi)(x)=\frac d{ds}\phi(xe^{sX})|_{s=0}$ for $X\in\mathfrak g$, extended
complex-linearly; then $R(Z_1)R(Z_2)-R(Z_2)R(Z_1)=R([Z_1,Z_2])$, and $\Omega=\sum_{a,b}E_{ab}E_{ba}$
acts by $R(\Omega)=\sum_{a,b}R(E_{ab})R(E_{ba})$. For a holomorphic function $\Phi$ on an open
subset of $\mathrm{GL}_n(\C)$ let $(\widetilde Z\Phi)(x)=d\Phi(x)[xZ]=\frac d{ds}\Phi(x(1+sZ))|_{s=0}$;
again $[\widetilde Z_1,\widetilde Z_2]=\widetilde{[Z_1,Z_2]}$. If $\phi=\Phi|_G$ near a point, then
$R(Z_1)\cdots R(Z_k)\phi=(\widetilde Z_1\cdots\widetilde Z_k\Phi)|_G$ there: for $Z\in\mathfrak g$ this is
the chain rule, and both sides are complex-linear in each $Z_i$.

\begin{lemma}\label{pco:local}
Let $g_0=k_0z_0\in G$ \textup(polar decomposition\textup) and let $W$ be the open set of
$x\in\mathrm{GL}_n(\C)$ with $k_0^{-1}x$ accretive. Then $\widetilde F(x)=d_\Lambda I(x;k_0)$ is
holomorphic on $W$, and $F=\widetilde F$ on $W\cap G$, a neighbourhood of $g_0$ in $G$. In particular
$F$ is real analytic on $G$, and $F(1)=1$.
\end{lemma}
\begin{proof}
For $x\in W\cap G$ with polar decomposition $x=kz$, both $k$ and $k_0$ lie in $A_x$, so
$F(x)=d_\Lambda I(x;k)=d_\Lambda I(x;k_0)$ by Proposition~\ref{pco:indep}. Holomorphy: the integrand
of $I(x;k_0)$ is jointly continuous on $W\times K$ and holomorphic in $x$, as in
Proposition~\ref{pco:indep}(iii). Finally $F(1)=d_\Lambda\int_K|\pcom_\Lambda|^2du/\pcom_{-\mu}(1)=1$.
\end{proof}

\begin{lemma}[$K$-behaviour]\label{pco:K}
Let $x\in G$, $a\in A_x$ and $k_1\in K$. Then
\begin{align}
F(xk_1)&=d_\Lambda\int_K\frac{\pcom_\Lambda(au)\,\pcom_\Lambda(u^*k_1)}{\pcom_{-\mu}(u^{-1}a^{-1}xu)}\,du,
\label{pco:eq-right}\\
F(k_1x)&=d_\Lambda\int_K\frac{\pcom_\Lambda(k_1au)\,\overline{\pcom_\Lambda(u)}}{\pcom_{-\mu}(u^{-1}a^{-1}xu)}
\,du.\label{pco:eq-left}
\end{align}
Consequently $F(xt)=t^\Lambda F(x)=F(tx)$ for $t\in T$, and if $\varepsilon_a=\varepsilon_b$,
$\lambda_a>\lambda_b$, then $R(E_{ab})F=0$.
\end{lemma}
\begin{proof}
$ak_1\in A_{xk_1}$ because $(ak_1)^{-1}xk_1=k_1^{-1}(a^{-1}x)k_1$; by Lemma~\ref{pco:local} and
Proposition~\ref{pco:indep} (applied at the point $xk_1$), $F(xk_1)=d_\Lambda I(xk_1;ak_1)$. Substitute
$u\mapsto k_1^{-1}u$ and use $\overline{\pcom_\Lambda(k_1^{-1}u)}=\pcom_\Lambda(u^*k_1)$
(Lemma~\ref{pco:minors}(a)). Similarly $k_1a\in A_{k_1x}$ and $(k_1a)^{-1}k_1x=a^{-1}x$, which gives
\eqref{pco:eq-left}. The weights follow from Corollary~\ref{pco:equiv}(b). For the derivative, fix
$x$ and $a\in A_x$; the right side of \eqref{pco:eq-right} is the restriction to $K$ of a
holomorphic function of $k_1\in\pcoH$, so $R(E_{ab})F(x)$ is its derivative along
$k_1=1+sE_{ab}$ at $s=0$, which vanishes by Corollary~\ref{pco:equiv}(b).
\end{proof}

\begin{proposition}[Casimir]\label{pco:casimir}
$R(\Omega)F=\kappa_\lambda F$ on $G$.
\end{proposition}
\begin{proof}
Fix $g_0$ and use Lemma~\ref{pco:local}: near $g_0$, $F=\widetilde F|_G$ with
$\widetilde F(x)=d_\Lambda\int_K\alpha(u)\,\psi(L_u(x))\,du$, where $\alpha(u)=\pcom_\Lambda(k_0u)
\overline{\pcom_\Lambda(u)}$, $L_u(x)=u^{-1}k_0^{-1}xu$ and $\psi=\pcom_{-\mu}^{-1}$ as in
Corollary~\ref{pco:equiv}(c). By Proposition~\ref{pco:indep}(iii) (applied to $\widetilde F$),
holomorphic derivatives of $\widetilde F$ may be taken under the integral sign. Since
$L_u(x(1+sZ))=L_u(x)(1+s\,u^{-1}Zu)$, $\widetilde Z(\psi\circ L_u)=(\widetilde{Z^u}\psi)\circ L_u$ with
$Z^u=u^{-1}Zu$, and so
\[
\widetilde\Omega(\psi\circ L_u)=\Bigl(\sum_{a,b}\widetilde{E_{ab}^u}\,\widetilde{E_{ba}^u}\,\psi\Bigr)\circ L_u
=(\widetilde\Omega\psi)\circ L_u,
\]
because $\sum_{a,b}c^{-1}E_{ab}c\otimes c^{-1}E_{ba}c=\sum_{a,b}E_{ab}\otimes E_{ba}$ for $c\in
\mathrm{GL}_n(\C)$ (as an operator on $\C^n\otimes\C^n$ this tensor is the flip, which commutes with
$c\otimes c$). It remains to show $\widetilde\Omega\psi=\kappa_\lambda\psi$ on the open set where
$\pcom_{-\mu}\ne0$. By Corollary~\ref{pco:equiv}(c), $\widetilde E_{aa}\psi=\mu_a\psi$ and
$\widetilde E_{ab}\psi=0$ if $\lambda_a<\lambda_b$. For a pair $\lambda_a<\lambda_b$,
$\widetilde E_{ba}\widetilde E_{ab}\psi=0$ and $\widetilde E_{ab}\widetilde E_{ba}\psi=\widetilde E_{ba}
\widetilde E_{ab}\psi+(\widetilde E_{aa}-\widetilde E_{bb})\psi=(\mu_a-\mu_b)\psi$. Summing,
\[
\widetilde\Omega\psi=\Bigl(\sum_a\mu_a^2-\sum_{\lambda_a<\lambda_b}(\mu_b-\mu_a)\Bigr)\psi
=(|\mu|^2-2\langle\mu,\rho\rangle)\psi=(|\lambda|^2-|\rho|^2)\psi,
\]
since $2\rho=\sum_{\lambda_a<\lambda_b}(e_b-e_a)$ and $\mu=\lambda+\rho$.
\end{proof}

\begin{lemma}[Decay]\label{pco:decay}
For every integer $0\le c\le h$ and $g\in G$, $|F(g)|\le(2/\|g\|)^c$. If $h\ge n^2+1$, then $F\in
L^1(G)\cap L^2(G)\cap L^\infty(G)$.
\end{lemma}
\begin{proof}
The proof of {\cite[Theorem~5.6]{P51}} uses {\cite[(59)]{P51}} only through its right side, which is
$F$, together with {\cite[Lemmas~5.4 and~5.5]{P51}}; these are stated for an
arbitrary $\auU(J)$, an arbitrary regular $\lambda\in(\frac{n-1}2+\Z)^n$ and its gap. So it proves
the bound for $F$. With $c=h$, $\int_G\|g\|^{-c}dg<\infty$ by the proof of
{\cite[Lemma~4.7(c)]{P51}}, which is stated there for $G_\tau$ but applies verbatim to any $\auU(p,q)$,
$n=p+q$: in $G=KA^+K$, with coordinates $t_1\ge\dots\ge t_{\min(p,q)}\ge0$ on $A^+$, one has
$\|kak'\|=e^{t_1}$ and Haar density $\delta(a)\le e^{n(n-1)t_1}$ (\cite[Ch.~I, \S5, no.~2]{Helgason}, \cite[Lemma~2.4.2]{WallachRRG1}: $\delta(a)=\prod_{\alpha>0}\sinh\alpha(\log a)$, with multiplicities, up to a constant), so the integral
is bounded by a constant times $\int_{A^+}e^{(n(n-1)-c)t_1}\,da<\infty$ for $c>n(n-1)$, in particular
for $c=h\ge n^2+1$. Also $|F|\le2^c$.
\end{proof}

\subsubsection{Identification with $c_\lambda$}\label{pco:ss-ident}
Let $(\rho_\Lambda,V_\Lambda)$ be the irreducible representation of $K$ of highest weight $\Lambda$
for the positive compact roots $e_a-e_b$, $\varepsilon_a=\varepsilon_b$, $\lambda_a>\lambda_b$, with
unit highest weight vector $v_\Lambda$, and $\beta(k)=\dim V_\Lambda\cdot\overline{\langle v_\Lambda,
\rho_\Lambda(k)v_\Lambda\rangle}$. We use two standard facts on a unitary representation $\sigma$ of
$K$, with $\sigma(\beta)=\int_K\beta(k)\sigma(k)\,dk$.
\begin{enumerate}
\item[(K1)] $\sigma(\beta)$ is the orthogonal projection onto the space of vectors of the
$\Lambda$-isotypic subspace that are highest weight vectors of weight $\Lambda$. (Peter--Weyl and
Schur orthogonality: on an irreducible summand $\cong V_\Lambda$, $\sigma(\beta)u=\langle v_\Lambda,u
\rangle v_\Lambda$ in the identification, and on other types $\sigma(\beta)=0$.)
\item[(K2)] If $\sigma$ is finite-dimensional and $y$ satisfies $\sigma(t)y=t^\Lambda y$ ($t\in T$) and
$d\sigma(E_{ab})y=0$ for the positive compact roots, then $\sigma(\beta)y=y$. (Decompose $y$ along
irreducible summands; in an irreducible module a vector killed by the positive root vectors is a
highest weight vector \cite[Theorem~5.5(c), p.~279]{KnappBeyond}; then apply (K1).)
\end{enumerate}

\begin{lemma}\label{pco:Kfin}
For every $x\in G$, $\int_K\beta(k)F(xk)\,dk=F(x)$ and $\int_K\beta(k)\,\overline{F(k^{-1}x)}\,dk=
\overline{F(x)}$.
\end{lemma}
\begin{proof}
Let $\mathcal P$ be the space of functions on $K$ of the form $(\det k)^{-e}P(k)$, with $P$ a
polynomial of degree $\le D$ in the entries of $k$ ($e,D$ fixed and large enough). It is
finite-dimensional and stable under right translation $\sigma(k_1)\phi=\phi(\cdot\,k_1)$, since the
entries of $kk_1$ are linear in those of $k$ and $\det(kk_1)=\det k\det k_1$. Let $\mathcal
V\subset\mathcal P$ be the span of the functions $k\mapsto\pcom_\Lambda(Yk)$, $Y\in
\mathrm{GL}_n(\C)$; it is finite-dimensional, hence closed in $\mathcal P$. By
Corollary~\ref{pco:equiv}(b) every $\phi\in\mathcal V$ satisfies $\sigma(t)\phi=t^\Lambda\phi$
($t\in T$) and $d\sigma(E_{ab})\phi=0$ (the derivative of the holomorphic extension along
$1+sE_{ab}$); these are pointwise properties of $\phi$. By (K2) applied to $\sigma$ on $\mathcal P$,
$\int_K\beta(k)\phi(k_1k)\,dk=\phi(k_1)$ for $\phi\in\mathcal V$. By \eqref{pco:eq-right} the function
$k_1\mapsto F(xk_1)$ is an integral of elements of $\mathcal V$ against a continuous function of
$u$, hence lies in $\mathcal V$; evaluating at $k_1=1$ gives the first identity. By
\eqref{pco:eq-left} and Lemma~\ref{pco:minors}(a), $\phi(k_1)=\overline{F(k_1^{-1}x)}=d_\Lambda\int_K
\pcom_\Lambda(u^*a^*k_1)\,(\cdots)\,du$ also lies in $\mathcal V$, which gives the second.
\end{proof}

\begin{proposition}[Rigidity for $F$]\label{pco:rigid}
Let $R$ be the right regular representation of $G$ on $L^2(G)$. Then $\langle F,R(g)F\rangle=
\|F\|^2c_\lambda(g)$ for $g\in G$; consequently $c_\lambda\in L^1(G)$, and with $e_\lambda=d_\lambda
\overline{c_\lambda}$ \textup($d_\lambda$ the formal degree\textup), $R(e_\lambda)F=F$.
\end{proposition}
\begin{proof}
We follow the proof of Lemma~\ref{par:rigid}, with its hypotheses on $y=F$ replaced by the
following three properties, which hold by Lemmas~\ref{pco:decay}, \ref{pco:Kfin} and
Proposition~\ref{pco:casimir}: (1) $F\in L^2(G)$ and $F=R(\beta)F$, where $R(\beta)=\int_K\beta(k)R(k)dk$
(by Lemma~\ref{pco:Kfin}, pointwise, hence in $L^2$); (2) for every $\varphi\in C_c^\infty(G)$,
$\kappa_\lambda R(\varphi)F=R(\Omega\varphi)F$, where $\Omega\varphi=R(\Omega)\varphi$. For (2):
$(R(\varphi)F)(x)=\int\varphi(h)F(xh)\,dh$, and for $X\in\mathfrak g$, substituting $h\mapsto he^{-sX}$
and differentiating under the integral sign (compact support, $F\in C^\infty$),
$R(\varphi)(R(X)F)=-R(R(X)\varphi)F$; iterating, $R(\varphi)(R(\Omega)F)=R(\sum_{a,b}R(E_{ba})R(E_{ab})
\varphi)F=R(\Omega\varphi)F$, and $R(\Omega)F=\kappa_\lambda F$.

Let $H_F$ be the closed span of $R(G)F$, $H_F=\int^\oplus H_\zeta\,d\nu(\zeta)$ a decomposition into
irreducible unitary representations \cite{Dixmier}, $F=\int^\oplus y_\zeta$. Fix a sequence
$\varphi_i\in C_c^\infty(G)$, $\varphi_i\ge0$, $\int\varphi_i=1$, with supports shrinking to $\{1\}$.
The identities $F=R(\beta)F$ and $\kappa_\lambda R(\varphi_i)F=R(\Omega\varphi_i)F$ ($i\ge1$) are
countably many and decompose; so for almost every $\zeta$, $y_\zeta=\sigma_\zeta(\beta)y_\zeta$ and
$\kappa_\lambda\sigma_\zeta(\varphi_i)y_\zeta=\sigma_\zeta(\Omega\varphi_i)y_\zeta$ for all $i$. By (K1),
$y_\zeta$ is a highest weight vector of weight $\Lambda$ in the $\Lambda$-isotypic subspace of
$H_\zeta$, which is finite-dimensional (admissibility, \cite[Theorem~8.1, p.~205]{Knapp}); so $y_\zeta$ is
$K$-finite, hence a smooth vector, and $\sigma_\zeta(\Omega\varphi_i)y_\zeta=\sigma_\zeta(\varphi_i)
d\sigma_\zeta(\Omega)y_\zeta$ by the same integration by parts. Letting $i\to\infty$,
$d\sigma_\zeta(\Omega)y_\zeta=\kappa_\lambda y_\zeta$. From here the proof of Lemma~\ref{par:rigid}
applies word for word: if $y_\zeta\ne0$, (D1) and (D2) give the infinitesimal character $\lambda$,
Lemma~\ref{par:D3} gives $H_\zeta\cong\pi_\lambda$ and (D3)(b) gives $y_\zeta=a(\zeta)U_\zeta v$, and $\langle F,R(g)F\rangle=
\int|a|^2d\nu\cdot c_\lambda(g)=\|F\|^2c_\lambda(g)$.

Since $F\in L^1$, $\int_G|\langle F,R(g)F\rangle|\,dg\le\int\!\!\int|F(x)||F(xg)|\,dx\,dg=\|F\|_1^2$,
and $\|F\|\ne0$ ($F$ is continuous with $F(1)=1$); so $c_\lambda\in L^1$. The last assertion is the
last step of the proof of Lemma~\ref{par:rigid}: by Schur orthogonality $\int e_\lambda(h)c_\lambda
(g^{-1}h)\,dh=c_\lambda(g^{-1})$, so $\langle R(g)F,R(e_\lambda)F-F\rangle=0$ for all $g$, and
$R(e_\lambda)F\in H_F$.
\end{proof}

\begin{proof}[Proof of Theorem~\ref{pco:FJ}]
Write $\pi=\pi_\lambda$, $c=c_\lambda$, $e=e_\lambda=d\bar c$ with $d=d_\lambda$.

\emph{Step 1: $F(g)=d\int_GF(x)\,c(x^{-1}g)\,dx$.} By Proposition~\ref{pco:rigid}, $F=R(e)F$ in
$L^2$. Since $F$ is bounded and continuous and $e\in L^1$, $g\mapsto\int e(h)F(gh)\,dh$ is
continuous (dominated convergence), so $F(g)=\int e(h)F(gh)\,dh$ for every $g$. Substitute $x=gh$
(the Haar measure is left invariant) and use $e(h)=d\,\overline{c(h)}=d\,c(h^{-1})$, which holds
because $c(h)=\langle v,\pi(h)v\rangle$ is a diagonal coefficient of a unitary representation:
\[
F(g)=\int e(g^{-1}x)F(x)\,dx=d\int F(x)\,c(x^{-1}g)\,dx ,
\]
the integral converging absolutely since $F\in L^1$ and $|c|\le1$.

\emph{Step 2: $F(g)=\langle w,\pi(g)v\rangle$ with $w=d\int_G\overline{F(x)}\,\pi(x)v\,dx$.} The
integral converges in norm since $F\in L^1$, and $c(x^{-1}g)=\langle\pi(x)v,\pi(g)v\rangle$, with inner
products conjugate-linear in the first variable.

\emph{Step 3: $w$ is a multiple of $v$.} By Step 2 and the second identity of Lemma~\ref{pco:Kfin},
for all $g$,
\[
\langle\pi(\beta)w,\pi(g)v\rangle=\int_K\overline{\beta(k)}\langle w,\pi(k^{-1}g)v\rangle dk
=\overline{\int_K\beta(k)\overline{F(k^{-1}g)}\,dk}=F(g)=\langle w,\pi(g)v\rangle .
\]
Since $\pi$ is irreducible, the $\pi(g)v$ span a dense subspace, so $\pi(\beta)w=w$. By (K1), $w$ is
a highest weight vector of weight $\Lambda$ in the $\Lambda$-isotypic subspace of $\pi$, which is
irreducible by (D3)(b); so $w=\alpha v$ for some $\alpha\in\C$.

\emph{Step 4.} Hence $F=\bar\alpha\,c$. Evaluating at $g=1$, $1=F(1)=\bar\alpha\,c(1)=\bar\alpha$
(Lemma~\ref{pco:local}). So $F=c$.
\end{proof}

\begin{remark}\label{pco:rem-FJ}
The proof uses: {\cite[Lemmas~5.4 and~5.5]{P51}} (Fischer's inequality, Jacobi's
identity), {\cite[Lemma~4.7(c)]{P51}}, the published results (D1)--(D3) and \cite{Dixmier} as in
Lemma~\ref{par:rigid}, Schur orthogonality for square-integrable representations, admissibility
and smoothness of $K$-finite vectors \cite[Theorem~8.1 and Proposition~8.5]{Knapp}, highest weight theory for $K$, and
uniqueness of invariant measures on a quotient by a compact subgroup. No property of the signature
or of the special parameters of {\cite[\S5.1]{P51}} is used beyond $h\ge n^2+1$. The
non-vanishing of $F$ and the normalisation at the identity are both $F(1)=1$, which is immediate
from \eqref{pco:eq-F}; the scalar $\alpha$ in Step~3 is not assumed nonzero.
\end{remark}

\subsubsection{The vacuum character (V)}\label{pco:ss-V}
Fix $0\le j<m$ and $\tau$, and let $V=Y_j\otimes_{E,\tau}\C$, $V'=Y_{j+1}\otimes_{E,\tau}\C$, of
signatures $(p,q)$, $(p',q')$ and ranks $n,n'$, and $W=\auHom(V,V')$ with the form
$\langle w_1,w_2\rangle=\pcotr(w_1^\dagger w_2)$ of \S\ref{ps:ss-weil}; it has signature
$(P,Q)=(pp'+qq',pq'+qp')$. Let $\iota\colon\auU(V)\times\auU(V')\to\auU(W)$, $\iota(g,h)w=hwg^{-1}$;
then $\det_W\iota(g,h)=\det(h)^n\det(g)^{-n'}$, and on Lie algebras
$\pcotr_W\iota(X,Y)=n\,\pcotr Y-n'\pcotr X$. Let $\omega_\infty$ be the component at $\tau$ of the
Weil representation $\omega_W$ of $\auU(W)$ defined with $\upsilon$ (\S\ref{ps:ss-weil}), acting on
the Fock space $\mathfrak F_\tau$ of the Heisenberg representation $\varrho_\tau$, so that
$\omega_{\auHom(V,V')}(g,h)=\omega_\infty(\iota(g,h))$ at $\tau$. Let $\mathrm{Mp}$ be the metaplectic
double cover of the symplectic group of $W$ viewed as a real symplectic space (symplectic form
$\Iml$ of the hermitian form, as in $\beta_\infty$), $\omega^{\mathrm{Mp}}$ the metaplectic
(oscillator) representation on $\mathfrak F_\tau$, characterised up to a scalar on each element by
$\omega^{\mathrm{Mp}}(\tilde u)\varrho_\tau(w)\omega^{\mathrm{Mp}}(\tilde u)^{-1}=\varrho_\tau(uw)$
\cite[Th\'eor\`eme~1 and no.~34]{WeilActa}, and $\widetilde{\auU}(W)$ the inverse image of $\auU(W)(\R)$ in $\mathrm{Mp}$.

We state (V) in the form in which Lemma~\ref{par:vac} uses it, the derived action on polynomial
vectors:
\begin{enumerate}
\item[(V$'$)] For $X\in\mathfrak u(V)$, $Y\in\mathfrak u(V')$, on polynomial vectors,
$d\omega_\infty(\iota(X,Y))=d\omega^{\mathrm{Mp}}(\iota(X,Y))+\tfrac{n'}2\pcotr X-\tfrac n2\pcotr Y$.
\end{enumerate}
Here $d\omega^{\mathrm{Mp}}$ on polynomial vectors is given by the Weyl-ordered quadratic operators
in the $a^\pm_{rc}$, in the following sense. Let $\mathcal V$ be the span of the operators $\frac12(a_ia_k+a_ka_i)$ on polynomials, $a_i,a_k$ running over the $a^\pm_{rc}$. By Howe \cite[p.~537, (2.2), and (2.3)]{HoweJAMS}, the polynomials are the $\widetilde K$-finite vectors of $\omega^{\mathrm{Mp}}$ for the maximal compact subgroup $\widetilde K$ of $\mathrm{Mp}$ (so they are smooth vectors, stable under $d\omega^{\mathrm{Mp}}$), and $d\omega^{\mathrm{Mp}}$ maps the complexified Lie algebra onto $\mathcal V$ (there: the spans of $\frac12(z_i\partial_j+\partial_jz_i)$, $z_iz_j$, $\partial_i\partial_j$, which together span $\mathcal V$). Differentiating Weil's relation on polynomial vectors (which are smooth for $\varrho_\tau$, with $d\varrho_\tau$ given by linear combinations of the $a^\pm_{rc}$ and scalars) gives $[d\omega^{\mathrm{Mp}}(X),d\varrho_\tau(w)]=d\varrho_\tau(Xw)$. This determines $d\omega^{\mathrm{Mp}}(X)$ within $\mathcal V$: if $Q\in\mathcal V$ commutes with all $a^\pm_{rc}$, then $Q$ commutes with all $z_{rc}$ and $\partial_{rc}$, so $Q1$ is a constant $c$ and $Qf=fQ1=cf$; and a scalar lies in $\mathcal V$ only if it is $0$ (writing $\mathcal V$ as the span of $z_iz_j$, $\partial_i\partial_j$, $z_i\partial_j+\frac12\delta_{ij}$, up to signs: if a combination of these equals $c$, applying it to the polynomials of degree $\le2$ shows that all coefficients vanish, hence $c=0$). So $d\omega^{\mathrm{Mp}}(X)$ is the unique element of $\mathcal V$ whose commutator with $d\varrho_\tau(w)$ is $d\varrho_\tau(Xw)$; this is what ``Weyl-ordered quadratic operator'' means here. Exponentiating, (V$'$) says that
$\omega_\infty(\iota(g,h))=\det(g)^{n'/2}\det(h)^{-n/2}\,\omega^{\mathrm{Mp}}(\widetilde{\iota(g,h)})$,
the product of the two genuine factors being a representation of the linear group; this is (V).

We use the following properties of the archimedean Weil representation, for every nondegenerate
hermitian space $W$ over $E_\tau$ (not only those of the form $\auHom(V,V')$); Proposition~\ref{pwf:W2}(c) and Lemma~\ref{pwf:sum} of \S\ref{pwf:sec}, on which they rest, are
stated for an arbitrary nondegenerate hermitian space.
\begin{enumerate}
\item[(W1$_\infty$)] $\omega_\infty(u)\varrho_\tau(w)\omega_\infty(u)^{-1}=\varrho_\tau(uw)$ for
$u\in\auU(W)(\R)$; and $\omega_\infty$ is strongly continuous. (The component at $\tau$ of (W1).)
\item[(W2$_\infty$)] Let $I\subset W$ be an isotropic complex line, $I^*$ an isotropic line with
$\langle I^*,I\rangle\ne0$, $W_I=(I+I^*)^\perp$, and $P_I$ the stabiliser of $I$ in $\auU(W)(\R)$.
There is a nonzero continuous linear map $\mathsf Z\colon\mathfrak F_\tau^{\mathrm{sm}}\to
\mathfrak F_\tau(W_I)^{\mathrm{sm}}$ (evaluation at $0$ in the mixed model $\mathcal S(I^*)\otimes
\mathfrak F_\tau(W_I)$) with $\mathsf Z\,\omega_\infty(p)=\chi_{I,\tau}(p)\,\omega_{W_I,\infty}(\bar p)
\mathsf Z$ for $p\in P_I$, where $\chi_{I,\tau}(p)=\upsilon_\tau(a)^{-1}|a|$ if $p$ acts on $I$ by
$a\in\C^\times$. (The component at $\tau$ of (W2) and (W2$'$); it is proved in
Proposition~\ref{pwf:W2}(c), local statement, from the Levi action of the Siegel parabolic in
Kudla's Schr\"odinger model \cite[Theorem~3.1 and \S5]{KudlaSplitting} (Lemma~\ref{pwf:siegel}) and
the compatibility with orthogonal sums (Lemma~\ref{pwf:sum}).)
\item[(M)] In the same mixed model, for $X$ in the Lie algebra of the subgroup of $P_I$ that acts
trivially on $W_I$ and preserves $I$ and $I^*$, $\mathsf Z\,d\omega^{\mathrm{Mp}}(X)=\pm\frac12
\pcotr_\R(X|_{I})\,\mathsf Z$ on smooth vectors, where $\pcotr_\R$ is the trace on $I$ viewed as a
real plane. (The Levi subgroup $\mathrm{GL}(I)$ of the Siegel parabolic of $\mathrm{Sp}(I+I^*)$ acts in
the Schr\"odinger model by $\phi(x)\mapsto|\det_\R a|^{\pm1/2}\phi(a^{\mp1}x)$ up to a scalar, these operators being intertwiners for Weil's relation, as in the proof of Lemma~\ref{pwf:siegel}. Infinitesimally: by the description of $d\omega^{\mathrm{Mp}}$ just given, transported to the Schr\"odinger model on the first factor (the commutation relations determine the element of the span of symmetrised quadratics in $x_i$, $\partial_i$ in the same way), if $X$ acts on the coordinates of $I^*$ by a real matrix $A$, then $d\omega^{\mathrm{Mp}}(X)=\mp\frac12\sum_{i,k}A_{ik}(x_i\partial_k+\partial_kx_i)=\mp(\sum_{i,k}A_{ik}x_i\partial_k+\frac12\pcotr A)$; evaluation at $0$ kills the vector field.)
\item[(S)] (Only for definite $W$.) For a hermitian line $W'$ over $E_\tau$, under
$\mathfrak F_\tau(W\oplus W')=\mathfrak F_\tau(W)\hat\otimes\mathfrak F_\tau(W')$,
$\omega_{W\oplus W',\infty}(u\oplus1)=\omega_{W,\infty}(u)\otimes1$ for $u\in\auU(W)(\R)$
(compatibility of the $\upsilon$-splittings with orthogonal sums, Lemma~\ref{pwf:sum}; cf.\ \cite[p.~389]{KudlaSplitting}).
\end{enumerate}

\begin{theorem}\label{pco:V}
Assume \textup{(W1$_\infty$)}, \textup{(W2$_\infty$)}, \textup{(M)} for every nondegenerate hermitian
space over $E_\tau$ (in particular for $W$ and, if $W$ is definite, for $W\oplus W'$ as in Step~3),
and, if $W$ is definite, \textup{(S)}. Then \textup{(V$'$)} holds; that is, \textup{(V)} of Lemma~\ref{par:vac} holds at
$\tau$, with $\upsilon_\tau|_{\auU(1)}(z)=z$.
\end{theorem}
\begin{proof}
\emph{Step 1: the ratio is a character.} For $\tilde u\in\widetilde{\auU}(W)$ over $u$, by
(W1$_\infty$) the operator $\omega_\infty(u)\omega^{\mathrm{Mp}}(\tilde u)^{-1}$ commutes with every
$\varrho_\tau(w)$; $\varrho_\tau$ is irreducible, so it is a scalar $\chi(\tilde u)$ of modulus one.
Then $\chi$ is a continuous character of $\widetilde{\auU}(W)$, and on smooth (in particular
polynomial) vectors $d\omega_\infty(X)=d\omega^{\mathrm{Mp}}(X)+d\chi(X)$ for $X\in\mathfrak u(W)$.
Since $d\chi$ vanishes on $[\mathfrak u(W),\mathfrak u(W)]=\mathfrak{su}(W)$, $d\chi=c\cdot\pcotr_W$
for some $c\in\C$. By the trace formula for $\iota$, (V$'$) is equivalent to $c=-\frac12$.

\emph{Step 2: $W$ indefinite.} Then $P,Q\ge1$ and there are $I,I^*$ as in (W2$_\infty$). The
elements $m_\theta\in\auU(W)(\R)$ acting by $e^{i\theta}$ on $I+I^*$ and trivially on $W_I$ lie in
$P_I$, with $a=e^{i\theta}$ and $\bar m_\theta=1$; their generator $X_0$ has $\pcotr_WX_0=2i$ and
$\pcotr_\R(X_0|_I)=0$ (multiplication by $i$ on a real plane). Differentiating
(W2$_\infty$) along $m_\theta$ at $\theta=0$ on a smooth vector $\Phi$ (continuity of $\mathsf Z$),
with $\upsilon_\tau(e^{i\theta})=e^{i\theta}$, gives $\mathsf Z\,d\omega_\infty(X_0)\Phi=-i\,\mathsf Z\Phi$;
and (M) gives $\mathsf Z\,d\omega^{\mathrm{Mp}}(X_0)\Phi=0$. By Step 1, $d\chi(X_0)\mathsf Z\Phi=-i
\,\mathsf Z\Phi$; choosing $\Phi$ with $\mathsf Z\Phi\ne0$, $2ic=-i$, so $c=-\frac12$.

\emph{Step 3: $W$ definite.} Let $W'$ be a line with form of the sign opposite to that of $W$; then
$W\oplus W'$ is indefinite, and Step~2 gives $d\chi_{W\oplus W'}=-\frac12\pcotr_{W\oplus W'}$. By
(S), $d\omega_{W,\infty}(X)\otimes1=d\omega_{W\oplus W',\infty}(X\oplus0)$, and the Weyl-ordered
operator of $X\oplus0$ is that of $X$ tensored with $1$ (it involves only the variables of $W$); so
$d\chi_W(X)=d\chi_{W\oplus W'}(X\oplus0)=-\frac12\pcotr_WX$.
\end{proof}

\begin{remark}\label{pco:rem-V}
(a) $W_\tau$ is indefinite as soon as $Y_j$ or $Y_{j+1}$ is indefinite at $\tau$: if
$p,q\ge1$ and $n'\ge1$, then $P,Q\ge1$; if $V$ is definite and $p',q'\ge1$, likewise. By
\S\ref{ps:ss-Y}, $Y_j$ contains the string $f_0,\dots,f_{s_j-1}$, which is indefinite for
$s_j\ge1$, i.e.\ $j\le m-2$; for $j=m-1$ both $Y_{m-1}=X_{m-1}$ and $Y_m$ are positive definite, so
among the pairs $(Y_j,Y_{j+1})$, (S) is used exactly for the last pair $j=m-1$. (In general (S) is
needed exactly when $W_\tau$ is definite.)
For definite pairs (V$'$) is the formula of \S\ref{pk:ss-conv} (``no further multiplier''); the
argument there identifies $\omega_\infty$ with $F\mapsto F(h^tz\bar g)$ only up to a character,
and Step~3 supplies the character.

(b) The alternative of \emph{defining} $\omega_\infty$ by the Fock-model formulas with the
determinant twists, as an explicit representation, does not remove the need for a comparison:
(W1) requires $\theta\circ\omega(\gamma)=\theta$ for $\gamma\in\auU(W)(E^+)$, and multiplying
$\omega_\infty$ by a character $\prod_\tau\det(g_\tau)^{a_\tau}$ multiplies $\theta\circ\omega(\gamma)$
by $\prod_\tau\tau(\det\gamma)^{a_\tau}$, which is not identically $1$ on $\auU(W)(E^+)$ unless all
$a_\tau=0$. So with that definition (W1) would have to be proved, and that is equivalent to
Theorem~\ref{pco:V}. We therefore keep the definition of \S\ref{ps:ss-weil} and derive (V) from
the local properties (W1$_\infty$), (W2$_\infty$), (M), (S).

(c) (W2$_\infty$) is the only input through which the value of $\upsilon_\tau$ enters the
computation of $c$; (M) and the Weyl-ordered form of $d\omega^{\mathrm{Mp}}$ are statements about the
metaplectic representation alone, and (S) only transports $c$ from $W\oplus W'$ to $W$. The signs
in (M) do not matter, because $\pcotr_\R(X_0|_I)=0$.
\end{remark}



\providecommand{\auO}{\mathcal O}
\providecommand{\auPl}{\Phi}
\providecommand{\auPa}{\Phi_a}
\providecommand{\auPb}{\Phi_b}
\providecommand{\auS}{\mathcal S}
\providecommand{\auSram}{\mathcal S^{\mathrm{ram}}}
\providecommand{\auDX}{D_X}
\providecommand{\augain}{\operatorname{gain}}
\providecommand{\auerr}{\operatorname{err}}
\providecommand{\auNm}{\operatorname{Nm}}
\providecommand{\auU}{\mathrm U}
\providecommand{\auGL}{\mathrm{GL}}
\providecommand{\auA}{\mathbb A}
\providecommand{\auF}{\mathbb F}
\providecommand{\aucol}{\operatorname{col}}
\providecommand{\aubar}[1]{\overline{#1}}
\providecommand{\auf}[1]{\texttt{#1}}
\providecommand{\pao}{\mathfrak o}
\providecommand{\pam}{\mathfrak m}
\providecommand{\paCp}{\C_p}
\providecommand{\paI}{I_p^{\circ}}
\providecommand{\pavp}{\varpi}
\providecommand{\paded}{\operatorname{ded}}
\providecommand{\paPred}{\operatorname{Pred}}
\providecommand{\paloss}{\operatorname{loss}}
\providecommand{\panot}[1]{\par\smallskip{\raggedright\noindent\textbf{[Not proved here:} #1\textbf{]}\par}\smallskip}

\section{Divisibility of the composite lift at every prime}\label{pa:sec}

This section gives a complete proof of {\cite[Proposition~14.10]{P46}} (\S\ref{pa:ss-count}); a
complete proof, from standard facts on elliptic curves with complex multiplication, that products of
theta functions of a CM curve along a resolution of the identity have the normal form
{\cite[(44)]{P46}} (\S\S\ref{pa:ss-curve}--\ref{pa:ss-prod}); and the deduction of
{\cite[Proposition~14.20]{P46}} from these (\S\ref{pa:ss-divis}). The deduction is complete for the
primes that are not exceptional and differ from $N$, given one structural statement on the
theta kernels (Lemma~\ref{pa:kernel}). This lemma is proved in \S\ref{pk:ss-kernel}
(Propositions~\ref{pk:K0}, \ref{pk:std}, \ref{pk:K12}) and in Corollary~\ref{pf:K1} (the integrality of the
jets in (K1) at the primes $\ell\mid2q_1M$). For the exceptional primes
and for $N$ the deduction rests on three further statements
(Lemmas~\ref{pa:hecke}, \ref{pa:eldiv}, \ref{pa:level}), proved in {\cite[\S\S2 and~3]{P50}} and~\S\ref{pl:sec}; the case of $N$ uses Lemma~\ref{pa:kernel} in the extended form (K1$_N$), which
is Lemma~\ref{pl:kernelN} and is proved in Corollary~\ref{pf:K1}(b).

\subsection{Conventions}\label{pa:ss-conv}
Throughout, $p$ is a prime that does not split in $K$; thus $p$ is odd ($2$ splits, as
$d\equiv7\pmod 8$), $p\nmid q_1$, and either $p$ is inert or $p=d$. We fix a ring homomorphism
$\C\to\paCp$, write $\pao$ for the valuation ring of $\paCp$, $\pam$ for its maximal ideal and $v$
for the valuation with $v(p)=1$; numbers and power series with coefficients in $\pao$ are called
\emph{integral}, and congruences are between integral objects. Let $\pi=\pi_p\in\auO_K$ be $p$ if
$p$ is inert and $\sqrt{-d}$ if $p=d$, and $q=q_p=\auNm(\pi)$, so $q=p^2$ or $q=p$, $v(\pi)=1$ or
$\tfrac12$, and $q/\pi^2$ is a unit. Put
\[
\rho=\rho_p=\frac{v(\pi)}{q-1},\qquad\gamma=\gamma_p=\frac1{p-1}-\rho,\qquad
\paI=I_p\cup(\{0\}\cap\{p=d\}),
\]
so that $\paI$ is the set of odd positive integers if $p$ is inert and $\N$ if $p=d$. In both
cases
\begin{equation}\label{pa:eq-const}
(p+1)\rho-1+v(\pi)=\frac1{p-1},\qquad\gamma=v(\pi)\sum_{e\in I_p}p^{-e}.
\end{equation}
We use the following elementary facts. For $n\ge1$, with $s_p(n)$ the sum of the digits of $n$ in
base $p$,
\begin{equation}\label{pa:eq-legendre}
\frac n{p-1}-1-\log_pn\ \le\ v(n!)=\frac{n-s_p(n)}{p-1}\ \le\ \frac{n-1}{p-1}.
\end{equation}
\begin{itemize}
\item[(E1)] If $h$ is an integral power series in any set of variables with $h(0)=0$ and
$v(c)\ge\frac1{p-1}$, then $\exp(ch)$ is integral: $v(c^k/k!)\ge0$ by \eqref{pa:eq-legendre}.
\item[(E2)] If moreover $\beta\in\pao$ has $v(\beta)\ge\tfrac12$, then
$\exp(\beta h)\in1+\beta h\,\pao[[\cdot]]$ and $\log(1+\beta h)\in\beta h\,\pao[[\cdot]]$: for
$k\ge1$, $v(\beta^{k-1}/k!)\ge(k-1)(\tfrac12-\tfrac1{p-1})\ge0$ and
$v(\beta^{k-1}/k)\ge\tfrac{k-1}2-v(k)\ge0$, because $p\ge3$.
\item[(E3)] In a commutative $\pao$-algebra, $(A+B)^{p^e}\equiv A^{p^e}+B^{p^e}\pmod p$.
\end{itemize}

\subsection{The count}\label{pa:ss-count}
Let $\auPl$ be a finite set with $f$ elements, $\mathsf C$ a finite set, and
$z=(z_{\tau,i})_{\tau\in\auPl,\,i\in\mathsf C}$ variables. Let $\mathrm{pr}_e\colon\auPl\to\auPl$
($e\in\paI$) be arbitrary maps (\emph{partner maps}); in the application
$\mathrm{pr}_e(\tau)=u(-p^e)\tau$. The \emph{couplings} are
$\Gamma_e=\sum_{\tau,i}z_{\tau,i}^{\,p^e}z_{\mathrm{pr}_e\tau,i}$, homogeneous of degree $p^e+1$.
A series $F\in\paCp[[z]]$ \emph{has the normal form} if, for some $\pavp$ with $v(\pavp)=\rho$,
\begin{equation}\label{pa:eq-nf}
F(\pavp z)=W(z)\exp\bigl(c_\pavp\,\Xi(z)\bigr),\qquad c_\pavp=\frac{\pavp^{p+1}}p,\qquad
\Xi=\sum_{e\in\paI}\sum_{j\ge0}\lambda_{e,j}\Gamma_e^{\,p^j},
\end{equation}
with $W$ integral and $\lambda_{e,j}\in\pao$; this is {\cite[(44)]{P46}}. By
\eqref{pa:eq-const}, $v(c_\pavp\pi)=\frac1{p-1}$. Hence, by (E1), \emph{it suffices to know the
exponent modulo $\pi$}: if $F(\pavp z)=W\exp(c_\pavp\Xi')$ with $W,\Xi'$ integral,
$\Xi'(0)=0$ and $\Xi'\equiv\Xi\pmod\pi$, then $F$ has the normal form (replace $W$ by
$W\exp(c_\pavp(\Xi'-\Xi))$). Also, if $F$ has the normal form and $\alpha\in\pao$, then so has
$F(\alpha z)$, with $\lambda_{e,j}$ replaced by $\alpha^{(p^e+1)p^j}\lambda_{e,j}$, the couplings
being homogeneous.

Now let $\mathsf C=[n']\times[n]$, so that for each $\tau$ the variables form an $n'\times n$
matrix $z_\tau=(z_{\tau,b,a})$. A \emph{shape} $\kappa$ gives numbers $\kappa_\tau(h)\ge0$ of
columns of height $h$ at $\tau$, for $1\le h<n'$. Put
\[
a_\tau=\kappa_\tau(1),\qquad a'_\tau=\sum_{h\ge2}\kappa_\tau(h),\qquad
l_\tau=\sum_{h\ge2}h\,\kappa_\tau(h),\qquad B=\sum_\tau(a_\tau+l_\tau),
\]
the numbers of columns of height one, of the other (\emph{low}) columns, of the boxes in these, and
of all boxes. A \emph{filling} $T$ of $\kappa$ assigns to each column $\mathfrak c=(\tau,h,c)$,
$1\le c\le\kappa_\tau(h)$, a sequence $T(\mathfrak c)=(b_1,\dots,b_h)$ of rows in $[n']$. For
matrices $y_\tau\in M_{n\times(n'-1)}(\pao)$ (\emph{directions}) let
\[
Q_{T,y}(z)=\prod_{\mathfrak c=(\tau,h,c)}\det\bigl((z_\tau y_\tau)_{b_i,k}\bigr)_{1\le i,k\le h},
\qquad(b_1,\dots,b_h)=T(\mathfrak c),
\]
the product over the columns of the minors of $z_\tau y_\tau$ on the rows of the column and on the
first $h$ columns; it is homogeneous of degree $B$ with integral coefficients. For $e\in\paI$ and
$\tau\in\auPl$, with $\tau'=\mathrm{pr}_e\tau$, the \emph{deduction} is
\[
\paded_e(\tau)=\min\Bigl(\frac{a_\tau}{p^e},\,a_{\tau'}+a'_{\tau'}\Bigr)+\frac{a'_\tau}{p^e}
+\min\Bigl(\frac{l_\tau}{p^e},\,l_{\tau'}\Bigr),
\]
and the \emph{predicted divisibility} of $\kappa$ is
\[
\paPred(\kappa)=v(\pi)\sum_{e\in I_p}\sum_{\tau}\Bigl(\bigl(\tfrac{a_\tau}{p^e}-a_{\tau'}-
a'_{\tau'}\bigr)^+ +\bigl(\tfrac{l_\tau}{p^e}-l_{\tau'}\bigr)^+-\tfrac{a'_\tau}{p^e}\Bigr)
-\mathbf 1_{p=d}\cdot\tfrac12\sum_\tau\paded_0(\tau).
\]
Since $u-\min(u,u')=(u-u')^+$, the second identity of \eqref{pa:eq-const} gives
\begin{equation}\label{pa:eq-pred}
\paPred(\kappa)=\gamma B-v(\pi)\sum_{e\in\paI}\sum_\tau\paded_e(\tau).
\end{equation}

\begin{proposition}[The count; {\cite[Proposition~14.10]{P46}}]\label{pa:count}
Let $F\in\paCp[[z]]$ have the normal form \eqref{pa:eq-nf}, let $\kappa$ be a shape with $B\ge1$
boxes, $T$ a filling of $\kappa$ and $y$ integral directions. Then
\[
v\bigl(Q_{T,y}(\partial)F(0)\bigr)\ \ge\ \gamma B-v(\pi)\!\!\sum_{e\in\paI,\ p^e\le B}\ \sum_\tau
\paded_e(\tau)-f\,n'^3\,(1+\log_pB)\ \ge\ \paPred(\kappa)-f\,n'^3\,(1+\log_pB).
\]
\end{proposition}

For $p$ inert the right-hand side is the expression of {\cite[Proposition~14.10]{P46}}; for $p=d$ the
term subtracted in addition is the $e=0$ term $v(\pi_d)\sum_\tau\paded_0(\tau)=\tfrac12\sum_\tau\paded_0(\tau)$. Together with
Lemma~\ref{pa:schur}, the bound holds for the jets along $Q_{T,y}$ and along the entries
$S_\kappa(z)_{T',T}$ of the Schur matrix; these are the only jets to which it is applied, and no
claim is made for products of minors on non-nested column sets. The proof needs one lemma of
linear algebra.

\begin{lemma}[Elimination]\label{pa:elim}
Let $h\ge1$, let $\ell_1,\dots,\ell_s\colon\pao^h\to\pao$ be linear forms and
$\alpha_1,\dots,\alpha_s\in\pao$. There is $g\in M_h(\pao)$ with $\det g\mid\alpha_1\cdots\alpha_s$
such that $\alpha_i\mid\ell_i(r)$ for every row $r$ of $g$ and every $i$.
\end{lemma}
\begin{proof}
Induction on $s$, starting from $g=1$. Let $g$ serve $\ell_1,\dots,\ell_{s-1}$, with rows
$r_1,\dots,r_h$. If all $\ell_s(r_k)$ vanish, $g$ serves. Otherwise let $k_0$ be an index for which
$v(\ell_s(r_{k_0}))$ is least, replace $r_k$ by $r_k-\bigl(\ell_s(r_k)/\ell_s(r_{k_0})\bigr)r_{k_0}$
for $k\ne k_0$, and $r_{k_0}$ by $\alpha_sr_{k_0}$. The new rows are integral combinations of the
old ones, so the earlier divisibilities persist; $\ell_s$ vanishes on the first kind and is
divisible by $\alpha_s$ on the last; and the determinant has been multiplied by $\alpha_s$.
\end{proof}

\begin{proof}[Proof of Proposition~\ref{pa:count}]
The second inequality is \eqref{pa:eq-pred}, the deductions being non-negative. For the first, let
$S_B=\{e\in\paI:p^e\le B\}$.

\emph{Step 1: a linear substitution.} Introduce new variables $s_{\tau,h,k,b}$ for $\tau\in\auPl$,
$1\le k\le h<n'$ and $b\in[n']$; there are at most $f\,n'^3$ of them. For matrices
$g_{\tau,h}\in M_h(\pao)$, to be chosen, define the \emph{new directions}
$\beta_{\tau,h,k}=\sum_{k'\le h}g_{\tau,h}(k,k')\,y_\tau(\cdot,k')\in\pao^n$ and the substitution
\[
z_{\tau,b,a}(s)=\sum_{h,k}\beta_{\tau,h,k}(a)\,s_{\tau,h,k,b}.
\]
Let $G(s)=F(\pavp z(s))$ and $P_*(s)=\prod_{\mathfrak c=(\tau,h,c)}\det(s_{\tau,h,k,b_i})_{1\le
i,k\le h}$, a homogeneous polynomial of degree $B$ with integer coefficients. For a linear
substitution $z=Ms$ and a polynomial $P$ one has $P(\partial_s)(F\circ M)=(\widetilde
P(\partial_z)F)\circ M$, where $\widetilde P$ is obtained from $P$ by replacing each $s_o$ by the
linear form $\sum_xM_{x,o}z_x$. Here $s_{\tau,h,k,b}$ is replaced by
$\sum_a\beta_{\tau,h,k}(a)z_{\tau,b,a}=\sum_{k'}g_{\tau,h}(k,k')(z_\tau y_\tau)_{b,k'}$, so the
matrix of a column $\mathfrak c$ becomes the corresponding matrix of $z_\tau y_\tau$ times the
transpose of $g_{\tau,h}$. Taking determinants, and using that $Q_{T,y}$ is homogeneous of degree
$B$,
\begin{equation}\label{pa:eq-chain}
P_*(\partial)G(0)=\pavp^{B}\prod_{\tau,h}\det(g_{\tau,h})^{\kappa_\tau(h)}\cdot Q_{T,y}(\partial)F(0).
\end{equation}

\emph{Step 2: the couplings after substitution.} For $\beta,\beta'\in\pao^n$ put
$[\beta,\beta']_e=\sum_a\beta(a)^{p^e}\beta'(a)$. By (E3), for $\tau'=\mathrm{pr}_e\tau$,
\[
\Gamma_e(z(s))\equiv\sum_{\tau,b}\ \sum_{h,k}\sum_{h',k'}[\beta_{\tau,h,k},\beta_{\tau',h',k'}]_e\;
s_{\tau,h,k,b}^{\,p^e}\,s_{\tau',h',k',b}\pmod p .
\]
We claim that the $g_{\tau,h}$ can be chosen such that
\begin{equation}\label{pa:eq-goal}
\pi\mid[\beta_{\tau,h,k},\beta_{\mathrm{pr}_e\tau,h',k'}]_e\ \ \text{for all }e\in S_B,\ \tau,\
h,k,h',k',\qquad
\sum_{\tau,h}\kappa_\tau(h)\,v(\det g_{\tau,h})\le v(\pi)\sum_{e\in S_B}\sum_\tau\paded_e(\tau).
\end{equation}
For $e\in S_B$ fix $\pi_e\in\pao$ with $\pi_e^{\,p^e}=\pi$, and for $w\in\pao^n$ a vector
$w^{1/p^e}$ of $p^e$-th roots of its coordinates. By (E3),
\begin{equation}\label{pa:eq-root}
\pi_e\mid\textstyle\sum_aw^{1/p^e}(a)\beta(a)\ \Longrightarrow\ \pi\mid[\beta,w]_e .
\end{equation}
For each pair $(\tau,e)$, with $\tau'=\mathrm{pr}_e\tau$, we impose the following congruences on
the rows of the matrices at $\tau$ (the \emph{source}) and at $\tau'$; recall that a new direction
of height $1$ at a place $\sigma$ is a multiple of $y_\sigma(\cdot,1)$.
\begin{itemize}
\item[(a)] On every row $r$ of $g_{\tau,h}$, $h\ge2$: one congruence modulo $\pi_e$, namely
$\pi_e\mid\sum_a y_{\tau'}(a,1)^{1/p^e}\beta_r(a)$, where $\beta_r=\sum_{k'}r_{k'}y_\tau(\cdot,k')$.
By \eqref{pa:eq-root} it gives $\pi\mid[\beta,\beta']_e$ for $\beta$ of height $h\ge2$ at $\tau$
and $\beta'$ of height $1$ at $\tau'$.
\item[(b)] If $a_\tau/p^e\le a_{\tau'}+a'_{\tau'}$: $\pi_e$ divides the $1\times1$ matrix
$g_{\tau,1}$. Otherwise: on every row $r$ of every $g_{\tau',h'}$, $h'\ge1$, one congruence modulo
$\pi$, namely $\pi\mid[y_\tau(\cdot,1),\beta'_r]_e$, which is linear in $r$. In both cases
$\pi\mid[\beta,\beta']_e$ for $\beta$ of height $1$ at $\tau$ and all $\beta'$ at $\tau'$.
\item[(c)] If $l_\tau/p^e\le l_{\tau'}$: $\pi_e$ divides all entries of $g_{\tau,h}$, $h\ge2$
($h$ congruences on each row). Otherwise: $\pi$ divides all entries of $g_{\tau',h'}$, $h'\ge2$. In
the first case $\pi\mid[\beta,\beta']_e$ for all $\beta$ of height $\ge2$ at $\tau$; in the second,
for all $\beta'$ of height $\ge2$ at $\tau'$.
\end{itemize}
Every pair of heights $(h,h')$ is covered: $h=1$ by (b); $h\ge2$, $h'=1$ by (a); $h,h'\ge2$ by
(c) or, in its second case, by the congruences on $\beta'$. The conditions imposed on the rows of
one matrix $g_{\sigma,h}$ by all pairs $(\tau,e)$ with $\tau=\sigma$ or $\mathrm{pr}_e\tau=\sigma$
are finitely many congruences of linear forms, so Lemma~\ref{pa:elim} gives $g_{\sigma,h}$ whose
determinant has valuation at most the sum of the valuations of the moduli. A modulus $\pi_e$ costs
$v(\pi)/p^e$ and a modulus $\pi$ costs $v(\pi)$. Counting each matrix $\kappa_\sigma(h)$ times, the
pair $(\tau,e)$ contributes: from (a), $v(\pi)a'_\tau/p^e$; from (b),
$v(\pi)\min(a_\tau/p^e,a_{\tau'}+a'_{\tau'})$; from (c), $v(\pi)\min(l_\tau/p^e,l_{\tau'})$. The
sum is $v(\pi)\paded_e(\tau)$, which proves \eqref{pa:eq-goal}.

\emph{Step 3: $G$ is integral in degrees at most $B$.} By \eqref{pa:eq-nf},
$G=W(z(s))\exp(c_\pavp\Xi(z(s)))$. For $e\in\paI\setminus S_B$ the series $\Gamma_e(z(s))^{p^j}$ has
no terms of degree $\le B$. For $e\in S_B$, Step~2 and $\pi\mid p$ give
$\pi\mid\Gamma_e(z(s))$. Hence $\Xi(z(s))=\pi h+h'$ with $h,h'$ integral, $h(0)=0$, and $h'$ of
order $>B$. As $v(c_\pavp\pi)=\frac1{p-1}$, (E1) shows that $\exp(c_\pavp\pi h)$ is integral; and
$G$ agrees with $W(z(s))\exp(c_\pavp\pi h)$ in degrees $\le B$.

\emph{Step 4: Legendre.} Write $P_*=\sum_\delta c_\delta s^\delta$ with $c_\delta\in\Z$,
$|\delta|=B$. Then $P_*(\partial)G(0)=\sum_\delta c_\delta\,\delta!\,G_\delta$, where $G_\delta$ is
the coefficient of $s^\delta$ in $G$, integral by Step~3. By \eqref{pa:eq-legendre}, applied to each
of the at most $fn'^3$ exponents of $\delta$,
$v(\delta!)\ge\frac B{p-1}-fn'^3(1+\log_pB)$. With \eqref{pa:eq-chain}, \eqref{pa:eq-goal} and
$v(\pavp^B)=\rho B$,
\[
v\bigl(Q_{T,y}(\partial)F(0)\bigr)\ge\frac B{p-1}-fn'^3(1+\log_pB)-\rho B-v(\pi)\sum_{e\in S_B}
\sum_\tau\paded_e(\tau),
\]
and $\frac1{p-1}-\rho=\gamma$.
\end{proof}

\begin{remark}[Marked places]\label{pa:charges}
A variant is used at the exceptional primes (Lemma~\ref{pa:hecke}); it is proved as
Proposition~\ref{pl:charges}. Let
$\mathsf E\subseteq\auPl$ and let integers $w_{\tau,h,b}\ge0$ (\emph{charges}) be such that, for
$e\in\paI$, $\tau'=\mathrm{pr}_e\tau$, every row $b$ and every height $h'$: if $\tau\in\mathsf E$,
or $h'=1$ and $\tau'\in\mathsf E$, then $w_{\tau,1,b}>0$ or $w_{\tau',h',b}>0$. Substitute
$z_{\tau,b,a}(s)=\sum_{h,k}p^{w_{\tau,h,b}}\beta_{\tau,h,k}(a)s_{\tau,h,k,b}$. The pairs of heights
$(1,h')$ of the hypothesis then need no congruence, since the coefficient of
$s_{\tau,1,k,b}^{p^e}s_{\tau',h',k',b}$ acquires the factor $p^{p^ew_{\tau,1,b}+w_{\tau',h',b}}$,
which is divisible by $p$; so the columns of height one at the places of $\mathsf E$ may be left out of the
deductions, and the result is: $v(Q_{T,y}(\partial)F(0))+\sum_{\text{boxes}}w\ge
\paPred(\kappa^\circ)+\gamma(B-B^\circ)-fn'^3(1+\log_pB)$, where $\kappa^\circ$ is $\kappa$
without its columns of height one at $\mathsf E$, $B^\circ$ its number of boxes, and the sum is
over the boxes of $T$, a box in row $b$ of a column of height $h$ at $\tau$ counting
$w_{\tau,h,b}$.
\end{remark}

\subsection*{From entries of the Schur matrix to products of minors}
Let $n\le n'$ and let $\kappa$ be a shape all of whose columns have height at most $n$. For a
generic $n\times\infty$ matrix $Y$ and a filling $T$ of a diagram by elements of $[n]$ let
$\nabla_T(Y)$ be the product over the columns of the minors of $Y$ on the rows of the column and
the first $h$ columns. By standard monomial theory~\cite[\S\S1--2, Thm.~2.1]{DEP}, \cite[\S8.1, Thm.~1 and its corollary]{FultonYT} the $\nabla_T$ with $T$
standard form a $\Z$-basis of the span of all $\nabla_T$ of a given shape, and they remain linearly
independent over every field; the lattice $W_\kappa(\Z)$ of {\cite[\S14.3]{P46}} is the dual of the tensor product over $\tau$ of
these spans, with the basis $(e_T)$ dual to the standard $\nabla_T$ (\S\ref{pl:ss-norm}, (N1)). For $z_\tau\in M_{n'\times n}$ and a standard filling $T'$ by
elements of $[n']$, the Cauchy--Binet formula and the straightening law give
\[
\prod_\tau\nabla_{T'_\tau}(z_\tau Y_\tau)=\sum_{T\text{ standard}}S_\kappa(z)_{T',T}\,\prod_\tau
\nabla_{T_\tau}(Y_\tau),\qquad S_\kappa(z)_{T',T}\in\Z[z];
\]
$S_\kappa(z)$ is the \emph{Schur matrix}, the matrix of the map $W_\kappa(\C^n)\to W_\kappa(\C^{n'})$,
$\mathrm{ev}_Y\mapsto\mathrm{ev}_{zY}$, in the bases $(e_T)$.

\begin{lemma}\label{pa:schur}
For $F\in\paCp[[z]]$ and $T'$ standard,
$\min_Tv\bigl(S_\kappa(\partial)_{T',T}F(0)\bigr)\ge\inf_yv\bigl(Q_{T',y}(\partial)F(0)\bigr)$,
the infimum over all integral directions $y$.
\end{lemma}
\begin{proof}
Put $J_T=S_\kappa(\partial)_{T',T}F(0)$. Specialising $Y_\tau$ to $y_\tau$ in the identity above,
$Q_{T',y}(\partial)F(0)=\sum_TJ_T\nabla_T(y)$. Let $T_0$ have $|J_{T_0}|$ largest; we may assume
$J_{T_0}\ne0$. The polynomial $\sum_T(J_T/J_{T_0})\nabla_T(Y)$ is integral and its reduction modulo
$\pam$ is not zero, by the linear independence of the $\nabla_T$ over the residue field. This
field is infinite, so the reduction does not vanish at some point, which lifts to an integral $y$
with $v(Q_{T',y}(\partial)F(0))=v(J_{T_0})$.
\end{proof}

\subsection{The theta function of a CM curve at a non-split prime}\label{pa:ss-curve}
Let $\Lambda\subset\C$ be a lattice with $\auO_K\Lambda=\Lambda$, of area $A$, with Weierstrass
functions $\sigma,\wp$, and let $t\in\tfrac12\Lambda\setminus\Lambda$ be the \emph{admissible} point: the one with
$\omega t\notin\Lambda$ and $\aubar\omega t\notin\Lambda$, $\omega=(1+\sqrt{-d})/2$ (\S\ref{pf:ss-adm}). After scaling
$\Lambda$ we may assume $\wp''(t)=2$. Put
\[
x(z)=\wp(z+t)-\wp(t),\qquad a_1^2=12\,\wp(t),\qquad y(z)=\tfrac12\bigl(\wp'(z+t)-a_1x(z)\bigr).
\]
With $e_1=\wp(t)$ and $e_2,e_3$ the other roots of the cubic of $\wp$ one has $e_2+e_3=-e_1$ and
$(e_1-e_2)(e_1-e_3)=\wp''(t)/2=1$, so $\wp'^2=4x(x^2+3e_1x+1)$ at $z+t$ and
\begin{equation}\label{pa:eq-curve}
y^2+a_1xy=x^3+x,\qquad x'=2y+a_1x,\qquad y'=1+3x^2-a_1y,
\end{equation}
the last from $\wp''=6\wp^2-g_2/2=6x^2+a_1^2x+2$. The point $z=0$ is the point $P_0=(0,0)$ of
order two of the curve $\mathcal E\colon y^2+a_1xy=x^3+x$, the point at infinity is $z\equiv t$,
and $y=z+O(z^2)$. All Taylor series at $z=0$ below have algebraic coefficients and are read in
$\paCp[[z]]$.

\begin{lemma}\label{pa:model}
$a_1$ is an algebraic integer and the discriminant $a_1^4-64$ of $\mathcal E$ is a unit at every
odd prime. The reduction $\widetilde{\mathcal E}$ modulo $\pam$ is a supersingular elliptic curve,
that is, it has no point of order $p$ over the algebraic closure of the residue field.
\end{lemma}
\begin{proof}
The invariants of $\mathcal E$ are $c_4=a_1^4-48$, $\Delta=a_1^4-64$, so
$(a_1^4-48)^3=j\cdot(a_1^4-64)$. As $j$ is an algebraic integer~\cite[Chapter~II, Theorem~6.1]{SilvermanAT},
so is $a_1$. If a prime above an odd prime divided $a_1^4-64$ it would divide $a_1^4-48$, hence
$16$. The curve has complex multiplication by $K$ and good reduction, and $p$ does not split in
$K$; by Deuring's criterion~\cite[Chapter~13, Theorem~12]{LangEF} the reduction is supersingular (for a curve over a field
of characteristic $p$, having no point of order $p$ over the algebraic closure is equivalent to
the other usual definitions of supersingularity~\cite[Chapter~V, Theorem~3.1(a)]{SilvermanAEC}).
\end{proof}
For the admissible point $t$, $a_1$ is a unit at the places above $2$ as well, so that $a_1^4-64$
is a unit at every prime (Lemma~\ref{pf:two}).

\emph{The ring $R$.} Let $R=\pao[[y]]\subset\paCp[[z]]$ be the ring of the series $\sum c_ny(z)^n$
with $c_n\in\pao$; it is closed under coefficientwise limits. Solving \eqref{pa:eq-curve} for $x$
by iteration, $x=y^2+a_1xy-x^3\in y^2(1+yR)$; hence $R$ contains $x$ and, by
\eqref{pa:eq-curve}, is stable under $d/dz$. Geometrically $R$ is the completion of
$\mathcal E_{/\pao'}$ along the section $P_0$, for $\pao'\subset\pao$ a discrete valuation ring
containing $a_1$, tensored with $\pao$.

\emph{Endomorphisms.} Let $\alpha\in\auO_K$, $\alpha\equiv1\pmod{2\auO_K}$, $\alpha\ne0$. Since
$(\alpha-1)t\in\Lambda$, multiplication by $\alpha$ on $\C/\Lambda$ fixes $P_0$ and is $z\mapsto
\alpha z$ in the coordinate $z$. The smooth Weierstrass scheme $\mathcal E_{/\pao'}$ is the N\'eron
model of its generic fibre~\cite[\S1.2, Definition~1, p.~12, and Proposition~8, p.~15]{BLR} (\cite[IV, Corollary~6.3, p.~329, and the N\'eron mapping property, p.~319]{SilvermanAT}), so $[\alpha]$ extends to it and preserves the
section $P_0$. Therefore
\begin{equation}\label{pa:eq-endo}
G\in R\ \Longrightarrow\ G(\alpha z)\in R\qquad(\alpha\in1+2\auO_K).
\end{equation}
Let $K_p$ be the closure of $K$ in $\paCp$. The coefficients of $G(\alpha z)$ in powers of $y$ are
polynomials in $\alpha$, and $1+2\auO_K$ is dense in $\auO_{K_p}$ because $p$ is odd; so
\eqref{pa:eq-endo} holds for all $\alpha\in\auO_{K_p}$. In particular it holds for the group
$\mu_{q-1}\subset\auO_{K_p}$ of roots of unity of order prime to $p$. Note that $\pi\equiv1
\pmod{2\auO_K}$ in both cases, since $(\sqrt{-d}-1)/2\in\auO_K$.

\begin{lemma}[The parameter]\label{pa:param}
Let $T(z)=\frac1{q-1}\sum_{\zeta\in\mu_{q-1}}\zeta^{-1}y(\zeta z)$, the part of $y$ in the degrees
$\equiv1\pmod{q-1}$.
\begin{enumerate}
\item[(a)] $T\in y+y^2R$, so $R=\pao[[T]]$; and $T(\pi z)=\pi T+\sum_{k\ge1}b_kT^{1+k(q-1)}$ with
$b_k\in\pao$ and $b_1\in\pao^\times$.
\item[(b)] Let $\pavp^{q-1}=-\pi/b_1$; then $v(\pavp)=\rho$. The series $Y(z)=T(\pavp z)/\pavp$ is
integral, and $Y\equiv\sum_{j\ge0}z^{q^j}\pmod\pi$.
\item[(c)] For every $G\in R$ the series $G(\pavp z)$ is integral, and the coefficients of $G$
itself in the degrees $<q$ in $z$ are integral. In particular this holds for $x$ and $y$.
\end{enumerate}
\end{lemma}
\begin{proof}
(a) By \eqref{pa:eq-endo}, $T\in R$, and $T=z+\dots=y+\dots$. Also $T(\zeta z)=\zeta T(z)$, and
$T(\pi z)\in R=\pao[[T]]$ by \eqref{pa:eq-endo}; writing $T(\pi z)=P(T)$ we get $P(\zeta T)=\zeta
P(T)$, which gives the shape of the expansion, with first coefficient $\pi$. Modulo $\pam$, $P$
is the action of the reduction of $[\pi]$ on a parameter at the point $\widetilde P_0$ of
$\widetilde{\mathcal E}$. This reduction is an isogeny of degree $\auNm(\pi)=q$ (reduction of homomorphisms between curves with good reduction preserves degrees~\cite[Chapter~II, Proposition~4.4]{SilvermanAT}), a power of $p$, and
$\widetilde{\mathcal E}$ has no point of order $p$ (Lemma~\ref{pa:model}). Let $\widetilde O$ be the origin of $\widetilde{\mathcal E}$. The kernel of
$\widetilde{[\pi]}$ consists of points killed by $\widetilde{[\aubar\pi]}\circ\widetilde{[\pi]}=[q]$;
if it contained a point $Q\ne\widetilde O$, then for the least $k$ with $p^kQ=\widetilde O$ the
point $p^{k-1}Q$ would have order $p$. So $\ker\widetilde{[\pi]}=\{\widetilde O\}$. Since
$\pi\equiv1\pmod2$ and $P_0$ has order $2$, $\widetilde{[\pi]}(\widetilde P_0)=\widetilde P_0$, so the
fibre of $\widetilde{[\pi]}$ over $\widetilde P_0$ is $\widetilde P_0+\ker\widetilde{[\pi]}=\{\widetilde P_0\}$. For a nonconstant morphism $\phi$ of smooth projective curves the
ramification indices over any point add up to $\deg\phi$~\cite[Chapter~II, Proposition~2.6(a)]{SilvermanAEC}; hence the
ramification index of $\widetilde{[\pi]}$ at $\widetilde P_0$ is $q$, which says exactly that
$P\bmod\pam$, the pull-back of the uniformiser $T$, has order exactly $q$ in $T$. Hence $b_1$ is a unit.

(b) Replacing $z$ by $\pavp z$ in (a) and dividing by $\pavp$,
\[
Y(\pi z)=\pi\bigl(Y-Y^q+\pi E(Y)\bigr),\qquad E(Y)=\sum_{k\ge2}b_k(-b_1)^{-k}\pi^{k-2}\,
Y^{1+k(q-1)}\in Y^{2q-1}\pao[[Y]].
\]
Write $Y=\sum_{n\ge1}c_nz^n$, $c_1=1$. Comparing the coefficients of $z^n$, $n\ge2$:
$(1-\pi^{n-1})c_n=[Y^q]_n-\pi[E(Y)]_n$, and the right side involves only $c_k$ with $k<n$. By
induction all $c_n$ are integral, and $c_n-[Y^q]_n\in\pi\pao$ for $n\ge2$, that is,
$Y\equiv z+Y^q\pmod\pi$. Iterating with (E3) gives $Y\equiv\sum_jz^{q^j}$.

(c) If $G=\sum g_nT^n$ with $g_n\in\pao$, then $G(\pavp z)=\sum g_n\pavp^nY^n$ is integral. By
(b), $T(z)=\sum_nc_n\pavp^{1-n}z^n$ with $\pi\mid c_n$ for $1<n<q$, and
$v(\pi)-(n-1)\rho=(q-n)\rho>0$; so $T\equiv z$ modulo $\pam$ and $z^q$, and the second claim
follows.
\end{proof}

\emph{The theta function.} Let $H(u,w)=\boldsymbol\pi\,\aubar uw/A$, where
$\boldsymbol\pi=3.14\ldots$ (elsewhere $\pi$ is $\pi_p$), and $s_2=\lim_{s\to0^+}\sum_{\omega\in\Lambda\setminus0}\omega^{-2}
|\omega|^{-2s}$. The function $\vartheta(z)=e^{-s_2z^2/2}\sigma(z)$ satisfies
$\vartheta(z+\omega)=\varepsilon(\omega)e^{H(\omega,z+\omega/2)}\vartheta(z)$ for $\omega\in\Lambda$,
with $\varepsilon(\omega)=1$ if $\omega\in2\Lambda$ and $-1$ otherwise~\cite[\S1.2]{BannaiKobayashi}.
Put
\[
\theta(z)=\frac{\vartheta(z+t)}{\vartheta(t)}\,e^{-H(t,z)},\qquad g=-\log\theta\in z\,\C[[z]].
\]
Then $\theta(0)=1$, the zeros of $\theta$ are simple and are the points of $t+\Lambda$, and
\begin{equation}\label{pa:eq-theta}
\theta(z+\omega)=\chi(\omega)\,e^{H(\omega,z+\omega/2)}\,\theta(z),\qquad
\chi(\omega)=\varepsilon(\omega)e^{H(\omega,t)-H(t,\omega)}\in\{\pm1\},
\end{equation}
where $\chi(\omega)$ depends only on $\omega$ modulo $2\Lambda$ (as $2t\in\Lambda$ and the
imaginary part of $H$ is $\boldsymbol\pi$ times an integral form on $\Lambda$). The function
$\theta(-z)$ satisfies \eqref{pa:eq-theta} and has the same zeros, so $\theta$ is even. Finally
\begin{equation}\label{pa:eq-g2}
g''=x+2g_2,\qquad g_2=\tfrac12(s_2+\wp(t))\ \text{the coefficient of }z^2\text{ in }g,
\end{equation}
because $-(\log\sigma)''=\wp$.

\begin{lemma}[Units]\label{pa:units}
\begin{enumerate}
\item[(a)] For $\alpha\in1+2\auO_K$, $\alpha\ne0$, the function $f_\alpha(z)=\theta(\alpha z)/
\theta(z)^{\auNm\alpha}$ is a polynomial $R_\alpha(x)$ in $x$ of degree $(\auNm\alpha-1)/2$ with
$R_\alpha\in\pao[x]$, $R_\alpha(0)=1$. In particular $f_\alpha\in1+yR$.
\item[(b)] For $\zeta\in\mu_{q-1}$, $\exp\bigl(\zeta^{p+1}g(z)-g(\zeta z)\bigr)\in1+yR$.
\item[(c)] $\pi g_2\in\pao$, and $\pi g_n\in\pao$ for all $n<q$, where $g=\sum g_nz^n$.
\end{enumerate}
\end{lemma}
\begin{proof}
(a) By \eqref{pa:eq-theta}, $\chi(\alpha\omega)=\chi(\omega)=\chi(\omega)^{\auNm\alpha}$
($\auNm\alpha$ is odd) and $H(\alpha\omega,\alpha w)=\auNm\alpha\cdot H(\omega,w)$; so $f_\alpha$ is
$\Lambda$-periodic. In the variable $u=z+t$ its zeros are the points of
$\alpha^{-1}\Lambda\setminus\Lambda$ (as $\alpha t\equiv t$), all simple, its only pole is $u=0$,
of order $\auNm\alpha-1$, and it is even. Hence $f_\alpha=\prod_s(1-x/x_s)$, the product over the
non-zero points $s$ of $\alpha^{-1}\Lambda/\Lambda$ up to sign, with $x_s=\wp(s)-\wp(t)$. We claim
$|x_s|\ge1$. If $|x_s|<1$, then the point $s$ of $\mathcal E$ has $|y(y+a_1x)|=|x(1+x^2)|<1$, hence
$|y|<1$, and reduces to $\widetilde P_0$; as reduction is a homomorphism, $s-t$ lies in its kernel
$E_1$, and so does $[\auNm\alpha](s-t)=-t$ ($\auNm\alpha$ is odd and kills $s$). But $t$ does not
reduce to the origin. This proves the claim, and (a) follows since $x\in y^2R$.

(b) Choose $\alpha_n\in1+2\auO_K$ with $\alpha_n\to\zeta$ in $\auO_{K_p}$. Then
$\auNm\alpha_n=\alpha_n\aubar{\alpha_n}\to\zeta\cdot\zeta^{p}$ if $p$ is inert (the conjugation of
$K_p$ is the Frobenius) and $\to\zeta^2=\zeta^{p+1}$ if $p=d$ ($\zeta\in\Z_p$). The coefficients of
$f_{\alpha_n}=\exp(\auNm\alpha_n\,g(z)-g(\alpha_nz))$ are polynomials in
$(\alpha_n,\auNm\alpha_n)$, and $1+yR$ is closed.

(c) By (a) and Lemma~\ref{pa:param}(c), the coefficient of $z^2$ in $\log f_\alpha=\auNm\alpha\,
g(z)-g(\alpha z)$ is integral; it is $\alpha(\aubar\alpha-\alpha)g_2$. For
$\alpha=2+\sqrt{-d}$ this is $-2\sqrt{-d}\,(2+\sqrt{-d})g_2$, and $2(2+\sqrt{-d})$ is a unit at
$p$: if $p=d$ its norm $4(4+d)$ is prime to $p$, and if $p$ is inert then $p$ is a prime element
of $\auO_K$ that does not divide $2+\sqrt{-d}$. So $\sqrt{-d}\,g_2$ is integral, and hence
$\pi g_2$ is. For $3\le n<q$, \eqref{pa:eq-g2} gives
$n(n-1)g_n=x_{n-2}$, the coefficient of $z^{n-2}$ in $x$, integral by Lemma~\ref{pa:param}(c); and
$v(n(n-1))\le v(\pi)$ for $n<q$. Finally $g_1=0$.
\end{proof}

\begin{proposition}[Frobenius congruence]\label{pa:frob}
$q\,g(z)-g(\pi z)=\pi H$ for some $H\in yR$.
\end{proposition}
\begin{proof}
By Lemma~\ref{pa:units}(a), $\exp(q\,g(z)-g(\pi z))=f_\pi=R_\pi(x)$ with $R_\pi=1+\sum_{i=1}^{(q-1)/2}
r_ix^i$, $r_i\in\pao$. We show $\pi\mid r_i$ for all $i$. The coefficient of $z^n$ in
$q\,g(z)-g(\pi z)$ is $(q-\pi^n)g_n=(q/\pi-\pi^{n-1})\cdot\pi g_n$; for $2\le n<q$ it is divisible
by $\pi$ by Lemma~\ref{pa:units}(c), and it vanishes for $n=1$. By (E2) with $\beta=\pi$,
$R_\pi(x(z))\equiv1$ modulo $\pi$ in all degrees $<q$ in $z$. By Lemma~\ref{pa:param}(c) and
$x=z^2+\dots$, the series $x(z)^i$ is integral in degrees $<q$ with lowest term $z^{2i}$, and
$2i\le q-1$. Comparing the coefficients of $z^2,z^4,\dots,z^{q-1}$ in turn gives $\pi\mid r_1$, then
$\pi\mid r_2$, and so on. Thus $R_\pi(x)=1+\pi H_0$ with $H_0\in x\pao[x]\subset yR$, and
$q\,g(z)-g(\pi z)=\log(1+\pi H_0)\in\pi\,yR$ by (E2).
\end{proof}

\begin{theorem}[Normal form on one curve]\label{pa:onecurve}
Let $\pavp$ be as in Lemma~\ref{pa:param}(b) and $c_\pavp=\pavp^{p+1}/p$. There are integral series
$W,\Xi_1\in\pao[[z]]$, $\Xi_1(0)=0$, and $\lambda\in\pao$ such that
\[
\theta(\pavp z)=W(z)\exp\bigl(c_\pavp\,\Xi_1(z)\bigr),\qquad
\Xi_1\equiv\sum_{e\in\paI}\sum_{j\ge0}\lambda_{e,j}\bigl(z^{p^e}\cdot z\bigr)^{p^j}\pmod\pi ,
\]
where $\lambda_{e,j}=\lambda$ if $p$ is inert, and $\lambda_{0,0}=\lambda/2$,
$\lambda_{0,j}=\lambda$ $(j\ge1)$, $\lambda_{e,j}=2\lambda$ $(e\ge1)$ if $p=d$.
\end{theorem}
\begin{proof}
\emph{Step 1: averaging.} For a series $G$ let $G^\flat=\frac1{q-1}\sum_{\zeta\in\mu_{q-1}}
\zeta^{-(p+1)}G(\zeta z)$ be its part in the degrees $\equiv p+1\pmod{q-1}$. Then
\[
\exp(g^\flat-g)=\prod_{\zeta}\Bigl(\exp\bigl(\zeta^{p+1}g(z)-g(\zeta z)\bigr)\Bigr)^{\epsilon_\zeta},
\qquad\epsilon_\zeta=-\frac{\zeta^{-(p+1)}}{q-1}\in\Z_p ,
\]
because $\zeta^{p+1}\in\mu_{p-1}\subset\Z_p$. A $p$-adic integral power of an element of $1+yR$ is
in $1+yR$, so $U:=\theta\exp(g^\flat)\in1+yR$ by Lemma~\ref{pa:units}(b). Put $W(z)=U(\pavp z)$,
integral by Lemma~\ref{pa:param}(c), and $\Xi_1(z)=-g^\flat(\pavp z)/c_\pavp$. Then
$\theta(\pavp z)=W\exp(c_\pavp\Xi_1)$.

\emph{Step 2: a functional equation.} The operation $\flat$ commutes with $z\mapsto\pi z$ and
preserves $R$, by \eqref{pa:eq-endo}; so Proposition~\ref{pa:frob} gives
$q\,g^\flat(z)-g^\flat(\pi z)=\pi H^\flat$ with $H^\flat\in R$. Expanding $H^\flat$ in powers of
$T$, which has only degrees $\equiv1\pmod{q-1}$, only the powers $T^n$ with $n\equiv p+1$ occur.
Let $i_0$ be the least positive integer $\equiv p+1\pmod{q-1}$: $i_0=p+1$ if $p$ is inert, $i_0=2$
if $p=d$ (then $q-1=d-1\ge6$). So $H^\flat=h\,T^{i_0}+T^{i_0+q-1}H_1$ with $h\in\pao$, $H_1\in R$.
Replace $z$ by $\pavp z$, use $T(\pavp z)=\pavp Y$ and divide by $-q\,c_\pavp$:
\[
\Xi_1(z)-q^{-1}\,\Xi_1(\pi z)=-u\bigl(h\,Y^{i_0}+\pavp^{q-1}Y^{i_0+q-1}H_1(\pavp z)\bigr),\qquad
u=\frac{\pi\,\pavp^{i_0}}{q\,c_\pavp}.
\]
One checks $v(u)=1-v(\pi)+(i_0-p-1)\rho=0$ in both cases, and $v(\pavp^{q-1})=v(\pi)$. With
$\lambda=-uh$, for every $n$,
\begin{equation}\label{pa:eq-fe}
\bigl(1-\pi^n/q\bigr)\,[\Xi_1]_n\equiv\lambda\,[Y^{i_0}]_n\pmod{\pi\pao}.
\end{equation}

\emph{Step 3: solution.} $\Xi_1$ has only terms of degree $n\ge i_0$, $n\equiv i_0\pmod{q-1}$. For
$n\ge3$, $v(\pi^n/q)=(n-2)v(\pi)\ge v(\pi)$, so \eqref{pa:eq-fe} shows that $[\Xi_1]_n$ is integral
and $\equiv\lambda[Y^{i_0}]_n$. If $p$ is inert, all $n$ are $\ge p+1\ge4$; thus
$\Xi_1\equiv\lambda Y^{p+1}$, and by Lemma~\ref{pa:param}(b) and (E3)
\[
Y^{p+1}\equiv\Bigl(\sum_{i\text{ odd}}z^{p^i}\Bigr)\Bigl(\sum_{i\text{ even}}z^{p^i}\Bigr)
=\sum_{e\text{ odd}}\ \sum_{j\ge0}\bigl(z^{p^e}z\bigr)^{p^j},
\]
each pair (odd $i$, even $i'$) being written once as $\{j,j+e\}$. If $p=d$, then for $n=2$ one has
$\pi^2/q=-1$ and $[Y^2]_2=1$, so $[\Xi_1]_2\equiv\lambda/2$; hence
$\Xi_1\equiv\lambda(Y^2-\tfrac12z^2)$ and
$Y^2\equiv\sum_iz^{2p^i}+2\sum_{i<i'}z^{p^i+p^{i'}}$, which gives the stated coefficients with
$j=i$, $e=i'-i$.
\end{proof}

\begin{remark}
For results of this kind for one curve at supersingular primes (radius of convergence of the
$p$-adic theta function, divisibility of its coefficients) see \cite{BannaiKobayashi} and
\cite{KatzDivisibilities}; we do not use them. Lemma~\ref{pa:param} says that $[\pi]$ is, in the
parameter $T$, an endomorphism of the kind studied by Lubin and Tate~\cite{LubinTate}.
\end{remark}

\subsection{Products of theta functions along a resolution of the identity}\label{pa:ss-prod}
We identify $\auPl$ with $\operatorname{Gal}(E/K)=(\Z/q_1)^\times$: $\tau_\nu=\tau_1\circ\sigma_\nu$,
where $\tau_1$ is the given embedding and $\sigma_\nu(\zeta_{q_1})=\zeta_{q_1}^{\nu}$; then
$u(\nu)\tau=\tau\circ\sigma_\nu$. Each $\tau$ gives $E\to\C\to\paCp$.

\begin{lemma}[Powers and partners]\label{pa:partner}
Let $e\in\paI$ and let $a\in E$ be integral at all primes above $p$. For $\tau\in\auPl$ and
$\tau'=u(-p^e)\tau$,
\[
\bigl(\aubar{\tau(a)}\bigr)^{p^e}\equiv\tau'(a)\pmod{\pi\pao}.
\]
\end{lemma}
\begin{proof}
Multiplying by an integer prime to $p$ we may assume $a\in\auO_E=\auO_K[\zeta_{q_1}]$ (the
discriminants of $K$ and of $\Q(\zeta_{q_1})$ are coprime). Write $a=\sum_ic_i\zeta_{q_1}^i$ with
$c_i\in\auO_K$. If $p$ is inert and $e$ is odd, $c_i^{p^e}\equiv\aubar{c_i}\pmod p$; if $p=d$,
$c_i^{p^e}\equiv c_i\equiv\aubar{c_i}\pmod{\sqrt{-d}}$. By (E3), in both cases
$a^{p^e}\equiv\sum_i\aubar{c_i}\,\zeta_{q_1}^{ip^e}=\aubar{\sigma_{-p^e}(a)}\pmod{\pi\auO_E}$.
Apply this to $\aubar a$ and then $\tau$: $\tau(\aubar a)^{p^e}\equiv\tau(\sigma_{-p^e}(a))=
\tau'(a)$, and $\tau(\aubar a)=\aubar{\tau(a)}$ because $E$ is a CM field.
\end{proof}

Let $\mathsf C$ be a finite set. For $u\in E^{\mathsf C}$ let $u^\dagger$ be the linear form
$u^\dagger(z)=\sum_{\tau\in\auPl,\,i\in\mathsf C}\aubar{\tau(u_i)}\,z_{\tau,i}$. A \emph{resolution
of the identity} is a finite family $u_1,\dots,u_J\in E^{\mathsf C}$ with multiplicities
$m_1,\dots,m_J\in\N$ such that
\begin{equation}\label{pa:eq-res}
\sum_{j}m_j\,\tau(u_{j,i})\cdot u_j^\dagger(z)=z_{\tau,i}\qquad\text{for all }\tau\in\auPl,\
i\in\mathsf C .
\end{equation}

\begin{proposition}[Normal form of theta products]\label{pa:product}
Let $(u_j,m_j)$ be a resolution of the identity by vectors whose coordinates are integral at all
primes above $p$, let $\mu_j=u_j^\dagger$, and let $G$ be an integral power series in the series
$y(\mu_j(z))$. Then, for every $c\in\pao$, the series
\[
F(z)=F_0(cz),\qquad F_0(z)=\prod_j\theta(\mu_j(z))^{m_j}\cdot G(z),
\]
has the normal form \eqref{pa:eq-nf}, with $\pavp$ as in Lemma~\ref{pa:param}(b), the index set
$\paI$ and the partner maps $\mathrm{pr}_e(\tau)=u(-p^e)\tau$.
\end{proposition}
\begin{proof}
By the remark after \eqref{pa:eq-nf} we may take $c=1$. The forms $\mu_j$ have integral
coefficients. By Theorem~\ref{pa:onecurve} and Lemma~\ref{pa:param}(c),
\[
F_0(\pavp z)=W_0(z)\exp\Bigl(c_\pavp\sum_jm_j\,\Xi_1(\mu_j(z))\Bigr),\qquad
W_0=\prod_jW(\mu_j(z))^{m_j}\cdot G(\pavp z)\ \text{integral}.
\]
It suffices to compute the exponent modulo $\pi$. By Theorem~\ref{pa:onecurve}, (E3) and
$m^{p}\equiv m\pmod p$,
\[
\sum_jm_j\,\Xi_1(\mu_j)\equiv\sum_{e,k}\lambda_{e,k}\Bigl(\sum_jm_j\,\mu_j^{\,p^e}\mu_j\Bigr)^{p^k}
\pmod\pi .
\]
By (E3) and Lemma~\ref{pa:partner}, $\mu_j^{\,p^e}\equiv\sum_{\tau,i}\tau'(u_{j,i})\,
z_{\tau,i}^{\,p^e}$ with $\tau'=u(-p^e)\tau$; hence, by \eqref{pa:eq-res} at the place $\tau'$,
\[
\sum_jm_j\,\mu_j^{\,p^e}\mu_j\equiv\sum_{\tau,i}z_{\tau,i}^{\,p^e}\sum_jm_j\,\tau'(u_{j,i})\,
u_j^\dagger(z)=\sum_{\tau,i}z_{\tau,i}^{\,p^e}\,z_{\tau',i}=\Gamma_e\pmod\pi .\qedhere
\]
\end{proof}

\subsection{The composite lift}\label{pa:ss-divis}
We now prove {\cite[Proposition~14.20]{P46}} as far as stated in the introduction of this section.
Let $\ell$ be a prime and $\C\to\C_\ell$ a homomorphism (when $\ell=p$ is non-split, the
conventions of \S\ref{pa:ss-conv} apply). Let $0\le j<m$. The weight $\lambda_{j+1}=\lambda(X_{j+1})$
has, at $\tau$ (see {\cite[\S14.3]{P46}}), $a\cdot\mathbf 1_{\tau\in\auPa}$ columns of height $1$,
$\kappa_1(\tau)-\lfloor n_1/2\rfloor$ columns of height $k$ and $\Delta-(t+1)$ columns of each of
the heights $n_1,\dots,n_j$; all heights are $<n_{j+1}$ and $\le n_j$. In the notation of
\S\ref{pa:ss-count}, with $n'=n_{j+1}$, $n=n_j$:
\begin{equation}\label{pa:eq-diag}
a_\tau=a\,\mathbf 1_{\tau\in\auPa},\qquad\kappa(\tau)\le a'_\tau\le\kappa(\tau)+x^{60},\qquad
k\,\kappa(\tau)\le l_\tau\le k\,\kappa(\tau)+x^{60},
\end{equation}
because $a'_\tau-\kappa(\tau)=\tfrac32\Delta-\lfloor n_1/2\rfloor+j(\Delta-t-1)$ and
$l_\tau-k\kappa(\tau)=k(\tfrac32\Delta-\lfloor n_1/2\rfloor)+(\Delta-t-1)\sum_{i\le j}n_i$ are
non-negative ($n_1=k+t+1\le\Delta$) and at most $x^{60}$ ($\Delta=2r^8\le2\cdot3^8x^{40}$,
$n_m\le40x^6$). We write $\paPred_j(p)$ for the
predicted divisibility of this shape, $B_j$ for its number of boxes, and
\[
\paloss_p=f\,(\log b)^2\,\frac ap\,\epsilon_m(p)\qquad(p\in\mathcal N),
\]
so that $\auerr_p=2\paloss_p$ in {\cite[(45)]{P46}}.

\begin{lemma}[The kernels]\label{pa:kernel}
Let $\ell\ne N$, let $\vartheta$ be a twist with $\ell\nmid\auNm Q_\vartheta$, $\Psi$ a chain datum
of $\vartheta$, and $\Psi'$ as in {\cite[Proposition~14.25]{P46}}, so that
$B_{\vartheta,\Psi}=\operatorname{res}\circ C_{\vartheta,\Psi'}$.
\begin{enumerate}
\item[(K0)] For $\phi\in\auS$ and $g_m\in\auU(X_m)(\auA_{E^+,f})$,
\[
(C_{\vartheta,\Psi'}\phi)(g_m)=\sum_{g_0,\dots,g_{m-1}}\ \sum_x\Psi'(x)\;\theta_{m-1}
(\phi^\circ_{m-1,x_{m-1}};g_{m-1},g_m)\cdots\theta_0(\phi^\circ_{0,x_0};g_0,g_1)\,\phi(g_0),
\]
where $g_i$ runs over representatives of the classes of the level of $X_i$ and $\theta_i$ is the
kernel of {\cite[\S14.8]{P46}}.
\item[(K1)] For every $x$ in the support of $\Psi'$, every $j<m$ and all $g,g'$ there is a power
series $F=F_{x,j,g,g'}$ in the variables $z_{\tau,b,a}$ ($\tau\in\auPl$, $b\le n_{j+1}$, $a\le n_j$)
such that the matrix of $g'^{-1}_\ell\circ\theta_j(\phi^\circ_{j,x_j};g,g')\circ g_\ell$ in the
bases $(e_T)$ of \S\ref{pl:ss-norm} is, up to multiplication on both sides by matrices that are
invertible over $\auO_{\C_\ell}$, the matrix
$\bigl(S_{\lambda_{j+1}}(\partial)_{T',T}F(0)\bigr)_{T',T}$ of jets along the entries of the Schur
matrix. If $\ell\ne d$, all jets $Q(\partial)F(0)$ with $Q\in\auO_{\C_\ell}[z]$ are integral
(Proposition~\ref{pk:K12} for $\ell\nmid2q_1M$; Corollary~\ref{pf:K1}(a) for all $\ell\ne d$).
\item[(K2)] If $\ell=p\in\mathcal N$, then $F$ is one of the series of
Proposition~\ref{pa:product}, for $\mathsf C=[n_{j+1}]\times[n_j]$, the resolution of the identity
$u_{k,i}=(\zeta_{q_1}^k/q_1)e_i$ $(k<q_1,\ i\in\mathsf C)$ with multiplicities $q_1$, the lattice
$\Lambda=\Omega_0\auO_K$ and $c=\pi_M\varpi_{\mathrm{vac}}/\mathsf u$.
\item[(K1$_N$)] Let $\ell=N$ and $\iota\colon\C\to\C_N$. Let $z\in\operatorname{Hom}(X_j,X_{j+1})(\auA_{E^+,f})$
be integral at the primes above $N$ that are conjugate to the primes induced by the $\iota\circ\tau$,
and such that $Nz$ is integral at the primes induced by the $\iota\circ\tau$. Then for all $g,g'$ the
matrix of $g'^{-1}_N\circ\theta_j(\phi^\circ_{j,z};g,g')\circ g_N$ in the bases $(e_T)$ has entries in
$\auO_\iota$. (This is Lemma~\ref{pl:kernelN}.)
\end{enumerate}
\end{lemma}
\begin{proof} (K0) is Proposition~\ref{pk:K0}; the form of the matrix in (K1), and (K2), are
Proposition~\ref{pk:K12}; the integrality in (K1) at every prime $\ell\ne d$ and (K1$_N$) are
Corollary~\ref{pf:K1}.\end{proof}


\begin{lemma}[Gauge of a chain]\label{pa:chain}
Assume Lemma~\ref{pa:kernel}. Let $\ell\ne N$, $\ell\nmid\auNm Q_\vartheta$, and let $r_j\in\R$
$(0\le j<m)$ be such that all the jets $S_{\lambda_{j+1}}(\partial)_{T',T}F_{x,j,g,g'}(0)$ of (K1)
have valuation at least $r_j$. Then for every $\phi\in\auS$ integral at the given place,
$B_{\vartheta,\Psi}\phi$ has gauge at most $\ell^{-\sum_jr_j}$.
\end{lemma}
\begin{proof}
That $\phi(g_0)$ has gauge at most $1$ means that $g_{0,\ell}^{-1}\phi(g_0)$ has integral
coordinates. Insert $g_{i,\ell}\,g_{i,\ell}^{-1}$ between consecutive kernels in (K0): by (K1) and
the hypothesis the $j$-th factor becomes a matrix with entries of valuation $\ge r_j$. The
coefficients $\Psi'(x)$ are integral at $\ell$ ({\cite[Proposition~14.25]{P46}}, as $\ell\nmid
N\auNm Q_\vartheta$), and the valuation is ultrametric; so $g_{m,\ell}^{-1}(C_{\vartheta,\Psi'}
\phi)(g_m)$ has coordinates of valuation $\ge\sum_jr_j$. The restriction is the composition of the
evaluation at $\operatorname{diag}(y,1)$, of the branching projection, which is
$\auGL_{n_m-1}$-equivariant and has an integral matrix in the standard bases
({\cite[Proposition~14.24]{P46}}), and of the multiplication by $\prod_\tau\tau(\aubar{\pi_N})^{-
n_\tau}$, an $\ell$-adic unit because $\ell\ne N$.
\end{proof}

\begin{lemma}[One lift gains its share]\label{pa:gain}
Let $p\in\mathcal N$ and $0\le j<m$. Call $(\tau,e)$ a \emph{collision} if $e\in I_p$ and both
$\tau$ and $u(-p^e)\tau$ are $a$-places, and let $C_p=\sum_{(\tau,e)\text{ collision}}a/p^e$. Then
\[
\paPred_j(p)-f\,n_{j+1}^3\,(1+\log_pB_j)\ \ge\ \frac{\augain_p}m-C_p-\frac{f\,a\,(\log b)^2}{p\,m}.
\]
\end{lemma}
\begin{proof}
By {\cite[Definition~14.8]{P46}}, $\augain_p/m=v(\pi)\sum_{e\in I_p}\bigl(\sum_{\tau\in\auPa}
(a/p^e-\kappa(\tau'))^++k\sum_\tau(\kappa(\tau)/p^e-\kappa(\tau'))^+\bigr)$, $\tau'=u(-p^e)\tau$.
All these terms vanish if $p^e\ge x_1$, since then $a/p^e\le a/x_1\le\kappa$ and
$\kappa(\tau)/p^e\le b/x_1<a/x_1$. We compare with the terms of $\paPred_j(p)$, using
\eqref{pa:eq-diag} and writing $s=x^{60}$. Let $e\in I_p$.

(1) Let $\tau\in\auPa$. If $\tau'\notin\auPa$, then $(a_\tau/p^e-a_{\tau'}-a'_{\tau'})^+\ge
(a/p^e-\kappa(\tau')-s)^+\ge(a/p^e-\kappa(\tau'))^+-s$. If $\tau'\in\auPa$ (a collision), the left
side is $\ge0\ge(a/p^e-\kappa(\tau'))^+-a/p^e$.

(2) For every $\tau$, $(l_\tau/p^e-l_{\tau'})^+\ge(k\kappa(\tau)/p^e-k\kappa(\tau')-s)^+\ge
k(\kappa(\tau)/p^e-\kappa(\tau'))^+-s$.

(3) By \eqref{pa:eq-const}, $v(\pi)\sum_{e\in I_p}a'_\tau/p^e=\gamma a'_\tau\le\frac2p(b+s)$.

\noindent In (1) and (2) the loss $s$ is needed only for the at most $2\log x_1$ exponents with
$p^e<x_1$. If $p=d$ there is the further term $\tfrac12\sum_\tau\paded_0(\tau)$, with
$\tau'=u(-1)\tau$. The places $\tau$ and $u(-1)\tau$ are neither both $a$-places nor both
$b$-places (their classes modulo $q_0$ are $\pm1$ and $\mp1$); so
$\min(a_\tau,a_{\tau'}+a'_{\tau'})\le b+s$, $a'_\tau\le b+s$ and $\min(l_\tau,l_{\tau'})\le
ka/x_1+s$. Altogether the left side of the lemma is at least $\augain_p/m-C_p$ minus
\[
f\Bigl(4s\log x_1+\frac{2(b+s)}p+\mathbf 1_{p=d}\bigl(b+2s+\tfrac{ka}{x_1}\bigr)+n_m^3(1+\log_pB_j)
\Bigr).
\]
Here $p<x_1$, $m\le2\sqrt x$, $a=xb$, $B_j\le b^3$, $d\le\log\log b$ by (D2), and $s$, $n_m$,
$k$ are polynomial in $x$; so each term is at most a quarter of $a(\log b)^2/(pm)\ge
\sqrt x\,b\,(\log b)^2/(2p)$, by a comparison of sizes.
\end{proof}

\begin{lemma}[Collisions]\label{pa:coll}
Let $p\in\mathcal N$ and let $o_0$ be the order of $-p$ in $(\Z/q_0)^\times$. If $p=d$ or $o_0$ is
even, there are no collisions. Otherwise the collisions are the pairs $(\tau,e)$ with
$\tau\in\auPa$, $e$ odd and $o_0\mid e$, and $C_p\le2fa/p^{o_0}$. If $p$ is not exceptional, then
$m\,C_p\le\paloss_p$. In all cases $m\,C_p\le j_p:=\lceil2mfa/p\rceil$.
\end{lemma}
\begin{proof}
For $\tau\in\auPa$ the place $u(-p^e)\tau$ is an $a$-place if and only if $-p^e\equiv1\pmod{q_0}$.
If $p=d$: $d\equiv q_0/2-1$ and $16\mid q_0$ give $d^2\equiv1\pmod{q_0}$, so $-d^e$ is $-1$ or $-d$
modulo $q_0$, and neither is $1$. If $p$ is inert, $e$ is odd and $-p^e=(-p)^e$, which is $1$
modulo $q_0$ exactly when $o_0\mid e$; this needs $o_0$ odd. Then
$C_p\le f\,a\sum_{i\ge0}p^{-o_0(1+2i)}\le2fa/p^{o_0}\le2fa/p$, which gives the last claim. Let $p$
be inert and not exceptional, with collisions. Either $p\equiv\pm1\pmod m$, and then
$\paloss_p=mf(\log b)^2a/p\ge2mfa/p$; or $o_0>\mathfrak r_p$, and then $p^{o_0}\ge p\,m$ and
$m\,C_p\le2fa/p\le\paloss_p$.
\end{proof}

\subsection*{Exceptional primes}
Let $p\in\mathcal N$ be exceptional. Each $\tau\in\auPl$ induces a prime $\mathfrak l_\tau$ of $E$
above $p$, the centre of the valuation of $\paCp$ on $\tau(E)$. A \emph{full mark} is a set
$\Lambda$ of primes of $E$ above $p$ with $\Lambda\cap\aubar\Lambda=\emptyset$, each of the form
$\mathfrak l_\tau$ with $\tau\in\auPa$, and containing $\mathfrak l_\tau$ or
$\aubar{\mathfrak l_\tau}$ for every $\tau\in\auPa$.

\begin{lemma}\label{pa:mark}
A full mark exists.
\end{lemma}
\begin{proof}
It suffices that no prime $\mathfrak l$ of $E$ above $p$ equals its conjugate; then one chooses
one prime from each pair $\{\mathfrak l,\aubar{\mathfrak l}\}$ that contains some
$\mathfrak l_\tau$, $\tau\in\auPa$, taking a member of that form. The decomposition group of
$\mathfrak l$ in $\operatorname{Gal}(E/\Q)$ is generated by the Frobenius $\varphi$, which is not
trivial on $K$. If $\aubar{\mathfrak l}=\mathfrak l$, then $c=\varphi^i$ with $i$ odd, so
$c\varphi^i\in\operatorname{Gal}(E/K)$ is trivial; it is the class of $-p^i$, so
$p^i\equiv-1\pmod{q_1}$. Now $m\mid q_1$ ($m$ is a power of $2$ with $m\le2\sqrt x$, while $2^{s+4}\mid q_1$ and $2^s\ge x^{11}/(3l_0(s))$ by (S3) of {\cite[Lemma~13.4]{P46}}), and $m$ is a power of $2$, so $(\Z/m)^\times$ is a $2$-group and raising to the
odd power $i$ is a bijection of it; as $(-1)^i=-1$, we get $p\equiv-1\pmod m$, which is excluded
for an exceptional prime.
\end{proof}

For $\mathfrak l\in\Lambda$ the group $\auU(X_0)$ at the place of $E^+$ below $\mathfrak l$ is
$\auGL_k(E_{\mathfrak l})$. Let $T_{\mathfrak l}$ be the Hecke operator on $\auS$ of the double coset of
$\operatorname{diag}(p1_{\mathfrak r_p},1_{k-\mathfrak r_p})$, $(T_{\mathfrak l}\varphi)(g)=\sum_h
\varphi(gh)$, in the identification by the component at $\mathfrak l$ (\S\ref{pl:ss-norm}, (N3)), and
\[
Y_{\mathfrak l}=p^{A_{\mathfrak l}}\,T_{\mathfrak l},\qquad
A_{\mathfrak l}=a\cdot\#\{\tau\in\auPa:\mathfrak l_\tau=\mathfrak l\},\qquad
Y_\Lambda=\prod_{\mathfrak l\in\Lambda}Y_{\mathfrak l}.
\]
Write $\auS_\iota$ for the module of source forms integral at the given place.

\begin{lemma}[Valuation on the image of the Hecke operators]\label{pa:hecke}
Let $p$ be exceptional, $\Lambda$ a full mark, $\vartheta$ a twist with $p\nmid\auNm Q_\vartheta$
and $\Psi$ a chain datum. For every $\phi\in\auS_\iota$ the form
$B_{\vartheta,\Psi}(Y_\Lambda^{\,m}\phi)$ has gauge at most $p^{\,2\paloss_p-\augain_p}$.
\end{lemma}
\noindent Proved in \S\ref{pl:ss-hecke}, from Lemma~\ref{pa:kernel} (K1), (K2).

\begin{lemma}[Elementary divisors]\label{pa:eldiv}
Let $p$ be exceptional, $\Lambda$ a full mark and $\mathfrak l\in\Lambda$. Then
$Y_{\mathfrak l}\auS_\iota\subseteq\auS_\iota$, and for every integer $j_0\ge0$ there is an integer
$i\le\bigl(12A_{\mathfrak l}/(mp)+j_0/(m^2f)\bigr)\auDX$ with
$\aucol\bigl(\auS_\iota,\,Y_{\mathfrak l}\auS_\iota+p^{j_0}\auS_\iota\bigr)\mid p^{\,i}$.
\end{lemma}
\noindent Proved in {\cite[\S2.7]{P50}}, from {\cite[Theorem~2.2]{P50}}; the three inputs of that theorem
({\cite[Propositions~2.10, 2.12 and~2.13]{P50}}) are proved in {\cite[\S3]{P50}}, the
normalisation used is justified in Remark~\ref{pl:N}, and weak approximation is
Corollary~\ref{pl:wacor}. Nothing used in that proof depends on the present lemma, which enters only case (iv) of the theorem below.

\subsection*{The level prime}
\begin{lemma}[The prime $N$]\label{pa:level}
Let $\ell=N$, and let $U_N$, $L_0$ and $\mathcal A'=\{\varphi\in\auS:U_N^{\,m/2}\varphi\text{ is
integral for }L_0\}$ be as in \S\ref{pl:ss-level}. Then {\rm(a)} and {\rm(b)} of Lemma~\ref{pl:level}
hold.
\begin{enumerate}
\item[(a)] for every twist $\vartheta$, every chain datum $\Psi$ and every
$\phi\in\mathcal A'\cap\auSram$, the form $B_{\vartheta,\Psi}\phi$ has gauge at most $N^{2mfb}$;
\item[(b)] $\aucol(\auS_\iota,\mathcal A'\cap\auS_\iota)\mid N^{\,i}$ for some
$i\le m\,f\,b\,\auDX$.
\end{enumerate}
\end{lemma}
\noindent Proved in \S\ref{pl:ss-level}, from Lemma~\ref{pl:kernelN} (which is
Corollary~\ref{pf:K1}(b)) and Lemma~\ref{ps:branching}(c) (which is Proposition~\ref{pk:branching}).

\subsection*{The theorem}

\begin{theorem}[{\cite[Proposition~14.20]{P46}}]\label{pa:divis}
(The statements used in the proof, Lemma~\ref{pa:kernel} including (K1$_N$),
Lemma~\ref{ps:branching}(c) and {\cite[Proposition~14.25]{P46}}, are proved in \S\ref{pk:ss-kernel}
and Corollary~\ref{pf:K1}, in Proposition~\ref{pk:branching} and in \S\ref{ps:ss-glue}.) Let $\ell$ be a prime with a
homomorphism $\C\to\C_\ell$, and let $\auS_\iota$ be the module of source forms integral there.
\begin{enumerate}
\item[(i)] If $\ell\ne N$ and $\ell\notin\mathcal N$, then for every twist $\vartheta$ with
$\ell\nmid\auNm Q_\vartheta$, every chain datum $\Psi$ and every $\phi\in\auS_\iota$ the form
$B_{\vartheta,\Psi}\phi$ is integral.
\item[(ii)] If $\ell=p\in\mathcal N$ is not exceptional, then under the same conditions
$B_{\vartheta,\Psi}\phi$ has gauge at most $p^{\,2\paloss_p-\augain_p}$.
\item[(iii)] If $\ell=N$, let $\mathcal A'$ be as in Lemma~\ref{pl:level}; then $S_0=\mathcal A'\cap\auS_\iota$ has
$\aucol(\auS_\iota,S_0)\mid N^{\,i}$ with $i\le mfb\,\auDX$, and $B_{\vartheta,\Psi}\phi$ has gauge
at most $N^{2mfb}$ for $\phi\in S_0\cap\auSram$.
\item[(iv)] If $\ell=p\in\mathcal N$ is exceptional, let $\Lambda$ be a full mark and $j_p=\lceil2mfa/p\rceil$. Then
$S_0=Y_\Lambda^{\,m}\auS_\iota+p^{j_p}\auS_\iota$ is contained in $\auS_\iota$,
$\aucol(\auS_\iota,S_0)\mid p^{\,i}$ with $i\le\paloss_p\auDX$, and $B_{\vartheta,\Psi}\phi$ has
gauge at most $p^{\,2\paloss_p-\augain_p}$ for $\phi\in S_0$.
\end{enumerate}
Consequently, in each case, with $S_0=\auS_\iota$ in cases (i), (ii), the module
$S=S_0\cap\auSram$ and the number $i$ satisfy the first assertion of
{\cite[Proposition~14.20]{P46}}: $S\subseteq\auS_\iota\cap\auSram$, $i\le\auerr_\ell\auDX$,
$\aucol(\auS_\iota\cap\auSram,S)\mid\ell^{\,i}$, and $B_{\vartheta,\Psi}\phi$ has gauge at most
$\ell^{\,\auerr_\ell-\augain_\ell}$ for $\phi\in S$.
\end{theorem}
\begin{proof}
(i) Here $\ell\ne d$, because $d\in\mathcal N$. By (K1) and Lemma~\ref{pa:chain} with $r_j=0$,
$B_{\vartheta,\Psi}\phi$ is integral. (And $\augain_\ell=\auerr_\ell=0$ by
{\cite[Lemma~14.9(a)]{P46}}.)

(ii) Let $x$, $j$, $g$, $g'$ be as in (K1). By (K2) and Proposition~\ref{pa:product} the series
$F_{x,j,g,g'}$ has the normal form. By Lemma~\ref{pa:schur} and Proposition~\ref{pa:count}, every
jet along an entry of the Schur matrix $S_{\lambda_{j+1}}$ has valuation at least
$\paPred_j(p)-fn_{j+1}^3(1+\log_pB_j)$, hence, by Lemma~\ref{pa:gain}, at least
$r_j:=\augain_p/m-C_p-\paloss_p/m$. By Lemma~\ref{pa:chain} the gauge of $B_{\vartheta,\Psi}\phi$
is at most $p^{-\sum r_j}=p^{\,mC_p+\paloss_p-\augain_p}$, for every $p\in\mathcal N$ and
$\phi\in\auS_\iota$; and $mC_p\le\paloss_p$ by Lemma~\ref{pa:coll}.

(iii) is Lemma~\ref{pl:level}; $\augain_N=0$ because $N\ge x_1$, and $i\le mfb\,\auDX\le\auerr_N\auDX$.

(iv) We first note a rule for colengths: if $\gamma$ is an endomorphism of a module $M$,
$\alpha$ a scalar, and $\aucol(M,\gamma M+\alpha M)\mid z$, $\aucol(M,M'+\alpha M)\mid z'$ for a
submodule $M'\subseteq M$, then $\aucol(M,\gamma M'+\alpha M)\mid zz'$. Indeed, for a family
$\varphi$ in $M$ there is $\mu$ with $\det\mu\mid z$ and $\mu\varphi=\gamma\psi+\alpha\varrho$, and
$\mu'$ with $\det\mu'\mid z'$ and $\mu'\psi=\nu+\alpha\varrho'$, where $\psi,\varrho,\varrho'$ are
families in $M$ and $\nu$ in $M'$; then $\mu'\mu\varphi=\gamma\nu+\alpha(\gamma\varrho'+\mu'
\varrho)$. Applying this $m\cdot\#\Lambda$ times with Lemma~\ref{pa:eldiv} for $j_0=j_p$ gives
$S_0\subseteq\auS_\iota$ and $\aucol(\auS_\iota,S_0)\mid p^{\,i}$ with
\[
i\le m\sum_{\mathfrak l\in\Lambda}\Bigl(\frac{12A_{\mathfrak l}}{mp}+\frac{j_p}{m^2f}\Bigr)\auDX
\le\Bigl(\frac{12af}p+\frac{j_p}m\Bigr)\auDX\le\Bigl(\frac{14af}p+\frac1m\Bigr)\auDX\le\paloss_p
\auDX,
\]
since $\sum_{\mathfrak l}A_{\mathfrak l}\le af$ and $\#\Lambda\le f$. For the gauge, let
$\phi=Y_\Lambda^{\,m}\phi_1+p^{j_p}\phi_2$ with $\phi_1,\phi_2\in\auS_\iota$. The first term is
treated by Lemma~\ref{pa:hecke}. For the second, the bound proved in (ii) for all $p\in\mathcal N$
gives gauge at most $p^{-j_p+mC_p+\paloss_p-\augain_p}\le p^{\,\paloss_p-\augain_p}$, by
Lemma~\ref{pa:coll}.

For the last statement: $i\le\auerr_\ell\auDX$ in each case, as $\auerr_p=2\paloss_p$ and
$\auerr_N=2mfb$. If a family lies in $\auS_\iota\cap\auSram$, the matrix $\mu$ given by
$\aucol(\auS_\iota,S_0)\mid\ell^{\,i}$ carries it into $S_0$ and, $\auSram$ being a subspace, into
$S_0\cap\auSram$.
\end{proof}

\begin{remark}[Use in the Short Forms Theorem]
The deduction of the Short Forms Theorem ({\cite[Theorem~14.13]{P46}}) in {\cite[\S14.7]{P46}} uses {\cite[Proposition~14.20]{P46}} only
through input (I2), which is the last statement of Theorem~\ref{pa:divis}; so it goes through
unchanged. Compared with the items of {\cite[Proposition~14.20]{P46}}: in (iv) the module is
$(Y_\Lambda^{\,m}\auS_\iota+p^{j_p}\auS_\iota)\cap\auSram$ rather than a sum of submodules of
$\auS_\iota\cap\auSram$ (the operators $Y_{\mathfrak l}$ are not claimed to preserve $\auSram$); in
(i) and (ii) nothing is used of the hypothesis that $\phi$ is ramified. What {\cite[Theorem~14.13]{P46}} then depends on, for this input, is exactly: Lemma~\ref{pa:kernel}
including (K1$_N$), Lemma~\ref{ps:branching}(c), and the properties (W4), (W5) of the Weil
representation in the form (N5) of \S\ref{pl:ss-norm}. Of these, Lemma~\ref{ps:branching}(c) is
Proposition~\ref{pk:branching}, and Lemma~\ref{pa:kernel} is proved in \S\ref{pk:ss-kernel} and in
Corollary~\ref{pf:K1} (the integrality of the jets in (K1) at the primes $\ell\mid2q_1M$, and
(K1$_N$)); nothing remains open for this input.
\end{remark}



\section{The seesaw identity and the upper bound}\label{A1:g4}
\providecommand{\auO}{\mathcal O}
\providecommand{\auU}{\mathrm U}
\providecommand{\auGL}{\mathrm{GL}}
\providecommand{\auHom}{\operatorname{Hom}}
\providecommand{\auA}{\mathbb A}
\providecommand{\auF}{\mathbb F}
\providecommand{\auPl}{\Phi}
\providecommand{\auPa}{\Phi_a}
\providecommand{\auS}{\mathcal S}
\providecommand{\auSram}{\mathcal S^{\mathrm{ram}}}
\providecommand{\auAz}{\mathcal A_0}
\providecommand{\auG}{\overline G}
\providecommand{\auP}{\mathcal P}
\providecommand{\auNm}{\operatorname{Nm}}
\providecommand{\auAut}{\operatorname{Aut}}
\providecommand{\aubar}[1]{\overline{#1}}
\providecommand{\ausinf}{s_\infty}
\providecommand{\aurd}{\mathrm{rd}_E}
\providecommand{\psM}{\mathsf M}
\providecommand{\psHS}{\mathscr H}
\providecommand{\psW}{\mathcal W}
\providecommand{\psZ}{Z}
\providecommand{\psvac}{\mathbf v}
\providecommand{\psrho}{\varrho}
\providecommand{\psE}{\mathcal E}
\providecommand{\psGlue}{\operatorname{Glue}}
\providecommand{\psfa}{\mathfrak a}
\providecommand{\psnot}[1]{\par\smallskip{\raggedright\noindent\textsc{Not proved here.} [#1]\par}\smallskip}
\renewcommand{\psHS}{\mathcal H}

\subsection{The upper bound: the seesaw identity, sizes and restriction}\label{ps:sec}

This section proves {\cite[Lemma~14.21]{P46}} with its constant and
{\cite[Propositions~14.27, 14.24, 14.23, 14.25 and~14.19]{P46}}. To do so it must first make precise several objects that
{\cite[\S14.3]{P46}} only names: the model of the Weil representation, the vacuum and the theta
kernels with their determinant twists (\S\ref{ps:ss-weil}), the lattices of the spaces $Y_j$
(\S\ref{ps:ss-Y}), the gluing and the chain data (\S\ref{ps:ss-glue}). Three statements are stated
here and proved in \S\ref{pk:sec}: the integrality of the branching map
(Lemma~\ref{ps:branching}(c); Proposition~\ref{pk:branching}), the identification of the definite
theta lift with the lift by jets (Lemma~\ref{ps:definite}; Proposition~\ref{pk:definite}), and the
law of Galois conjugation of jets (Lemma~\ref{ps:conj}), which is deduced in
Proposition~\ref{pk:conj} from a reciprocity law on one elliptic curve
(Proposition~\ref{pk:curve}(C3)(ii)), which is Theorem~\ref{pf:recip}. The
standard properties of the Weil representation are collected in Lemma~\ref{ps:weil} with references.

\subsubsection{The Weil representation, the Fock model, the vacuum and the kernels}\label{ps:ss-weil}

\emph{Heisenberg representation.} Put $\jmath=-1/\sqrt{-d}$, so that
$\tau(\jmath)=i/\sqrt d$ for $\tau\in\auPl$, and let $\psi=\psi_\Q\circ
\operatorname{Tr}_{E/\Q}$ be the standard character of $\auA_E/E$, with
$\psi_{\Q,\infty}(t)=e^{-2\boldsymbol\pi it}$; fix a finite idele $\alpha$ such that $\widehat\auO_E$ is its
own dual for $(z,u)\mapsto\psi(\jmath\alpha\aubar\alpha\,\aubar zu)$, and a unitary Hecke
character $\upsilon$ of $E$ whose restriction to $\auA_{E^+}^\times$ is the quadratic character of
$E/E^+$; $\upsilon$ is unramified at every finite place ($\upsilon(x)=\prod_\tau x_\tau/|x_\tau|$
when the finite part of $x$ is a unit; such a character exists by Lemma~\ref{ps:upsilon} below, and
we fix one), and $\alpha$ is $1/P'(\zeta_{q_1})$ at the places dividing
$(1+\sqrt{-d})/2$ and $1$ elsewhere ($P$ the minimal polynomial of $\zeta_{q_1}$ over $K$, as in
\S\ref{pk:ss-conv}). We also fix the finite idele $a$ with $a_w=\alpha_w$ if $\alpha_w$ is not a unit
and $a_w=1$ otherwise; so $a\widehat\auO_E=\alpha\widehat\auO_E$, $a/\alpha$ is a unit idele, and
$a_w\ne1$ only at the primes $w\mid\omega$ above $q_1$ ({\cite[Proposition~6.9]{P53}}). The
coordinates $\tilde x=a^{-1}x$ in which the data at split places are written are described in
{\cite[Lemma~6.32]{P53}}. All hermitian forms $\langle\,,\rangle$ are conjugate-linear in the first variable. For
a hermitian space $W$ over $E$ put $\beta(x,y)=\psi(\tfrac12\jmath\langle y,x\rangle)$ on
$W(\auA)=W\otimes_{E^+}\auA_{E^+}$ and $[x,y]=\beta(x,y)\beta(y,x)^{-1}$. Let $(\psrho,\psHS(W))$
be the smooth irreducible representation of the Heisenberg group of $W(\auA)$ with
$\psrho(x)\psrho(y)=\beta(x,y)\psrho(x+y)$; it is the restricted tensor product
$\psHS_\infty(W)\otimes\psHS_f(W)$, and $\omega=\omega_W$ denotes the Weil representation of
$\auU(W)(\auA)$ on it defined with $\upsilon$ (Weil~\cite{WeilActa}, Kudla~\cite{KudlaSplitting}),
normalised at every place as in Definition~\ref{pwf:def} (at the places of $E^+$ split in $E$, which
the cited works do not treat, by (D3) there).
For hermitian spaces $V,V'$ the space $\auHom_E(V,V')$ carries the form
$\langle w_1,w_2\rangle=\operatorname{tr}(w_1^\dagger w_2)$, where $w^\dagger\colon V'\to V$ is the
adjoint, and $\auU(V)\times\auU(V')$ acts on it by $(g,h)w=hwg^{-1}$. For a character $\mu$ of
$\auU(1)(E^+)\backslash\auU(1)(\auA_{E^+})$ we write $\omega_\mu(g,h)=\mu(\det h)\mu(\det g)^{-1}
\upsilon_{n}(\det h)\upsilon_{n'}(\det g)^{-1}\,\omega_{\auHom(V,V')}(g,h)$, where $n,n'$ are the
ranks of $V,V'$ and $\upsilon_k=(\upsilon|_{\auU(1)})^{\lfloor k/2\rfloor}$. (These are the
\emph{determinant twists}.)

\begin{lemma}[The splitting character]\label{ps:upsilon}
Let $E^+$ be the maximal real subfield of $E$ and $\epsilon$ the quadratic character of
$\auA_{E^+}^\times/E^{+\times}$ attached to $E/E^+$.
\begin{enumerate}
\item[(a)] $E/E^+$ is unramified at every finite place, the group of roots of unity of $E$ is
$\mu_{q_1}$, and it is generated by $-\zeta_{q_1}=(1-\zeta_{q_1})/(1-\aubar{\zeta_{q_1}})$, where
$1-\zeta_{q_1}$ is a unit of $E$ with $\prod_{\tau\in\auPl}\tau(1-\zeta_{q_1})=1$.
\item[(b)] There is a continuous character $\upsilon\colon\auA_E^\times/E^\times\to\{z\in\C:|z|=1\}$
such that $\upsilon|_{\auA_{E^+}^\times}=\epsilon$ and $\upsilon(x)=\prod_{\tau\in\auPl}x_\tau/|x_\tau|$
for every idele $x$ whose finite part lies in $\widehat\auO_E^\times$; here $x_\tau\in\C^\times$ is the
component of $x$ at the place $\tau$, read through the embedding of $E$ which induces $\tau$ and
extends the fixed embedding of $K$.
\item[(c)] For such a character: $\upsilon_w$ is unramified at every finite place $w$ of $E$;
$\upsilon_\tau(z)=z/|z|$ for $\tau\in\auPl$, so that $\upsilon_\tau(z)=z$ for $|z|=1$;
$\upsilon(x\aubar x)=1$ for $x\in\auA_E^\times$; if a place $v$ of $E^+$ splits into $w,\aubar w$,
then $\upsilon_{\aubar w}=\upsilon_w^{-1}$ on $E^{+\times}_v$, and $\upsilon(t,t^{-1})=\upsilon_w(t)^2$
for the element $(t,t^{-1})$ of $\auU(1)(E^+_v)\subset E_w^\times\times E_{\aubar w}^\times$; and
$\upsilon|\cdot|^{-1/2}$ is an algebraic Hecke character, with component $z\mapsto\aubar z^{\,-1}$ at
every $\tau\in\auPl$.
\end{enumerate}
\end{lemma}
\begin{proof}
(a) The field $E$ contains $i$ (as $4\mid q_1$) and $\sqrt{-d}$, and neither lies in $E^+$; so
$E=E^+(i)=E^+(\sqrt{-d})$. The polynomial $X^2+1$ has discriminant $-4$, so $E/E^+$ is unramified
outside $2$; and $\frac12(1+\sqrt{-d})$ is a root of $X^2-X+\frac14(1+d)$, which has integral
coefficients and odd discriminant $-d$ because $d\equiv7\pmod 8$, so $E/E^+$ is unramified above $2$.
Let $\mu_n$ be the group of roots of unity of $E$; then $q_1\mid n$ and $\Q(\mu_n)\subseteq E$. Let $p$
be a prime, $p^c$ and $p^e$ the exact powers of $p$ dividing $n$ and $q_1$, and suppose $c>e$. The
ramification index of $p$ in $E$ is at least $\varphi(p^c)$. If $p\ne d$, then $p$ is unramified in $K$
and its ramification index in $E=K\Q(\zeta_{q_1})$ is $\varphi(p^e)$; but $\varphi(p^c)>\varphi(p^e)$,
because $e\ge1$ if $p=2$. If $p=d$, then $e=0$, $d$ is unramified in $\Q(\zeta_{q_1})$ and its
ramification index in $E$ is $2$; but $\varphi(d^c)\ge d-1\ge6$. So $n=q_1$. Write $\zeta=\zeta_{q_1}$.
We have $-\zeta=\zeta^{1+q_1/2}$, and $1+q_1/2$ is prime to $q_1$ because $q_1/2$ is even and is
divisible by every prime factor of $q_1$; so $-\zeta$ generates $\mu_{q_1}$. Clearly
$(1-\zeta)/(1-\zeta^{-1})=-\zeta$. The places $\tau\in\auPl$ correspond to the embeddings of $E$
extending that of $K$, that is, to $\operatorname{Gal}(E/K)$, which restricts isomorphically to
$\operatorname{Gal}(\Q(\zeta)/\Q)$; so $\prod_\tau\tau(1-\zeta)=\prod_{j\in(\Z/q_1)^\times}(1-\zeta^j)=
\Phi^{\mathrm{cyc}}_{q_1}(1)$, the value at $1$ of the $q_1$-th cyclotomic polynomial, which is $1$ because $q_1$ is divisible by $2$
and by $3$ and so is not a power of a prime. In particular $1-\zeta$ is a unit.

(b) Let $\mathsf U=(E\otimes\R)^\times\times\widehat\auO_E^\times$, an open subgroup of $\auA_E^\times$,
and $\upsilon_{\mathsf U}(x)=\prod_\tau x_\tau/|x_\tau|$ for $x\in\mathsf U$. By (a), $\epsilon$ is
unramified at every finite place, so it is trivial on $\widehat\auO_{E^+}^\times$; at the real place of
$E^+$ below $\tau$ it is the sign. As $\mathsf U\cap\auA_{E^+}^\times=(E^+\otimes\R)^\times\times
\widehat\auO_{E^+}^\times$ and $x_\tau/|x_\tau|$ is the sign of $x_\tau$ for real $x_\tau$, the
characters $\upsilon_{\mathsf U}$ and $\epsilon$ agree on $\mathsf U\cap\auA_{E^+}^\times$. Hence there
is a unique character $\chi$ of the open subgroup $\mathsf U\cdot\auA_{E^+}^\times$ with
$\chi(ua)=\upsilon_{\mathsf U}(u)\epsilon(a)$; it is continuous, being continuous on the open subgroup
$\mathsf U$.

We claim that $\chi(e)=1$ for every $e\in E^\times\cap\mathsf U\auA_{E^+}^\times$. Write $e=ua$. Then
$e/\aubar e=u/\aubar u$ is an element of $E^\times$ which is a unit at every finite place and has
absolute value $1$ at every infinite place; so it is a root of unity (Kronecker), and by (a)
$e/\aubar e=(-\zeta)^j$ for an integer $j$. Then $e^+=e\,(1-\zeta)^{-j}$ satisfies $\aubar{e^+}=e^+$,
that is, $e^+\in E^{+\times}$. As $(1-\zeta)^j\in\auO_E^\times\subset\mathsf U$ and
$e^+\in\auA_{E^+}^\times$,
\[
\chi(e)=\upsilon_{\mathsf U}\bigl((1-\zeta)^j\bigr)\,\epsilon(e^+)=\Bigl(\prod_\tau\frac{\tau(1-\zeta)}
{|\tau(1-\zeta)|}\Bigr)^{j}=1,
\]
because $\epsilon$ is trivial on $E^{+\times}$ and $\prod_\tau\tau(1-\zeta)=1$ by (a). This proves
the claim.

Consequently $\chi_1(eua)=\chi(ua)$ ($e\in E^\times$, $u\in\mathsf U$, $a\in\auA_{E^+}^\times$) is a
well-defined character of the open subgroup $G_1=E^\times\mathsf U\auA_{E^+}^\times$ of
$\auA_E^\times$, trivial on $E^\times$ and continuous. The quotient $\auA_E^\times/G_1$ is a quotient of
the ideal class group $\auA_E^\times/E^\times\mathsf U$ of $E$, hence finite. A character of a
subgroup of an abelian group with values in the divisible group $\{|z|=1\}$ extends to the group; an
extension $\upsilon$ of $\chi_1$ to $\auA_E^\times$ is continuous because $G_1$ is open, and it has the
properties stated.

(c) The first two assertions are the formula for $\upsilon$ on $\mathsf U$. Next,
$x\aubar x\in\auA_{E^+}^\times$ is a norm from $\auA_E^\times$, and $\epsilon$ is trivial on norms. If $v$
splits into $w,\aubar w$, then $t\in E^{+\times}_v$ is the idele of $E$ with components $t$ at $w$ and at
$\aubar w$, and $\epsilon_v=1$; so $\upsilon_w(t)\upsilon_{\aubar w}(t)=1$, and $\upsilon(t,t^{-1})=
\upsilon_w(t)\upsilon_{\aubar w}(t)^{-1}=\upsilon_w(t)^2$. Finally the idele norm of $E$ is $|z|^2$ at a
complex place, so the component of $\upsilon|\cdot|^{-1/2}$ at $\tau$ is $z|z|^{-2}=\aubar z^{\,-1}$.
\end{proof}

The character $\upsilon$ is determined by the conditions of Lemma~\ref{ps:upsilon}(b) up to a
character of the finite group $\auA_E^\times/E^\times\mathsf U\auA_{E^+}^\times$, where
$\mathsf U=(E\otimes\R)^\times\times\widehat\auO_E^\times$; any choice will do,
and it is fixed once and for all.

\begin{lemma}[Standard facts]\label{ps:weil}
Let $W$ be a hermitian space over $E$.
\begin{enumerate}
\item[(W1)] $\omega(g)\psrho(w)\omega(g)^{-1}=\psrho(gw)$. There is a functional $\theta=\theta_W$
on $\psHS(W)$, unique up to a scalar, with $\theta\circ\psrho(w)=\theta$ for $w\in W(E)$; it
satisfies $\theta\circ\omega(\gamma)=\theta$ for $\gamma\in\auU(W)(E^+)$.
\item[(W2)] Let $I\subset W$ be an isotropic $E$-subspace, $W_I=I^\perp/I$, and $P_I$ the stabiliser
of $I$ in $\auU(W)$. There is a continuous linear map $\psZ_I\colon\psHS(W)\to\psHS(W_I)$ (the
\emph{zero term}) with $\theta_{W_I}(\psZ_IF)=\int_{I(E)\backslash I(\auA)}\theta(\psrho(i)F)\,di$,
$\psZ_I\psrho(w)=\psrho(\bar w)\psZ_I$ for $w\in I^\perp(\auA)$ with image $\bar w$, and
$\psZ_I\omega(p)=\chi_I(p)\,\omega_{W_I}(\bar p)\psZ_I$ for $p\in P_I(\auA)$, where $\chi_I(p)=
(\upsilon|\cdot|^{-1/2})(\det p|_I)^{-1}=\upsilon(\det p|_I)^{-1}|\det p|_I|^{1/2}$, with
$\det p|_I\in\auA_E^\times$ the determinant on the $\auA_E$-module $I(\auA)$ and $|\cdot|$ the
idele norm of $E$; so $\chi_I=1$ on unipotent elements and on $P_I(E^+)$.
For isotropic $I\subset I'$ one has $\psZ_{I'}=\psZ_{I'/I}\psZ_I$, and $\psZ_I=\psZ_{I,\infty}
\otimes\psZ_{I,f}$.
\item[(W3)] For $F\in\psHS(W)$, $\theta(F)=\sum_{w\in W(E)/I^\perp(E)}\theta_{W_I}(\psZ_I\psrho(w)F)$,
the series converging absolutely. If $W=\auHom(V,V')$, $F=\Psi\otimes f$ with $\Psi$ smooth, and
$\mathcal C$ is a compact set of finite adelic points of $\auU(V)\times\auU(V')$, there are $C,A$ with
$\sum_w|\theta_{W_I}(\psZ_I\psrho(w)\omega(g,h)F)|\le C(\|g_\infty\|\|h_\infty\|)^A$ whenever
$(g_f,h_f)\in\mathcal C$.
\item[(W4)] Let $L\subset W(\auA_f)$ be a lattice that is its own dual for $[\,,]$. The vectors of
$\psHS_f(W)$ on which $\psrho(l)$, $l\in L$, acts by a fixed character form a line. If $\psvac_L$
spans the line of the character $\epsilon_L$ with $\epsilon_L(l)\epsilon_L(l')=\beta(l,l')
\epsilon_L(l+l')$ fixed below, and $\psrho^\circ(x)=\epsilon(x)\psrho(x)$ for a fixed function
$\epsilon$ with values in roots of unity extending $\epsilon_L^{-1}$, then
$\psrho^\circ(l)\psvac_L=\psvac_L$, $\psrho^\circ(x+l)\psvac_L=\zeta\psrho^\circ(x)\psvac_L$ with a
root of unity $\zeta$, the $\psrho^\circ(x)\psvac_L$, $x\in W(\auA_f)/L$, form a basis of
$\psHS_f(W)$, and $\omega(g)\psvac_L\in\C^\times\psvac_{gL}$ for $g\in\auU(W)(\auA_f)$. If
$g\in\auU(W)(\auA_f)$ stabilises $L$ then $\omega(g)\psvac_L=\psvac_L$.
\item[(W5)] At a place $v=w\aubar w$ of $E^+$ split in $E$, $\psHS_v(\auHom(V,V'))$ is the space
of Schwartz--Bruhat functions on $\auHom(V_w,V'_w)$, with $(\omega_\mu(g,h)\Psi)(x)=\mu_w(\det h_w)
\mu_w(\det g_w)^{-1}\upsilon'_w(g_w,h_w)\,\Psi(h_w^{-1}xg_w)$, where, for even ranks $n,n'$, $\mu_w\upsilon'_w$ is given by \eqref{pl:eq-w5}: the factor is
$\chi_w(\det h_w)|\det h_w|_w^{-n/2}\chi_w(\det g_w)^{-1}|\det g_w|_w^{n'/2}$ with $\chi$ the base
change of $\mu$; for arbitrary ranks the untwisted operator is
\[\Psi\mapsto\upsilon_w(\det h_w)^{-n}\upsilon_w(\det g_w)^{n'}|\det h_w|_w^{-n/2}|\det g_w|_w^{n'/2}\Psi(h_w^{-1}xg_w)\]
(Proposition~\ref{pwf:W5}), so that all factors multiplying $\Psi(h_w^{-1}xg_w)$ are built from $\mu_w$, $\upsilon_w$ and $|\cdot|_w$ evaluated at $\det g_w,\det h_w$; $\psrho(a_w,0)\Psi(x)=\Psi(x+a_w)$
and $\psrho(0,a_{\bar w})$ is multiplication by the character $x\mapsto[x,a_{\bar w}]$.
\item[(W0), (W6)] Two further facts are used; they are stated in full in \S\ref{pk:ss-conv}.
(W0): $\theta$ is defined on the smooth vectors of the Fock space tensored with $\psHS_f(W)$ and is
continuous in the first factor. (W6): for hermitian spaces $V,V'$ of ranks $n,n'$ the map
$w\mapsto w^\dagger$ identifies $\psHS(\auHom(V,V'))$ with $\psHS(\auHom(V',V^-))$, carrying
$\omega_{\auHom(V,V')}(g,g')$ to $\upsilon(\det g)^{n'}\upsilon(\det g')^{-n}\,
\omega_{\auHom(V',V^-)}(g',g)$. It carries the self-dual lattice $\alpha\auHom(\psM,\psM')$ to
$\aubar\alpha\auHom(\psM',\psM)$, which differs from $\alpha\auHom(\psM',\psM)$ at the places above $q_1$
(Remark~\ref{pf:rem-flip-lattice}).
\end{enumerate}
\end{lemma}
\begin{proof}[References]
(W1): the commutation relation is part of the definition (Definition~\ref{pwf:def}, with
\cite[Th\'eor\`eme~1 and no.~34]{WeilActa}); existence and uniqueness of $\theta$ are
Lemma~\ref{pwf:unique}, and its invariance under $\auU(W)(E^+)$ for the splitting by $\upsilon$ is
Theorem~\ref{pwf:inv} (by Corollary~\ref{pwf:rk}, $\omega$ agrees on the rational points with Weil's
lift $r_k$, for which the invariance is \cite[Th\'eor\`eme~6]{WeilActa}). (W2), (W3) are the
properties of the mixed model attached to $W=I\oplus I^*\oplus W_I$: $\psHS(W)=\auS(I^*(\auA))
\otimes\psHS(W_I)$, $\psZ_I$ is evaluation at $0$, $\theta(F)=\sum_{\xi\in I^*(E)}\theta_{W_I}(F(\xi))$;
they are proved in Propositions~\ref{pwf:W2} and~\ref{pwf:W3} (Kudla's explicit formulas \cite[\S5]{KudlaSplitting} are
for a Lagrangian $I$); the absolute
convergence and the bound in (W3) are Lemma~\ref{pz:absgrowth} (which uses (W0), (W1), (W2), (W4) and
Lemma~\ref{pz:growth}, not (W3)). (Weil~\cite[Th\'eor\`eme~6]{WeilActa} gives the continuity of the theta
distribution and its invariance under the rational points, used in (W1), but no bound of this kind.)
(W4) is the lattice model; for the last assertion, the vector
$\omega(g)\psvac_L$ lies on the line of $\psvac_L$ and the scalar is $1$ for the splitting used
(it is $\upsilon(\det u)|\det u|^{-1/2}$ with $uL=L\cap g^{-1}L$). (W5) is Proposition~\ref{pwf:W5}
(the cited works treat only places where $E^+_v\otimes E$ is a field); see \S\ref{pl:ss-norm}, (N5). For (W0) see \S\ref{pk:ss-conv}; (W6) is proved in
Proposition~\ref{pf:flip} from (W1), (W4) and the local form of (W2).
\end{proof}

\emph{The Fock model and the theta function.} Let $X,X'$ be positive definite of ranks $n<n'$ and
$W=\auHom(X,X')$, so that $W\otimes_{E^+,\tau}\R=M_{n'\times n}(\C)$ for $\tau\in\auPl$. We use
the coordinates $x_\tau=c_\psi\cdot\tau(w)\in M_{n'\times n}(\C)$ with
$c_\psi=\sqrt{2\boldsymbol\pi/\sqrt d}$, in which $\beta_\infty(x,y)=\exp(+i\,\Iml\sum_\tau
\operatorname{tr}(y_\tau^*x_\tau))$; the variable $z$ of the Fock space is the complex conjugate of
$x$. Then $\psHS_\infty(W)$ is the Fock space of entire functions $F$ on
$\prod_\tau M_{n'\times n}(\C)$ with $\int|F|^2e^{-|z|^2}<\infty$, with
\[
(\psrho_\infty(y)F)(z)=e^{z\cdot y-|y|^2/2}F(z-\aubar y),\qquad z\cdot y=\sum_\tau
\operatorname{tr}(y_\tau^{\,t}z_\tau),
\]
the \emph{real vacuum} is $F=1$, and $\auU(X)(\R)\times\auU(X')(\R)$ acts by
$F\mapsto F(h^{t}z\,\aubar g)$ times the determinant twists
$\prod_\tau\det(h_\tau)^{e_{j+1,\tau}}\det(g_\tau)^{-(e_{j,\tau}-c_\tau(n_j))}$. For $f\in\psHS_f(W)$ the \emph{theta function} of $f$ is the entire
function
\begin{equation}\label{ps:eq-theta}
\Theta_f(z)=\theta\bigl(e^{\langle\,\cdot\,,\bar z\rangle}\otimes f\bigr)=e^{|z|^2/2}\,\theta\bigl(
\psrho_\infty(z)(1\otimes f)\bigr)\qquad(z\in\textstyle\prod_\tau M_{n'\times n}(\C)),
\end{equation}
and for a polynomial $Q$ in the entries of $z$ and $c\in\C^\times$ the \emph{jet of $f$ along $Q$
at scale $c$} is $Q(\partial)[\Theta_f(c\,\cdot)](0)$. By \eqref{ps:eq-theta} and the reproducing
property of the Fock space, with $d\mu(y)=\pi^{-fnn'}dy$,
\begin{equation}\label{ps:eq-bargmann}
Q(\partial)[\Theta_f(c\,\cdot)](0)=\theta\bigl(Q(c\,\cdot)\otimes f\bigr)=\int Q(c\,\bar y)\,
e^{-|y|^2/2}\,\theta\bigl(\psrho_\infty(y)(1\otimes f)\bigr)\,d\mu(y);
\end{equation}
indeed $\partial^\alpha H(0)=\int\bar y^\alpha H(y)e^{-|y|^2}d\mu(y)$ for entire $H$ with
$H(y)e^{-|y|^2/2}$ bounded, and the theta values are bounded.

\emph{The vacuum and its normalisation.} The lattice of a standard space $X$ is
$\widehat\auO_E^{\,n}$, and $L=\alpha\auHom(\widehat\auO_E^{\,n},\widehat\auO_E^{\,n'})$ is
self-dual in $W(\auA_f)$. The character $\epsilon_L$ is $l\mapsto\psi_f(-\tfrac12\omega\,\jmath\langle l,l\rangle)$
with $\omega=(1+\sqrt{-d})/2$, an element of trace one (then $\epsilon_L(l)\epsilon_L(l')=\beta(l,l')
\epsilon_L(l+l')$, and $\epsilon_L=\pm1$: \S\ref{pf:ss-adm}, Lemma~\ref{pf:eps}), the \emph{vacuum} is the vector $\psvac=\psvac_L$ of (W4) normalised by
$\theta(1\otimes\psvac)$ having mean square $1$ over $W(E)\backslash W(\auA)$, and the
\emph{integral vacuum} is $\psvac^\circ=c_{\mathrm{vac}}^{\,nn'}\psvac$, where $c_{\mathrm{vac}}\psvac$, for the space of rank one, is the vector \eqref{pk:eq-vac}: a sum of
$q_1^{2f}$ translates, with coefficients roots of unity, of the vector $F_0$ whose theta function
is the product \eqref{pk:eq-product} of theta functions of the curve $\C/\Omega_0\auO_K$. (No
closed formula for $c_{\mathrm{vac}}$ is needed; \eqref{pk:eq-vac} defines it.) The translates $\phi^\circ_x=\psrho^\circ(x)\psvac^\circ$, $x\in W(\auA_f)$, are
the ``translates of the vacuum'' of {\cite[\S14.3]{P46}}.

\emph{Weights and kernels.} At $\tau\in\auPl$ let $c_\tau(h)$ be the number of columns of height
$h$ fixed in {\cite[\S14.3]{P46}}: $c_\tau(1)=a\mathbf 1_{\tau\in\auPa}$, $c_\tau(k)=\kappa_1(\tau)
-n_1/2$, $c_\tau(n_i)=\Delta-(t+1)$ for $i\ge1$, and $c_\tau(h)=0$ otherwise. The weight
$\lambda(X_j)$ has the columns of height $<n_j$ and the determinant exponent $e_{j,\tau}=0$ for
$j=0$ and $e_{j,\tau}=n_{j-1}/2-\kappa_j(\tau)$ for $j\ge1$; thus its rows at $\tau$ are
$\lambda(X_j)_{\tau,p}=e_{j,\tau}+\sum_{p\le h<n_j}c_\tau(h)$ ($1\le p\le n_j$), and
$e_{j,\tau}-c_\tau(n_j)=n_{j+1}/2-\kappa_{j+1}(\tau)=e_{j+1,\tau}+\tfrac12(n_{j+1}-n_j)$. For a matrix $z_\tau\colon\C^{n_j}\to\C^{n_{j+1}}$ let
$S_\tau(z_\tau)\colon W_{\lambda(X_j)_\tau}\to W_{\lambda(X_{j+1})_\tau}$ be the Schur functor of
the diagram with $c_\tau(h)$ columns of each height $h<n_{j+1}$, applied to $z_\tau$ (on $\C^{n_j}$
the columns of height $n_j$ are powers of the determinant), and $S(z)=\bigotimes_\tau S_\tau
(z_\tau)$; its matrix entries in the standard bases are polynomials in $z$ with integer
coefficients, homogeneous of degree $\sum_pl_{\tau,p}$ in $z_\tau$, where $l_{\tau,p}=\sum_{p\le h<
n_{j+1}}c_\tau(h)$, $1\le p\le n_j$, are the rows of the diagram of the $j$-th lift. With these
coordinates the rational points act on $W_\lambda$ through the conjugate embeddings,
$\gamma\cdot\xi=\rho_\lambda(\aubar{\tau(\gamma)})\xi$, and $g_\infty\cdot\xi=\rho_\lambda(
\aubar{g_\infty})\xi$; the expressions $g_\infty^{-1}\psi(g_f)$ below are to be read with this
action. Let
$s_0=(\varpi_{\mathrm{vac}}/\mathsf u)\,c_\psi/\Omega_0$ be the scale \eqref{pk:eq-s0};
$|s_0|^2=1/\Omega$. For $\Phi\in\psHS_f(\auHom(X_j,X_{j+1}))$
and finite adelic points $g,g'$ the \emph{kernel} is
\begin{equation}\label{ps:eq-kernel}
\theta_j(\Phi;g,g')=S(\partial)\bigl[\Theta_{\omega_{\mu_{j+1}}(g,g')\Phi}(\pi_Ms_0\,\cdot)\bigr](0)
\ \in\auHom(W_{\lambda(X_j)},W_{\lambda(X_{j+1})}),
\end{equation}
that is, the matrix of the jets at scale $\pi_Ms_0$ along the entries of $S$. The kernels have the
two properties (T1), (T2) of \S\ref{pl:ss-norm}, (N4): equivariance under the rational points,
$\theta_j(\Phi;\gamma g,\gamma'g')=\gamma'\circ\theta_j(\Phi;g,g')\circ\gamma^{-1}$, and under the
finite adelic points, $\theta_j(\omega_{\mu_{j+1}}(h,h')\Phi;g,g')=\theta_j(\Phi;gh,g'h')$. The lift is
$(\Theta^X_j(\Phi)\varphi)(g')=\int_{\Gamma\backslash G_f}\theta_j(\Phi;g,g')\varphi(g)\,dg$ for the
Haar measure giving volume $1$ to $J_{X_j}$; this is $[J_{X_j}:J]^{-1}\sum_{g\in\Gamma\backslash
G_f/J}\theta_j(\Phi;g,g')\varphi(g)$ for every free level $J\subseteq J_{X_j}$ fixing $\varphi$ and
$\omega_{\mu_{j+1}}(\cdot,1)\Phi$. This is the definition used in the rest of the paper; that it has
values in algebraic modular forms, and its relation to theta integrals, is the following.

\begin{lemma}[The definite lift]\label{ps:definite}
Let $X,X',X''$ be consecutive spaces $X_{j},X_{j+1},X_{j+2}$. For every smooth vector $u$ of the Fock
space of $\auHom(X',X^-\oplus X'')\otimes\R$ and every linear form $l$ on $W_{\lambda(X)}$ there is
a linear form $l''$ on $W_{\lambda(X'')}$ such that for all $\vartheta$, all
$\Phi\in\psHS_f(\auHom(X,X'))\otimes\psHS_f(\auHom(X',X''))$, all algebraic modular forms $\psi$ on
$\auU(X)$ of weight $\lambda(X)$ and all $k\in\auU(X'')(\auA)$,
\[
\int_{[\auU(X')]}\int_{[\auU(X)]}l\bigl(g_\infty^{-1}\psi(g_f)\bigr)\,\theta\bigl(\omega(g,r,k)(u\otimes
\Phi)\bigr)\,dg\,dr=l''\Bigl(k_\infty^{-1}\bigl(\Theta^X_{j+1}\Theta^X_j(\Phi)\psi\bigr)(k_f)\Bigr),
\]
where $[G]=G(E^+)\backslash G(\auA)$ has volume $1$, $\omega(g,r,k)=\omega_{\mu_{j+1}}(g,r)\otimes
\omega_{\mu_{j+2}}(r,k)$ on $\psHS(\auHom(X,X'))\otimes\psHS(\auHom(X',X''))=\psHS(\auHom(X',
X^-\oplus X''))$, and $\Theta^X_{j+1}\Theta^X_j(\Phi)$ is defined by linearity from pure tensors.
\end{lemma}
\begin{proof} See Proposition~\ref{pk:definite}. \end{proof}

\subsubsection{Short bases: proof of \texorpdfstring{{\cite[Lemma~14.21]{P46}}}{[Row D, Lemma~14.21]}}\label{ps:ss-short}

\begin{lemma}\label{ps:transfer}
Let $\Lambda$ be a lattice in a euclidean space of dimension $m\ge1$ whose dual lattice has no
non-zero vector shorter than $\rho>0$. Then every point of the space is within $m^2/\rho$ of
$\Lambda$, and the vectors of $\Lambda$ of length at most $3m^2/\rho$ span the space.
\end{lemma}
\begin{proof}
By Minkowski's convex body theorem applied to the cube $[-t,t]^m$, $t=(\det\Lambda)^{1/m}$, $\Lambda$ has a non-zero vector $v$ with $|v|\le\sqrt m\,
(\det\Lambda)^{1/m}$, and $\det\Lambda=(\det\Lambda^*)^{-1}\le(\sqrt m/\rho)^m$ by the same bound
for $\Lambda^*$; so $|v|\le m/\rho$. We prove the first claim by induction on $m$. For $m=1$,
$\Lambda=\Z v$, $1/|v|\ge\rho$, and the covering radius is $|v|/2$. For $m>1$ let $v$ be a
shortest vector and $\Lambda'$ the projection of $\Lambda$ to $v^\perp$; its dual is
$\Lambda^*\cap v^\perp$, so the hypothesis holds for $\Lambda'$, and the covering radius of
$\Lambda$ is at most that of $\Lambda'$ plus $|v|/2$, hence at most $(m-1)^2/\rho+m/(2\rho)\le
m^2/\rho$. For the second claim put $r=m^2/\rho$: if the vectors of length $\le3r$ lay in a
hyperplane $H$, a point at distance $2r$ from $H$ on the normal through $0$ would be within $r$ of
some $\lambda\in\Lambda$, with $|\lambda|\le3r$, so $\lambda\in H$, which is absurd.
\end{proof}

\begin{proof}[Proof of {\cite[Lemma~14.21]{P46}}, with $C_{\mathrm{sb}}=6c_0(fn)^2$]
Let $L=E^n\cap g\widehat\auO_E^{\,n}$, a lattice in $E^n$ that is unimodular for the standard
hermitian form $h$, and embed $E^n$ in $\C^{\auPl\times n}=\R^{2fn}$ by $y\mapsto(\tau(y_i))$, with
$\|y\|^2=\sum_\tau\tau(h(y,y))$; the real inner product is $(x,y)=\tfrac12\operatorname{Tr}_{E/\Q}
h(x,y)$. If $y\in L$ is not zero then $h(y,y)$ is a totally positive integer of $E^+$, so
$\|y\|^2\ge f\ge1$. Put $\Lambda=(2/c_0)L$. If $x$ lies in the dual of $\Lambda$, then
$\operatorname{Tr}_{E/\Q}h(\xi,x/c_0)\in\Z$ for all $\xi\in L$; since $L$ is unimodular this means
$x/c_0\in\mathfrak d_E^{-1}L$, so $x\in c_0\mathfrak d_E^{-1}L\subseteq L$ because $c_0$ lies in
the different $\mathfrak d_E$. Hence the dual of $\Lambda$ has no non-zero vector shorter than $1$,
and by Lemma~\ref{ps:transfer} with $m=2fn$ the vectors of $\Lambda$ of length at most $3(2fn)^2$
span $\R^{2fn}$. Multiplying by $c_0/2$: the vectors of $L$ of length at most $6c_0(fn)^2$ span
$\R^{2fn}$ over $\R$, hence $E^n$ over $E$, and $n$ of them form a basis; let $Y$ be the matrix of
its columns. Then $g^{-1}Y$ is integral, and every entry $y$ of $Y$ lies in $E$ and satisfies
$|\tau(y)|\le6c_0(fn)^2$ for $\tau\in\auPl$, hence at every embedding of $E$ in $\C$, hence
$|\sigma y|\le6c_0(fn)^2$ for every $\sigma\in\auAut(\C)$.
\end{proof}
The sharper transference theorem of Banaszczyk~\cite[Theorem~2.1]{Banaszczyk} would give a constant linear in
$fn$; the lemma is stated with the constant that the elementary argument gives.

\subsubsection{The spaces $Y_j$, their lattices and data; proof of \texorpdfstring{{\cite[Proposition~14.27]{P46}}}{[Row D, Proposition~14.27]}}
\label{ps:ss-Y}

Let $0\le j<m$ and $s_j=m-1-j$. The space $Y_j$ has a basis consisting of the standard basis of
$X_j$, with last vector $y=e_{\mathrm{last}}$, and of a \emph{string} $f_0,\dots,f_{s_j-1}$ orthogonal
to $X_j$ with $\langle f_a,f_b\rangle=0$ unless $a+b=s_j-1$, $\langle f_a,f_b\rangle=1$ if
$a+b=s_j-1$, $a\ne b$, and $\langle f_a,f_a\rangle=-N$ for the \emph{middle vector} $z=f_a$,
$2a=s_j-1$ (which exists when $j$ is even). The \emph{grade} is $0$ on $X_j$ and $s_j-1-2a$ on
$f_a$. Put $Y_m=e_{\mathrm{last}}^\perp\subset X_m$, with no string. Let $c_M$ be the idele that
generates $(\pi_M)$ and is $1$ away from the chosen places, $\mathsf a=\alpha/\aubar{c_M}$,
$\psfa$ the ideal of $\mathsf a$, $u_j=(\mathsf a/\aubar{\mathsf a})^{\lceil(m-j)/2\rceil}$, and
$c=-(\jmath M\delta_E)^{-1}\in E^+$ with $\delta_E=\varpi$ the purely imaginary generator of the different of $E$ fixed in {\cite[Lemma~5.10]{P51}} (the one used in {\cite[(4)]{P51}});
$c$ has the valuations of $\mathsf a\aubar{\mathsf a}$ at all finite places. The lattice
$\psM_j$ of $Y_j$ is spanned over $\widehat\auO_E$ by
\begin{equation}\label{ps:eq-lattice}
\begin{cases}
e_i\ (i<n_j),\quad y/\aubar{\pi_N}+u_jz/N,\quad u_jf_a\ (\text{all }a)&(j\text{ even}),\\
e_i\ (i\le n_j),\quad u_j\mathsf a^{-1}f_a\ (2a<s_j),\quad u_j\aubar{\mathsf a}f_a\ (2a>s_j)&(j
\text{ odd}),
\end{cases}
\end{equation}
and $\psM_m=\widehat\auO_E^{\,n_m-1}$. It is unimodular: for $j$ even the vector
$y/\aubar{\pi_N}+u_jz/N$ is isotropic and pairs to $-1$ with $u_jz$. The idele $u_j$ satisfies
$u_j\aubar{u_j}=1$ and is not a unit: with $e=\lceil(m-j)/2\rceil$, at a chosen place $w$ it is
$c_{M,w}^{\,e}$, a generator of $\pi_M^{\,e}\auO_w$, and at $\aubar w$ it is $\aubar{c_{M,w}}^{\,-e}$, at a place $w\mid\omega$ above $q_1$
it is $\alpha_w^{\,e}$ and at $\aubar w$ it is $\aubar{\alpha_w}^{\,-e}$, and it is $1$ elsewhere
(in particular above $N$). For $j=0$ this is the lattice $L_{u_0}=u_0\cdot\mathsf h_{u_0}L$ of
{\cite[Lemma~2.14]{P51}}, where $L$ is the lattice of $Y_0$ used in {\cite[\S2]{P51}} (the same basis
with $u_0$ replaced by $1$) and $\mathsf h_{u_0}$ is the central element of $\auU(X_0)(\auA_f)$ acting by $u_0^{-1}$;
all the objects on $Y_0$ of {\cite[\S14.3]{P46}} (the level $J_M$, the space $\auAz$, the action
of $\auG$, the map $\auP$) are formed with $\psM_0$, and {\cite[Lemma~2.14]{P51}} carries the
statements of {\cite[Proposition~14.32]{P46}} from $L$ to $\psM_0$. For $j$ even, $\psM_j\cap X_j=
\widehat\auO^{\,n_j-1}\oplus\pi_N\widehat\auO y$, and $\psM_j$ is not the sum of its intersections
with $X_j$ and with the string at the places above $N$ (it is \emph{glued} there). We write
$\psM_j^-$ for the intersection of $\psM_j$ with the span of the string and $\psM_j^+$ for its
projection to that span; they are dual to each other, and equal except that $\psM_j^+\ni u_jz/N$
for $j$ even. Let $G_j=\auU(Y_j)$, $K_j\subset G_j(\auA_f)$ the stabiliser of $\psM_j$,
$J_j\subseteq K_j$ the subgroup of elements $\equiv1$ at the chosen places and at
the places of $Q=Q_\vartheta$ (the principal congruence subgroup modulo $(\pi_M)\prod_{w\in Q}w$; for
$j=m$ it is contained in, and for $Q\ne\emptyset$ strictly smaller than, the level $J_\vartheta$ of
$\auAm(\vartheta)$, which is block upper triangular modulo $Q$, Lemma~\ref{pz:lastlevel}), and $J_M\supseteq J_0$ the level of $\auAz$. The frame of $\psM_j$
used to write reductions is \eqref{ps:eq-lattice} up to an element of $\auGL(\widehat\auO)$ that
is $\equiv1$ at the chosen places, numbered so that the frame of $Y_j$ is the beginning of the
frame of $Y_{j+1}$; $\iota_j$ denotes the corresponding inclusion matrix. (Any numbering with
this property will do; nothing below depends on it.)

\emph{The string map.} Let $T_j\colon Y_j\to Y_{j+1}$ be zero on $X_j$ and $T_jf_a=t_af'_a$
($a\le s_j-2$), $T_jf_{s_j-1}=0$, where $f'_a$ is the string of $Y_{j+1}$ and, for $j$ even,
$t_a=1,\,-Nc,\,c$ according as the grade of $f_a$ is positive, zero, negative, and for $j$ odd
$t_a=c,\,1$ according as it is positive, negative. Then $T_j$ lowers grades by $1$, and a direct
computation with \eqref{ps:eq-lattice} gives
\begin{equation}\label{ps:eq-TM}
T_j(\psM_j^+)=\psfa\,\psM_{j+1}^-,\qquad T_j^\dagger T_j=\mathfrak e_j,\qquad T_jT_j^\dagger=
\mathfrak e_{j+1},
\end{equation}
where $\mathfrak e_j$ is a self-adjoint nilpotent endomorphism of $Y_j$ that kills $X_j$ and maps
$f_a$ to a non-zero multiple of $f_{a+1}$. (For $j$ even both sides of the first identity are
$u_j(\bigoplus_{a<p}\widehat\auO f'_a\oplus c\bigoplus_{a\ge p}\widehat\auO f'_a)$, $2p=s_j-1$,
because $u_{j+1}=u_j$; for $j$ odd both are $u_j\aubar{\mathsf a}\bigoplus_a\widehat\auO f'_a$,
because $u_j\aubar{\mathsf a}=u_{j+1}\mathsf a$.)

\emph{Bessel groups.} Let $N_j\subset G_j$ be the group of $u$ such that $u-1$ raises grades,
$N_j^{\ge2}$ the subgroup where it raises them by at least $2$ (so $N_j^{\ge2}=N_j$ for $j$ even),
and $R_j=N_j\rtimes\auU(X_j)$, where $\auU(X_j)$ acts on $X_j$ and fixes the string; put
$R_m=G_m$. Equivalently $R_j$ is the group of $g$ with $gf_a\in f_a+\langle f_b:b<a\rangle$ for
all $f_a$ of grade $\ge0$. Let $\psi_j(u)=\psi(\tfrac12\jmath\operatorname{tr}(\mathfrak e_j(u-1)))$,
a character of $N_j^{\ge2}(\auA)$ trivial on $N_j^{\ge2}(E^+)$ and invariant under conjugation by
$R_j(\auA)$. A continuous function $\psW$ on $R_j(E^+)\backslash R_j(\auA)$ is
\emph{equivariant} if $\psW(ur)=\psi_j(u)\psW(r)$ for $u\in N_j^{\ge2}(\auA)$. The group $R_0$ is
the Bessel subgroup of {\cite[\S14.3]{P46}}, $\psi_0=\psi_U$, and $\psW_\varphi$ is equivariant.
(For the character {\cite[(4)]{P51}}: with $p=(s_0-1)/2$ and $u_{a,a+1}$ the coefficient of $f_a$ in
$uf_{a+1}$, \eqref{ps:eq-TM} gives $\mathfrak e_0f_a=cf_{a+1}$ for $a\ne p$ and $\mathfrak e_0f_p=-Ncf_{p+1}$,
and the unitarity of $u$ gives $u_{s_0-2-a,s_0-1-a}=-\aubar{u_{a,a+1}}$ for $a<p-1$ and
$u_{p,p+1}=\aubar{u_{p-1,p}}/N$; hence $\operatorname{tr}(\mathfrak e_0(u-1))=c\,(\sigma(u)-\aubar{\sigma(u)})$
with $\sigma(u)=\sum_{a<p}u_{a,a+1}$ the sum in {\cite[(4)]{P51}}, and $\psi_0(u)=\psi(\jmath c\,\sigma(u))=
\psi(-\sigma(u)/(M\delta_E))$.)

\emph{The data.} The lift $\Theta^Y_{j,\vartheta}$ is $[K_j:J_j]$ times the theta lift for
$\omega_{\mu_{j+1}}$ on $\psHS(\auHom(Y_j,Y_{j+1}))$ with the datum $\Upsilon_j=P_j\otimes
\Upsilon_{j,f}$: $P_j$ is the joint harmonic of {\cite[\S14.3]{P46}}, and $\Upsilon_{j,f}$ is the
vacuum $\psvac_{L_j}$, $L_j=\alpha\auHom(\psM_j,\psM_{j+1})$, at the places outside the level
places, and at a level place $w$ (chosen, or in $Q$) the function
$\mathbf 1_{a_w(\iota_j+\varpi_w\auHom(\psM_{j,w},\psM_{j+1,w}))}$ in the model (W5), written in the
frames, with $a_w$ the component of the idele $a$ of \S\ref{ps:ss-weil} (so $a_w=1$ except at the
primes of $Q$ above $q_1$ dividing $\omega$; in the coordinates $\tilde x=a_w^{-1}x$ of
{\cite[Lemma~6.32]{P53}} this is $\mathbf 1_{\iota_j+\varpi_w\auHom}$, the datum $\delta_\iota^\natural$ of
{\cite[Definition~6.31]{P53}}). By (W5), {\cite[Lemma~6.30]{P53}} ($L_{j,w}=a_w\auHom(\psM_{j,w},\psM_{j+1,w})$ and
$L_{j,\bar w}$ are dual for $[\,,]$) and finite Fourier inversion on $L_{j,w}/\varpi_wL_{j,w}$, with
$\iota^\natural_j=a_w\iota_j$ at each level place,
\begin{equation}\label{ps:eq-upsilon}
\Upsilon_{j,f}=\prod_w\frac1{\#_w}\sum_{a}[\iota^\natural_j,a]^{-1}\,\psrho^\circ(a)\psvac_{L_j},\qquad
\psrho(l)\Upsilon_{j,f}\in\C^\times\Upsilon_{j,f}\ \ (l\in L'_j),
\end{equation}
where $a$ runs over $\varpi_{\bar w}^{-1}L_{j,\bar w}/L_{j,\bar w}$ for each level place $w$,
$\#_w$ is the number of such $a$, and $L'_j$ is the self-dual lattice equal to $L_j$ outside the
level places and to $\varpi_wL_{j,w}\oplus\varpi_{\bar w}^{-1}L_{j,\bar w}$ at them. The forms on
$G_j$ \emph{of the archimedean type of the chain} are defined in Definition~\ref{pz:def-type}. We
use the following facts, which are Proposition~\ref{pz:types} (it is proved in \S\ref{pz:ss-types}
from three statements (A1)--(A3) on discrete series of real unitary groups): $\Theta^Y_{j,\vartheta}$
maps the forms on $G_j$ of the archimedean type of the chain to those on $G_{j+1}$; the
transposed lift $\Theta^{Y,*}_{j,\vartheta}z(g)=[K_j:J_j]\int_{[G_{j+1}]}\theta(\omega_{\mu_{j+1}}
(g,h)\Upsilon_j)z(h)\,dh$ maps the complex conjugates of these forms on $G_{j+1}$ to those on
$G_j$; and all these forms are cusp forms, hence rapidly decreasing.

\begin{lemma}\label{ps:level}
Let $k\in K_j$, $k'\in K_{j+1}$ have reductions $\bar k,\bar k'$ at the level places with
$\bar k'\iota_j=\iota_j\bar k$. Then $\omega_{\mu_{j+1}}(k,k')\Upsilon_{j,f}=\prod_w\psi_{j+1,w}
(\det\bar k'_w)\psi_{j+1,w}(\det\bar k_w)^{-1}\cdot\Upsilon_{j,f}$, where $\psi_{j+1,w}$ is the
tame part of $\mu_{j+1}$ at $w$. In particular $J_j\times J_{j+1}$ fixes $\Upsilon_{j,f}$.
\end{lemma}
\begin{proof}
Outside the level places $k,k'$ stabilise $L_j$ and all characters are unramified, so (W4)
applies. At a level place $w$, by (W5) the function $\mathbf 1_{a_w(\iota+\varpi\auHom)}$ is carried
to $x\mapsto\mathbf 1(k_w'^{-1}xk_w\in a_w(\iota+\varpi\auHom))=\mathbf 1(x\in a_w(k'_w\iota k_w^{-1}+\varpi
\auHom))$ ($a_w$ is a scalar), which is the same function by hypothesis, times $\mu_{j+1,w}(\det k'_w)\mu_{j+1,w}(\det
k_w)^{-1}$; the factor $\upsilon'_w$ is trivial on units: by Proposition~\ref{pwf:W5} (valid for ranks of either parity) it is built from $\upsilon_w$ and $|\cdot|_w$ evaluated at $\det k_w,\det k'_w\in\mathcal O_w^\times$, and $\upsilon$ is unramified at every
finite place.
\end{proof}

\begin{proof}[Proof of {\cite[Proposition~14.27]{P46}}]
The statement is proved in the following form, which is what the proof of
{\cite[Proposition~14.28]{P46}} uses: \emph{for every tame character $\vartheta$ and every $u\in\auG$
there is $\beta_\vartheta(u)\in K_m$ with $(\iota_\vartheta(\varrho(u)z))(g)=(\iota_\vartheta z)
(g\beta_\vartheta(u))$ for all $z\in\auAz$ and $g$.} (The element depends on $\vartheta$ through
the places of $Q_\vartheta$.) The level places are split, so the reduction $K_i\to\prod_w\auGL
(\psM_i/w)$ is onto. Choose $k_0\in K_0$ with reduction $u$ at the chosen places and $1$ at $Q$, and
$k_i\in K_i$ ($1\le i\le m$) with reduction $\operatorname{diag}(u,1)$ at the chosen places and $1$
at $Q$, in the frames; then $\bar k_{i+1}\iota_i=\iota_i\bar k_i$. By (W1) and
Lemma~\ref{ps:level}, for a form $F$ on $G_i$,
\[
(\Theta^Y_{i,\vartheta}F)(hk_{i+1})=[K_i:J_i]\!\int_{[G_i]}\!\theta\bigl(\omega(g,h)\omega(k_i,k_{i+1})
\omega(k_i^{-1},1)\Upsilon_i\bigr)F(g)\,dg=\varepsilon_i\,\bigl(\Theta^Y_{i,\vartheta}(F(\,\cdot\,k_i))
\bigr)(h),
\]
after the change of variable $g\mapsto gk_i$, with $\varepsilon_i=\prod_w\psi_{i+1,w}(\det\bar
k_{i+1,w})\psi_{i+1,w}(\det\bar k_{i,w})^{-1}=1$, because $\det\operatorname{diag}(u_w,1)=\det u_w$.
By induction $(\iota_\vartheta z)(gk_m)=\iota_\vartheta(z(\,\cdot\,k_0))(g)$, and
$z(\,\cdot\,k_0)=\varrho(u)z$ by the definition of $\varrho$. Put $\beta_\vartheta(u)=k_m$.
\end{proof}

\subsubsection{Branching and restriction: proof of \texorpdfstring{{\cite[Proposition~14.24]{P46}}}{[Row D, Proposition~14.24]}}
\label{ps:ss-restrict}

\emph{Conventions for $W_\lambda$.} For $\auGL_n(\C)$ and a diagram with $c(h)$ columns of height
$h<n$, the module $W_\lambda$ is the span, in the polynomials on $M_n(\C)$, of the products of
minors $\prod_{\text{columns}}\Delta_{I,[h]}$ (rows $I$, $|I|=h$, columns $1,\dots,h$), with
$(gp)(x)=p(g^tx)$, twisted by $\det^e$. It is the image of $\bigotimes_{\text{columns}}\Lambda^h
\C^n$ under the multiplication map $\mathsf m$, which is the orthogonal projection onto the Cartan
component when $W_\lambda$ carries the invariant form for which the highest weight vector
$\prod\Delta_{[h],[h]}^{c(h)}$ has norm $1$; this is the hermitian form of {\cite[\S14.3]{P46}}.
Hence a product of minors has norm at most $1$. The products of minors span a lattice
$\Sigma_\lambda$, with the standard ones as basis; $W_\lambda(\Z)$ is the dual lattice
$\Sigma_\lambda^\#$, with the dual basis, whose Gram matrix is integral (it is the matrix of
coefficients of the reproducing kernel $\prod\Delta_{[h],[h]}(\bar y^tx)^{c(h)}$ in the products
$b_i(x)\overline{b_j(y)}$). For a linear map $Z\colon\C^n\to\C^{n'}$ the Schur functor
$S(Z)\colon W_\lambda(\C^n)\to W_\lambda(\C^{n'})$ is the restriction of $\bigotimes\Lambda^hZ$; its
weight spaces for a diagonal $Z$ are orthogonal, so, with $s_1\ge s_2\ge\dots$ the singular values
of $Z$ and $l_p$ the rows of the diagram,
\begin{equation}\label{ps:eq-schurnorm}
\|S(Z)\|_{\mathrm{op}}=\prod_ps_p^{\,l_p},\qquad\sum_ps_p^2=\|Z\|_{\mathrm{HS}}^2 .
\end{equation}

\emph{The weight $\lambda^\circ$.} Put $\mathrm{kept}(h)=\tfrac12(\Delta-(t+1))$ if $h\in\{k,n_1,
\dots,n_{m-1}\}$ and $0$ otherwise; then $\mathrm{kept}(h)\le c_\tau(h)$. The weight
$\lambda^\circ$ on $\auU(Y_m)$ has, at $\tau$, $\mathrm{kept}(h)+c_\tau(h+1)-\mathrm{kept}(h+1)$
columns of height $h<n_m-1$ and the determinant exponent $e_{m,\tau}$; that is,
$\lambda^\circ_{\tau,p}=\lambda(X_m)_{\tau,p+1}+\mathrm{kept}(p)$: of the columns of
$\lambda(X_m)_\tau$ of height $h$, $\mathrm{kept}(h)$ keep their height and the others lose a box.
The two weights interlace, and
\begin{equation}\label{ps:eq-ntau}
n_\tau=|\lambda(X_m)_\tau|-|\lambda^\circ_\tau|=\sum_h\bigl(c_\tau(h)-\mathrm{kept}(h)\bigr)+e_{m,\tau}
=a\mathbf 1_{\tau\in\auPa}-\frac{m\Delta}2 .
\end{equation}
Indeed the sum is $a\mathbf 1_{\auPa}(\tau)+\kappa_1(\tau)-\tfrac{n_1}2+(m-1)(\Delta-(t+1))-\tfrac m2
(\Delta-(t+1))$ and $e_{m,\tau}=\tfrac{n_{m-1}}2-\kappa_1(\tau)-(m-1)\Delta$, with
$n_{m-1}-n_1=(m-2)(t+1)$. Summing over $\tau$ gives {\cite[(41)]{P46}}, with equality.

\begin{lemma}[The branching map]\label{ps:branching}
Let $L$ be the number of boxes of $\lambda^\circ$, summed over $\tau$, in the columns of height at
least $k$. There is a non-zero linear map $\mathrm{br}\colon W_{\lambda(X_m)}\to W_{\lambda^\circ}$
with:
\begin{enumerate}
\item[(a)] $\mathrm{br}\circ\rho(\operatorname{diag}(g,\zeta))=\prod_\tau\zeta_\tau^{\,n_\tau}\cdot
\rho^\circ(g)\circ\mathrm{br}$ for $g\in\prod_\tau\auGL_{n_m-1}(\C)$, $\zeta\in(\C^\times)^{\auPl}$;
\item[(b)] $\|\mathrm{br}\|_{\mathrm{op}}\le c_{\mathrm{br}}$ for a positive integer
$c_{\mathrm{br}}$;
\item[(c)] the matrix of $\mathrm{br}$ in the bases of $W_{\lambda(X_m)}(\Z)$ and
$W_{\lambda^\circ}(\Z)$ has integer entries, and $c_{\mathrm{br}}\le L^{3L}$.
\end{enumerate}
\end{lemma}
\begin{proof}[Proof of (a), (b)]
Let $v\in W_{\lambda(X_m)}$ be the product, over $\tau$ and over the columns, of the minors
$\Delta_{[h],[h]}$ for the columns that keep their height and $\Delta_{\{n_m\}\cup[h-1],[h]}$ for
those that lose a box. It is a highest weight vector of weight $\lambda^\circ$ for
$\prod_\tau\auGL_{n_m-1}$, of weight $\prod\zeta_\tau^{n_\tau}$ for the last diagonal entries by
\eqref{ps:eq-ntau} (the exponent $e_{m,\tau}$ comes from the twist), and $\|v\|\le1$. Let
$i\colon W_{\lambda^\circ}\to W_{\lambda(X_m)}$ be the equivariant map sending the highest weight
vector to $v$; it has norm $\|v\|\le1$, since both modules are irreducible and unitary. Put
$\mathrm{br}=c_{\mathrm{br}}\,i^*$ with $c_{\mathrm{br}}=\prod_\tau H_{\lambda^{\circ\flat}_\tau}$,
the product of the hook numbers (Lemma~\ref{pk:hook}) of the diagrams $\lambda^\circ_\tau$ without
their first $k-1$ rows.
\end{proof}
\noindent Part (c) is Proposition~\ref{pk:branching}. Parts (a) and (b) have just been proved without it, from Lemma~\ref{pk:hook}, whose proof does not use the present lemma; the proof of part (c) then starts from the vector $v$ constructed here.

\begin{proof}[Proof of {\cite[Proposition~14.24]{P46}}, given Lemma~\ref{ps:branching}]
The restriction is $(\operatorname{res}\psi)(y)=\prod_\tau\tau(\aubar{\pi_N})^{-n_\tau}\cdot
\mathrm{br}\bigl(\psi(\operatorname{diag}(y,1))\bigr)$; by (a) it is an algebraic modular form of
weight $\lambda^\circ$. Let $\sigma\in\auAut(\C)$. As $\mathrm{br}$ has an integral matrix,
$\sigma(\mathrm{br}\,\xi)=\mathrm{br}(\sigma\xi)$, so $\|\sigma(\mathrm{br}\,\xi)\|\le
c_{\mathrm{br}}\|\sigma\xi\|$. For every embedding $\tau'$ of $E$ in $\C$,
$|\tau'(\aubar{\pi_N})|^2=\tau'(\pi_N\aubar{\pi_N})=N$, because complex conjugation induces $c$ on
the CM field $E$; hence $|\sigma\prod_\tau\tau(\aubar{\pi_N})^{-n_\tau}|^2=N^{-\sum n_\tau}$.
Finally $L$ is small (at most $\Delta$ columns of each of at most $2m$ heights
$\ge k$, by the definition of $\lambda^\circ$; here $L\le2f\,n_m\,m\Delta$), so $L\le x^{60}$ and $c_{\mathrm{br}}^2\le
L^{6L}\le\exp(x^{70})$.
\end{proof}

\subsubsection{The size of one kernel: proof of \texorpdfstring{{\cite[Proposition~14.23]{P46}}}{[Row D, Proposition~14.23]}}\label{ps:ss-kernel}

\begin{lemma}[Conjugates of jets]\label{ps:conj}
Let $X,X'$ be standard of ranks $n<n'$, $W=\auHom(X,X')$, and $\sigma\in\auAut(\C)$ with
$\sigma|_K=1$. There are $t_\sigma\in\C$ with $|t_\sigma|^2=\Omega^{-1}$, a permutation
$\varpi_\sigma$ of $\auPl$, and a $\sigma$-semilinear map $\Lambda_\sigma$ of $\psHS_f(W)$ such
that for all $f\in\psHS_f(W)$:
\begin{enumerate}
\item[(i)] the Taylor series of $\Theta_{\Lambda_\sigma f}(t_\sigma\,\cdot)$ is obtained from that
of $\Theta_f(s_0\,\cdot)$ by applying $\sigma$ to the coefficients and $\varpi_\sigma$ to the
variables $z_\tau$;
\item[(ii)] if $f=\psrho^\circ(x)\omega(g)\psvac^\circ$ with $x\in W(\auA_f)$, $g\in\auU(W)
(\auA_f)$, then $|\theta(\psrho_\infty(y)(1\otimes\Lambda_\sigma f))|\le\exp\bigl(C_0(q_1nn'+1)^2
\bigr)$ for all $y$, where $C_0$ depends only on $q_0$ and $d$.
\end{enumerate}
\end{lemma}
\begin{proof} Proposition~\ref{pk:conj}, with $t_\sigma=\sigma(\varpi_{\mathrm{vac}}/\mathsf u)c_\psi/
\Omega_\sigma$ and an explicit $C_0(q_0,d)$, and Theorem~\ref{pf:recip}.\end{proof}

\begin{lemma}\label{ps:monomial}
Let $L\ge1$, $u\in\R$, $R\ge0$, let $\rho_i\ge0$ and $l_i\in\N$ ($i$ in a finite set) with
$\sum_i\rho_i\le e^u(2L+2)R$. Then $\prod_i\rho_i^{\,l_i}\le\exp\bigl(\sum_i(l_i/L+\log l_i!+l_iu)
+2LR\bigr)$.
\end{lemma}
\begin{proof}
With $v=u+1/L$ and $x^l\le l!\,e^x$ for $x=\rho_ie^{-v}$: $\rho_i^{l_i}\le l_i!\,e^{l_iv}\exp
(\rho_ie^{-v})$, and $\sum_i\rho_ie^{-v}\le(2L+2)R\,e^{-1/L}\le(2L+2)R/(1+1/L)=2LR$.
\end{proof}

\begin{proof}[Proof of {\cite[Proposition~14.23]{P46}}, given Lemma~\ref{ps:conj}]
Here $X=X_j$, $X'=X_{j+1}$, $n=n_j$, $n'=n_{j+1}$, $\phi=\phi^\circ_x$, the twist is trivial, and
$\theta_j(\phi;g,g')$ is \eqref{ps:eq-kernel}. Since $|\sigma w|=|\sigma'w|$ for all $w$, where
$\sigma'$ is $\sigma$ or its composite with complex conjugation, and one of them fixes $K$, we may
assume $\sigma|_K=1$.

\emph{Step 1: the datum.} The element $(g,g')$ acts on $W$ through $k\in\auU(W)(\auA_f)$, and
$\omega_{\mu_{j+1}}(g,g')\phi^\circ_x=\epsilon\,\psrho^\circ(kx)\omega(k)\psvac^\circ$ by (W1), where
$\epsilon$ is the product of a root of unity and of the values of the characters of
\S\ref{ps:ss-weil} at $\det g$, $\det g'$. These characters have algebraic values on finite adeles
and $|\sigma\epsilon|=1$: the values are products of roots of unity and of numbers
$\prod_\tau\tau(\nu)^{e_\tau}$ with $\nu\in E$, $\nu\aubar\nu=1$, all of whose conjugates have
absolute value $1$. So it suffices to bound $\sigma$ of the jets of $f=\psrho^\circ(x')\omega(k)
\psvac^\circ$ at scale $\pi_Ms_0$.

\emph{Step 2: a conjugate jet is a Gaussian integral.} Let $Q$ be a polynomial, $Q^\sigma$ the
polynomial obtained by applying $\sigma$ to its coefficients and $\varpi_\sigma$ to its variables.
By Lemma~\ref{ps:conj}(i) and \eqref{ps:eq-bargmann} for $\Lambda_\sigma f$ at scale $t_\sigma$,
\[
\sigma\bigl(Q(\partial)[\Theta_f(\pi_Ms_0\,\cdot)](0)\bigr)=\int Q^\sigma\bigl(\sigma(\pi_M)\,t_\sigma
\bar y\bigr)\,e^{-|y|^2/2}\,\theta\bigl(\psrho_\infty(y)(1\otimes\Lambda_\sigma f)\bigr)\,d\mu(y).
\]
\emph{Step 3: growth of the Schur matrix.} Let $\xi\in W_{\lambda(X_j)}$ and let $\eta$ be a
covector on $W_{\lambda(X_{j+1})}$ with $\|\sigma\eta\|\le1$ for the dual norm, and $Q=\eta\cdot
S(z)\xi$. The entries of $S$ have integer coefficients, so $Q^\sigma(z)=\sigma\eta\cdot S(z')
\sigma\xi$ with $z'$ a permutation of $z$, and by \eqref{ps:eq-schurnorm} $|Q^\sigma(z)|\le
\|\sigma\xi\|\prod_{\tau,p}s_{\tau,p}^{\,l_{\tau,p}}$ with $\sum_{\tau,p}s_{\tau,p}^2=|z|^2$. For
$z=\sigma(\pi_M)t_\sigma\bar y$ one has $|z|^2=(M/\Omega)|y|^2$, as $|\sigma\pi_M|^2=M$. Apply
Lemma~\ref{ps:monomial} to $\rho=s^2$ with $u=\log(M/\Omega)$, $R=|y|^2/(2L+2)$:
\[
|Q^\sigma(\sigma(\pi_M)t_\sigma\bar y)|\le\|\sigma\xi\|\exp\Bigl(\tfrac12\sum_{\tau,p}\bigl(
\tfrac{l_{\tau,p}}L+\ausinf(l_{\tau,p})\bigr)+\frac{L|y|^2}{2L+2}\Bigr),
\]
since $\log l!+l\log(M/\Omega)=\ausinf(l)$.

\emph{Step 4: conclusion.} By Steps 2, 3 and Lemma~\ref{ps:conj}(ii) the integrand is at most
$\|\sigma\xi\|\exp(C_0(q_1nn'+1)^2+\tfrac12\sum(l/L+\ausinf(l)))\cdot\exp(-|y|^2/(2L+2))$, and
$\int\exp(-|y|^2/(2L+2))\,d\mu=(2L+2)^{fnn'}$. Taking the supremum over $\eta$,
\[
\|\sigma(\theta_j(\phi;g,g')\xi)\|\le\exp\Bigl(C_0(q_1nn'+1)^2+fnn'\log(2L+2)+\frac1{2L}\sum_{\tau,p}
l_{\tau,p}+\tfrac12\mathrm{rs}_j\Bigr)\|\sigma\xi\| .
\]
Take $L=\sum_{\tau,p}l_{\tau,p}\le fn_m\cdot2a$. Then the third term is $\tfrac12$, and by
{\cite[Lemma~13.4]{P46}} ($f\le q_1\le x^{12}$, $nn'\le n_m^2\le1600x^{12}$, $\log(2L+2)\le3\log
b$) and $C_0\le\log\log b$ (Lemma~\ref{pz:C0}) the first two are at most $x^{50}-1$; see \S\ref{pz:ss-const}.
\end{proof}
The rows $l_{\tau,p}$ are: $a\mathbf 1_{\auPa}(\tau)+\kappa_1(\tau)-\tfrac{n_1}2+j(\Delta-(t+1))$
for $p=1$, the same without $a$ for $1<p\le k$, and at most $m\Delta$ for $p>k$; this is the
description of the rows in {\cite[\S14.8]{P46}}, with slack at most $\kappa_1(\tau)-\kappa(\tau)+j\Delta\le(m+1)\Delta$ per row.

\subsubsection{Gluing and chain data; proof of \texorpdfstring{{\cite[Proposition~14.25]{P46}}}{[Row D, Proposition~14.25]}}\label{ps:ss-glue}

Write $W_j=\auHom(X_j,X_{j+1})$, $L^X_j=\alpha\auHom(\widehat\auO^{\,n_j},\allowbreak \widehat\auO^{\,n_{j+1}})$,
and let $\mathsf P$ be the set of the $f$ primes of $E$ dividing $\pi_N$. The operators
$\psrho^\circ$ of (W4) are products of local operators, so $\psrho^\circ(x+x')=\psrho^\circ(x)
\psrho^\circ(x')$ when $x,x'$ have disjoint supports. A \emph{translation tuple} is
$x=(x_0,\dots,x_{m-1})$, $x_j\in W_j(\auA_f)$; it is \emph{admissible for $q$} if
$\aubar{\pi_M}x_j\in L^X_j$ at all places not dividing $q$. For $0\le i<m/2$ and a family
$\mathsf y=(\mathsf y_{b,w})$ of integers in $[0,N)$ indexed by the basis vectors $e'_b$ of
$X_{2i+1}$ and $w\in\mathsf P$, let $p_{\mathsf y}\in W_{2i}(\auA_f)$ be the map that kills
$e_1,\dots,e_{n_{2i}-1}$ and sends $e_{\mathrm{last}}$ to $-\alpha N^{-1}\sum_b\mathsf y_{b,w}e'_b$
at each $w\in\mathsf P$ and to $0$ at the other places, and $p'_{\mathsf y}\in W_{2i+1}(\auA_f)$
the map $e'_b\mapsto-\alpha N^{-1}\mathsf y_{b,w}\,e_{\mathrm{last}}$ at $\aubar w$, $w\in\mathsf P$,
and $0$ elsewhere. The \emph{gluing} of the $i$-th pair of lifts is
\[
\psGlue_i=\sum_{\mathsf y}\psrho(p_{\mathsf y})\otimes\psrho(p'_{\mathsf y})\quad
\text{on }\psHS_f(W_{2i})\otimes\psHS_f(W_{2i+1}),
\]
and $\psGlue=\bigotimes_i\psGlue_i$; the points $p_{\mathsf y}$, $p'_{\mathsf y}$ are supported on
$\mathsf P$, respectively $\overline{\mathsf P}$, and $\psrho$ of such points is given by
\eqref{pl:eq-heis}. For a tuple $x$ put $\phi^\circ_x=\bigotimes_j\phi^\circ_{j,
x_j}$, and for a finitely supported $\Psi$ put $C_{\vartheta,\Phi}=\Theta^X_{m-1}\cdots\Theta^X_0
(\Phi)$ for a tensor $\Phi$ (by linearity). A \emph{chain datum} for $\vartheta$, with
$q=\auNm Q_\vartheta$, is a finitely supported $\Psi$ such that: (C1) every $x$ in its support is
admissible for $q$ and vanishes at the places not dividing $Mq$; (C2) every $\Psi(x)$ is an
algebraic number integral at all primes not dividing $q$; (C3) if $Q_\vartheta=\emptyset$ then
$\sum_x|\sigma\Psi(x)|\le M^{3fmn_m^2}$ for all $\sigma$; (C4) the tensor $\psGlue(\sum_x\Psi(x)
\phi^\circ_x)$ is fixed by $\bigotimes_j\omega_{\mu_{j+1}}(b_j,b_{j+1})$ for all $b_j\in J_{X_j}$
($j<m$), $b_m=1$. The composite lift is $B_{\vartheta,\Psi}=\operatorname{res}\circ C_{\vartheta,
\psGlue(\sum\Psi(x)\phi^\circ_x)}$.

\begin{proof}[Proof of {\cite[Proposition~14.25]{P46}}]
Let $x$ be in the support of $\Psi$. By (C1) $x$ vanishes above $N$, while $p_{\mathsf y}$,
$p'_{\mathsf y}$ are supported above $N$; so $\psGlue_i(\phi^\circ_{2i,x_{2i}}\otimes\phi^\circ_{2i+1,
x_{2i+1}})=\sum_{\mathsf y}\zeta_{\mathsf y}\,\phi^\circ_{2i,x_{2i}+p_{\mathsf y}}\otimes
\phi^\circ_{2i+1,x_{2i+1}+p'_{\mathsf y}}$ with roots of unity $\zeta_{\mathsf y}$ ($\psrho^\circ$
and $\psrho$ differ by roots of unity). Hence $\psGlue(\sum_x\Psi(x)\phi^\circ_x)=\sum_{x'}\Psi'(x')
\phi^\circ_{x'}$ with
$\Psi'(x')=\sum\zeta\,\Psi(x)$ over the pairs $(x,(\mathsf y^{(i)})_i)$ with $x'=x+(p_{\mathsf y^{(i)}},
p'_{\mathsf y^{(i)}})_i$, $\zeta$ the product of the $\zeta_{\mathsf y^{(i)}}$, and
$B_{\vartheta,\Psi}=\operatorname{res}\circ C_{\vartheta,\Psi'}$. The
values of $\Psi'$ are sums of roots of unity times values of $\Psi$, so integral outside $q$; the tensor is the same,
so it is invariant under the levels; the tuples $x'$ are admissible for $qN$. The number of
families $(\mathsf y^{(i)})$ is $N^{f\sum_in_{2i+1}}\le N^{fmn_m/2}$, so for $Q_\vartheta=\emptyset$,
$\sum_{x'}|\sigma\Psi'(x')|\le M^{3fmn_m^2}N^{fmn_m/2}\le\exp(x^{60})$ by {\cite[(38)]{P46}}.
\end{proof}

\subsubsection{The seesaw identity: proof of \texorpdfstring{{\cite[Proposition~14.19]{P46}}}{[Row D, Proposition~14.19]}}\label{ps:ss-seesaw}

We keep the notation of \S\ref{ps:ss-Y}. For gradings on $V$ and $V'$ a map $w\in\auHom(V,V')$
\emph{has degree $\ge d$} if it maps the part of grade $a$ into the parts of grade $\ge a+d$;
$\auHom_{>0}$, $\auHom_{\ge0}$ denote the maps of degree $\ge1$, $\ge0$. The space
$I=\auHom_{>0}(V,V')$ is isotropic and $I^\perp=\auHom_{\ge0}(V,V')$. For $F\in\psHS(\auHom(V,V'))$
and $w\in\auHom(V,V')(\auA)$ put
\[
c_w(F)=\theta_{I^\perp/I}\bigl(\psZ_I\psrho(w)F\bigr)=\int_{[I]}\theta(\psrho(i)\psrho(w)F)\,di .
\]
By (W1), (W2): $c_w$ depends only on $w$ modulo $I(\auA)+I^\perp(E)$ up to the factors of
$\beta$; $c_{w+\delta}(F)=\beta(\delta,w)^{-1}c_w(F)$ for $\delta\in I(\auA)$; and if $p$ is a
unipotent element of the stabiliser of $I$ acting trivially on $I^\perp/I$, then
$c_w(\omega(p)F)=c_{p^{-1}w}(F)$.

\begin{proposition}[One step]\label{ps:onestep}
Let $0\le j<m$, $V=Y_j$, $V'=Y_{j+1}$, $T=T_j$, $F\in\psHS(\auHom(V,V'))$, $\psW$ an equivariant
function on $[R_j]$, and
\[
\psW'(h)=\int_{[R_j]}\psW(r)\,c_T\bigl(\omega_{\mu_{j+1}}(r,h)F\bigr)\,dr\qquad(h\in G_{j+1}(\auA)).
\]
Then $\psW'$ is continuous, left invariant under $R_{j+1}(E^+)$, and $\psW'(n'h)=\psi_{j+1}(n')
\psW'(h)$ for $n'\in N_{j+1}^{\ge2}(\auA)$; and for every rapidly decreasing automorphic function
$z$ on $G_{j+1}$, the integral $\int_{R_{j+1}(E^+)\backslash G_{j+1}(\auA)}\psW'z$ converges
absolutely and
\[
\int_{[R_j]}\psW(r)\int_{[G_{j+1}]}\theta\bigl(\omega_{\mu_{j+1}}(r,h)F\bigr)z(h)\,dh\,dr=
\int_{R_{j+1}(E^+)\backslash G_{j+1}(\auA)}\psW'(h)z(h)\,dh .
\]
\end{proposition}
\begin{proof}
For $j=m-1$ there are no strings on $V'$ and at most the vector $z$ on $V$: $I=0$, $T=0$, $c_T=
\theta$, and there is nothing to prove. Let $j<m-1$, $s=s_j$. Write $V=V^+\oplus V^0\oplus V^-$
by the sign of the grade, $p=\dim V^+$, and $T_+=T|_{V^+}$; $T_+$ is injective, and by
\eqref{ps:eq-TM} $R_{j+1}$ is the group of $h\in G_{j+1}$ that preserve $T(V^+)$ and act on it by a
unitriangular matrix in the basis $Tf_0,\dots,Tf_{p-1}$. Let $N_L\subset N_j$ be the unitriangular
group of $\auGL(V^+)$, $U^+$ the group of $u\in G_j$ acting trivially on $V^+$, on
$V^{+\perp}/V^+$ and on $V/V^{+\perp}$, and $\mathcal Z\subset U^+$ the subgroup acting trivially on
$V^0$; then $N_j=N_L\ltimes U^+$, and $N_L,\mathcal Z\subseteq N_j^{\ge2}$.

\emph{(1) Invariance.} For $n\in N_j^{\ge2}(\auA)$, $Tn-T=T(n-1)\in I(\auA)$, so
$c_T(\omega(n,1)F')=c_{Tn}(F')=\beta(T(n-1),T)^{-1}c_T(F')=\psi_j(n)^{-1}c_T(F')$, since
$\langle T,T(n-1)\rangle=\operatorname{tr}(\mathfrak e_j(n-1))$. Likewise $c_T(\omega(1,n')F')=
\psi_{j+1}(n')c_T(F')$ for $n'\in N^{\ge2}_{j+1}(\auA)$, and for $\delta\in R_{j+1}(E^+)$,
$\gamma\in R_j(E^+)$ the map $\delta^{-1}T\gamma-T$ has degree $\ge0$, so
$c_T(\omega(\gamma,\delta)F')=c_T(F')$ by (W1). This gives the properties of $\psW'$ (continuity
is clear, $[R_j]$ being compact).

\emph{(2) Coarse expansion.} Let $I_0\subseteq I$ be the maps vanishing on $V^+\oplus V^0$; then
$I_0^\perp$ is the space of maps vanishing on $V^+$, it is stable under $G_j'=P^+\times G_{j+1}$
($P^+$ the stabiliser of $V^+$), and by (W3)
$\theta(F')=\sum_{w\in\auHom(V^+,V')(E)}c^0_w(F')$, $c^0_w=\theta\psZ_{I_0}\psrho(\tilde w)$,
$\tilde w=w$ on $V^+$ and $0$ on $V^0\oplus V^-$, with $c^0_w(\omega(\gamma,\delta)F')=
c^0_{\delta^{-1}w\gamma}(F')$ for rational $\gamma\in R_j$, $\delta\in G_{j+1}$. For $u\in U^+
(\auA)$, $\tilde wu=\tilde w+\tilde w(u-1)$ and $\tilde w(u-1)\in I_0^\perp(\auA)$ has image
$w\circ\mathrm b(u)\in\auHom(V^0,V')$ in $I_0^\perp/I_0$, where $\mathrm b(u)\colon V^0\to V^+$ is
the component of $u-1$; so by (W2)
\begin{equation}\label{ps:eq-coarse}
c^0_w(\omega(u,1)F')=\psi\bigl(\tfrac12\jmath\operatorname{tr}(\tilde w^\dagger\tilde w(u-1))
\bigr)^{-1}\,\theta\bigl(\psrho(w\mathrm b(u))\,\psZ_{I_0}\psrho(\tilde w)F'\bigr),
\end{equation}
and $c^0_w(\omega(a,1)F')=c^0_w(F')$ for $a\in N_L(\auA)$ with $wa=w$.

\emph{(3) Only the orbit of $T_+$ contributes} to the integral over $[R_j]$ of the coarse
expansion; see below.

\emph{(4) The orbit of $T_+$.} The stabiliser $\mathsf M$ of $T_+$ in $G_{j+1}$ is contained in
$R_{j+1}$, and $\delta\mapsto a_\delta$, $\delta T_+=T_+a_\delta$, identifies
$\mathsf M\backslash R_{j+1}$ with $N_L$ (Witt). Hence $\{\delta^{-1}T_+:\delta\in\mathsf M(E^+)
\backslash R_{j+1}(E^+)\}=T_++J(E)$, $J=\auHom_{>0}(V^+,V')$ (this is correct only because the non-zero grades of $V$ and
$V'$ have different parities, remark (G1) of \S\ref{pk:ss-onestep}), and by (3)
\[
\int_{[R_j]}\psW(r)\theta(\omega(r,h)F)\,dr=\sum_{\gamma\in R_{j+1}(E^+)\backslash G_{j+1}(E^+)}
\int_{[R_j]}\psW(r)\sum_{w\in T_++J(E)}c^0_w(\omega(r,\gamma h)F)\,dr .
\]
\emph{(3), (5), (6).} Steps (3), (5) and (6) are carried out in \S\ref{pk:ss-onestep}
(Lemmas~\ref{pk:vanish}, \ref{pk:coarse}, \ref{pk:fine} and the proof given there).
\end{proof}

The statement is the global transfer of Bessel and Fourier--Jacobi periods under a theta lift for
a hook-type nilpotent orbit; compare Gomez--Zhu~\cite{GomezZhu} for the local statement. (See also B.~Wang, \emph{Global theta lifting and automorphic periods associated to
nilpotent orbits}, J.\ Number Theory 271 (2025), for related global statements.)

\emph{The functions $\psW_j$.} Fix $\vartheta$, $\varphi\in\auS$, and points $g_i\in G_i(\auA)$,
$0\le i<m$. Put $\psW_0=\psW_\varphi$ on $R_0(\auA)$ and, for $0\le j<m$,
\[
\psW_{j+1}[g_0,\dots,g_j](h)=\int_{[R_j]}\psW_j[g_0,\dots,g_{j-1}](rg_j)\,c_{T_j}\bigl(
\omega_{\mu_{j+1}}(rg_j,h)\Upsilon_j\bigr)\,dr\qquad(h\in G_{j+1}(\auA)),
\]
where for $j=0$ the first factor is $\psW_\varphi(r)$. By Proposition~\ref{ps:onestep}(1),
$r\mapsto\psW_j[\dots](rg_j)$ is equivariant on $[R_j]$, and by Lemma~\ref{ps:level} $\psW_{j+1}$ is
right invariant under $J_{j+1}$ and depends on $g_{i,f}$ only through $g_{i,f}J_i$.

\begin{lemma}[Two steps]\label{ps:double}
Let $j=2i<m$, $X=X_j$, $X'=X_{j+1}$, $X''=X_{j+2}$. There is a linear map $\psE_i\colon
\psHS_f(\auHom(Y_j,Y_{j+1}))\otimes\psHS_f(\auHom(Y_{j+1},Y_{j+2}))\to\psHS_f(W_j)\otimes
\psHS_f(W_{j+1})$, independent of $\vartheta$, with the following properties.
\begin{enumerate}
\item[(a)] For every linear form $l$ on $W_{\lambda(X)}$ and all $\xi\in G_j(\R)\times G_{j+1}(\R)
\times G_{j+2}(\R)$ there is a linear form $l''$ on $W_{\lambda(X'')}$ such that: if
$\psW_j[\dots](\iota(g)g_j)=l(g_\infty^{-1}\psi(g_f))$ for $g\in\auU(X)(\auA)$ and an algebraic
modular form $\psi$, then for all $k\in\auU(X'')(\auA)$, with $\iota(k)=k$ if $j+2=m$,
\[
\psW_{j+2}[\dots,g_j,g_{j+1}](\iota(k)g_{j+2})=l''\Bigl(k_\infty^{-1}\bigl(\Theta^X_{j+1}\Theta^X_j
(\Phi)\psi\bigr)(k_f)\Bigr),\quad\Phi=\psE_i\bigl(\omega(g_j,g_{j+1})\Upsilon_{j}\otimes\omega(g_{j+1},
g_{j+2})\Upsilon_{j+1}\bigr)_f,
\]
where $\xi$ is the archimedean part of $(g_j,g_{j+1},g_{j+2})$ and $g_m=1$.
\item[(b)] $\psE_i\circ(\omega(\iota b,\iota b')\otimes\omega(\iota b',\iota b''))=(\omega(b,b')
\otimes\omega(b',b''))\circ\psE_i$ for finite adelic points $b,b',b''$ of $\auU(X)$, $\auU(X')$,
$\auU(X_{j+2})$ (of $\auU(Y_m)$ if $j+2=m$).
\end{enumerate}
\end{lemma}
\begin{proof}
Let $V=Y_j$, $V'=Y_{j+1}$, $V''=Y_{j+2}$, $\mathbb V=V^-\oplus V''$ ($V^-$ is $V$ with the opposite
form), with the grading of $V$ and $V''$; $s=s_j$ is odd. By (W6) the map $w\mapsto w^\dagger$ identifies
$\psHS(\auHom(V,V'))$ with $\psHS(\auHom(V',V^-))$, carrying $\omega_{\auHom(V,V')}(g,g')$ to
$\upsilon(\det g)^{N'}\upsilon(\det g')^{-N}\omega_{\auHom(V',V^-)}(g',g)$ ($N,N'$ the ranks of
$V,V'$); so $\psHS(\auHom(V,V'))\otimes\psHS
(\auHom(V',V''))=\psHS(\auHom(V',\mathbb V))$ with this character, $\theta\otimes\theta$ is a non-zero multiple of
$\theta$, and $c_T(F_1)c_{T'}(F_2)$ is a multiple of $c_S(F_1\otimes F_2)$ for the isotropic space
$\mathbb I=\auHom_{>0}(V',\mathbb V)$ and $S=(T^\dagger,T')$, $T=T_j$, $T'=T_{j+1}$. By
\eqref{ps:eq-TM}, $S^\dagger S=-TT^\dagger+T'^\dagger T'=0$. The maps of degree $0$ are
$\auHom(X',\mathbb V_0)$; if $s\ge3$, $\mathbb V_0=X^-\oplus X_{j+2}\oplus H$ with $H$ the
hyperbolic plane spanned by the middle vectors $z,z''$, and $\ell=E\cdot Sf'_{p-1}$ ($f'_{p-1}$
the vector of grade $1$ of $V'$) is an isotropic line of $H$; if $s=1$, $\mathbb V_0=X^-\oplus
X_m$, where $X_m=Y_m\oplus E\,z/\pi_N$ and $z/\pi_N$, of norm $1$ in $V^-$, is
$e_{\mathrm{last}}$. Let $\mathbb J=\mathbb I\oplus\auHom(X',\ell)$ (resp.\ $\mathbb I$), so that
$\mathbb J^\perp/\mathbb J=\auHom(X',X^-\oplus X'')=W_j\oplus W_{j+1}$, and define
$\psE_i=\psZ_{\mathbb J,f}\circ\psrho_f(S)$.

By definition, $\psW_{j+2}(h)=\int_{[R_{j+1}]}\int_{[R_j]}\psW_j(rg_j)c_T(\omega(rg_j,r'g_{j+1})
\Upsilon_j)c_{T'}(\omega(r'g_{j+1},h)\Upsilon_{j+1})\,dr\,dr'$. Put $F=\omega(g_j,g_{j+1})\Upsilon_j
\otimes\omega(g_{j+1},g_{j+2})\Upsilon_{j+1}$ and $h=\iota(k)g_{j+2}$. By
Proposition~\ref{ps:onestep}(1) the integrand is invariant under $r\mapsto nr$, $n\in N_j(\auA)$
(here $N_j=N_j^{\ge2}$), so the integral over $[R_j]$ is one over $[\auU(X)]$. For $n'\in
N_{j+1}(\auA)$, $Sn'-S=S(n'-1)$ has degree $\ge0$, its component of degree $0$ is the map
$\lambda(n')\in\auHom(X',\ell)$, $x\mapsto(\text{coefficient of }f'_{p-1}\text{ in }n'x)\,Sf'_{p-1}$,
and $\beta(S(n'-1),S)=1$ because $S^\dagger S=0$. Hence, by (W2), the product of the two constant
terms at $(g,n'r',k)$, $g\in\auU(X)$, $r'\in\auU(X')$, is a constant times $\theta(\psrho(\lambda
(n'))\psZ_{\mathbb I}\psrho(S)\omega(g,r',k)F)$; $n'\mapsto\lambda(n')$ maps $[N_{j+1}]$ onto
$[\auHom(X',\ell)]$ with the Haar measure as image, so integration over $[N_{j+1}]$ gives
$\theta_{W_j\oplus W_{j+1}}(\psZ_{\mathbb J}\psrho(S)\omega(g,r',k)F)$. The elements $g,r',k$ fix $S$
and lie in the Levi subgroup of the stabiliser of $\mathbb J$; by (W2), $\psZ_{\mathbb J}\omega(\iota
g,\iota r',\iota k)=\chi_0(g,r',k)\,\omega(g,r',k)\psZ_{\mathbb J}$, where $\chi_0$ is a character of
the three determinants, built from $\upsilon$ and the ranks only. One has $\chi_0=1$ (Lemma~\ref{pk:chi0}). By (W2), $\psZ_{\mathbb J}
\psrho(S)F=u_\xi\otimes\psE_i(F_f)$ with a Fock vector $u_\xi$ depending only on $\xi$ and the
$P_j$. So
\[
\psW_{j+2}(\iota(k)g_{j+2})=\mathrm{const}\int_{[\auU(X')]}\int_{[\auU(X)]}l(g_\infty^{-1}\psi(g_f))\,
\theta\bigl(\omega(g,r',k)(u_\xi\otimes\Phi)\bigr)\,dg\,dr',
\]
and (a) is Lemma~\ref{ps:definite}. (b) is the same intertwining property of $\psZ_{\mathbb J,f}$,
with $\iota(b)S\iota(b')^{-1}=S$.
\end{proof}

\begin{remark}\label{ps:chi0}
The identity $\chi_0=1$ is Lemma~\ref{pk:chi0}; it uses that all $n_j$ are even and that $s_j$ is
odd for $j$ even.
\end{remark}

\begin{lemma}[Support]\label{ps:support}
Let $\varphi\in\auSram$, $g_{0,f}\in J_M$, $1\le j<m$, and let $\mathfrak B_j$ be the set of
$x\in G_j(\auA_f)$ for which there is $r\in R_j(\auA_f)$ such that $\Lambda'=rx\psM_j$ is the
direct sum of its intersections with the grades and $\mathsf q^{\,j}\Lambda'^-\subseteq\psM_j^-$,
$\aubar{\mathsf q}^{\,j}\psM_j^-\subseteq\Lambda'^-$, where $\mathsf q$ is an idele generating
$\prod_{w\in Q}w$ and $\Lambda'^-$ is the intersection with the span of the string. If
$\psW_j[g_0,\dots,g_{j-1}](h)\ne0$ then $h_f\in\mathfrak B_j$. Moreover $\mathfrak B_j=R_j(\auA_f)
\mathsf C_jK_j$ for a finite set $\mathsf C_j$ of elements that are $1$ at all places not dividing
$q$, and $\mathfrak B_j=R_j(\auA_f)K_j$ if $Q=\emptyset$.
\end{lemma}
\begin{proof}
Induction on $j$; put $\mathfrak B_0=J_M$. If $\psW_j(h)\ne0$ there is $r$ with
$\psW_{j-1}(rg_{j-1})\ne0$ and $c_T(\omega(rg_{j-1},h)\Upsilon_{j-1})\ne0$, $T=T_{j-1}$. By induction
$(rg_{j-1})_f=r_1g^\circ$ with $r_1\in R_{j-1}(\auA_f)$ and $L=g^\circ\psM_{j-1}$ graded with the
bounds. By the Iwasawa decomposition for the parabolic subgroup of $G_j$ with unipotent radical
$N_j$ (see (R1) of \S\ref{pk:ss-rest}; $K_{j,v}$ is hyperspecial because $E/E^+$ is unramified and
$\psM_j$ is unimodular), $h_f=n\,l\,k$ with $n\in N_j(\auA_f)$, $k\in K_j$ and $l$ preserving the grades; so
$\Lambda'=l\psM_j$ is graded. As in Proposition~\ref{ps:onestep}(1), $r_1$ and $n$ change $T$ by
maps of degree $\ge0$ and act by invertible operators after $\psZ_I$, so
$\psZ_{I,f}\psrho(T)\omega(g^\circ,lk)\Upsilon_{j-1,f}\ne0$. By Lemma~\ref{ps:level} and
\eqref{ps:eq-upsilon}, $\omega(g^\circ,lk)\Upsilon_{j-1,f}$ is a multiple of $\psrho(x_0)
\psvac_{\mathcal L}$ with $\mathcal L=lL'_{j-1}g^{\circ-1}$ and $x_0\in\alpha\auHom(L,\Lambda')$; and
$\psZ_{I,f}\psrho(y)\psvac_{\mathcal L}=0$ unless $y\in\mathcal L+I^\perp(\auA_f)$ (for $i\in
I\cap\mathcal L$, $\psrho(i)$ fixes $\psvac_{\mathcal L}$ up to a scalar and $\psZ_I\psrho(i)=\psZ_I$,
while $\psZ_I\psrho(i)\psrho(y)=[i,y]\psZ_I\psrho(y)\psrho(i)$; so $[i,y]=1$ on $I\cap\mathcal L$,
and $I\cap\mathcal L$ is the annihilator of $\mathcal L+I^\perp$). Hence $T\in\alpha\,\aubar{c_M
\mathsf q}^{-1}\auHom(L,\Lambda')+\auHom_{\ge0}$; since $L$, $\Lambda'$ are graded and $T$ is
homogeneous, $T(L)\subseteq\psfa\,\aubar{\mathsf q}^{-1}\Lambda'$, that is $T(L^+)\subseteq\psfa\,
\aubar{\mathsf q}^{-1}\Lambda'^-$. Now $L^+=(L^-)^\#\supseteq\aubar{\mathsf q}^{\,j-1}\psM_{j-1}^+$
by the induction hypothesis, and \eqref{ps:eq-TM} gives $\aubar{\mathsf q}^{\,j}\psM_j^-\subseteq
\Lambda'^-$. Dually $\Lambda'^-\subseteq\Lambda'^+\subseteq\mathsf q^{-j}\psM_j^+$, and
$\psM_j^+=\psM_j^-$ unless $j$ is even, where they differ by $u_jz/N$ at the places above $N$. Let
$j$ be even and $v$ a place of $E^+$ below $w\in\mathsf P$. The line $\Lambda'\cap E_vz$ contains
$\widehat\auO u_jz$ and is integral; if it is larger, it is unimodular, so the part of grade $0$ of
$\Lambda'_v$ is the sum of its intersections with $X_j$ and with $E_vz$, and the scalar
$\iota(\mathsf z_w)$ ($\mathsf z_w$ the idele $z_w$ of {\cite[\S14.3]{P46}}) stabilises
$\Lambda'$; then $u=h_f^{-1}n\iota(\mathsf z_w)n^{-1}h_f\in J_j$, and by right invariance
$\psW_j(n\iota(\mathsf z_w^{\,e})n^{-1}h)=\psW_j(h)$ for all $e$. On the other hand, by
Lemma~\ref{ps:double}(a) applied $j/2$ times, $\psW_j(\gamma\iota(g)h)=\psi_j(\gamma)l(g_\infty^{-1}
(C\varphi)(g_f))$ for $\gamma\in N_j(\auA)$, $g\in\auU(X_j)(\auA)$, with $C$ a composite of lifts;
the central element $\mathsf z_w$ commutes with the lifts (it acts on $\auHom(X_i,X_{i+1})$
trivially, and the characters are unramified above $N$), so $\sum_{e=0}^{N-2}\psW_j(n\iota(\mathsf
z_w^{\,e})n^{-1}h)=\psi\cdot l((C\sum_e\mathsf z_w^{\,e}\varphi)(1))=0$ for ramified $\varphi$
(see (R2) of \S\ref{pk:ss-rest}: $n\iota(\mathsf z_w^{\,e})n^{-1}=n'\iota(\mathsf z_w^{\,e})$ with
$n'\in N_j(\auA)=N_j^{\ge2}(\auA)$ and $\psi_j(n')=1$). So $(N-1)\psW_j(h)=0$, a contradiction; hence $\Lambda'\cap E_vz=\widehat\auO u_jz$
and $\mathsf q^{\,j}\Lambda'^-\subseteq\psM_j^-$. The last assertion is Lemma~\ref{pf:orbits}.
\end{proof}

\begin{lemma}[Finite data]\label{ps:fin}
There is $C_1\in\C$, independent of $\vartheta$, with the following property. Let
$q=\auNm Q_\vartheta$ and let $\mathsf y=(\mathsf y_0,\dots,\mathsf y_m)$, $\mathsf y_j\in
G_j(\auA_f)$, be integral at all places not dividing $q$, with $\mathsf y_m=1$. Then
\[
\Phi(\mathsf y):=\bigotimes_{i<m/2}\psE_i\bigl(\omega(\mathsf y_{2i},\mathsf y_{2i+1})\Upsilon_{2i,f}
\otimes\omega(\mathsf y_{2i+1},\mathsf y_{2i+2})\Upsilon_{2i+1,f}\bigr)=C_1\sum_xr_{\mathsf y}(x)\,
\psGlue(\phi^\circ_x),
\]
a finite sum over tuples $x$ as in (C1), where $B^{2m}r_{\mathsf y}(x)$, $B=M^{fn_m^2}$, is an
algebraic number integral at the primes not dividing $q$, and $\sum_x|\sigma r_{\mathsf y}(x)|\le1$
for all $\sigma\in\auAut(\C)$ if $Q_\vartheta=\emptyset$.
\end{lemma}
\begin{proof}
It suffices to treat one pair $i$, with $j=2i$ and the notation of the proof of
Lemma~\ref{ps:double}. \emph{(1)} By \eqref{ps:eq-upsilon}, (W1), (W4), each
$\omega(\mathsf y_j,\mathsf y_{j+1})\Upsilon_{j,f}$ is a combination $\sum_ar(a)\psrho^\circ(a)
\psvac_{L_j}$ with $\aubar{\pi_M}a\in L_j$ outside $q$: at the places outside $q$ the $\mathsf y$
stabilise $L_j$ and permute the translates up to roots of unity and unit values of characters; at
the places of $q$, $\omega(g)\psvac_L$ is a multiple of $\psvac_{gL}=\sum_{r\in gL/(L\cap gL)}
\zeta_r\psrho^\circ(r)\psvac_L$ with a factor whose denominator divides a power of $q$. The
coefficients $r(a)$ are roots of unity times $\prod_w\#_w^{-1}$ outside $q$; $\prod_{w\text{
chosen}}\#_w$ divides $B$; and if $Q=\emptyset$ then $\sum|\sigma r(a)|\le1$. Under the
identification (W6) of the first factor, $w\mapsto w^\dagger$ carries $L_j$ to
$L_j^\dagger=\aubar\alpha\auHom(\psM_{j+1},\psM_j^-)=\alpha\auHom(\psM_{j+1},u\psM_j^-)$, $u=\aubar a/a$
(Remark~\ref{pf:rem-flip-lattice}; $u_v=1$ at the places not above $q_1$), and $\mathrm{Fl}$, normalised
locally to carry $\psvac_{L_{j,v}}$ to $\psvac_{L^\dagger_{j,v}}$ (possible by that remark), carries
translates to translates up to roots of unity. So the tensor of two translates is a root of unity
times a translate $\psrho^\circ(a)\psvac_{\mathbb L}$,
$\mathbb L=L_j^\dagger\oplus L_{j+1}=\alpha\auHom(\psM_{j+1},u\psM_j^-\oplus\psM_{j+2})$.

\emph{(2)} $\psE_i(\psrho^\circ(a)\psvac_{\mathbb L})=\psZ_{\mathbb J,f}\psrho(S+a)\psvac_{\mathbb L}$
up to a root of unity. As in the proof of Lemma~\ref{ps:support}, this vanishes unless
$S+a\in\mathbb L+\mathbb J^\perp$, and then it is a root of unity times
$\psrho^\circ(a_0)\psZ_{\mathbb J,f}\psvac_{\mathbb L}$, $a_0$ the image in $W_j\oplus W_{j+1}$ of
an element of $(S+a+\mathbb L)\cap\mathbb J^\perp$. The vector $\psZ_{\mathbb J,f}\psvac_{\mathbb L}$
is not zero and is fixed, up to scalars, by $\psrho(\bar l)$, $l\in\mathbb L\cap\mathbb J^\perp$; so
it is $C\,\psvac_{\mathbb L_0}$ with $\mathbb L_0$ the image of $\mathbb L\cap\mathbb J^\perp$ (a
self-dual lattice), and $C$ depends on the pair $i$ only.

\emph{(3) The lattice $\mathbb L_0$.} Since $\psM_{j+1}$ is the sum of its parts in $X'$ and in
the string, $\mathbb L_0=\alpha\auHom(\widehat\auO^{\,n'},\Lambda_0)$ with $\Lambda_0\subset
(X^-\oplus X'')(\auA_f)$ the image of the vectors of $u\psM_j^-\oplus\psM_{j+2}$ without components
of negative grade and with component of grade $0$ in $X^-\oplus X_{j+2}\oplus\ell$. At the places
above $q_1$ (where $u_v\ne1$ is possible) there is no prime above $N$ and no chosen place
({\cite[Lemma~6.32(a)]{P53}}), the parts of $\psM_j$, $\psM_{j+2}$ in $X_j$, $X_{j+2}$ are the standard
lattices (and for $j+2=m$ the factor $u\,u_{m-2}$ of the vector $z/\pi_N$ is a unit) and $\Lambda_{0,v}$ is
$u_v$ times the standard lattice of $X^-$ plus the standard lattice of $X''$, so that $\mathbb
L_{0,v}=(L^X_{j,v})^\dagger\oplus L^X_{j+1,v}$ is exactly the image of $L^X_{j,v}\oplus L^X_{j+1,v}$
under $w\mapsto w^\dagger$ on the first factor. Elsewhere $u_v=1$ and, by
\eqref{ps:eq-lattice}, with $y,y''$ the last vectors of $X,X''$: outside $N$ and the chosen
places, $\Lambda_0$ is the standard lattice; at the chosen places it is the standard lattice
except that for $j+2=m$ the vector $y''=z/\pi_N$ is multiplied by the norm-one idele $u_{m-2}$
(whose denominator divides $\aubar{\pi_M}$) (see (R3) of \S\ref{pk:ss-rest}: at the chosen places
$\psM_j$ is the orthogonal sum of its parts in $X_j$ and in the string, so the image is standard for
$j+2<m$ whatever $u_j,u_{j+2}$ are); and at a place $v=w\aubar w$, $w\in\mathsf P$, where
$u_j=u_{j+2}=1$ and $Sf'_{p-1}=c(z+z'')$ (resp.\ $\psM_{m-2}\ni(y+z/\pi_N)/\aubar{\pi_N}$, $z$),
\[
\Lambda_{0,v}=\bigl\{\textstyle\sum a_ie_i+\beta y+\sum a''_ie''_i+\beta''y'':\ \beta''\equiv\beta\
(\varpi_w)\text{ at }w,\ \ (\beta,\beta'')\in\auO_{\bar w}^2+\varpi_{\bar w}^{-1}\auO_{\bar w}(1,1)
\text{ at }\aubar w\bigr\},
\]
the neighbour of the standard lattice along the isotropic vector $y+y''$. Hence
$\mathbb L_0/(\mathbb L_0\cap L^X)$, $L^X=L^X_j\oplus L^X_{j+1}$ (read in $\auHom(X',X^-\oplus X'')$ through $w\mapsto w^\dagger$ on the first factor), is represented by the sums
$t+\mathsf p_{\mathsf y}$, where $t$ runs over a finite set $\mathsf T_i$ of elements vanishing
above $N$ (trivial unless $j+2=m$, of order dividing $B$), and $\mathsf p_{\mathsf y}$ is the map
$x'\mapsto-\alpha N^{-1}\langle\mathsf y,x'\rangle(y+y'')$ at the $\aubar w$; under
$W_j\oplus W_{j+1}=\auHom(X',X^-\oplus X'')$, $w\mapsto w^\dagger$, $\mathsf p_{\mathsf y}$ is the
pair $(p_{\mathsf y},p'_{\mathsf y})$ of \S\ref{ps:ss-glue}. The $\mathsf p_{\mathsf y}$ are
supported at the $\aubar w$, so they are mutually isotropic, $\psrho^\circ=\psrho$ on them, and
they commute with the $\psrho^\circ(t)$. Therefore $\sum_{t,\mathsf y}\zeta_t\psrho^\circ(t)
\psrho^\circ(\mathsf p_{\mathsf y})\psvac_{L^X}$, for suitable roots of unity $\zeta_t$, is fixed up
to characters by $\mathbb L_0$ and is not zero (the translates are independent): it is a multiple
of $\psvac_{\mathbb L_0}$, and
\[
\psZ_{\mathbb J,f}\psvac_{\mathbb L}=C'\,\frac1{\#\mathsf T_i}\sum_{t\in\mathsf T_i}\zeta_t\,
\psGlue_i\bigl(\phi^\circ_t\bigr),\qquad\phi^\circ_t=\psrho^\circ(t)(\psvac^\circ\otimes\psvac^\circ).
\]
\emph{(4)} In (2), $a_0$ is integral for $\mathbb L_0$ above $N$, so $\psrho^\circ(a_0)
\psvac_{\mathbb L_0}$ is a root of unity times $\psrho^\circ(a_0')\psvac_{\mathbb L_0}$ with $a_0'$
vanishing at all places not dividing $Mq$; $\psrho^\circ(a_0')$ commutes with the
$\psrho^\circ(\mathsf p_{\mathsf y})$, and $\psrho^\circ(a'_0)\psrho^\circ(t)=\zeta\psrho^\circ
(a_0'+t)$. Altogether $\psE_i(\dots)$ is $C'$ times a combination of the $\psGlue_i(\phi^\circ_x)$
with coefficients that are roots of unity times products of the $r(a)$ and of $1/\#\mathsf T_i$;
the product of the denominators over the $m/2$ pairs divides $B^{m+1}$, hence $B^{2m}$ (see (R4)
of \S\ref{pk:ss-rest}), and the sizes multiply.
\end{proof}

\begin{proof}[Proof of {\cite[Proposition~14.19]{P46}}]
Fix $g_0\in G_0(\R)$ and $\mathsf t$. Let $\eta$ be a form on $G_m$ of level $J_m$ and weight
$\lambda^\circ$, $z_m$ the complex conjugate of the corresponding function on $[G_m]$, and
$z_j=\Theta^{Y,*}_{j,\vartheta}z_{j+1}$. By Fubini (the forms are rapidly decreasing and the theta
kernels of moderate growth) $\int_{[G_m]}z_m\,\iota_\vartheta(\auP\varphi)=\int_{[G_0]}z_0\,\auP
\varphi$. \emph{Step 1.} The function $z_0$ has level $J_0$ (Lemma~\ref{ps:level}); its trace
$\sum_{\gamma\in J_M/J_0}z_0(\,\cdot\,\gamma)$ is the conjugate of an element of $\auAz$, so by the
definition of $\auP$
\[
[J_M:J_0]\int_{[G_0]}z_0\,\auP\varphi=\sum_{\gamma\in J_M/J_0}\int_{[R_0]}\psW_\varphi(r)z_0(rg_0
\gamma)\,dr .
\]
\emph{Step 2.} Put $\Pi_j(x_0,\dots,x_j)=\int_{[R_j]}\psW_j[x_0,\dots,x_{j-1}](rx_j)z_j(rx_j)\,dr$
for $j<m$. By Proposition~\ref{ps:onestep} with $F=\omega(x_j,1)\Upsilon_j$ and
$\psW=\psW_j(\,\cdot\,x_j)$,
\begin{align*}
\Pi_j(x_0,\dots,x_j)&=[K_j:J_j]\int_{R_{j+1}(E^+)\backslash G_{j+1}(\auA)}\psW_{j+1}[x_0,\dots,x_j]
\,z_{j+1}\\&=[K_j:J_j]\int_{R_{j+1}(\auA)\backslash G_{j+1}(\auA)}\Pi_{j+1}(x_0,\dots,x_{j+1})\,dx_{j+1}
\end{align*}
for $j+1<m$, absolutely convergent, and for $j=m-1$ the first integral is over $[G_m]$.
\emph{Step 3 (the top).} By Lemma~\ref{ps:double}(a), applied $m/2$ times starting from
$\psW_\varphi(\iota(g))=\mathsf t(g_\infty^{-1}\varphi(g_f))$, for $h\in G_m(\auA)$
\[
\psW_m[x_0,\dots,x_{m-1}](h)=l_\xi\bigl(h_\infty^{-1}(C_{\vartheta,\Phi(\mathsf y)}\varphi)
(\operatorname{diag}(h_f,1))\bigr),
\]
where $\xi$, $\mathsf y$ are the archimedean and finite parts of $(x_0,\dots,x_{m-1},1)$ and
$l_\xi$ is a linear form on $W_{\lambda(X_m)}$ depending on $\xi$, $g_0$, $\mathsf t$ only. The
group $G_m(\R)$ is compact and $W_{\lambda^\circ}$ occurs once in $W_{\lambda(X_m)}$ restricted to
it; by Schur orthogonality there is a linear form $\mathsf c$ on the dual of $W_{\lambda(X_m)}$,
independent of everything else, with $\int_{[G_m]}\psW_mz_m=\mathsf c(l_\xi)\,\langle
\operatorname{res}C_{\vartheta,\Phi(\mathsf y)}\varphi,\eta\rangle$ (the factor
$\prod\tau(\aubar{\pi_N})^{n_\tau}$ and the dimension are put in $\mathsf c$). \emph{Step 4 (the
outer integrals).} Thus $\Pi_{m-1}=[K_{m-1}:J_{m-1}]\,\mathfrak a_{m-1}(\xi)\,S_{m-1}(\mathsf y)$ with
$\mathfrak a_{m-1}=\mathsf c(l_\xi)$ and $S_{m-1}(\mathsf y)=\langle\operatorname{res}C_{\vartheta,
\Phi(\mathsf y)}\varphi,\eta\rangle$. Suppose $\Pi_{j+1}=[K_{j+1}:J_{j+1}]\mathfrak a_{j+1}(\xi)
S_{j+1}(\mathsf y)$, where $\mathfrak a_{j+1}$ does not depend on $\vartheta,\varphi,\eta$. The
function $\Pi_{j+1}$ is left invariant under $R_{j+1}(\auA)$ in $x_{j+1}$, right invariant under
$J_{j+1}$, and vanishes unless $x_{j+1,f}\in\mathfrak B_{j+1}$ (Lemma~\ref{ps:support}; for
$x_{0,f}\in J_M$). Put $H_{j+1}=R_{j+1}(\auA_f)\cap K_{j+1}$. The group
$R_{j+1}(\auA_f)$ has finitely many orbits in $\mathfrak B_{j+1}/J_{j+1}=R_{j+1}(\auA_f)\mathsf
C_{j+1}K_{j+1}/J_{j+1}$, because $\mathsf C_{j+1}$ and $K_{j+1}/J_{j+1}$ are finite. In each orbit
choose a point of $\mathsf C_{j+1}K_{j+1}/J_{j+1}$, and let $\Omega_{j+1}$ be the union of the
$H_{j+1}$-orbits of the chosen points: a finite subset of $\mathfrak B_{j+1}/J_{j+1}$, stable under
left multiplication by $H_{j+1}$, contained in $H_{j+1}\mathsf C_{j+1}K_{j+1}/J_{j+1}$, which every
$R_{j+1}(\auA_f)$-orbit meets in exactly one $H_{j+1}$-orbit. (For $Q=\emptyset$, $\mathsf
C_{j+1}=\{1\}$ and $\Omega_{j+1}$ is a subset of $K_{j+1}/J_{j+1}$, in general not contained in
$H_{j+1}J_{j+1}/J_{j+1}$.) For $t\in\Omega_{j+1}$ let $m(t)=[\mathsf S_t:\mathsf S_t\cap K_{j+1}]$,
$\mathsf S_t$ the stabiliser of $t$ in $R_{j+1}(\auA_f)$. With the measures giving volume $1$ to
$K_{j+1}$ and to $H_{j+1}$, every function $\mathsf b$ on $R_{j+1}(\auA_f)\backslash
\mathfrak B_{j+1}/J_{j+1}$ has $\int\mathsf b=[K_{j+1}:J_{j+1}]^{-1}\sum_{t\in\Omega_{j+1}}\mathsf
b(t)/m(t)$: the orbit of $t$ has volume $[K:J]^{-1}\operatorname{vol}(\mathsf S_t)^{-1}$ (the
groups $G_{j+1}(\auA_f)$ and $R_{j+1}(\auA_f)$ are unimodular and $tJ_{j+1}$ has volume
$[K:J]^{-1}$), and $\operatorname{vol}(\mathsf S_t)^{-1}=[H_{j+1}:\mathsf S_t\cap K_{j+1}]/m(t)$ is
the number of elements of $\Omega_{j+1}$ in the orbit, which is the number of elements of
$H_{j+1}t$, divided by $m(t)$. Therefore
\[
\Pi_j=[K_j:J_j]\Bigl(\int_{R_{j+1}(\R)\backslash G_{j+1}(\R)}\mathfrak a_{j+1}\,d\xi_{j+1}\Bigr)
\sum_{t\in\Omega_{j+1}}\frac{S_{j+1}(\mathsf y_0,\dots,\mathsf y_j,t)}{m(t)}
\]
whenever the sum is not zero, the integral being then absolutely convergent by Step~2; define
$\mathfrak a_j(\xi_0,\dots,\xi_j)$ as this integral when it converges absolutely and as $0$
otherwise, and $S_j$ as the sum. Then $\Pi_j=[K_j:J_j]\mathfrak a_jS_j$ in all cases. By descending
induction and Step~1,
\[
\int_{[G_m]}z_m\,\iota_\vartheta(\auP\varphi)=c_0\,\bigl\langle\operatorname{res}C_{\vartheta,
\Phi_\vartheta}\varphi,\eta\bigr\rangle,\quad\Phi_\vartheta=\sum_{\gamma\in J_M/J_0}\,\sum_{t_1\in
\Omega_1}\!\cdots\!\sum_{t_{m-1}\in\Omega_{m-1}}\frac{\Phi(\gamma,t_1,\dots,t_{m-1},1)}{m(t_1)\cdots
m(t_{m-1})},
\]
with $c_0=[K_0:J_M]\,\mathfrak a_0(g_0)$ independent of $\vartheta$, $\varphi$, $\eta$. As $\eta$
is arbitrary, $\iota_\vartheta(\auP\varphi)=c_0\operatorname{res}C_{\vartheta,\Phi_\vartheta}\varphi$
for all $\varphi\in\auSram$.

\emph{Step 5 (the chain datum).} The elements of $\Omega_j$ are integral outside $q$, and $m(t)$
divides a power of $q$: $\mathsf S_t$ is a product of local groups, contained in $K_j$ outside $q$,
and at $v\mid q$ it is conjugate to a subgroup of the pro-$p$ group $J_{j,v}$. If
$Q=\emptyset$ then $\Omega_j\subseteq K_j/J_j$, $m=1$ and $\#\Omega_j\le B$. Indeed, for
$Q=\emptyset$ the group $J_j$ is the kernel of the reduction $K_j\to\auGL(\psM_j/\pi_M\psM_j)$, so
$[K_j:J_j]\le\#\operatorname{End}(\psM_j/\pi_M\psM_j)=(\#\auO_E/\pi_M\auO_E)^{(\dim Y_j)^2}=
M^{f(\dim Y_j)^2}$, since $|\sigma\pi_M|^2=M$ for every embedding $\sigma$; and
$\dim Y_j\le\dim Y_m=n_m-1$, the frame of $Y_j$ being the beginning of that of $Y_{j+1}$. Hence
$\#\Omega_j\le[K_j:J_j]\le B$, and likewise $[J_M:J_0]\le[K_0:J_0]\le B$, so $\Phi_\vartheta$ is a
sum of at most $B^m$ terms $\Phi(\mathsf y)/(m(t_1)\cdots m(t_{m-1}))$ with all $m(t_j)=1$. By
Lemma~\ref{ps:fin}, applied to each term, $\Phi_\vartheta=C_1\sum_xr_\vartheta(x)\psGlue(\phi^\circ_x)$ with
$B^{2m}r_\vartheta$ integral outside $q$ and $\sum_x|\sigma r_\vartheta(x)|\le B^{m}$ if
$Q=\emptyset$. Put $\Psi=B^{2m}r_\vartheta$. Then (C1)--(C3) hold, since $B^{3m}=M^{3fmn_m^2}$. For
(C4), let $b_j\in J_{X_j}$, $b_m=1$. The element $\iota(b_j)$ lies in $R_j(\auA_f)\cap K_j$ and is
$\equiv1$ at the chosen places: for $j$ odd this is clear from \eqref{ps:eq-lattice}; for $j$ even,
the last row of $b_j$ is $\equiv e_{\mathrm{last}}$ modulo $\pi_N$ and, $b_j^{-1}=b_j^*$ being in
$J_{X_j}$ too, its last column is $\equiv e_{\mathrm{last}}$ modulo $\aubar{\pi_N}$; so
$b_je_i\in\widehat\auO^{\,n_j-1}\oplus\pi_N\widehat\auO y\subset\psM_j$ and $(b_j-1)y/\aubar{\pi_N}
\in\widehat\auO^{\,n_j-1}\oplus\pi_N\widehat\auO y$, whence $\iota(b_j)\psM_j=\psM_j$. By
Lemma~\ref{ps:double}(b), $\bigotimes_j\omega(b_j,b_{j+1})\Phi(\mathsf y)=\Phi((\iota(b_j)\mathsf
y_j)_j)$; and $t\mapsto\iota(b_j)t$ permutes $\Omega_j$ (and $J_M/J_0$) and preserves $m$. So
$\Phi_\vartheta$ is invariant, and $\psGlue(\sum\Psi(x)\phi^\circ_x)=B^{2m}C_1^{-1}\Phi_\vartheta$
if $C_1\ne0$. If $c_0C_1\ne0$ put $c_{\mathrm{ss}}=c_0C_1B^{-2m}$: then $\iota_\vartheta(\auP
\varphi)=c_{\mathrm{ss}}B_{\vartheta,\Psi}\varphi$. If $c_0C_1=0$ then $\iota_\vartheta(\auP\varphi)
=0$ for all $\vartheta$ and $\varphi\in\auSram$, and the statement holds with $c_{\mathrm{ss}}=1$
and $\Psi=0$.
\end{proof}



\providecommand{\auO}{\mathcal O}
\providecommand{\auU}{\mathrm U}
\providecommand{\auGL}{\mathrm{GL}}
\providecommand{\auHom}{\operatorname{Hom}}
\providecommand{\auA}{\mathbb A}
\providecommand{\auF}{\mathbb F}
\providecommand{\auPl}{\Phi}
\providecommand{\auPa}{\Phi_a}
\providecommand{\auS}{\mathcal S}
\providecommand{\auSram}{\mathcal S^{\mathrm{ram}}}
\providecommand{\auAz}{\mathcal A_0}
\providecommand{\auNm}{\operatorname{Nm}}
\providecommand{\auAut}{\operatorname{Aut}}
\providecommand{\aubar}[1]{\overline{#1}}
\providecommand{\aurd}{\mathrm{rd}_E}
\providecommand{\psM}{\mathsf M}
\providecommand{\psHS}{\mathcal H}
\providecommand{\psW}{\mathcal W}
\providecommand{\psZ}{Z}
\providecommand{\psvac}{\mathbf v}
\providecommand{\psrho}{\varrho}
\providecommand{\psE}{\mathcal E}
\providecommand{\pao}{\mathfrak o}
\providecommand{\pam}{\mathfrak m}
\providecommand{\pkI}{\mathsf I}
\providecommand{\pkm}{\mathsf m}
\providecommand{\pkev}{\operatorname{ev}}
\providecommand{\pktr}{\operatorname{tr}}
\providecommand{\pkLW}{\mathrm{LW}}

\subsection{The upper bound: definite lifts, unfolding and kernels}\label{pk:sec}

This section follows \S\ref{ps:sec} and uses its notation. It proves the statements of \S\ref{ps:sec}
whose proofs are not given there, and Lemma~\ref{pa:kernel}, as follows.
\begin{itemize}
\item \emph{Proved here}, from Lemma~\ref{ps:weil} (with the conventions of
\S\ref{pk:ss-conv}) and the classical results named at the place of use: the definite lift
(Lemma~\ref{ps:definite}; Proposition~\ref{pk:definite}); the identity $\chi_0=1$ of
Remark~\ref{ps:chi0} (Lemma~\ref{pk:chi0}); Proposition~\ref{ps:onestep}, with complete proofs of
its steps (3), (5) and (6) (\S\ref{pk:ss-onestep}); Lemma~\ref{ps:branching}(c)
(Proposition~\ref{pk:branching}); Lemma~\ref{pa:kernel}(K0) (Proposition~\ref{pk:K0}).
\item \emph{Proved here from published results}: Lemma~\ref{pa:kernel}(K2), and (K1) at the
primes $\ell\nmid2q_1Md$ (all primes $\ell\ne d$ in \S\ref{pf:sec}) (Propositions~\ref{pk:std}, \ref{pk:K12}), from the properties (C1), (C2),
(C3)(i), (C4) of one curve with complex multiplication (Proposition~\ref{pk:curve}).
\item \emph{Proved in \S\ref{pf:sec}}: Proposition~\ref{pk:curve}(C3)(ii) (Theorem~\ref{pf:recip}), hence
Lemma~\ref{ps:conj}; the integrality in Lemma~\ref{pa:kernel}(K1) at the primes $\ell\mid2q_1M$ and at
$\ell=N$ (Corollary~\ref{pf:K1}); the fact (W6) (Proposition~\ref{pf:flip}).
\end{itemize}


\subsubsection{Conventions}\label{pk:ss-conv}

We fix the following conventions for \S\ref{ps:ss-weil} and this section; they are the ones
for which the statements of \S\ref{pa:ss-prod} on the linear forms $u^\dagger$ hold.

\emph{The additive character and $\jmath$.} $\psi=\psi_\Q\circ\pktr_{E/\Q}$ with
$\psi_{\Q,\infty}(t)=e^{-2\boldsymbol\pi it}$ and $\psi_{\Q,p}(t)=e^{2\boldsymbol\pi i\{t\}_p}$,
where $\boldsymbol\pi=3.14\ldots$; and $\jmath=-1/\sqrt{-d}$, so that $\tau(\jmath)=i/\sqrt d$
for every $\tau\in\auPl$ (the elements of $\auPl$ extend the given embedding of $K$). The idele
$\alpha$ is $1/P'(\zeta_{q_1})$ at the places of $E$ dividing
$\omega=(1+\sqrt{-d})/2$, and $1$ elsewhere, where $P$ is the minimal polynomial of $\zeta_{q_1}$
over $K$; every prime dividing $q_1$ divides $\auNm\omega=(1+d)/4$ and splits in $K$, so that
$\mathfrak a_0=E\cap\alpha\widehat\auO_E$ satisfies $\mathfrak a_0\aubar{\mathfrak a_0}=
\mathfrak d_{E/K}^{-1}$ and $\auO_E\subseteq\mathfrak a_0\subseteq\mathfrak d_{E/K}^{-1}$. The
character $\upsilon$ is unramified at \emph{every} finite place: $\upsilon(x)=\prod_\tau
x_\tau/|x_\tau|$ if the finite part of $x$ is a unit. Its existence, and the properties of it that
are used ($\upsilon_\tau(z)=z/|z|$; $\upsilon_{\aubar w}=\upsilon_w^{-1}$ on $E^{+\times}_v$ at a split
place; $\upsilon|\cdot|^{-1/2}$ algebraic), are Lemma~\ref{ps:upsilon}.

\emph{Fock coordinates.} For $\tau\in\auPl$ and $x,y\in W\otimes_{E^+,\tau}\R$,
\[
\psi_\tau\bigl(\tfrac12\jmath\langle y,x\rangle\bigr)=\exp\bigl(-2\boldsymbol\pi i\,\Rel(i\,d^{-1/2}
\langle y,x\rangle)\bigr)=\exp\bigl(i\,c_\psi^2\,\Iml\langle y,x\rangle\bigr),\qquad
c_\psi=\sqrt{2\boldsymbol\pi/\sqrt d}\,.
\]
Thus, in the coordinates $x_\tau=c_\psi\,\tau(w)$, $\beta_\infty(x,y)=\exp(+i\,\Iml\sum_\tau
\pktr(y_\tau^*x_\tau))$, the sign used in \S\ref{ps:ss-weil}. The Fock
space $\mathfrak F$ of $W=\auHom(X,X')$ consists of the entire functions $F$ of
$z=(z_\tau)_\tau\in\prod_\tau M_{n'\times n}(\C)$ with $\int|F|^2e^{-|z|^2}<\infty$; the variable
$z$ is the \emph{complex conjugate} of the coordinate $x$, and
\begin{equation}\label{pk:eq-fock}
(\psrho_\infty(y)F)(z)=e^{z\cdot y-|y|^2/2}F(z-\aubar y),\qquad z\cdot y=\sum_{\tau}\pktr(y_\tau^{\,t}
z_\tau),
\end{equation}
\[
(\omega_\infty(g,h)F)(z)=F(h^{t}z\,\aubar g)\qquad(g\in\auU(X)(\R),\ h\in\auU(X')(\R)),
\]
times the determinant twists (there is no further multiplier, $W$ being positive definite).
Indeed the exponents of $\psrho_\infty(x)\psrho_\infty(y)$ and of $\psrho_\infty(x+y)$ differ by
$-\aubar x\cdot y+\Rel(\aubar x\cdot y)=i\,\Iml(y^*x)$, and for a unitary matrix $u$ acting on $W
\otimes\R$ one has $(\aubar u^{-1}z)\cdot y=z\cdot(uy)$, so that $F\mapsto F(\aubar u^{-1}z)$
conjugates $\psrho_\infty(y)$ into $\psrho_\infty(uy)$; for $u\colon w\mapsto hwg^{-1}$,
$\aubar u^{-1}z=h^tz\aubar g$. With $E_z(w)=e^{w\cdot z}=e^{|z|^2/2}(\psrho_\infty(z)1)(w)$ the
theta function of $f\in\psHS_f(W)$ is
\begin{equation}\label{pk:eq-theta}
\Theta_f(z)=\theta(E_z\otimes f)=e^{|z|^2/2}\,\theta\bigl(\psrho_\infty(z)(1\otimes f)\bigr),
\end{equation}
which agrees with \eqref{ps:eq-theta}, and \eqref{ps:eq-bargmann} reads
\begin{equation}\label{pk:eq-bargmann}
Q(\partial)[\Theta_f(c\,\cdot)](0)=\theta\bigl(Q(c\,\cdot)\otimes f\bigr)=\int Q(c\,\bar y)\,
e^{-|y|^2/2}\,\theta\bigl(\psrho_\infty(y)(1\otimes f)\bigr)\,d\mu(y).
\end{equation}
The first equality follows from
\eqref{pk:eq-theta} and (W0) below, because $E_z=\sum_\alpha z^\alpha w^\alpha/\alpha!$ in
$\mathfrak F$.

\emph{The action of $\Gamma$ on the weights.} With these coordinates the rational points act on
$W_\lambda$ through the conjugate embeddings: $\gamma\cdot\xi=\rho_\lambda(\aubar{\tau(\gamma)})\xi$,
and for a real point $g_\infty\cdot\xi=\rho_\lambda(\aubar{g_\infty})\xi$ (see
\eqref{pk:eq-schuract}). The sentence ``$\Gamma$ acts on $W_\lambda$ through the embeddings $\tau$''
of {\cite[\S14.3]{P46}}, and the gauge there, refer to the embeddings $\aubar\tau=\tau\circ c$; this
is harmless, the set of data being stable under $\auPl\leftrightarrow\aubar\auPl$.

\emph{Three additions to Lemma~\ref{ps:weil}.}
\begin{itemize}
\item[(W0)] $\theta$ is defined on $\mathfrak F^{\mathrm{sm}}\otimes\psHS_f(W)$, where
$\mathfrak F^{\mathrm{sm}}$ is the space of smooth vectors, and is continuous in the first factor
for its Fr\'echet topology. In particular, for a compact subgroup $C$ of $\auU(W)(\R)$ and a
continuous function $a$ on $C$, $\theta\bigl(\int_Ca(c)\omega(c)F\,dc\bigr)=\int_Ca(c)\theta(\omega(c)F)\,dc$.
\item[(W2$'$)] In (W2) the character is $\chi_I(p)=\bigl(\upsilon|\cdot|^{-1/2}\bigr)(\det p|_I)^{-1}=
\upsilon(\det p|_I)^{-1}|\det p|_I|^{1/2}$, where $\det p|_I\in\auA_E^\times$ is the determinant of
$p$ on the $\auA_E$-module $I(\auA)$ and $|\cdot|$ is the idele norm of $E$. (This is Proposition~\ref{pwf:W2}(c).)
\item[(W6)] Let $V,V'$ be hermitian spaces of ranks $n,n'$ and $V^-$ the space $V$ with the opposite
form. The map $w\mapsto w^\dagger$ is an $E$-antilinear isomorphism $\auHom(V,V')\to\auHom(V',V^-)$
with $\langle w_1^\dagger,w_2^\dagger\rangle=-\langle w_2,w_1\rangle$; hence it preserves $\beta$,
and there is an isomorphism $\mathrm{Fl}\colon\psHS(\auHom(V,V'))\to\psHS(\auHom(V',V^-))$,
unique up to a scalar, with $\mathrm{Fl}\,\psrho(w)=\psrho(w^\dagger)\mathrm{Fl}$. It satisfies
\[
\mathrm{Fl}\circ\omega_{\auHom(V,V')}(g,g')=\upsilon(\det g)^{n'}\upsilon(\det g')^{-n}\cdot
\omega_{\auHom(V',V^-)}(g',g)\circ\mathrm{Fl}\qquad(g\in\auU(V)(\auA),\ g'\in\auU(V')(\auA)),
\]
and $\theta\circ\mathrm{Fl}$ is a multiple of $\theta$. The map $\mathrm{Fl}$ carries the vacuum of
$\alpha\auHom(\psM,\psM')$ to a multiple of the vacuum of its image $\aubar\alpha\auHom(\psM',\psM)$, which
is $\omega(1,u)$ applied to the vacuum of $\alpha\auHom(\psM',\psM)$, up to a scalar, with the norm-one
central idele $u=\aubar a/a$ of $\auU(V)$, trivial at the places not above $q_1$
(Remark~\ref{pf:rem-flip-lattice}).
\end{itemize}
(W0) is part of the construction of $\theta$ in \cite[no.~39, p.~189, and Th\'eor\`eme~6, p.~193]{WeilActa}. For (W6): the first assertions
are immediate ($(w^\dagger)^\dagger=-w$ for the pair $(V',V^-)$, and $\psi(\aubar a)=\psi(a)$); the
two sides of the displayed identity are operators with the same commutation relations (W1), so
they differ by a character of $\auU(V)(\auA)\times\auU(V')(\auA)$, and the assertion is that this
character is $\upsilon(\det g)^{n'}\upsilon(\det g')^{-n}$. 
(W6) is proved in Proposition~\ref{pf:flip}.

\emph{The lifts of the upper chain.} The measure on $\auU(X_j)(\auA_f)$ gives volume $1$ to the
level $J_{X_j}$, and for $\Phi\in\psHS_f(W_j)$
\begin{equation}\label{pk:eq-lift}
(\Theta^X_j(\Phi)\varphi)(g')=\int_{\Gamma\backslash\auU(X_j)(\auA_f)}\theta_j(\Phi;g,g')\varphi(g)\,dg;
\end{equation}
this is the class sum over $\Gamma\backslash\auU(X_j)(\auA_f)/J$ divided by $[J_{X_j}:J]$, for
every free level $J\subseteq J_{X_j}$ fixing $\varphi$ and $\omega_{\mu_{j+1}}(\cdot,1)\Phi$; it
makes precise the words ``for a free level $J$ fixing $\Phi$'' after \eqref{ps:eq-kernel}.

\emph{The constants.} $c_\psi=\sqrt{2\boldsymbol\pi/\sqrt d}$. Let $\Omega_0$ be a complex number
such that the lattice $\Lambda=\Omega_0\auO_K$ admits a point $t$ as in \S\ref{pa:ss-curve} with
$\wp''(t)=2$ (the \emph{period}) and $t$ admissible; let $\mathsf u$ be the number of
Proposition~\ref{pk:curve}(C4) for this lattice (the \emph{Chowla--Selberg unit}); let
$\varpi_{\mathrm{vac}}=(\upsilon(\alpha)|\alpha|^{-1/2})^{-1/f}$ (principal branch). Then
\begin{equation}\label{pk:eq-s0}
s_0=\frac{\varpi_{\mathrm{vac}}}{\mathsf u}\cdot\frac{c_\psi}{\Omega_0},\qquad|\sigma
\varpi_{\mathrm{vac}}|^2=\frac{\aurd}{\sqrt d}\ \ (\sigma\in\auAut(\C/K)),\qquad|s_0|^2=\frac1\Omega .
\end{equation}
(For the second formula: $|\alpha|=\auNm(\mathfrak d_{E/K})^{1/2}$, the discriminant of $E$ is
$d^f\auNm(\mathfrak d_{E/K})$, and the conjugates over $K$ of the value $\upsilon(\alpha)$ of the
algebraic Hecke character $\upsilon$ have absolute value $1$. The third follows from the second
and (C4), since the area of $\Lambda$ is $|\Omega_0|^2\sqrt d/2$.) The integral vacuum of rank one
is defined in \eqref{pk:eq-vac}; it is $c_{\mathrm{vac}}\psvac$ for a constant
$c_{\mathrm{vac}}\ne0$ which is thereby determined but for which we have no closed formula.

\subsubsection{The definite lift: proof of Lemma~\ref{ps:definite}}\label{pk:ss-definite}

Let $X,X'$ be consecutive spaces $X_j,X_{j+1}$, of ranks $n<n'$, $\lambda=\lambda(X)$,
$\lambda'=\lambda(X')$, $e=e_{j,\tau}$, $e'=e_{j+1,\tau}$, $c(n)=c_\tau(n_j)$. Let $\rho_X$ be the
representation of $\prod_\tau\auGL_n(\C)$ on $W_\lambda$ and $\rho_{X'}$ that on $W_{\lambda'}$. The
determinant twists of $\omega_{\mu_{j+1}}$ at the infinite places are $\prod_\tau\det(h_\tau)^{e'}
\det(g_\tau)^{-(e-c(n))}$: indeed $\mu_{j+1}$ has weight $-\kappa_{j+1}$, $\upsilon|_{\auU(1)}$ has
weight $1$, $n,n'$ are even, and $e'=n/2-\kappa_{j+1}$, $e-c(n)=n'/2-\kappa_{j+1}$ (for $j=0$:
$e=0$, $c_\tau(k)=\kappa_1-n_1/2$). All $n_j$ are even, because $k=r+1-m$ and $t+1=5r+1$ are. By
the multiplicativity of Schur functors, $S(AzB)=\rho^{\circ}_{X'}(A)S(z)\rho^\circ_X(B)$, where
$\rho^\circ_{X'}=\rho_{X'}\otimes\det^{-e'}$ and $\rho_X^\circ=\rho_X\otimes\det^{-(e-c(n))}$ (the
columns of height $n$ give $\det^{c(n)}$). Hence, for a covector $A$ on $W_{\lambda'}$ and
$\xi\in W_\lambda$, as $\det\aubar g=\det g^{-1}$ and $h^t=\aubar h^{-1}$,
\begin{equation}\label{pk:eq-schuract}
\omega_{\mu_{j+1},\infty}(g,h)\bigl[A\cdot S(z)\xi\bigr]=A\cdot\rho_{X'}(\aubar h)^{-1}S(z)\,\rho_X(\aubar
g)\xi .
\end{equation}

\begin{lemma}[Cauchy decomposition]\label{pk:cauchy}
Let $\mathcal P$ be the space of polynomials in $z$. As a module over $\prod_\tau\auU(n)$ acting by
$\omega_{\mu_{j+1},\infty}(g,1)$, $\mathcal P$ is the direct sum of the spaces $\mathcal P_\nu$
spanned by the polynomials $A\cdot S_\nu(z)\xi$, where $\nu=(\nu_\tau)$ runs over the families of
diagrams with at most $n$ rows and $S_\nu$ is the Schur functor of $\nu$; the modules
$\mathcal P_\nu$ are isotypic of pairwise different irreducible types, and $(A,\xi)\mapsto A\cdot
S_\nu(z)\xi$ is injective on $W^*_{\nu}(\C^{n'})\otimes W_\nu(\C^n)$.
\end{lemma}
\begin{proof}
This is the decomposition of the polynomials on $M_{n'\times n}(\C)$ under $\auGL_{n'}\times\auGL_n$
into the irreducible modules $W_\nu(\C^{n'})^*\otimes W_\nu(\C^n)^*$, each once, the summand of
$\nu$ being the span of the matrix coefficients of $S_\nu(z)$
(\cite{HoweTheta}, \cite[Theorem~5.6.7]{GoodmanWallach}; for the unitary groups see
\cite[Ch.~III, Proposition~5.6]{KashiwaraVergne}), applied at each $\tau$.
\end{proof}

\begin{lemma}[One lift]\label{pk:single}
For every $P\in\mathcal P$ and every linear form $l$ on $W_\lambda$ there is a linear form $l'$ on
$W_{\lambda'}$ such that for all twists, all $\Phi\in\psHS_f(\auHom(X,X'))$, all algebraic modular
forms $\psi$ on $\auU(X)$ of weight $\lambda$ and all $h\in\auU(X')(\auA)$,
\[
\int_{[\auU(X)]}l\bigl(g_\infty^{-1}\cdot\psi(g_f)\bigr)\,\theta\bigl(\omega_{\mu_{j+1}}(g,h)(P\otimes
\Phi)\bigr)\,dg=l'\Bigl(h_\infty^{-1}\cdot\bigl(\Theta^X_j(\Phi)\psi\bigr)(h_f)\Bigr).
\]
The form $l'$ depends linearly on $P$, and $l'=0$ if $P$ is orthogonal to $\mathcal P_{\nu^0}$,
where $\nu^0$ is the diagram of $\lambda$ with its $c(n)$ columns of height $n$.
\end{lemma}
\begin{proof}
Write $G_\infty=\auU(X)(\R)$, compact, and $[\auU(X)]=\Gamma\backslash(G_\infty\times G_f)$. Since
$[\auU(X)]$ has volume $1$ and $J_X$ is free, the measure is the product of the measure of $G_f$
with volume $1$ on $J_X$ and of the Haar measure of $G_\infty$ of total mass
$m_\infty=\#(\Gamma\backslash G_f/J_X)^{-1}$. For fixed $g_f$ put $v=\psi(g_f)$ and
$\Phi'=\omega_{\mu_{j+1},f}(g_f,h_f)\Phi$. By (W0) the integral over $G_\infty$ is
$\theta(P_v\otimes\Phi')$ with
\[
P_v=\int_{G_\infty}l\bigl(\rho_X(\aubar g)^{-1}v\bigr)\,\omega_\infty(g,1)P'\,dg\in\mathcal P,\qquad
P'=\omega_{\mu_{j+1},\infty}(1,h_\infty)P .
\]
The map $v\mapsto P_v$ is linear and $P_{\rho_X(\aubar g_0)v}=\omega_\infty(g_0,1)P_v$. The module
$(W_\lambda,\ g\mapsto\rho_X(\aubar g))$ is irreducible, of the type of $\mathcal P_{\nu^0}$ by
\eqref{pk:eq-schuract}. By Lemma~\ref{pk:cauchy} and Schur's lemma there is a unique covector
$A=A(l,P')$ on $W_{\lambda'}$ with $P_v=A\cdot S(z)v$ for all $v$; it is linear in $P'$ and
vanishes if $P'\perp\mathcal P_{\nu^0}$. As $\omega_\infty(1,h_\infty)$ commutes with
$\omega_\infty(g,1)$, \eqref{pk:eq-schuract} gives $A(l,P')=A(l,P)\circ\rho_{X'}(\aubar h_\infty)^{-1}$.
The entries of $S$ are homogeneous of degree $L=\sum_{\tau,p}l_{\tau,p}$, so by
\eqref{pk:eq-bargmann} and \eqref{ps:eq-kernel}
\[
\theta(P_v\otimes\Phi')=A(l,P')\bigl(\theta(S(z)\otimes\Phi')v\bigr)=(\pi_Ms_0)^{-L}A(l,P)\Bigl(
h_\infty^{-1}\cdot\theta_j(\Phi;g_f,h_f)\psi(g_f)\Bigr).
\]
Integrating over $\Gamma\backslash G_f$ and using \eqref{pk:eq-lift} gives the lemma with
$l'=m_\infty(\pi_Ms_0)^{-L}A(l,P)$.
\end{proof}

\begin{proposition}[Lemma~\ref{ps:definite}]\label{pk:definite}
Lemma~\ref{ps:definite} holds, for every $u\in\mathfrak F^{\mathrm{sm}}(\auHom(X',X^-\oplus X''))$,
with $g_\infty^{-1}\psi(g_f)$ and $k_\infty^{-1}(\cdots)$ read as $g_\infty^{-1}\cdot\psi(g_f)$,
$k_\infty^{-1}\cdot(\cdots)$, and the tensor product of the two Weil representations transported
by (W6).
\end{proposition}
\begin{proof}
Denote the left side by $\mathcal L(u,l)$; it is linear in $u$ and $l$ and, by (W0) and the
moderate growth of theta functions, continuous in $u$. \emph{Reduction to polynomials.} Replacing
$u$ by $\omega_\infty(g_0,r_0,1)u$ with $g_0\in\auU(X)(\R)$, $r_0\in\auU(X')(\R)$ and substituting
$g\mapsto gg_0^{-1}$, $r\mapsto rr_0^{-1}$, we get $\mathcal L(\omega_\infty(g_0,r_0,1)u,l)=
\mathcal L(u,l\circ g_0)$. Apply this to the scalar matrices $g_0=(e^{i\vartheta_\tau})$,
$r_0=(e^{i\vartheta'_\tau})$, which act on $W_\lambda$ by characters: by (W6), $\mathfrak F
(\auHom(X',X^-\oplus X''))$ is the completed tensor product of the Fock spaces of $\auHom(X,X')$
and $\auHom(X',X'')$, with variables $z_1,z_2$, and by \eqref{pk:eq-fock} these scalars multiply
the part of $u$ of bidegree $(d_1(\tau),d_2(\tau))_\tau$ in $(z_1,z_2)$ by a character of
$(\vartheta,\vartheta')$ from which $d_1$ and $d_1-d_2$ can be read. Hence $\mathcal L(u,l)=
\mathcal L(u^\flat,l)$, where $u^\flat$ is the component of $u$ of one bidegree, a polynomial.
\emph{Polynomials.} By linearity let $u=P_1\otimes P_2$ and $\Phi=\Phi_1\otimes\Phi_2$. Then
$\theta(\omega(g,r,k)(u\otimes\Phi))$ is a constant multiple of $\theta(\omega_{\mu_{j+1}}(g,r)
(P_1\otimes\Phi_1))\cdot\theta(\omega_{\mu_{j+2}}(r,k)(P_2\otimes\Phi_2))$ by (W6). Integrate over
$g$ by Lemma~\ref{pk:single}, which gives $l'(r_\infty^{-1}\cdot(\Theta^X_j(\Phi_1)\psi)(r_f))$,
and then over $r$ by the same lemma for the pair $(X',X'')$, the form $l'$ and the form
$\Theta^X_j(\Phi_1)\psi$. The resulting form $l''$ depends only on $(l,P_1,P_2)$.
\end{proof}

\subsubsection{The character $\chi_0$ is trivial}\label{pk:ss-chi0}

\begin{lemma}\label{pk:chi0}
In the proof of Lemma~\ref{ps:double}, $\chi_0=1$: with the determinant twists of
\S\ref{ps:ss-weil} on the spaces $Y$ and on the spaces $X$, and the identifications {\rm(W6)},
\[
\psZ_{\mathbb J}\circ\bigl(\omega_{\mu_{j+1}}(\iota g,\iota r')\otimes\omega_{\mu_{j+2}}(\iota r',
\iota k)\bigr)=\bigl(\omega_{\mu_{j+1}}(g,r')\otimes\omega_{\mu_{j+2}}(r',k)\bigr)\circ\psZ_{\mathbb J}
\]
for $g\in\auU(X)(\auA)$, $r'\in\auU(X')(\auA)$, $k\in\auU(X'')(\auA)$ (resp.\ $\auU(Y_m)(\auA)$).
Hence {\cite[Proposition~14.19]{P46}} holds with $\chi_0=1$, the upper chain being defined by $\omega_\mu$.
\end{lemma}
\begin{proof}
Let $n,n',n''$ be the ranks of $X,X',X''$ and $N=n+s$, $N'=n'+s-1$, $N''$ the ranks of
$V,V',V''$; so $N''=n''+s-2$ if $s\ge3$ and $N''=n''-1$ if $s=1$. Here $s$ is odd and
$n,n',n''$ are even; thus $N,N''$ are odd and $N'$ is even. Write $\upsilon(g)$ for
$\upsilon(\det g)$. By the definition of $\omega_\mu$ and (W6), the left operator, before
$\psZ_{\mathbb J}$, is
\[
\frac{\mu_{j+1}(r')\mu_{j+2}(k)}{\mu_{j+1}(g)\mu_{j+2}(r')}\cdot\upsilon(r')^{\frac{N-1}2}\upsilon(g)^{-
\frac{N'}2}\cdot\upsilon(k)^{\frac{N'}2}\upsilon(r')^{-\frac{N''-1}2}\cdot\upsilon(g)^{N'}\upsilon(r')^{-N}
\cdot\omega_{\auHom(V',\mathbb V)}(r',g\oplus k),
\]
and the right operator is the same expression with $(n,n',n'')$ in place of $(N,N',N'')$, the
floors being $n/2,n'/2,n''/2$, and $\auHom(X',X^-\oplus X'')$ in place of $\auHom(V',\mathbb V)$.
By (W2$'$), $\psZ_{\mathbb J}\,\omega_{\auHom(V',\mathbb V)}(p)=\upsilon(\det p|_{\mathbb J})^{-1}
\omega_{\auHom(X',X^-\oplus X'')}(p)\,\psZ_{\mathbb J}$ for $p=(r',g\oplus k)$; the absolute value
is $1$ because $\det g$, $\det r'$, $\det k$ are ideles of norm one to $E^+$.

\emph{The determinant on $\mathbb J$.} For $s=1$ all grades are $0$, $\mathbb J=0$ and the
determinant is $1$. Let $s\ge3$. The element $p$ acts on $w\in\auHom(V',\mathbb V)$ by
$w\mapsto(g\oplus k)wr'^{-1}$; it acts by $r'$ on $X'$, by $g\oplus k$ on $X\oplus X''$ and
trivially on all vectors of the strings. On a block $\auHom(A,B)$ the determinant is
$\det(p|_B)^{\dim A}\det(p|_A)^{-\dim B}$. The string of $V'$ has $s-1$ vectors, of odd grades,
$(s-1)/2$ of them negative; the strings of $V$ and $V''$ have $(s-1)/2$ and $(s-3)/2$ vectors of
positive grade. The blocks of $\mathbb J=\auHom_{>0}(V',\mathbb V)\oplus\auHom(X',\ell)$ on which
$p$ is not the identity are therefore: $\auHom(V'_a,X\oplus X'')$ for the $(s-1)/2$ negative
grades $a$; $\auHom(X',\mathbb V_b)$ for the $(s-1)/2+(s-3)/2=s-2$ string vectors of positive
grade $b$; and $\auHom(X',\ell)$. So
\[
\det p|_{\mathbb J}=(\det g\det k)^{\frac{s-1}2}(\det r')^{-(s-1)}\ \ (s\ge3),\qquad\det p|_{\mathbb
J}=1\ \ (s=1).
\]
\emph{Comparison of exponents} (the characters $\mu$ agree on both sides). For $s\ge3$: the
exponent of $\upsilon(g)$ on the left is $-\frac{N'}2+N'-\frac{s-1}2=\frac{n'}2$, and on the right
$-\frac{n'}2+n'=\frac{n'}2$. The exponent of $\upsilon(k)$ is $\frac{N'}2-\frac{s-1}2=\frac{n'}2$
on the left and $\frac{n'}2$ on the right. The exponent of $\upsilon(r')$ on the left is
$\frac{N-1}2-\frac{N''-1}2-N+(s-1)=\frac{n-n''}2+1-n-1=-\frac{n+n''}2$, and on the right
$\frac n2-\frac{n''}2-n=-\frac{n+n''}2$. For $s=1$: $N=n+1$, $N'=n'$, $N''=n''-1$, and the exponents
on the left are $\frac{n'}2$ for $\upsilon(g)$ and $\upsilon(k)$, and $\frac n2-(\frac{n''}2-1)-
(n+1)=-\frac{n+n''}2$ for $\upsilon(r')$; they agree with those on the right ($\det\operatorname{diag}
(k,1)=\det k$).
\end{proof}

\subsubsection{One step of the unfolding: proof of Proposition~\ref{ps:onestep}}\label{pk:ss-onestep}

We give a complete proof of Proposition~\ref{ps:onestep}. Steps (1), (2) and (4) of the proof in
\S\ref{ps:ss-seesaw} are kept; steps (3), (5), (6) are replaced by
Lemmas~\ref{pk:vanish}--\ref{pk:fine} and the argument at the end. The notation is that of
Proposition~\ref{ps:onestep}: $0\le j<m-1$, $V=Y_j$, $V'=Y_{j+1}$, $s=s_j\ge1$,
$W=\auHom(V,V')$, $G=G_j$, $G'=G_{j+1}$, $R=R_j$, $R'=R_{j+1}$, $N=N_j$, $T=T_j$,
$\mathfrak e=\mathfrak e_j=T^\dagger T$, $\omega=\omega_{\mu_{j+1}}$. We write $V_g$ for the part of
grade $g$ of $V$, $g_a=s-1-2a$ for the grade of $f_a$, $g'_b=s-2-2b$ for that of $f'_b$, and
$V^+,V^0,V^-$ for the parts of positive, zero, negative grade; $p=\dim V^+=\lfloor s/2\rfloor$.
Recall $Tf_a=t_af'_a$ with $t_a\ne0$ ($a\le s-2$), $\mathfrak ef_a=\mathsf e_af_{a+1}$ with
$\mathsf e_a\ne0$ ($a\le s-2$), and $T(X_j)=0$.

\emph{Three remarks on grades.} (G1) The non-zero grades of $V$ are $\equiv s-1$ and those of $V'$
are $\equiv s\pmod2$. Hence a map $V\to V'$ of degree $\ge0$ has degree $\ge1$ on $V^+\oplus V^-$,
and $\auHom_{\ge0}(V^+,V')=\auHom_{>0}(V^+,V')=:J$; the maps in $J$ have image in
$V'^{+}\subseteq T(V^+)$. (G2) $T(V^+)$ is the span of $f'_0,\dots,f'_{p-1}$; the form
$Q_T(v,v')=\langle Tv,Tv'\rangle$ on $V^+$ is $0$ if $s$ is odd, and if $s$ is even its radical is
$\langle f_0,\dots,f_{p-2}\rangle$ and $Q_T(f_{p-1},f_{p-1})\ne0$ (because $\langle f'_a,f'_b\rangle
\ne0$ only for $a+b=s-2$, and $a,b\le p-1$). (G3) For $y\in V$, $y'\in V'$ let $y'\otimes y^\flat$
be the map $x\mapsto\langle y,x\rangle y'$; then $\pktr(w^\dagger(y'\otimes y^\flat))=\langle wy,y'
\rangle$ for $w\in W$, and $f_a^\flat$ vanishes outside $V_{-g_a}$.

\emph{Constant terms along an isotropic space.} Let $\pkI\subseteq W$ be one of the two isotropic
spaces
\[
I_0=\{w:\ w|_{V^+\oplus V^0}=0\},\qquad I=\auHom_{>0}(V,V'),
\]
so that $I_0^\perp=\{w:w|_{V^+}=0\}$ and $I^\perp=\auHom_{\ge0}(V,V')$. Both are stable under
$w\mapsto w\gamma$, $\gamma\in R$, and $I_0$ is stable under $w\mapsto\delta w$, $\delta\in G'$.
Moreover
\begin{equation}\label{pk:eq-H}
\pkI^\perp\cdot(n-1)\subseteq\pkI\qquad(n\in N),
\end{equation}
because $n-1$ has degree $\ge1$ and maps $V^+\oplus V^0$ into $V^+$. For $w\in W(\auA)$ and
$F\in\psHS(W)$ put $c^{\pkI}_w(F)=\theta_{W_{\pkI}}(\psZ_{\pkI}\psrho(w)F)$; thus $c^I=c$ and
$c^{I_0}=c^0$ in the notation of \S\ref{ps:ss-seesaw}. By (W1), (W2):
\begin{itemize}
\item[(c1)] $c^{\pkI}_{w+i}=\beta(i,w)^{-1}c^{\pkI}_w$ for $i\in\pkI(\auA)$, and $c^{\pkI}_{w+\delta}=
c^{\pkI}_w$ for $w\in W(E)$, $\delta\in\pkI^\perp(E)$;
\item[(c2)] $c^{\pkI}_w(\omega(\gamma,1)F)=c^{\pkI}_{w\gamma}(F)$ for $\gamma\in R(E^+)$, and
$c^{\pkI}_w(\omega(n,1)F)=c^{\pkI}_{wn}(F)$ for $n\in N(\auA)$; for $\pkI=I_0$ also
$c^{0}_w(\omega(1,\delta)F)=c^{0}_{\delta^{-1}w}(F)$ for $\delta\in G'(E^+)$.
\end{itemize}
(For (c2): $\psrho(w)\omega(g,h)=\omega(g,h)\psrho(h^{-1}wg)$; $\psZ_{\pkI}\omega(p)=\chi_{\pkI}(p)
\omega_{W_{\pkI}}(\bar p)\psZ_{\pkI}$; for rational $p$, $\chi_{\pkI}(p)=1$ and
$\theta_{W_{\pkI}}\circ\omega_{W_{\pkI}}(\bar p)=\theta_{W_{\pkI}}$; for $p=(n,1)$, $\chi_{\pkI}(p)=1$
and $\bar p=1$ by \eqref{pk:eq-H}; and the determinant twists are trivial on rational and on
unipotent elements.)

For $w\in W(E)$ let
\[
A^{\pkI}_w=\{n\in N^{\ge2}:\ w(n-1)\in\pkI\},\qquad\chi^{\pkI}_w(n)=\beta\bigl(w(n-1),w\bigr)^{-1}=
\psi\bigl(\tfrac12\jmath\,\pktr(w^\dagger w(n-1))\bigr)^{-1}.
\]

\begin{lemma}\label{pk:char}
{\rm(a)} $A^{\pkI}_w$ is a unipotent algebraic subgroup of $N^{\ge2}$ over $E^+$, and
$\chi^{\pkI}_w$ is a character of $A^{\pkI}_w(\auA)$, trivial on $A^{\pkI}_w(E^+)$; both depend
only on the class of $w$ modulo $\pkI^\perp(E)$. {\rm(b)} $c^{\pkI}_w(\omega(n,1)F)=
\chi^{\pkI}_w(n)\,c^{\pkI}_w(F)$ for $n\in A^{\pkI}_w(\auA)$. {\rm(c)} For $\gamma\in R(E^+)$,
$A^{\pkI}_{w\gamma}=\gamma^{-1}A^{\pkI}_w\gamma$ and $\chi^{\pkI}_{w\gamma}(\gamma^{-1}n\gamma)=
\chi^{\pkI}_w(n)$; for $n\in A^{\pkI}_w(E^+)$, $wn\equiv w$ modulo $\pkI^\perp(E)$.
\end{lemma}
\begin{proof}
Write $\nu=n-1$, $\nu'=n'-1$; $\nu^\dagger=n^{-1}-1$. (a) $w(nn'-1)=w\nu+w\nu'+(w\nu)\nu'$ and
$(w\nu)\nu'\in\pkI\nu'\subseteq\pkI$ by \eqref{pk:eq-H}; $w(n^{-1}-1)=-w\nu-(w\nu)\nu^\dagger\in
\pkI$. So $A^{\pkI}_w$ is a group; it is defined by linear conditions on $\nu$. Since
$\langle w,u\nu'\rangle=\pktr(\nu'w^\dagger u)=\langle w\nu'^\dagger,u\rangle$, we get $\langle w,
(w\nu)\nu'\rangle=\langle w(n'^{-1}-1),w\nu\rangle=0$, both entries being in $\pkI$; so
$\chi^{\pkI}_w$ is multiplicative. It is trivial on rational points because $\psi$ is trivial on
$E$. If $\delta\in\pkI^\perp(E)$ then $\delta\nu\in\pkI$ by \eqref{pk:eq-H}, so
$A^{\pkI}_{w+\delta}=A^{\pkI}_w$, and $\langle w+\delta,(w+\delta)\nu\rangle-\langle w,w\nu\rangle=
\langle\delta,w\nu\rangle+\langle w\nu^\dagger,\delta\rangle+\langle\delta,\delta\nu\rangle=0$, each
term pairing $\pkI$ with $\pkI^\perp$. (b) By (c2) and (c1), $c^{\pkI}_w(\omega(n,1)F)=
c^{\pkI}_{w+w\nu}(F)=\beta(w\nu,w)^{-1}c^{\pkI}_w(F)$. (c) $w\gamma(\gamma^{-1}n\gamma-1)=(w\nu)
\gamma$ and $\pkI\gamma=\pkI$; $\langle w\gamma,(w\nu)\gamma\rangle=\langle w,w\nu\rangle$ as
$\gamma$ is unitary; and $wn-w=w\nu\in\pkI(E)$.
\end{proof}

A class $w\in W(E)/\pkI^\perp(E)$ is \emph{dead} (for $\pkI$) if the character
$\psi_j\chi^{\pkI}_w$ of $A^{\pkI}_w(\auA)$ is not trivial. By Lemma~\ref{pk:char}(c) and the
invariance of $\psi_j$ under conjugation by $R(\auA)$, the set of dead classes is stable under
$R(E^+)$.

\begin{lemma}[Vanishing]\label{pk:vanish}
Let $\mathcal O$ be a set of dead classes for $\pkI$, stable under $R(E^+)$, $\psW$ an equivariant
function on $[R]$ and $F\in\psHS(W)$. Then $\sum_{w\in\mathcal O}|c^{\pkI}_w(\omega(r,1)F)|$
converges uniformly for $r$ in compact sets and
\[
\int_{[R]}\psW(r)\sum_{w\in\mathcal O}c^{\pkI}_w(\omega(r,1)F)\,dr=0 .
\]
\end{lemma}
\begin{proof}
The convergence is (W3). Since $[R]$ is compact it suffices to treat one orbit $\mathcal O=w_0
R(E^+)$. Let $S\subseteq R(E^+)$ be the stabiliser of the class $w_0$, $A=A^{\pkI}_{w_0}$,
$\chi=\chi^{\pkI}_{w_0}$. By (c2), $\sum_{w\in\mathcal O}c^{\pkI}_w(\omega(r,1)F)=\sum_{\gamma\in
S\backslash R(E^+)}\phi(\gamma r)$ with $\phi(r)=c^{\pkI}_{w_0}(\omega(r,1)F)$, and $\psW$ is left
invariant under $R(E^+)$; so the integral is $\int_{S\backslash R(\auA)}\psW\phi$, absolutely
convergent. By Lemma~\ref{pk:char}(c), $A(E^+)\subseteq S$, $S\cap A(\auA)=A(E^+)$, and $S$
normalises $A$. The group $S\,A(\auA)$ is closed in $R(\auA)$ (it is $S\cdot\mathsf K$ with $S$
discrete and $\mathsf K\subset A(\auA)$ compact, $[A]$ being compact), and it is unimodular, like
$R(\auA)$, since both contain discrete subgroups with compact quotient. Hence
\[
\int_{S\backslash R(\auA)}\psW\phi=\int_{SA(\auA)\backslash R(\auA)}\int_{[A]}\psW(ar)\phi(ar)\,da\,
d\dot r,
\]
and $\psW(ar)\phi(ar)=\psi_j(a)\chi(a)\psW(r)\phi(r)$ by the equivariance of $\psW$ and
Lemma~\ref{pk:char}(b). The inner integral is $0$ because $\psi_j\chi$ is a non-trivial character
of the compact group $[A]$.
\end{proof}

\begin{lemma}[Coarse classes]\label{pk:coarse}
Let $w\in\auHom(V^+,V')(E)=W(E)/I_0^\perp(E)$. If $w\notin G'(E^+)\,T|_{V^+}$, then $w$ is dead
for $I_0$.
\end{lemma}
\begin{proof}
Let $\tilde w$ be the extension of $w$ by $0$ on $V^0\oplus V^-$, and $Q_w(v,v')=\langle wv,wv'
\rangle$ on $V^+$. For $n\in N^{\ge2}$, $n\in A^{I_0}_w$ means $w(n-1)|_{V^+\oplus V^0}=0$.

(a) \emph{$Q_w=Q_T$ unless $w$ is dead.} Let $\nu\colon V\to V$ be zero on $V^+\oplus V^0$, map
$V^-$ into $V^+$, and satisfy $\nu^\dagger=-\nu$; then $n=1+\nu\in N^{\ge2}$ and $\tilde w\nu$
vanishes on $V^+\oplus V^0$, so $n\in A^{I_0}_w$. In the basis $f_{s-1-a}$ ($a<p$) of $V^-$, dual
to the basis $f_a$ of $V^+$, write $\nu f_{s-1-a}=\sum_{b<p}\nu_{ba}f_b$; the condition
$\nu^\dagger=-\nu$ is $\nu_{ab}=-\aubar{\nu_{ba}}$, that is, $(\nu_{ba})=\jmath\mathsf H$ with
$\mathsf H$ hermitian. Now $\pktr(\tilde w^\dagger\tilde w\nu)=\sum_{a<p}\langle f_a,\tilde
w^\dagger\tilde w\nu f_{s-1-a}\rangle=\sum_{a,b}\nu_{ba}Q_w(f_a,f_b)$, and the same with $T$. Hence
$\psi_j\chi^{I_0}_w(n)=\psi\bigl(\tfrac12\jmath^2\pktr(\mathsf H\mathsf D)\bigr)$ with
$\mathsf D=(Q_T-Q_w)(f_a,f_b)$, a hermitian matrix over $E$. If $\mathsf D\ne0$, $\mathsf
H\mapsto\jmath^2\pktr(\mathsf H\mathsf D)$ is a surjective $\auA_{E^+}$-linear map from the
hermitian adelic matrices to $\auA_{E^+}$, and $x\mapsto\psi(\tfrac12x)$ is a non-trivial character
of $\auA_{E^+}$; so $w$ is dead.

(b) \emph{$w$ is injective unless it is dead.} Let $Q_w=Q_T$ and $\mathsf K=\ker w\ne0$; then
$\mathsf K$ lies in the radical of $Q_T$. \emph{First case:} $\mathsf K$ contains a non-zero
vector of $\langle f_0,\dots,f_{p-2}\rangle$. Choose such a vector $\kappa=\sum_{b<i}\kappa_bf_b$
with $\kappa_{i-1}\ne0$, $1\le i\le p-1$. For $\mathsf t\in\auA_E$ let $\nu_+$ be the map
$f_i\mapsto\mathsf t\kappa$, zero on the other basis vectors of $V$, and $n=1+\nu_+-\nu_+^\dagger$.
As $\nu_+^2=0=\nu_+\nu_+^\dagger=\nu_+^\dagger\nu_+$, $n$ is unitary; $\nu_+$ raises grades by
$2(i-b)\ge2$, so $n\in N^{\ge2}$; $\nu_+^\dagger$ vanishes on $V^+\oplus V^0$ and $\tilde w\nu_+=0$,
so $\tilde w(n-1)=0$: thus $n\in A^{I_0}_w$ and $\chi^{I_0}_w(n)=1$. Further
$\pktr(\mathfrak e\nu_+)=\mathsf t\,\mathsf e_{i-1}\kappa_{i-1}=:\varsigma$ and
$\pktr(\mathfrak e\nu_+^\dagger)=\aubar\varsigma$ ($\mathfrak e$ is self-adjoint), so
$\psi_j(n)=\psi(\tfrac12\jmath(\varsigma-\aubar\varsigma))=\psi(\jmath\varsigma)$, which is not $1$
for a suitable $\mathsf t$. \emph{Second case:} $\mathsf K\cap\langle f_0,\dots,f_{p-2}\rangle=0$.
By (G2) $s$ is odd, $V^0=X_j\oplus Ez$ with $z=f_p$, and $\mathsf K=E\kappa$ with
$\kappa_{p-1}\ne0$. Let $b\colon V^0\to V^+$ be $z\mapsto\mathsf t\kappa$, $X_j\mapsto0$, extended
by $0$, and $n=\exp(b-b^\dagger)=1+b-b^\dagger-\tfrac12bb^\dagger$. All grades being even,
$n\in N^{\ge2}$; $\tilde wb=0$ and $\tilde wb^\dagger=0$ ($b^\dagger$ maps $V^-$ to $V^0$), so
$\tilde w(n-1)=0$, $n\in A^{I_0}_w$, $\chi^{I_0}_w(n)=1$. And $\pktr(\mathfrak eb)=\mathsf t\,
\mathsf e_{p-1}\kappa_{p-1}=:\varsigma$, $\pktr(\mathfrak eb^\dagger)=\aubar\varsigma$,
$\pktr(\mathfrak ebb^\dagger)=0$ ($\mathfrak ebb^\dagger$ maps $V^-$ into $V^+\oplus V^0$ and kills
the rest). So $\psi_j(n)=\psi(\jmath\varsigma)\ne1$ for a suitable $\mathsf t$.

(c) If $w$ is not dead, it is injective with $Q_w=Q_T$; then $Tv\mapsto wv$ is an isometry between
subspaces of $V'$, which by Witt's theorem is induced by an element of $G'(E^+)$.
\end{proof}

\begin{lemma}[Fine classes]\label{pk:fine}
Let $w=T+v$ with $v\in W(E)$, $v|_{V^+}=0$, $v\notin\auHom_{\ge0}(V,V')$. Then the class of $w$ is
dead for $I$.
\end{lemma}
\begin{proof}
Write $v$ in the bases: for a basis vector $\kappa$ of $V^0\oplus V^-$ (a vector of the basis of
$X_j$, or $f_a$ with $a\ge p$), of grade $g(\kappa)\le0$, $v\kappa=\sum_bv_{b,\kappa}f'_b+
(\text{a vector of }X_{j+1})$. An entry $v_{b,\kappa}$ has negative degree if $g'_b<g(\kappa)$;
by hypothesis some entry of negative degree is not zero (an entry with values in $X_{j+1}$ has
degree $-g(\kappa)\ge0$). Among the non-zero entries of negative degree choose one,
$v_{c,\kappa_0}$, with $g'_c$ least, and put $g'=g'_c<g(\kappa_0)\le0$ and $a=s-2-c$. Then
$g_a=1-g'\ge2$, so $f_a\in V^+$; and $Tf_a=t_af'_a$ pairs to zero with all basis vectors of $V'$
except $f'_c$.

For $\mathsf t\in\auA_E$ let $\kappa=\mathsf t\kappa_0$ and
\[
n=1+f_a\otimes\kappa^\flat-\kappa\otimes f_a^\flat-\tfrac12\langle\kappa,\kappa\rangle f_a\otimes
f_a^\flat .
\]
Since $\langle f_a,f_a\rangle=0=\langle f_a,\kappa_0\rangle$ (the only basis vector not orthogonal
to $f_a$ is $f_{s-1-a}$, of grade $g'-1<g(\kappa_0)$), a direct computation gives $\langle nx,ny
\rangle=\langle x,y\rangle$. By (G3) the three terms of $n-1$ map $V_{-g(\kappa_0)}\to V_{g_a}$,
$V_{-g_a}\to V_{g(\kappa_0)}$, $V_{-g_a}\to V_{g_a}$, of degrees $g_a+g(\kappa_0)=1-g'+g(\kappa_0)
\ge2$ (twice) and $2g_a\ge2$; so $n\in N^{\ge2}(\auA)$. Next, $wf_a=Tf_a$ and
\[
w(n-1)=Tf_a\otimes\kappa^\flat-w\kappa\otimes f_a^\flat-\tfrac12\langle\kappa,\kappa\rangle
Tf_a\otimes f_a^\flat .
\]
The first and third terms have degrees $g_a-1+g(\kappa_0)\ge1$ and $2g_a-1\ge1$. The second maps
$V_{-g_a}$ to the line of $w\kappa_0=T\kappa_0+v\kappa_0$; $T\kappa_0$ has grade
$g(\kappa_0)-1\ge g'=1-g_a$, the components of $v\kappa_0$ of negative degree have grades $\ge g'$
by the choice of $c$, and the others have grades $\ge g(\kappa_0)>g'$. So the second term has
degree $\ge g'+g_a=1$, and $n\in A^I_w(\auA)$. By (G3),
\[
\pktr(w^\dagger w(n-1))=\langle w\kappa,Tf_a\rangle-\langle Tf_a,w\kappa\rangle-\tfrac12\langle
\kappa,\kappa\rangle\langle Tf_a,Tf_a\rangle ,
\]
and the same holds with $w$ replaced by $T$. Subtracting,
$\pktr((\mathfrak e-w^\dagger w)(n-1))=\langle Tf_a,v\kappa\rangle-\langle v\kappa,Tf_a\rangle=
\varsigma-\aubar\varsigma$ with $\varsigma=\mathsf t\,\aubar{t_a}\,v_{c,\kappa_0}\langle f'_a,f'_c
\rangle$, and $\langle f'_a,f'_c\rangle\ne0$ as $a+c=s-2$. So $\psi_j\chi^I_w(n)=\psi(\jmath
\varsigma)$, which is not $1$ for a suitable $\mathsf t$.
\end{proof}

\begin{proof}[Proof of Proposition~\ref{ps:onestep}]
The case $j=m-1$ and step (1) (the properties of $\psW'$) are as in \S\ref{ps:ss-seesaw}. If
$s=1$ then $V^+=0$, $I=0$, $T=0$, $R'=G'$ and there is nothing to prove; let $s\ge2$. Fix
$h\in G'(\auA)$, write $T_+=T|_{V^+}$, and let $\mu(h)$ be a function of moderate growth on
$G'(\auA)$ with
\begin{equation}\label{pk:eq-major}
\sum_{w\in\auHom(V^+,V')(E)}\bigl|c^0_w(\omega(r,h)F)\bigr|\le\mu(h)\qquad(r\in\mathsf K_R),
\end{equation}
where $\mathsf K_R\subset R(\auA)$ is a compact set with $R(E^+)\mathsf K_R=R(\auA)$; it exists by
(W3) applied to $I_0$. In detail: by the law $c^0_w(\omega(\gamma,\delta)F')=c^0_{\delta^{-1}w\gamma}(F')$ of
\S\ref{ps:ss-seesaw} (with $\gamma=1$), the left side of \eqref{pk:eq-major} is unchanged when $h$ is
replaced by $\delta h$, $\delta\in G'(E^+)$, the twisting characters being trivial on rational points.
By {\cite[Lemma~3.28]{P53}} write $h=\delta^{-1}x$ with $\delta$ rational, $x_f$ in a fixed compact set and
$\|x_\infty\|\le C\,\mathrm{ht}(h)$ ({\cite[Definition~3.27]{P53}}); by Lemma~\ref{pz:absgrowth}(b) (for
$F$ a sum of tensors $\Psi\otimes f$ with $\Psi$ smooth, $(r,x)$ with $(r_f,x_f)$ in a compact set and
$r_\infty$ bounded) the left side is at most $C'\,\mathrm{ht}(h)^{A_0}$, where $A_0=4(D_W+1)$ and $2D_W=\dim_\R\auHom(V,V')_\infty$ (the exponent $4s$ of that lemma; we rename it, $s$ being the rank here). So $\mu=C'\,\mathrm{ht}^{A_0}$ will do.

\emph{(i) Coarse expansion.} By (W3), $\theta(\omega(r,h)F)=\sum_wc^0_w(\omega(r,h)F)$. The set of
$w\notin G'(E^+)T_+$ is stable under $R(E^+)$, which acts through its restriction to $V^+$ (the
orbit $G'(E^+)T_+$ is stable by (ii) below), and consists of dead classes
(Lemma~\ref{pk:coarse}). By Lemma~\ref{pk:vanish} and (c2),
\[
\int_{[R]}\psW(r)\theta(\omega(r,h)F)\,dr=\sum_{\delta\in\mathsf M(E^+)\backslash G'(E^+)}\Psi(\delta h),
\qquad\Psi(h)=\int_{[R]}\psW(r)\,c^0_{T_+}(\omega(r,h)F)\,dr,
\]
where $\mathsf M$ is the stabiliser of $T_+$ in $G'$, and $\sum_\delta|\Psi|(\delta h)\le\|\psW
\|_\infty\mu(h)$, where $|\Psi|$ is defined with absolute values inside the integral.

\emph{(ii) The orbit of $T_+$ under $R'$.} By the description of $R'$ in \S\ref{ps:ss-Y} and (G2),
$R'$ is the group of $\epsilon\in G'$ with $\epsilon^{-1}Tf_i\in Tf_i+\langle Tf_b:b<i\rangle$ for
$i<p$; so $\mathsf M\subseteq R'$ and $\epsilon^{-1}T_+-T_+\in J(E)$ for $\epsilon\in R'(E^+)$.
Conversely, for $\eta\in J(E)$ the map $T_++\eta$ is injective and, by (G1), (G2),
$Q_{T_++\eta}=Q_T$; by Witt's theorem $T_++\eta=\epsilon^{-1}T_+$ with $\epsilon\in G'(E^+)$, and
then $\epsilon\in R'$. Hence $\epsilon\mapsto\epsilon^{-1}T_+$ is a bijection $\mathsf M(E^+)
\backslash R'(E^+)\to T_++J(E)$, and by (c2)
\begin{equation}\label{pk:eq-ii}
\sum_{\epsilon\in\mathsf M(E^+)\backslash R'(E^+)}\Psi(\epsilon h)=\int_{[R]}\psW(r)\sum_{\eta\in
J(E)}c^0_{T_++\eta}(\omega(r,h)F)\,dr .
\end{equation}
\emph{(iii) From coarse to fine classes.} Let $F'\in\psHS(W)$, let $\tilde T_+$ and $\tilde\eta$
be the extensions by $0$, and $\tilde J=\{\tilde\eta\}\subseteq I$. We claim
\begin{equation}\label{pk:eq-poisson}
\sum_{\eta\in J(E)}c^0_{T_++\eta}(F')=\sum_{w\in\mathcal C}c_w(F'),\qquad\mathcal C=\bigl(T+I_0^\perp
(E)\bigr)\big/\bigl(I_0^\perp(E)\cap I^\perp(E)\bigr),
\end{equation}
the set of classes modulo $I^\perp(E)$ of the maps $T+v$, $v|_{V^+}=0$. First, $I_0\cap\tilde
J^\perp=I_0\cap I$: for $i\in I_0$, $\langle\tilde\eta,i\rangle$ pairs the components of $\eta$
from $V_g$ to $V'_{g'}$, $g'>g>0$, with the components of $i$ from $V_{-g}$ to $V'_{-g'}$, and
these are the components of $i$ of negative degree; by (G1) a map of degree $\ge0$ on $V^-$ has
degree $>0$. Put $I_0'=I_0\cap I$. Since $\tilde J\cap I_0^\perp=0$, the pairing identifies
$\tilde J(E)$ with the character group of the compact group $\mathsf B=I_0(\auA)/(I_0(E)+
I_0'(\auA))$. The function $\phi(i)=\theta(\psrho(i)\psrho(\tilde T_+)F')$ on $[I_0]$ is
continuous, and by (W1), as $\psrho(i)\psrho(\tilde\eta)=[i,\tilde\eta]\psrho(\tilde\eta)\psrho(i)$
and $\theta\circ\psrho(\tilde\eta)=\theta$, one has $c^0_{T_++\eta}(F')=\int_{[I_0]}[i,\tilde\eta]
\phi(i)\,di$. These are the Fourier coefficients of the continuous function
$\phi_1(\mathsf b)=\int_{[I_0']}\phi(\mathsf b+i')\,di'$ on $\mathsf B$; they are absolutely
summable by (W3), so their sum is $\phi_1(0)=\theta_{W_{I_0'}}(\psZ_{I_0'}\psrho(\tilde T_+)F')$.
Now apply (W3) to the space $W_{I_0'}=I_0'^\perp/I_0'$ and its isotropic subspace $I/I_0'$, and
(W2) ($\psZ_I=\psZ_{I/I_0'}\psZ_{I_0'}$): the sum is $\sum_uc_{\tilde T_++u}(F')$ over
$u\in I_0'^\perp(E)/I^\perp(E)$. Finally $I_0'^\perp=I_0^\perp+I^\perp$ and $T-\tilde T_+\in
I_0^\perp$, which gives \eqref{pk:eq-poisson}.

\emph{(iv) Only the class of $T$ survives.} The set $\mathcal C\smallsetminus\{T\}$ is stable
under $R(E^+)$: for $\gamma\in R(E^+)$, $T\gamma-T$ has degree $\ge0$, and $v\gamma$ vanishes on
$V^+$ and has degree $\ge0$ if and only if $v$ has. Its elements are dead for $I$
(Lemma~\ref{pk:fine}). By Lemma~\ref{pk:vanish}, \eqref{pk:eq-ii} and \eqref{pk:eq-poisson} with
$F'=\omega(r,h)F$,
\[
\sum_{\epsilon\in\mathsf M(E^+)\backslash R'(E^+)}\Psi(\epsilon h)=\int_{[R]}\psW(r)\,c_T(\omega(r,h)F)
\,dr=\psW'(h).
\]
\emph{(v) Conclusion and convergence.} By (i) and (iv), $\int_{[R]}\psW(r)\theta(\omega(r,h)F)\,dr=
\sum_{\gamma\in R'(E^+)\backslash G'(E^+)}\psW'(\gamma h)$ and $\sum_\gamma|\psW'(\gamma h)|\le
\sum_\delta|\Psi|(\delta h)\le\|\psW\|_\infty\mu(h)$. Let $z$ be rapidly decreasing on $[G']$. Then
$\int_{R'(E^+)\backslash G'(\auA)}|\psW'z|=\int_{[G']}\sum_\gamma|\psW'(\gamma h)|\,|z(h)|\,dh\le
\|\psW\|_\infty\int_{[G']}\mu|z|<\infty$. (Here ``rapidly decreasing'' means: on a Siegel set
$\mathfrak S=\mathfrak S_\infty\mathsf cK'$ of $G'$ with $G'(\auA)=G'(E^+)\mathfrak S$, as in the proof of
Proposition~\ref{pz:types}(a), $|z(g)|\le C_N\|g_\infty\|^{-N}$ for every $N$. On $\mathfrak S$,
$\mathrm{ht}(g)\le c_0^{-1}\|g_\infty\|$: for $g=g_\infty ck'$ and $0\ne\gamma\in\Lambda_g$, the vector
$c^{-1}\gamma$ lies in one of finitely many lattices, so $\|\gamma\|\ge c_0>0$ and $|\gamma|_g=
\|g_\infty^{-1}\gamma\|\ge c_0\|g_\infty\|^{-1}$; this is the paragraph ``Heights on Siegel sets'' in the
proof of {\cite[Proposition~4.10]{P53}}. Hence $\mu|z|=C'\mathrm{ht}^{A_0}|z|$ is bounded on $\mathfrak S$, hence
on $G'(\auA)$ by left $G'(E^+)$-invariance, and $[G']$ has finite volume.) And since $|\theta(\omega(r,h)F)|\le\mu(h)$ for
$r\in\mathsf K_R$, Fubini's theorem gives
\[
\int_{[R]}\psW(r)\int_{[G']}\theta(\omega(r,h)F)z(h)\,dh\,dr=\int_{[G']}z(h)\sum_\gamma\psW'(\gamma h)
\,dh=\int_{R'(E^+)\backslash G'(\auA)}\psW'z .\qedhere
\]
\end{proof}

\subsubsection{Support and finite data in the seesaw identity}\label{pk:ss-rest}

We give the verifications that complete Lemmas~\ref{ps:support} and~\ref{ps:fin} and the proof
of {\cite[Proposition~14.23]{P46}}.

\emph{(R1) Iwasawa decomposition (used in Lemma~\ref{ps:support}).} Let $v$ be a finite place of
$E^+$, $\Lambda$ a unimodular lattice in $Y_j\otimes E^+_v$, $K_v$ its stabiliser and $P$ the
stabiliser of the flag of the spaces $\langle f_0,\dots,f_i\rangle$, $i<p$. Then
$G_j(E^+_v)=P(E^+_v)K_v$. Indeed $E/E^+$ is unramified at $v$, so $K_v$ is hyperspecial and this
is the Iwasawa decomposition \cite[\S3.3.2, p.~51]{TitsCorvallis}. (Directly: $K_v$ is transitive on the
isotropic flags of a given type, because for a totally isotropic direct summand $\Lambda_1$ of
$\Lambda$ there is a totally isotropic $\Lambda_1^*\subset\Lambda$ with $\Lambda=(\Lambda_1\oplus
\Lambda_1^*)\perp\Lambda_2$, the trace $\auO_{E,v}\to\auO_{E^+,v}$ being surjective, and
$\auGL(\Lambda_1)\subset K_v$ is transitive on the full flags of $\Lambda_1$ with saturated
members.)

\emph{(R2) The step with $\mathsf z_w$ in Lemma~\ref{ps:support}.} Let $j$ be even, $h=n\,l\,k$,
$\Lambda'=l\psM_j$, and suppose that $\Lambda'\cap E_vz$ is larger than $\widehat\auO u_jz$ at a
place $v=w\aubar w$, $w\mid\pi_N$. As shown there, the part of grade $0$ of $\Lambda'_v$ is then the
orthogonal sum of its intersections with $X_j$ and $E_vz$, so $\iota(\mathsf z_w)$, which is a
unit scalar on $X_j$ and the identity elsewhere, stabilises $\Lambda'$. Hence
$\mathsf u_e=h_f^{-1}n\iota(\mathsf z_w^{\,e})n^{-1}h_f=k^{-1}(l^{-1}\iota(\mathsf z_w^{\,e})l)k$ lies
in $K_j$; it is $1$ outside $v$, so $\mathsf u_e\in J_j$, and $\psW_j(n\iota(\mathsf z_w^{\,e})n^{-1}
h)=\psW_j(h\mathsf u_e)=\psW_j(h)$. On the other hand $n\iota(\mathsf z_w^{\,e})n^{-1}=n'\iota(
\mathsf z_w^{\,e})$ with $n'=n\cdot\iota(\mathsf z_w^{\,e})n^{-1}\iota(\mathsf z_w^{\,e})^{-1}\in
N_j(\auA)=N_j^{\ge2}(\auA)$ ($j$ is even) and $\psi_j(n')=\psi_j(n)\psi_j(n)^{-1}=1$, because
$\psi_j$ is invariant under conjugation by $R_j(\auA)$; so $\psW_j(n\iota(\mathsf z_w^{\,e})n^{-1}h)=
\psW_j(\iota(\mathsf z_w^{\,e})h)$. By Lemma~\ref{ps:double}(a), applied $j/2$ times,
$\psW_j(\iota(g)h)=l_h(g_\infty^{-1}\cdot(C_h\varphi)(g_f))$ for $g\in\auU(X_j)(\auA)$, where $C_h$
is a composite of lifts \eqref{pk:eq-lift}. Each lift commutes with the translation by the central
element $\mathsf z_w$: the pair $(\mathsf z_w,\mathsf z_w)$ acts trivially on $\auHom(X_i,X_{i+1})$,
and the determinant twists are trivial on it since $\mu_{i+1}$ and $\upsilon$ are unramified above
$N$. Therefore $(N-1)\psW_j(h)=\sum_{e=0}^{N-2}\psW_j(\iota(\mathsf z_w^{\,e})h)=l_h\bigl((C_h
\sum_e\mathsf z_w^{\,e}\varphi)(1)\bigr)=0$ for $\varphi\in\auSram$.

\emph{(R3) The lattice $\Lambda_0$ at the chosen places (Lemma~\ref{ps:fin}(3)).} Let $w$ be a
chosen place. There $N$ is a unit, $\alpha=1$, $\mathsf a_w=1$,
$\mathsf a_{\bar w}=\varpi^{-1}$, so $u_j=\varpi^{e_j}$ at $w$ and $\varpi^{-e_j}$ at $\aubar w$,
$e_j=\lceil(m-j)/2\rceil$. By \eqref{ps:eq-lattice}, for $j$ even $y/\aubar{\pi_N}+u_jz/N$ and
$u_jz$ span the same module as $y$ and $u_jz$; so $\psM_j=\widehat\auO^{\,n_j}\oplus u_j\bigoplus_a
\widehat\auO f_a$ at $w$ and $\aubar w$, the orthogonal sum of its parts in $X_j$ and in the
string, and likewise for $\psM_{j+2}$ if $j+2<m$. The lattice $\Lambda_0$ is the image in
$X^-\oplus X''$ of the vectors of $\psM_j\oplus\psM_{j+2}$ without components of negative grade and
with component of grade $0$ in $X^-\oplus X''\oplus\ell$; for a split lattice this image is
$\widehat\auO^{\,n_j}\oplus\widehat\auO^{\,n_{j+2}}$, whatever $u_j,u_{j+2}$ are. Thus $\Lambda_0$
is standard at the chosen places for all pairs with $j+2<m$. For $j+2=m$: $\psM_{m-2}=\widehat
\auO^{\,n_{m-2}}\oplus\widehat\auO u_{m-2}z$ and $z=\pi_Ne_{\mathrm{last}}$ with $\pi_N$ a unit, so
$\Lambda_0=\widehat\auO^{\,n_{m-2}}\oplus\widehat\auO^{\,n_m-1}\oplus u_{m-2}\widehat\auO
e_{\mathrm{last}}$, as stated in Lemma~\ref{ps:fin}.

\emph{(R4) The denominators in Lemma~\ref{ps:fin}.} For a chosen place $w$ above $p\in\{p_0,p_1\}$,
$\#_w=p^{\,\operatorname{rk}Y_j\operatorname{rk}Y_{j+1}}$ divides $p^{n_m^2}$; there are $f$ chosen
places above each of $p_0,p_1$, so $\prod_{w\text{ chosen}}\#_w$ divides $M^{fn_m^2}=B$ for each of
the $m$ data $\Upsilon_j$. By (R3), $\mathsf T_i$ is trivial unless $2i+2=m$, and then
$\#\mathsf T_i=\prod_{w}p^{\,n_{m-1}}$ divides $B$. So the product of all denominators at the primes
not dividing $q$ divides $B^{m+1}$, hence $B^{2m}$.

\emph{(R5) Constants in the proof of {\cite[Proposition~14.23]{P46}}.} $C_0$ is the constant of
Proposition~\ref{pk:conj}(ii); it depends on $q_0,d$ only, and it is a term of the constant $C_P$ in
the definition of admissible data, so that $C_0\le\log\log b$ (\S\ref{pz:ss-const},
Lemma~\ref{pz:C0}; there is no circularity, $C_0$ being defined from $q_0$ and $K$ alone). One
may assume $\sigma|_K=1$: if $\sigma|_K\ne1$, then $c\circ\sigma$ fixes $K$ and $|c\sigma(w)|=|\sigma(w)|$. The
bound $q_1\le x^{12}$ is {\cite[Lemma~13.4(S3)]{P46}}.

\emph{(R6) Lemma~\ref{ps:level}.} The factor $\upsilon'_w$ is trivial on units because $\upsilon$
is unramified at every finite place (\S\ref{pk:ss-conv}).

\subsubsection{Integrality of the branching map: proof of Lemma~\ref{ps:branching}(c)}
\label{pk:ss-branch}

All modules below are tensor products over $\tau\in\auPl$ and all statements are multiplicative
in $\tau$; we write them for one $\tau$ and one group $\auGL_N$. Let $\mu$ be a diagram with
columns $\mathfrak c$ of heights $h_{\mathfrak c}\le N$ and $|\mu|=\sum h_{\mathfrak c}$ boxes.

\emph{The tensor model.} Let $\mathbb T_\mu(R)=\bigotimes_{\mathfrak c}\Lambda^{h_{\mathfrak c}}R^N$
for a ring $R$, with the basis $e_U=\bigotimes_{\mathfrak c}e_{U_{\mathfrak c}}$ indexed by the
fillings $U$ (a subset $U_{\mathfrak c}\subseteq[N]$ of $h_{\mathfrak c}$ rows for each column), and
the symmetric bilinear form $(\ ,\,)$ for which this basis is orthonormal. For a matrix $y$ with
columns $y_1,\dots,y_N\in R^N$ put $\pkev_y=\bigotimes_{\mathfrak c}y_1\wedge\dots\wedge
y_{h_{\mathfrak c}}$; then $(e_U,\pkev_y)=b_U(y):=\prod_{\mathfrak c}\Delta_{U_{\mathfrak c},
[h_{\mathfrak c}]}(y)$, the product of minors of \S\ref{ps:ss-restrict}. So the multiplication map
is $\pkm(u)(y)=(u,\pkev_y)$, and it is equivariant: $\pkm(gu)(y)=\pkm(u)(g^ty)$. Let
$\mathsf C_\mu\subseteq\mathbb T_\mu(\C)$ be the span of the $\pkev_y$, $y\in M_N(\C)$. It is
defined over $\Q$, so $\mathbb T_\mu(\C)=\mathsf C_\mu\oplus\mathsf C_\mu^\perp$, and
$\ker\pkm=\mathsf C_\mu^\perp$ by the formula for $\pkm$; let $\pi_\mu$ be the orthogonal projection
onto $\mathsf C_\mu$. Thus $\pkm=\pkm|_{\mathsf C_\mu}\circ\pi_\mu$ and $\pkm\colon\mathsf C_\mu\to
W_\mu$ is an isomorphism. The pull-back of the invariant form of $W_\mu$ to $\mathsf C_\mu$ is
invariant under $\auU(N)$, and $W_\mu$ is irreducible under $\auGL_N(\C)$~\cite[\S8.1, Thm.~1 and its corollary; \S8.2, Thm.~2]{FultonYT}, hence under its Zariski-dense subgroup $\auU(N)$; so it is a multiple of
the standard hermitian form, and the multiple is $1$ because $e_{U_0}=\pkev_1$ ($U_0$ the filling
by the leading rows) has norm $1$ and $\pkm(e_{U_0})=f_\mu:=\prod\Delta_{[h_{\mathfrak c}],
[h_{\mathfrak c}]}$ has norm $1$. Consequently, for $w\in\mathsf C_\mu$,
$\langle\pkm w,b_U\rangle=\langle w,\pi_\mu e_U\rangle=(w,e_U)$, and
\begin{equation}\label{pk:eq-dual}
\Sigma_\mu=\pkm(\mathbb T_\mu(\Z)),\qquad W_\mu(\Z)=\Sigma_\mu^\#=\pkm\bigl(\mathsf C_\mu\cap
\mathbb T_\mu(\Z)\bigr).
\end{equation}
We also use that the standard products of minors are a $\Z$-basis of $\Sigma_\mu$ and remain
linearly independent modulo every prime~\cite[Thm.~2.1]{DEP}, \cite[\S8.1, Thm.~1 and its corollary]{FultonYT}; hence
\begin{equation}\label{pk:eq-sat}
\Sigma_\mu=(\Sigma_\mu\otimes\Q)\cap\Z[y_{ab}] .
\end{equation}

\begin{lemma}[The hook number]\label{pk:hook}
Let $\langle P,Q\rangle_F=(P(\partial)\aubar Q)(0)=\sum_\alpha\alpha!\,P_\alpha\aubar{Q_\alpha}$ be the
Fischer form on polynomials in the entries of $y$, and $H_\mu=\langle f_\mu,f_\mu\rangle_F$. Then
$H_\mu$ is a positive integer, $H_\mu\le|\mu|^{3|\mu|}$ $(0^0=1)$, the invariant form of $W_\mu$ is
$H_\mu^{-1}\langle\ ,\,\rangle_F$, and $H_\mu\Sigma_\mu\subseteq\Sigma_\mu^\#$.
\end{lemma}
\begin{proof}
The Fischer form is positive definite and invariant under unitary substitutions of the variables,
in particular under $P\mapsto P(g^ty)$, $g\in\auU(N)$; on the irreducible module $W_\mu$ it is
therefore $H_\mu$ times the invariant form. Products of minors have integer coefficients, so
$\langle b,b'\rangle_F\in\Z$ and $\langle H_\mu b,b'\rangle=\langle b,b'\rangle_F\in\Z$ for
$b,b'\in\Sigma_\mu$. Finally $f_\mu$ is homogeneous of degree $L=|\mu|$ and the sum of the absolute
values of its coefficients is at most $\prod_{\mathfrak c}h_{\mathfrak c}!\le\prod h_{\mathfrak
c}^{h_{\mathfrak c}}\le L^L$; so $H_\mu\le L!\,(L^L)^2\le L^{3L}$.
\end{proof}

\begin{lemma}[Leading rows]\label{pk:leading}
Let $0\le p<N$ and let $\mu^\flat$ be the diagram, for $\auGL_{N-p}$, with the columns of heights
$\max(h_{\mathfrak c}-p,0)$. Let $(r_1,\dots,r_p)$ be distinct rows and $U$ a filling of $\mu$ such
that $U_{\mathfrak c}\supseteq\{r_1,\dots,r_p\}$ if $h_{\mathfrak c}\ge p$ and $U_{\mathfrak c}=
\{r_1,\dots,r_{h_{\mathfrak c}}\}$ if $h_{\mathfrak c}<p$. Then $H_{\mu^\flat}\,b_U\in\Sigma_\mu^\#$.
\end{lemma}
\begin{proof}
Permutation matrices preserve $\Sigma_\mu$, $\Sigma_\mu^\#$ and the form, and map $b_U$ to $\pm$
the product of minors of the permuted filling; so let $r_i=i$. Then $U=\pkLW(T)$ for a unique
filling $T$ of $\mu^\flat$: $U_{\mathfrak c}=[p]\cup(p+T_{\mathfrak c})$ if $h_{\mathfrak c}\ge p$.
Define $\pkLW\colon\mathbb T_{\mu^\flat}(\Z^{N-p})\to\mathbb T_\mu(\Z^N)$ by $e_T\mapsto
e_{\pkLW(T)}$, for all fillings $T$ of $\mu^\flat$. (1) For $y\in M_{N-p}$,
$\pkLW(\pkev_y)=\pkev_{\operatorname{diag}(1_p,y)}$, since the first $h$ columns of
$\operatorname{diag}(1_p,y)$ have exterior product $e_{[h]}$ for $h\le p$ and $e_{[p]}\wedge(y_1
\wedge\dots\wedge y_{h-p})$ for $h>p$. Hence $\pkLW(\mathsf C_{\mu^\flat})\subseteq\mathsf C_\mu$.
(2) For $x=\bigl(\begin{smallmatrix}A&B\\C&D\end{smallmatrix}\bigr)$ with $A\in\auGL_p$ and
$S(x)=D-CA^{-1}B$, the formula for the determinant of a block matrix gives
$\Delta_{[p]\cup(p+I),[p+k]}(x)=\det A\cdot\Delta_{I,[k]}(S(x))$ for $I\subseteq[N-p]$, $|I|=k$ (the
first $k$ columns of $S(x)$ depend only on the first $p+k$ columns of $x$). Hence, for every
filling $T$ of $\mu^\flat$, $b_{\pkLW(T)}(x)=\mathsf F(x)\,b^\flat_T(S(x))$ with $\mathsf F(x)=
\prod_{h_{\mathfrak c}<p}\Delta_{[h_{\mathfrak c}],[h_{\mathfrak c}]}(x)\cdot\det A^{\,\#\{\mathfrak
c:h_{\mathfrak c}\ge p\}}$ independent of $T$; by linearity $\pkm(\pkLW u)(x)=\mathsf F(x)\,
\pkm^\flat(u)(S(x))$ for all $u$, on a dense open set. So $\pkLW(\ker\pkm^\flat)\subseteq\ker\pkm$,
that is $\pkLW(\mathsf C_{\mu^\flat}^\perp)\subseteq\mathsf C_\mu^\perp$; with (1),
$\pi_\mu\circ\pkLW=\pkLW\circ\pi_{\mu^\flat}$. (3) By Lemma~\ref{pk:hook} and \eqref{pk:eq-dual} for
$\mu^\flat$, $H_{\mu^\flat}\pi_{\mu^\flat}(e_T)\in\mathsf C_{\mu^\flat}\cap\mathbb T_{\mu^\flat}(\Z)$.
Applying $\pkLW$, which preserves integrality: $H_{\mu^\flat}\pi_\mu(e_U)\in\mathsf C_\mu\cap\mathbb
T_\mu(\Z)$, and $\pkm$ of this vector is $H_{\mu^\flat}b_U$; conclude by \eqref{pk:eq-dual}.
\end{proof}

\begin{proposition}[Lemma~\ref{ps:branching}(c)]\label{pk:branching}
Let $\lambda=\lambda(X_m)$, $n=n_m$, and let $\lambda^{\circ\flat}_\tau$ be $\lambda^\circ_\tau$
without its first $k-1$ rows. Put $c_{\mathrm{br}}=\prod_\tau H_{\lambda^{\circ\flat}_\tau}$ and
$\mathrm{br}=c_{\mathrm{br}}\,i^*$, with $i$ as below. Then {\rm(a), (b), (c)} of
Lemma~\ref{ps:branching} hold, and $c_{\mathrm{br}}\le L^{3L}$.
\end{proposition}
\begin{proof}
We write the proof for one $\tau$. Call a column of $\lambda$ \emph{kept} or \emph{lost}
according as it keeps its height in $\lambda^\circ$ or loses a box; the kept columns have heights
$\ge k$. Let $\varepsilon\colon\mathbb T_{\lambda^\circ}(\Z^{n-1})\to\mathbb T_\lambda(\Z^n)$ send
$e_{U^\circ}$ to the tensor with factor $e_{U^\circ_{\mathfrak c}}$ in a kept column and
$e_n\wedge e_{U^\circ_{\mathfrak c}}$ in a lost one. It is isometric, integral and equivariant for
$\auGL_{n-1}$. Put $i=\pkm\circ\varepsilon\circ(\pkm^\circ|_{\mathsf C_{\lambda^\circ}})^{-1}\colon
W_{\lambda^\circ}\to W_\lambda$. Then $i$ is $\auGL_{n-1}$-equivariant, $i(f_{\lambda^\circ})=\pm v$
with $v$ the vector of the proof of Lemma~\ref{ps:branching}(a), and under the identifications
\eqref{pk:eq-dual} $i=\pi_\lambda\varepsilon|_{\mathsf C_{\lambda^\circ}}$, $i^*=\pi_{\lambda^\circ}
\varepsilon^t|_{\mathsf C_\lambda}$; so $\|i^*\|\le1$. This gives (a) and (b) as in
\S\ref{ps:ss-restrict}. By duality, (c) is equivalent to
\begin{equation}\label{pk:eq-c}
c_{\mathrm{br}}\cdot i(b^\circ_{T'})\in\Sigma_\lambda\qquad\text{for every filling }T'\text{ of }
\lambda^\circ .
\end{equation}
\emph{The kernel.} For $x\in M_{n-1}(R)$, whose columns are regarded as vectors of $R^n$ with last
coordinate $0$, and $y\in M_n(R)$ let $\mathcal K(x,y)=(\varepsilon\,\pkev^\circ_x,\pkev_y)$. Then
\begin{equation}\label{pk:eq-kernelbr}
\mathcal K(x,y)=(\pkev^\circ_x,\pi_{\lambda^\circ}\varepsilon^t\pkev_y),\qquad
(\pi_{\lambda^\circ}\varepsilon^t\pkev_y,e_{T'})=(\pkev_y,\varepsilon\pi_{\lambda^\circ}e_{T'})=
i(b^\circ_{T'})(y).
\end{equation}
\emph{Special $y$.} All columns of $\lambda$ have height $\le n_{m-1}\le n-2$. Let the last row
of $y$ be $(\mathsf t,0,\dots,0)$ and let $g\in M_{n-1}(R)$ be the matrix of the first $n-1$ rows
and columns of $y$, with columns $g_1,\dots,g_{n-1}$ (the last column of $y$ does not occur in
$\pkev_y$). In a kept column of height $h$ the factor of
$\mathcal K$ is $\det(x_a\cdot g_b)_{a,b\le h}=\Delta_{[h],[h]}(g^tx)$; in a lost column of height
$h$ it is the determinant of the matrix with first row $(\mathsf t,0,\dots,0)$ and rows
$(x_a\cdot g_b)_{b\le h}$, $a<h$, that is $\mathsf t\,\Delta_{\{2,\dots,h\},[h-1]}(g^tx)$. So
\[
\mathcal K(x,y)=\mathsf t^{\,\ell}\,\mathsf P(g^tx),\qquad\mathsf P=\prod_{\text{kept}}\Delta_{[h],[h]}
\cdot\prod_{\text{lost}}\Delta_{\{2,\dots,h\},[h-1]}\ \in\Sigma_{\lambda^\circ},
\]
with $\ell$ the number of lost columns. The product $\mathsf P$ is $b^\circ_U$ for a filling $U$
of $\lambda^\circ$ to which Lemma~\ref{pk:leading} applies with $p=k-1$ and the rows
$(2,3,\dots,k)$: a kept column has height $h\ge k$ and rows $[h]\supseteq\{2,\dots,k\}$; a lost
column, of height $h-1$ in $\lambda^\circ$, has rows $\{2,\dots,h\}$, which contain $\{2,\dots,k\}$
if $h-1\ge k-1$ and are the first $h-1$ of the rows $(2,\dots,k)$ otherwise. Hence $c\,\mathsf P=
\pkm^\circ(w_0)$ with $c=H_{\lambda^{\circ\flat}}$ and $w_0\in\mathsf C_{\lambda^\circ}\cap\mathbb
T_{\lambda^\circ}(\Z)$, and $c\,\mathcal K(x,y)=\mathsf t^{\,\ell}(w_0,\pkev^\circ_{g^tx})=(\mathsf
t^{\,\ell}\,g\,w_0,\pkev^\circ_x)$, where $g$ acts on the tensors.

\emph{General $y$.} Work over $R=\Z[y_{ab}]$ and let $y'$ have the columns $y'_1=y_1$,
$y'_b=y_{n1}y_b-y_{nb}y_1$ $(b\ge2)$. Its last row is $(y_{n1},0,\dots,0)$, and
$\pkev_{y'}=y_{n1}^{\,\mathsf M}\pkev_y$ with $\mathsf M=\sum_{\mathfrak c}(h_{\mathfrak c}-1)$. By
the special case and \eqref{pk:eq-kernelbr}, with $g'$ the matrix of the first $n-1$ rows and
columns of $y'$, $y_{n1}^{\,\mathsf M}\,c\,\pi_{\lambda^\circ}
\varepsilon^t\pkev_y$ and $y_{n1}^{\,\ell}g'w_0$ are two vectors of $\mathsf C_{\lambda^\circ}
\otimes\Q(y)$ with the same pairing against all $\pkev^\circ_x$; so they are equal, and the second
has coordinates in $R$. Therefore $y_{n1}^{\,\mathsf M}\cdot c\,i(b^\circ_{T'})(y)\in\Z[y_{ab}]$,
and as $c\,i(b^\circ_{T'})\in\Q[y_{ab}]$ and $y_{n1}$ is a prime element of $\Z[y_{ab}]$,
$c\,i(b^\circ_{T'})\in\Z[y_{ab}]$. Since $i(b^\circ_{T'})\in W_\lambda=\Sigma_\lambda\otimes\C$ has
rational coordinates, \eqref{pk:eq-sat} gives \eqref{pk:eq-c}.

\emph{The size.} $|\lambda^{\circ\flat}_\tau|$ is at most the number $L_\tau$ of boxes of
$\lambda^\circ_\tau$ in the columns of height $\ge k$, so $c_{\mathrm{br}}\le\prod_\tau
L_\tau^{3L_\tau}\le L^{3L}$ by Lemma~\ref{pk:hook}.
\end{proof}

\subsubsection{One curve with complex multiplication}\label{pk:ss-curve}

We use the notation of \S\ref{pa:ss-curve}: $\Lambda$ with $\auO_K\Lambda=\Lambda$, of area $A$,
the point $t$ with $\wp''(t)=2$, $a_1$, the curve $\mathcal E$, $x,y$, $\theta$, $g=-\log\theta$,
$H(u,w)=\boldsymbol\pi\aubar uw/A$, and $s_2$. The transformation law of $\vartheta$ quoted there
is equivalent to $\zeta(z+\omega)=\zeta(z)+\eta(\omega)$ with
$\eta(\omega)=s_2\omega+(\boldsymbol\pi/A)\aubar\omega$. For $u\in\C$ the \emph{twisted translate} is
\[
\theta_u(z)=\exp\bigl(-\tfrac{\boldsymbol\pi}A(z\aubar u+\tfrac12|u|^2)\bigr)\,\theta(z+u);
\]
$|\theta_u(0)|$ is $\Lambda$-periodic in $u$, and $\theta_{u+\omega}=\chi(\omega)e^{i(\boldsymbol\pi/A)
\Iml(\aubar\omega u)}\theta_u$ for $\omega\in\Lambda$ by \eqref{pa:eq-theta}. We call
$(\Lambda,t,a_1)$ a \emph{normalised lattice} if $t$ is the admissible point of order two
(\S\ref{pf:ss-adm}). A power series $F(z)$ in one or several variables
is \emph{$\partial$-integral} at a homomorphism $\C\to\C_\ell$ if all its derivatives at $0$ are in
$\pao$; sums and products of $\partial$-integral series are $\partial$-integral, and so is
$\Phi\circ F$ if $F(0)=0$ and $\Phi\in\pao[[X]]$ or $\Phi=\exp$.

\begin{proposition}\label{pk:curve}
Let $(\Lambda,t,a_1)$ be a normalised lattice.
\begin{enumerate}
\item[(C1)] {\rm(Conjugation.)} For $\sigma\in\auAut(\C/K)$ there is a normalised lattice
$(\Lambda^\sigma,t^\sigma,\sigma a_1)$, unique, such that $\sigma$ maps the Taylor series at $0$ of
$x$, $y$, $\theta$ to those of $\Lambda^\sigma$. One has $g_i(\Lambda^\sigma)=\sigma g_i(\Lambda)$
$(i=2,3)$ and $s_2(\Lambda^\sigma)=\sigma s_2(\Lambda)$.
\item[(C2)] {\rm(Integrality.)} For every prime $\ell\ne d$ and every $\C\to\C_\ell$, $x$, $y$ and
$\theta$ are $\partial$-integral.
\item[(C3)] {\rm(Translates.)} Let $u\in\C$ with $\mathsf Nu\in\Lambda$, $\mathsf N\ge1$,
$u\notin t+\Lambda$.
\begin{itemize}
\item[(i)] Let $\ell$ be an odd prime not dividing $\mathsf N$ and $\C\to\C_\ell$. Then
$\theta_u(0)\in\pao$ (this holds at every prime, Proposition~\ref{pf:values}) and $\theta_u(z)=\theta_u(0)\,\theta(z)\,\mathsf G_u(z)$ with
$\mathsf G_u=\Gamma_u(y(z))$, $\Gamma_u\in1+Y\pao[[Y]]$; and $\mathsf G_u$ is $\partial$-integral.
If $\ell\ne d$, $\theta_u$ is $\partial$-integral.
\item[(ii)] Let $\Lambda=\Omega_0\auO_K$ and $\sigma\in\auAut(\C/K)$. There are $\Omega_\sigma\in\C^\times$ and an ideal $\mathfrak c$ of $\auO_K$ with
$\Lambda^\sigma=\Omega_\sigma\mathfrak a$, $\mathfrak a=\mathfrak c\aubar{\mathfrak c}^{-1}$, and a finite idele
$\beta_\sigma$ of $K$ generating $\mathfrak a\widehat\auO_K$, such that for every $\mathsf x\in K$, every $n\ge1$
with $n\mathsf x\in\auO_K$ and every $\mathsf x'\in K$ with $\mathsf x'-\beta_\sigma\mathsf x\in2n\mathfrak a\widehat\auO_K$,
$\sigma$ maps the Taylor series of $\theta_{\mathsf x\Omega_0}$ to that of the twisted translate
$\theta^\sigma_{\mathsf x'\Omega_\sigma}$ of the theta function of $\Lambda^\sigma$.
\end{itemize}
\item[(C4)] {\rm(Chowla--Selberg.)} There is an algebraic number $\mathsf u\ne0$, a unit at every prime, with $\operatorname{area}(\Lambda^\sigma)\,|\sigma\mathsf u|^2=2\boldsymbol\pi^2e^{c_{CS}}/
\sqrt d$ for all $\sigma\in\auAut(\C/K)$, where $c_{CS}=\gamma-\Rel\frac{L'}L(1,\chi_K)$.
\end{enumerate}
\end{proposition}
\begin{proof}
(C1). Let $\Lambda^\sigma$ be the lattice with invariants $\sigma g_2,\sigma g_3$ (it exists, as
$\sigma(g_2^3-27g_3^2)\ne0$). The number $e_1=\wp(t)$ is a root of $4X^3-g_2X-g_3$ with
$6e_1^2-g_2/2=\wp''(t)=2$; so $\sigma e_1=\wp_{\Lambda^\sigma}(t^\sigma)$ for a point $t^\sigma$ of
order two with $\wp''(t^\sigma)=2$, and $(\sigma a_1)^2=12\sigma e_1$. For $\alpha\in\auO_K$ there
is a rational function $R_\alpha$ with $\wp(\alpha z)=R_\alpha(\wp(z))$; it is determined by the
Laurent series of $\wp$, whose coefficients are polynomials in $g_2,g_3$ over $\Q$, so
$\sigma R_\alpha$ satisfies the same identity for $\Lambda^\sigma$ ($\sigma\alpha=\alpha$), and
$\alpha\Lambda^\sigma\subseteq\Lambda^\sigma$. The point $t^\sigma$ is admissible
(\S\ref{pf:ss-adm}): for $\beta\in\{\omega,\aubar\omega\}\subset\auO_K$, $\beta t\notin\Lambda$
means $\wp(\beta t)\ne\infty$, i.e.\ $R_\beta(e_1)\ne\infty$ ($\wp(t)=e_1$ is finite as $t\notin\Lambda$);
applying $\sigma$ to the numerator and denominator of $R_\beta$, this is equivalent to
$(\sigma R_\beta)(\sigma e_1)\ne\infty$, i.e.\ $\wp_{\Lambda^\sigma}(\beta t^\sigma)\ne\infty$, i.e.\
$\beta t^\sigma\notin\Lambda^\sigma$. By \eqref{pa:eq-curve}, $x$ and $y$ are the solutions
vanishing at $0$ of a differential system with coefficients in $\Z[a_1]$; so their Taylor
coefficients lie in $\Q[a_1]$ and are mapped by $\sigma$ to those of $\Lambda^\sigma$. By
\eqref{pa:eq-g2}, $\theta=\exp(-g)$ with $g''=x+s_2+a_1^2/12$, $g(0)=g'(0)=0$, and it remains to
show $\sigma s_2(\Lambda)=s_2(\Lambda^\sigma)$. This follows from the identity
\begin{equation}\label{pk:eq-s2}
\alpha(\alpha-\aubar\alpha)\,s_2=\sum_{0\ne u\in\alpha^{-1}\Lambda/\Lambda}\wp(u)\qquad(\alpha\in
\auO_K,\ \alpha\ne0),
\end{equation}
because the right side is minus the second coefficient of the monic denominator
$\prod_{u\ne0}(X-\wp(u))$ of $R_\alpha$, and $\alpha(\alpha-\aubar\alpha)\ne0$ for
$\alpha=\sqrt{-d}$. To prove \eqref{pk:eq-s2}: the function $\alpha\zeta(\alpha z)-\sum_{u\in
\alpha^{-1}\Lambda/\Lambda}\zeta(z+u)$ is entire, and changes under $z\mapsto z+\omega$ by
$\alpha\eta(\alpha\omega)-\auNm\alpha\cdot\eta(\omega)=s_2(\alpha^2-\alpha\aubar\alpha)\omega$; so
it is $s_2\alpha(\alpha-\aubar\alpha)z$ plus a constant. Differentiate and let $z\to0$; the terms
$-\alpha^2\wp(\alpha z)+\wp(z)$ tend to $0$. Uniqueness of $(\Lambda^\sigma,t^\sigma)$ is clear
from the series of $x$.

(C2). $a_1$ is an algebraic integer (Lemma~\ref{pa:model}), and $d/dz$ preserves $\Z[a_1][x,y]$ by
\eqref{pa:eq-curve}; so $x$, $y$ are $\partial$-integral at every prime. Since $g^{(n)}(0)=
x^{(n-2)}(0)$ for $n\ge3$ and $\theta=\exp(-g)$, it suffices to show $g''(0)=s_2+a_1^2/12\in\pao$.
By \eqref{pk:eq-s2} with $\alpha=\sqrt{-d}$, $\wp(u)=x(u-t)+a_1^2/12$ and $\auNm\alpha=d$,
\[
-2d\,g''(0)=\sum_{0\ne u\in\alpha^{-1}\Lambda/\Lambda}x(u-t)\ -\ \frac{d+1}{12}\,a_1^2 .
\]
The numbers $x(u-t)$ are the $x$-coordinates of the non-zero points of order $d$ of $\mathcal E$;
they are $\ell$-integral for $\ell\ne d$, because the points with non-integral $x$ form a group
isomorphic to the group of points of a formal group over the integers of a finite extension of
$\Q_\ell$, which has no torsion of order prime to $\ell$~\cite[Chapter~VII, Proposition~2.2, and Chapter~IV, Proposition~3.2(b)]{SilvermanAEC}.
They are equal in pairs ($x(-u-t)=x(u-t)$, $u\not\equiv-u$), and $24\mid d+1$ because $d\equiv
q_0/2-1\pmod{q_0}$ and $48\mid q_0$. So $-d\,g''(0)\in\pao$, and $g''(0)\in\pao$ for $\ell\ne d$.

(C3)(i). Put $R_u=\theta_u/\theta$. By \eqref{pa:eq-theta}, $R_u(z+\omega)=e^{2i(\boldsymbol\pi/A)
\Iml(\aubar\omega u)}R_u(z)$, and $\Iml(\aubar\omega\mathsf Nu)\in A\Z$; so $R_u^{\mathsf N}$ is
$\Lambda$-periodic. Its only poles are at $z\in t+\Lambda$, the origin $O$ of $\mathcal E$, of order $\mathsf N$, and its only zeros are at $z\in t-u+\Lambda$, a point $Q_1$ of $\mathcal E$
whose order divides $\mathsf N$, and they have order $\mathsf N$ (if $u\in\Lambda$, $R_u$ is
constant). Thus $R_u^{\mathsf N}=\theta_u(0)^{\mathsf N}\mathsf f(x,y)/
\mathsf f(0,0)$ for a polynomial $\mathsf f$ with divisor $\mathsf N[Q_1]-\mathsf N[O]$; here
$\mathsf f(0,0)\ne0$ because $z=0$ is the point $P_0=(0,0)\ne Q_1$ ($u\notin t+\Lambda$). Scale
$\mathsf f$ so that its coefficients are integral and not all in $\pam$. The curve has good
reduction at $\ell$ (Lemma~\ref{pa:model}), so the divisor of the reduction $\tilde{\mathsf f}\ne0$
is the reduction $\mathsf N[\tilde Q_1]-\mathsf N[\tilde O]$ of the divisor of $\mathsf f$. The
point $Q_1-P_0$ is not zero and has order prime to $\ell$ (as $\ell\nmid2\mathsf N$), so it does
not reduce to the origin; hence $\tilde Q_1\ne\tilde P_0$ and $\mathsf f(0,0)$ is a unit. As $x\in
y^2\pao[[y]]$, $\mathsf f(x,y)/\mathsf f(0,0)\in1+y\pao[[y]]$, and $\Gamma_u$ is its $\mathsf N$-th
root with constant term $1$, which has coefficients in $\pao$ because $\mathsf N$ is a unit at
$\ell$. It is $\partial$-integral in $z$ because $y$ is (C2). Finally $\theta_u(0)\in\pao$ by Proposition~\ref{pf:values}.
The last assertion of (i) follows from (C2).

(C4). Let $h$ be the class number of $K$; it is odd, so every ideal class contains an ideal
$\mathfrak c\aubar{\mathfrak c}^{-1}$. Choose fractional ideals $\mathfrak c_1,\dots,\mathfrak c_h$
representing the classes with $\prod_i\mathfrak c_i=(1)$, numbers $\mu_i\in\C^\times$ such that
$\Lambda_i=\mu_i\mathfrak c_i\Lambda$ is the lattice of a normalised lattice $(\Lambda_i,t_i,a_{1,i})$,
and let $\mathsf u$ be a root of $\mathsf u^{12h}=\prod_i\mu_i^{12}\Delta(\Lambda_i)$, where
$\Delta=g_2^3-27g_3^2=a_1^4-64$. The numbers $\mu_i^2$ are algebraic. For $\sigma\in\auAut(\C/K)$,
$\Lambda_i^\sigma=\sigma(\mu_i)\mathfrak c_i\Lambda^\sigma$: if $\lambda\in\mu_i\mathfrak c_i$,
multiplication by $\lambda$ maps $\Lambda$ into $\Lambda_i$, so $\wp_{\Lambda_i}(\lambda z)$ is a
rational function of $\wp_\Lambda(z)$ determined by the Laurent series, and conjugating as in
(C1) gives $\sigma(\lambda)\Lambda^\sigma\subseteq\Lambda_i^\sigma$; the same for the inverse
ideal gives equality. The lattices $\Lambda_i^\sigma$ represent the $h$ classes of lattices with
multiplication by $\auO_K$. For such a lattice $\Z\omega_1+\Z\omega_2$, $\tau=\omega_2/\omega_1$,
one has $\operatorname{area}\cdot|\Delta|^{1/6}=4\boldsymbol\pi^2\,\Iml\tau\,|\eta(\tau)|^4$, and
Kronecker's limit formula together with $\zeta_K=\zeta\cdot L(\cdot,\chi_K)$ and the class number
formula gives the following: by \cite[\S12, the display preceding~(2)]{ChowlaSelberg}, with
$\tau_j=(b_j+i\sqrt d)/(2a_j)$ running over the reduced forms of discriminant $-d$ and $w=2$,
$\sum_j\log|\Delta(\tau_j)|=6\bigl(h\gamma+\log(\prod_ja_j/d^h)\bigr)-\frac{3w}{\boldsymbol\pi}\sqrt d\,
L'(1,\chi_K)$, where $\Delta=\eta^{24}$; since $\Iml\tau_j=\sqrt d/(2a_j)$ and
$L(1,\chi_K)=2\boldsymbol\pi h/(w\sqrt d)$, dividing by $6h$ and adding
$\frac1h\sum_j\log\Iml\tau_j=\log(\sqrt d/2)-\frac1h\log\prod_ja_j$ gives $\gamma-\frac{w\sqrt d}{2\boldsymbol\pi h}L'(1,\chi_K)-\log d+\log(\sqrt d/2)=\gamma-\frac{L'}{L}(1,\chi_K)-\log(2\sqrt d)$, the terms in $a_j$ cancelling; here $\log|\eta(\tau)|^4=\frac16\log|\Delta(\tau)|$, and $\Iml\tau\,|\eta(\tau)|^4$ is invariant under $\mathrm{SL}_2(\Z)$ and under $\tau\mapsto-\aubar\tau$, so it depends only on the homothety class of $\Z+\Z\tau$, and the reduced forms $[a_j,\pm b_j,c_j]$ of discriminant $-d$ correspond bijectively to the $h$ homothety classes of lattices with multiplication by $\auO_K=\Z[\frac{1+\sqrt{-d}}2]$ via $\tau_j\mapsto\Z+\Z\tau_j$. This proves (for Kronecker's limit formula see also
\cite[Chapter~20, \S4]{LangEF})
\[
\frac1h\sum_{\text{classes}}\log\bigl(\Iml\tau\,|\eta(\tau)|^4\bigr)=c_{CS}-\log(2\sqrt d).
\]
Hence $\prod_i\operatorname{area}(\Lambda_i^\sigma)|\Delta(\Lambda_i^\sigma)|^{1/6}=(2\boldsymbol\pi^2
e^{c_{CS}}/\sqrt d)^h$; the left side is $\prod_i|\sigma\mu_i|^2\auNm(\mathfrak c_i)\operatorname{area}
(\Lambda^\sigma)|\sigma\Delta(\Lambda_i)|^{1/6}=(\operatorname{area}(\Lambda^\sigma)|\sigma\mathsf
u|^2)^h$. For the unit property let $\ell$ be any prime. The $\Delta(\Lambda_i)$ are units
(Lemmas~\ref{pa:model}, \ref{pf:two}), the $t_i$ being admissible. If $\lambda\Lambda\subseteq\Lambda'$ for normalised lattices, then
$z\mapsto\lambda z$ is a homomorphism between elliptic curves with good reduction which multiplies
the invariant differentials $dx/(2y+a_1x)=dz$ by $\lambda$; these differentials generate the
modules of invariant differentials of the smooth models, and the homomorphism extends to these
models, which are N\'eron models~\cite[\S1.2, Definition~1, p.~12, and Proposition~8, p.~15]{BLR} (\cite[IV, Corollary~6.3, p.~329, and the N\'eron mapping property, p.~319]{SilvermanAT}), so $\lambda\in\pao$. Applying this to
$\lambda\in\mu_i\mathfrak c_i$ and to $\lambda\in\mu_i^{-1}\mathfrak c_i^{-1}$ gives
$|\mu_i|_\ell=|\mathfrak c_i|_\ell^{-1}$ in the obvious sense, and $\prod_i\mathfrak c_i=(1)$ gives
$\prod_i|\mu_i|_\ell=1$.
\end{proof}

(C3)(ii) is Theorem~\ref{pf:recip}; that $\mathsf u$ is a unit at $2$ is Corollary~\ref{pf:CS}.

\subsubsection{Theta products, the integral vacuum, and conjugate jets: Lemma~\ref{ps:conj}}
\label{pk:ss-conj}

Let $W$ be a standard space $E^{\mathsf C}$ with the form $\sum\aubar x_ix_i$ ($\mathsf C$ a finite
set; for $W=\auHom(X,X')$, $\mathsf C=[n']\times[n]$), $L=\alpha\widehat\auO_E^{\,\mathsf C}$,
$\Lambda_L=W(E)\cap L=\mathfrak a_0^{\mathsf C}$. A point $\gamma\in W(E)$ has the coordinate
$\gamma^\natural=(c_\psi\,\tau(\gamma_i))_{\tau,i}$. Note that the argument $z$ of $\Theta_f$ in
\eqref{pk:eq-theta} is a point of $W\otimes\R$ in these coordinates (it enters through
$\psrho_\infty(z)$), while the variable of the Fock space is the conjugate coordinate.

\begin{lemma}[Dictionary]\label{pk:dict}
The map $f\mapsto\Theta_f$ is a bijection from $\psHS_f(W)$ onto the space of entire functions
$\Theta$ on $\C^{\auPl\times\mathsf C}$ for which there are an integer $\mathsf N\ge1$ and a function
$\epsilon$ on $\mathsf N\Lambda_L$, with values in the roots of unity, with
\[
\Theta(z+\lambda^\natural)=\epsilon(\lambda)\,e^{\langle z,\lambda^\natural\rangle+|\lambda^\natural|^2
/2}\,\Theta(z)\qquad(\lambda\in\mathsf N\Lambda_L),\qquad\langle z,w\rangle=\sum\aubar w z .
\]
For $\gamma\in W(E)$, $\Theta_{\psrho_f(\gamma)f}(z)$ is a root of unity times
$e^{\langle z,\gamma^\natural\rangle-|\gamma^\natural|^2/2}\,\Theta_f(z-\gamma^\natural)$.
\end{lemma}
\begin{proof}
For $\gamma\in W(E)$, $\theta\circ\psrho(-\gamma)=\theta$ and $\psrho(-\gamma)=\psrho_\infty(
-\gamma^\natural)\psrho_f(-\gamma)$. So $\theta(\psrho_\infty(z)(1\otimes\psrho_f(\gamma)f))$ is a
root of unity (a value of $\beta_f$) times $\beta_\infty(-\gamma^\natural,z)\,\theta(\psrho_\infty(z-
\gamma^\natural)(1\otimes f))$, and $\beta_\infty(-x,z)e^{|z|^2/2-|z-x|^2/2}=e^{\langle z,x\rangle-
|x|^2/2}$; with \eqref{pk:eq-theta} this is the last formula. If $f$ is fixed up to characters by
$\psrho_f(l)$, $l\in\mathsf NL$ (every $f$ is, for some $\mathsf N$, by (W4)), the formula with
$\gamma=-\lambda$ gives the functional equation, with $\epsilon(\lambda)$ the inverse of the
eigenvalue of $\psrho_f(\lambda)$ on $f$ (the root of unity $\beta_f(\lambda,-\lambda)$ is $1$,
because $\pktr_{E/E^+}(\jmath\langle\lambda,\lambda\rangle)=0$); this eigenvalue is a root of unity,
the action of $\psrho_f$ on $\mathsf NL$ being smooth and $\beta_f$ having values in the roots of
unity. The map is injective: if $\Theta_f=0$ then, by
the last formula and the density of $W(E)$ in $W(\auA_f)$, $\theta(\psrho(w)(1\otimes f))=0$ for
all $w\in W(\auA)$, so $1\otimes f$ lies in the kernel of $F\mapsto\theta(\psrho(\cdot)F)$, a
proper closed invariant subspace of an irreducible representation. It is surjective: for fixed
$\mathsf N$ the functions with a given $\epsilon$ form a space of dimension
$[\mathsf N^{-1}\Lambda_L^{*}:\mathsf N\Lambda_L]^{1/2}=\mathsf N^{2f\#\mathsf C}$, the hermitian form
$|z|^2$ being unimodular on $\Lambda_L^\natural$ by the self-duality of
$L$~\cite[\S3, Proposition, p.~26]{MumfordAV}; the admissible $\epsilon$ with values in the roots
of unity differ from one another by characters of the finitely generated group
$\mathsf N\Lambda_L$ with values in the roots of unity, that is, by the characters of
$\mathsf N\Lambda_L$ of finite order (without the condition on the values the statement would
fail: the functions $z\mapsto e^{-\langle z,c\rangle}\Theta_f(z+c)$, $c\in\C^{\auPl\times\mathsf C}$,
satisfy the functional equation with $\epsilon$ multiplied by the unitary character
$\lambda\mapsto e^{2i\,\Iml\langle c,\lambda^\natural\rangle}$, which is arbitrary);
these are produced by the translations
$\psrho_f(x)$, $x\in W(\auA_f)$; and the space of $f$ on which $\psrho_f(\mathsf NL)$ acts by a
given character has the same dimension, by (W4).
\end{proof}

\emph{The theta product.} Let $\Lambda=\Omega_0\auO_K$, $\theta$ its theta function, $u_k=
\zeta_{q_1}^k/q_1$ $(0\le k<q_1)$ and $u_k^\dagger(z_i)=\sum_\tau\aubar{\tau(u_k)}z_{\tau,i}$. Since
$\sum_k\tau(\zeta^k)\aubar{\tau'(\zeta^k)}=q_1\delta_{\tau\tau'}$, the family $(u_k)$ with
multiplicities $q_1$ is a resolution of the identity \eqref{pa:eq-res}. Put
\begin{equation}\label{pk:eq-product}
\Pi(z)=\prod_{i\in\mathsf C}\prod_{k<q_1}\theta\bigl(\tfrac{\Omega_0}{c_\psi}\,u_k^\dagger(z_i)\bigr)^{q_1}.
\end{equation}
For $\lambda\in q_1\Lambda_L$ the arguments change by $\Omega_0\,\pktr_{E/K}(\aubar{u_k}\lambda_i)\in
\Omega_0\auO_K$ (as $\mathfrak a_0\subseteq\mathfrak d_{E/K}^{-1}$), and by \eqref{pa:eq-theta},
\eqref{pa:eq-res} and $A=|\Omega_0|^2\sqrt d/2$, $\Pi$ satisfies the functional equation of
Lemma~\ref{pk:dict} with $\mathsf N=q_1$, $\epsilon=\pm1$. So $\Pi=\Theta_{F_0}$ for a unique
$F_0\in\psHS_f(W)$, and $\theta(1\otimes F_0)=\Pi(0)=1$. By Lemma~\ref{pk:dict}, for $\gamma\in
W(E)$,
\begin{equation}\label{pk:eq-transl}
\Theta_{\psrho_f(\gamma)F_0}(z)=\zeta_\gamma\prod_{i,k}\theta_{\mathsf x_{k,i}\Omega_0}\bigl(\tfrac
{\Omega_0}{c_\psi}u_k^\dagger(z_i)\bigr)^{q_1},\qquad\mathsf x_{k,i}=-\pktr_{E/K}(\aubar{u_k}\gamma_i)
\in K,
\end{equation}
with a root of unity $\zeta_\gamma$: by \eqref{pa:eq-res}, $\langle z,x\rangle=\sum_{i,k}q_1\,H(a_{k,i}
(x),a_{k,i}(z))$ for $a_{k,i}(z)=\frac{\Omega_0}{c_\psi}u_k^\dagger(z_i)$, since
$\boldsymbol\pi|\Omega_0|^2/(Ac_\psi^2)=1$, and $a_{k,i}(\gamma^\natural)=-\mathsf x_{k,i}\Omega_0$. The vectors $\psrho_f(x)F_0$, $x\in W(\auA_f)$, span
$\psHS_f(W)$ (irreducibility), and each is a root of unity times $\psrho_f(\gamma)F_0$ for a
rational $\gamma$ close to $x$.

\emph{The integral vacuum.} The character $\epsilon_L$ is trivial on $q_1L$. On the space of the vectors fixed by
$\psrho(q_1L)$ the operator $\sum_{l\in L/q_1L}\epsilon_L(l)^{-1}\psrho(l)$ is a multiple of the
projector onto the line of $\psvac_L$, and this space is spanned by the translates of $F_0$ that
it contains, which are of the form $\psrho(b)F_0$ with $q_1b\in L$ (the function $\epsilon$ of
$F_0$ is $\epsilon_L$ on $q_1\Lambda_L$, Lemma~\ref{pf:semi}). So there are $b\in q_1^{-1}L$ and roots of unity $\zeta_l$ such that
\begin{equation}\label{pk:eq-vac}
\psvac^\circ:=\sum_{l\in L/q_1L}\zeta_l\,\psrho(l)\psrho(b)F_0\ \in\C^\times\psvac_L
\end{equation}
(some such translate has a non-zero projection). For
$\mathsf C$ a point this defines the integral vacuum of rank one and $c_{\mathrm{vac}}=\psvac^\circ
/\psvac$; in general $\Pi$, $F_0$ and $\psvac_L$ are tensor products over $\mathsf C$, and
$\psvac^\circ=\bigotimes_{\mathsf C}\psvac^\circ_1=c_{\mathrm{vac}}^{\#\mathsf C}\psvac$. It is a sum
of $q_1^{2f\#\mathsf C}$ translates of $F_0$ with coefficients roots of unity.

\begin{proposition}[Lemma~\ref{ps:conj}]\label{pk:conj}
Lemma~\ref{ps:conj} holds, with
$t_\sigma=\sigma(\varpi_{\mathrm{vac}}/\mathsf u)\,c_\psi/\Omega_\sigma$, $\varpi_\sigma$ the
permutation $\tau\mapsto\sigma\circ\tau$ of $\auPl$, and $C_0=C_0(q_0,d)$ given at the end of the proof.
\end{proposition}
\begin{proof}
\emph{(i).} Let $F_0^\sigma$ be the vector attached by \eqref{pk:eq-product} to the lattice
$\Lambda^\sigma=\Omega_\sigma\mathfrak c\aubar{\mathfrak c}^{-1}$ and the number $\Omega_\sigma$; the
functional equation holds with $\mathsf N=q_1\auNm\mathfrak c$, because $\auNm\mathfrak c\in
\mathfrak c\aubar{\mathfrak c}^{-1}$ and $\operatorname{area}(\Lambda^\sigma)=|\Omega_\sigma|^2\sqrt
d/2$. At the scale $s_0$ the arguments of $\theta$ in \eqref{pk:eq-transl} are
$(\varpi_{\mathrm{vac}}/\mathsf u)\,u_k^\dagger(z_i)$, and at the scale $t_\sigma$ those of the
product for $\Lambda^\sigma$ are $\sigma(\varpi_{\mathrm{vac}}/\mathsf u)\,u_k^\dagger(z_i)$. Since
$\sigma$ fixes $K$ and commutes with the complex conjugation of the CM field $E$,
$\sigma(\aubar{\tau(u_k)})=\aubar{(\sigma\tau)(u_k)}$; so $\sigma$ maps the linear form
$(\varpi_{\mathrm{vac}}/\mathsf u)u_k^\dagger$ to $\sigma(\varpi_{\mathrm{vac}}/\mathsf u)u_k^\dagger$
with the variables permuted by $\varpi_\sigma$. Let $\mathsf B$ be $\beta_\sigma$ regarded in
$\auA_{E,f}$; it is an idele and $\mathsf B\aubar{\mathsf B}$ is a unit (Theorem~\ref{pf:recip}(a)), so
$\mathsf BL$ is self-dual; for $\gamma'\in W(E)$ close to $\mathsf B\gamma$ the numbers
$-\pktr_{E/K}(\aubar{u_k}\gamma'_i)$ are close to $\beta_\sigma\mathsf x_{k,i}$. By (C1) and
(C3)(ii), for $\gamma\in W(E)$ and $\gamma'$ close enough to $\mathsf B\gamma$, the Taylor series
of $\Theta_{\psrho_f(\gamma')F_0^\sigma}(t_\sigma\,\cdot)$ is $\zeta'$ times the image under
$\sigma$ and $\varpi_\sigma$ of that of $\Theta_{\psrho_f(\gamma)F_0}(s_0\,\cdot)$, with $\zeta'$ a
root of unity. Define $\Lambda_\sigma$ on the spanning family by $\Lambda_\sigma(\sum_ic_i
\psrho_f(\gamma_i)F_0)=\sum_i\sigma(c_i)\zeta_i'^{-1}\psrho_f(\gamma'_i)F_0^\sigma$. It is well
defined and satisfies (i), because $f\mapsto$ (Taylor series of $\Theta_f$) is injective
(Lemma~\ref{pk:dict}) and both sides of (i) are $\sigma$-semilinear in $f$. Moreover
$|t_\sigma|^2=1/\Omega$ by \eqref{pk:eq-s0}, (C4) and the area of $\Lambda^\sigma$.

\emph{Commutation rules.} Comparing Taylor series with Lemma~\ref{pk:dict},
$\Lambda_\sigma\psrho_f(x)=\zeta_x\,\psrho_f(\mathsf Bx)\Lambda_\sigma$ with roots of unity
$\zeta_x$ $(x\in W(\auA_f))$; in particular $\Lambda_\sigma$ is injective and maps the line of
$\psvac_L$ to the line of a character of the self-dual lattice $L'=\mathsf BL$. For
$g\in\auU(W)(\auA_f)$ the linear operator $\mathsf D=\omega(g)^{-1}\Lambda_\sigma\omega(g)
\Lambda_\sigma^{-1}$ on $\Lambda_\sigma\psHS_f$ satisfies $\mathsf D\psrho_f(y)\mathsf D^{-1}\in
\C^\times\psrho_f(y)$ for all $y$ by (W1) ($\mathsf B$ is a scalar and commutes with $g$); by the
irreducibility of the Heisenberg representation, $\mathsf D=\mathsf c_g\psrho_f(y_g)$ for some
$y_g$. Thus $\Lambda_\sigma\omega(g)\psvac^\circ=\mathsf c_g\,\omega(g)\psrho_f(y_g)\Lambda_\sigma
\psvac^\circ$, and $|\mathsf c_g|=1$: by Lemma~\ref{pf:delta}, $\omega(g)\psvac_L=\delta_g\sum_r\zeta_r\psrho^\circ(r)
\psvac_L$, a sum over the $\mathsf I_g=[gL:L\cap gL]$ classes $r$, with roots of unity $\zeta_r$ and
$|\sigma\delta_g|=|\delta_g|=\mathsf I_g^{-1/2}$ (part (i) of the lemma); so the left
side is a sum of $\mathsf I_g$ orthogonal vectors of norm $\mathsf I_g^{-1/2}\|\Lambda_\sigma
\psvac^\circ\|$, and the right side has norm $|\mathsf c_g|\,\|\Lambda_\sigma\psvac^\circ\|$.

\emph{(ii).} We normalise $\theta$ by $\int_{W(E)\backslash W(\auA)}|\theta(\psrho(w)F)|^2dw=\|F
\|^2$ (self-dual measure; this is the normalisation of the vacuum in \S\ref{ps:ss-weil}). Then
$\sum_{\gamma\in W(E)}\langle F',\psrho(\gamma)F\rangle=\theta(F)\aubar{\theta(F')}$ for
$F,F'\in1\otimes\psHS_f$: the left side converges, is invariant under $W(E)$ in $F$ and in $F'$,
hence is a multiple of the right side by (W1), and the multiple is $1$ by the normalisation. With
$F'=F=\psrho(w)(1\otimes\psvac_{L''})$, $L''$ self-dual, this gives
\begin{equation}\label{pk:eq-sup}
\sup_{w\in W(\auA)}\bigl|\theta\bigl(\psrho(w)(1\otimes\psvac_{L''})\bigr)\bigr|^2\le\sum_{\gamma\in
W(E)\cap L''}e^{-|\gamma^\natural|^2/2},
\end{equation}
because $|\langle\psrho(\gamma)\psrho(w)\mathbf v'',\psrho(w)\mathbf v''\rangle|=|\langle\psrho(\gamma)
\mathbf v'',\mathbf v''\rangle|=e^{-|\gamma^\natural|^2/2}\mathbf 1_{L''}(\gamma)$ for $\mathbf v''=1\otimes
\psvac_{L''}$. For $\gamma\ne0$ in $W(E)\cap L''$, $\nu=\langle\gamma,\gamma\rangle$ is a totally
positive element of $\mathfrak d_{E/K}^{-1}$ (self-duality), so $\sum_\tau\tau(\nu)\ge
f\auNm(\mathfrak d_{E/K})^{-1/(2f)}$ by the inequality of the means ($\nu\in E^+$ is totally
positive and $N_{E^+/\Q}(\nu)\ge\auNm(\mathfrak d_{E/K})^{-1/2}$). Here $E=K\otimes\Q(\zeta_{q_1})$
with $f=\varphi(q_1)$, and $d\nmid q_1$ (every prime dividing $q_1$ divides $(1+d)/4$), so the
discriminants of $K$ and $\Q(\zeta_{q_1})$ are coprime, the discriminant of $E$ is
$d^f\operatorname{disc}(\Q(\zeta_{q_1}))^2$, and $\auNm(\mathfrak d_{E/K})^{1/(2f)}$ is the root
discriminant $q_1\prod_{p\mid q_1}p^{-1/(p-1)}$ of $\Q(\zeta_{q_1})$. Since $q_1$ and $q_0$ have the same
prime divisors ($q_0\mid q_1$, and {\cite[Lemma~13.4(S7)]{P46}}), this gives
$f\auNm(\mathfrak d_{E/K})^{-1/(2f)}=\prod_{p\mid q_1}(1-\tfrac1p)p^{1/(p-1)}=\mathsf m:=\prod_{p\mid q_0}
(1-\tfrac1p)p^{1/(p-1)}$; hence the
points $\gamma^\natural$ are separated by $\delta=c_\psi\sqrt{\mathsf m}$, the number of them with
$k\delta\le|\gamma^\natural|<(k+1)\delta$ is at most $(2k+3)^{\mathsf D}$, $\mathsf D=2f\#\mathsf
C\le2q_1\#\mathsf C$, and the right side of \eqref{pk:eq-sup} is at most $\sum_{k\ge0}(2k+3)^{
\mathsf D}e^{-\delta^2k^2/2}\le(1+4/\delta)\exp(2\mathsf D+\mathsf D^2/\delta^2)$.

Let $\mathsf M_d=\max_{\Lambda'}\sup_u|\theta^{\Lambda'}_u(0)|$, the maximum over the finitely many
normalised lattices $\Lambda'$ with multiplication by $\auO_K$; it is finite because
$|\theta_u(0)|$ is continuous and periodic in $u$. By \eqref{pk:eq-vac}, (i) and
\eqref{pk:eq-transl} for $\Lambda^\sigma$, $\Lambda_\sigma\psvac^\circ$ is a sum of $q_1^{2f\#
\mathsf C}$ translates of $F_0^\sigma$ with coefficients roots of unity, and every value at $0$ of
the theta function of a translate of $F_0^\sigma$ is a product of $q_1^2\#\mathsf C$ numbers of
absolute value $\le\mathsf M_d$. By Lemma~\ref{pk:dict} and continuity,
$\sup_w|\theta(\psrho(w)(1\otimes\Lambda_\sigma\psvac^\circ))|\le(q_1^{2f}\mathsf M_d^{q_1^2})^{\#
\mathsf C}$, and by the normalisation of $\theta$ the norm $\|\Lambda_\sigma\psvac^\circ\|$ is at
most this number. Now let $f=\psrho^\circ(x)\omega(g)\psvac^\circ$. By the commutation rules,
$\Lambda_\sigma f$ is a number of absolute value $\|\Lambda_\sigma\psvac^\circ\|$ times
$\psrho_f(x')\omega(g)$ applied to a unit vector on the line of a character of $L'$, that is,
times a translate of a unit vector on the line of $\psvac_{gL'}$; so by \eqref{pk:eq-sup}, for
all $y$,
\[
\bigl|\theta\bigl(\psrho_\infty(y)(1\otimes\Lambda_\sigma f)\bigr)\bigr|\le(q_1^{2f}\mathsf
M_d^{q_1^2})^{\#\mathsf C}\,(1+4/\delta)^{1/2}\exp(\mathsf D+\mathsf D^2/2\delta^2)\le\exp\bigl(C_0
(q_1\#\mathsf C+1)^2\bigr),
\]
with $C_0=4+\log^+\mathsf M_d+\sqrt d/(\boldsymbol\pi\mathsf m)+\log(1+4/\delta)$, which depends
on $q_0$ and $d$ only.
\end{proof}

\subsubsection{The kernels: Lemma~\ref{pa:kernel}}\label{pk:ss-kernel}

\begin{proposition}[Lemma~\ref{pa:kernel}(K0)]\label{pk:K0}
With the lifts defined by \eqref{pk:eq-lift}, {\rm(K0)} holds.
\end{proposition}
\begin{proof}
For a pure tensor $\Phi=\bigotimes_j\Phi_j$ and $\phi\in\auS$, \eqref{pk:eq-lift} gives
$(C_{\vartheta,\Phi}\phi)(g_m)=\int\cdots\int\mathsf F_\Phi(g_0,\dots,g_{m-1})\,dg_0\cdots dg_{m-1}$
over $\prod_j\Gamma_j\backslash\auU(X_j)(\auA_f)$, with $\mathsf F_\Phi=\theta_{m-1}(\Phi_{m-1};
g_{m-1},g_m)\cdots\theta_0(\Phi_0;g_0,g_1)\phi(g_0)$; both sides are linear in $\Phi$. By
\eqref{ps:eq-kernel}, $\theta_j(\Phi_j;g_jb_j,g_{j+1}b_{j+1})=\theta_j(\omega_{\mu_{j+1}}(b_j,
b_{j+1})\Phi_j;g_j,g_{j+1})$. Hence, for the tensor $\Phi=\sum_x\Psi'(x)\phi^\circ_x$, which is
fixed by $\bigotimes_j\omega_{\mu_{j+1}}(b_j,b_{j+1})$ for $b_j\in J_{X_j}$, $b_m=1$ ((C4) of
\S\ref{ps:ss-glue}; the tensor is the same for $\Psi$ and $\Psi'$), and for $\phi$ of level
$J_{X_0}$, $\mathsf F_\Phi$ is invariant under $g_j\mapsto g_jb_j$ in all variables independently.
The levels are free and have volume $1$, so the iterated integral is the sum over the
representatives of the classes, which is (K0).
\end{proof}

A \emph{standard datum at $\ell$} on $W=\auHom(X_j,X_{j+1})$ is a combination, with coefficients
in $\pao\cap\C$, of the vectors $\psrho^\circ(x)\psvac^\circ$ with $\aubar{\pi_M}x$ integral at
the places above $\ell$.

\begin{proposition}\label{pk:std}
Let $\ell\ne N$ be a prime and $\C\to\C_\ell$.
\begin{enumerate}
\item[(S1)] If $g\in\auU(W)(\auA_f)$ is integral at the places above $\ell$, then $\omega(g)$ maps
standard data at $\ell$ to standard data at $\ell$.
\item[(S2)] If $\ell\ne d$ and $f$ is a standard datum at $\ell$, then
$\Theta_f(\pi_Ms_0\,\cdot)$ is $\partial$-integral.
\item[(S3)] If $\ell=p$ does not split in $K$ (so $p\nmid2q_1M$) and $f$ is a standard datum at
$p$, then $\Theta_f(\pi_Ms_0\,z)=\prod_{i,k}\theta(\mu_{k,i}(z))^{q_1}\cdot G(z)$, where
$\mu_{k,i}=\mathsf c\,u_k^\dagger(z_i)$ with $\mathsf c=\pi_M\varpi_{\mathrm{vac}}/\mathsf u\in
\pao$, and $G$ is an integral power series in the series $y(\mu_{k,i}(z))$; so it is one of the
series of Proposition~\ref{pa:product}, for the resolution of the identity $(u_ke_i)_{k,i}$ of
$E^{\mathsf C}$ with multiplicities $q_1$.
\end{enumerate}
\end{proposition}
\begin{proof}
(S1). By (W1), $\omega(g)\psrho^\circ(x)\psvac^\circ$ is a root of unity times $\psrho^\circ(gx)
\omega(g)\psvac^\circ$, and $\aubar{\pi_M}gx$ is integral above $\ell$. By Lemma~\ref{pf:delta},
$\omega(g)\psvac_L=\delta_g\sum_{r\in gL/(L\cap gL)}\zeta_r\psrho^\circ(r)\psvac_L$ with roots of
unity $\zeta_r$, and, since $g$ stabilises $L$ at the places above $\ell$, $\delta_g$ is a unit above
$\ell$ and the $r$ can be taken integral above $\ell$ (part (ii) of the lemma).

(S2), (S3). By \eqref{pk:eq-vac} and \eqref{pk:eq-transl}, and since a translation by $x$ agrees
up to a root of unity with the translation by a rational $\gamma$ close to $x$, $f$ is a
combination with coefficients in $\pao$ of vectors $\psrho_f(\gamma)F_0$ with $\gamma\in W(E)$,
$2q_1\aubar{\pi_M}\gamma$ integral above $\ell$, and
\[
\Theta_{\psrho_f(\gamma)F_0}(\pi_Ms_0z)=\zeta_\gamma\prod_{i,k}\theta_{\mathsf x_{k,i}\Omega_0}\bigl(
\mathsf c\,u_k^\dagger(z_i)\bigr)^{q_1},\qquad\mathsf c=\tfrac{\Omega_0}{c_\psi}\pi_Ms_0=
\frac{\pi_M\varpi_{\mathrm{vac}}}{\mathsf u}.
\]
The number $\mathsf c$ is integral at $\ell\nmid2q_1$: $\mathsf u$ is a unit (C4), and
$\varpi_{\mathrm{vac}}^{2f}=\upsilon(\alpha)^{-2}\auNm(\mathfrak d_{E/K})^{1/2}$ is a unit at the
primes not dividing $q_1$. The coefficients $\aubar{\tau(u_k)}$ are integral at $\ell\nmid q_1$,
and the points $\mathsf x_{k,i}\Omega_0$ are torsion points of order prime to $\ell$, as
$\ell\nmid2q_1M$. Let $u$ be one of them. If $u\notin t+\Lambda$, (C3)(i) gives
$\theta_u(\mu)=\theta_u(0)\,\theta(\mu)\,\Gamma_u(y(\mu))$ with $\theta_u(0)\in\pao$,
$\Gamma_u\in\pao[[Y]]$. If $u\in t+\Lambda$, then $\theta_u$ is a constant multiple of $\theta_t$
and $(\theta_t/\theta)^2$ is $\Lambda$-periodic with divisor $2[P_0]-2[O]$, hence equal to $\mathsf
c_tx$; as $x=y^2\mathsf e(y)$ with $\mathsf e\in1+y\pao[[y]]$, $\theta_t=\mathsf c_t^{1/2}\,\theta
\cdot y\,\mathsf e(y)^{1/2}$, and $\mathsf c_t^{1/2}=\theta_t'(0)$ is a unit (this is shown in the proof of
Corollary~\ref{pf:K1}(c) from Lemma~\ref{pf:chi}(c) and Proposition~\ref{pf:values} alone, without the present proposition). In both cases $\theta_u(\mu)=\theta(\mu)\cdot\Gamma(y(\mu))$ with
$\Gamma\in\pao[[Y]]$. Hence $\Theta_{\psrho_f(\gamma)F_0}(\pi_Ms_0z)=\prod_{i,k}\theta(\mu_{k,i}(z))^{
q_1}\cdot G_\gamma(z)$ with $G_\gamma$ an integral power series in the $y(\mu_{k,i}(z))$; the first
factor does not depend on $\gamma$, so the same holds for $f$, with $G=\sum_\gamma c_\gamma
G_\gamma$. This is (S3). For (S2): $\theta$ and $y$ are $\partial$-integral at $\ell\ne d$ by (C2),
the $\mu_{k,i}$ have coefficients in $\pao$, and compositions and products preserve
$\partial$-integrality. This argument gives (S2) for $\ell\nmid2q_1Md$. For $\ell\mid2q_1M$, (S2) is Theorem~\ref{pf:std} with
$\delta=\aubar{\pi_M}$.
\end{proof}


\begin{proposition}[Lemma~\ref{pa:kernel}(K1), (K2)]\label{pk:K12}
{\rm(K1)} and {\rm(K2)} of Lemma~\ref{pa:kernel} hold.
The integrality of the jets in (K1) holds for every $\ell\ne d$, by (S2) and Theorem~\ref{pf:std}.
\end{proposition}
\begin{proof}
Let $x$ be in the support of $\Psi'$; then $\aubar{\pi_M}x_j$ is integral at the places above
$\ell$ ($\ell\nmid N\auNm Q_\vartheta$), so $\phi^\circ_{j,x_j}$ is a standard datum at $\ell$. Let
$g\in\auU(X_j)(\auA_f)$, $g'\in\auU(X_{j+1})(\auA_f)$. Unitary groups are rational varieties (Cayley transform) and therefore satisfy weak
approximation~\cite[\S7.1, Propositions~7.3 and~7.4]{PlatonovRapinchuk}, so there are $\gamma\in\auU(X_j)(E^+)$, $\gamma'\in\auU(X_{j+1})(E^+)$ with $\mathsf k=\gamma^{-1}g$,
$\mathsf k'=\gamma'^{-1}g'$ so close to $1$ at the places above $\ell$ that they lie in the
stabilisers of the standard lattices there. By the equivariance of the kernels for the rational
points (which follows from \eqref{pk:eq-schuract} and $\theta\circ\omega(\gamma,\gamma')=\theta$),
$\theta_j(\Phi;g,g')=\gamma'\cdot\theta_j(\Phi;\mathsf k,\mathsf k')\cdot\gamma^{-1}$, and
therefore
\[
g'^{-1}_\ell\circ\theta_j(\Phi;g,g')\circ g_\ell=\mathsf k'^{-1}_\ell\circ\theta_j(\Phi;\mathsf k,
\mathsf k')\circ\mathsf k_\ell ,
\]
where $\mathsf k_\ell,\mathsf k'_\ell$ act on $W_\lambda(\pao)$ by matrices invertible over
$\pao$. Thus the two invertible matrices of (K1) are $\rho(\mathsf k'_\ell)^{-1}$ on the left and
$\rho(\mathsf k_\ell)$ on the right; if the components of $g$ and $g'$ above $\ell$ are trivial one
may take $\gamma=\gamma'=1$, so that $\mathsf k=g$, $\mathsf k'=g'$, both matrices are the identity,
and the matrix of $\theta_j(\Phi;g,g')$ in the bases $(e_T)$ is exactly the matrix of jets below. By \eqref{ps:eq-kernel}, $\theta_j(\Phi;\mathsf k,\mathsf k')$ is the matrix of jets along
$S$ of $F=\Theta_f(\pi_Ms_0\,\cdot)$ with $f=\omega_{\mu_{j+1}}(\mathsf k,\mathsf k')
\phi^\circ_{j,x_j}$, which is a unit multiple of a standard datum at $\ell$ by (S1) (the values of
$\mu_{j+1}$ and $\upsilon$ at ideles that are units above $\ell$ are units at $\ell$). This is the
series $F$ of (K1); its jets are integral by (S2) if $\ell\ne d$, and (K2) is (S3): a prime
that does not split in $K$ is odd and does not divide $q_1M$.
\end{proof}

\subsubsection{The prime $N$: a second proof of Lemma~\ref{pl:kernelN} under a condition}\label{pk:ss-N}

We use the normalisations (N1)--(N5) of \S\ref{pl:ss-norm}; they agree with those of
\S\ref{pk:ss-conv} ($W_\lambda(\Z)=\Sigma_\lambda^\#$ with the dual basis, the action through
$\tau\circ c$, the lifts \eqref{pk:eq-lift}; Lemma~\ref{pl:sum} is Proposition~\ref{pk:K0}). Lemma~\ref{pl:kernelN} is proved in Corollary~\ref{pf:K1}(b); this section gives a second proof under the condition (Gen) of Proposition~\ref{pk:K1N}, and is not used elsewhere. Let
$\ell$ be an odd prime, $\iota\colon\C\to\C_\ell$, and $E_1\subset\C/\Lambda$ the group of the
torsion points that reduce to the origin of $\mathcal E$ (the kernel of reduction). Put
$\mathrm E_1^*(v)=\zeta(v)-s_2v-(\boldsymbol\pi/A)\aubar v$; it is $\Lambda$-periodic and odd, and
$\theta_u'(0)/\theta_u(0)=\mathrm E_1^*(u+t)$ by the definitions of $\theta$ and $\theta_u$.

\begin{lemma}\label{pk:E1}
Let $v\in\C$ be a torsion point, $v\notin E_1$, and suppose that there is $\alpha\in\auO_K$ with
$\alpha v\equiv v$ modulo $\Lambda$, $\iota(\alpha(1-\aubar\alpha))$ a unit and $\auNm\alpha$ prime
to $\ell$. Then $\mathrm E_1^*(v)\in\pao$. If $\ell$ splits in $K$ and $\mathfrak n$ is the prime
of $K$ induced by $\iota$, such $\alpha$ exists whenever the annihilator of $v$ in $\auO_K$ is
prime to $\aubar{\mathfrak n}$.
\end{lemma}
\begin{proof}
As in the proof of \eqref{pk:eq-s2}, $\alpha\mathrm E_1^*(\alpha z)-\sum_{w\in\alpha^{-1}\Lambda/
\Lambda}\mathrm E_1^*(z+w)$ is constant, and it is odd, so it is $0$. For $w\notin\Lambda$,
$\mathrm E_1^*(z+w)-\mathrm E_1^*(z)-\mathrm E_1^*(w)=\tfrac12(\wp'(z)-\wp'(w))/(\wp(z)-\wp(w))=:
\lambda(z,w)$, and $\sum_{w\ne0}\mathrm E_1^*(w)=0$. Hence $\alpha\mathrm E_1^*(\alpha z)-\auNm
\alpha\,\mathrm E_1^*(z)=\sum_{w\ne0}\lambda(z,w)$, and for $z=v$:
$\alpha(1-\aubar\alpha)\mathrm E_1^*(v)=\sum_{0\ne w\in\alpha^{-1}\Lambda/\Lambda}\lambda(v,w)$. In
the coordinates of $\mathcal E$, $\lambda(v,w)=(Y_v-Y_w)/(X_v-X_w)+a_1/2$. The points $v$ and $w$
do not reduce to the origin ($w$ has order prime to $\ell$), so their coordinates are integral;
and $X_v-X_w$ is a unit unless $v\mp w\in E_1$. This is impossible: $E_1$ is stable under
$\auO_K$ and contains no non-zero point of order prime to $\ell$, and $(\alpha-1)(v\mp w)\equiv
\pm w$ has order prime to $\ell$; so $v\mp w\in E_1$ would give $w=0$. For the last assertion:
choose $\alpha\equiv1$ modulo $\mathfrak n$ and modulo the annihilator of $v$, and $\alpha\equiv2$
modulo $\aubar{\mathfrak n}$; then $\iota(\alpha)\equiv1$ and $\iota(\aubar\alpha)\equiv2$ modulo
$\pam$, so $\iota(\alpha(1-\aubar\alpha))$ is a unit, and $\auNm\alpha$ is prime to $\ell$.
\end{proof}

\begin{lemma}\label{pk:formal}
Let $\ell\ne d$ be odd, $u$ a torsion point with $u\notin t+\Lambda+E_1$, and suppose that
$v=u+t$ satisfies the hypothesis of Lemma~\ref{pk:E1}. Then $\theta_u$ is $\partial$-integral.
\end{lemma}
\begin{proof}
As in the proof of Proposition~\ref{pk:curve}(C3)(i), $R_u=\theta_u/\theta$ satisfies
$R_u'/R_u=\mathsf C+\mathsf h$ with $\mathsf h=(y-y_Q)/(x-x_Q)-(\text{its value at }z=0)$,
$(x_Q,y_Q)$ the point of parameter $u-t$, and $\mathsf C=\mathrm E_1^*(u+t)\in\pao$
(Lemma~\ref{pk:E1}). If $Q$ does not reduce to the origin, $x_Q$ is a unit (it reduces to $\tilde
P_0$ only if $u-t\in E_1$) and $\mathsf h$ is $\partial$-integral, $d/dz$ preserving the local
ring of $\pao[x,y]$ at $x-x_Q$. If $Q$ reduces to the origin, with parameter $\mathsf t_Q=-x_Q/y_Q
\in\pam$, then $x_Q^{-1}\in\mathsf t_Q^2\pao$, $y_Q/x_Q\in\mathsf t_Q^{-1}\pao^\times$ and
$(y-y_Q)/(x-x_Q)=(y_Q/x_Q-y/x_Q)\sum_{k\ge0}(x/x_Q)^k$, whose terms other than the constant
$y_Q/x_Q$ have coefficients in $\pao$; so $\mathsf h$ is $\partial$-integral again. Hence
$R_u=\theta_u(0)\exp(\mathsf Cz+\int_0^z\mathsf h)$ is $\partial$-integral, $\theta_u(0)$ being
integral as in (C3)(i), and so is $\theta_u=\theta R_u$ by (C2).
\end{proof}

\begin{proposition}[Lemma~\ref{pl:kernelN}, with a restriction]\label{pk:K1N}
Let $\iota\colon\C\to\C_N$, and let $z$ be as in Lemma~\ref{pl:kernelN}: integral at the primes
induced by $\iota\circ\tau\circ c$ and with $Nz$ integral at the primes $\mathfrak l_\tau$ induced
by $\iota\circ\tau$ $(\tau\in\auPl)$, which are the primes of $E$ above the prime $\mathfrak n$ of
$K$ induced by $\iota$. Assume the following condition {\rm(Gen)}: in the expansion of
$\omega_{\mu_{j+1}}(\mathsf k,\mathsf k')\phi^\circ_{j,z}$ in translates $\psrho_f(\gamma)F_0$
used in the proof of Proposition~\ref{pk:std}, no point $\mathsf x_{k,i}\Omega_0$ lies in
$t+\Lambda+E_1$ outside $t+\Lambda$. Then the conclusion of Lemma~\ref{pl:kernelN} holds.
\end{proposition}
\begin{proof}
The reduction of Proposition~\ref{pk:K12} applies with $\mathsf k,\mathsf k'$ integral above $N$
((G2) of \S\ref{pl:ss-norm}), and (S1) holds for the data that are combinations of translates by
$x$ with $x$ integral above $\aubar{\mathfrak n}$ and $Nx$ integral above $\mathfrak n$ (the
elements integral above $N$ preserve these conditions). As in the proof of (S2), such a datum is
a combination with coefficients in $\pao$ of translates $\psrho_f(\gamma)F_0$ whose theta
functions at the scale $\pi_Ms_0$ are $\zeta_\gamma\prod\theta_{\mathsf x\Omega_0}(\mathsf c\,
u_k^\dagger(z_i))^{q_1}$ with $\mathsf c$ a unit at $N$ ($N\nmid2q_1M$) and $\mathsf x=-\pktr_{E/K}(
\aubar{u_k}\gamma_i)\in K$ with denominator prime to $\aubar{\mathfrak n}$: the trace of an element
of $E$ that is integral above $\aubar{\mathfrak n}$ is integral at $\aubar{\mathfrak n}$. So the
annihilator of $\mathsf x\Omega_0+t$ is prime to $\aubar{\mathfrak n}$ ($N$ is odd). By (Gen) each
factor is either a multiple of $\theta_t$, treated in the proof of (S3), or satisfies the
hypotheses of Lemma~\ref{pk:formal}; so each factor, and the product, is $\partial$-integral.
\end{proof}
Lemma~\ref{pl:kernelN} is proved in full, without the condition (Gen), in Corollary~\ref{pf:K1}(b):
the translates $\theta_u$ with $u\in t+E_1$ are $\partial$-integral as well (Remark~\ref{pf:gen}). So
Lemmas~\ref{pk:E1}, \ref{pk:formal} and Proposition~\ref{pk:K1N} are not used in the sequel.



\providecommand{\auO}{\mathcal O}
\providecommand{\auPl}{\Phi}
\providecommand{\auPa}{\Phi_a}
\providecommand{\auPb}{\Phi_b}
\providecommand{\auS}{\mathcal S}
\providecommand{\auSram}{\mathcal S^{\mathrm{ram}}}
\providecommand{\auDX}{D_X}
\providecommand{\augain}{\operatorname{gain}}
\providecommand{\auerr}{\operatorname{err}}
\providecommand{\auNm}{\operatorname{Nm}}
\providecommand{\auU}{\mathrm U}
\providecommand{\auGL}{\mathrm{GL}}
\providecommand{\auHom}{\operatorname{Hom}}
\providecommand{\psrho}{\varrho}
\providecommand{\psGlue}{\operatorname{Glue}}
\providecommand{\auA}{\mathbb A}
\providecommand{\auF}{\mathbb F}
\providecommand{\aucol}{\operatorname{col}}
\providecommand{\aubar}[1]{\overline{#1}}
\providecommand{\auf}[1]{\texttt{#1}}
\providecommand{\pao}{\mathfrak o}
\providecommand{\paCp}{\C_p}
\providecommand{\paI}{I_p^{\circ}}
\providecommand{\paded}{\operatorname{ded}}
\providecommand{\paPred}{\operatorname{Pred}}
\providecommand{\paloss}{\operatorname{loss}}
\providecommand{\plone}{\mathbf 1}
\providecommand{\plD}{\mathsf D}
\providecommand{\pld}{\mathsf d}
\providecommand{\ple}{\mathsf e}
\providecommand{\plP}{\mathsf P}
\providecommand{\plSi}{\Sigma_N}
\providecommand{\plMir}{\mathrm{Mir}}
\providecommand{\pldir}{\plone^{\mathrm{dir}}}
\providecommand{\plrow}{\plone^{\mathrm{row}}}
\providecommand{\plcol}{\plone^{\mathrm{col}}}
\providecommand{\pltr}{\operatorname{tr}}
\providecommand{\plvol}{\operatorname{vol}}
\providecommand{\plr}{\mathfrak r}
\providecommand{\plE}{\mathsf E}
\providecommand{\plH}{\mathcal H}
\providecommand{\plgb}[2]{\|#2\|_{#1}}

\section{Hecke operators through the lifts, and the level prime}\label{pl:sec}

This section proves Lemma~\ref{pa:hecke} (\S\S\ref{pl:ss-count}--\ref{pl:ss-hecke}) and
Lemma~\ref{pa:level} (\S\ref{pl:ss-level}); the latter is first given a closed statement
(Lemma~\ref{pl:level}), for which the lattices at $N$, the operator $U_N$ and the gluing are
defined. Both proofs use Lemma~\ref{pa:kernel} only as stated, except that at $\ell=N$, which
Lemma~\ref{pa:kernel} excludes, the same integrality statement is needed and is isolated as
Lemma~\ref{pl:kernelN}; it is proved in Corollary~\ref{pf:K1}(b). \S\ref{pl:ss-norm} fixes the normalisations that
{\cite[\S14.3]{P46}} leaves open and checks that they are the ones under which
{\cite[Theorem~2.2]{P50}} was applied in {\cite[\S2.7]{P50}}; \S\ref{pl:ss-wa} proves weak approximation
in the form used.

\subsection{Normalisations}\label{pl:ss-norm}
\emph{{\rm(N1)} Weight modules.} For $\auGL_n$ and a diagram with $c(h)$ columns of height $h<n$
let $\Sigma_\lambda$ be the lattice spanned by the products of minors $\nabla_T(Y)$ of a generic
$n\times\infty$ matrix $Y$, with the standard ones as basis (\S\ref{pa:ss-count},
\S\ref{ps:ss-restrict}). We put
\[
W_\lambda=\operatorname{Hom}_\Z(\Sigma_\lambda,\C),\qquad W_\lambda(\Z)=\operatorname{Hom}_\Z
(\Sigma_\lambda,\Z),\qquad(e_T)\ \text{the basis dual to }(\nabla_T)_{T\text{ standard}},
\]
and $g\in\auGL_n(\C)$ acts by $(\rho(g)\xi)(P)=(\det g)^{e}\,\xi\bigl(P(g\,Y)\bigr)$. Thus the
evaluations satisfy $\rho(g)\mathrm{ev}_Y=(\det g)^e\mathrm{ev}_{gY}$, and the matrix of $\rho(g)$
in the basis $(e_T)$ is $(\det g)^eS_\lambda(g)$, where $\nabla_{T'}(gY)=\sum_TS_\lambda(g)_{T',T}
\nabla_T(Y)$ is the Schur matrix of \S\ref{pa:ss-count}; the same formula defines $S_\lambda(z)$
for a rectangular matrix $z$ over any ring, with entries that are integral polynomials in $z$. For
$g=\operatorname{diag}(u_1,\dots,u_n)$,
\begin{equation}\label{pl:eq-diag}
\rho(g)\,e_T=(u_1\cdots u_n)^{e}\prod_{\text{boxes }\square\text{ of }T}u_{T(\square)}\cdot e_T .
\end{equation}
For a single row of length $a$ ($a$ columns of height $1$), $\Sigma_\lambda$ is the lattice of
the forms of degree $a$ in the first column $y$ of $Y$ with its basis of monomials, and the matrix
of $\rho(g)$ in $(e_\beta)$ is the transpose of the matrix, in the monomials, of the substitution
$P(y)\mapsto P(gy)$. This is the lattice $\Sigma_\lambda^{\#}$ of \S\ref{ps:ss-restrict} and the
basis in which Lemma~\ref{pa:kernel}(K1) is stated.

\emph{{\rm(N2)} The action of the groups, and the gauge.} For $G=\auU(X)$, $X$ standard of rank
$n$, the group $\Gamma$ acts on $W_\lambda=\bigotimes_\tau W_{\lambda_\tau}$ through the embeddings
$\tau\circ c$: $\gamma$ acts by $\bigotimes_\tau\rho_\tau(\tau(\aubar\gamma))$, and
$\tau(\aubar\gamma)={}^t\tau(\gamma)^{-1}$. Let $\ell$ be a prime and $\iota\colon\C\to\C_\ell$ a
homomorphism, with valuation ring $\pao$. For $\tau\in\auPl$ let $\mathfrak l_\tau$ be the prime of
$E$ induced by $\iota\circ\tau$ and $w_\tau=\aubar{\mathfrak l_\tau}$ the prime induced by
$\iota\circ\tau\circ c$, and let $\hat\iota_\tau\colon E_{w_\tau}\to\C_\ell$ be the continuous
extension of $\iota\circ\tau\circ c$. An element $g\in G_f$ acts on $W_\lambda\otimes_{\C,\iota}
\C_\ell$ by $\rho(g_\ell)=\bigotimes_\tau\rho_\tau(\hat\iota_\tau(g_{w_\tau}))$, where $g_w\in\auGL_n
(E_w)$ is the component of $g$ at $w$; this extends the action of $\Gamma$. With $\|\cdot\|$ the
largest absolute value of a coordinate in $(e_T)$, the gauge at $(\iota,g)$ of $\xi$ is
$\plgb g\xi:=\|\rho(g_\ell)^{-1}\iota(\xi)\|$. A form $\varphi$ \emph{has gauge at most $C$ at
the base point $\pld\in G_f$} if $\plgb{g\pld}{\varphi(g)}\le C$ for all $g$; for $\pld=1$ this is
the notion of {\cite[\S14.3]{P46}}. We use:
\begin{itemize}
\item[(G1)] $\plgb{\gamma g}{\gamma\xi}=\plgb g\xi$ for $\gamma\in\Gamma$; the gauge is
ultrametric in $\xi$;
\item[(G2)] if all components of $k\in G_f$ above $\ell$ are integral with integral inverse
(``$k$ is integral above $\ell$''), then $\rho(k_\ell)$ has integral entries and integral inverse,
so $\plgb{gk}\xi=\plgb g\xi$ and $\|\rho(k_\ell)\eta\|=\|\eta\|$.
\end{itemize}

\emph{{\rm(N3)} Hecke operators.} Let $\mathfrak l\ne\aubar{\mathfrak l}$ be a prime of $E$ with
uniformiser $\varpi$. The component at $\mathfrak l$ identifies $G$ at the place below $\mathfrak
l$ with $\auGL_n(E_{\mathfrak l})$ (the component at $\aubar{\mathfrak l}$ being ${}^t\aubar
g^{-1}$); for $h\in\auGL_n(E_{\mathfrak l})$ we denote by the same letter the element of $G_f$ that
is $h$ there and $1$ at the other places. Let $K_{\mathfrak l}=\auGL_n(\auO_{\mathfrak l})$ and
$d_i=\operatorname{diag}(\varpi1_i,1_{n-i})$. On a representation $(\pi,V)$ of $G_f$ and a vector
fixed by $K_{\mathfrak l}$,
\[
T_{\mathfrak l,i}\,v=\sum_{h\in K_{\mathfrak l}d_iK_{\mathfrak l}/K_{\mathfrak l}}\pi(h)v ;\qquad
\text{on forms,}\quad(T_{\mathfrak l,i}\varphi)(g)=\sum_h\varphi(gh).
\]
The operator $T_{\mathfrak l}$ of \S\ref{pa:ss-divis} is $T_{\mathfrak l,\plr_p}$ on $\auS$, with
$\varpi=p$.

\emph{{\rm(N4)} Lifts with arbitrary data.} Let $G=\auU(X_j)$, $G'=\auU(X_{j+1})$ and let $dg$ be
the Haar measure of $G_f$ for which $J_{X_j}$ has measure $1$. For $\Phi\in\plH_f(\auHom(X_j,
X_{j+1}))$ and a smooth form $\varphi$ on $X_j$ of weight $\lambda(X_j)$ put
\[
(\Theta^X_j(\Phi)\varphi)(g')=\int_{\Gamma\backslash G_f}\theta_j(\Phi;g,g')\,\varphi(g)\,dg ,
\]
a finite sum; here $\theta_j$ is \eqref{ps:eq-kernel}. We use the following properties of the
kernels, which express that $\theta$ is invariant under rational points (Lemma~\ref{ps:weil}(W1))
and that $\theta_j(\Phi;g,g')$ is defined through $\omega_{\mu_{j+1}}(g,g')\Phi$:
\begin{itemize}
\item[(T1)] $\theta_j(\Phi;\gamma g,\gamma'g')=\gamma'\circ\theta_j(\Phi;g,g')\circ\gamma^{-1}$ for
$\gamma\in\Gamma$, $\gamma'\in\Gamma'$;
\item[(T2)] $\theta_j(\omega_{\mu_{j+1}}(h,h')\Phi;g,g')=\theta_j(\Phi;gh,g'h')$ for $h\in G_f$,
$h'\in G'_f$.
\end{itemize}
By (T1) the integrand is a function on $\Gamma\backslash G_f$ and $\Theta^X_j(\Phi)\varphi$ is a
form on $X_{j+1}$; by (T2), with $R$ the right translation,
\begin{equation}\label{pl:eq-equiv}
R_{h'}\circ\Theta^X_j(\Phi)=\Theta^X_j\bigl(\omega_{\mu_{j+1}}(h,h')\Phi\bigr)\circ R_h .
\end{equation}
If $\Phi$ and $\varphi$ are fixed by the free level $J_{X_j}$, this is the sum of
{\cite[\S14.8]{P46}}. For a tensor of data, $C_\Phi=\Theta^X_{m-1}\cdots\Theta^X_0(\Phi)$ is defined
by linearity from pure tensors, as in \S\ref{ps:ss-glue}.

\begin{lemma}[Chains are plain sums]\label{pl:sum}
Let $\Phi_{x,j}\in\plH_f(\auHom(X_j,X_{j+1}))$ $(x$ in a finite set, $0\le j<m)$ and $c_x\in\C$ be
such that the tensor $\sum_xc_x\bigotimes_j\Phi_{x,j}$ is fixed by $\bigotimes_j\omega_{\mu_{j+1}}
(b_j,b_{j+1})$ for all $b_j\in J_{X_j}$ $(j<m)$, $b_m=1$. Then for $\varphi\in\auS$ and $g_m$,
\[
\Bigl(\sum_xc_x\,\Theta^X_{m-1}(\Phi_{x,m-1})\cdots\Theta^X_0(\Phi_{x,0})\varphi\Bigr)(g_m)=
\sum_{g_0,\dots,g_{m-1}}\sum_xc_x\,\theta_{m-1}(\Phi_{x,m-1};g_{m-1},g_m)\cdots\theta_0(\Phi_{x,0};
g_0,g_1)\,\varphi(g_0),
\]
where $g_i$ runs over representatives of $\Gamma_i\backslash\auU(X_i)(\auA_{E^+,f})/J_{X_i}$.
\end{lemma}
\begin{proof}
The left side is the integral over $\prod_{j<m}\Gamma_j\backslash\auU(X_j)(\auA_f)$ of the function
$F(g_0,\dots,g_{m-1})$ given by the inner sum on the right. By (T1) it is invariant on the left
under $\prod\Gamma_j$. Replacing every $g_j$ by $g_jb_j$ with $b_j\in J_{X_j}$ replaces, by (T2),
each $\Phi_{x,j}$ by $\omega(b_j,b_{j+1})\Phi_{x,j}$ and $\varphi(g_0)$ by $\varphi(g_0b_0)=
\varphi(g_0)$; the sum over $x$ is the value of a multilinear function at the tensor, which is
unchanged. So $F$ is right invariant under $\prod J_{X_j}$, and each double coset has measure
$1$ because the levels are free ({\cite[Lemma~14.5]{P46}}).
\end{proof}
Formula (K0) of Lemma~\ref{pa:kernel} is the case $\Phi_{x,j}=\phi^\circ_{j,x_j}$, $c_x=\Psi'(x)$.

\emph{{\rm(N5)} The Weil representation at split places.} Let $\mathsf S$ be a finite set of primes
of $E$ with $\mathsf S\cap\aubar{\mathsf S}=\emptyset$, $E_{\mathsf S}=\prod_{w\in\mathsf S}E_w$,
and let $n=n_j$, $n'=n_{j+1}$ (both are even). By Lemma~\ref{ps:weil}(W5), $\plH_f(\auHom(X_j,
X_{j+1}))=\auS(M_{n'\times n}(E_{\mathsf S}))\otimes\plH^{\mathsf S}$, where the first factor is the
space of Schwartz--Bruhat functions of $T=(a_w^{-1}x_w)_{w\in\mathsf S}$ ($x_w$ the matrix in the frames,
$a$ the idele of \S\ref{ps:ss-weil}: $a_w=\alpha_w$ at the primes $w\mid\omega$ above $q_1$ and $a_w=1$
elsewhere; these are the coordinates $\tilde x$ of {\cite[Lemma~6.32]{P53}}) and the second is the
representation at the other places. In these coordinates the vacuum at $\mathsf S\cup\aubar{\mathsf
S}$ is the indicator $\plone$ of $M_{n'\times n}(\auO_{\mathsf S})$, and we use (W5) in the
following explicit form: for $g\in\auGL_n(E_{\mathsf S})$, $h\in\auGL_{n'}(E_{\mathsf S})$, read as
elements of the unitary groups supported on $\mathsf S\cup\aubar{\mathsf S}$,
\begin{equation}\label{pl:eq-w5}
\omega_{\mu_{j+1}}(g,h)(\Psi\otimes v)=\frac{\chi_{n}(h)}{\chi_{n'}(g)}\,\Psi(h^{-1}Tg)\otimes v,
\qquad\chi_r(g)=\prod_{w\in\mathsf S}\chi_{j+1,w}(\det g_w)\,|\det g_w|_w^{-r/2},
\end{equation}
where $\chi_{j+1}=\mu_{j+1}\circ(x\mapsto x/\aubar x)$ is the base change of $\mu_{j+1}$ to a Hecke
character of $E$; and for $a$ supported on $\mathsf S$, $a'$ supported on $\aubar{\mathsf S}$ (in
the coordinates $T$, and with $E_{\aubar w}$ identified with $E_w$ by $c$),
\begin{equation}\label{pl:eq-heis}
\psrho(a)(\Psi\otimes v)=\Psi(T+a)\otimes v,\qquad\psrho(a')(\Psi\otimes v)=\lambda\bigl(\pltr
({}^ta'T)\bigr)\Psi\otimes v,\qquad\lambda=\prod_{w\in\mathsf S}\lambda_w,
\end{equation}
with $\lambda_w(t)=\psi_w(\jmath(a\aubar a)_wt)$, a character of $E_w$ that is trivial on
$\auO_w$ and not on $\varpi_w^{-1}\auO_w$ (self-duality of $\widehat\auO_E$ for
$\psi(\jmath\alpha\aubar\alpha\,\cdot)$, \S\ref{ps:ss-weil}, and $a\aubar a/\alpha\aubar\alpha$ a unit);
here $a'$ is written in the coordinate $a_{\aubar w}^{-1}x_{\aubar w}$ of the $\aubar w$-component. Formula \eqref{pl:eq-heis} is $[x,y]=\psi(\jmath\pltr(y^\dagger x))$ for the
polarisation $\auHom(X_{j,w},X_{j+1,w})\oplus\auHom(X_{j,\aubar w},X_{j+1,\aubar w})$; these
operators do not depend on the model, both summands being isotropic. Formula \eqref{pl:eq-w5} is
Proposition~\ref{pwf:W5}, for the normalisation (D3) of Definition~\ref{pwf:def} at split places;
for even ranks the powers of $\upsilon$ in the splitting and in the determinant twists $\upsilon_n$
of \S\ref{ps:ss-weil} cancel, leaving the base change of $\mu_{j+1}$. (Only two properties of $\chi_{j+1}$
are used below: it is unramified where $\mu_{j+1}$ is, and Lemma~\ref{pl:char}.)

\begin{remark}[The normalisation (N) of {\cite[\S2.7]{P50}}]\label{pl:N}
For $\lambda^X$, (N1)--(N3) are the normalisation (N): $V_\tau=W_{\lambda^X_\tau}$ has the basis
$(e_\beta)$ dual to the monomials, there is no power of the determinant, $\gamma$ acts by
$\rho_\tau(\tau(\aubar\gamma))$, whose matrix is the transpose of that of the substitution
$x\mapsto x\cdot{}^t\tau(\aubar\gamma)$ on row vectors, and $T_{\mathfrak l}$ is the operator of
(N3). With this choice the module $\auS_\iota$ of \S\ref{pa:ss-divis} is the image of
$f\colon\mathcal F_X\to\auS$, $F\mapsto(\varphi\colon\varphi(x)=\sum F(x)_{\eta,\beta}\,e_\eta
\otimes e_\beta)$, by (G1), (G2) and the choice of representatives $x$ integral above $p$
(Corollary~\ref{pl:wacor}(a)); and for $\gamma_e=\beta_\varkappa\mathsf d\gamma(e)^{-1}$ the
substitution is by $\tau_s(\gamma_e)^{-1}$, that is, by $\iota_s(\gamma(e))\cdot\iota_s(\mathsf
d^{-1})\cdot\iota_s(\beta_\varkappa^{-1})$ for even $s$ and by $\iota_s({}^t\gamma(e)^{-1})\cdot
\mathsf d\cdot\iota_s({}^t\beta_\varkappa)$ for odd $s$; substitutions compose in the order
written, and $p\,\mathsf d^{-1}$, $\mathsf d$ multiply the deep variables of an even, respectively
odd, letter by $p$; so $p^{A_{\mathfrak l}}$ times its matrix on $\bigotimes_sV_{\tau_s}$ is
$\rho(\gamma(e))\Delta^2\rho(t(\varkappa))$, and $\gamma_e$ acts on coordinates by the transpose. This is {\cite[(4)]{P50}}; so {\cite[Theorem~2.2]{P50}}
is applied in {\cite[\S2.7]{P50}} to the lattice $\auS_\iota$ of Theorem~\ref{pa:divis}. Had
$W_\lambda(\Z)$ been the lattice $\Sigma_\lambda$ itself, or had $\Gamma$ acted through $\tau$,
the operator would not have the form {\cite[(4)]{P50}}.
\end{remark}

\subsection{Weak approximation}\label{pl:ss-wa}
\begin{proposition}\label{pl:wa}
Let $F$ be a number field, $E/F$ a quadratic extension with conjugation $c$, $V$ a non-degenerate
hermitian space over $E$, $G=\auU(V)$ and $\mathsf T$ a finite set of places of $F$. Then $G(F)$ is
dense in $\prod_{v\in\mathsf T}G(F_v)$.
\end{proposition}
\begin{proof}
For an $F$-algebra $R$ let $\mathfrak s(R)$ be the $R$-module of the $X\in\operatorname{End}_{E
\otimes R}(V\otimes R)$ with $X^\dagger=-X$, and $\mathfrak s^\circ$, $G^\circ$ the subsets of
$\mathfrak s$, $G$ where $1+X$, respectively $1+u$, is invertible. The Cayley transform
$X\mapsto(1-X)(1+X)^{-1}$ is a bijection $\mathfrak s^\circ(R)\to G^\circ(R)$ with inverse
$u\mapsto(1-u)(1+u)^{-1}$: if $X^\dagger=-X$ then $u^\dagger u=(1-X)^{-1}(1+X)(1-X)(1+X)^{-1}=1$
and $1+u=2(1+X)^{-1}$; if $u^\dagger=u^{-1}$ then $X^\dagger=(1+u^{-1})^{-1}(1-u^{-1})=-X$. For
$R=F_v$ it is a homeomorphism. Now $\mathfrak s(F)$ is a finite-dimensional $F$-vector space, so it
is dense in $\prod_{v\in\mathsf T}\mathfrak s(F_v)$ (weak approximation in $F$), and
$\prod_v\mathfrak s^\circ(F_v)$ is open; hence $G^\circ(F)$ is dense in $\prod_vG^\circ(F_v)$. It
remains to see that $G^\circ(F_v)$ is dense in $G(F_v)$. Let $u\in G(F_v)$. The group $\{\zeta\in
(E\otimes F_v)^\times:\zeta\aubar\zeta=1\}$ contains the elements $x/\aubar x$, which form an
infinite set accumulating at $1$; all but finitely many of them satisfy $\det(1+\zeta u)\ne0$, and
$\zeta u\in G(F_v)$.
\end{proof}

\begin{corollary}\label{pl:wacor}
Let $X$ be a standard space, $G=\auU(X)$, $\ell$ a prime.
{\rm(a)} For $g,d\in G_f$ there is $\gamma\in\Gamma$ such that $\gamma g$ and $d\,\gamma g\,d^{-1}$
are integral above $\ell$. In particular every class of $\Gamma\backslash G_f/J$ has a
representative that is integral above $\ell$.
{\rm(b)} If $\varphi$ is a form of level $J\subseteq K^\circ$ and $(x)$ a system of representatives
that are integral above $\ell$, and if $J$ is the full group $\auGL_n(\auO_w)$ at the places above
$\ell$, then $\varphi$ is integral at $\iota$ if and only if all $\iota(\varphi(x))$ have
coordinates in $\pao$.
\end{corollary}
\begin{proof}
(a) The set of the $y\in\prod_{v\mid\ell}G(E^+_v)$ with $yg_\ell$ and $d_\ell yg_\ell d_\ell^{-1}$
integral is open and contains $g_\ell^{-1}$; apply Proposition~\ref{pl:wa}. (b) By (G1), (G2),
$\plgb{\gamma xk}{\varphi(\gamma xk)}=\plgb{xk}{\varphi(x)}=\|\iota(\varphi(x))\|$.
\end{proof}
This is the form of weak approximation used in {\cite[\S2.7]{P50}}.

\subsection{Values of the characters}\label{pl:ss-char}
\begin{lemma}\label{pl:char}
Let $1\le i\le m$, $\chi=\chi_i$ the base change of $\mu_i=\mu_i^0\vartheta^i$, and $w\ne\aubar w$
a prime of $E$ above $\ell$, unramified over $\Q$, not chosen and not in $Q_\vartheta\cup\aubar{
Q_\vartheta}$. Then $\chi_w$ is unramified, $\chi_w(\ell)$ is an algebraic number, and for
$\iota\colon\C\to\C_\ell$, with $v$ the valuation normalised by $v(\ell)=1$,
\[
\bigl|v\bigl(\iota\,\chi_w(\ell)\bigr)\bigr|=\Bigl|\sum_{\tau:\,\mathfrak l_\tau=w}\kappa_i(\tau)-
\sum_{\tau:\,\mathfrak l_\tau=\aubar w}\kappa_i(\tau)\Bigr| .
\]
\end{lemma}
\begin{proof}
$\mu_i$ is unramified at the place below $w$, so $\chi_w$ is trivial on $\auO_w^\times$. Let $h$ be
the class number of $E$ and $w^h=(\beta)$, $\beta\in\auO_E$. As $\chi$ is trivial on $E^\times$,
$\chi_w(\ell)^h=\chi_w(\beta)=\chi_\infty(\beta)^{-1}\prod_{w'\ne w}\chi_{w'}(\beta)^{-1}$. The
element $\beta$ is a unit at every $w'\ne w$, and a continuous character of the compact totally
disconnected group $\prod_{w'\ne w}\auO_{w'}^\times$ has finite image; so the last product is a root
of unity (its factors are $1$ at almost all $w'$). The character $\mu_i$ has weight $-\kappa_i$ at the infinite
places, so $\chi_\infty(\beta)=\prod_\tau(\tau(\beta)/\tau(\aubar\beta))^{\mp\kappa_i(\tau)}$ with
one sign for all $\tau$. Finally $v(\iota\tau(\beta))$ is $h$ if $\mathfrak l_\tau=w$ and $0$
otherwise, and $v(\iota\tau(\aubar\beta))$ is $h$ if $\mathfrak l_\tau=\aubar w$ and $0$
otherwise.
\end{proof}

\subsection{The count with charges, and the gain with marked places}\label{pl:ss-count}
We prove the variant of Proposition~\ref{pa:count} announced in Remark~\ref{pa:charges}. Notation
and hypotheses are those of \S\ref{pa:ss-count}.

\begin{proposition}[Count with charges]\label{pl:charges}
Let $F$ have the normal form \eqref{pa:eq-nf}, let $\kappa$ be a shape with $B\ge1$ boxes, $T$ a
filling and $y$ integral directions. Let $\plE\subseteq\auPl$ and let integers $w_{\tau,h,b}\ge0$
$(\tau\in\auPl$, $1\le h<n'$, $b\in[n'])$ be such that for all $e\in\paI$, $\tau$, $b$, $h'$, with
$\tau'=\mathrm{pr}_e\tau$:
\begin{equation}\label{pl:eq-chargehyp}
\tau\in\plE,\ \text{or}\ \ (h'=1\ \text{and}\ \tau'\in\plE)\quad\Longrightarrow\quad
w_{\tau,1,b}>0\ \ \text{or}\ \ w_{\tau',h',b}>0 .
\end{equation}
Let $\kappa^\circ$ be $\kappa$ without its columns of height one at the places of $\plE$, with
$B^\circ$ boxes, and $w(T)=\sum w_{\tau,h,b}$, the sum over the boxes of $T$, a box in row $b$ of a
column of height $h$ at $\tau$ counting $w_{\tau,h,b}$. Then
\[
v\bigl(Q_{T,y}(\partial)F(0)\bigr)+w(T)\ \ge\ \paPred(\kappa^\circ)+\gamma\,(B-B^\circ)-f\,n'^3\,
(1+\log_pB).
\]
\end{proposition}
\begin{proof}
We follow the four steps of the proof of Proposition~\ref{pa:count}, with the substitution
\[
z_{\tau,b,a}(s)=\sum_{h,k}p^{\,w_{\tau,h,b}}\,\beta_{\tau,h,k}(a)\,s_{\tau,h,k,b}.
\]
\emph{Step 1.} Now $s_{\tau,h,k,b}$ is replaced by $p^{w_{\tau,h,b}}\sum_{k'}g_{\tau,h}(k,k')(z_\tau
y_\tau)_{b,k'}$, so the matrix of a column $(\tau,h,c)$ with rows $b_1,\dots,b_h$ is multiplied in
addition on the left by $\operatorname{diag}(p^{w_{\tau,h,b_i}})$, and \eqref{pa:eq-chain} becomes
$P_*(\partial)G(0)=\varpi^{B}p^{\,w(T)}\prod_{\tau,h}\det(g_{\tau,h})^{\kappa_\tau(h)}\cdot Q_{T,y}
(\partial)F(0)$.

\emph{Step 2.} Modulo $p$, the coefficient of $s_{\tau,h,k,b}^{\,p^e}s_{\tau',h',k',b}$ in
$\Gamma_e(z(s))$ is $p^{\,p^ew_{\tau,h,b}+w_{\tau',h',b}}[\beta_{\tau,h,k},\beta_{\tau',h',k'}]_e$.
By \eqref{pl:eq-chargehyp} it is divisible by $p$, hence by $\pi$, for the pairs of heights
$(1,h')$ if $\tau\in\plE$, and for the pair $(1,1)$ if $\tau'\in\plE$; call these pairs
\emph{free}. For the others we impose the congruences (a), (b), (c) of the proof of
Proposition~\ref{pa:count}, with the following changes in (b), for a pair $(\tau,e)$: if
$\tau\in\plE$ nothing is imposed; if $\tau\notin\plE$ and $\tau'\in\plE$, either $\pi_e$ divides
$g_{\tau,1}$, or the congruence modulo $\pi$ is imposed on the rows of the $g_{\tau',h'}$ with
$h'\ge2$ only, whichever is cheaper. Every pair of heights that is not free is covered as before.
The cost of (b) for $(\tau,e)$ is now $v(\pi)\min(a^\circ_\tau/p^e,a^\circ_{\tau'}+a'_{\tau'})$,
where $a^\circ_\sigma=a_\sigma$ for $\sigma\notin\plE$ and $a^\circ_\sigma=0$ for
$\sigma\in\plE$; the costs of (a) and (c) are unchanged. So \eqref{pa:eq-goal} holds with the
deductions $\paded^\circ_e(\tau)$ of the shape $\kappa^\circ$ (whose numbers are $a^\circ_\tau$,
$a'_\tau$, $l_\tau$).

\emph{Steps 3, 4} are unchanged and give $v(P_*(\partial)G(0))\ge\frac B{p-1}-fn'^3(1+\log_pB)$.
Hence
\[
v\bigl(Q_{T,y}(\partial)F(0)\bigr)+w(T)\ge\gamma B-v(\pi)\sum_{e\in\paI}\sum_\tau\paded^\circ_e
(\tau)-fn'^3(1+\log_pB),
\]
and $\gamma B^\circ-v(\pi)\sum_{e,\tau}\paded_e^\circ(\tau)=\paPred(\kappa^\circ)$ by
\eqref{pa:eq-pred}.
\end{proof}

For $p\in\mathcal N$, $0\le j<m$ and $\plE\subseteq\auPl$ let $\kappa^\circ$ be the shape of the
$j$-th lift (\S\ref{pa:ss-divis}) without its columns of height one at $\plE$, and
\[
D_j(\plE)=\paPred(\kappa^\circ)+\gamma_p\,a\,\#(\auPa\cap\plE)-f\,n_{j+1}^3(1+\log_pB_j),\qquad
C_p(\plE)=\sum_{(\tau,e)}\frac a{p^e},
\]
the last sum over the pairs with $e\in I_p$ and $\tau,u(-p^e)\tau\in\auPa\setminus\plE$.

\begin{lemma}\label{pl:gain}
Let $p\in\mathcal N$ be inert in $K$. Then
$D_j(\plE)\ge\dfrac{\augain_p}m-C_p(\plE)+\#(\auPa\cap\plE)\,\dfrac a{x_1}-\dfrac{f\,a\,(\log b)^2}
{p\,m}$.
\end{lemma}
\begin{proof}
This is the proof of Lemma~\ref{pa:gain} with two changes; recall $v(\pi)=1$, $\gamma_p=\sum_{e\in
I_p}p^{-e}$, and that the numbers of $\kappa^\circ$ are $a^\circ_\tau$, $a'_\tau$, $l_\tau$. In
(1), for $\tau\in\auPa\setminus\plE$: if $\tau'\notin\auPa\setminus\plE$ then $a^\circ_{\tau'}=0$
and $(a/p^e-a^\circ_{\tau'}-a'_{\tau'})^+\ge(a/p^e-\kappa(\tau'))^+-s$; otherwise $(\tau,e)$ is one
of the pairs of $C_p(\plE)$ and the loss is at most $a/p^e$. For $\tau\in\auPa\cap\plE$ the shape
$\kappa^\circ$ contributes nothing, but the term $\gamma_pa$ of $D_j(\plE)$ gives
\[
\sum_{e\in I_p}\frac a{p^e}\ \ge\ \sum_{e\in I_p}\Bigl(\frac a{p^e}-\kappa(\tau')\Bigr)^++\min\Bigl(
\frac ap,\kappa(u(-p)\tau)\Bigr)\ \ge\ \sum_{e\in I_p}\Bigl(\frac a{p^e}-\kappa(\tau')\Bigr)^++\frac
a{x_1},
\]
because $1\in I_p$, $p<x_1$ and $\kappa\ge a/x_1$. Steps (2), (3) and the comparison of sizes are
those of Lemma~\ref{pa:gain}.
\end{proof}

\subsection{Hecke operators through one lift and through the chain}\label{pl:ss-transfer}
Let $\mathfrak l\ne\aubar{\mathfrak l}$ be a prime of $E$, with residue field of $q$ elements and
uniformiser $\varpi$, at which $\mu_{j+1}$ is unramified. We use (N5) for $\mathsf S=\{\mathfrak
l\}$; write $K=\auGL_n(\auO_{\mathfrak l})$, $K'=\auGL_{n'}(\auO_{\mathfrak l})$, $n'=n+k_0$ with
$k_0=t+1$, and on data fixed by $\omega(K,K')$ put
\[
T'_i=\sum_{h\in K'd_iK'/K'}\omega_{\mu_{j+1}}(1,h),\qquad T^-_i=\sum_{g\in Kd_i^{-1}K/K}
\omega_{\mu_{j+1}}(g,1),\qquad P'(X)=\sum_{i\ge0}(-X)^iq^{\binom i2}T'_i,
\]
and $P^-(X)$ likewise with $T^-_i$. For a representation as in (N3) let $P_{\mathfrak l}(X)=\sum_i
(-X)^iq^{\binom i2}T_{\mathfrak l,i}$.

\begin{lemma}[The vacuum]\label{pl:vac}
Let $\alpha=\chi_{j+1,\mathfrak l}(\varpi)q^{n/2}$ and $\beta=\chi_{j+1,\mathfrak l}(\varpi)
q^{n'/2}$. For every $v$, the vector $\plone\otimes v$ is fixed by $\omega(K,K')$, and
\[
P'(X)(\plone\otimes v)=\prod_{l<k_0}(1-q^l\alpha X)\cdot P^-\bigl(q^{k_0}\alpha\beta^{-1}X\bigr)
(\plone\otimes v).
\]
\end{lemma}
\begin{proof}
The first claim is \eqref{pl:eq-w5}, $\chi_{j+1,\mathfrak l}$ being unramified. For $h\in
K'd_iK'$, \eqref{pl:eq-w5} gives $\omega(1,h)\plone=\alpha^i\,\plone(h^{-1}T)$, and $\plone(h^{-1}
T)=1$ if and only if $T$ is integral and the image of $T$ modulo $\varpi$ lies in the image of
$h\auO^{n'}$ in $\auF_q^{n'}$, a subspace of codimension $i$; every such subspace occurs for
exactly one coset. So $T'_i\plone$ is $\alpha^i\plone$ times the number of subspaces of dimension
$i$ of a space of dimension $n'-\operatorname{rk}(T\bmod\varpi)$. Similarly, for $g\in
Kd_i^{-1}K$, $\omega(g,1)\plone=\beta^i\plone(Tg)$, and $\plone(Tg)=1$ if and only if $T$ is
integral and $T$ modulo $\varpi$ kills the image of $\varpi g\auO^n$ in $\auF_q^n$, a subspace of
dimension $i$; so $T^-_i\plone$ is $\beta^i\plone$ times the number of subspaces of dimension $i$
of $\ker(T\bmod\varpi)$, of dimension $n-\operatorname{rk}(T\bmod\varpi)$. For a space $V$ of
dimension $d$ over $\auF_q$,
\begin{equation}\label{pl:eq-qbin}
\sum_{U\subseteq V}(-X)^{\dim U}q^{\binom{\dim U}2}=\prod_{l<d}(1-q^lX),
\end{equation}
by induction on $d$: of the subspaces of dimension $i$, those in a fixed hyperplane $H$ are
counted by the case $d-1$, each subspace of dimension $i-1$ of $H$ is $U\cap H$ for exactly
$q^{d-i}$ subspaces $U\not\subseteq H$ of dimension $i$, and $\binom i2-\binom{i-1}2=i-1$; so the
sum for $V$ is $(1-q^{d-1}X)$ times the sum for $H$. Hence, with
$r=\operatorname{rk}(T\bmod\varpi)$, $P'(X)\plone=\prod_{l<n'-r}(1-q^l\alpha X)\plone$ and
$P^-(Y)\plone=\prod_{l<n-r}(1-q^l\beta Y)\plone$ as functions of $T$, and $n'-r=k_0+(n-r)$.
\end{proof}

\begin{proposition}[One lift]\label{pl:onelift}
Let $\Phi=\plone\otimes v$ be a datum that is the vacuum at $\mathfrak l,\aubar{\mathfrak l}$, and
let $\varphi$ be a smooth form on $X_j$ fixed by $K$. Then $\Theta^X_j(\Phi)\varphi$ is fixed by
$K'$, and with $z=q^{1/2}$ and $\alpha$ as in Lemma~\ref{pl:vac},
\[
P_{\mathfrak l}(X)\,\Theta^X_j(\Phi)\varphi=\prod_{l<t+1}(1-q^l\alpha X)\cdot\Theta^X_j(\Phi)\bigl(
P_{\mathfrak l}(z^{\,t+1}X)\varphi\bigr).
\]
\end{proposition}
\begin{proof}
By \eqref{pl:eq-equiv}, $R_{h'}\Theta^X_j(\Phi)\varphi=\Theta^X_j(\omega(1,h')\Phi)\varphi$; this
gives the first claim and $T_{\mathfrak l,i}\Theta^X_j(\Phi)\varphi=\Theta^X_j(T'_i\Phi)\varphi$.
Also by \eqref{pl:eq-equiv}, $\Theta^X_j(\omega(g,1)\Phi)\varphi=\Theta^X_j(\Phi)(R_{g^{-1}}
\varphi)$, and $\Theta^X_j(\Phi)\circ R_k=\Theta^X_j(\Phi)$ for $k\in K$. Therefore, with $dy$ the
Haar measure of $\auGL_n(E_{\mathfrak l})$ giving $K$ the measure $1$,
\[
\Theta^X_j(T^-_i\Phi)\varphi=\Theta^X_j(\Phi)\Bigl(\sum_{x\in K\backslash Kd_iK}\int_KR_{kx}\varphi
\,dk\Bigr)=\Theta^X_j(\Phi)\Bigl(\int_{Kd_iK}R_y\varphi\,dy\Bigr)=\Theta^X_j(\Phi)(T_{\mathfrak l,i}
\varphi).
\]
Apply $\Theta^X_j(\,\cdot\,)\varphi$ to Lemma~\ref{pl:vac}; $q^{k_0}\alpha/\beta=q^{k_0/2}$.
\end{proof}

From now on $p$ is an exceptional prime, $\Lambda$ a full mark, $\plr=\plr_p$, $\vartheta$ a twist
with $p\nmid\auNm Q_\vartheta$, $\Psi$ a chain datum and $\Psi'$ its expansion
({\cite[Proposition~14.25]{P46}}). Every $x$ in the support of $\Psi'$ is integral above $p$, since
$p\nmid MN\auNm Q_\vartheta$; so $\phi^\circ_{j,x_j}$ is the vacuum at all places above $p$, and
\begin{equation}\label{pl:eq-fix}
\omega_{\mu_{j+1}}(k,k')\,\phi^\circ_{j,x_j}=\phi^\circ_{j,x_j}\qquad\text{for }k,k'\text{ in }
K^\circ\text{ and supported above }p,
\end{equation}
by Lemma~\ref{ps:weil}(W4) ($\omega(k,k')$ fixes the vacuum of a lattice that it stabilises, and
commutes with the translation by the part of $x_j$ away from $p$) and because $\mu_{j+1}$ is
unramified above $p$. We recall from {\cite[\S2.7]{P50}}, for $\mathfrak l\in\Lambda$, $\mathfrak
l=\mathfrak l_{\tau_0}$ with $\tau_0\in\auPa$: $p$ is a uniformiser at $\mathfrak l$; the places
inducing $\mathfrak l$ or $\aubar{\mathfrak l}$ are the $u((-p)^i)\tau_0$, those with even $i$
inducing $\mathfrak l$; the order $o$ of $-p$ modulo $q_1$ is even; and the residue field of
$\mathfrak l$ has $q=p^{o}$ elements, $o\le f$.

\begin{lemma}[Balance]\label{pl:balance}
Let $\mathfrak l\in\Lambda$. No place inducing $\mathfrak l$ or $\aubar{\mathfrak l}$ is a
$b$-place; so $\kappa(\tau)=a/x_1$, $\kappa_i(\tau)$ and $e_{i,\tau}$ are the same for all these
places, and as many of them induce $\mathfrak l$ as $\aubar{\mathfrak l}$. Consequently
$\iota(\chi_{i,\mathfrak l}(p))$ is a unit for $1\le i\le m$. Moreover, for odd $e$ the place
$u(-p^e)\tau$ induces $\aubar{\mathfrak l_\tau}$.
\end{lemma}
\begin{proof}
The classes modulo $q_0$ of these places are powers of $-p$, which has odd order modulo $q_0$; so
none is $-1$, the class of the $b$-places. The second assertion holds because $o$ is even, and the
third follows from Lemma~\ref{pl:char}. The last is the description just recalled, as $-p^e=
(-p)^e$.
\end{proof}

\begin{corollary}[The chain]\label{pl:chainT}
Let $\mathfrak l\in\Lambda$ and $\nu_p=f\,\plr\,(n_m+\plr)$. For $x$ in the support of $\Psi'$ and
$0\le e\le m$ let $C_{x,e}=\Theta^X_{e-1}(\phi^\circ_{e-1,x_{e-1}})\cdots\Theta^X_0(\phi^\circ_{0,
x_0})$. For each $1\le e\le m$ there are $c_i\in\pao$ $(0\le i\le\plr)$, independent of $x$, such
that for every smooth form $\varphi$ on $X_0$ fixed by $\auGL_k(\auO_{\mathfrak l})$,
\[
\sum_{i\le\plr}c_i\,T_{\mathfrak l,i}\,C_{x,e}\varphi=p^{\,\nu_p}\,C_{x,e}\bigl(T_{\mathfrak l,
\plr}\varphi\bigr).
\]
\end{corollary}
\begin{proof}
\emph{One lift.} In Proposition~\ref{pl:onelift}, $R(X)=\prod_{l<t+1}(1-q^l\alpha X)$ has
coefficients in $\pao$ by Lemma~\ref{pl:balance} and $R(0)=1$. Dividing by $R$ in $\pao[[X]]$ and
comparing the coefficients of $X^i$: for $i\le n_j$ there are $c_{i,i'}\in\pao$, depending only on
$j$ and $\mathfrak l$, with
\[
z^{(t+1)i}q^{\binom i2}\,\Theta^X_j(\Phi)(T_{\mathfrak l,i}\varphi)=\sum_{i'\le i}c_{i,i'}\,
T_{\mathfrak l,i'}\,\Theta^X_j(\Phi)\varphi .
\]
Here $z^{(t+1)i}q^{\binom i2}=p^{\,oi(t+i)/2}$ divides $p^{\,f\plr t}$ for $i\le\plr$, because
$o\le f$ and $\plr\le t$. \emph{Induction on $e$:} we claim $\sum_{i\le\plr}c^{(e)}_iT_{\mathfrak
l,i}C_{x,e}\varphi=p^{\,ef\plr t}C_{x,e}(T_{\mathfrak l,\plr}\varphi)$ with $c_i^{(e)}\in\pao$.
For $e=0$ take $c^{(0)}_i=\delta_{i\plr}$. The form $C_{x,e}\varphi$ is fixed by the integral
points at $\mathfrak l$ (Proposition~\ref{pl:onelift}), so applying $p^{f\plr t}\Theta^X_e
(\phi^\circ_{e,x_e})$ to the identity for $e$ and using the case of one lift for each $i\le\plr\le
k$ gives the identity for $e+1$. Finally multiply by $p^{\,\nu_p-ef\plr t}$, an integer because
$e\,t\le m\,t\le n_m$.
\end{proof}

\subsection{A lift at a neighbour; proof of Lemma~\ref{pa:hecke}}\label{pl:ss-hecke}
Let $\plE=\{\tau:\mathfrak l_\tau\in\Lambda\}$ and $\aubar\plE=\{\tau:\aubar{\mathfrak l_\tau}\in
\Lambda\}$; they are disjoint and $\Phi_a\subseteq\plE\cup\aubar\plE$, $\Lambda$ being a full
mark. With the numbers $a'_\tau$ of \eqref{pa:eq-diag} for the $j$-th lift put
\[
\mathrm L_j=\plr\sum_{\tau\in\plE}a'_\tau-D_j(\plE),\qquad\sum_{\mathfrak l\in\Lambda}A_{\mathfrak l}=a\cdot
\#(\auPa\cap\plE).
\]

\begin{lemma}[One lift at a neighbour]\label{pl:neigh}
Assume Lemma~\ref{pa:kernel}. Let $0\le j<m$, let $x$ be in the support of $\Psi'$, and let $h\in
\auU(X_{j+1})(\auA_{E^+,f})$ be supported on $\Lambda\cup\aubar\Lambda$ with component in
$K'd_{i_{\mathfrak l}}K'$, $i_{\mathfrak l}\le\plr$, at each $\mathfrak l\in\Lambda$. Then for all
$g,g'$ and $\xi$,
\[
\plgb{g'}{\theta_j\bigl(\omega_{\mu_{j+1}}(1,h)\phi^\circ_{j,x_j};g,g'\bigr)\xi}\ \le\ p^{\,\sum_{
\mathfrak l\in\Lambda}A_{\mathfrak l}+\mathrm L_j}\ \plgb g\xi .
\]
\end{lemma}
\begin{proof}
\emph{Reduction to a diagonal element.} By the theory of elementary divisors $h=u\,d\,u'$ with
$u,u'$ integral and supported on $\Lambda\cup\aubar\Lambda$, and $d$ diagonal: at $\mathfrak
l\in\Lambda$ its entries are $p$ on a set $J(\mathfrak l)\subseteq[n_{j+1}]$ of $i_{\mathfrak l}$
indices and $1$ elsewhere, and at $\aubar{\mathfrak l}$ they are the inverses. By (T2), the
left side concerns $\theta_j(\phi^\circ;g,g'h)$, where $\phi^\circ=\phi^\circ_{j,x_j}$. By
Corollary~\ref{pl:wacor}(a) there are $\gamma\in\Gamma$, $\varepsilon\in\Gamma'$ such that
$g_0=\gamma g$, $g_1=\varepsilon g'ud$ and $dg_1d^{-1}$ are integral above $p$. By (T1),
$\theta_j(\phi^\circ;g,g'h)=\varepsilon^{-1}\theta_j(\phi^\circ;g_0,g_1u')\gamma$. Let $\dot g_0$,
$\dot g_1$ be $g_0$, $g_1u'$ with their components above $p$ replaced by $1$. By (T2) and
\eqref{pl:eq-fix}, $\theta_j(\phi^\circ;g_0,g_1u')=\theta_j(\phi^\circ;\dot g_0,\dot g_1)=:\theta$.
Since $(\varepsilon g')_p^{-1}=u_p\,d_p\,g_{1,p}^{-1}$, and $\rho(d)\rho(g_1)^{-1}=\rho(dg_1^{-1}
d^{-1})\rho(d)$, (G1) and (G2) give
\[
\plgb{g'}{\theta_j(\phi^\circ;g,g'h)\xi}=\bigl\|\rho(d_p)\,\theta\,\rho(g_{0,p})\,\eta\bigr\|,
\qquad\eta=\rho(g_{0,p})^{-1}\iota(\gamma\xi),\quad\|\eta\|=\plgb g\xi .
\]
The components of $\dot g_0,\dot g_1$ above $p$ are trivial. Hence, by the definition
\eqref{ps:eq-kernel} of the kernel, the matrix of $\theta$ in the bases $(e_T)$ is the matrix
$(J_{T',T})=(S_{\lambda_{j+1}}(\partial)_{T',T}F(0))$ of the jets of the series $F=\Theta_f(\pi_M
s_0\,\cdot)$, $f=\omega_{\mu_{j+1}}(\dot g_0,\dot g_1)\phi^\circ$, itself, with no further factors:
for such points one takes $\gamma=\gamma'=1$ in the proof of Proposition~\ref{pk:K12}, so that the
two invertible matrices of (K1), which are the matrices of the components above $p$, are the
identity. (This is needed here, because $\rho(d_p)$ is not integral and could not be moved across an
arbitrary invertible matrix.) The datum $f$ is a unit multiple of a standard datum at $p$ by
Proposition~\ref{pk:std}(S1), so $F$ has the normal form by Proposition~\ref{pk:std}(S3) and
Proposition~\ref{pa:product}, as in (K2). As $\rho(g_{0,p})$ is integral and $\rho(d_p)$
is diagonal, it suffices to show, for all standard $T'$ and all $T$,
\begin{equation}\label{pl:eq-entry}
v\bigl(\rho(d_p)_{T'}\bigr)+v(J_{T',T})\ \ge\ -\sum_{\mathfrak l\in\Lambda}A_{\mathfrak l}-\mathrm L_j ,
\end{equation}
where $\rho(d_p)_{T'}$ is the eigenvalue \eqref{pl:eq-diag} on $e_{T'}$.

\emph{The eigenvalue.} At $\tau$ the action is through the component at $w_\tau=\aubar{\mathfrak
l_\tau}$. So $\hat\iota_\tau(d_{w_\tau})$ is diagonal with entries $p^{-1}$ on $J(\mathfrak l_\tau)$
if $\tau\in\plE$, $p$ on $J(\aubar{\mathfrak l_\tau})$ if $\tau\in\aubar\plE$, and it is $1$ for
the other $\tau$. Call an index $b$ \emph{marked at $\tau$} if it lies in this set $J$. The
determinants contribute $\sum_{\mathfrak l\in\Lambda}\#J(\mathfrak l)\bigl(\sum_{\tau:\aubar{
\mathfrak l_\tau}=\mathfrak l}e_{j+1,\tau}-\sum_{\tau:\mathfrak l_\tau=\mathfrak l}e_{j+1,\tau}
\bigr)=0$ to the valuation, by Lemma~\ref{pl:balance}. Hence
\[
v\bigl(\rho(d_p)_{T'}\bigr)=\#\{\text{boxes of }T'\text{ at }\aubar\plE\text{ with marked entry}\}-
\#\{\text{boxes of }T'\text{ at }\plE\text{ with marked entry}\}.
\]
\emph{The jets.} Define charges: $w_{\tau,h,b}=1$ if $\tau\in\plE$, $h=1$ and $b$ is not marked at
$\tau$, or if $\tau\in\aubar\plE$ and $b$ is marked at $\tau$; and $w_{\tau,h,b}=0$ otherwise.
They satisfy \eqref{pl:eq-chargehyp} for $\plE$ (an exceptional prime $p$ is inert in $K$, so
$p\ne d$ and $\paI=I_p$ is the set of odd positive integers): let $e\in I_p$, $\tau'=u(-p^e)\tau$; by
Lemma~\ref{pl:balance} one of $\tau,\tau'$ is in $\plE$ exactly when the other is in $\aubar\plE$,
and then the same indices are marked at both. If $\tau\in\plE$: a marked $b$ has $w_{\tau',h',b}=
1$ and an unmarked one has $w_{\tau,1,b}=1$. If $\tau'\in\plE$ and $h'=1$: a marked $b$ has
$w_{\tau,1,b}=1$ and an unmarked one has $w_{\tau',1,b}=1$. By Lemma~\ref{pa:schur} and
Proposition~\ref{pl:charges}, $v(J_{T',T})\ge D_j(\plE)-w(T')$.

\emph{Count of boxes.} A box at $\aubar\plE$ with marked entry counts $+1$ in $v(\rho(d_p)_{T'})$
and $1$ in $w(T')$; a box at $\plE$ in a column of height one counts $-1$ or $0$ in the first and
$0$ or $1$ in the second, according as its entry is marked or not; in a column of height $h\ge2$
at $\plE$ the entries are distinct, so at most $\plr$ of them are marked. Therefore
\[
v\bigl(\rho(d_p)_{T'}\bigr)-w(T')\ \ge\ -\sum_{\tau\in\plE}\bigl(a_\tau+\plr\,a'_\tau\bigr)=-\sum_{
\mathfrak l\in\Lambda}A_{\mathfrak l}-\plr\sum_{\tau\in\plE}a'_\tau ,
\]
which with the bound for the jets is \eqref{pl:eq-entry}.
\end{proof}

\begin{lemma}\label{pl:arith}
For $0\le j<m$: $\ \#\Lambda\cdot\nu_p+\mathrm L_j\le(2\paloss_p-\augain_p)/m$.
\end{lemma}
\begin{proof}
Recall that $p$ is inert in $K$, being exceptional, so $\paI=I_p$ is the set of odd positive
integers. There is no pair $(\tau,e)$ in $C_p(\plE)$: if $\tau\in\auPa\setminus\plE$ then $\tau\in\aubar
\plE$, and $u(-p^e)\tau\in\plE$ by Lemma~\ref{pl:balance}. So Lemma~\ref{pl:gain} gives
$D_j(\plE)\ge\augain_p/m+\#(\auPa\cap\plE)a/x_1-\paloss_p/m$ (here $\epsilon_m(p)=1$). It remains
to show
\[
\#\Lambda\,\nu_p+\plr\sum_{\tau\in\plE}a'_\tau\ \le\ \frac{\paloss_p}m+\#(\auPa\cap\plE)\frac
a{x_1}.
\]
By \eqref{pa:eq-diag} and Lemma~\ref{pl:balance}, $a'_\tau\le a/x_1+s$ for $\tau\in\plE$,
$s=x^{60}$. If $p\le m$, then $\plr(a/x_1+s)\le a/(pm)$, because $\plr\le2\log m$, $pm\le4x$ and
$x_1\ge x^{10^6}$. If $p>m$, then $\plr=1$, so the order of $-p$ modulo $q_0$ is $1$, all places of
$\plE$ are $a$-places, and $s\le a/(pm)$. In both cases $\plr\sum_{\tau\in\plE}a'_\tau\le
fa/(pm)+\#(\auPa\cap\plE)a/x_1$. Finally $\#\Lambda\,\nu_p\le f^2\plr(n_m+\plr)\le fa/(pm)$ and
$2fa/(pm)\le fa(\log b)^2/(pm)=\paloss_p/m$, by a comparison of sizes ($p<x_1$, $m\le2\sqrt x$,
$f\le x^{12}$, $n_m\le40x^6$, $a\ge x_1^{10^6}$).
\end{proof}

\begin{proof}[Proof of Lemma~\ref{pa:hecke}, given Lemma~\ref{pa:kernel}]
Let $\mathcal V_j$ be the space of the data of the $j$-th lift fixed by $\omega(1,K'_{\mathfrak
l})$ for all $\mathfrak l\in\Lambda$. For $\mathfrak l\in\Lambda$ let $c^{(j+1,\mathfrak l)}_i$
be the numbers of Corollary~\ref{pl:chainT} for $e=j+1$, and on $\mathcal V_j$ put
\[
t_{j,\mathfrak l}=p^{\,A_{\mathfrak l}-\nu_p}\sum_{i\le\plr}c^{(j+1,\mathfrak l)}_i\,T'_{\mathfrak
l,i},\qquad t_j=\prod_{\mathfrak l\in\Lambda}t_{j,\mathfrak l},
\]
with $T'_{\mathfrak l,i}$ as in \S\ref{pl:ss-transfer}; these operators preserve $\mathcal V_j$
and commute.

\emph{{\rm(i)} $\Theta^X_j(t_j\phi^\circ_{j,x_j})\,C_{x,j}\,y=C_{x,j+1}(Y_\Lambda y)$ for
$y\in\auS$ and $x$ in the support of $\Psi'$.} For $\Phi\in\mathcal V_j$ and any $\psi$,
$\Theta^X_j(T'_{\mathfrak l,i}\Phi)\psi=T_{\mathfrak l,i}\Theta^X_j(\Phi)\psi$ by
\eqref{pl:eq-equiv}. So if $\Theta^X_j(\Phi)C_{x,j}y=C_{x,j+1}(y')$ with $y'\in\auS$, then by
Corollary~\ref{pl:chainT} ($y'$ is fixed by $\auGL_k(\auO_{\mathfrak l})\subseteq J_X$)
\[
\Theta^X_j(t_{j,\mathfrak l}\Phi)\,C_{x,j}\,y=p^{\,A_{\mathfrak l}-\nu_p}\sum_ic_i\,T_{\mathfrak
l,i}\,C_{x,j+1}(y')=C_{x,j+1}\bigl(p^{A_{\mathfrak l}}T_{\mathfrak l,\plr}\,y'\bigr)=C_{x,j+1}
(Y_{\mathfrak l}\,y').
\]
Start from $\Phi=\phi^\circ_{j,x_j}$, $y'=y$, and apply this for the primes of $\Lambda$ in turn;
the operators $Y_{\mathfrak l}$ preserve $\auS$.

\emph{{\rm(ii)}} By (i) and induction on $j$, $\Theta^X_{j-1}(t_{j-1}\phi^\circ_{j-1,x_{j-1}})
\cdots\Theta^X_0(t_0\phi^\circ_{0,x_0})\,\varphi=C_{x,j}(Y_\Lambda^{\,j}\varphi)$. Hence
\[
C_{\vartheta,\Psi'}(Y_\Lambda^{\,m}\varphi)=\sum_x\Psi'(x)\,\Theta^X_{m-1}(t_{m-1}\phi^\circ_{m-1,
x_{m-1}})\cdots\Theta^X_0(t_0\phi^\circ_{0,x_0})\,\varphi .
\]
\emph{{\rm(iii)} The tensor $\sum_x\Psi'(x)\bigotimes_jt_j\phi^\circ_{j,x_j}$ is invariant under
the levels.} The tensor $\sum_x\Psi'(x)\bigotimes_j\phi^\circ_{j,x_j}$ is (proof of
{\cite[Proposition~14.25]{P46}}) and lies in $\bigotimes_j\mathcal V_j$ by \eqref{pl:eq-fix}. For
$b=(b_j,b_{j+1})$ in the levels, $\omega(b)$ preserves $\mathcal V_j$ and $\omega(b)T'_{\mathfrak
l,i}=T'_{\mathfrak l,i}\omega(b)$ on $\mathcal V_j$: the part of $b_{j+1}$ away from $\mathfrak l$
commutes with the elements $h$, and conjugation by its component at $\mathfrak l$, which lies in
$K'_{\mathfrak l}$, permutes the cosets $hK'_{\mathfrak l}$ of the double coset.

\emph{{\rm(iv)} Conclusion.} By (iii), Lemma~\ref{pl:sum} applies. The datum
$t_j\phi^\circ_{j,x_j}$ is $p^{\sum_{\mathfrak l}A_{\mathfrak l}-\#\Lambda\nu_p}$ times a
combination with coefficients in $\pao$ of data $\omega(1,h)\phi^\circ_{j,x_j}$ with $h$ as in
Lemma~\ref{pl:neigh}. By that lemma and the ultrametric inequality,
\[
\plgb{g'}{\theta_j(t_j\phi^\circ_{j,x_j};g,g')\xi}\le p^{\,\#\Lambda\nu_p+\mathrm L_j}\,\plgb g\xi .
\]
The numbers $\Psi'(x)$ are integral at $p$. As in the proof of Lemma~\ref{pa:chain}, the formula of
Lemma~\ref{pl:sum} and the restriction to $Y_m$ give, for $\varphi\in\auS_\iota$, that
$B_{\vartheta,\Psi}(Y_\Lambda^{\,m}\varphi)$ has gauge at most $p^{\sum_{j<m}(\#\Lambda\nu_p+\mathrm L_j)}
\le p^{\,2\paloss_p-\augain_p}$, by Lemma~\ref{pl:arith}.
\end{proof}

\subsection{The level prime}\label{pl:ss-level}
Let $\ell=N$ and $\iota\colon\C\to\C_N$. Let $\plP$ be the set of the $f$ primes of $E$
dividing $\pi_N$; the primes above $N$ are those of $\plP$ and of $\aubar\plP$.

\emph{The places $\plSi$.} The primes $\mathfrak l_\tau$ $(\tau\in\auPl)$ lie above one prime
$\mathfrak n$ of $K$, and $\mathfrak n\ne\aubar{\mathfrak n}$; as $N$ splits completely, $\tau
\mapsto\mathfrak l_\tau$ is a bijection of $\auPl$ onto the primes above $\mathfrak n$. Let
\[
\plSi=\{w_\tau=\aubar{\mathfrak l_\tau}:\tau\in\auPl\},
\]
the primes above $\aubar{\mathfrak n}$: it contains one prime of each pair $w,\aubar w$ above $N$;
for $w\in\plSi$ we write $\tau_w$ for the place with $w_{\tau_w}=w$.
Put $E_\Sigma=\prod_{w\in\plSi}E_w=\Q_N^{\plSi}$, $\auO_\Sigma=\Z_N^{\plSi}$, $\auF=\auO_\Sigma/N=
\auF_N^{\plSi}$. By (N2) the component at $\plSi$ identifies $\auU(X_j)(E^+\otimes\Q_N)$ with
$G_j=\auGL_{n_j}(E_\Sigma)$, and $g\in G_j$ acts on $W_{\lambda(X_j)}\otimes\C_N$ by $\bigotimes_\tau
\rho_\tau(g_{w_\tau})$. We write $e$ for the last basis vector, of any $X_j$.

\emph{Lattices.} For $S\subseteq\plSi$ let $\plD_S\in G_j$ be $\operatorname{diag}(1,\dots,1,N)$ at
the places of $S$ and $1$ at the others, and
\[
\plD=\plD_{\plSi},\qquad\pld_j=\plD_{\plSi\cap\plP},\qquad\ple_j=\plD_{\plSi\setminus\plP},\qquad
\plD=\pld_j\ple_j .
\]
The \emph{lattice of $X_j$ at $N$} is $\mathcal L_j=\pld_j\auO_\Sigma^{\,n_j}$, the integral vectors whose
last coordinate is divisible by $\pi_N$ (at $w\in\plSi\cap\plP$ the number $N$ is a unit times
$\pi_N$, at the other $w$ the number $\pi_N$ is a unit); it is \emph{marked} by the class
$\varepsilon_j$ of $\pld_je$ in $\mathcal L_j/N\mathcal L_j$; and $\mathcal L_j^+=\pld_j\plD^{-1}\auO_\Sigma^{\,n_j}=\mathcal L_j+
N^{-1}\auO_\Sigma\,\pld_je$. Let
\[
\plMir=\{k\in\auGL_{n_j}(\auO_\Sigma):k\,e\equiv e\pmod N\},
\]
the stabiliser of the standard lattice marked by $e$.

\begin{lemma}\label{pl:mir}
The group of the components at $\plSi$ of the elements of $J_{X_j}$ is $\pld_j\,\plMir\,
\pld_j^{-1}$, the stabiliser of $(\mathcal L_j,\varepsilon_j)$. For $k\in\plMir$, $\plD k\plD^{-1}$ is
integral with integral inverse and has last row $\equiv{}^te$ modulo $N$.
\end{lemma}
\begin{proof}
Above $N$, $J_{X_j}$ is defined by: $g$ integral with integral inverse, and the last row of $g_u$
is $\equiv{}^te$ modulo $u$ for $u\in\plP$. For $w\in\plSi\cap\plP$ this is a condition on $g_w$.
For $w\in\plSi\setminus\plP$ we have $\aubar w\in\plP$ and $g_{\aubar w}={}^t\aubar{g_w}^{-1}$, so
the condition is that the last column of $g_w^{-1}$, equivalently of $g_w$, is $\equiv e$ modulo
$w$; there $\pld_j=1$. For $w\in\plSi\cap\plP$, conjugation
$k\mapsto\operatorname{diag}(1,\dots,1,N)\,k\,\operatorname{diag}(1,\dots,1,N)^{-1}$
multiplies the last row by $N$ and divides the last column by $N$, outside the corner: it
exchanges ``integral with integral inverse, last row $\equiv{}^te$'' and ``integral with integral
inverse, last column $\equiv e$''. The last claim is the same computation.
\end{proof}

\emph{The representations $\pi_j$.} Let $\mathcal A(X_j)$ be the space of smooth forms on $X_j$ of
weight $\lambda(X_j)$. For $a\in G_j$, read as the element of $\auU(X_j)(\auA_{E^+,f})$ supported
above $N$ with these components at $\plSi$, put $\pi_j(a)=R_{\pld_ja\pld_j^{-1}}$ on $\mathcal
A(X_j)$ (``$a$ is read in a basis of $\mathcal L_j$''). By Lemma~\ref{pl:mir} the forms of level $J_{X_j}$
are those fixed by $\pi_j(\plMir)$ and by the part of $J_{X_j}$ away from $N$. A vector $\varphi$
of $\pi_j$ is \emph{ramified} if for each $w\in\plSi$ there is $x\in\auO_\Sigma^\times$ with
$x_{w'}\equiv1\pmod N$ for $w'\ne w$ and $\varphi\in(1-\pi_j(x\cdot1))\mathcal A(X_j)$. The forms
of $\auSram$ are ramified: for $w\in\plSi$ one of $w,\aubar w$ is a prime $u\in\plP$, and
$\varphi\in(1-z_u)\auS$, where $z_u$ is a scalar whose component at $w$ generates $\auF_N^\times$
and which is $1$ at the other places of $\plSi$.

\emph{The character $\psi_N$ and the operator $U_N$.} Let $\psi_N\colon\auF\to\C^\times$,
$\psi_N(c)=\lambda(-\tilde c/N)$ for a lift $\tilde c\in\auO_\Sigma$, with $\lambda$ as in
\eqref{pl:eq-heis} for $\mathsf S=\plSi$; it is a character, non-trivial on each factor $\auF_N$.
For $y=(y',u)\in\auF^{n_j}$ with last entry $u\in\auF^\times$ let $k_y\in\auGL_{n_j}(\auO_\Sigma)$
be a lift of
\[
\begin{pmatrix}1+u^{-1}y'\,{}^ty'&y'\\{}^ty'&u\end{pmatrix}=\begin{pmatrix}1&y'\\0&u\end{pmatrix}
\begin{pmatrix}1&0\\u^{-1}\,{}^ty'&1\end{pmatrix},
\]
whose last column and last row are both $y$. The operator on $\mathcal A(X_j)$
\begin{equation}\label{pl:eq-U}
U=U^{(j)}=\pi_j(\plD^{-1})\sum_{y\in\auF^{n_j},\ y_{\mathrm{last}}\in\auF^\times}\psi_N(y_{
\mathrm{last}})\,\pi_j(k_y),\qquad U_N:=U^{(0)},
\end{equation}
replaces $\mathcal L_j$ by $\mathcal L_j^+$ and sums against $\psi_N$ over the vectors of $\mathcal L_j/N\mathcal L_j$ with unit
coordinate along $\varepsilon_j$. On vectors fixed by $\pi_j(\plMir)$ the term of $y$ depends only
on $k_y\plMir$, that is, on $y$.

\begin{lemma}\label{pl:Uprop}
{\rm(a)} $U$ commutes with $\pi_j(x\cdot1)$, $x\in\auO_\Sigma^\times$, and with the right
translations by elements that are trivial above $N$; it preserves the ramified vectors.
{\rm(b)} $U$ preserves the vectors fixed by $\pi_j(\plMir)$. In particular $U_N\auS\subseteq\auS$
and $U_N\auSram\subseteq\auSram$.
\end{lemma}
\begin{proof}
(a) is clear. (b) Let $\varphi$ be fixed by $\plMir$ and $g\in\plMir$. By Lemma~\ref{pl:mir},
$k=\plD g\plD^{-1}$ is integral with last row $\equiv{}^te$. Then $kk_y$ has last column $ky$, and
$(ky)_{\mathrm{last}}=y_{\mathrm{last}}$; so $kk_y\in k_{ky}\plMir$ and
\[
\pi_j(g)U\varphi=\pi_j(\plD^{-1})\sum_y\psi_N(y_{\mathrm{last}})\pi_j(kk_y)\varphi=\pi_j(\plD^{-1})
\sum_y\psi_N((ky)_{\mathrm{last}})\pi_j(k_{ky})\varphi=U\varphi .
\]
For $\auSram$ use (a): $U_N$ commutes with the $z_u$.
\end{proof}

\emph{Data at $\plSi$ and the lift law.} We use (N5) for $\mathsf S=\plSi$. Let $x$ be a
translation tuple vanishing above $N$; then $\phi^\circ_{j,x_j}=\plone\otimes v_{x,j}$. For a
function $\Psi$ on $M_{n_{j+1}\times n_j}(E_\Sigma)$ put
\[
L_j(\Psi)=L^x_j(\Psi)=\Theta^X_j\bigl(\omega_{\mu_{j+1}}(\pld_j,\pld_{j+1})(\Psi\otimes v_{x,j})
\bigr)\colon\mathcal A(X_j)\to\mathcal A(X_{j+1}).
\]
We consider only functions $\Psi$ supported on the integral matrices and depending on $T$ modulo
$N$; with $e,e'$ the last basis vectors of $X_j$, $X_{j+1}$, and $T$ integral:
\[
\plcol_y(T)=[\,Te\equiv y\,],\quad\plrow(T)=[\,{}^te'T\equiv{}^te\,],\quad\psi_y(T)=\psi_N({}^te'Ty),
\quad\pldir=\plcol_{e'},
\]
for $y\in\auF^{n_{j+1}}$, respectively $y\in\auF^{n_j}$ in $\psi_y$. Thus $\pldir$, the
\emph{direct datum}, is the indicator of the integral $T$ with $Te\equiv e'$; read in bases of
$\mathcal L_j$ and $\mathcal L_{j+1}$, as in the definition of $L_j(\Psi)$, it is the indicator of the maps
$(\mathcal L_j,\varepsilon_j)\to(\mathcal L_{j+1},\varepsilon_{j+1})$ of marked lattices.

\begin{lemma}[Lift law]\label{pl:law}
For $g\in G_j$, $h\in G_{j+1}$ and $\Psi$, with $\chi_r$ as in \eqref{pl:eq-w5} for $\plSi$,
\[
L_j\bigl(\Psi(h^{-1}\,\cdot\,g)\bigr)=\frac{\chi_{n_{j+1}}(g)}{\chi_{n_j}(h)}\ \pi_{j+1}(h)\circ
L_j(\Psi)\circ\pi_j(g)^{-1}.
\]
The characters $\chi_r$ are trivial on integral elements. In particular, with $\gamma_j=\chi_{n_{
j+1}}(\plD)/\chi_{n_j}(\plD)=N^{f(t+1)/2}\in\Z$ (the first $\plD$ in $G_j$, the second in
$G_{j+1}$),
\begin{equation}\label{pl:eq-rowdir}
\pi_{j+1}(\plD^{-1})\,L_j(\plrow)=\gamma_j\,L_j(\pldir)\,\pi_j(\plD^{-1}).
\end{equation}
\end{lemma}
\begin{proof}
By \eqref{pl:eq-w5}, $\Psi(h^{-1}Tg)\otimes v=\chi_{n_{j+1}}(g)\chi_{n_j}(h)^{-1}\omega(g,h)(\Psi
\otimes v)$, and
\[
\omega(\pld_j,\pld_{j+1})\,\omega(g,h)=\omega\bigl(\pld_jg\pld_j^{-1},\pld_{j+1}h\pld_{j+1}^{-1}
\bigr)\,\omega(\pld_j,\pld_{j+1});
\]
conclude by \eqref{pl:eq-equiv}. The characters $\chi_{j+1,
w}$ are unramified at $w\mid N$. For \eqref{pl:eq-rowdir} take $g=\plD^{-1}$, $h=\plD^{-1}$:
$\plrow(\plD T\plD^{-1})=\pldir(T)$, as in the proof of Lemma~\ref{pl:mir}; and $\chi_r(\plD)=
\prod_w\chi_{j+1,w}(N)\cdot N^{fr/2}$, with $t+1$ even.
\end{proof}

\begin{lemma}[Calculus of direct data]\label{pl:calc}
Let $\varphi\in\mathcal A(X_j)$ be ramified. Write $L=L_j$, $L'=L_{j+1}$, $\pi=\pi_j$, etc.
\begin{enumerate}
\item[(a)] If a function $\mathsf f$ of $T$ modulo $N$ satisfies $\mathsf f(cT)=\mathsf f(T)$
for all $c\in\auF^\times$ that are $1$ outside one place $w$, then $L(\mathsf f)\varphi=0$. For
every $\mathsf f$, $L(\mathsf f)\varphi$ is ramified.
\item[(b)] $\pi'(\plD^{-1})L(\psi_e)\varphi=U'\bigl(L(\pldir)\varphi\bigr)=\gamma_j\,L(\pldir)
(U\varphi)$.
\item[(c)] $\sum_{y\in\auF^{n_{j+1}}}L'(\psi_y)L(\plcol_y)\varphi=\#\mathrm{Prim}\cdot L'(\psi_{e'})
L(\pldir)\varphi$, where $\mathrm{Prim}$ is the set of the $y$ that are non-zero at every place.
\item[(d)] $\pi''(\plD^{-1})\sum_yL'(\psi_y)L(\plcol_y)\varphi=\#\mathrm{Prim}\cdot\gamma_{j+1}
\gamma_j\cdot L'(\pldir)L(\pldir)(U\varphi)$.
\end{enumerate}
\end{lemma}
\begin{proof}
(a) Let $\varphi=\psi-\pi(x)\psi$ with $x$ as in the definition, for the place $w$. By
Lemma~\ref{pl:law} with $g=x^{-1}\cdot1$, $L(\mathsf f(\,\cdot\,x^{-1}))=L(\mathsf f)\pi(x)$, and
$\mathsf f(Tx^{-1})=\mathsf f(T)$; so $L(\mathsf f)\pi(x)\psi=L(\mathsf f)\psi$. For the second
claim, Lemma~\ref{pl:law} with $g=x$, $h=x$ gives $L(\mathsf f)\pi(x)=\pi'(x)L(\mathsf f)$.

(b) Sort $\psi_e$ by the corner entry $c$ of $T$: $\psi_e=\sum_{c\in\auF}\psi_N(c)[T_{\mathrm{
last,last}}\equiv c]$. If $c$ vanishes at a place $w$, its term is invariant as in (a) and
contributes nothing. Sorting the remaining terms by the last column,
\[
L(\psi_e)\varphi=\sum_{y_{\mathrm{last}}\in\auF^\times}\psi_N(y_{\mathrm{last}})\,L(\plcol_y)
\varphi,\qquad\plcol_y(T)=\pldir(k_y^{-1}T),\qquad L(\plcol_y)=\pi'(k_y)L(\pldir),
\]
by Lemma~\ref{pl:law}; applying $\pi'(\plD^{-1})$ gives $U'(L(\pldir)\varphi)$. Sorting instead by
the last row, $L(\psi_e)\varphi=\sum_y\psi_N(y_{\mathrm{last}})L([{}^te'T\equiv{}^ty])\varphi$, and
$[{}^te'T\equiv{}^ty]=\plrow(Tk_y^{-1})$ because the last row of $k_y$ is ${}^ty$; so
$L([{}^te'T\equiv{}^ty])=L(\plrow)\pi(k_y)$ and, by \eqref{pl:eq-rowdir},
\[
\pi'(\plD^{-1})L(\psi_e)\varphi=\gamma_j\,L(\pldir)\,\pi(\plD^{-1})\sum_y\psi_N(y_{\mathrm{last}})
\pi(k_y)\varphi=\gamma_jL(\pldir)(U\varphi).
\]
(c) If $y$ vanishes at a place, $\plcol_y$ is invariant as in (a) and its term is zero. If
$y\in\mathrm{Prim}$, there is $k\in\auGL_{n_{j+1}}(\auO_\Sigma)$ with $ke'\equiv y$; then
$\plcol_y(T)=\pldir(k^{-1}T)$ and $\psi_y(R)=\psi_{e'}(Rk)$, so $L(\plcol_y)=\pi'(k)L(\pldir)$ and
$L'(\psi_y)=L'(\psi_{e'})\pi'(k)^{-1}$ by Lemma~\ref{pl:law}. (d) By (a), $L(\pldir)\varphi$ is
ramified; apply (b) to $L'$ and then (b) to $L$.
\end{proof}

\emph{The gluing.} The gluing of \S\ref{ps:ss-glue} is $\psGlue_i=\sum_{\mathsf y}\psrho(p_{\mathsf
y})\otimes\psrho(p'_{\mathsf y})$, with the Heisenberg operators $\psrho$ themselves (the points
$p_{\mathsf y}$, $p'_{\mathsf y}$ are supported on $\plP$, respectively $\aubar\plP$, where
$\psrho$ is given by \eqref{pl:eq-heis} without any choice); on data that are the vacuum above
$N$ the term of $\mathsf y$ depends only on $\mathsf y$ modulo $N$. In the coordinates $T$, the
point $p_{\mathsf y}$ has last column $-\mathsf y/N$ at the places of $\plP$, and $p'_{\mathsf y}$
has last row $-\mathsf y/N$ at the places of $\aubar\plP$; all other entries vanish.

\begin{lemma}[The glued double lift]\label{pl:glue}
Let $x$ vanish above $N$ and $0\le i<m/2$. Then
\[
\Theta^X_{2i+1}\Theta^X_{2i}\bigl(\psGlue_i(\phi^\circ_{2i,x_{2i}}\otimes\phi^\circ_{2i+1,x_{2i+1}})
\bigr)=\mathsf g_i\cdot\pi_{2i+2}(\plD^{-1})\circ\sum_{z\in\auF^{n_{2i+1}}}L^x_{2i+1}(\psi_z)\circ
L^x_{2i}(\plcol_z),
\]
where $\mathsf g_i=\chi^{(2i+2)}_{n_{2i+1}}(\ple_{2i+2})\cdot\chi^{(2i+1)}_{n_{2i+1}}(\pld_{2i})$,
the upper index being that of the character $\mu$.
\end{lemma}
\begin{proof}
We compute $\sum_{\mathsf y}\psrho(p_{\mathsf y})\plone\otimes\psrho(p'_{\mathsf y})\plone$ as a
function of $(T,R)$, at $\plSi$, by \eqref{pl:eq-heis}. At $w\in\plSi\cap\plP$: $\psrho(p_{\mathsf
y})\plone(T)=\plone(T+p_{\mathsf y})$ is $1$ if and only if $T\pld_{2i}$ is integral with last
column $\equiv\mathsf y$; and $\psrho(p'_{\mathsf y})$ multiplies $\plone(R)$ by $\lambda_w(-{}^te
R\mathsf y/N)=\psi_N({}^teR\mathsf y)$. At $w\in\plSi\setminus\plP$: $\aubar w\in\plP$, so
$\psrho(p_{\mathsf y})$ multiplies $\plone(T)$ by $\psi_N({}^t\mathsf yTe)$, and $\psrho(p'_{
\mathsf y})\plone(R)=\plone(R+p'_{\mathsf y})$ is $1$ if and only if $\ple_{2i+2}R$ is integral
with last row $\equiv{}^t\mathsf y$; the sum over $\mathsf y$ at this place is $\psi_N$ of the
corner entry of $(\ple_{2i+2}R)T$. In both cases the sum is
\[
\sum_{z\in\auF^{n_{2i+1}}}\plcol_z(T\pld_{2i})\,\psi_z(\ple_{2i+2}R).
\]
By \eqref{pl:eq-w5} and \eqref{pl:eq-equiv}, as in the proof of Lemma~\ref{pl:law},
\begin{align*}
\Theta^X_{2i}\bigl(\plcol_z(\,\cdot\,\pld_{2i})\otimes v\bigr)&=\chi^{(2i+1)}_{n_{2i+1}}(\pld_{2i})\,
R_{\pld_{2i+1}}^{-1}\,L_{2i}(\plcol_z),\\
\Theta^X_{2i+1}\bigl(\psi_z(\ple_{2i+2}\,\cdot\,)\otimes v'\bigr)&=\chi^{(2i+2)}_{n_{2i+1}}
(\ple_{2i+2})\,R_{\plD^{-1}}\,L_{2i+1}(\psi_z)\,R_{\pld_{2i+1}},
\end{align*}
because $\ple_{2i+2}\pld_{2i+2}=\plD$; and $R_{\plD^{-1}}=\pi_{2i+2}(\plD^{-1})$.
\end{proof}

\begin{proposition}[All lifts become direct]\label{pl:direct}
There is $c\in\C^\times$, algebraic, with $|\iota(c)|\le N^{mfb}$, such that for every $x$
vanishing above $N$ and every $\varphi\in\auSram$,
\[
C_{\psGlue(\phi^\circ_x)}\varphi=c\cdot L^x_{m-1}(\pldir)\circ\cdots\circ L^x_0(\pldir)\bigl(U_N^{\,
m/2}\varphi\bigr).
\]
\end{proposition}
\begin{proof}
Write $\mathsf L_j=L^x_j(\pldir)$. By Lemmas~\ref{pl:glue} and~\ref{pl:calc}(d), on ramified vectors of
$\pi_{2i}$ the $i$-th glued double lift is
\[
\mathsf g_i\cdot\#\mathrm{Prim}_i\cdot\gamma_{2i+1}\gamma_{2i}\cdot\mathsf L_{2i+1}\mathsf L_{2i}\,
U^{(2i)} .
\]
By Lemma~\ref{pl:calc}(a),(b), for ramified $\psi$ the vector $U^{(2i)}\mathsf L_{2i-1}\cdots\mathsf L_0
\psi$ is $\prod_{l<2i}\gamma_l$ times $\mathsf L_{2i-1}\cdots\mathsf L_0U^{(0)}\psi$; and $U^{(0)}$
preserves the ramified vectors. Induction on the number of double lifts gives the formula with
\[
c=\prod_{i<m/2}\Bigl(\mathsf g_i\cdot\#\mathrm{Prim}_i\cdot\prod_{l<2i+2}\gamma_l\Bigr).
\]
The $\gamma_l$ and $\#\mathrm{Prim}_i$ are integers. Let $w\in\plSi$; then $\mathfrak l_\tau=\aubar
w$ for $\tau=\tau_w$ only, and $\mathfrak l_\tau=w$ for no $\tau$, so by Lemma~\ref{pl:char}
$|\iota\chi^{(l)}_w(N)|\le N^{\kappa_l(\tau_w)}\le N^{2b}$ for $l\le m$. Each $w\in\plSi$ occurs
in exactly one of the two factors of $\mathsf g_i$, and the factors $|N|_w^{-n_{2i+1}/2}=N^{n_{2i+
1}/2}$ are integers. So $|\iota(\mathsf g_i)|\le N^{2fb}$ and $|\iota(c)|\le N^{2fb\cdot m/2}$.
The constant $c$ does not depend on $x$; it depends on the twist $\vartheta$, through the numbers
$\chi^{(l)}_w(N)$ in the $\mathsf g_i$ only (a twist is unramified above $N$ and of finite order, so
it changes these numbers by roots of unity), and the bound for $|\iota(c)|$, which is all that is
used below, is the same for all twists.
\end{proof}

\begin{lemma}[The kernels at $N$]\label{pl:kernelN}
Let $0\le j<m$ and let $z\in\auHom(X_j,X_{j+1})(\auA_{E^+,f})$ be integral at the places of
$\plSi$ and such that $Nz$ is integral at the places of $\aubar{\plSi}$ $($the primes $\mathfrak
l_\tau)$. Then for all $g,g'$ the matrix of $\rho(g'_N)^{-1}\circ\theta_j(\phi^\circ_{j,z};g,g')
\circ\rho(g_N)$ in the bases $(e_T)$ has entries in $\auO_\iota$.
\end{lemma}
\begin{proof} This is Corollary~\ref{pf:K1}(b): with $\delta\in\mathfrak n\setminus(\mathfrak n^2\cup\aubar{\mathfrak n})$,
$\mathfrak n$ the prime of $K$ induced by $\iota$, the datum $\phi^\circ_{j,z}$ is a standard datum of level $\delta$ at $N$;
$\iota(\aubar\delta)$ is a unit, and Theorem~\ref{pf:std} applies with $c=\iota(\pi_M/\aubar\delta)$.\end{proof}

\begin{lemma}[The prime $N$; closed form of Lemma~\ref{pa:level}]\label{pl:level}
Let $\ell=N$, $\iota\colon\C\to\C_N$, and let $\plSi$, $\mathcal L_0$, $\pld_0$ and $U_N$ be as above. Say
that $\psi\in\auS$ is \emph{integral for $\mathcal L_0$} if it has gauge at most $1$ at the base point
$\pld_0$, that is, $\iota(\psi(g))\in\rho(g_N\pld_0)W_{\lambda^X}(\auO_\iota)$ for all $g$. Let
$\mathcal A'=\{\varphi\in\auS:U_N^{\,m/2}\varphi\text{ is integral for }\mathcal L_0\}$. Then:
\begin{enumerate}
\item[(a)] for every twist $\vartheta$, every chain datum $\Psi$
and every $\varphi\in\mathcal A'\cap\auSram$, the form $B_{\vartheta,\Psi}\varphi$ has gauge at
most $N^{2mfb}$;
\item[(b)] $\aucol(\auS_\iota,\mathcal A'\cap\auS_\iota)\mid N^{\,i}$ for an integer $i\le
m\,f\,b\,\auDX$.
\end{enumerate}
\end{lemma}
\begin{proof}[Proof of {\rm(a)}]
\emph{Step 1.} By (C1) every $x$ in the support of $\Psi$ vanishes above $N$. Let $\psi=U_N^{\,
m/2}\varphi$; it lies in $\auS$ (Lemma~\ref{pl:Uprop}) and is integral for $\mathcal L_0$. By
Proposition~\ref{pl:direct},
\[
C_{\psGlue(\sum_x\Psi(x)\phi^\circ_x)}\varphi=c\sum_x\Psi(x)\,\Theta^X_{m-1}(A_{m-1}\phi^\circ_{m-1,
x_{m-1}})\cdots\Theta^X_0(A_0\phi^\circ_{0,x_0})\,\psi,\qquad A_j(\plone\otimes v)=\omega(\pld_j,
\pld_{j+1})(\pldir\otimes v).
\]
\emph{Step 2: the tensor $\sum_x\Psi(x)\bigotimes_jA_j\phi^\circ_{j,x_j}$ is invariant under the
levels.} First, $\sum_x\Psi(x)\bigotimes_j\phi^\circ_{j,x_j}$ is: let $\mathrm{pr}_j$ be the
projection of $\plH_f$ onto the vectors that are the vacuum above $N$ (the mean of $\psrho^\circ
(l)$ over the integral $l$ supported above $N$). It commutes with $\omega(b_j,b_{j+1})$ for $b$ in
the levels, which are integral; and $\bigotimes_j\mathrm{pr}_j$ maps $\psGlue(\phi^\circ_x)$ to
$\phi^\circ_x$, because $\mathrm{pr}(\psrho(p_{\mathsf y})\phi^\circ)=0$ unless $p_{\mathsf y}$ is
integral, that is, $\mathsf y\equiv0$. So the claim follows from (C4). Secondly, for $b_j\in
J_{X_j}$, $b_{j+1}\in J_{X_{j+1}}$ the operator $\omega(b_j,b_{j+1})$ commutes with $A_j$ on the
vectors that are the vacuum above $N$: by Lemma~\ref{pl:mir} the components at $\plSi$ of $\pld_j^{
-1}b_j\pld_j$ and $\pld_{j+1}^{-1}b_{j+1}\pld_{j+1}$ lie in $\plMir$, and $\pldir(k'^{-1}Tk)=
\pldir(T)$ for $k,k'\in\plMir$; use \eqref{pl:eq-w5}.

\emph{Step 3: the direct datum loses at most $N^{fn_{j+1}}$.} We claim that for $x$ vanishing
above $N$ and all $g,g'$ the matrix of $\rho(g'_N)^{-1}\theta_j(\pldir\otimes v_{x,j};g,g')\rho
(g_N)$ has entries of absolute value at most $N^{fn_{j+1}}$. Indeed, by Fourier inversion on
$\auF^{n_{j+1}}$ and \eqref{pl:eq-heis},
\[
\pldir(T)=N^{-fn_{j+1}}\sum_{\mathsf y\in\auF^{n_{j+1}}}\psi_N(-{}^t\mathsf ye')\,\psi_N({}^t\mathsf
yTe)\,\plone(T),\qquad\psi_N({}^t\mathsf yTe)\,\plone=\psrho(a_{\mathsf y})\plone,
\]
where $a_{\mathsf y}$ is supported on $\aubar{\plSi}$ and has last column $-\mathsf y/N$. So
$\pldir\otimes v_{x,j}$ is $N^{-fn_{j+1}}$ times a sum of roots of unity times data $\phi^\circ_{j,
x_j+a_{\mathsf y}}$, to which Lemma~\ref{pl:kernelN} applies.

\emph{Step 4.} By (T2), $\theta_j(A_j\phi^\circ_{j,x_j};g,g')=\theta_j(\pldir\otimes v_{x,j};g
\pld_j,g'\pld_{j+1})$; so by Step~3 this kernel maps vectors of gauge $C$ at $g\pld_j$ to vectors
of gauge at most $N^{fn_{j+1}}C$ at $g'\pld_{j+1}$. By Step~2 and Lemma~\ref{pl:sum}, and since the
$\Psi(x)$ are integral at $N$ and $\psi$ has gauge at most $1$ at $\pld_0$, the form
$F=\sum_x\Psi(x)\Theta^X_{m-1}(A_{m-1}\phi^\circ)\cdots\Theta^X_0(A_0\phi^\circ)\psi$ has gauge at
most $N^{f\sum_jn_{j+1}}\le N^{mfb}$ at the base point $\pld_m$ (comparison of sizes: $n_m\le b$).

\emph{Step 5: restriction.} $(\operatorname{res}F)(y)=\prod_\tau\tau(\aubar{\pi_N})^{-n_\tau}
\mathrm{br}(F(\operatorname{diag}(y,1)))$. The element $(\operatorname{diag}(y,1)\pld_m)^{-1}$ is
$\operatorname{diag}(y^{-1},s)$ with $s_\tau=N^{-1}$ if $w_\tau\in\plP$ and $s_\tau=1$ otherwise,
so by Lemma~\ref{ps:branching}(a),(c)
\[
\Bigl\|\prod_\tau s_\tau^{\,n_\tau}\cdot\rho^\circ(y_N)^{-1}\mathrm{br}\,F(\operatorname{diag}(y,1))
\Bigr\|=\bigl\|\mathrm{br}\,\rho(\operatorname{diag}(y,1)\pld_m)^{-1}F(\operatorname{diag}(y,1))
\bigr\|\le N^{mfb}.
\]
On the other hand $|\iota\tau(\aubar{\pi_N})|=N^{-1}$ if $\pi_N\in\aubar{\mathfrak l_\tau}=w_\tau$,
that is, if $w_\tau\in\plP$, and $=1$ otherwise; so $|\iota\tau(\aubar{\pi_N})^{-n_\tau}|=
|s_\tau^{\,n_\tau}|$ for every $\tau$, and $\operatorname{res}F$ has gauge at most $N^{mfb}$ at
the base point $1$. With $|\iota(c)|\le N^{mfb}$ this proves (a).
\end{proof}

\begin{proof}[Proof of {\rm(b)}]
Let $\pld=\pld_0$, $\pld^+=\pld_0\plD^{-1}=\ple_0^{-1}$, and let $\mathcal A^L$, $\mathcal A^+$ be
the modules of the $\psi\in\auS$ with gauge at most $1$ at the base points $\pld$, $\pld^+$
(integral for $\mathcal L_0$, for $\mathcal L_0^+$).

\emph{{\rm(i)} $\mathcal A^L\subseteq\auS_\iota\subseteq\mathcal A^+$}, because $\rho(\pld)$ and
$\rho(\pld^+)^{-1}=\rho(\ple_0)$ have integral entries: the weight $\lambda^X$ has no power of the
determinant.

\emph{{\rm(ii)} $\aucol(\mathcal A^+,\mathcal A^L)\mid N^{\,l}$ with $l\le\auDX fa/k$.} By
Corollary~\ref{pl:wacor}(a) (for $g=g_0\pld^+$, $g_0$ any representative of the class, $d=\plD^{-1}$;
put $r=\gamma g_0$) every class has a representative $r$ such that $r\pld^+$ and $\plD^{-1}r\pld$ are integral above $N$. For $b\in J_X$,
$\pld^{-1}b\pld\in\plMir$ and $(\pld^+)^{-1}b\pld^+=\plD(\pld^{-1}b\pld)\plD^{-1}$ are integral
above $N$ (Lemma~\ref{pl:mir}). So by (G1), (G2), for $\psi\in\auS$ and $g=\gamma rb$,
\[
\plgb{g\pld^+}{\psi(g)}=\|\iota\psi(r)\|,\qquad\plgb{g\pld}{\psi(g)}=\bigl\|\rho(\plD)^{-1}\iota
\psi(r)\bigr\| .
\]
Thus evaluation at the representatives identifies $\mathcal A^+$ with $\prod_rW_{\lambda^X}
(\auO_\iota)$ and $\mathcal A^L$ with $\prod_r\rho(\plD)W_{\lambda^X}(\auO_\iota)$. By
\eqref{pl:eq-diag}, $\rho(\plD)$ is diagonal with entries $N^{\deg_e\beta}$, where $\deg_e\beta$
is the total degree of the basis vector $e_\beta$, $\beta=(\beta_\tau)_{\tau\in\auPa}$, in the
last variable. By {\cite[Lemma~2.3(a)]{P50}} the colength divides $N^{\,l}$, $l=\#\{r\}\cdot\sum_\beta
\deg_e\beta$; by the symmetry of the $k$ variables the mean of $\deg_e\beta$ is $a\#\auPa/k$.

\emph{{\rm(iii)} $U_N\mathcal A^L\subseteq\mathcal A^+$.} For $\psi\in\mathcal A^L$, $U_N\psi\in
\auS$ and $(U_N\psi)(g)=\sum_y\psi_N(y_{\mathrm{last}})\psi(gg_y)$ with $g_y=\pld^+k_y\pld^{-1}$.
As $gg_y\pld=g\pld^+k_y$ and $k_y$ is integral, $\plgb{g\pld^+}{\psi(gg_y)}=\plgb{gg_y\pld}{\psi(g
g_y)}\le1$.

\emph{{\rm(iv)}} Let $\mathcal A'_i=\{\varphi\in\auS:U_N^{\,i}\varphi\in\mathcal A^L\}$. We show
$\aucol(\auS_\iota,\mathcal A'_i)\mid N^{\,l(i+1)}$ by induction. For $i=0$ this is (i), (ii). If
a family in $\auS_\iota$ has been carried into $\mathcal A'_i$ by a matrix of determinant dividing
$N^{l(i+1)}$, the image $(\varphi_k)$ has $U_N^{\,i+1}\varphi_k\in\mathcal A^+$ by (iii); by (ii)
a matrix $\mu$ with $\det\mu\mid N^{\,l}$ carries the family $(U_N^{\,i+1}\varphi_k)$ into
$\mathcal A^L$, that is, $(\varphi_k)$ into $\mathcal A'_{i+1}$. For $i=m/2$: the matrices are
integral, so the families are carried into $\mathcal A'\cap\auS_\iota$, and the exponent is
$l(m/2+1)\le\auDX fa(m/2+1)/k\le mfb\,\auDX$, because $a=xb$ and $x(m/2+1)\le mk$.
\end{proof}

\emph{Case (iii) of Theorem~\ref{pa:divis}.} With Lemma~\ref{pl:level} in place of
Lemma~\ref{pa:level}, the deduction in the proof of Theorem~\ref{pa:divis} stands: $S_0=\mathcal A'
\cap\auS_\iota$ has $\aucol(\auS_\iota,S_0)\mid N^{\,i}$, $i\le mfb\,\auDX\le\auerr_N\auDX$, the
form $B_{\vartheta,\Psi}\varphi$ has gauge at most $N^{2mfb}=N^{\auerr_N-\augain_N}$ for
$\varphi\in S_0\cap\auSram$, and the last paragraph of that proof gives $\aucol(\auS_\iota\cap
\auSram,S_0\cap\auSram)\mid N^{\,i}$. The hypotheses of case (iii) are: Lemma~\ref{pl:kernelN} and
the integrality of the branching map (Lemma~\ref{ps:branching}(c)); neither
{\cite[Proposition~14.25]{P46}} nor (K0)--(K2) is used.

\begin{remark}[Dependencies]\label{pl:status}
Lemma~\ref{pa:hecke} is proved from Lemma~\ref{pa:kernel} (K1), (K2),
{\cite[Proposition~14.25]{P46}}, the facts (W4), (W5) on the Weil representation in the form (N5),
the equivariance (T1) of the kernels under rational points (which is part of their definition as
maps to algebraic modular forms, \S\ref{ps:ss-weil}), and Lemma~\ref{ps:branching}(c) (through the
restriction, as in Lemma~\ref{pa:chain}). The transfer
of Hecke operators (Proposition~\ref{pl:onelift}) is the unramified theta correspondence for
$(\auGL_n,\auGL_{n'})$ at a split place, in the form of Rallis~\cite{RallisFunctoriality}; the
proof given is self-contained. Lemma~\ref{pa:level} is proved, in the form of
Lemma~\ref{pl:level}, from Lemma~\ref{pl:kernelN} and the same facts.
\end{remark}



\providecommand{\auO}{\mathcal O}
\providecommand{\auU}{\mathrm U}
\providecommand{\auGL}{\mathrm{GL}}
\providecommand{\auHom}{\operatorname{Hom}}
\providecommand{\auA}{\mathbb A}
\providecommand{\auF}{\mathbb F}
\providecommand{\auPl}{\Phi}
\providecommand{\auS}{\mathcal S}
\providecommand{\auNm}{\operatorname{Nm}}
\providecommand{\auAut}{\operatorname{Aut}}
\providecommand{\aubar}[1]{\overline{#1}}
\providecommand{\psM}{\mathsf M}
\providecommand{\psHS}{\mathcal H}
\providecommand{\psZ}{Z}
\providecommand{\psvac}{\mathbf v}
\providecommand{\psrho}{\varrho}
\providecommand{\pao}{\mathfrak o}
\providecommand{\pam}{\mathfrak m}
\providecommand{\pktr}{\operatorname{tr}}
\providecommand{\pfT}{\operatorname{Tr}}
\providecommand{\pfv}{\operatorname{ord}}
\providecommand{\pfbr}[1]{\theta[#1]}

\section{The reciprocity law on one curve and integrality at the remaining primes}\label{pf:sec}

This section follows \S\ref{pk:sec} and uses its notation and that of \S\ref{pa:ss-curve}. It
proves the following statements used there: Proposition~\ref{pk:curve}(C3)(ii)
(Theorem~\ref{pf:recip}); the integrality in Lemma~\ref{pa:kernel}(K1) at every prime $\ell\ne d$,
and Lemma~\ref{pl:kernelN} without any restriction (Theorem~\ref{pf:std},
Corollary~\ref{pf:K1}); the fact (W6) (Proposition~\ref{pf:flip}); the classification of orbits at
the end of Lemma~\ref{ps:support} (Lemma~\ref{pf:orbits}); the unit property at $2$ of the
Chowla--Selberg unit (Corollary~\ref{pf:CS}); and the bound for the number of classes used in
{\cite[Propositions~14.26 and~14.28]{P46}} (Proposition~\ref{pf:classes}). Two conventions, which matter
only at the prime $2$, are recalled in \S\ref{pf:ss-adm}: the point $t$ of order two is the
\emph{admissible} one, and the character $\epsilon_L$ of the vacuum is defined with the element
$\omega=(1+\sqrt{-d})/2$. Remark~\ref{pf:gen} records that the translates $\theta_u$ with
$u\in t+E_1$ are also $\partial$-integral.

Throughout, $\ell$ is a prime, $\iota\colon\C\to\C_\ell$ a ring homomorphism, $\pao$ the valuation
ring of $\C_\ell$, $\pam$ its maximal ideal; a complex number $c$ is \emph{integral} (a
\emph{unit}) if $\iota(c)\in\pao$ ($\in\pao^\times$), and $|c|$ means $|\iota(c)|$. We write
$\omega=(1+\sqrt{-d})/2$, so that $\auO_K=\Z[\omega]$, $\omega+\aubar\omega=1$ and
$\omega\aubar\omega=(1+d)/4$, an even number divisible by every prime dividing $q_1$; and
$E(u,v)=\Iml(\aubar uv)/A$, so that $H(u,v)-H(v,u)=2\boldsymbol\pi i\,E(u,v)$ and $E$ is integral
on $\Lambda$ and unimodular.

\subsection{The admissible point of order two}\label{pf:ss-adm}

Let $\Lambda$ be a lattice with $\auO_K\Lambda=\Lambda$. The prime $2$ splits in $K$,
$2\auO_K=\mathfrak p_2\aubar{\mathfrak p_2}$, and $\tfrac12\Lambda/\Lambda$ is free of rank one over
$\auO_K/2=\auF_2\times\auF_2$. Its three non-zero points are the generator of the part killed by
$\mathfrak p_2$, the generator of the part killed by $\aubar{\mathfrak p_2}$, and their sum. The
last one is the \emph{admissible point}: it is the only $t\in\tfrac12\Lambda\setminus\Lambda$ with
$\omega t\notin\Lambda$ and $\aubar\omega t\notin\Lambda$ (exactly one of $\omega,\aubar\omega$ lies
in each of the two primes). \emph{From now on a normalised lattice $(\Lambda,t,a_1)$ is one in which
$t$ is the admissible point}; this is needed at the prime $2$ only. For $\Lambda=\Omega_0\auO_K$ this fixes $\Omega_0$ up to a fourth root of unity.

\begin{lemma}\label{pf:chi}
Let $(\Lambda,t,a_1)$ be a normalised lattice and $\chi$ the function of \eqref{pa:eq-theta}.
\begin{enumerate}
\item[(a)] For $\lambda\in\Lambda$: $\chi(\lambda)=-1$ if $\lambda/2\in t+\Lambda$, and
$\chi(\lambda)=1$ otherwise. For any choice of $t\in\frac12\Lambda\setminus\Lambda$ this holds.
\item[(b)] $\chi(\aubar\delta\lambda)=\chi(\lambda)^{\auNm\delta}$ for all $\delta\in\auO_K$. If
$\Lambda=\Omega_0\auO_K$, then $\chi(\Omega_0a)=(-1)^{\auNm a}$ for $a\in\auO_K$.
\item[(c)] For $u\in\C$ and $\lambda\in\Lambda$, $\theta_u(z+\lambda)=\chi(\lambda)e^{2\boldsymbol
\pi iE(\lambda,u)}e^{H(\lambda,z+\lambda/2)}\theta_u(z)$ and $\theta_{u+\lambda}=\chi(\lambda)
e^{\boldsymbol\pi iE(\lambda,u)}\theta_u$. In particular $\theta_{u+2n\lambda}=\theta_u$ if
$n\ge1$ and $nu\in\Lambda$, and $\theta_\lambda=\chi(\lambda)\theta$.
\item[(d)] Let $\delta\in\auO_K$, $\delta\ne0$, and $\delta u\in\Lambda$. Then $z\mapsto\theta_u
(\aubar\delta z)$ satisfies $G(z+\lambda)=\chi(\lambda)^{m}e^{mH(\lambda,z+\lambda/2)}G(z)$
$(\lambda\in\Lambda)$ with $m=\auNm\delta$.
\end{enumerate}
\end{lemma}
\begin{proof}
(a) $\chi(\lambda)=\varepsilon(\lambda)e^{2\boldsymbol\pi iE(\lambda,t)}$. If $\lambda\in2\Lambda$
both factors are $1$. Otherwise $\varepsilon(\lambda)=-1$ and $e^{2\boldsymbol\pi iE(\lambda,t)}=
(-1)^{E(\lambda,2t)}$; the form $E$ modulo $2$ is the non-degenerate alternating form on
$\Lambda/2\Lambda\cong\auF_2^2$, so $E(\lambda,2t)$ is even exactly when $\lambda\equiv2t$ modulo
$2\Lambda$. (b) Identify $\Lambda/2\Lambda$ with $\auF_2\times\auF_2$ as an $\auO_K/2$-module so
that $2t=(1,1)$, and let $\aubar\delta$ act by $(a,b)$. By (a), $\chi(\aubar\delta\lambda)=-1$ if
and only if $a=b=1$ and $\lambda\equiv2t$; and $a=b=1$ means that $\auNm\delta$ is odd. For
$\Lambda=\Omega_0\auO_K$ the class of $\Omega_0a$ is $(1,1)$ exactly when $a\notin\mathfrak p_2
\cup\aubar{\mathfrak p_2}$, that is, when $\auNm a$ is odd. (c) is a direct computation from
\eqref{pa:eq-theta} and the definition of $\theta_u$; for the last assertions, $\chi(2n\lambda)=1$
by (a) and $E(2n\lambda,u)=2E(\lambda,nu)\in2\Z$. (d) By (c), $G(z+\lambda)=\chi(\aubar\delta
\lambda)e^{2\boldsymbol\pi iE(\aubar\delta\lambda,u)}e^{H(\aubar\delta\lambda,\aubar\delta z+
\aubar\delta\lambda/2)}G(z)$; here $E(\aubar\delta\lambda,u)=E(\lambda,\delta u)\in\Z$,
$H(\aubar\delta\lambda,\aubar\delta w)=\auNm\delta\,H(\lambda,w)$, and (b) applies.
\end{proof}

\begin{lemma}[Good reduction everywhere]\label{pf:two}
Let $(\Lambda,t,a_1)$ be a normalised lattice. Then $a_1$ is a unit at every place above $2$.
Consequently $a_1^4-64$ is a unit at every prime, the scheme $y^2+a_1xy=x^3+x$ over the ring of
integers of $\Q(a_1)$ is an elliptic curve, and at every $\iota$ the point $P_0=(0,0)$ does not
reduce to the origin.
\end{lemma}
\begin{proof}
Let $\ell=2$ and let $L\subset\C_2$ be a finite extension of $\Q_2$ over which $\mathcal E$ has
all its points of order $2$ and has good reduction; it exists because $j$ is integral
\cite[Chapter~VII, Proposition~5.5]{SilvermanAEC}. Let $y'^2+a_1'x'y'+a_3'y'=x'^3+a_2'x'^2+a_4'x'+a_6'$ be a model with
integral coefficients and unit discriminant $\Delta'$, with formal parameter $\mathsf t=-x'/y'$. As
$2$ splits in $K$, the reduction is ordinary (Deuring's criterion~\cite[Chapter~13, Theorem~12]{LangEF}), i.e.\ the reduced
curve has a point of order $2$ over the algebraic closure of the residue field; this is the
definition of ``ordinary'' used here, and we prove the direction needed. Enlarge $L$ so that it
contains $K$ and the fields of definition of the endomorphisms $[\alpha]$, $\alpha\in\auO_K$ (as
$\auO_K\Lambda=\Lambda$, these are the maps $z\mapsto\alpha z$ of $\C/\Lambda$). The smooth model
is the N\'eron model of its generic fibre~\cite[\S1.2, Definition~1, p.~12, and Proposition~8, p.~15]{BLR} (\cite[IV, Corollary~6.3, p.~329, and the N\'eron mapping property, p.~319]{SilvermanAT}), so each $[\alpha]$ extends to it and
reduces to an endomorphism $\widetilde{[\alpha]}$ of $\widetilde{\mathcal E}$ of the same degree
$\auNm\alpha$ (the curve and the $[\alpha]$ are defined over a number field and have good
reduction at the prime above $2$ in question, and reduction of homomorphisms then preserves
degrees~\cite[Chapter~II, Proposition~4.4]{SilvermanAT}). The invariant differential $\omega'=dx'/(2y'+a_1'x'+a_3')$ is a
nonzero multiple of $dz$, so $[\alpha]^*\omega'=\alpha\omega'$; it reduces to the invariant
differential $\tilde\omega'=(1+\dots)d\tilde{\mathsf t}$ of $\widetilde{\mathcal E}$, so
$\widetilde{[\alpha]}^*\tilde{\mathsf t}=\tilde\alpha\tilde{\mathsf t}+O(\tilde{\mathsf t}^2)$, with
$\tilde\alpha$ the residue of $\alpha$. Let $\mathfrak p=\{\alpha\in\auO_K:\tilde\alpha=0\}$, one of
the two primes above $2$, and $h$ the class number of $K$; then $\aubar{\mathfrak p}^h=(\beta)$ with
$\beta\notin\mathfrak p$ (else $\aubar{\mathfrak p}\subseteq\mathfrak p$), so $\tilde\beta\neq0$ and
$\widetilde{[\beta]}$ has ramification index $1$ at the origin $O$. For $R$ in its kernel,
$\widetilde{[\beta]}\circ\tau_R=\widetilde{[\beta]}$ with $\tau_R$ the translation, so the index is
$1$ at every point of the kernel; as the indices over $O$ add up to the degree~\cite[Chapter~II, Proposition~2.6(a)]{SilvermanAEC},
the kernel has $\auNm\beta=2^h\ge2$ points. They are killed by $[2^h]$, so some multiple of a nonzero
one has order $2$. In characteristic $2$ a point
of order $2$ satisfies $\tilde y=-\tilde y-\tilde a_1'\tilde x-\tilde a_3'$, i.e.\ $\tilde a_1'\tilde
x+\tilde a_3'=0$; if $\tilde a_1'=0$ then $\tilde b_2'=\tilde b_4'=0$, $\tilde b_6'=\tilde a_3'^2$,
$\tilde\Delta'=-27\tilde a_3'^4=\tilde a_3'^4\ne0$, and no point satisfies this. So $a_1'$ is a unit. A
point $T\ne O$ reduces to the origin iff $v(x'(T))<0$ (if $v(x')\ge0$ then $y'$, a root of a monic
polynomial over the integers, is integral; if $v(x')<0$ then $v(y')=\tfrac32v(x')$ and
$[x':y':1]=[x'/y':1:1/y']$ reduces to $[0:1:0]$). The $x'$-coordinates of the points of order $2$ are the
roots of $g(X)=4X^3+b_2'X^2+2b_4'X+b_6'$, with $v(4)=2$, $v(b_2')=v(a_1'^2+4a_2')=0$, $v(2b_4')\ge1$,
$v(b_6')\ge0$ (normalising $v(2)=1$). The Newton polygon of $g$ has the segment from $(2,0)$ to $(3,2)$ and
all other slopes $\le0$; so exactly one root has negative valuation, namely $-2$. Hence the kernel of reduction
contains exactly one point $T_c$ of order two, and $v(x'(T_c))=-2$. This kernel is stable under $\auO_K$, so
$\{O,T_c\}$ is the part of $\mathcal E[2]$ killed by one of the two primes above $2$; let $T_e$
generate the other part. The admissible point is $P_0=T_c+T_e$ (the point $z=0$ is the point $t$
of $\C/\Lambda$). So $P_0$ and $T_e$ do not reduce to the origin, their coordinates $x'$ are
integral, and $v(x'(T_c)-x'(P_0))=-2$. The numbers $x'(T)$, $T\in\mathcal E[2]\setminus O$, are the
roots of $4X^3+b_2'X^2+2b_4'X+b_6'$, and the square of the product of their differences is
$\Delta'/16$ (for a cubic $aX^3+bX^2+cX+d$ with roots $e_i$, $a^4\prod_{i<j}(e_i-e_j)^2=b^2c^2-4ac^3-4b^3d
-27a^2d^2+18abcd$; with $(a,b,c,d)=(4,b_2',2b_4',b_6')$ this is $4(b_2'^2b_4'^2-32b_4'^3-b_2'^3b_6'-108b_6'^2
+36b_2'b_4'b_6')$, and by the definition $\Delta'=-b_2'^2b_8'-8b_4'^3-27b_6'^2+9b_2'b_4'b_6'$ with
$4b_8'=b_2'b_6'-b_4'^2$ the bracket is $4\Delta'$, so the right side is $16\Delta'$; dividing by $a^4=256$ gives $\Delta'/16$; see also
\cite[Chapter~III, \S1]{SilvermanAEC}); hence $2(-2)+2(-2)+2v(x'(T_e)-x'(P_0))=-4$ and
$v(x'(T_e)-x'(P_0))=2$. The number
\[
\varrho=\frac{x'(T_c)-x'(P_0)}{x'(T_e)-x'(P_0)},\qquad v(\varrho)=-4,
\]
does not change under the changes of coordinates $x'=u^2x+r$ between Weierstrass models. In the
model $\mathcal E$ we have $x(P_0)=0$, and $x(T_c),x(T_e)$ are the roots of $X^2+\tfrac14a_1^2X+1$;
so $\varrho+2+\varrho^{-1}=(x(T_c)+x(T_e))^2/(x(T_c)x(T_e))=a_1^4/16$, and $v(a_1^4)=4+v(\varrho)=0$.
The rest follows from Lemma~\ref{pa:model} (the discriminant is $a_1^4-64$; at odd primes $P_0$,
of order two, does not reduce to the origin).
\end{proof}

\begin{corollary}[Proposition~\ref{pk:curve}(C1), (C2), (C4) at every prime]\label{pf:CS}
With normalised lattices in the sense above, {\rm(C1)} holds, {\rm(C2)} holds at every prime
$\ell\ne d$ including $\ell=2$, and the number $\mathsf u$ of {\rm(C4)} is a unit at every prime,
including $2$.
\end{corollary}
\begin{proof}
(C1): the point $t^\sigma$ constructed there is admissible, by Lemma~\ref{pf:transport}(a) below
(the bijection $S_\sigma$ is $\auO_K$-linear and maps $t$ to $t^\sigma$; part (a) of that lemma is proved from the construction of the conjugate curve alone and does not use (C1), which enters the lemma only in part (b)). The proof of (C2) uses
only that $a_1$ is an algebraic integer and that torsion points of order prime to $\ell$ have
integral coordinates; it is valid for $\ell=2$. In the proof of (C4), with the lattices
$(\Lambda_i,t_i,a_{1,i})$ normalised in the present sense (this fixes $\mu_i^{12}$), the only
properties of an odd prime that are used are: $\Delta(\Lambda_i)=a_{1,i}^4-64$ is a unit, and the
curves have good reduction, so that homomorphisms extend to the smooth models and multiply the
invariant differentials $dx/(2y+a_1x)$ by integers. By Lemma~\ref{pf:two} both hold at $2$.
\end{proof}

\emph{The character of the vacuum.} Recall (\S\ref{ps:ss-weil}) that the character $\epsilon_L$ is
\begin{equation}\label{pf:eq-eps}
\epsilon_L(l)=\psi_f\bigl(-\tfrac12\,\omega\,\jmath\,\langle l,l\rangle\bigr)\qquad(l\in L),
\end{equation}
that is, the element ``$t$ with $t+\aubar t=1$'' is $\omega$, and there is a factor $-\frac12$. We check the cocycle relation: $\epsilon_L(l)\epsilon_L(l')
\epsilon_L(l+l')^{-1}=\psi(\frac12\omega\jmath(a+\aubar a))$ with $a=\langle l',l\rangle$, and
$\psi(\frac12\omega\jmath\aubar a)=\psi(-\frac12\aubar\omega\jmath a)$ because $\psi\circ c=\psi$;
so the quotient is $\psi(\frac12\jmath a(\omega-\aubar\omega))=\beta(l,l')\psi(\jmath a(\omega-1))=
\beta(l,l')$, as $\psi(\jmath\langle l',l\rangle\widehat\auO_E)=1$ by self-duality. Another
choice of the element of trace one replaces $\psvac_L$ by its translate by a point of
$\frac12L$; this changes nothing at odd primes, but at $\ell=2$ the statement (K1) would then be
false (the theta function of such a translate contains factors $\theta_s$ with $s$ of order two,
which are not $\partial$-integral at $2$).

\begin{lemma}\label{pf:eps}
Let $W$ be a hermitian space and $L\subset W(\auA_f)$ a self-dual lattice. Then $\epsilon_L$,
defined by \eqref{pf:eq-eps}, has values $\pm1$. If $W=E^{\mathsf C}$ is standard and
$L=\alpha\widehat\auO_E^{\,\mathsf C}$, then $\epsilon_L(\lambda)=(-1)^{\pfT_{E^+/\Q}\langle\lambda,
\lambda\rangle}$ for $\lambda\in\Lambda_L=W(E)\cap L=\mathfrak a_0^{\mathsf C}$.
\end{lemma}
\begin{proof}
Let $\nu=\langle l,l\rangle\in\auA_{E^+,f}$. Since $\omega\jmath=(\jmath-1)/2$ and
$\pfT_{E/E^+}(\jmath\nu)=0$, $\epsilon_L(l)=\psi_f(\nu/4)=\psi_{E^+,f}(\nu/2)$, where
$\psi_{E^+}=\psi_\Q\circ\pfT_{E^+/\Q}$. Since $[x,y]=\beta(x,y)\beta(y,x)^{-1}=\psi(\jmath\langle y,x\rangle)$ and $L$ is a
self-dual $\widehat\auO_E$-module, $\psi_f(\jmath\nu\widehat\auO_E)=1$, that is, $\jmath\nu\in
\mathfrak d_E^{-1}\widehat\auO_E$; as $\jmath^{-1}=-\sqrt{-d}$ is integral, $\nu\in\mathfrak
d_E^{-1}\widehat\auO_E\cap\auA_{E^+,f}=\mathfrak d_{E^+}^{-1}\widehat\auO_{E^+}$, because $E/E^+$ is
unramified at the finite places ($\mathfrak d_E=\mathfrak d_{E^+}\auO_E$).
Hence $\psi_{E^+,f}(\nu)=1$ and $\epsilon_L(l)^2=1$. In the standard case, for
$\lambda\in\Lambda_L$ the number $n=\pfT_{E^+/\Q}\langle\lambda,\lambda\rangle=\pfT_{E/K}
\langle\lambda,\lambda\rangle$ lies in $\pfT_{E/K}(\mathfrak a_0\aubar{\mathfrak a_0})\cap\R=
\pfT_{E/K}(\mathfrak d_{E/K}^{-1})\cap\R\subseteq\Z$, and $\psi_{\Q,f}(n/2)=e^{2\boldsymbol\pi in/2}$.
\end{proof}

\subsection{Values at torsion points and the integrality principle}\label{pf:ss-tors}

Let $(\Lambda,t,a_1)$ be a normalised lattice, $\mathcal E$ its curve. We write $Q(z)$ for the
point of $\mathcal E$ with parameter $z$: $Q(z)=(x(z),y(z))$ for $z\notin t+\Lambda$ and $Q(z)=O$
for $z\in t+\Lambda$; thus $z\mapsto Q(z)-P_0$ is the group isomorphism $\C/\Lambda\to\mathcal
E(\C)$ of the Weierstrass uniformisation, in the coordinates $X=x+\wp(t)$, $2Y=2y+a_1x$. Torsion
points have algebraic coordinates; $E_1$ denotes the group of the torsion points $v\in\C/\Lambda$
whose point $Q(v)-P_0$ reduces to the origin at $\iota$ (the kernel of reduction). It is stable
under $\auO_K$ and has no element $\ne0$ of order prime to $\ell$~\cite[Chapter~VII, Proposition~3.1(a)]{SilvermanAEC}; by
Lemma~\ref{pf:two}, $t\notin E_1$. For $u\in\C$ put
\[
\pfbr u=\theta_u(0)=e^{-H(u,u)/2}\theta(u);\qquad\text{then}\qquad e^{-H(w,w)/2}\theta_u(w)=
e^{\boldsymbol\pi iE(w,u)}\,\pfbr{w+u}\quad(u,w\in\C).
\]
For $m\ge0$ let $\mathcal S_m=\mathcal S_m(\Lambda)$ be the space of entire functions $G$ with
$G(z+\lambda)=\chi(\lambda)^me^{mH(\lambda,z+\lambda/2)}G(z)$ for $\lambda\in\Lambda$. If
$G\in\mathcal S_m$, then $G/\theta^m$ is $\Lambda$-periodic with poles only on $t+\Lambda$, of
order at most $m$; so
\begin{equation}\label{pf:eq-sections}
\mathcal S_m=\theta^m\cdot\Bigl\{P(x,y)=\sum_{2a+3e\le m,\ e\le1}c_{a,e}\,x^ay^e\Bigr\},
\end{equation}
the monomials $x^ay^e$, $e\le1$, being a basis of the functions on $\mathcal E$ regular outside
$O$, with a pole of order $2a+3e$ at $O$. In particular $\mathcal S_0=\C$.

\begin{proposition}[Values at torsion points]\label{pf:values}
Let $u\in\Q\Lambda$. If $u\in t+\Lambda$ then $\pfbr u=0$. Otherwise $\pfbr u$ is an algebraic
number and, at every prime $\ell$ and every $\iota$,
\[
|\pfbr u|=1\ \text{ if }u-t\notin E_1,\qquad|\pfbr u|=|x(u)|^{-1/2}<1\ \text{ if }u-t\in E_1 .
\]
In particular $\theta_u(0)$ is an algebraic integer for every torsion point $u$.
\end{proposition}
\begin{proof}
Let $u\notin t+\Lambda$ have order $\mathsf N$ modulo $\Lambda$ and put $n=1+2\mathsf N\ell$. By
Lemma~\ref{pf:chi}(d) the function $\theta(nz)$ lies in $\mathcal S_{n^2}$, and as in
Lemma~\ref{pa:units}(a)
\[
\theta(nz)=\theta(z)^{n^2}R_n(x(z)),\qquad R_n(X)=\prod_s\bigl(1-X/x_s\bigr),
\]
the product over the $(n^2-1)/2$ pairs $\pm s$ of non-zero points of $n^{-1}\Lambda/\Lambda$, with
$x_s=x(s-t)$ the $x$-coordinate of the point $Q(s-t)$; the point $Q(s-t)-P_0$ has parameter $s-t$,
and $Q(s-t)$ reduces to the origin $O=Q(t)$ only if $s=(s-t)-t$ lies in $E_1$, which is impossible since
$s\ne0$ has order dividing $n$, prime to $\ell$. So these points do not reduce to
the origin, so the $x_s$ are integral; and $|x_s|=1$: if $|x_s|<1$ then
$|y(y+a_1x)|=|x^3+x|<1$ at this point, so it reduces to $\widetilde P_0$, that is, $s-t\in E_1$; then $n(s-t)\equiv t$ ($n$ is odd)
would lie in $E_1$, contrary to Lemma~\ref{pf:two}. Since $H(nz,nz)=n^2H(z,z)$, we get
$\pfbr{nz}=\pfbr z^{n^2}R_n(x(z))$ for all $z\notin t+\Lambda$. Take $z=u$. As $nu=u+\lambda$ with
$\lambda=2\ell\mathsf Nu\in2\Lambda$, Lemma~\ref{pf:chi}(c) gives $\pfbr{nu}=\zeta\pfbr u$ with
$\zeta=e^{\boldsymbol\pi iE(\lambda,u)}$ a root of unity, and $\pfbr u\ne0$; so
\[
\pfbr u^{\,n^2-1}=\zeta\cdot\prod_s\frac{x_s}{x_s-x(u)} .
\]
This shows that $\pfbr u$ is algebraic. The point $Q(u)-P_0$ has parameter $u$, so $Q(u)$ reduces to the origin
if and only if $u+t\equiv u-t$ lies in $E_1$, and $x(u)$ is integral if and only if $u-t\notin E_1$. If $u-t\in E_1$, then $|x(u)|>1$, every
factor has absolute value $|x(u)|^{-1}$, and $|\pfbr u|^{n^2-1}=|x(u)|^{-(n^2-1)/2}$. If
$u-t\notin E_1$, every factor has absolute value $\ge1$, with equality unless $x_s\equiv x(u)$
modulo $\pam$, that is, unless $Q(u)$ and $\pm Q(s-t)$ have the same reduction, that is,
$u+t\mp s\in E_1$. Then $(n-1)(u+t\mp s)\equiv\pm s$ (as $(n-1)u,(n-1)t\in\Lambda$, $ns\in\Lambda$)
would be a non-zero point of
$E_1$ of order prime to $\ell$, which is impossible. So $|\pfbr u|=1$.
\end{proof}

\begin{proposition}[Integrality principle]\label{pf:pull}
Let $G\in\mathcal S_m$ and suppose that $e^{-mH(z,z)/2}G(z)$ is integral for all torsion points $z$
of odd order prime to $\ell$. Then $G=\theta^m\cdot\Gamma(y)$ with $\Gamma\in\pao[[Y]]$, as power
series in $z$. In particular, if $\delta\in\auO_K$, $\delta\ne0$, and $\delta u\in\Lambda$, then
\[
\theta_u(\aubar\delta z)=\theta(z)^{\auNm\delta}\,\Gamma_{u,\delta}(y(z)),\qquad\Gamma_{u,\delta}\in
\pao[[Y]],
\]
at every prime $\ell$, and for $\ell\ne d$ the series $\theta_u(\aubar\delta z)$ is
$\partial$-integral. If moreover $\iota(\aubar\delta)$ is a unit and $\ell\ne d$, then $\theta_u$
itself is $\partial$-integral.
\end{proposition}
\begin{proof}
Write $G=\theta^mP(x,y)$ as in \eqref{pf:eq-sections}. Suppose that $P$ is not integral, let
$c_0$ be a coefficient of $P$ of largest absolute value, $|c_0|>1$, and $P_1=P/c_0$. The reduction
$\widetilde P_1$ of $P_1$ modulo $\pam$ is a non-zero function on the reduced curve (the monomials
are linearly independent there), regular outside $\widetilde O$ with a pole of order $\le m$, so
it has at most $m$ zeros. Let $z$ be a torsion point of odd order $n'$
prime to $\ell$. Then $z\notin E_1+t$, since otherwise $n'(z-t)\equiv t$ would lie in $E_1$; so
$Q(z)$ has integral coordinates, and by Proposition~\ref{pf:values} $\pfbr z$ is a unit. Hence
$P(Q(z))=G(z)/\theta(z)^m=e^{-mH(z,z)/2}G(z)\cdot\pfbr z^{-m}$ is integral, and $P_1(Q(z))=P(Q(z))/c_0\in\pam$.
Thus $\widetilde P_1$ vanishes at the reductions of all these points; they are pairwise distinct
(differences of such $z$ have odd order prime to $\ell$) and infinite in number, a contradiction.
So $P$ has coefficients in $\pao$, and $P(x,y)\in\pao[[y]]$ because $x\in y^2\pao[[y]]$
(\S\ref{pa:ss-curve}; this uses only that $a_1$ is integral).

For $G(z)=\theta_u(\aubar\delta z)$ and $m=\auNm\delta$: $G\in\mathcal S_m$ by
Lemma~\ref{pf:chi}(d), and $e^{-mH(z,z)/2}G(z)=e^{-H(w,w)/2}\theta_u(w)$ with $w=\aubar\delta z$ is
a root of unity times $\pfbr{w+u}$, which is integral by Proposition~\ref{pf:values}. The series
$\theta$ and $y$ are $\partial$-integral for $\ell\ne d$ (Corollary~\ref{pf:CS}), and
$\partial$-integral series are stable under products, under substitution into series of
$\pao[[Y]]$, and under $z\mapsto cz$ with $c\in\pao$; for the last assertion take
$c=\iota(\aubar\delta)^{-1}$.
\end{proof}

\begin{remark}[The case excluded in \S\ref{pk:ss-N}]\label{pf:gen}
Proposition~\ref{pf:pull} has no exception: for $u=t+e$ with $0\ne e\in E_1$ of order $\ell$, and
$\delta$ with $\delta e=0$, $\iota(\aubar\delta)\in\pao^\times$, the series $\theta_u$ \emph{is}
$\partial$-integral. Note that $|\theta_u(0)|<1$ (Proposition~\ref{pf:values}) and that $\theta_u'(0)/\theta_u(0)$
is not integral, but this does not prevent integrality of the derivatives: $\theta_u=\theta\cdot
\Gamma(y(\cdot/\aubar\delta))$ with $\Gamma\in\pao[[Y]]$ a series that has a zero in the open unit
disc (at the parameter of the point $-e$), like the series $Y-c$ with $0<|c|<1$. This is why the condition
(Gen) of Proposition~\ref{pk:K1N} is not needed, and Lemmas~\ref{pk:E1}, \ref{pk:formal} are not used. At $\ell=N$ the excluded case is the typical one: the kernel of reduction contains
$\mathfrak n^{-1}\Lambda/\Lambda$ for the prime $\mathfrak n$ induced by $\iota$, which is where the
translation points have their denominators.
\end{remark}

\subsection{Conjugation: proof of Proposition~\ref{pk:curve}{\rm(C3)(ii)}}\label{pf:ss-recip}

Let $(\Lambda,t,a_1)$ be a normalised lattice, $\sigma\in\auAut(\C/K)$, and
$(\Lambda^\sigma,t^\sigma,\sigma a_1)$ the lattice of (C1), with functions $x^\sigma,y^\sigma,
\theta^\sigma$, curve $\mathcal E^\sigma$, and $Q^\sigma(z)$. For $G=\theta^mP(x,y)\in\mathcal
S_m(\Lambda)$ put $G^\sigma=(\theta^\sigma)^mP^\sigma(x^\sigma,y^\sigma)\in\mathcal S_m
(\Lambda^\sigma)$, where $P^\sigma$ is $P$ with $\sigma$ applied to the coefficients.

\begin{lemma}[Transport]\label{pf:transport}
\begin{enumerate}
\item[(a)] There is a unique bijection $S_\sigma\colon\C/\Lambda\to\C/\Lambda^\sigma$ with
$Q^\sigma(S_\sigma z)=\sigma(Q(z))$. It is an isomorphism of $\auO_K$-modules, and
$S_\sigma(t)=t^\sigma$.
\item[(b)] $G\mapsto G^\sigma$ is $\sigma$-semilinear, multiplicative, and the Taylor series of
$G^\sigma$ at $0$ is the image under $\sigma$ of that of $G$. If $G\in\mathcal S_m(\Lambda)$ and
$G'$ is an entire function whose Taylor series is the image under $\sigma$ of that of $G$, then
$G'=G^\sigma$.
\item[(c)] For $z\in\C$ and $z'\in\C$ with $z'\equiv S_\sigma(z)$ there is $\mu=\mu(z,z')\in
\C^\times$ such that $G^\sigma(z')=\mu^m\,\sigma(G(z))$ for all $m\ge0$ and all $G\in\mathcal S_m
(\Lambda)$.
\end{enumerate}
\end{lemma}
\begin{proof}
(a) The map $\widetilde S=(Q^\sigma-P_0^\sigma)^{-1}\circ\sigma\circ(Q-P_0)$ is a group isomorphism
$\C/\Lambda\to\C/\Lambda^\sigma$, since $\sigma\colon\mathcal E(\C)\to\mathcal E^\sigma(\C)$ is one.
As $\sigma(P_0)=P_0^\sigma$, $Q^\sigma(\widetilde Sz)=\sigma(Q(z)-P_0)+P_0^\sigma=\sigma Q(z)$; so
$S_\sigma=\widetilde S$, and $S_\sigma(t)=t^\sigma$ because $Q(t)=O$. For $a\in\auO_K$ let $[a]$ be
the endomorphism of $\mathcal E$ induced by $z\mapsto az$; it is characterised by $[a]^*\varpi=
a\varpi$ for the invariant differential $\varpi=dx/(2y+a_1x)=dz$. The conjugate endomorphism
$\sigma[a]$ of $\mathcal E^\sigma$ satisfies $(\sigma[a])^*\varpi^\sigma=\sigma(a)\varpi^\sigma=
a\varpi^\sigma$, so it is $[a]$ on $\mathcal E^\sigma$, and $S_\sigma(az)=aS_\sigma(z)$.
(b) By (C1), in the form of Corollary~\ref{pf:CS} (whose proof of it uses only part (a) above), the series of $x^\sigma,y^\sigma,\theta^\sigma$ are the conjugates of those of
$x,y,\theta$. The last assertion holds because an entire function is determined by its Taylor
series. (c) If $z\notin t+\Lambda$, then $P^\sigma(Q^\sigma(z'))=\sigma(P(Q(z)))$ and
$\mu=\theta^\sigma(z')/\sigma(\theta(z))$. If $z\in t+\Lambda$, let $\mathsf t=x/y$, a function
with a simple zero at $O$; for $2a+3e=m$ the function $x^ay^e\mathsf t^m=(x^3/y^2)^{(m-e)/2}$
takes the value $1$ at $O$, so $G(z)=c\,\lambda_z^m$, where $c$ is the coefficient in $P$ of the
monomial with $2a+3e=m$ (or $0$) and $\lambda_z=\lim_{w\to z}\theta(w)/\mathsf t(w)\ne0$; likewise
$G^\sigma(z')=\sigma(c)\lambda_{z'}'^{\,m}$, and $\mu=\lambda'_{z'}/\sigma(\lambda_z)$.
\end{proof}

\begin{theorem}[Proposition~\ref{pk:curve}{\rm(C3)(ii)}]\label{pf:recip}
Let $\Lambda=\Omega_0\auO_K$ be normalised and $\sigma\in\auAut(\C/K)$. There are $\Omega_\sigma
\in\C^\times$, an ideal $\mathfrak c\ne0$ of $\auO_K$ and a finite idele $\beta_\sigma$ of $K$ such
that, with $\mathfrak a=\mathfrak c\aubar{\mathfrak c}^{-1}$ and $\widehat{\mathfrak a}=\mathfrak
a\widehat\auO_K$:
\begin{enumerate}
\item[(a)] $\Lambda^\sigma=\Omega_\sigma\mathfrak a$, $\beta_\sigma\widehat\auO_K=\widehat{
\mathfrak a}$ (so $\beta_\sigma\aubar{\beta_\sigma}$ is a unit), and $S_\sigma(\mathsf x\Omega_0)
\equiv\mathsf x'\Omega_\sigma$ modulo $\Lambda^\sigma$ whenever $\mathsf x,\mathsf x'\in K$ and
$\mathsf x'-\beta_\sigma\mathsf x\in\widehat{\mathfrak a}$;
\item[(b)] if $\mathsf x\in K$, $n\ge1$, $n\mathsf x\in\auO_K$, and $\mathsf x'\in K$ satisfies
$\mathsf x'-\beta_\sigma\mathsf x\in2n\widehat{\mathfrak a}$, then $\sigma$ maps the Taylor series
at $0$ of $\theta_{\mathsf x\Omega_0}$ to that of $\theta^\sigma_{\mathsf x'\Omega_\sigma}$.
\end{enumerate}
\end{theorem}
\begin{proof}
(a) $\Lambda^\sigma$ is stable under $\auO_K$, so $\Lambda^\sigma=\Omega\mathfrak b$ for a
fractional ideal $\mathfrak b$. The class number of $K$ is odd. (Proof: only the prime $d$ ramifies in $K$, and
$(\sqrt{-d})$ is the prime above it. Let a class $c$ satisfy $c^2=1$ and let $\mathfrak a\in c$; since
$\mathfrak a\aubar{\mathfrak a}=(\auNm\mathfrak a)$, $\aubar{\mathfrak a}\in c^{-1}=c$, so $\aubar{\mathfrak
a}=\gamma\mathfrak a$ with $\gamma\in K^\times$, and $\gamma\aubar\gamma$ is a positive unit, i.e.\ $1$.
Then $\gamma=\beta/\aubar\beta$ with $\beta=1+\gamma$ if $\gamma\ne-1$ and $\beta=\sqrt{-d}$ if
$\gamma=-1$, and $\mathfrak b=\beta\mathfrak a\in c$ satisfies $\aubar{\mathfrak b}=\aubar\beta\gamma
\mathfrak a=\mathfrak b$. In the prime factorisation of a conjugation-stable fractional ideal, split
primes $\mathfrak p\ne\aubar{\mathfrak p}$ occur with equal exponents and their product is $(p)$, inert
primes are $(p)$, and the ramified prime is $(\sqrt{-d})$; so $\mathfrak b$ is principal and $c=1$. A
finite abelian group without elements of order $2$ has odd order. See also \cite{CoxPrimes}.) So $\mathfrak b=\eta\mathfrak c^2$ with $\eta\in K^\times$, and
$\mathfrak c^2=\auNm(\mathfrak c)\cdot\mathfrak c\aubar{\mathfrak c}^{-1}$. By
Lemma~\ref{pf:transport}(a), $S_\sigma$ induces an isomorphism of $\auO_K$-modules
$K/\auO_K\to K/\mathfrak a$ on the torsion points. Now
$\auHom_{\auO_K}(K/\auO_K,K/\mathfrak a)=\varprojlim_n\auHom(n^{-1}\auO_K/\auO_K,n^{-1}\mathfrak
a/\mathfrak a)=\varprojlim_n\mathfrak a/n\mathfrak a=\widehat{\mathfrak a}$, acting by
multiplication on $K/\auO_K=\auA_{K,f}/\widehat\auO_K$ with values in $\auA_{K,f}/\widehat{
\mathfrak a}=K/\mathfrak a$; the element $\beta_\sigma\in\widehat{\mathfrak a}$ of an isomorphism
generates $\widehat{\mathfrak a}$.

(b) Put $w=\Omega_0/(2n^2)$, choose $\mathsf b\in K$ with $\mathsf b-\beta_\sigma\in2n^2\widehat{
\mathfrak a}$, and put $w'=\mathsf b\Omega_\sigma/(2n^2)$, $u=\mathsf x\Omega_0=2n^2\mathsf xw$,
$u'=2n^2\mathsf xw'=\mathsf x\mathsf b\Omega_\sigma$. By (a), $w'\equiv S_\sigma(w)$, hence
$aw'\equiv S_\sigma(aw)$ for $a\in\auO_K$ and $u'\equiv S_\sigma(u)$. Also $\mathsf b\in\mathfrak
a$, so $nu'\in\Lambda^\sigma$; and $\mathsf x'-\mathsf x\mathsf b\in K\cap2n\widehat{\mathfrak a}=
2n\mathfrak a$, so $\theta^\sigma_{\mathsf x'\Omega_\sigma}=\theta^\sigma_{u'}$ by
Lemma~\ref{pf:chi}(c). Let
\[
g(z)=\theta_u(nz)\in\mathcal S_{n^2}(\Lambda),\qquad g'(z)=\theta^\sigma_{u'}(nz)\in\mathcal S_{n^2}
(\Lambda^\sigma)
\]
(Lemma~\ref{pf:chi}(d) with $\delta=n$). We show $g^\sigma=g'$; by Lemma~\ref{pf:transport}(b)
this is the assertion, since $\sigma(n)=n$.

\emph{$g^\sigma=c\,g'$ with $c\in\C$.} The zeros of $g'$ are simple and are the $z'$ with
$nz'+u'-t^\sigma\in\Lambda^\sigma$. Let $z'$ be one of them and $z\in\C$ with $S_\sigma(z)\equiv
z'$. Then $S_\sigma(nz+u-t)\equiv nz'+u'-t^\sigma\equiv0$, so $nz+u\in t+\Lambda$, $g(z)=0$, and
$g^\sigma(z')=0$ by Lemma~\ref{pf:transport}(c). Hence $g^\sigma/g'$ is entire and lies in
$\mathcal S_0(\Lambda^\sigma)=\C$.

\emph{$c=1$.} Put $p=u/2=n^2\mathsf xw$, $p'=u'/2$, $z_0=-n\mathsf xw$, $z_0'=-n\mathsf xw'$; then
$z_0'\equiv S_\sigma(z_0)$, $nz_0=-p$, $nz_0'=-p'$. Let $\mu=\mu(z_0,z_0')$. From the definition,
$\theta_u(-p)=e^{H(u,p)-H(u,u)/2}\theta(p)=\theta(p)$; so $g(z_0)=\theta(p)$ and $g'(z_0')=
\theta^\sigma(p')$. Let $h(z)=\theta(nz)\in\mathcal S_{n^2}(\Lambda)$; by (C1) and
Lemma~\ref{pf:transport}(b), $h^\sigma(z)=\theta^\sigma(nz)$, and $\theta$ is even, so
Lemma~\ref{pf:transport}(c) gives
\[
\theta^\sigma(p')=h^\sigma(z_0')=\mu^{n^2}\sigma(\theta(p)),\qquad g^\sigma(z_0')=\mu^{n^2}\sigma
(g(z_0))=\mu^{n^2}\sigma(\theta(p))=\theta^\sigma(p')=g'(z_0').
\]
If $\theta(p)\ne0$ this gives $c=1$. If $\theta(p)=0$, then $\theta^\sigma(p')=0$, so
$u=2p\in2t+2\Lambda$ and $u'\in2t^\sigma+2\Lambda^\sigma$; by Lemma~\ref{pf:chi}(a),(c),
$\theta_u=\chi(u)\theta=-\theta$ and $\theta^\sigma_{u'}=-\theta^\sigma$, so $g=-h$,
$g'=-h^\sigma=g^\sigma$.
\end{proof}

\begin{corollary}\label{pf:conj}
Lemma~\ref{ps:conj} and {\cite[Proposition~14.23]{P46}} hold.
\end{corollary}
\begin{proof}
Proposition~\ref{pk:conj} deduces Lemma~\ref{ps:conj} from (C3)(ii), which is
Theorem~\ref{pf:recip} (``close enough'' meaning the congruence of (b)); in its proof $\mathsf
B=\beta_\sigma$ is an idele with $\mathsf B\aubar{\mathsf B}$ a unit by (a), so that $\mathsf BL$
is self-dual, as is used there. {\cite[Proposition~14.23]{P46}} was proved in \S\ref{ps:ss-kernel} from
Lemma~\ref{ps:conj}.
\end{proof}

\subsection{Tight frames, the scale, and standard data at every prime}\label{pf:ss-frames}

Let $h(\mathsf x,\mathsf y)=\pfT_{E/K}(\aubar{\mathsf x}\mathsf y)$ on $E$; the lattice
$\mathfrak a_0$ is unimodular for $h$ over $\auO_K$, because $\mathfrak a_0\aubar{\mathfrak a_0}=
\mathfrak d_{E/K}^{-1}$. It is stable under $\operatorname{Gal}(E/K)$ and under $\auO_E$. A
\emph{tight frame of length $\mathsf b\ge1$} is a finite family $(v_i)_{i\in\mathsf I}$ in
$\mathfrak a_0$ with
\[
\sum_{i\in\mathsf I}\tau(v_i)\,\aubar{\tau'(v_i)}=\mathsf b\,\delta_{\tau\tau'}\qquad(\tau,\tau'\in
\auPl);
\]
then the vectors $u_i=v_i/\mathsf b$ with multiplicities $\mathsf b$, placed in each coordinate
of $E^{\mathsf C}$, form a resolution of the identity \eqref{pa:eq-res}. The family
$(\zeta_{q_1}^k)_{0\le k<q_1}$ is a tight frame of length $q_1$; it is the one of
\eqref{pk:eq-product}.

\begin{lemma}[A frame of length prime to $q_1$]\label{pf:frame}
There is a tight frame of a length $\mathsf h$ prime to $q_1$.
\end{lemma}
\begin{proof}
Let $\zeta=\zeta_{q_1}$, $P$ its minimal polynomial over $K$ (the cyclotomic polynomial, of degree
$f$), and $P(X)/(X-\zeta)=\sum_{i<f}b_iX^i$ with $b_i\in\auO_E$. By Euler's
formula~\cite[III, \S6, Lemma~2, p.~56]{SerreLF}, $\pfT_{E/K}(\zeta^k\,b_i/P'(\zeta))=\delta_{ik}$ for $0\le i,k<f$
(compare the coefficients of $X^i$ in the Lagrange interpolation identity
$\sum_{\sigma}\sigma(\zeta)^kP(X)/\bigl((X-\sigma\zeta)P'(\sigma\zeta)\bigr)=X^k$, $\sigma$ running over the
$K$-embeddings of $E$, which holds since both sides have degree $<f$ and agree at the $f$ roots of $P$),
and $P'(\zeta)$ generates $\mathfrak d_{E/K}$, which divides $q_1$. Choose $J\ge1$ with
$q_1\mid(\omega\aubar\omega)^J$ and put $e_i=\aubar\omega^{\,J}b_i/P'(\zeta)$. Then
$e_i\in\mathfrak a_0$: at the places dividing $\omega$, where $\alpha=1/P'(\zeta)$, the element
$P'(\zeta)e_i$ is integral; at the other places dividing $q_1$, $\omega$ is a unit
($\omega+\aubar\omega=1$) and $\aubar\omega^{\,J}/q_1=(\omega\aubar\omega)^J/(q_1\omega^J)$ is
integral; elsewhere $P'(\zeta)$ is a unit. Transposing the matrix identity that expresses the
duality of the bases $(\zeta^k)$ and $(b_i/P'(\zeta))$, and using that the $\tau\in\auPl$ fix $K$
and commute with $c$,
\begin{equation}\label{pf:eq-cross}
\sum_{i<f}\tau(e_i)\,\aubar{\tau'(\zeta^{-i})}=\sum_{i<f}\tau(e_i)\tau'(\zeta^{i})=\aubar\omega^{\,J}
\delta_{\tau\tau'} .
\end{equation}
For $\mathsf w\in\mathfrak a_0$ the family $(\zeta^k\rho(\mathsf w))_{k<q_1,\ \rho\in\operatorname{
Gal}(E/K)}$ has the matrix $\sum_{k,\rho}\tau(\zeta)^k\tau'(\zeta)^{-k}\tau\rho(\mathsf w)\aubar{
\tau'\rho(\mathsf w)}=q_1\nu(\mathsf w)\delta_{\tau\tau'}$, where $\nu(\mathsf w)=h(\mathsf w,
\mathsf w)\in\auO_K\cap\R_{\ge0}=\Z_{\ge0}$; so it is a tight frame of length $q_1\nu(\mathsf w)$
if $\mathsf w\ne0$. Consider the union of the frame $(\zeta^{-k})_{k<q_1}$ and of these frames for
$\mathsf w=e_0,\dots,e_{f-1}$; in it, for each $i<f$, replace the two members $\zeta^{-i}$ and
$e_i$ $(k=0$, $\rho=1)$ by the single member $\zeta^{-i}+e_i$. By \eqref{pf:eq-cross} and its
conjugate transpose the matrix of the new family is
\[
\Bigl(q_1\bigl(1+\sum_{i<f}\nu(e_i)\bigr)+\omega^J+\aubar\omega^{\,J}\Bigr)\delta_{\tau\tau'}=:
\mathsf h\,\delta_{\tau\tau'} .
\]
The diagonal entries are sums of squares of absolute values, so $\mathsf h$ is a positive integer.
For a prime $p\mid q_1$, $p$ divides $\omega\aubar\omega$; the integers $T_J=\omega^J+\aubar
\omega^{\,J}$ satisfy $T_J=T_{J-1}-\omega\aubar\omega T_{J-2}\equiv T_{J-1}$ and $T_1=1$, so
$\mathsf h\equiv1\pmod p$.
\end{proof}

\begin{lemma}[The scale]\label{pf:scale}
For every prime $\ell$, every $\iota$, every $\tau\in\auPl$ and every $\mathsf v\in\mathfrak a_0$,
the number $\varpi_{\mathrm{vac}}\cdot\aubar{\tau(\mathsf v)}$ is integral. More precisely
$|\varpi_{\mathrm{vac}}|=|\alpha_{w_\tau}|_{w_\tau}^{-1}$, where $w_\tau$ is the place of $E$
induced by $\iota\circ\tau\circ c$ and $|\cdot|_{w_\tau}=|\iota\tau c(\cdot)|$.
\end{lemma}
\begin{proof}
Let $\chi'=\upsilon|\cdot|^{-1/2}$, so that $\varpi_{\mathrm{vac}}^{-f}=\chi'(\alpha)$. It is a
Hecke character of $E$, trivial on $\widehat\auO_E^\times$ ($\upsilon$ is unramified at every
finite place), with $\chi'_\infty(\mathsf a)=\prod_\tau\tau(\mathsf a)/|\tau(\mathsf a)|^2=\prod_\tau
\aubar{\tau(\mathsf a)}^{-1}$. Let $k\ge1$ and $\mathsf a\in E^\times$ with $\alpha^k\in\mathsf a
\widehat\auO_E^\times$. Then $\chi'(\alpha)^k=\chi'_f(\mathsf a)=\chi'_\infty(\mathsf a)^{-1}=
\prod_\tau\aubar{\tau(\mathsf a)}$, and $|\aubar{\tau(\mathsf a)}|=|\mathsf a|_{w_\tau}=|\alpha|_{
w_\tau}^k$. The places $w_\tau$, $\tau\in\auPl$, are the places of $E$ above one place of $K$, they
are permuted transitively by $\operatorname{Gal}(E/K)$, and $\alpha$ is $1$ or a generator of the
Galois-stable ideal $\mathfrak d_{E/K}^{-1}$ at all of them; so $|\alpha|_{w_\tau}$ does not depend
on $\tau$, and $|\varpi_{\mathrm{vac}}|^{-fk}=|\alpha|_{w_\tau}^{fk}$. Finally $|\aubar{\tau(\mathsf
v)}|=|\mathsf v|_{w_\tau}\le|\alpha|_{w_\tau}$ for $\mathsf v\in\mathfrak a_0$.
\end{proof}

Let $W=E^{\mathsf C}$, $L$, $\Lambda_L$ be as in \S\ref{pk:ss-conv}, $\Lambda=\Omega_0\auO_K$
normalised, and let $(v_i)_{i\in\mathsf I}$ be a tight frame of length $\mathsf b$,
$u_i=v_i/\mathsf b$. Put
\[
\Pi^v(z)=\prod_{\mathsf c\in\mathsf C}\prod_{i\in\mathsf I}\theta\bigl(a_{i,\mathsf c}(z)\bigr)^{
\mathsf b},\qquad a_{i,\mathsf c}(z)=\tfrac{\Omega_0}{c_\psi}\,u_i^\dagger(z_{\mathsf c}) ,
\]
so that $a_{i,\mathsf c}(\gamma^\natural)=\Omega_0\pfT_{E/K}(\aubar{u_i}\gamma_{\mathsf c})$ for
$\gamma\in W(E)$, and, by \eqref{pa:eq-res}, as in \S\ref{pk:ss-conj},
\[
\langle z,x\rangle=\sum_{i,\mathsf c}\mathsf b\,H\bigl(a_{i,\mathsf c}(x),a_{i,\mathsf c}(z)\bigr),
\qquad|x|^2=\sum_{i,\mathsf c}\mathsf b\,H\bigl(a_{i,\mathsf c}(x),a_{i,\mathsf c}(x)\bigr).
\]

\begin{lemma}[The product and the vacuum have the same character]\label{pf:semi}
$\Pi^v=\Theta_{F^v}$ for a unique $F^v\in\psHS_f(W)$; $\theta(1\otimes F^v)=1$;
$\psrho_f(l)F^v=\epsilon_L(l)F^v$ for all $l\in\mathsf bL$; and for $\gamma\in W(E)$
\begin{equation}\label{pf:eq-transl}
\Theta_{\psrho_f(\gamma)F^v}(z)=\zeta_\gamma\prod_{i,\mathsf c}\theta_{\mathsf x_{i,\mathsf c}
\Omega_0}\bigl(a_{i,\mathsf c}(z)\bigr)^{\mathsf b},\qquad\mathsf x_{i,\mathsf c}=-\pfT_{E/K}(
\aubar{u_i}\gamma_{\mathsf c}),
\end{equation}
with a root of unity $\zeta_\gamma$.
\end{lemma}
\begin{proof}
For $\lambda\in\mathsf b\Lambda_L$ the numbers $\omega_{i,\mathsf c}=a_{i,\mathsf c}(\lambda^\natural)
=\Omega_0\pfT_{E/K}(\aubar{v_i}\lambda_{\mathsf c}/\mathsf b)$ lie in $\Omega_0\pfT_{E/K}(\mathfrak
d_{E/K}^{-1})\subseteq\Lambda$. By \eqref{pa:eq-theta} and the two identities above,
$\Pi^v(z+\lambda^\natural)=\epsilon(\lambda)e^{\langle z,\lambda^\natural\rangle+|\lambda^\natural
|^2/2}\Pi^v(z)$ with $\epsilon(\lambda)=\prod_{i,\mathsf c}\chi(\omega_{i,\mathsf c})^{\mathsf b}$.
By Lemma~\ref{pf:chi}(b) and \eqref{pa:eq-res},
\[
\epsilon(\lambda)=(-1)^{\sum_{i,\mathsf c}\mathsf b\,|\pfT_{E/K}(\aubar{u_i}\lambda_{\mathsf c})|^2}
=(-1)^{\sum_{\mathsf c}\sum_\tau|\tau(\lambda_{\mathsf c})|^2}=(-1)^{\pfT_{E^+/\Q}\langle\lambda,
\lambda\rangle}=\epsilon_L(\lambda)
\]
(Lemma~\ref{pf:eps}). So Lemma~\ref{pk:dict} gives $F^v$, and \eqref{pf:eq-transl} is proved like
\eqref{pk:eq-transl}. The computation in the proof of Lemma~\ref{pk:dict} shows that
$\Theta_f(z+\lambda^\natural)=\kappa(\lambda)^{-1}e^{\langle z,\lambda^\natural\rangle+|\lambda^
\natural|^2/2}\Theta_f(z)$ whenever $\psrho_f(\lambda)f=\kappa(\lambda)f$, and conversely by the
injectivity of $f\mapsto\Theta_f$; so $\psrho_f(\lambda)F^v=\epsilon_L(\lambda)^{-1}F^v=\epsilon_L
(\lambda)F^v$ for $\lambda\in\mathsf b\Lambda_L$, and for $l\in\mathsf bL$ by continuity.
\end{proof}

\begin{lemma}[Values and expansion of the integral vacuum]\label{pf:fourier}
{\rm(a)} For every $x\in W(\auA_f)$ the number $\theta(1\otimes\psrho^\circ(x)\psvac^\circ)$ is an
algebraic integer.
{\rm(b)} Let $(v_i)$ be a tight frame of length $\mathsf b$. Then $\psvac^\circ=\sum_ja_j\,
\psrho_f(y_j)F^v$, a finite sum with $y_j\in\mathsf b^{-1}L$ and $\mathsf b^{2f\#\mathsf C}a_j$
algebraic integers.
\end{lemma}
\begin{proof}
(a) By \eqref{pk:eq-vac}, $\psrho^\circ(x)\psvac^\circ$ is a sum, with coefficients roots of unity,
of vectors $\psrho_f(\gamma)F_0$, $\gamma\in W(E)$ ($F_0$ is an eigenvector of $\psrho_f(q_1L)$
with eigenvalues $\pm1$, and $W(E)+q_1L=W(\auA_f)$). By \eqref{pk:eq-transl} the value of $\theta$
on $1\otimes\psrho_f(\gamma)F_0$ is a root of unity times a product of numbers $\theta_u(0)$,
$u\in\Q\Lambda$, which are algebraic integers by Proposition~\ref{pf:values}.

(b) Let $\mathsf V$ be the space of the $f\in\psHS_f(W)$ with $\psrho_f(l)f=\epsilon_L(l)f$ for
$l\in\mathsf bL$, that is, $\psrho^\circ(l)f=f$. By (W4) it has the basis $\psrho^\circ(x)\psvac_L$,
$x\in\mathsf b^{-1}L/L$ (these vectors lie in $\mathsf V$ because $[l,x]=1$, and they are
eigenvectors of $\psrho(L)$ for $\mathsf b^{2f\#\mathsf C}=\dim\mathsf V$ distinct characters). The
map $\psrho^\circ$ is a homomorphism on $L$, and $\mathsf P=\mathsf b^{-2f\#\mathsf C}\sum_{l\in
L/\mathsf bL}\psrho^\circ(l)$ is the projector of $\mathsf V$ onto $\C\psvac_L$ along the other
basis vectors. Write $F^v=\sum_xc_x\psrho^\circ(x)\psvac_L$ (Lemma~\ref{pf:semi}). Then $\mathsf
P\psrho^\circ(-x)F^v=\eta_xc_x\psvac_L$ with the root of unity $\eta_x=\epsilon(x)\epsilon(-x)
\beta(-x,x)$. Let $\psvac^\circ=\kappa\psvac_L$ and $V^\circ(x)=\theta(1\otimes\psrho^\circ(x)
\psvac^\circ)$. Applying $\theta$ to the expansion of $F^v$ gives $1=\sum_xc_x\kappa^{-1}V^\circ
(x)$, and therefore
\[
\psvac^\circ=\sum_xV^\circ(x)\,c_x\psvac_L=\sum_{x\in\mathsf b^{-1}L/L}\ \sum_{l\in L/\mathsf bL}
\frac{V^\circ(x)}{\eta_x\,\mathsf b^{2f\#\mathsf C}}\;\psrho^\circ(l)\psrho^\circ(-x)F^v .
\]
Each $\psrho^\circ(l)\psrho^\circ(-x)$ is a root of unity times $\psrho_f(l-x)$, and $V^\circ(x)$
is an algebraic integer by (a).
\end{proof}

For $\delta\in\auO_K$, $\delta\ne0$, a \emph{standard datum of level $\delta$ at $\ell$} is a
combination, with coefficients in $\iota^{-1}(\pao)$, of vectors $\psrho^\circ(x)\psvac^\circ$ with
$\delta x$ integral at the places above $\ell$ (that is, $\delta x_w\in L_w$ for $w\mid\ell$). The
standard data of \S\ref{pk:ss-kernel} are those of level $\aubar{\pi_M}$. \emph{If $g\in\auU(W)(\auA_f)$ is integral with integral inverse at the places above $\ell$, then
$\omega(g)$ preserves the standard data of level $\delta$ at $\ell$, for every $\delta$ and every $\ell$.}
Indeed, by (W1), $\omega(g)\psrho^\circ(x)\psvac^\circ$ is a root of unity times $\psrho^\circ(gx)\,
\omega(g)\psvac^\circ$, and $\delta gx$ is integral above $\ell$ when $\delta x$ is. By Lemma~\ref{pf:delta},
$\omega(g)\psvac_L=\delta_g\sum_{r\in gL/(L\cap gL)}\zeta_r\psrho^\circ(r)\psvac_L$ with roots of unity
$\zeta_r$; as $g_w L_w=L_w$ for $w\mid\ell$, the coefficient $\delta_g$ is a unit above $\ell$ and the $r$ may be taken in $L_w$ above
$\ell$ (part (ii) of the lemma), so that $\delta r$ is integral there ($\delta\in\auO_K$). This is the proof of
Proposition~\ref{pk:std}(S1), which does not use the level.

\begin{theorem}[Standard data at every prime]\label{pf:std}
Let $\ell$ be a prime, $\delta\in\auO_K$, $\delta\ne0$, and $f$ a standard datum of level $\delta$
at $\ell$. There are a resolution of the identity $(u'_j,m_j)_j$ of $E^{\mathsf C}$ with
$\alpha^{-1}u'_j$ integral above $\ell$, and an integral power series $\mathsf H$ in the series
$y(\mu_j(z))$, $\mu_j=(\varpi_{\mathrm{vac}}/\mathsf u)\,u_j'^{\dagger}$, such that
\[
\Theta_f\bigl(\aubar\delta\,s_0\,z\bigr)=\Bigl(\prod_j\theta(\mu_j(z))^{m_j}\Bigr)^{\auNm\delta}
\cdot\mathsf H(z),
\]
and the linear forms $\mu_j$ have integral coefficients. If $\ell\ne d$, the series
$\Theta_f(\aubar\delta s_0\,\cdot)$ is $\partial$-integral, and so is $\Theta_f(c\,\aubar\delta s_0
\,\cdot)$ for every $c\in\pao$.
\end{theorem}
\begin{proof}
Let $(v_i)$ be the frame $(\zeta_{q_1}^k)$, $\mathsf b=q_1$, if $\ell\nmid q_1$, and the frame of
Lemma~\ref{pf:frame}, $\mathsf b=\mathsf h$, if $\ell\mid q_1$; so $\mathsf b$ is a unit at
$\ell$. By Lemma~\ref{pf:fourier}(b), $f$ is a combination with integral coefficients of vectors
$\psrho_f(x')F^v$ with $\delta x'$ integral above $\ell$; each of them is a root of unity times
$\psrho_f(\gamma)F^v$ with $\gamma\in W(E)$, $\gamma-x'\in\mathsf bL$ (Lemma~\ref{pf:semi}), and
then $\delta\gamma$ is integral above $\ell$. Let $\mathsf m\ge1$ be an integer prime to $\ell$
with $\mathsf m\delta\gamma\in\Lambda_L$ for all the $\gamma$ that occur, and $\delta'=\mathsf m
\mathsf b\delta$. In \eqref{pf:eq-transl}, $\delta'\mathsf x_{i,\mathsf c}=-\pfT_{E/K}(\aubar{v_i}
\cdot\mathsf m\delta\gamma_{\mathsf c})\in\auO_K$, so $\delta'\cdot\mathsf x_{i,\mathsf c}\Omega_0
\in\Lambda$. Since $s_0\Omega_0/c_\psi=\varpi_{\mathrm{vac}}/\mathsf u=:\mathsf p$, we have
$a_{i,\mathsf c}(\aubar\delta s_0z)=\aubar{\delta'}\cdot\mathsf p\,u_i'^{\dagger}(z_{\mathsf c})$
with $u'_i=v_i/(\mathsf m\mathsf b^2)$, and Proposition~\ref{pf:pull} gives, for each factor,
\[
\theta_{\mathsf x_{i,\mathsf c}\Omega_0}\bigl(a_{i,\mathsf c}(\aubar\delta s_0z)\bigr)=\theta
\bigl(\mu_{i,\mathsf c}(z)\bigr)^{\mathsf m^2\mathsf b^2\auNm\delta}\cdot\Gamma\bigl(y(\mu_{i,
\mathsf c}(z))\bigr),\qquad\mu_{i,\mathsf c}(z)=\mathsf p\,u_i'^{\dagger}(z_{\mathsf c}),\quad
\Gamma\in\pao[[Y]].
\]
Taking the product over $(i,\mathsf c)$ with the exponents $\mathsf b$ we obtain the formula for
$\psrho_f(\gamma)F^v$, with the multiplicities $m_{i,\mathsf c}=\mathsf m^2\mathsf b^3$; the first
factor does not depend on $\gamma$, so the formula holds for $f$. The vectors $u'_i$ with these
multiplicities are a resolution of the identity: $\sum_i\mathsf m^2\mathsf b^3\tau(v_i)\aubar{\tau'
(v_i)}/(\mathsf m\mathsf b^2)^2=\delta_{\tau\tau'}$. The coefficients of $\mu_{i,\mathsf c}$ are
$\varpi_{\mathrm{vac}}\aubar{\tau(v_i)}/(\mathsf u\,\mathsf m\mathsf b^2)$; they are integral by
Lemma~\ref{pf:scale} and Corollary~\ref{pf:CS}. For $\ell\ne d$ the last assertion follows as in
Proposition~\ref{pf:pull}.
\end{proof}

\begin{corollary}[{\rm(K1)} at every prime; the prime $N$]\label{pf:K1}
\begin{enumerate}
\item[(a)] The integrality assertion of Lemma~\ref{pa:kernel}{\rm(K1)} holds for every prime
$\ell\ne d,N$ with $\ell\nmid\auNm Q_\vartheta$, in particular for $\ell=2$, for the primes
dividing $q_1$, and for $p_0,p_1$.
\item[(b)] Lemma~\ref{pl:kernelN} $(=$ {\rm(K1$_N$)}$)$ holds as stated, without the condition
{\rm(Gen)}.
\item[(c)] Proposition~\ref{pk:std}{\rm(S3)}, hence Lemma~\ref{pa:kernel}{\rm(K2)}, holds without
appeal to \cite{KubertLang}.
\end{enumerate}
\end{corollary}
\begin{proof}
(a) By the proof of Proposition~\ref{pk:K12}, $F=\Theta_f(\pi_Ms_0\,\cdot)$ for a unit multiple $f$
of a standard datum of level $\aubar{\pi_M}$ at $\ell$; apply Theorem~\ref{pf:std} with
$\delta=\aubar{\pi_M}$. (b) Let $\ell=N$ and let $\mathfrak n$ be the prime of $K$ induced by
$\iota$; it is the prime below the $\mathfrak l_\tau$, and $\mathfrak n\ne\aubar{\mathfrak n}$.
Choose $\delta\in\mathfrak n\setminus(\mathfrak n^2\cup\aubar{\mathfrak n})$ (possible by the Chinese remainder theorem), so that $\operatorname{ord}_{\mathfrak n}\delta=1$. Then $\iota(\delta)\in\pam$ and
$\iota(\aubar\delta)$ is a unit. The point $z$ of the lemma is integral at the primes above
$\aubar{\mathfrak n}$ and $Nz$ is integral at the primes above $\mathfrak n$, where $N$ and $\delta$
are both uniformisers ($N$ splits completely in $E$, and $\operatorname{ord}_{\mathfrak n}\delta=1$); so $\delta z$ is integral above $N$ and
$\phi^\circ_{j,z}$ is a standard datum of level $\delta$ at $N$. The reduction of
Proposition~\ref{pk:K12}, with $\mathsf k,\mathsf k'$ integral above $N$ ((G2) of
\S\ref{pl:ss-norm}) and (S1) at the level $\delta$, shows that the matrix in question is, up to
matrices invertible over $\pao$, the matrix of the jets along the entries of $S$ of
$\Theta_f(\pi_Ms_0\,\cdot)$ for a unit multiple $f$ of a standard datum of level $\delta$. By
Theorem~\ref{pf:std} with $c=\iota(\pi_M/\aubar\delta)\in\pao$ ($N\ne d$) this series is
$\partial$-integral, and the entries of $S$ are integral polynomials. (c) In the proof of
Proposition~\ref{pk:std}, \cite{KubertLang} is used only for the integrality of $\theta_u(0)$
(through (C3)(i)) and of $\theta_t'(0)$. The first is Proposition~\ref{pf:values}. For the second:
by Lemma~\ref{pf:chi}(c), $\theta_{-t}=\pm\theta_t$, so $\theta_t^2/\theta^2$ is $\Lambda$-periodic
with divisor $2[P_0]-2[O]$ and $\theta_t^2=\mathsf c_t\,x\,\theta^2$, $\mathsf c_t=\theta_t'(0)^2$.
Let $z\ne0$ be a torsion point of odd order prime to $\ell$. Multiplying by $e^{-H(z,z)}$ gives
$\zeta\,\pfbr{z+t}^2=\mathsf c_t\,x(z)\,\pfbr z^2$ with a root of unity $\zeta$. Here $z-t\notin
E_1$ and $z\notin E_1$, so $\pfbr z$ and $\pfbr{z+t}$ are units (Proposition~\ref{pf:values}), and
$x(z)$ is a unit because $Q(z)$ reduces neither to the origin nor to $\widetilde P_0$. Hence
$\theta_t'(0)$ is a unit, at every prime.
\end{proof}

\subsection{The fact {\rm(W6)}}\label{pf:ss-flip}

We use (W1), (W4) and (W2$'$) of \S\ref{pk:ss-conv}; the last one also in its local form: for a
finite place $v$ of $E^+$ and an isotropic $E_v$-subspace $I$ of $W_v=W\otimes E^+_v$, the local
representation $\psHS_v(W)$ has the zero term $\psZ_I\colon\psHS_v(W)\to\psHS_v(W_I)$, with
$\psZ_I\psrho(w)=\psrho(\bar w)\psZ_I$ for $w\in I^\perp$ and $\psZ_I\omega(p)=\chi_I(p)\omega_{W_I}
(\bar p)\psZ_I$, $\chi_I(p)=(\upsilon_v|\cdot|^{-1/2})(\det p|_I)^{-1}$, for $p$ in the stabiliser
$P_I(E^+_v)$. (This is how (W2) is obtained: the mixed model $\auS(I^*)\otimes\psHS_v(W_I)$ and
the formulas for $P_I$ are local, see the references given for Lemma~\ref{ps:weil}.) The map
$\psZ_I$ is, up to a scalar, the only non-zero continuous linear map with the first property:
such a map kills $(\psrho(i)-1)\psHS_v(W)$ for $i\in I$, hence factors through the evaluation at
$0$ of the mixed model, and $\psHS_v(W_I)$ is irreducible.

\begin{proposition}[{\rm(W6)}]\label{pf:flip}
Let $W$ be a hermitian space over $E$ that has a self-dual lattice $L\subset W(\auA_f)$, and let
$W'$ be the $E$-vector space $\aubar W$ $(W$ with the conjugate scalar multiplication$)$ with the
form $\langle w_1,w_2\rangle'=-\langle w_2,w_1\rangle$. Then $\beta'=\beta$ on $W(\auA)=W'(\auA)$,
$\auU(W)=\auU(W')$ as groups acting on this set, and if $\mathrm{Fl}\colon\psHS(W)\to\psHS(W')$ is
an isomorphism with $\mathrm{Fl}\,\psrho(w)=\psrho'(w)\,\mathrm{Fl}$ (it exists and is unique up
to a scalar), then
\[
\mathrm{Fl}\circ\omega_W(g)=\upsilon(\det\nolimits_Wg)^{-1}\cdot\omega_{W'}(g)\circ\mathrm{Fl}
\qquad(g\in\auU(W)(\auA)),
\]
where $\det_W$ is the determinant on the $E$-module $W$, and $\theta_{W'}\circ\mathrm{Fl}$ is a
multiple of $\theta_W$. For $W=\auHom(V,V')$ the map $w\mapsto w^\dagger$ identifies $W'$ with
$\auHom(V',V^-)$, and $\det_W(g,g')=\det(g')^{n}\det(g)^{-n'}$; this is {\rm(W6)} of
\S\ref{pk:ss-conv}.
\end{proposition}
\begin{proof}
The form $\langle\,,\rangle'$ is hermitian on $W'$, and $\beta'(x,y)=\psi(\tfrac12\jmath\langle
y,x\rangle')=\psi(-\tfrac12\jmath\langle x,y\rangle)=\psi(\aubar{\tfrac12\jmath\langle y,x\rangle})
=\beta(x,y)$, because $\psi\circ c=\psi$. So $\psrho'$ is another
realisation of the same Heisenberg representation and $\mathrm{Fl}$ exists; $\theta_{W'}\circ
\mathrm{Fl}$ is invariant under $\psrho(W(E))$, hence equals $c_\theta\theta_W$ with $c_\theta\ne0$.
By (W1) for $W$ and $W'$ the operator $\mathrm{Fl}^{-1}\omega_{W'}(g)^{-1}\mathrm{Fl}\,\omega_W(g)$
commutes with all $\psrho(w)$, so it is a scalar $\kappa(g)$, and $\kappa$ is a continuous
character of $\auU(W)(\auA)$. Put $\kappa_0(g)=\kappa(g)\upsilon(\det_Wg)$; we show $\kappa_0=1$.

\emph{Rational points.} For $\gamma\in\auU(W)(E^+)$, $c_\theta\theta_W(F)=c_\theta\theta_W(\omega_W
(\gamma)F)=\kappa(\gamma)\theta_{W'}(\omega_{W'}(\gamma)\mathrm{Fl}F)=\kappa(\gamma)c_\theta\theta_W
(F)$, so $\kappa(\gamma)=1$; and $\upsilon(\det\gamma)=1$ since $\upsilon$ is a Hecke character.

\emph{The stabiliser of $L$.} $L$ is self-dual in $W'$ as well, and $\epsilon'_L(l)=\psi_f(-\frac12\omega\jmath\langle l,l\rangle')=\epsilon_L(l)^{-1}=\epsilon_L(l)$
by Lemma~\ref{pf:eps}. So $\mathrm{Fl}\,\psvac_L\in\C^\times\psvac'_L$, and by (W4), $\kappa(k)=1$ for $k$
in the stabiliser $K_L$ of $L$; and $\upsilon(\det k)=1$ because $\det k$ is a unit idele and
$\upsilon$ is unramified.

\emph{Parabolic elements.} Let $v$ be a finite place, $I\subset W_v$ isotropic and $p\in P_I(E^+_v)$
acting trivially on $I^\perp/I$. The same subset $I$ is an isotropic subspace of $W'_v$, with
$\det(p|_I)$ replaced by its conjugate $\aubar{\mathsf d}$, $\mathsf d=\det\nolimits_W(p|_I)$. By
the uniqueness of the zero term, $\psZ'_I\mathrm{Fl}=c\,\mathrm{Fl}_I\psZ_I$ with $c\ne0$, where
$\mathrm{Fl}_I$ is the map $\mathrm{Fl}$ of $W_I$. Hence, by (W2$'$) for $W$ and $W'$,
\[
\chi_I(p)\,\psZ'_I\mathrm{Fl}=\psZ'_I\mathrm{Fl}\,\omega_W(p)=\kappa(p)\,\psZ'_I\omega_{W'}(p)
\mathrm{Fl}=\kappa(p)\chi'_I(p)\,\psZ'_I\mathrm{Fl},
\]
so $\kappa(p)=\chi_I(p)/\chi'_I(p)=\upsilon(\aubar{\mathsf d})/\upsilon(\mathsf d)=\upsilon(\mathsf
d)^{-2}$, because $\upsilon(\mathsf d\aubar{\mathsf d})=1$ ($\upsilon$ restricted to $\auA_{E^+}^
\times$ is the character of $E/E^+$, trivial on norms). On the other hand $p$ acts on $W/I^\perp$,
the conjugate dual of $I$, with determinant $\aubar{\mathsf d}^{-1}$; so $\det_Wp=\mathsf d/\aubar{
\mathsf d}$ and $\upsilon(\det_Wp)=\upsilon(\mathsf d)^2$. Thus $\kappa_0(p)=1$.

\emph{All finite places.} Let $v$ be a finite place; $E/E^+$ is unramified at $v$. If $v$ is
inert, $L_v$ has an orthonormal basis: if $\langle x,y\rangle$ is a unit and $\langle x,x\rangle$,
$\langle y,y\rangle$ are not, then $x+\lambda y$ has unit norm when $\pfT(\lambda\langle x,y
\rangle)$ is a unit, and the trace of $\auO_{E,v}$ is surjective; a vector of unit norm can be
scaled to norm $1$, the norm of units being surjective; split it off and induct. Let $I_L$ be a
maximal totally isotropic direct summand of $L_v$. By (R1) of \S\ref{pk:ss-rest},
$L_v=(I_L\oplus I_L^*)\perp L_0$ with $I_L^*$ totally isotropic, and $L_0$ has rank $\le1$: a
unimodular lattice of rank $\ge2$ contains the primitive isotropic vector $e_1+\lambda e_2$,
$\lambda\aubar\lambda=-1$, for orthonormal $e_1,e_2$ (if $v$ splits, $L_0=0$). Let $I=E_vI_L$. By
the Iwasawa decomposition ((R1); \cite[\S3.3.2, p.~51]{TitsCorvallis}), $\auU(W)(E^+_v)=P_I(E^+_v)K_{L,v}$,
and $P_I(E^+_v)$ is generated by the elements acting trivially on $I^\perp/I$ and by
$\auU(E_vL_0)$, which consists of unit scalars on $L_0$ and lies in $K_{L,v}$. So $\kappa_0=1$ on
$\auU(W)(E^+_v)$, and by continuity on $\auU(W)(\auA_f)$.

\emph{The infinite places.} $\kappa_0$ is now a continuous character of $\auU(W)(E^+\otimes\R)$
that is trivial on the image of $\auU(W)(E^+)$, which is dense by Proposition~\ref{pl:wa}. So
$\kappa_0=1$. The last assertions are those of \S\ref{pk:ss-conv}: $\langle w_1^\dagger,w_2^\dagger
\rangle=-\langle w_2,w_1\rangle$ and $w\mapsto w^\dagger$ is antilinear.
\end{proof}
The hypothesis on $L$ holds for all spaces $\auHom(V,V')$ to which (W6) is applied (all spaces of
{\cite[\S14.3]{P46}} carry unimodular lattices $\psM$, and $\alpha\auHom(\psM,\psM')$ is
self-dual).

\begin{remark}[The lattice preserved by the flip]\label{pf:rem-flip-lattice}
The proof shows that $\mathrm{Fl}$ carries the vacuum of $L$ to a multiple of the vacuum of the
\emph{same} set $L$, viewed in $W'$. For $W=\auHom(V,V')$ and $L=\alpha\auHom(\psM,\psM')$ the
identification $w\mapsto w^\dagger$ of $W'$ with $\auHom(V',V^-)$ carries $L$ to
$L^\dagger=\aubar\alpha\auHom(\psM',\psM)=u\,\alpha\auHom(\psM',\psM)$, $u=\aubar a/a$ (with $a$ as in
\S\ref{ps:ss-weil}; {\cite[Lemma~6.32(f)]{P53}}), and not, in general, to the standard lattice
$\alpha\auHom(\psM',\psM)$ of $\auHom(V',V^-)$. Here $u$ is a finite idele with $u\aubar u=1$, acting on
$\auHom(V',V^-)$ through the centre of $\auU(V^-)=\auU(V)$, and $u_v=1$ at every place $v$ not above
$q_1$. By {\cite[Lemma~6.32(e)]{P53}},
\[
\mathrm{Fl}\,\psvac_{\alpha\auHom(\psM,\psM')}\in\C^\times\,\omega_{\auHom(V',V^-)}(1,u)\,
\psvac_{\alpha\auHom(\psM',\psM)}
\]
($\omega(1,u)$ meaning $u$ acting through $\auU(V^-)$). So $\mathrm{Fl}$ may be normalised to send
vacuum to vacuum at the places not above $q_1$, but at the places above $q_1$ it sends the standard
vacuum to the translate by $\omega(1,u)$ of the standard vacuum. In an integral over $[\auU(V)]$
this translate changes the integral only by the factor of modulus one of {\cite[Lemma~6.32(e)]{P53}}; in a
kernel at a fixed point it shifts the argument in $\auU(V)$ by $u$ (as in {\cite[Proposition~4.40]{P53}} and
{\cite[Lemma~5.37]{P53}}).
\end{remark}

\subsection{The orbits at the end of Lemma~\ref{ps:support}}\label{pf:ss-orbits}

\begin{lemma}\label{pf:orbits}
Let $1\le j<m$ and let $v$ be a finite place of $E^+$. Let $\mathcal X_v$ be the set of the
lattices $\Lambda'_v=y\psM_{j,v}$, $y\in G_j(E^+_v)$, that are the direct sum of their intersections
with the grades and satisfy the bounds of Lemma~\ref{ps:support} at $v$ $(\mathsf q^{\,j}\Lambda'^-\subseteq
\psM_j^-$, $\aubar{\mathsf q}^{\,j}\psM_j^-\subseteq\Lambda'^-$, with $\Lambda'^-$ the intersection
with the span of the whole string$)$. Then $\auU(X_j)(E^+_v)$ has finitely many orbits in
$\mathcal X_v$, and exactly one, that of $\psM_{j,v}$, if $v$ does not divide $q=\auNm Q$. Hence
$\mathfrak B_j=R_j(\auA_f)\mathsf C_jK_j$ with $\mathsf C_j$ finite, consisting of elements that are
$1$ outside $q$, and $\mathsf C_j=\{1\}$ if $Q=\emptyset$.
\end{lemma}
\begin{proof}
Let $\Lambda'=\Lambda'_v\in\mathcal X_v$; it is unimodular. For a grade $g\ne0$ its part
$\Lambda'_g$ is a lattice in the line $E_vf_a$, between $\aubar{\mathsf q}^{\,j}\psM_{j,g}$ and
$\mathsf q^{-j}\psM_{j,g}$: finitely many possibilities, and $\Lambda'_g=\psM_{j,g}$ if $v\nmid q$.
The sum of the parts of non-zero grade is unimodular (it is so for $\psM_j$ and, in general, it is
a sum of pairs of dual lines; unimodularity of $\Lambda'$ forces $\Lambda'_{-g}$ to be the dual of
$\Lambda'_g$), so $\Lambda'=\Lambda'_0\perp\bigoplus_{g\ne0}\Lambda'_g$ with $\Lambda'_0$
unimodular in the space $Y_{j,0}$ of grade $0$. The group $\auU(X_j)$ acts on $Y_{j,0}$ and fixes
the string; so we must classify the $\Lambda'_0$.

\emph{$j$ odd.} $Y_{j,0}=X_j$, and the unimodular lattices of $X_{j,v}$ form one orbit: if $v$
splits, $v=w\aubar w$, they are the pairs $(B,B^\#)$ with $B$ a lattice in $X_{j,w}$, and
$\auU(X_j)(E^+_v)=\auGL(X_{j,w})$ is transitive on lattices; if $v$ is inert, every unimodular
lattice has an orthonormal basis (proof of Proposition~\ref{pf:flip}), and two such lattices in
the same space differ by an isometry.

\emph{$j$ even.} $Y_{j,0}=X_j\perp E z$, $\langle z,z\rangle=-N$. Put $\mathfrak z=\Lambda'_0\cap
E_vz$; the projection of the unimodular lattice $\Lambda'_0$ to $E_vz$ is the dual $\mathfrak
z^\#$ of $\mathfrak z$ in $E_vz$. The bounds leave finitely many possibilities for $\mathfrak z$,
and $\mathfrak z=\psM_j\cap E_vz$ if $v\nmid q$ (the bounds then say $\Lambda'^-=\psM_j^-$). If $v\nmid Nq$, then $\mathfrak z$ is unimodular, so
$\Lambda'_0=(\Lambda'_0\cap X_j)\perp\mathfrak z$ with $\Lambda'_0\cap X_j$ unimodular, and we
conclude as for $j$ odd. Let $v=w\aubar w$ be split (the places above $N$ and above $q$ are). Then
$\Lambda'_0=(B,B^\#)$ with $B$ a lattice in $X_{j,w}\oplus E_wz$, and the conditions are $B\cap
E_wz=\mathfrak z_w$, $\operatorname{pr}_z(B)=\mathfrak z'_w:=(\mathfrak z_{\bar w})^\#$. Such a
$B$ is the same as a triple $(B_1,B_2,\phi)$: lattices $B_1=B\cap X_{j,w}\subseteq B_2=
\operatorname{pr}_X(B)$ and an isomorphism $\phi\colon B_2/B_1\to\mathfrak z'_w/\mathfrak z_w$, with
$B=\{(b,\xi)\in B_2\oplus\mathfrak z'_w:\phi(b)\equiv\xi\}$. The module $\mathfrak z'_w/\mathfrak
z_w$ is cyclic, $\cong\auO_w/\varpi^e$, and $\auGL(X_{j,w})$ is transitive on the pairs
$B_1\subseteq B_2$ with $B_2/B_1\cong\auO_w/\varpi^e$ (elementary divisors); the stabiliser of a
pair maps onto $\operatorname{Aut}(B_2/B_1)=(\auO_w/\varpi^e)^\times$ (choose a basis of $B_2$ whose
last vector generates $B_2/B_1$ and whose other vectors lie in $B_1$, and multiply the last vector
by a unit). So the triples with given $\mathfrak z_w\subseteq\mathfrak z'_w$ form one orbit. For
$v\mid N$, $v\nmid q$, the pair $(\mathfrak z_w,\mathfrak z'_w)$ is that of $\psM_j$ (and $e=1$);
for $v\mid q$ there are finitely many pairs.

\emph{The set $\mathfrak B_j$.} For each $v\mid q$ choose representatives $c_v\psM_{j,v}$,
$c_v\in G_j(E^+_v)$, of the orbits in $\mathcal X_v$, and let $\mathsf C_j$ be the set of the
elements $(c_v)_{v\mid q}$, completed by $1$. If $x\in\mathfrak B_j$ and $r\in R_j(\auA_f)$ is as
in the definition, then $rx\psM_j\in\prod_v\mathcal X_v$, so $u\,rx\psM_j=c\psM_j$ for some $u\in\auU(X_j)(\auA_f)\subset R_j(\auA_f)$ and
$c\in\mathsf C_j$, and $x\in R_j(\auA_f)cK_j$. Conversely $rck\psM_j=r(c\psM_j)$.
\end{proof}

\subsection{The number of classes}\label{pf:ss-classes}

\begin{proposition}\label{pf:classes}
Let $X=E^n$ be a standard space with $n\le n_m$, $G=\auU(X)$, $\Gamma=G(E^+)$, $K^\circ$ the
stabiliser of $\widehat\auO_E^{\,n}$ in $G_f$, and for an integer $\mathsf G\ge1$ let $K(\mathsf G)$
be the group of the $k\in K^\circ$ with $k\equiv1$ modulo $\mathsf G$. Then
\[
\#\bigl(\Gamma\backslash G_f/K^\circ\bigr)\le(2\mathsf B+1)^{4fn^2},\qquad\mathsf B=n^2C_{\mathrm{sb}}^2
\cdot2f\sqrt d,\quad C_{\mathrm{sb}}=6q_1d(fn)^2,
\]
and $\#(\Gamma\backslash G_f/J)\le(2\mathsf B+1)^{4fn^2}\mathsf G^{2fn^2}$ for every subgroup
$J\supseteq K(\mathsf G)$ of $G_f$. For $\mathsf G=MN$ this is at most $\exp(x^{30})$. In
particular the levels $J_{X_j}$ $(j\le m)$, and the level $J_1$ of the target forms for the trivial
twist, have at most $\exp(x^{30})$ classes.
\end{proposition}
\begin{proof}
\emph{Two elementary facts.} (i) Every $\mathsf y\in\auA_{E,f}$ is congruent modulo
$\widehat\auO_E$ to an $\mathsf a\in E$ with $|\tau(\mathsf a)|\le\mathsf B_0:=2f\sqrt d$ for all
embeddings $\tau$: indeed $\auA_{E,f}=E+\widehat\auO_E$, and $\auO_E=\auO_K[\zeta_{q_1}]$ has the
$\Z$-basis $\omega^e\zeta_{q_1}^k$ $(e\le1$, $k<f)$, whose conjugates have absolute value at most
$|\omega|\le\sqrt d$; reduce the coordinates of $\mathsf a$ to $[0,1)$. (ii) For $B>0$ the number
of $\mathsf a\in\auO_E$ with $|\tau(\mathsf a)|\le B$ for all $\tau\in\auPl$ is at most
$(2B+1)^{2f}$: two distinct such elements have $\prod_\tau|\tau(\mathsf a-\mathsf a')|^2=\auNm(
\mathsf a-\mathsf a')\ge1$, so the polydiscs of radius $\frac12$ around the points $(\tau(\mathsf
a))_\tau\in\C^{\auPl}$ are disjoint, and they lie in the polydisc of radius $B+\frac12$.

\emph{The invariant of a class.} Let $g\in G_f$. By {\cite[Lemma~14.21]{P46}} with $c_0=q_1d$ there
is $Y\in\auGL_n(E)$ with $A:=g^{-1}Y\in M_n(\widehat\auO_E)$ and all conjugates of all entries
bounded by $C_{\mathrm{sb}}$. By (i) there is $X\in M_n(E)$ with $A^{-1}-X\in M_n(\widehat\auO_E)$
and entries bounded by $\mathsf B_0$. Put $\mathsf G_g=Y^*Y$ and $\mathsf I(g)=(\mathsf G_g,\mathsf
G_gX)$. As $g$ is unitary, $\mathsf G_g=A^*A$ and $\mathsf G_gX=A^*\bigl(1-A(A^{-1}-X)\bigr)$ have
entries in $E\cap\widehat\auO_E=\auO_E$, and their conjugates are bounded by $nC_{\mathrm{sb}}^2$
and by $n\cdot nC_{\mathrm{sb}}^2\cdot\mathsf B_0=\mathsf B$. By (ii) the number of possible
$\mathsf I(g)$ is at most $(2\mathsf B+1)^{2f\cdot2n^2}$.

\emph{The invariant determines the class.} Let $g,g'$ with data $(Y,A,X)$, $(Y',A',X')$ and
$\mathsf I(g)=\mathsf I(g')$. Then $Y^*Y=Y'^*Y'$, so $\gamma=YY'^{-1}\in\Gamma$; and $X=X'$, so
$A'A^{-1}=A'(A^{-1}-X)+A'X'$ is integral, since $A'X'=1-A'(A'^{-1}-X')$. But $A'A^{-1}=g'^{-1}\gamma^{-1}g=:k$ is unitary, so
$k^{-1}=k^*$ is integral too, $k\in K^\circ$, and $g=\gamma g'k$. This proves the first bound.

\emph{Levels.} The map $\Gamma\backslash G_f/(J\cap K^\circ)\to\Gamma\backslash G_f/J$ is onto, and
the fibres of $\Gamma\backslash G_f/(J\cap K^\circ)\to\Gamma\backslash G_f/K^\circ$ have at most
$[K^\circ:J\cap K^\circ]\le[K^\circ:K(\mathsf G)]\le\#M_n(\auO_E/\mathsf G)=\mathsf G^{2fn^2}$
elements.

\emph{Sizes.} By {\cite[(38)]{P46}} and the definitions, $f\le x^{12}$, $n\le n_m\le40x^6$, and
$6q_1d(fk)^2\le x^{100}$ (the inequality used in the proof of {\cite[Proposition~14.22]{P46}}), and
$C_{\mathrm{sb}}=6q_1d(fn)^2\le x^{102}$ (Lemma~\ref{pz:sizes}(b), a numerical inequality proved there from {\cite[Lemma~13.4]{P46}} without the present proposition, which only the last sentence of that lemma applies), $\mathsf B_0\le C_{\mathrm{sb}}$, $n\le
C_{\mathrm{sb}}$ and $2\mathsf B+1\le C_{\mathrm{sb}}^6$; hence $4fn^2\log(2\mathsf B+1)\le4\cdot
x^{12}\cdot1600x^{12}\cdot612\log x\le x^{26}$. Also $\log(MN)\le22\log b+30m\log b\le x$ by
{\cite[(36)]{P46}}, so $2fn^2\log(MN)\le3200x^{25}\le x^{26}$. The levels named contain $K(MN)$:
they are defined inside $K^\circ$ by congruences modulo the chosen places, which divide $M$, and
modulo primes dividing $N$ ({\cite[\S14.3]{P46}}).
\end{proof}

\subsection{The published results used in this section}\label{pf:ss-status}
Proposition~\ref{pk:curve}(C3)(ii) is Theorem~\ref{pf:recip}; Lemma~\ref{ps:conj} and
{\cite[Proposition~14.23]{P46}} follow (Corollary~\ref{pf:conj}); Lemma~\ref{pa:kernel} holds in full,
(K1) at every prime $\ell\ne d,N$ and (K1$_N$), that is, Lemma~\ref{pl:kernelN}, without restriction
(Corollary~\ref{pf:K1}); (W6) is Proposition~\ref{pf:flip}; the number of classes is bounded by
Proposition~\ref{pf:classes}. The published results used are: integrality of
$j$ and potential good reduction, Deuring's criterion, the formal group of a Weierstrass equation
and the kernel of reduction~\cite{SilvermanAEC}, \cite[Chapter~13, Theorem~12]{LangEF}; the oddness of the class number of
\cite{CoxPrimes}; Euler's formula for the dual basis~\cite[III, \S6, Lemma~2]{SerreLF}; the Iwasawa
decomposition~\cite[\S3.3.2, p.~51]{TitsCorvallis}; and the facts (W1), (W2$'$) (also in its local form), (W4) on
the Weil representation.




\providecommand{\auO}{\mathcal O}
\providecommand{\auU}{\mathrm U}
\providecommand{\auGL}{\mathrm{GL}}
\providecommand{\auHom}{\operatorname{Hom}}
\providecommand{\auA}{\mathbb A}
\providecommand{\auF}{\mathbb F}
\providecommand{\auPl}{\Phi}
\providecommand{\auPa}{\Phi_a}
\providecommand{\auS}{\mathcal S}
\providecommand{\auSram}{\mathcal S^{\mathrm{ram}}}
\providecommand{\auAz}{\mathcal A_0}
\providecommand{\auG}{\overline G}
\providecommand{\auP}{\mathcal P}
\providecommand{\auNm}{\operatorname{Nm}}
\providecommand{\auAut}{\operatorname{Aut}}
\providecommand{\aubar}[1]{\overline{#1}}
\providecommand{\aueu}{\varepsilon_{\mathrm u}}
\providecommand{\aurd}{\mathrm{rd}_E}
\providecommand{\psM}{\mathsf M}
\providecommand{\psHS}{\mathcal H}
\providecommand{\psrho}{\varrho}
\providecommand{\pzF}{\mathfrak F}

\emph{The coefficient $\delta_g$.} The following lemma is used in the proofs of Propositions~\ref{pk:conj} and~\ref{pk:std}; it rests on (W1), (W4), the local form of (W2$'$) recalled in \S\ref{pf:ss-flip}, and the Cartan decomposition.

\begin{lemma}[The coefficient of a translate of the vacuum]\label{pf:delta}
Let $L\subset W(\auA_f)$ be a self-dual lattice, $K_L$ its stabiliser, $g\in\auU(W)(\auA_f)$ and
$\mathsf I_g=[gL:L\cap gL]$. Then
\[
\omega(g)\psvac_L=\delta_g\sum_{r\in gL/(L\cap gL)}\zeta_r\,\psrho^\circ(r)\psvac_L ,
\]
with roots of unity $\zeta_r$, $\zeta_0=1$, and
\[
\delta_g=\mathsf I_g^{-1/2}\,\upsilon(\mathsf d_g)^{-1}
\]
for a finite idele $\mathsf d_g$ of $E$ whose component is $1$ at every place $w$ with $g_wL_w=L_w$.
Consequently: \textup{(i)} $\delta_g$ is algebraic and $|\sigma\delta_g|=\mathsf I_g^{-1/2}$ for every
automorphism $\sigma$ of $\C$; \textup{(ii)} if $g_wL_w=L_w$ at the places $w$ above a prime $\ell$, then
$\delta_g$ is a unit at every place above $\ell$, and the representatives $r$ can be taken in $L_w$ at
the places $w$ above $\ell$.
\end{lemma}
\begin{proof}
\emph{The expansion.} By (W1) and (W4), $\psrho(x)\omega(g)\psvac_L=\epsilon_L(g^{-1}x)\,\omega(g)\psvac_L$ for
$x\in gL$. Write $\omega(g)\psvac_L=\sum_xc_x\psrho^\circ(x)\psvac_L$ in the basis of (W4), $x\in W(\auA_f)/L$.
For $l\in L\cap gL$, $\psrho(l)\psrho^\circ(x)\psvac_L=[l,x]\epsilon_L(l)\psrho^\circ(x)\psvac_L$, so
$c_x\ne0$ implies $[l,x]=\epsilon_L(g^{-1}l)\epsilon_L(l)^{-1}$ for all $l\in L\cap gL$. By the proof of
Lemma~\ref{pf:eps}, $\epsilon_L(l)$ depends only on $\langle l,l\rangle$, which $g$ preserves; so
$[l,x]=1$ for all $l\in L\cap gL$, that is, $x$ lies in the dual $L+gL$ of $L\cap gL$. Thus the sum is over
$(L+gL)/L=gL/(L\cap gL)$. For $r\in gL$ the operator $\psrho(r)$ maps the line of
$\psrho^\circ(x)\psvac_L$ to that of $\psrho^\circ(x+r)\psvac_L$, multiplying the basis vectors by roots of
unity (values of $\epsilon$ and $\beta$), and it multiplies $\omega(g)\psvac_L$ by $\epsilon_L(g^{-1}r)=\pm1$;
comparing the components, $c_r$ is a root of unity times $c_0$. This is the expansion, with
$\delta_g=c_0$. The $\mathsf I_g$ vectors $\psrho^\circ(r)\psvac_L$ are orthogonal of norm $\|\psvac_L\|$
and $\omega(g)$ is unitary, so $|\delta_g|=\mathsf I_g^{-1/2}$.

\emph{Reductions.} For $k\in K_L$, $\omega(k)\psvac_L=\psvac_L$ (W4), so $\delta_{gk}=\delta_g$; and
$\omega(k)\psrho^\circ(r)\psvac_L$ is a root of unity times $\psrho^\circ(kr)\psvac_L$ (W1), with the class
of $0$ fixed, so $\delta_{kg}=\delta_g$; also $\mathsf I_{kgk'}=\mathsf I_g$. If $g$ and $g'$ are supported at
disjoint sets of places, then $\omega(gg')\psvac_L=\delta_{g'}\sum_{r'}\zeta_{r'}\,\omega(g)\psrho^\circ(r')
\psvac_L$ and $\omega(g)\psrho^\circ(r')\psvac_L$ is a root of unity times $\psrho^\circ(r')\omega(g)\psvac_L$
($gr'=r'$); expanding $\omega(g)\psvac_L$ and noting that $r+r'\in L$ only if $r,r'\in L$, the coefficient of
$\psvac_L$ is $\delta_g\delta_{g'}$. So $\delta_{gg'}=\delta_g\delta_{g'}$ and $\mathsf I_{gg'}=\mathsf I_g
\mathsf I_{g'}$, and it suffices to treat $g\in\auU(W)(E^+_v)$ for one finite place $v$.

\emph{The Cartan decomposition.} As in the proof of Proposition~\ref{pf:flip}, $L_v=(I_L\oplus I_L^*)\perp
L_0$ with $I_L$ a maximal totally isotropic direct summand, $I_L^*$ totally isotropic, in duality with
$I_L$, and $L_0$ of rank $\le1$ (if $v$ splits into $w,\aubar w$, then $I_L=L_w$, $I_L^*=L_{\aubar w}$,
$L_0=0$). Put $I=E_vI_L$, $I^*=E_vI_L^*$. Since $E/E^+$ is unramified at $v$, $K_{L,v}$ is hyperspecial,
and by the Cartan decomposition \cite[\S3.3.3]{TitsCorvallis} every element of $\auU(W)(E^+_v)$ is in
$K_{L,v}\,a\,K_{L,v}$ for an $a$ that is diagonal, with entries powers of a uniformiser, in a basis of
$I_L$, acts on $I^*$ by the inverse of the adjoint, and is $1$ on $E_vL_0$ (at a split place this is the
theorem of elementary divisors for $\auGL(L_w)$). By the reductions we may take $g=a$. Then
$aL_v=aI_L\oplus aI_L^*\oplus L_0$.

\emph{The zero term.} Let $\psZ_I$ be the local zero term (\S\ref{pf:ss-flip}). Since $a\in P_I(E^+_v)$
acts trivially on $I^\perp/I$,
\[
\psZ_I\,\omega(a)\psvac_L=\chi_I(a)\,\psZ_I\psvac_L,\qquad\chi_I(a)=\upsilon_v(\det a|_I)^{-1}\,|\det a|_I|^{1/2}.
\]
On the other side take representatives $r=i+u$ with $i\in aI_L$, $u\in aI_L^*$. For $l\in I_L$ we have
$\psZ_I\psrho(l)=\psZ_I$ ($l\in I$) and $\psrho(l)\psvac_L=\epsilon_L(l)\psvac_L=\psvac_L$ ($l$ is isotropic,
Lemma~\ref{pf:eps}), so $\psZ_I\psrho(r)\psvac_L=\psZ_I\psrho(r)\psrho(l)\psvac_L=[r,l]\,\psZ_I\psrho(l)
\psrho(r)\psvac_L=[u,l]\,\psZ_I\psrho(r)\psvac_L$ (as $[i,l]=1$). If $u\notin I_L^*$ there is $l\in I_L$ with
$[u,l]\ne1$ ($I_L^*$ is the dual of $I_L$ in $I^*$), and the term vanishes. If $u\in I_L^*$ the class of
$r$ is that of $i\in aI_L$; then $\psrho(i)\omega(a)\psvac_L=\epsilon_L(a^{-1}i)\omega(a)\psvac_L=\omega(a)
\psvac_L$, and comparing the components on the line of $\psrho(i)\psvac_L$ gives $\zeta_i\psrho^\circ(i)
\psvac_L=\psrho(i)\psvac_L$, whence $\psZ_I(\zeta_i\psrho^\circ(i)\psvac_L)=\psZ_I\psvac_L$. Therefore
\[
\chi_I(a)\,\psZ_I\psvac_L=\delta_a\,\mathsf N\,\psZ_I\psvac_L,\qquad\mathsf N=[aI_L:I_L\cap aI_L].
\]
Finally $\psZ_I\psvac_L\ne0$: in the local mixed model $\psZ_I\psrho(u)F$, $u\in I^*$, is, up to a factor of absolute value $1$, the value $F(u)$ ($I^*$ acts by translation of the variable); by the computation just made (with $i=0$) $\psZ_I\psrho(u)\psvac_L$ is $0$ for
$u\notin I_L^*$ and $\psZ_I\psvac_L$ for $u\in I_L^*$; so $\psZ_I\psvac_L=0$ would give $\psvac_L=0$. Hence
$\delta_a=\upsilon_v(\det a|_I)^{-1}\,|\det a|_I|^{1/2}\,\mathsf N^{-1}$, and since $|\delta_a|=\mathsf I_a^{-1/2}$ and $|\upsilon_v|=1$, the
positive number $|\det a|_I|^{1/2}\mathsf N^{-1}$ is $\mathsf I_a^{-1/2}$. This is the formula, with
$\mathsf d_g$ the idele whose component at $v$ is $\det a|_I$ (at a split place, $\det a|_{L_w}$ at $w$ and
$1$ at $\aubar w$), which is $1$ when $g_v\in K_{L,v}$.

\emph{Consequences.} Let $\mathsf d$ be a finite idele and $h$ the class number of $E$; then
$\mathsf d^h=zu$ with $z\in E^\times$ and $u$ a unit idele, and as $\upsilon$ is unramified and trivial on
$E^\times$, $\upsilon(\mathsf d)^h=\upsilon_\infty(z)^{-1}=\prod_{\tau}|\tau(z)|/\tau(z)=\mathrm N_{E/\Q}(z)^{1/2}\big/\prod_{\tau\in\Phi}\tau(z)$, the product over the
embeddings of $E$ fixed in \S\ref{ps:ss-weil}, which form a CM type $\Phi$. This is an algebraic number. An automorphism $\sigma$ of $\C$ carries $\Phi$
to a CM type $\sigma\Phi$ ($E$ is a CM field) and $\mathrm N_{E/\Q}(z)^{1/2}$ to $\pm\mathrm N_{E/\Q}(z)^{1/2}$,
and $|\prod_{\tau\in\Phi'}\tau(z)|^2=\mathrm N_{E/\Q}(z)$ for every CM type $\Phi'$; so every conjugate of $\upsilon(\mathsf d)$ has absolute
value $1$, which with $\sigma(\mathsf I_g^{-1/2})=\pm\mathsf I_g^{-1/2}$ gives (i). For (ii): $\mathsf I_g$ is a
product of local indices at places not above $\ell$, hence prime to $\ell$; $\mathsf d_g$ is $1$ above
$\ell$, so $z$ is a unit above $\ell$ and $\upsilon(\mathsf d_g)^h$, a quotient of two numbers that are
units above $\ell$, is a unit above $\ell$; and $(gL/(L\cap gL))_w=0$ for $w\mid\ell$.
\end{proof}

\section{Archimedean types, constants and levels}\label{pz:sec}

This section proves the remaining inputs of the Short Forms Theorem ({\cite[Theorem~14.13]{P46}}) not treated in
\S\S\ref{pa:sec}--\ref{pf:sec} and in~\cite{P50}. (1) In \S\ref{ps:ss-Y} it was taken ``as part of the definition of the lower
chain'' that the lifts $\Theta^Y_{j,\vartheta}$ and their transposes preserve the archimedean type
and that the forms concerned are cusp forms. In \S\ref{pz:ss-types} the types are defined, and the
assumption is \emph{reduced} to three statements about discrete series of real unitary groups,
(A1)--(A3) below; everything global is proved (Proposition~\ref{pz:types}). Statements (A1) and
(A2) are theorems of Harish-Chandra, Hecht--Schmid and Mili\v ci\'c; (A3) is a case of the
correspondence of discrete series under the theta correspondence. They are treated in \S\ref{par:sec} and \S\ref{pco:sec}: formula {\cite[(59)]{P51}} is proved (Theorem~\ref{pco:FJ}) and the normalisation (V) is proved from the local properties (W1$_\infty$), (W2$_\infty$), (M), (S) of the archimedean Weil representation (Theorem~\ref{pco:V}). (2) \S\ref{pz:ss-const} makes the constant $C_0$ of
{\cite[Proposition~14.23]{P46}} explicit and settles its relation to the definition of admissible data.
(3) \S\ref{pz:ss-levels} defines the levels of the spaces $X_j$, checks the two facts quoted in the
proof of Proposition~\ref{pf:classes}, and proves the estimate used in the proof of
{\cite[Proposition~14.22]{P46}}. (4) \S\ref{pz:ss-audit} lists every statement on
which {\cite[Theorem~14.13]{P46}} depends, with the place of its proof.

\subsection{Forms of a discrete series type and the lower chain}\label{pz:ss-types}

\emph{The groups and the parameters.} Let $0\le j\le m$, $G_j=\auU(Y_j)$, $N_j=r+jt$ its rank,
$G_{j,\infty}=G_j(E^+\otimes\R)=\prod_{\tau\in\auPl}G_{j,\tau}$, and $K_{j,\infty}$ the maximal compact
subgroup that stabilises, at every $\tau$, the positive and the negative part of an orthogonal
basis of $Y_j\otimes_{E,\tau}\C$. Through $\tau$ the group $G_{j,\tau}$ is the unitary group of a
hermitian form with $N_j-Q_j$ signs $+$ and $Q_j=\lceil(m-1-j)/2\rceil$ signs $-$
({\cite[\S14.3]{P46}}). For $x\in G_{j,\infty}$ let $\|x\|=\prod_\tau\|x_\tau\|$, the product of the
operator norms on $Y_j\otimes_{E,\tau}\C$ for the positive form attached to that basis; then
$\|x\|\ge1$ and $\|x^{-1}\|=\|x\|$. Put $\kappa_l=\kappa+(l+\tfrac12)\Delta$ as in
{\cite[\S14.1]{P46}}. The \emph{parameter} of $Y_j$ at $\tau$ is the following set
$\Lambda^{(j)}_\tau$ of $N_j$ numbers, each with a sign:
\begin{equation}\label{pz:eq-param}
\begin{aligned}
&\tfrac k2-1,\ \tfrac k2-2,\ \dots,\ 1-\tfrac k2&&(k-1\text{ numbers, sign }+);\\
&-\kappa(\tau)-i\Delta\quad(1\le i\le m)&&(\text{sign }-\text{ if }i\ge j+2\text{ and }i\equiv j\!\!\pmod 2,\text{ else }+);\\
&-\kappa_l(\tau)+\tfrac{t-1}2-b\quad(1\le l\le j,\ 0\le b<t)&&(\text{sign }+);
\end{aligned}
\end{equation}
all numbers being lowered by $\tfrac12$ if $j$ is odd. For $j=0$ this is the parameter
$\varsigma_\tau$ of {\cite[\S14.3]{P46}}; the number of signs $-$ is $Q_j$; and the passage from
$j$ to $j+1$ adds a string of $t$ numbers centred at $-\kappa_{j+1}(\tau)$, shifts by a half, and
changes the sign of the numbers below the centre. Two numbers of $\Lambda^{(j)}_\tau$ with
different signs differ by at least $\tfrac12(\Delta-t)-1\ge\tfrac12r^8$, and any two numbers
differ by a non-zero integer. So $\Lambda^{(j)}_\tau$ is the Harish-Chandra parameter of a discrete
series representation $\pi_{j,\tau}$ of $G_{j,\tau}$ (the signs say which coordinates belong to the
positive part), and for every non-compact root $\beta$
\begin{equation}\label{pz:eq-gap}
|\langle\Lambda^{(j)}_\tau,\beta\rangle|\ \ge\ \tfrac12r^8\ >\ x^{25}+N_j .
\end{equation}
For $j=m$ the group is compact, and $\Lambda^{(m)}_\tau=\lambda^\circ_\tau+\rho$ with $\rho=
(\tfrac{N_m-1}2,\dots,\tfrac{1-N_m}2)$: by \S\ref{ps:ss-restrict} the first $k-1$ rows of
$\lambda^\circ_\tau$ are equal to $-m(t+1)/2$, the $k$-th is $-m(t+1)/2+k/2-\kappa(\tau)-\Delta$,
the next $t$ rows are lower by $\tfrac12(\Delta-t-1)$, and so on; adding $\rho$ gives
\eqref{pz:eq-param}. Thus $\pi_{m,\tau}$ is the representation of highest weight $\lambda^\circ_\tau$.

Let $\pi_j=\bigotimes_\tau\pi_{j,\tau}$, $v_j$ a unit vector which, at every $\tau$, is a highest
weight vector of the minimal $K$-type of $\pi_{j,\tau}$, $d_j$ the formal degree of $\pi_j$, and
\[
c_j(x)=\langle v_j,\pi_j(x)v_j\rangle,\qquad e_j(x)=d_j\,\overline{c_j(x)}\qquad(x\in G_{j,\infty})
\]
(inner products conjugate-linear in the first variable, as in \S\ref{ps:ss-weil}; with this reading $\sigma(e_j)$ is the projector onto the type $\pi_j$; see \eqref{par:eq-coef}).
By the orthogonality relations for square-integrable representations, $e_j(x^{-1})=
\overline{e_j(x)}$, $e_j*e_j=e_j$, and for a unitary representation $\sigma$ of $G_{j,\infty}$ the
operator $\int e_j(x)\sigma(x)\,dx$, when $e_j$ is integrable, is the orthogonal projector onto
the vectors that span with their translates a multiple of $\pi_j$ and correspond to $v_j$.

\begin{definition}\label{pz:def-type}
Let $e=e_j$ or $\overline{e_j}$. For a function $F$ on $G_j(\auA)$ put $(R(e)F)(g)=
\int_{G_{j,\infty}}e(x)F(gx)\,dx$ when the integral converges absolutely. A \emph{form of type $e$}
on $G_j$ is a bounded continuous function $y$ on $G_j(\auA)$, left invariant under $G_j(E^+)$, right
invariant under an open subgroup of $G_j(\auA_f)$ (a \emph{level}), with $R(e)y=y$. The
\emph{forms of the archimedean type of the chain} on $G_j$ are the forms of type $e_j$; their
complex conjugates are the forms of type $\overline{e_j}$. The space $\auAz$ of
{\cite[\S14.3]{P46}} is the space of the forms of type $e_0$ and level $J_M$; this is the precise
meaning of ``the fixed vector of the discrete series with parameter $\varsigma_\tau$''. For $j=m$
the forms of type $e_m$ are the functions $g\mapsto l\bigl(g_\infty^{-1}\varphi(g_f)\bigr)$ with
$\varphi$ an algebraic modular form of weight $\lambda^\circ$ and $l$ the linear form given by
$v_m$ (Peter--Weyl); in this way the last lift has values in algebraic modular forms.
\end{definition}

\emph{The Fock space.} Let $0\le j<m$, $W=\auHom(Y_j,Y_{j+1})$, $W_\infty=W\otimes_{E^+}(E^+\otimes\R)$,
of real dimension $2D$, $D=fN_jN_{j+1}$. We realise $\psHS_\infty(W)$ as the Fock space $\pzF$ of
entire functions on $\C^D$ with $\|F\|^2=\boldsymbol\pi^{-D}\int|F(z)|^2e^{-|z|^2}dz$, for a
positive complex structure on $W_\infty$ commuting with $K_{j,\infty}\times K_{j+1,\infty}$, with
$\psrho_\infty$ given by \eqref{pk:eq-fock}; $|\cdot|$ is then a euclidean norm on $W_\infty$, and
$(g,h)\in G_{j,\infty}\times G_{j+1,\infty}$ acts on $W_\infty$ with operator norm at most
$\|g\|\,\|h\|$. Let $\omega_\infty$ be the archimedean part of $\omega_{\mu_{j+1}}$ on $\pzF$. For
$w\in W_\infty$ let $\partial_w$ be the derivative at $s=0$ of $\psrho_\infty(sw)$, an operator on the
space $\pzF^{\mathrm{sm}}$ of smooth vectors. Fix a basis $(\mathsf e_k)_{k\le D}$ of $\C^D$,
orthonormal for the hermitian form, so that $(\mathsf e_k,i\mathsf e_k)_k$ is an orthonormal real
basis $\mathsf B$ of $W_\infty$, and for an integer $s\ge0$ put
\[
\|\Psi\|_s=\max\bigl\{\|\partial_{w_1}\cdots\partial_{w_n}\Psi\|:\ 0\le n\le2s,\ w_1,\dots,w_n\in
\mathsf B\bigr\},
\]
and let $\pzF_s$ be the completion of $\pzF^{\mathrm{sm}}$ for $\|\cdot\|_s$, a Banach space
contained in $\pzF$.

\begin{lemma}[Growth]\label{pz:growth}
Let $s=D+1$.
\begin{enumerate}
\item[(a)] For $g\in G_{j,\infty}$, $h\in G_{j+1,\infty}$ and $\Psi\in\pzF^{\mathrm{sm}}$,
$\|\omega_\infty(g,h)\Psi\|_s\le(2D\|g\|\|h\|)^{2s}\|\Psi\|_s$. Hence $\omega_\infty$ extends to a
strongly continuous representation on $\pzF_s$ with the same bound.
\item[(b)] $(1+|z|)^{2s}\,|\Psi(z)|\,e^{-|z|^2/2}\le(D+1)^{2s}\|\Psi\|_s$ for $\Psi\in\pzF_s$ and
$z\in\C^D$.
\item[(c)] For every $f\in\psHS_f(W)$ there is $C_f$ with $|\theta(\Psi\otimes f)|\le C_f\|\Psi\|_s$
for $\Psi\in\pzF^{\mathrm{sm}}$; so $\theta(\,\cdot\otimes f)$ extends continuously to $\pzF_s$.
\end{enumerate}
\end{lemma}
\begin{proof}
(a) By (W1), $\omega_\infty(g,h)^{-1}\partial_w\,\omega_\infty(g,h)=\partial_{u^{-1}w}$, where $u$ is the
action of $(g,h)$ on $W_\infty$. The coordinates of $u^{-1}w$, $w\in\mathsf B$, in the basis
$\mathsf B$ have absolute value at most $\|u^{-1}\|\le\|g\|\|h\|$. So $\partial_{w_1}\cdots\partial_{w_n}
\omega_\infty(g,h)\Psi$ is $\omega_\infty(g,h)$ applied to a sum of $(2D)^n$ terms $c\,\partial_{w_1'}
\cdots\partial_{w'_n}\Psi$ with $|c|\le(\|g\|\|h\|)^n$, and $\omega_\infty(g,h)$ is unitary on $\pzF$.
The map $(g,h)\mapsto\omega_\infty(g,h)\Psi$ is continuous into $\pzF^{\mathrm{sm}}$ for $\Psi$
smooth (the smooth vectors of the Heisenberg group form the Schwartz space, on which the Weil
representation acts continuously~\cite[no.~35]{WeilActa}), and the bound extends this to $\pzF_s$.
(b) By \eqref{pk:eq-fock}, $(\partial_w\Psi)(z)=(z\cdot w)\Psi(z)-(\mathsf D_{\bar w}\Psi)(z)$, where
$\mathsf D_u$ is the derivative in the direction $u$; as $\Psi$ is holomorphic, $\mathsf D_{iu}=
i\mathsf D_u$, so $(z\cdot w)\Psi=\tfrac12(\partial_w-i\partial_{iw})\Psi$. Hence
$\|z^\beta\Psi\|\le\|\Psi\|_s$ for every monomial $z^\beta$ of degree at most $2s$ in the
coordinates $z_k=z\cdot\mathsf e_k$. The reproducing property gives $|F(z)|\le\|F\|e^{|z|^2/2}$ for
$F\in\pzF$; apply it to $F=z^\beta\Psi$ and expand $(1+\sum_k|z_k|)^{2s}$ into $(D+1)^{2s}$ monomials.
(c) Fix $w_0\in W_\infty$ and $f'\in\psHS_f(W)$ with $\theta(\psrho_\infty(w_0)1\otimes f')\ne0$. They
exist: the vectors $\psrho_\infty(w)1$, $w\in W_\infty$, span a dense subspace of $\pzF^{\mathrm{sm}}$,
and $\theta\ne0$ is continuous in the first factor (W0). There is
an integer $\mathsf N\ge1$ such that $\psrho_f(\mathsf NL)$ acts on $f$ and on $f'$ by characters
((W4), for a self-dual lattice $L$). If $\langle\psrho_f(\gamma)f,f'\rangle\ne0$, comparing the
actions of $\psrho_f(l)$, $l\in\mathsf NL$, on both arguments shows that $l\mapsto[l,\gamma]$ is a
fixed character of $\mathsf NL$; so $\gamma$ lies in one coset $x_0+\mathsf N^{-1}L$ ($x_0\in W(\auA_f)$ depending only on $f,f'$). For $\gamma\in
W(E)$ with coordinate $\gamma^\natural\in W_\infty$, $\langle\psrho_\infty(\gamma^\natural)\Psi,\psrho_\infty(w_0)1\rangle$ is a number of absolute value $1$
times $(\psrho_\infty(v)\Psi)(0)=e^{-|v|^2/2}\Psi(-\aubar v)$, $v=\gamma^\natural-w_0$,
which by (b) is at most $(D+1)^{2s}(1+|v|)^{-2s}\|\Psi\|_s$ in absolute value. The points
$\gamma^\natural$, $\gamma\in W(E)\cap(x_0+\mathsf N^{-1}L)$, form a coset of a lattice in
$W_\infty\cong\R^{2D}$, and $2s>2D$; so
\[
T(\Psi\otimes f)=\sum_{\gamma\in W(E)}\bigl\langle\psrho_\infty(w_0)1\otimes f',\psrho(\gamma)(\Psi\otimes f)\bigr\rangle
\]
converges absolutely (each term is the complex conjugate of the term estimated above), with
$|T(\Psi\otimes f)|\le C\|f\|\|f'\|\,\|\Psi\|_s$. Since the hermitian forms are conjugate-linear in
the first variable, $T$ is linear in $\Psi\otimes f$. The functional $T$ is
invariant under $\psrho(W(E))$, hence a multiple of $\theta$ by (W1); the multiple is
$\overline{\theta(\psrho_\infty(w_0)1\otimes f')}\ne0$, as in the proof of Proposition~\ref{pk:conj}(ii)
(for the normalisation of $\theta$ made there).
\end{proof}

\emph{Other spaces.} The proof of Lemma~\ref{pz:growth} uses only the Fock model of
$\psHS_\infty(W)$ for a positive complex structure, (W0), (W1), (W4), the unitarity of $\omega_\infty$
and the continuity of the Weil representation on smooth vectors; nothing specific to the chain
enters, and (a) holds with every integer $s'\ge0$ in place of $s$. So, for an arbitrary hermitian
space $W$ over $E$ with $W_\infty$ of real dimension $2D$ (and a positive complex structure on
$W_\infty$, $\mathsf B$ and $\|\cdot\|_{s'}$ defined as above), (b) and (c) hold, and (a) holds with
$s'$ for $W=\auHom(X,X')$ and $(g,h)\in\auU(X)(E^+\otimes\R)\times\auU(X')(E^+\otimes\R)$ (operator norms
in coordinates in which the forms are $\mathrm{diag}(\pm1)$).

\begin{lemma}[Absolute convergence along an isotropic subspace]\label{pz:absgrowth}
Let $W$ be a hermitian space over $E$, $2D=\dim_\R W_\infty$, $s=D+1$, with $\|\cdot\|_{s'}$ as in the
preceding paragraph. Let $I\subseteq W$ be an isotropic $E$-subspace and, for a datum $F$ and $w\in
W(E)$, $c_w(F)=\theta_{W_I}(\psZ_I\psrho(w)F)$ \textup{(}notation of \textup{(W2))}; $c_w$ depends only on
the class of $w$ in $W(E)/I^\perp(E)$.
\begin{enumerate}
\item[\textup{(a)}] For every $f\in\psHS_f(W)$ there is $C_f$ with $\sum_{w\in W(E)/I^\perp(E)}|c_w(\Psi\otimes
f)|\le C_f\|\Psi\|_{2s}$ for every smooth $\Psi$.
\item[\textup{(b)}] Let $W=\auHom(X,X')$, $f\in\psHS_f(W)$ and $\mathcal C$ a compact subset of
$\auU(X)(\auA_f)\times\auU(X')(\auA_f)$. There is $C$ such that
\[
\sum_{w\in W(E)/I^\perp(E)}\bigl|c_w\bigl(\omega(g,h)(\Psi\otimes f)\bigr)\bigr|\le C\,(2D\|g_\infty\|\|h_\infty\|)^{4s}
\|\Psi\|_{2s}
\]
for all smooth $\Psi$ and all $(g,h)\in\auU(X)(\auA)\times\auU(X')(\auA)$ with $(g_f,h_f)\in\mathcal C$;
the same holds with $\omega$ replaced by $\omega_{\alpha,\beta}$ or by the twisted $\omega_\mu$, which differ
from it by characters of absolute value $1$.
\item[\textup{(c)}] \textup{(Uniformity under archimedean translates.)} For every $f\in\psHS_f(W)$ there is
$C_f$ with
\[
\sum_{w\in W(E)/I^\perp(E)}|c_w(\psrho_\infty(z)\Psi\otimes f)|\le C_f\|\Psi\|_{2s}
\]
for every
$z\in W_\infty$ and every smooth $\Psi$.
\end{enumerate}
\end{lemma}
\begin{proof}
\emph{Step 1: translates.} For every $f$ there is $C_f$ with $|\theta(\psrho_\infty(z)G\otimes f)|\le
C_f\|G\|_s$ for all $z\in W_\infty$ and all smooth $G$. Indeed, in the proof of Lemma~\ref{pz:growth}(c),
replacing $\Psi$ by $\psrho_\infty(z)G$ replaces $v=\gamma^\natural-w_0$ by $\gamma^\natural+z-w_0$ (up to a
phase, as $\psrho_\infty(\gamma^\natural)\psrho_\infty(z)$ is a phase times $\psrho_\infty(\gamma^\natural+z)$) and
leaves the finite factors $\langle\psrho_f(\gamma)f,f'\rangle$ unchanged; so $|T(\psrho_\infty(z)G\otimes
f)|\le(D+1)^{2s}\|f\|\|f'\|\,\|G\|_s\sup_{u\in\R^{2D}}\sum_\lambda(1+|\lambda+u|)^{-2s}$, the sum over a coset
of a fixed lattice of $W_\infty\cong\R^{2D}$; the supremum is finite because $2s>2D$.

\emph{Step 2: the terms are Fourier coefficients.} Let $F=\Psi\otimes f$ and
$\phi(i)=\theta(\psrho(i)F)$ for $i\in I(\auA)$. On $I(\auA)\times I(\auA)$ the cocycle $\beta$ is $1$, so
$\psrho$ is additive on $I(\auA)$; hence $\phi(\gamma+i)=\phi(i)$ for $\gamma\in I(E)$ (W1), and, $f$ being
smooth, there is an open compact subgroup $K_f\subseteq I(\auA_f)$ with $\psrho_f(\kappa)f=f$, so that
$\phi(i+\kappa)=\phi(i)$ for $\kappa\in K_f$. For $w\in W(E)$, $\psrho(i)\psrho(w)=[i,w]\psrho(w)\psrho(i)$
and $\theta\circ\psrho(w)=\theta$, so by (W2)
\[
c_w(F)=\int_{[I]}\theta(\psrho(i)\psrho(w)F)\,di=\int_{[I]}[i,w]\,\phi(i)\,di .
\]
The character $i\mapsto[i,w]$ of $I(\auA)$ is trivial on $I(E)$, depends only on $w$ modulo
$I^\perp(\auA)$, and determines the class of $w$ modulo $I^\perp(E)$: if $[i,w]=1$ for all $i\in I(\auA)$
then, replacing $i$ by $ai$ ($a\in\auA_E$) and using the non-degeneracy of $\psi$, $\langle w,i\rangle=0$
for all $i$, i.e.\ $w\in I^\perp(E)$. If this character is non-trivial on $K_f$ then $c_w(F)=0$. Since $E$
is dense in $\auA_{E,f}$, $I(\auA_f)=I(E)+K_f$, so $i_\infty\mapsto(i_\infty,0)$ induces an isomorphism of
$I_\infty/\Lambda$ onto $[I]/K_f$, $\Lambda=\{\gamma\in I(E):\gamma\in K_f\}$ (a lattice in $I_\infty$),
compatible with the normalised Haar measures. Write $[i_\infty,w]=e^{2\boldsymbol\pi i\lambda_w(i_\infty)}$
with $\lambda_w$ a real linear form on $I_\infty$. Then, for the $w$ with $c_w(F)\ne0$, $\lambda_w$ lies in the
dual lattice $\Lambda^*$, distinct classes $w$ give distinct $\lambda_w$, and $c_w(F)=\hat\phi(\lambda_w)$, where
$\hat\phi(\lambda)=\int_{I_\infty/\Lambda}e^{2\boldsymbol\pi i\lambda(i)}\phi(i)\,di$. So
$\sum_w|c_w(F)|\le\sum_{\lambda\in\Lambda^*}|\hat\phi(\lambda)|$.

\emph{Step 3: decay of the coefficients.} For $i\in I_\infty$, $\phi(i)=\theta(\psrho_\infty(i)\Psi\otimes
f)$. The orbit map of the smooth vector $\Psi$ is smooth for the topology of smooth vectors, in which
$\|\cdot\|_s$ is continuous, $\psrho_\infty$ is additive on $I_\infty$, and $\theta(\,\cdot\otimes f)$ is
continuous for $\|\cdot\|_s$ (Lemma~\ref{pz:growth}(c)); so $\phi$ is smooth on $I_\infty$ and
$\mathsf D_{u_1}\cdots\mathsf D_{u_n}\phi(i)=\theta(\psrho_\infty(i)\partial_{u_1}\cdots\partial_{u_n}\Psi\otimes f)$
for $u_k\in I_\infty$. Let $\varkappa=\dim_\R I_\infty$; as $I$ is isotropic, $\varkappa\le D<2s$. Let
$u_1,\dots,u_\varkappa$ be an orthonormal basis of $I_\infty$ and $\Delta=\sum_k\mathsf D_{u_k}^2$. Each $u_k$ is a
combination of the elements of $\mathsf B$ with coefficients of absolute value $\le1$, and $w\mapsto
\partial_w$ is linear, so a product of $2s$ operators $\partial_{u_k}$ is a sum of at most $(2D)^{2s}$
products of $2s$ operators $\partial_b$, $b\in\mathsf B$, with coefficients of absolute value $\le1$; by
Step~1,
\[
\sup_{i}|\Delta^s\phi(i)|\le\varkappa^s(2D)^{2s}C_f\max_{b_1,\dots,b_{2s}\in\mathsf B}\|\partial_{b_1}\cdots
\partial_{b_{2s}}\Psi\|_s\le\varkappa^s(2D)^{2s}C_f\|\Psi\|_{2s}.
\]
For $0\ne\lambda\in\Lambda^*$, integration by parts on the torus gives $\hat\phi(\lambda)=(-4\boldsymbol\pi^2
|\lambda|^2)^{-s}\int e^{2\boldsymbol\pi i\lambda(i)}\Delta^s\phi(i)\,di$, and $|\hat\phi(0)|\le\sup|\phi|\le C_f
\|\Psi\|_s$ by Step~1. Since $2s>\varkappa$, $\sum_{0\ne\lambda\in\Lambda^*}|\lambda|^{-2s}<\infty$, and (a)
follows with $C_f(1+\varkappa^s(2D)^{2s}(2\boldsymbol\pi)^{-2s}\sum_{\lambda\ne0}|\lambda|^{-2s})$ (note
$\|\Psi\|_s\le\|\Psi\|_{2s}$).

(b) The vector $f$ is fixed, up to scalars of absolute value $1$, by an open subgroup $K'$ of
$\auU(X)(\auA_f)\times\auU(X')(\auA_f)$; $\mathcal C$ is covered by finitely many cosets $c_jK'$, so
$\omega_f(g_f,h_f)f$ lies, up to such scalars, in the finite set $\{\omega_f(c_j)f\}$. Now
$\omega(g,h)(\Psi\otimes f)=\omega_\infty(g_\infty,h_\infty)\Psi\otimes\omega_f(g_f,h_f)f$, the vector
$\omega_\infty(g_\infty,h_\infty)\Psi$ is smooth, and $\|\omega_\infty(g_\infty,h_\infty)\Psi\|_{2s}\le
(2D\|g_\infty\|\|h_\infty\|)^{4s}\|\Psi\|_{2s}$ by Lemma~\ref{pz:growth}(a) with $s'=2s$. Apply (a) to the
finitely many $\omega_f(c_j)f$.

(c) Fix $z\in W_\infty$ and put $F_z=\psrho_\infty(z)\Psi\otimes f$, $\phi_z(i)=\theta(\psrho(i)F_z)$. Steps~1
and~2 apply to $F_z$ (they use only that $F_z$ is a datum on which $\theta$ is continuous, and
$\psrho_\infty(z)$ preserves $\pzF_s$ by {\cite[Lemma~3.30(a)]{P53}}): with the same $K_f$, $\Lambda$ and
$\lambda_w$ (which depend on $f$ only), $\sum_w|c_w(F_z)|\le\sum_{\lambda\in\Lambda^*}|\hat\phi_z(\lambda)|$.
Step~3 applied to $F_z$ would bound the coefficients by $\|\psrho_\infty(z)\Psi\|_{2s}$, which grows with
$z$ (the derivatives of $\phi_z$ pick up factors $\kappa(u,z)$); we remove this dependence by a unitary
character. For $i\in I_\infty$,
$\psrho_\infty(i)\psrho_\infty(z)=e^{\kappa(i,z)}\psrho_\infty(z)\psrho_\infty(i)$ with $\kappa(i,z)\in i\R$
real bilinear ({\cite[Lemma~3.30]{P53}}); write $e^{\kappa(i,z)}=e^{2\boldsymbol\pi i\mu_z(i)}$, $\mu_z$ a real
linear form on $I_\infty$. Then
\[
\phi_z(i)=e^{2\boldsymbol\pi i\mu_z(i)}\varphi_z(i),\qquad\varphi_z(i)=\theta\bigl(\psrho_\infty(z)\psrho_\infty(i)
\Psi\otimes f\bigr).
\]
The orbit map $i\mapsto\psrho_\infty(i)\Psi$ is smooth into the smooth vectors, hence into $\pzF_s$, with
derivatives $\psrho_\infty(i)\partial_{u_1}\cdots\partial_{u_n}\Psi$; $\psrho_\infty(z)$ is bounded on
$\pzF_s$ and $\theta(\,\cdot\otimes f)$ is continuous on $\pzF_s$. So $\varphi_z$ is smooth and
$\mathsf D_{u_1}\cdots\mathsf D_{u_n}\varphi_z(i)=\theta(\psrho_\infty(z)\psrho_\infty(i)\partial_{u_1}\cdots
\partial_{u_n}\Psi\otimes f)$; as $\psrho_\infty(z)\psrho_\infty(i)$ is a phase times
$\psrho_\infty(z+i)$, Step~1 gives, exactly as in Step~3 and now uniformly in $z$,
\[
\sup_i|\varphi_z(i)|\le C_f\|\Psi\|_s,\qquad\sup_i|\Delta^s\varphi_z(i)|\le\varkappa^s(2D)^{2s}C_f\|\Psi\|_{2s}.
\]
Since $\phi_z$ is $\Lambda$-periodic, $\varphi_z(i+\gamma)=e^{-2\boldsymbol\pi i\mu_z(\gamma)}\varphi_z(i)$ for
$\gamma\in\Lambda$, and the same holds for every derivative of $\varphi_z$. For $\lambda\in\Lambda^*$ put
$\nu=\lambda+\mu_z$, so that $\hat\phi_z(\lambda)=\int_{I_\infty/\Lambda}e^{2\boldsymbol\pi i\nu(i)}\varphi_z(i)
\,di$, the integrand being $\Lambda$-periodic. If $\nu\ne0$, $e^{2\boldsymbol\pi i\nu(i)}=(-4\boldsymbol\pi^2
|\nu|^2)^{-1}\Delta e^{2\boldsymbol\pi i\nu(i)}$, and Green's formula on a fundamental parallelotope has no
boundary terms, because every product of a derivative of $e^{2\boldsymbol\pi i\nu(\cdot)}$ with a derivative
of $\varphi_z$ is $\Lambda$-periodic (the multipliers $e^{\pm2\boldsymbol\pi i\mu_z(\gamma)}$ cancel, and
$\lambda(\gamma)\in\Z$); iterating, $\hat\phi_z(\lambda)=(-4\boldsymbol\pi^2|\nu|^2)^{-s}\int e^{2\boldsymbol\pi
i\nu(i)}\Delta^s\varphi_z(i)\,di$. Hence for all $\lambda$ (covolume $1$ for the normalised measure)
\[
|\hat\phi_z(\lambda)|\le\min\bigl(C_f\|\Psi\|_s,\ (2\boldsymbol\pi|\lambda+\mu_z|)^{-2s}\varkappa^s(2D)^{2s}
C_f\|\Psi\|_{2s}\bigr)\le C'_f\,(1+|\lambda+\mu_z|)^{-2s}\|\Psi\|_{2s},
\]
with $C'_f=2^{2s}C_f\max(1,\varkappa^s(2D)^{2s}(2\boldsymbol\pi)^{-2s})$ (use $\|\Psi\|_s\le\|\Psi\|_{2s}$ and
$(1+r)^{2s}\le2^{2s}\max(1,r^{2s})$). Finally $\sup_{\mu}\sum_{\lambda\in\Lambda^*}(1+|\lambda+\mu|)^{-2s}<\infty$,
the supremum over real linear forms $\mu$, since $2s>\varkappa$ (compare the sum with an integral over
$\R^\varkappa$, as in Step~1). This gives (c), with a constant depending on $f$ only.

All three bounds extend to $\Psi\in\pzF_{2s}$: by the formula of Step~2 and Step~1, each $c_w(\,\cdot\otimes
f)$ is continuous for $\|\cdot\|_s$, and so is $\Psi\mapsto c_w(\psrho_\infty(z)\Psi\otimes f)$ for fixed
$z$ ({\cite[Lemma~3.30(a)]{P53}}), so the bounds for the finite partial sums extend by density, and
$\omega_\infty$ extends to $\pzF_{2s}$ with the same bound.
\end{proof}

\emph{The three archimedean statements.} Put $B=x^{25}$; then $2D+2\le B$, since $f\le x^{12}$ and
$N_j\le n_m\le40x^6$. Let $0\le j<m$. The \emph{joint harmonic} $P_j$ of \S\ref{ps:ss-Y} is a
non-zero polynomial in $\pzF$, harmonic for $G_{j,\infty}$ and for $G_{j+1,\infty}$, which is a
highest weight vector for $K_{j,\infty}$ and for $K_{j+1,\infty}$ with the weights of the minimal
$K$-types of the two types (for the action $\omega_\infty$, which contains the determinant twists
of \S\ref{ps:ss-weil}); we write $\Phi_j$ for this vector of $\pzF$.
\begin{enumerate}
\item[(A1)] \emph{(Integrability.)} For $0\le j\le m$ the function $c_j$ is smooth, left and right
$K_{j,\infty}$-finite and an eigenfunction of the centre of the enveloping algebra, and
$\int_{G_{j,\infty}}|\mathsf Xc_j(x)|\,\|x\|^{B}dx<\infty$ for every differential
operator $\mathsf X$ that is invariant under left or under right translations (including
$\mathsf X=1$).
\item[(A2)] \emph{(Cusp property.)} For $0\le j<m$, every proper parabolic subgroup of
$G_{j,\infty}$ of the form $\prod_\tau\mathsf P_\tau$ with $\mathsf P_\tau$ proper for all $\tau$,
with unipotent radical $\mathsf N$, and all $a,b\in G_{j,\infty}$: $\int_{\mathsf N}c_j(a\,n\,b)\,dn=0$,
the integral being absolutely convergent.
\item[(A3)] \emph{(Correspondence.)} For $0\le j<m$ such a polynomial $P_j$ exists, it is unique up
to a scalar, and in $\pzF$
\[
\int_{G_{j,\infty}}e_j(x)\,\omega_\infty(x^{-1},1)\Phi_j\,dx=\int_{G_{j+1,\infty}}e_{j+1}(x)\,
\omega_\infty(1,x)\Phi_j\,dx .
\]
\end{enumerate}
In words, (A3) says that the component of $\Phi_j$ of the type of $Y_j$ (in the first variable)
equals its component of the type of $Y_{j+1}$ (in the second). It holds if the closed subspace of
$\pzF$ on which $G_{j,\infty}$ acts by multiples of the contragredient of $\pi_j$ is, under
$G_{j+1,\infty}$, a multiple of $\pi_{j+1}$ in which the vector attached to $v_j$ and to $\Phi_j$
is the vector $v_{j+1}$: this is the correspondence of the discrete series with the parameters
\eqref{pz:eq-param} together with the correspondence of their minimal $K$-types in the space of
joint harmonics.

\smallskip\noindent (A1) is Theorem~\ref{par:A1} and (A2) is Corollary~\ref{par:A2}; both are proved from the
decay bound \eqref{par:eq-decay}, which holds by {\cite[Theorem~5.6]{P51}} and formula
{\cite[(59)]{P51}} (Theorem~\ref{pco:FJ}). (A3) is Theorem~\ref{par:A3}, proved modulo the published
results (D1)--(D3) of \S\ref{par:ss-rigid} and the normalisation (V) (Lemma~\ref{par:vac}), which is
Theorem~\ref{pco:V}, proved modulo (W1$_\infty$), (W2$_\infty$), (M), and, for $j=m-1$, (S). The
identification of the type on $Y_m$ with $W_{\lambda^\circ}$ is Remark~\ref{par:Ym}. Everything below
is proved from (A1)--(A3).\par\smallskip

\begin{proposition}[The lower chain is well defined]\label{pz:types}
Assume {\rm(A1)--(A3)}. Let $0\le j\le m$.
\begin{enumerate}
\item[(a)] Every form $y$ of type $e_j$ or $\overline{e_j}$ on $G_j$ is smooth, $K_{j,\infty}$-finite
and finite under the centre of the enveloping algebra; it is a cusp form; it is rapidly
decreasing on Siegel sets, in the form: for every Siegel set $\mathfrak S_\infty\subset G_{j,\infty}$, finite set $\mathsf c\subset G_j(\auA_f)$, open subgroup $K'$ of $G_j(\auA_f)$ and $N\ge0$ there is $C_N$ with $|y(g)|\le C_N\|g_\infty\|^{-N}$ for $g\in\mathfrak S_\infty\mathsf cK'$; and the forms of a given level form a finite-dimensional space.
\item[(b)] Let $j<m$, $f\in\psHS_f(W)$ and let $y$ be a form of type $e_j$ on $G_j$. Then the
integral $F(h)=\int_{[G_j]}\theta\bigl(\omega_{\mu_{j+1}}(g,h)(\Phi_j\otimes f)\bigr)y(g)\,dg$ converges
absolutely and $F$ is a form of type $e_{j+1}$ on $G_{j+1}$.
\item[(c)] Let $j<m$, $f\in\psHS_f(W)$ and let $z$ be a form of type $\overline{e_{j+1}}$ on
$G_{j+1}$. Then $F^*(g)=\int_{[G_{j+1}]}\theta\bigl(\omega_{\mu_{j+1}}(g,h)(\Phi_j\otimes f)\bigr)z(h)
\,dh$ converges absolutely and $F^*$ is a form of type $\overline{e_j}$ on $G_j$.
\item[(d)] The space $\auAz$ is finite-dimensional and stable under the right translations by
the stabiliser of the lattice at the chosen places, so $\auG$ acts on it; the Poincar\'e map
$\auP$ of {\cite[\S14.3]{P46}} is well defined; and $\iota_\vartheta$ maps $\auAz$ to algebraic
modular forms of weight $\lambda^\circ$ on $\auU(Y_m)$.
\end{enumerate}
In particular the assumption made in \S\ref{ps:ss-Y} holds: the lifts $\Theta^Y_{j,\vartheta}$ and
their transposes $\Theta^{Y,*}_{j,\vartheta}$ preserve the types, and all the forms concerned are
cusp forms and rapidly decreasing.
\end{proposition}
\begin{proof}
We write $G=G_j$, $e=e_j$; the case of $\overline{e_j}$ is the same.

(a) Let $y=R(e)y$ be bounded. For $g=(g_\infty,g_f)$, $y(g)=\int e(g_\infty^{-1}x)\,y(x,g_f)\,dx$.
By (A1) the derivatives of $e$ under left translations are integrable, so $y$ is smooth in
$g_\infty$ and its derivatives are bounded; the left translates of $e$ under $K_{j,\infty}$ span a finite-dimensional
space, so $y$ is right $K_{j,\infty}$-finite; and an element of the centre of the enveloping
algebra acts on $y$ through its action on $e$, by a scalar. So $y$ is an automorphic form, and it
is bounded. For $j=m$ the quotient $[G]$ is compact, there are no proper parabolic subgroups, and
the finiteness follows from the Peter--Weyl theorem. Let $j<m$,
let $\mathsf P$ be a proper parabolic subgroup of $G$ over $E^+$ with unipotent radical
$\mathsf N$, and $y_{\mathsf P}(g)=\int_{\mathsf N(E^+)\backslash\mathsf N(\auA)}y(ng)\,dn$. It is
bounded, left invariant under $\mathsf N(\auA)$, and $R(e)y_{\mathsf P}=y_{\mathsf P}$ (Fubini).
Fix $g_f$ and put $\phi(x)=y_{\mathsf P}(x,g_f)$ on $G_{j,\infty}$; it is left invariant under
$\mathsf N_\infty=\mathsf N(E^+\otimes\R)$, the unipotent radical of $\prod_\tau\mathsf P(E^+_\tau)$,
and
\[
y_{\mathsf P}(g_\infty,g_f)=\int_{G_{j,\infty}}e(g_\infty^{-1}x)\phi(x)\,dx=\int_{\mathsf N_\infty
\backslash G_{j,\infty}}\phi(x)\Bigl(\int_{\mathsf N_\infty}e(g_\infty^{-1}nx)\,dn\Bigr)dx=0
\]
by (A2), the function $e$ being integrable and $\phi$ bounded. So $y$ is a cusp form. A cuspidal
automorphic form is rapidly decreasing on Siegel sets: on $\mathfrak S_\infty=\omega A_tK$ ($\omega$, $K$ compact) one has $|y(g)|\le C_NH(a)^{-N}$ for $g=uak\,ck'$, where $H(a)=\max_\alpha a^\alpha$ over the simple roots; and $\|g_\infty\|\le C\|a\|\le C''H(a)^{c_0}$ for a fixed $c_0>0$, because $\omega,K$ are compact, $\|a\|$ is the largest value of a weight of the standard representation, a non-negative combination of the simple roots on $A_t$, and $a^\alpha\ge t$; since $\|g_\infty\|\ge1$ ($|\det g_\infty|=1$), this gives the bound $|y(g)|\le C_N\|g_\infty\|^{-N}$ stated in (a). Moreover, the cusp forms of a given level,
$K_\infty$-type and infinitesimal character form a finite-dimensional space (Gelfand and
Piatetski-Shapiro, Harish-Chandra; see \cite[\S1.8, p.~192]{BorelJacquet}, \cite[Chapter~I, Theorem~1, p.~8]{HarishChandraAF}).

(b) Let $s=D+1$, so $2s\le B$. By reduction theory (finiteness of the class number \cite[Theorem~5.1]{BorelFin}; Siegel sets as fundamental sets for the arithmetic subgroups, \cite[Theorem~6.5, and Theorem~12.3 over a number field]{BorelHC1962}; finite Haar measure of Siegel sets, \cite[Proposition~4.3]{BorelHC1962}, stated there for an open subgroup of a real algebraic semisimple group; it applies to the derived group $G^1$ of $G_{j,\infty}$, a product of groups $\mathrm{SU}(p,q)$, and hence to $G_{j,\infty}=Z\cdot G^1$, $Z$ the compact centre: an Iwasawa decomposition $KAN$ of $G_{j,\infty}$ has $A,N\subset G^1$ and $K=Z\cdot(K\cap G^1)$, so a Siegel domain of $G_{j,\infty}$ is $Z$ times a Siegel domain of $G^1$, and its Haar measure is at most $\operatorname{vol}(Z)$ times that of the latter for compatible normalisations, the multiplication map $Z\times G^1\to G_{j,\infty}$ being a surjective homomorphism with finite kernel) there are a Siegel set
$\mathfrak S_\infty\subset G_{j,\infty}$ of finite measure, a finite set $\mathsf c\subset G(\auA_f)$
and an open subgroup $K'$ of $G(\auA_f)$ fixing $f$ under $\omega_{\mu_{j+1}}(\cdot,1)$ and fixing
$y$, such that $G(\auA)=G(E^+)\,\mathfrak S_\infty\,\mathsf c\,K'$. For $g=g_\infty ck'$ in this set
and $\Psi\in\pzF_s$, Lemma~\ref{pz:growth} gives
\[
\bigl|\theta\bigl(\omega(g,h)(\Psi\otimes f)\bigr)\bigr|\le C(h_f)\,(2D\|g_\infty\|\|h_\infty\|)^{2s}
\|\Psi\|_s,\qquad C(h_f)=\max_{c\in\mathsf c}C_{\omega_f(c,h_f)f},
\]
and by (a), with $N=2s$, $|y(g)|\le C'\|g_\infty\|^{-2s}$ on this set. The integrand is left invariant under
$G(E^+)$, by (W1). Hence
\[
L_h(\Psi)=\int_{[G]}\theta\bigl(\omega(g,h)(\Psi\otimes f)\bigr)y(g)\,dg
\]
converges absolutely, and $|L_h(\Psi)|\le C''(h_f)\|h_\infty\|^{2s}\|\Psi\|_s$. The function
$F=L_\bullet(\Phi_j)$ is continuous, left invariant under $G_{j+1}(E^+)$ by (W1), and right
invariant under an open subgroup fixing $f$. For $x\in G_{j+1,\infty}$, $F(hx)=L_h(\omega_\infty(1,x)
\Phi_j)$, so $|F(hx)|\le C(h)\|x\|^{2s}$, and by (A1) the following integrals converge absolutely,
the vector-valued ones in $\pzF_s$:
\[
\mathsf P'\Psi=\int_{G_{j+1,\infty}}e_{j+1}(x)\,\omega_\infty(1,x)\Psi\,dx,\qquad\mathsf P\Psi=
\int_{G_{j,\infty}}e_j(x)\,\omega_\infty(x^{-1},1)\Psi\,dx,\qquad R(e_{j+1})F=L_\bullet(\mathsf P'\Phi_j).
\]
Since $y=R(e_j)y$, Fubini and the change of variable
$g\mapsto gx^{-1}$ give, for every $\Psi\in\pzF_s$,
\[
L_h(\Psi)=\int_{G_{j,\infty}}e_j(x)\int_{[G]}\theta\bigl(\omega(gx^{-1},h)(\Psi\otimes f)\bigr)y(g)\,dg
\,dx=L_h(\mathsf P\Psi).
\]
By (A3), $\mathsf P\Phi_j=\mathsf P'\Phi_j$. Therefore
\[
R(e_{j+1})F=L_\bullet(\mathsf P'\Phi_j)=L_\bullet(\mathsf P\Phi_j)=L_\bullet(\Phi_j)=F .
\]
It remains to see that $F$ is bounded. Let $\mathcal V$ be the space of the forms of type
$e_{j+1}$ on $G_{j+1}$ of a level $K''$ fixing $F$; it is finite-dimensional by (a). For $T>0$ let
$F_T=F$ where $|F|\le T$ and $F_T=0$ elsewhere. It is bounded, measurable, left invariant under
$G_{j+1}(E^+)$ and right invariant under $K''$; so $R(e_{j+1})F_T$ is bounded, continuous
(translation is continuous in $L^1$), has the same invariance, and is fixed by $R(e_{j+1})$,
because $e_{j+1}*e_{j+1}=e_{j+1}$. Thus $R(e_{j+1})F_T\in\mathcal V$. For $T\to\infty$,
$R(e_{j+1})F_T\to R(e_{j+1})F=F$ pointwise, by dominated convergence. A finite-dimensional space of
functions is closed for pointwise convergence; so $F\in\mathcal V$.

(c) The same proof, with the two groups exchanged. With $L^*_g(\Psi)=\int_{[G_{j+1}]}\theta(\omega
(g,h)(\Psi\otimes f))z(h)\,dh$ one finds, using $\overline{e}(x^{-1})=e(x)$: $L^*_g(\Psi)=L^*_g(
\mathsf P'\Psi)$ from $z=R(\overline{e_{j+1}})z$, and $R(\overline{e_j})F^*=L^*_\bullet(\mathsf P
\Phi_j)$. So $R(\overline{e_j})F^*=L^*_\bullet(\mathsf P'\Phi_j)=F^*$ by (A3), and $F^*$ is bounded by
the truncation argument.

(d) The first assertion is (a); right translation by an element of $G_0(\auA_f)$ commutes with
$R(e_0)$. The group $R_0=U\rtimes\auU(X_0)$ has compact quotient $[R_0]$ ($\auU(X_0)$ is
anisotropic), the function $\mathcal W_\varphi$ is bounded on it and the elements of $\auAz$ are
bounded and continuous; so $y\mapsto\int_{[R_0]}\mathcal W_\varphi(r)\overline{y(rg_0)}\,dr$ is an
antilinear form on the finite-dimensional inner product space $\auAz$, which defines $\auP\varphi
\in\auAz$. The last assertion is (b) applied $m$ times with $f=\Upsilon_{j,f}$
(Lemma~\ref{ps:level} gives the level) and Definition~\ref{pz:def-type} for $j=m$.
\end{proof}

\begin{lemma}[Level and character on $Y_m$]\label{pz:lastlevel}
Assume {\rm(A1)--(A3)}. Let $\vartheta$ be a twist, $z\in\auAz$, and let $k\in K_m$ have, at every
level place $w$, a reduction $\bar k_w$ that is block upper triangular for the flag
$\psM_0\subset\psM_1\subset\dots\subset\psM_m$ of the frames, with diagonal blocks
$\bar k_w^{(0)},\dots,\bar k_w^{(m)}$, and $\bar k_w^{(0)}=1$ at the chosen places. Then
$(\iota_\vartheta z)(gk)=\prod_w\prod_{i=1}^m\psi_{i,w}(\det\bar k_w^{(i)})\cdot(\iota_\vartheta z)(g)$,
where $\psi_{i,w}$ is the tame part of $\mu_i$ at $w$. So $\iota_\vartheta z$ has the level
$J_\vartheta$ and the character $\chi_\vartheta$ of {\cite[\S14.3]{P46}}.
\end{lemma}
\begin{proof}
Choose $k_i\in K_i$ $(0\le i\le m)$, $k_m=k$, whose reduction at each level place is the upper
left block of $\bar k_w$ of the size of $Y_i$ (the reduction maps are onto, the level places being
split). Then $\bar k_{i+1}\iota_i=\iota_i\bar k_i$, and as in the proof of
{\cite[Proposition~14.27]{P46}}, by (W1) and Lemma~\ref{ps:level},
$(\Theta^Y_{i,\vartheta}F)(hk_{i+1})=\varepsilon_i\,\Theta^Y_{i,\vartheta}(F(\,\cdot\,k_i))(h)$ with
$\varepsilon_i=\prod_w\psi_{i+1,w}(\det\bar k_{i+1,w})\psi_{i+1,w}(\det\bar k_{i,w})^{-1}=\prod_w
\psi_{i+1,w}(\det\bar k_w^{(i+1)})$. By induction $(\iota_\vartheta z)(gk)=\prod_i\varepsilon_i\cdot
\iota_\vartheta(z(\,\cdot\,k_0))(g)$, and $z(\,\cdot\,k_0)=z$ because $k_0\in J_M$.
\end{proof}

\begin{lemma}[The characters of the lifts exist]\label{pz:mu}
Let $S$ be a finite set of primes of $E$, no two above the same place of $E^+$, each above a place
split in $E$; let $(c_\tau)_{\tau\in\auPl}$ be integers and $\psi_w$ ($w\in S$) characters of
$\auF_w^\times$ with $\prod_\tau\tau(\xi)^{c_\tau}\prod_{w\in S}\psi_w(\xi\bmod w)=1$ for every root of
unity $\xi\in E$. There is a unitary character $\mu$ of $\auU(1)(E^+)\backslash\auU(1)(\auA_{E^+})$ with
$\mu_\tau(u)=\tau(u)^{c_\tau}$ at the infinite places, with tame part $\psi_w$ at $w\in S$, and
unramified at all other finite places. If the reduction of the roots of unity of $E$ modulo some
$w_0\in S$ is injective (for instance a chosen place $w_0$), then for all $(c_\tau)$ and all $\psi_w$,
$w\ne w_0$, there is $\psi_{w_0}$ satisfying the condition. In particular the characters $\chi_M$,
$\mu^\flat$ and $\mu_j^0$ of {\cite[\S14.3]{P46}} exist ({\cite[Definition~6.14]{P54}},
{\cite[Proposition~6.15]{P54}}).
\end{lemma}
\begin{proof}
Let $\mathsf U=\auU(1)(E^+\otimes\R)\times\prod_v\auU(1)(\auO_v)$, an open subgroup of
$\auU(1)(\auA_{E^+})$. Its intersection with $\auU(1)(E^+)$ consists of the $\xi\in E$ with
$\xi\aubar\xi=1$ integral at every finite place; all conjugates of such $\xi$ have absolute value
$1$, so (Kronecker) it is the group $\mu_E$ of roots of unity of $E$. At $w\in S$, above the split
place $v$, $\auU(1)(\auO_v)=\auO_w^\times\to\auF_w^\times$. The formula
$\mu_{\mathsf U}(u)=\prod_\tau\tau(u_\tau)^{c_\tau}\prod_{w\in S}\psi_w(u_w\bmod w)$ defines a continuous
unitary character of $\mathsf U$, and the hypothesis says that it is trivial on $\mu_E$ (a rational
point $\xi$ has components $\tau(\xi)$ at the infinite places and $\xi$ at the finite ones). Hence
$\gamma u\mapsto\mu_{\mathsf U}(u)$ is a well-defined character of the open subgroup
$\auU(1)(E^+)\mathsf U$, trivial on $\auU(1)(E^+)$. A unitary character of an open (hence closed)
subgroup of a locally compact abelian group extends to the whole group (Pontryagin duality: the
restriction map of dual groups is onto); the extension is trivial on $\auU(1)(E^+)$ and is the
required $\mu$. For the last assertion, $\xi\mapsto\prod_\tau\tau(\xi)^{-c_\tau}\prod_{w\ne
w_0}\psi_w(\xi\bmod w)^{-1}$ is a character of the subgroup $\mu_E$ of the cyclic group
$\auF_{w_0}^\times$, and extends to it; at a chosen place $w_0$ the reduction of $\mu_E$ is injective
because $p_0,p_1>b^{10}$ do not divide $\#\mu_E$. The characters of {\cite[\S14.3]{P46}}: the
condition is checked in {\cite[Definition~6.14]{P54}} and {\cite[Proposition~6.15(T1)]{P54}}.
\end{proof}

\subsection{The constant $C_0$ and the definition of admissible data}\label{pz:ss-const}

{\cite[Proposition~14.23]{P46}} was proved in \S\ref{ps:ss-kernel} with the inequality $C_0\le\log\log
b$, justified in (R5) of \S\ref{pk:ss-rest} by ``taking the constant in the definition of
admissible data at least $C_0$''. We make this precise. Let $(q_0,d)$ be parameters in the sense
of {\cite[Definition~13.1]{P46}}, $K=\Q(\sqrt{-d})$, and
\begin{equation}\label{pz:eq-C0}
C_0(q_0,d)=4+\log^+\mathsf M_d+\frac{\sqrt d}{\boldsymbol\pi\,\mathsf m}+\log\Bigl(1+\frac4\delta\Bigr),
\qquad\mathsf m=\prod_{p\mid q_0}\Bigl(1-\frac1p\Bigr)p^{1/(p-1)},\quad\delta=\Bigl(\frac{2\boldsymbol
\pi\,\mathsf m}{\sqrt d}\Bigr)^{1/2},
\end{equation}
where $\mathsf M_d=\max_{\Lambda'}\sup_{u\in\C}|\theta^{\Lambda'}_u(0)|$, the maximum over the
normalised lattices $(\Lambda',t',a_1')$ with $\auO_K\Lambda'=\Lambda'$ (\S\ref{pk:ss-curve}, with
$t'$ admissible).

\begin{lemma}\label{pz:C0}
\begin{enumerate}
\item[(a)] $C_0(q_0,d)$ is a finite real number that depends only on $q_0$ and $d$: not on the
scales $x,x_1,b$, nor on $s$, $q_1$, $E$ or any other part of admissible data.
\item[(b)] It is the constant $C_0$ of Proposition~\ref{pk:conj} and of Lemma~\ref{ps:conj}{\rm(ii)}.
\item[(c)] With the constant $C_P$ of {\cite[(35)]{P46}} in condition {\rm(D2)} of
{\cite[Definition~13.3]{P46}}, which contains $d$ and $C_0(q_0,d)$ as summands (all summands are
non-negative), all admissible data satisfy
$C_0\le10^{-12}\log\log b$ and $d\le10^{-12}\log\log b$.
\end{enumerate}
\end{lemma}
\begin{proof}
(a) The lattices $\Lambda'$ with $\auO_K\Lambda'=\Lambda'$ form $h_K$ classes up to homothety. In
each class the admissible point is determined modulo the lattice (\S\ref{pf:ss-adm}), and the
condition $\wp''(t')=2$ fixes the lattice up to a fourth root of unity, since
$\wp''_{c\Lambda}(ct)=c^{-4}\wp''_\Lambda(t)$; and $a_1'$ is determined up to sign. So there are
finitely many normalised lattices, and they depend on $d$ only. For each of them $u\mapsto|\theta_u
(0)|$ is continuous and $\Lambda'$-periodic (\S\ref{pk:ss-curve}), hence bounded. So $\mathsf M_d$
is finite and depends on $d$ only. The primes dividing $q_1$ are those dividing $q_0$
({\cite[\S13.1]{P46}}), so the number $\mathsf m$ of the proof of Proposition~\ref{pk:conj},
defined there as a product over the primes dividing $q_0$, depends on $q_0$ only, and $\delta=
c_\psi\sqrt{\mathsf m}$ with $c_\psi=\sqrt{2\boldsymbol\pi/\sqrt d}$. (b) This is the last line of
the proof of Proposition~\ref{pk:conj}. (c) The number $c_1=-1-\gamma+\Rel\frac{L'}L(1,\chi_K)$
depends on $d$ only. {\cite[Definition~13.1]{P46}} and~\eqref{pz:eq-C0} precede and do not involve
the data, so there is no circularity: first $(q_0,d)$, then $C_P$, then the conditions (D1)--(D8).
By (D2), $10^{12}C_P\le\log\log b$, and $C_0\le C_P$, $d\le C_P$.
\end{proof}

Consequently, in the proof of {\cite[Proposition~14.23]{P46}} the terms $C_0(q_1nn'+1)^2+fnn'\log(2L+2)$
are at most $x^{50}-1$: by {\cite[Lemma~13.4]{P46}}, $q_1\le x^{12}$, $f\le x^{12}$, $nn'\le
n_m^2\le1600x^{12}$, $\log\log b\le\tfrac1{100}\log x$ and $\log(2L+2)\le3\log b\le3x^{1/100}$. The
existence theorem for admissible data ({\cite[Theorem~13.9]{P46}}) uses of $C_P$ only that it
depends on $(q_0,d)$ alone, and is not affected.

\subsection{The levels of the upper chain; sizes; short vectors}\label{pz:ss-levels}

\emph{The levels.} For $0\le j\le m$ the level $J_{X_j}$ of $X_j$ is defined as $J_X$ is for
$j=0$: it is the group of the $g$ in the stabiliser $K^\circ$ of $\widehat\auO_E^{\,n_j}$ in
$\auU(X_j)(\auA_{E^+,f})$ that are $\equiv1$ modulo every chosen place and whose last row is
$\equiv(0,\dots,0,1)$ modulo $\pi_N$. This is the definition used in Step~5 of the proof of
{\cite[Proposition~14.19]{P46}} and in \S\ref{pl:ss-level}. For the trivial twist the level of
$\mathcal A_m(1)$ is the group $J_1$ of the elements of the stabiliser of $\psM_m=\widehat\auO_E^{\,
n_m-1}$ in $\auU(Y_m)(\auA_{E^+,f})$ that are $\equiv1$ modulo every chosen place.

\begin{lemma}\label{pz:sizes}
\begin{enumerate}
\item[(a)] The levels $J_{X_j}$ $(0\le j\le m)$ and $J_1$ contain the principal congruence subgroup
$K(MN)$ of the corresponding standard space, of rank $n_j\le n_m$, respectively $n_m-1$; they are
free.
\item[(b)] $n_m\le40x^6$, $6q_1d(fk)^2\le x^{100}$, and $6q_1d(fn)^2\le x^{102}$ for $n\le n_m$.
\item[(c)] $\log(MN)\le x$.
\end{enumerate}
So the hypotheses of Proposition~\ref{pf:classes} hold for these levels, and each of them has at
most $\exp(x^{30})$ classes.
\end{lemma}
\begin{proof}
(a) The chosen places divide $\pi_M$, which divides $M$, and $\pi_N$ divides $N$; so an element
$\equiv1$ modulo $MN$ satisfies all the congruences. The levels are free by {\cite[Lemma~14.5]{P46}}.
(b) By {\cite[Lemma~13.4]{P46}}, $q_1\le x^{12}$, $f\le x^{12}$, $m\le2\sqrt x$, $x^5/2\le k\le r\le
3x^5$, and $d\le\log\log b\le\log x$. So $n_m=k+m(5r+1)\le3x^5+2\sqrt x\,(15x^5+1)\le40x^6$,
$6q_1d(fk)^2\le54\,x^{46}\log x\le x^{100}$, and for $n\le n_m$, $6q_1d(fn)^2\le6\cdot1600\,x^{48}
\log x\le x^{102}$. (c) $\log M\le22\log b$ and $\log N\le25m\log b$ ({\cite[Lemma~13.4(S5)]{P46}}),
and $(22+50\sqrt x)\log b\le x$ since $(\log b)^{100}\le x$.
\end{proof}

\emph{The estimate in the proof of {\cite[Proposition~14.22]{P46}}.} The proof of
{\cite[Proposition~14.22]{P46}} uses the following bound.

\begin{lemma}\label{pz:short}
Let $Y\in\auGL_k(E)$ have entries all of whose conjugates are at most $C\ge k$ in absolute value,
$\zeta=(\zeta_\tau)_{\tau\in\auPa}$ vectors of $\C^k$ whose coordinates are roots of unity, and
$\xi=Y\mathrm{ev}(\zeta)\in W_{\lambda^X}$ as in the proof of {\cite[Proposition~14.22]{P46}}. Then
$\|\sigma\xi\|^2\le C^{\,5fa}$ for all $\sigma\in\auAut(\C)$.
\end{lemma}
\begin{proof}
$W_{\lambda^X}=\bigotimes_{\tau\in\auPa}V_\tau$, where $V_\tau$ is the dual of the space of forms of
degree $a$ in $k$ variables, with the basis $(e_\beta)$ dual to the monomials. For $w\in\C^k$ the
evaluation $\mathrm{ev}_w$ at $w$ is $\sum_\beta w^\beta e_\beta$. By \S\ref{ps:ss-restrict} the
Gram matrix of $(e_\beta)$ is the matrix of the coefficients of the reproducing kernel
$(\sum_ix_i\bar y_i)^a$ in the products of monomials $x^\beta\bar y^{\beta'}$; hence the value of the
hermitian form at a vector with coordinates $(w^\beta)_\beta$ is $(\sum_i|w_i|^2)^a$. An element
$g$ of $\auGL_k$ acts on $V_\tau$ by $\mathrm{ev}_w\mapsto\mathrm{ev}_{g'w}$, where $g'$ is $g$, its
transpose, or the complex conjugate of one of these, according to the conventions (N1), (N2) of
\S\ref{pl:ss-norm}. So the coordinates of $\xi$ are the numbers $\prod_\tau w_\tau^{\beta_\tau}$ with
$w_\tau=\tau(Y)'\zeta_\tau$, the coordinates of $\sigma\xi$ are $\prod_\tau(\sigma w_\tau)^{\beta_\tau}$,
and the entries of $\sigma w_\tau$ are sums of $k$ products of a conjugate of an entry of $Y$ and a
root of unity, of absolute value at most $kC$. Therefore
\[
\|\sigma\xi\|^2=\prod_{\tau\in\auPa}\Bigl(\sum_i|(\sigma w_\tau)_i|^2\Bigr)^a\le(k\cdot k^2C^2)^{a\#\auPa}
\le C^{\,5fa}.\qedhere
\]
\end{proof}
With $C=C_{\mathrm{sb}}=6q_1d(fk)^2\ge k$ this is the estimate needed: $4^fC_{\mathrm{sb}
}^{\,5fa}\le\exp(f\log4+500fa\log x)\le e^{\aueu}$ by Lemma~\ref{pz:sizes}(b), since
$\aueu=fa(\log b)^3$ and $\log x\le\log b$.

\subsection{The inputs of the Short Forms Theorem}\label{pz:ss-audit}

We list the statements used in the proof of {\cite[Theorem~14.13]{P46}}, from the top, with the place of
the written proof. ``Published'' means that the proof quotes results from the literature, named at
the place of use.
\begin{itemize}
\item \emph{The deduction} ({\cite[\S14.7]{P46}}) uses {\cite[Lemmas~14.9 and~14.18]{P46}},
which are proved in {\cite[\S14]{P46}}, and the inputs (I1)--(I6). The objects in the
statement are defined in {\cite[\S14.3]{P46}}, with the complements of \S\ref{ps:ss-weil},
\S\ref{ps:ss-Y}, \S\ref{ps:ss-restrict}, \S\ref{ps:ss-glue}, \S\ref{pl:ss-norm} and of
Definition~\ref{pz:def-type}, Lemma~\ref{pz:mu} and \S\ref{pz:ss-levels}; that the lower chain and
the Poincar\'e map are well defined is Proposition~\ref{pz:types}, which rests on (A1)--(A3).
\item \emph{{\rm(I1)}, {\cite[Proposition~14.19]{P46}}}: \S\ref{ps:ss-seesaw}, from
Lemma~\ref{ps:weil} (W0)--(W5) (published: Weil, Kudla, Harris--Kudla--Sweet), (W6)
(Proposition~\ref{pf:flip}), Propositions~\ref{ps:onestep} (\S\ref{pk:ss-onestep}; Witt's theorem),
\ref{pk:definite} (Cauchy decomposition, published), Lemmas~\ref{pk:chi0}, \ref{ps:level},
\ref{ps:double}, \ref{ps:support} (with (R1), (R2) and Lemma~\ref{pf:orbits}; Iwasawa decomposition,
published), \ref{ps:fin} (with (R3), (R4)), and Proposition~\ref{pz:types} for the types, the
cuspidality and the rapid decrease of the forms $z_j$.
\item \emph{{\rm(I2)}, {\cite[Proposition~14.20]{P46}}}: Theorem~\ref{pa:divis}, from
Propositions~\ref{pa:count}, \ref{pa:product}, Theorem~\ref{pa:onecurve} (published results on
elliptic curves with complex multiplication), Lemma~\ref{pa:kernel} (Propositions~\ref{pk:K0},
\ref{pk:std}, \ref{pk:K12}, Theorem~\ref{pf:std}, Corollary~\ref{pf:K1}), Lemma~\ref{pa:hecke}
(\S\ref{pl:ss-hecke}), Lemma~\ref{pa:eldiv} ({\cite[Theorem~2.2]{P50}}, {\cite[\S3]{P50}}),
Lemma~\ref{pl:level} with Lemma~\ref{pl:kernelN} (Corollary~\ref{pf:K1}(b)),
{\cite[Proposition~14.25]{P46}} (\S\ref{ps:ss-glue}) and Proposition~\ref{pk:branching}.
\item \emph{{\rm(I3)}, {\cite[Proposition~14.22]{P46}}}: {\cite[\S14.8]{P46}}, from
{\cite[Lemma~14.21]{P46}} (\S\ref{ps:ss-short}) and Lemma~\ref{pz:short}.
\item \emph{{\rm(I4)}, {\cite[Proposition~14.26]{P46}}}: {\cite[\S14.8]{P46}}, from
{\cite[Proposition~14.23]{P46}} (\S\ref{ps:ss-kernel}; Lemma~\ref{ps:conj} is Proposition~\ref{pk:conj}
with Theorem~\ref{pf:recip}; the constant is Lemma~\ref{pz:C0}; Chowla--Selberg formula,
published), {\cite[Proposition~14.24]{P46}} (\S\ref{ps:ss-restrict}, Proposition~\ref{pk:branching}),
{\cite[Proposition~14.25]{P46}}, and Proposition~\ref{pf:classes} with Lemma~\ref{pz:sizes}.
\item \emph{{\rm(I5)}, {\cite[Proposition~14.28]{P46}}}: from {\cite[Proposition~14.27]{P46}} (\S\ref{ps:ss-Y})
and Proposition~\ref{pf:classes} with Lemma~\ref{pz:sizes}.
\item \emph{{\rm(I6)}, {\cite[Lemma~14.29]{P46}}}: proved in {\cite[\S14.9]{P46}}.
\end{itemize}
\emph{What the proof of {\cite[Theorem~14.13]{P46}} rests on.} After this section the proof
of {\cite[Theorem~14.13]{P46}} is written out: the archimedean statements (A1), (A2), (A3) of
\S\ref{pz:ss-types} are proved in \S\ref{par:sec} and \S\ref{pco:sec}
(Theorems~\ref{par:A1}, \ref{par:A3}, \ref{pco:FJ}, \ref{pco:V}, Corollary~\ref{par:A2}) from the
published results named there: (D1)--(D3) on discrete series, and the local properties
(W1$_\infty$), (W2$_\infty$), (M), (S) of the archimedean Weil representation. They enter through
Proposition~\ref{pz:types} only: in the definition of
$\iota_\vartheta$ and $\auP$, and in {\cite[Proposition~14.19]{P46}}. All statements have a
written proof, complete modulo the published results named in them. The explicit form
\eqref{pl:eq-w5} of (W5), with its factor, is Proposition~\ref{pwf:W5}. The absolute convergence and the bound in (W3) are proved here
(Lemma~\ref{pz:absgrowth}).


\appendix

\section{Statements used from \texorpdfstring{\cite{P46}}{[P46]}, \texorpdfstring{\cite{P50}}{[P50]}, \texorpdfstring{\cite{P51}}{[P51]}, \texorpdfstring{\cite{P53}}{[P53]}, \texorpdfstring{\cite{P54}}{[P54]}}\label{A1:quoted}

For the convenience of the reader we reproduce here, without proofs, the statements of other papers that are used in this one, in the notation of the paper quoted. Each is headed by its place in that paper.

\begin{quoteddef}{\cite[Definition~13.1]{P46}}{}
\emph{Parameters} are a pair $(q_0,d)$, where $q_0$ is a positive multiple of
$18480=2^4\cdot3\cdot5\cdot7\cdot11$ and $d$ is a prime with
$d\equiv q_0/2-1\pmod{q_0}$.
\end{quoteddef}

\begin{quoteddef}{\cite[Definition~13.3]{P46}}{ (admissible data)}
Let $(q_0,d)$ be parameters. \emph{Admissible data} $D$ for $(q_0,d)$ consist
of the following objects, subject to the conditions \textup{(D1)--(D8)}.
\begin{enumerate}
\item[(D1)] Integers $1\le x\le x_1\le b$, the \emph{scales}.
\item[(D2)] They satisfy
\begin{gather*}
 10^{12}\,\anCP\le\log\log b,\qquad 100\log\log b\le\log x,\qquad 10^{6}\log x\le\log x_1,\\
 10^{6}\log x_1\le\log b,\qquad 6\log\log x_1\le\log\log b .
\end{gather*}
\item[(D3)] Integers $m$, $r$, $s$: $m$ is a power of $2$ with
$\sqrt x\le m\le2\sqrt x$; $r$ is a power of $3$ with $x^{5}\le r\le3x^{5}$;
$2^{s}=m^{22}$; and $s\le32\log x$. The field of the data is $E=E_s$, with
$q_1=q_1(s)$, $q_2=q_2(s)$, $f=[E:K]$, $\anPl=\anPl_s$.
\item[(D4)] $q_1x_1$ divides $b$.
\item[(D5)] A set $\anPb\subseteq\anPl$ (the \emph{$b$-places}) with
$\ancl_{q_0}(\sigma)=-1$ for all $\sigma\in\anPb$, with
$T(s)\le\#\anPb\cdot\log b$, and sparse for the primes up to $5r(r+1-m)$.
\item[(D6)] The weight $\kappa\colon\anPl\to\mathbb N$: $\kappa(\tau)=b$ for
$\tau\in\anPb$ and $\kappa(\tau)=xb/x_1$ otherwise.
\item[(D7)] Two distinct primes $p_0,p_1\equiv1\pmod{q_2^{2}}$ in the interval
$[b^{10},b^{11}]$ and an integer $\pi_M$ of $K$ with
$\pi_M\overline{\pi_M}=p_0p_1=:M$.
\item[(D8)] A prime $N\equiv1\pmod{q_2^{2}}$ with
\[
 \log N_-\le\log N\le\log N_-+\sqrt{\log b},\qquad
 \log N_-:=\frac mf\sum_{\tau\in\anPl}\log\kappa(\tau)+(m-1)\log\frac{M\,\anrd_E}{2\pi}-\frac{15}2\log b ,
\]
where $\anrd_E=|\operatorname{disc}E|^{1/[E:\mathbb Q]}$ is the root
discriminant, and an integer $\pi_N$ of $E$ with $\pi_N\pi_N^{*}=N$.
\end{enumerate}
We write $a=xb$ and $k=r+1-m$; the sets $\anPa$ and $\anPb$ are disjoint.
\end{quoteddef}

\begin{quotedthm}{\cite[Lemma~13.4]{P46}}{}
For admissible data the following hold.
\begin{enumerate}
\item[(S1)] $x^{10^{6}}\le x_1$, $x_1^{10^{6}}\le b$,
$4\log x_1\le(\log x_1)^{6}\le\log b$,
$(\log b)^{100}\le x$, $\log b\ge128$, and
$\log x\ge10^{14}(q_0+2d+|\aneu|)$.
\item[(S2)] $x^{11}\le2^{s}\le2^{22}x^{11}$; $l_0(s)\le2592\log x$;
$2m\le r$, so $r/2\le k\le r$.
\item[(S3)] $x^{11}\le T(s)=\#\anPa\le f\le q_1\le x^{12}$, and
$f=\varphi(q_0)\,\#\anPa$; moreover $1\le\#\anPa/\#\anPb\le\log b$.
\item[(S4)] $\anrd_E=\sqrt d\;q_1\prod_{p\mid q_0}p^{-1/(p-1)}$, and
$x^{10}\le\anrd_E\le x^{12}$.
\item[(S5)] $1\le xb/x_1\le\kappa(\tau)\le b\le a\le b^{2}$ for all $\tau$;
$b^{20}\le M\le b^{22}$;
$20\log b\le\log(M\anrd_E/2\pi)\le23\log b$; and
$m\log b\le\log N\le25\,m\log b$.
\item[(S6)] With $\Omega:=e^{\aneu}\cdot2\pi/\anrd_E$: $0<\Omega\le1$ and
$20\log b\le\log(M/\Omega)\le23\log b$.
\item[(S7)] With $n=m/2$ and $\varrho=n-1$ (so that $r=k+2\varrho+1$): $r$ is odd; every prime
dividing $q_1$ divides $q_0$; $q_0^{\,q_0}\le x$; $\varrho\ge10^{5}$;
$k\ge512000\,n^{9}$ and $k\ge1000\,(2\varrho+2)^{9}$; $k\ge2\varrho^{2}+3$; $18\varrho f\le b^{8}$;
$f\le b/1000$; $\anrd_E\le b$; and $N\le b^{60n}$. Moreover, with $\Delta=2r^{8}$:
\begin{itemize}
\item[(a)] $\varrho\ge10^{8}$ (so $n\ge100$); $n^{7}\le k$; $12n^{3}+500\le k$; $k-3n-6>10^{9}n^{3}$;
\item[(b)] $r\ge100$ (so $2r^{8}\ge500r^{3}+10r^{5}$); $3r^{3}f\le r^{7}$; $8nr\le b$; $kf\ge40$; $f\ge4000\,n$;
\item[(c)] $2f\le2^{2\varrho}$ and $2^{-\varrho}\le b^{-6}$;
\item[(d)] $(200\,f\,k\,\anrd_E)^{1000}\le b$ and $(25\,n^{3}f\,\anrd_E)^{1000}\le b$;
\item[(e)] $\sqrt d\le\anrd_E$, and $n,r,f,\sqrt d\le b$; $M\le b^{23}$;
\item[(f)] $\binom{a+k-1}{k-1}^{f}\le b^{2kf}$ (this bounds the dimension $d_\infty=\dim W_{\lambda^X}$ of
{\cite[\S14.3]{P46}}, a product of at most $f$ such binomial coefficients);
\item[(g)] $m\Delta+r\le\kappa(\tau)$ for all $\tau$; $k+2b+2m\Delta\le b^{2}$ (this bounds every
entry of $2\varsigma_\tau$, so that $L_\tau:=\max_i|2\varsigma_{\tau,i}|\le b^{2}$); and
$(r+1)^{4}(b^{2}+r+1+2n^{2})\le b^{5}$.
\end{itemize}
\end{enumerate}
\end{quotedthm}

\begin{quotedthm}{\cite[Theorem~13.9]{P46}}{ (existence of admissible data)}
Let $(q_0,d)$ be parameters and $0<\varepsilon<\frac1{100}$. There is
$x_0=x_0(q_0,d,\varepsilon)$ such that for every integer $x\ge x_0$ there are
admissible data $D$ for $(q_0,d)$ with scale $x$ and with
\[
 x_1=\bigl\lfloor\exp(x^{\varepsilon/6})\bigr\rfloor,\qquad x^{\varepsilon}\le\log b\le x^{\varepsilon}+1,\qquad\sqrt x\le m\le2\sqrt x .
\]
\end{quotedthm}

\begin{quotedthm}{\cite[Lemma~14.5]{P46}}{ (Free levels)}
Let $n\le2n_m$ and let $J\subseteq K^\circ$ consist of elements congruent to $1$ modulo every
chosen place. Then $\Gamma\cap gJg^{-1}=1$ for all $g\in G_f$.
\end{quotedthm}

\begin{quoteddef}{\cite[Definition~14.8]{P46}}{ (Prime sums, gain, main amount, remainder)}
Let $y\ge0$ be real and $\sigma\in\auPl$.
\begin{enumerate}
\item[(a)] The \emph{local prime sum} at $p$ is
$s_p(y,\sigma)=v(\pi_p)\sum_{e\in I_p}\bigl(y/p^e-\kappa(u(-p^e)\sigma)\bigr)^+$, and the
\emph{prime sum} is
$\ausfin(y,\sigma)=\sum_{1\le\nu\le a}\Lambda_K(\nu)\bigl(y/\nu-\kappa(u(-\nu)\sigma)\bigr)^+$.
\item[(b)] The \emph{gain} at $p$ is
\[
\augain_p=m\Bigl(\sum_{\tau\in\auPa}s_p(a,\tau)+k\sum_{\sigma\in\auPl}s_p(\kappa(\sigma),\sigma)
\Bigr).
\]
\item[(c)] The size of a derivative of order $j$ is $\ausinf(j)=\log j!+j\log(M/\Omega)$; the
\emph{direction amount} is $\auhdir=\sum_{\sigma\in\auPl}\bigl(\ausinf(\kappa(\sigma))-
\ausfin(\kappa(\sigma),\sigma)\bigr)$; the \emph{main amount} is $\auhmain=m\,k\,\auhdir$.
\item[(d)] The \emph{expected value} is
$\ausmain=a\bigl(\log a-\tfrac1f\sum_{\tau\in\auPl}\log\kappa(\tau)+c_1\bigr)$, and the
\emph{total remainder} is $\aurem=\sum_{\tau\in\auPa}\bigl(\ausfin(a,\tau)-\ausmain\bigr)$.
\end{enumerate}
\end{quoteddef}

\begin{quotedthm}{\cite[Lemma~14.9]{P46}}{ (The gains and the prime sums)}
Let $0\le y\le a$ and $\sigma\in\auPl$.
\begin{enumerate}
\item[(a)] $\ausfin(y,\sigma)=\sum_p2\log p\cdot s_p(y,\sigma)$, and $s_p(y,\sigma)=0$ if $p$
splits in $K$ or $p\ge x_1$. Hence $\augain_p=0$ for such $p$.
\item[(b)] With $\mathcal N$ the set of primes $p<x_1$ that do not split in $K$,
\[
2\sum_{p\in\mathcal N}\augain_p\log p=m\sum_{\tau\in\auPa}\ausfin(a,\tau)+m\,k\sum_{\sigma\in\auPl}
\ausfin(\kappa(\sigma),\sigma),
\]
and
\begin{equation*}\tag*{{\cite[(42)]{P46}}}
\auhmain-m\cdot\aurem+2\sum_{p\in\mathcal N}\augain_p\log p=m\,k\sum_{\sigma\in\auPl}\ausinf
(\kappa(\sigma))+m\cdot\#\auPa\cdot\ausmain .
\end{equation*}
\item[(c)] $\ausinf(a)-\ausmain\le a\bigl(\tfrac1f\sum_\tau\log\kappa(\tau)+\log\frac{M\aurd}
{2\pi}\bigr)+1+\tfrac12\log a$, and
\begin{equation*}\tag*{{\cite[(43)]{P46}}}
m\bigl(\ausinf(a)-\ausmain\bigr)-a\log N\le32\,a\log b .
\end{equation*}
\end{enumerate}
\end{quotedthm}

\begin{quotedthm}{\cite[Proposition~14.10]{P46}}{ (The count)}
Let $F$ have the normal form. Let $\lambda$ be a weight whose columns have height less than $n'$, with $a_\tau$ columns of
height $1$, $a'_\tau$ other columns and $l_\tau$ boxes in these, and $B\ge1$ boxes in all. Let $Q$ be either
the product $Q_{T,y}$, for a filling $T$ of $\lambda$ and integral directions $y$, of the minors of
$z_\tau y_\tau$ on the rows of each column and the first $h$ columns ($h$ the height of the column),
as in Proposition~\ref{pa:count}, or an entry $S_\lambda(z)_{T',T}$ of the Schur matrix with $T'$
standard (Lemma~\ref{pa:schur}). Then the jet
$Q(\partial)F(0)$ has valuation at least
\[
v(\pi_p)\sum_{e\in I_p}\sum_{\tau}\Bigl(\bigl(\tfrac{a_\tau}{p^e}-a_{\tau'}-a'_{\tau'}\bigr)^+
+\bigl(\tfrac{l_\tau}{p^e}-l_{\tau'}\bigr)^+-\tfrac{a'_\tau}{p^e}\Bigr)-f\,n'^3\bigl(1+\log_p
B\bigr),\qquad\tau'=u(-p^e)\tau.
\]
For $p=d$ the same bound holds after subtracting in
addition the term $\tfrac12\sum_\tau\mathrm{ded}_0(\tau)=v(\pi_d)\sum_\tau\mathrm{ded}_0(\tau)$ for $e=0$, with
$\mathrm{ded}_0$ as in Proposition~\ref{pa:count}, which is the precise statement in all cases.
\end{quotedthm}

\begin{quotedthm}{\cite[Theorem~14.13]{P46}}{ (the Short Forms Theorem)}
There are $c_{\mathrm{ss}}\in\C^\times$ and a positive rational number $\aucB$, a product of
powers of the primes in $\mathcal N\cup\{N\}$, such that the map
\[
\aupp=\aucB\,c_{\mathrm{ss}}^{-1}\,\auP_1|_{\auSram}\colon\auSram\to\auT
\]
has the following properties.
\begin{enumerate}
\item[(a)] \emph{(Size.)} $\auSram$ is spanned over $\C$ by the forms
$\varphi\in\auSram_{\mathrm{int}}$ such that, for all $s\in\auG$ and $\sigma\in\auAut(\C)$,
\[
\|\Theta_1(s\aupp\,\varphi)\|_\sigma^2\le\exp\bigl(\auhmain-m\cdot\aurem+f\,a\,(\log b)^4\bigr).
\]
\item[(b)] \emph{(Index.)} For every finite set $\mathsf t\subseteq\auG$, every finite family of
target forms each of which lies in $s\aupp(\auSram_{\mathrm{int}})$ for some $s\in\mathsf t$ has
index at most $\#\mathsf t\cdot\dim\auSram\cdot2\aueu$ in $\auLat_D$.
\end{enumerate}
\end{quotedthm}

\begin{quotedthm}{\cite[Lemma~14.18]{P46}}{ (The losses are affordable)}
$\sum_{p\in T}\auerr_p\log p\le\aueu$ for every finite set $T$ of primes, and
$\log\aucB\le-\sum_{p\in\mathcal N}\augain_p\log p+\aueu+b$.
\end{quotedthm}

\begin{quotedthm}{\cite[Proposition~14.19]{P46}}{ (Seesaw identity for the two chains)}
For every point $g_0\in G_0(E^+\otimes\R)$ and every linear form $\mathsf t$ on $W_{\lambda^X}$
there is $c_{\mathrm{ss}}\in\C^\times$ such that every twist $\vartheta$ has a chain datum $\Psi$
with
\[
\iota_\vartheta(\auP\varphi)=c_{\mathrm{ss}}\cdot B_{\vartheta,\Psi}\varphi\qquad\text{for all }
\varphi\in\auSram .
\]
The constant does not depend on $\vartheta$. The statement does not assert that either side is
non-zero.
\end{quotedthm}

\begin{quotedthm}{\cite[Proposition~14.20]{P46}}{ (Divisibility of the composite lift)}
Let $\ell$ be a prime and $j\colon\C\to\C_\ell$. There are $i\in\N$ and an $\auO_j$-submodule
$S\subseteq\auS_j\cap\auSram$ with $i\le\auerr_\ell\cdot\auDX$ and
$\aucol(\auS_j\cap\auSram,S)\mid\ell^{\,i}$ such that for every twist $\vartheta$ with
$\ell\nmid\auNm Q_\vartheta$, every chain datum $\Psi$ of $\vartheta$ and every $\varphi\in S$ the
form $B_{\vartheta,\Psi}\varphi$ has gauge at most $\ell^{\,\auerr_\ell-\augain_\ell}$ at $j$,
with $\auerr$ as in {\cite[(45)]{P46}}. More precisely:
\begin{enumerate}
\item[(i)] if $\ell\ne N$ and $\ell$ splits in $K$ or $\ell\ge x_1$, then $S=\auS_j\cap\auSram$
and the composite lift preserves integrality (on all of $\auS_j$: the hypothesis that $\varphi$ is
ramified is not used);
\item[(ii)] if $\ell\in\mathcal N$ is not exceptional, then $S=\auS_j\cap\auSram$ (and the bound
holds on all of $\auS_j$);
\item[(iii)] if $\ell=N$, the composite lift loses at most $2mfb$ powers of $N$ on the ramified
forms $\varphi$ such that $U_N^{m/2}\varphi$ is integral for the lattice $L_0$; the operator $U_N$
\eqref{pl:eq-U} and the lattice $L_0$ are defined in \S\ref{pl:ss-level}, at the primes of $E$
above $N$ conjugate to those induced by $j$ (Lemma~\ref{pl:level});
\item[(iv)] if $\ell\in\mathcal N$ is exceptional, $S=(Y_\Lambda^{\,m}\auS_j+\ell^{\,j_\ell}\auS_j)
\cap\auSram$, where $j_\ell=\lceil2mfa/\ell\rceil$ and $Y_\Lambda$ is the product of the normalised
Hecke operators $Y_{\mathfrak l}$ at the primes $\mathfrak l$ of a full mark $\Lambda$ of primes of
$E$ above $\ell$ (\S\ref{pa:ss-divis}, Theorem~\ref{pa:divis}(iv)); the operators $Y_{\mathfrak l}$
are not claimed to preserve $\auSram$.
\end{enumerate}
\end{quotedthm}

\begin{quotedthm}{\cite[Lemma~14.21]{P46}}{ (Short bases)}
Let $n\ge1$ and let $c_0$ be a positive integer in the different of $E$. For every
$g\in\auU(E^n)(\auA_{E^+,f})$ there is $Y\in\auGL_n(E)$ with $g^{-1}Y\in M_n(\widehat{\auO}_E)$
all of whose entries have all conjugates of absolute value at most $C_{\mathrm{sb}}=6c_0(fn)^2$.
\end{quotedthm}

\begin{quotedthm}{\cite[Proposition~14.22]{P46}}{ (Short ramified forms span)}
$\auSram$ is spanned by the forms $\varphi\in\auSram$ that are integral at every
$j\colon\C\to\C_\ell$, for every $\ell$, and satisfy $\|\sigma(\varphi(g))\|^2\le e^{\aueu}$ for
all $\sigma\in\auAut(\C)$ and all $g$.
\end{quotedthm}

\begin{quotedthm}{\cite[Proposition~14.23]{P46}}{ (One kernel)}
Let $0\le j<m$, let $\phi$ be a translate of the vacuum and $\sigma\in\auAut(\C)$. Then for all
$g,g'$ and all $\xi$,
\[
\|\sigma(\theta_j(\phi;g,g')\xi)\|\le\exp\bigl(x^{50}+\tfrac12\mathrm{rs}_j\bigr)\,\|\sigma(\xi)\|,
\qquad\mathrm{rs}_j=\sum_{\tau\in\auPl}\ \sum_{\text{rows }i}\ausinf(\text{length of row }i
\text{ at }\tau).
\]
\end{quotedthm}

\begin{quotedthm}{\cite[Proposition~14.24]{P46}}{ (Restriction to $Y_m$)}
For every form $\psi$ on $\auU(X_m)$ of weight $\lambda(X_m)$, every $\sigma\in\auAut(\C)$ and
every $y\in\auU(Y_m)(\auA_{E^+,f})$,
$\|\sigma((\operatorname{res}\psi)(y))\|^2\le\exp\bigl(x^{70}-(\sum_\tau n_\tau)\log N\bigr)\cdot
\|\sigma(\psi(\operatorname{diag}(y,1)))\|^2$.
\end{quotedthm}

\begin{quotedthm}{\cite[Proposition~14.25]{P46}}{ (Expansion of a chain datum)}
Let $\Psi$ be a chain datum of the twist $\vartheta$. There is a finitely supported function
$\Psi'$ on tuples of translations with $B_{\vartheta,\Psi}=\operatorname{res}\circ C_{\vartheta,
\Psi'}$ (no gluing), whose values are integral at the primes not dividing $N\auNm Q_\vartheta$,
and with $\sum_x|\sigma\Psi'(x)|\le\exp(x^{60})$ for all $\sigma\in\auAut(\C)$ if
$Q_\vartheta=\emptyset$.
\end{quotedthm}

\begin{quotedthm}{\cite[Proposition~14.26]{P46}}{ (Size of the composite lift)}
Let $\Psi$ be a chain datum of the trivial twist, $\sigma\in\auAut(\C)$, $\varphi\in\auS$ and
$B\in\R$ with $\|\sigma(\varphi(g))\|^2\le e^B$ for all $g$. Then for all $y$,
\[
\|\sigma((B_{1,\Psi}\varphi)(y))\|^2\le\exp\Bigl(B+\auhmain-m\cdot\aurem+2\sum_{p\in\mathcal N}
\augain_p\log p+\aueu\Bigr).
\]
\end{quotedthm}

\begin{quotedthm}{\cite[Proposition~14.27]{P46}}{ (Equivariance of the lower chain)}
For every tame character $\vartheta$ and every $u\in\auG$ there is $\beta_\vartheta(u)$ in the
stabiliser $K^\circ$ of the lattice of $Y_m$ such
that $(\iota_\vartheta(\varrho(u)z))(g)=(\iota_\vartheta z)(g\beta_\vartheta(u))$ for all
$z\in\auAz$ and all $g$.
\end{quotedthm}

\begin{quotedthm}{\cite[Proposition~14.28]{P46}}{ (Passage to the good part)}
Let $z\in\auAz$, $s\in\auG$, and $z'=\#\auG\,\varrho(e_0)\varrho(s)z$.
\begin{enumerate}
\item[(a)] If $c\in\C$ and $c\cdot\iota_\vartheta z$ is integral at $j\colon\C\to\C_\ell$, so is
$c\cdot\iota_\vartheta z'$.
\item[(b)] If $\sigma\in\auAut(\C)$ and all values of $\sigma(\iota_1z)$ have squared norm at most
$e^B$, then $\|\iota_1z'\|_\sigma^2\le\exp(x^{99}+B)$.
\end{enumerate}
\end{quotedthm}

\begin{quotedthm}{\cite[Lemma~14.29]{P46}}{ (Count of the ramified forms)}
$\dim\auSram\ge\bigl(1-f/(N-1)\bigr)\auDX\ge(1-b^{-8})\auDX$.
\end{quotedthm}

\begin{quotedthm}{\cite[Proposition~14.32]{P46}}{ (Second moment of the Poincar\'e map)}
Let $\deg$ be the product over $\tau\in\auPl$ of the formal degrees of the discrete series
$\varsigma_\tau$, and $\mathrm{mass}_r$ the mass of the genus of the standard hermitian lattice of
rank $r$. There are an inner product on $\auS$ and a scalar $c_P>0$ such that $P=c_P\auP$
satisfies:
\begin{enumerate}
\item[(i)] $0<\audT\le\#\auG\cdot\deg\cdot\mathrm{mass}_r$;
\item[(ii)] $\|\varrho(e_0)P\,\mathrm{pr}_U\|^2\ge\dim U-5b^{-4}\auDX$ for every subspace
$U\subseteq\auS$;
\item[(iii)] $\sum_{s\in\auG}\|P^*\varrho(s)P\|^2\le(1+9b^{-4})\,\auDX^2/(\deg\cdot
\mathrm{mass}_r)$.
\end{enumerate}
\end{quotedthm}

\begin{quotedthm}{\cite[Theorem~2.2]{P50}}{ (Elementary divisors of an edge Hecke operator)}
Assume {\cite[Propositions~2.10, 2.12 and~2.13]{P50}}. Let $0<\theta\le1$ and
\begin{equation*}\tag*{{\cite[(5)]{P50}}}
n^4\le q,\qquad\bigl(Nk\log q/\theta\bigr)^5P^{3r}\le n .
\end{equation*}
Let $A$ be as in {\cite[(4)]{P50}}, let $f\colon\mathcal F_X\to M$ be an $\pao$-linear map and $Y$
an endomorphism of $M$ with $Yf=fA$, and $\Lambda=f(\mathcal F_X)$. Then $Y\Lambda\subseteq
\Lambda$, and for every integer $j\ge0$ there is an integer
\[
l\le\bigl(3\phcd+\theta j\bigr)\,|X|\,|H|\,|\phB|,\qquad\phcd=Nn\Bigl(\frac1{P^{2r}}+\frac{r^2}{kP}
\Bigr),\qquad\text{with}\quad\aucol(\Lambda,\,Y\Lambda+p^j\Lambda)\mid p^{\,l}.
\]
\end{quotedthm}

\begin{quotedthm}{\cite[Lemma~2.3]{P50}}{}
Let $\nu\ge1$. {\rm(a)} For $B\in M_\nu(\pao)$, $\aucol(\pao^\nu,B\pao^\nu)\mid\det B$.
{\rm(b)} If $M'\subseteq\pao^\nu$ is a submodule, $s$ a finite subset and $z\in\pao$, then
$\aucol(M'+\langle s\rangle,\,M'+z\pao^\nu)\mid z^{|s|}$.
{\rm(c)} Let $Z,d,d'\in M_\nu(\pao)$, $a,z\in\pao$ and $s\subset\pao^\nu$ finite with
$a\pao^\nu\subseteq Z\pao^\nu+\langle s\rangle$. Then
$\aucol\bigl(\pao^\nu,\ d'Zd\,\pao^\nu+z\pao^\nu\bigr)\mid\det d'\cdot a^\nu\cdot\det d\cdot
z^{|s|}$.
\end{quotedthm}

\begin{quotedthm}{\cite[Proposition~2.10]{P50}}{ (Root lemma)}
Let $R$ be $\pao$ or $\phk$ (with $w=0$ in $\phk$), let $c$ be real weights of the variables with
$c(Y_s)-c(X_s)\in[\tfrac12,1]$ for all $s$, $f$ a polynomial over $R$, $\alpha$ an exponent and
$\lambda\ge0$ an integer. Suppose that for every $g\in\mathrm{SL}_2^{a,a'}(W)$ and every monomial
$\beta$, the coefficient of $x^\beta$ in $g\cdot f$ is divisible by $w^{\lceil\lambda+c\cdot\alpha
-c\cdot\beta\rceil^+}$, and that $w^{\lambda+1}$ does not divide the coefficient of $x^\alpha$ in
$f$. Then $\phdef_i^{a,a'}(\alpha)\le0$ for every letter $i$.
\end{quotedthm}

\begin{quotedthm}{\cite[Proposition~2.12]{P50}}{ (Count of the forms of bounded defect)}
Let $c\ge1$, $0<\theta\le1$ and $24\,Nk^2c\,P^{2r}\le\theta n$. Then
$[G:K_r]\cdot\dim_{\phk} S_c\le(1+\theta)|\phB|$.
\end{quotedthm}

\begin{quotedthm}{\cite[Proposition~2.13]{P50}}{ (Averaging)}
Let $N=2\ell$ and $Nk(\ell+2)n+n+k<P^\ell$, and $c_0=k^2eN$. For every $c\ge1$ there is
$\mathrm{Av}\in M_\phB(\pao)$ such that $\mathrm{Av}\,h\in C_c$ for all $h\in\pao^\phB$, and
\[
{}^th'\,\mathsf b\,\mathrm{Av}\,h\equiv p^{Nc_0}\;{}^th'\,\mathsf b\,h\pmod{w^c}\qquad(h\in
\pao^\phB,\ h'\in\phA).
\]
\end{quotedthm}

\begin{quotedthm}{\cite[Lemma~2.14]{P51}}{ (Change of the lattice)}
Let $\mathsf u$ be a finite idele of $E$ with $\mathsf u\aubar{\mathsf u}=1$, let $\mathsf h_{\mathsf u}$
be the element of the centre of $H(\auA_{E^+,f})$ that acts on $X_0$ by the scalar $\mathsf u^{-1}$
(and trivially on $A$, on $z$ and on the $e_a^*$), and let $L_{\mathsf u}$ be the lattice of $Y_0$
spanned over $\widehat\auO_E$ by $x_1,\dots,x_{k-1}$, $x_k/\aubar{\pi_N}+\mathsf uz/N$, $\mathsf uz$,
$\mathsf ue_a$, $\mathsf ue_a^*$ $(1\le a\le\varrho)$.
\begin{enumerate}
\item[(a)] $L_{\mathsf u}=\mathsf u\cdot\mathsf h_{\mathsf u}L$. So $L_{\mathsf u}$ is unimodular, its
stabiliser is $K_{\mathsf u}=\mathsf h_{\mathsf u}K\mathsf h_{\mathsf u}^{-1}$, the principal congruence
subgroup of level $\pi_M$ in it is $J_{\mathsf u}=\mathsf h_{\mathsf u}J\mathsf h_{\mathsf u}^{-1}$, and
the reduction of $k\in K_{\mathsf u}$ at the chosen places, in the basis of $L_{\mathsf u}$ just given, is
the image of $\mathsf h_{\mathsf u}^{-1}k\,\mathsf h_{\mathsf u}\in K$ in $\auG$.
\item[(b)] Let $\auAz^{\mathsf u}$ be the space $\auAz$ of {\cite[\S14.3]{P46}} formed with
$J_{\mathsf u}$ in place of $J$, with the action $\varrho_{\mathsf u}$ of $\auG=K_{\mathsf u}/J_{\mathsf u}$
by right translation, and let $\auP_{\mathsf u}\colon\auS\to\auAz^{\mathsf u}$ be the Poincar\'e map
defined by the same $g_0$, $\mathsf t$, $\psi_U$ and Haar measure of $[R_0]$ as $\auP$. Then
$(\mathsf j_{\mathsf u}y)(g)=y(g\,\mathsf h_{\mathsf u}^{-1})$ defines a unitary isomorphism
$\mathsf j_{\mathsf u}\colon\auAz^{\mathsf u}\to\auAz$ with $\mathsf j_{\mathsf u}\varrho_{\mathsf u}(s)=
\varrho(s)\mathsf j_{\mathsf u}$ for $s\in\auG$, the map $(R_{\mathsf u}\varphi)(h)=\varphi(h\,\mathsf
h_{\mathsf u}^{-1})$ is a unitary automorphism of $\auS$, and
\[
\auP_{\mathsf u}=\mathsf j_{\mathsf u}^{-1}\circ\auP\circ R_{\mathsf u} .
\]
\item[(c)] If {\rm(i)}, {\rm(ii)}, {\rm(iii)} of {\cite[Proposition~14.32]{P46}} hold for $P=c_P\auP$,
the space $\auAz$ and $\auT=\varrho(e_0)\auAz$, then they hold, with the same inner product on $\auS$
and the same $c_P$, for $P_{\mathsf u}=c_P\auP_{\mathsf u}$, the space $\auAz^{\mathsf u}$ and
$\auT^{\mathsf u}=\varrho_{\mathsf u}(e_0)\auAz^{\mathsf u}$; and $\dim\auT^{\mathsf u}=\dim\auT$.
\end{enumerate}
Another basis of $L_{\mathsf u}$ changes the identification of $K_{\mathsf u}/J_{\mathsf u}$ with $\auG$ by
an inner automorphism of $\auG$; the central idempotent $e_0$, the space $\auT^{\mathsf u}$ and the
statements {\rm(i)--(iii)} do not depend on it.
\end{quotedthm}

\begin{quotedthm}{\cite[Lemma~4.7]{P51}}{}
Let $e_\tau$ be as in {\cite[(42)]{P51}}, on $G_\tau\cong G_s$.
\begin{enumerate}
\item[(a)] $e_\tau$ is continuous, $e_\tau^*=e_\tau$, $|e_\tau(g)|\le C\|g\|^{-B}$,
$e_\tau\in L^1$, and $e_\tau*e_\tau=e_\tau$.
\item[(b)] There is $\alpha_\tau\in C_c^\infty(G_\tau)$ with $\alpha_\tau*e_\tau=e_\tau$. Hence
$e_\tau=e_\tau*\alpha_\tau^*$, $e_\tau$ is smooth, and $|(De_\tau)(g)|\le C_D\|g\|^{-B}$ for every
left-invariant differential operator $D$.
\item[(c)] $\int_{G_\tau}\|g\|^{-B'}dg<\infty$ for $B'>n^2$.
\end{enumerate}
The same holds for $e$ on $G_\infty$ with $\alpha=\bigotimes\alpha_\tau$, with $\pmP(g)^{-B}$ in place
of $\|g\|^{-B}$.
\end{quotedthm}

\begin{quotedthm}{\cite[Lemma~4.8]{P51}}{ ($e$ is a cusp form)}
Let $N'$ be the unipotent radical of a proper parabolic subgroup of $G_\tau$. Then
$\int_{N'}|e_\tau(x\nu y)|\,d\nu<\infty$ and $\int_{N'}e_\tau(x\nu y)\,d\nu=0$ for all $x,y\in G_\tau$.
\end{quotedthm}

\begin{quotedthm}{\cite[Corollary~4.43]{P51}}{ (Formal degree)}
Under {\cite[(49)]{P51}}, $d=\dim V_\sigma=\prod_{i<j}(\lambda_i-\lambda_j)/(j-i)$ (entries listed in
decreasing order), and $\sigma_\lambda=(-1)^{pq}$.
\end{quotedthm}

\begin{quotedthm}{\cite[Proposition~5.2]{P51}}{}
Let $dg$ be the Haar measure on $G$ compatible with the Haar measure of mass $1$ on $\auU(N)$. Then
$\int_G|c_\tau|^2dg=\pbdeg_\tau^{-1}$ with $\pbdeg_\tau=\prod_{i<j}(\lambda_i-\lambda_j)/(j-i)$
(entries in decreasing order). Thus the measures of \textup{(DS1)(ii)} and \textup{(DS1)(iii)} agree.
\end{quotedthm}

\begin{quotedthm}{\cite[Lemma~5.4]{P51}}{}
For $g\in G$: $a(g)$ is invertible, $|\det a(g)|\ge(\|g\|/2)^2$, and $|\det a(kgk')|=|\det a(g)|$ for
$k,k'\in\pdsK$. If $g$ is positive definite, so is $a(g)=\frac12(g+g^{-1})$.
\end{quotedthm}

\begin{quotedthm}{\cite[Lemma~5.5]{P51}}{}
Let $w\in G$ be positive definite and $0\le c\le h$ an integer. Then $\pdsm_{-\mu}(w)=\pdsm_\mu(w)>0$
and $\det a(w)^c\le\pdsm_\mu(w)^2$.
\end{quotedthm}

\begin{quotedthm}{\cite[Theorem~5.6]{P51}}{ (\textup{(DS2)})}
For every $g\in G$ and every integer $0\le c\le h$ (formula {\cite[(59)]{P51}} holds by Theorem~\ref{pco:FJ}),
\[
|c_\tau(g)|\le|\det a(g)|^{-c/2}\le\bigl(2/\|g\|\bigr)^c .
\]
In particular $|c_\tau(g)|\le(2/\|g\|)^{j+\pbeq}$ for $0\le j\le r^7$, $\pbeq=9r^2$, and for $g\in
G_\tau$ with the norm of {\cite[\S3.2.1]{P51}}.
\end{quotedthm}

\begin{quotedthm}{\cite[Lemma~5.10]{P51}}{}
Let $\zeta\in E$ be a primitive $q_1$-th root of unity and $\varpi=\gamma\,\Phi_{q_1}'(\zeta)\sqrt{-d}$,
where $\gamma\in\auO_E^\times$ satisfies $\gamma\aubar\gamma=1$ and $\aubar{\gamma\Phi'_{q_1}(\zeta)}=
\gamma\Phi'_{q_1}(\zeta)$ \textup{(}such $\gamma$ exist, e.g.\ $\gamma=\zeta^s$ with
$s=q_1/4+(q_1-\varphi(q_1))/2+1$; see the first paragraph of the proof\textup{)}. Then $\varpi$ is a purely imaginary generator of the different of $E$, and
for every $\tau\in\auPl$, $2\le|\tau(\varpi)|\le b$ and $1\le M|\tau(\varpi)|/2\pi\le b^{23}$.
\end{quotedthm}

\begin{quotedthm}{\cite[Lemma~2.14]{P53}}{ (Clean data)}
Let $S,T$ be as in \textup{(SW)}, $\pcW=S\oplus S^-$, and let $F\in\psHS(\auHom(S,T))$ be a datum
whose component at one split place $v=w\aubar w$ is, in the model \textup{(W5)}, a Schwartz--Bruhat
function supported on injective matrices. Let $F'\in\psHS(\auHom(S,T))$ and
$\Phi=F\otimes\mathsf CF'\in\psHS(\auHom(S,T))\otimes\psHS(\auHom(S^-,T))=\psHS(\auHom(\pcW,T))$.
\begin{enumerate}
\item[(a)] If the Lagrangian $L$ of $\pcW$ is not of the form $L_\delta=\{(\delta s,-s):s\in S\}$
for an element $\delta\in\auU(S)(E^+)$, then $\psZ_L(\Phi)=0$.
\item[(b)] There is a constant $c'\ne0$, depending only on the normalisations, such that for all
$\delta\in\auU(S)(E^+)$ and all $F,F'$
\[
\psZ_{L_\delta}(F\otimes\mathsf CF')=c'\,\bigl\langle F',\omega_{\auHom(S,T)}(\delta^{-1},1)F\bigr
\rangle .
\]
\end{enumerate}
\end{quotedthm}

\begin{quoteddef}{\cite[Definition~3.27]{P53}}{ (Lattice and height at a point)}
Let $S_1$ be a space of rank $n$ with frame giving $S_1(\auA)\cong\auA_E^n$. For $g\in\auU(S_1)(\auA)$
the \emph{lattice at $g$} is $\Lambda_g=\{\gamma\in S_1(E): g^{-1}\gamma\text{ is integral at all
finite places}\}$, with the length $|\gamma|_g=\|(g^{-1}\gamma)_\infty\|$ (the euclidean norm of the
archimedean coordinates). Let $\lambda_1(g)=\min\{|\gamma|_g:0\ne\gamma\in\Lambda_g\}$ and
$\mathrm{ht}(g)=\max(1,\lambda_1(g)^{-1})$.
\end{quoteddef}

\begin{quotedthm}{\cite[Lemma~3.28]{P53}}{ (Reduction to a Siegel set)}
Let $S_1$ be a space of rank $n$ with a frame, $G=\auU(S_1)$, viewed as a connected reductive group over
$E^+$, and $d=[E^+:\Q]$. There are a compact set $C_f\subset G(\auA_f)$ and a constant $C$ such that every
$g\in G(\auA)$ can be written $g=\gamma^{-1}x$ with $\gamma\in G(E^+)$, $x=x_\infty x_f$, $x_f,x_f^{-1}\in C_f$
and $\|x_\infty\|,\|x_\infty^{-1}\|\le C\,\mathrm{ht}(g)$ (operator norms in the
coordinates of the frame).
\end{quotedthm}

\begin{quotedthm}{\cite[Lemma~3.30]{P53}}{ (Fock vectors are superpositions of translated vacua)}
Let $\mathcal F$ be the Fock space of a real hermitian space $\C^m$ (real dimension $2m$), realised as in
{\cite[\S5.1.4]{P54}} as the space of functions $v\in L^1\cap L^2$ on $\C^m$ with $v(x)=\int v(y)\beta(y,x)
e(x-y)\,dy$, where $e(x)=e^{-|x|^2/2}$ is the vacuum, $|\beta|=1$, and the Heisenberg translations are
$(\rho(y)v)(x)=\beta(y,x)v(x-y)$; the measure is the normalised Lebesgue measure. For $N\ge0$ put
$\|v\|_N=\int(1+|x|)^N|v(x)|\,dx\in[0,\infty]$. Let $s\ge0$ be an integer, and let $\|\cdot\|_s$ and
$\mathcal F_s\subseteq\mathcal F$ be the norm and the Banach space of Lemma~\ref{pz:growth} (with $D=m$),
transported by the unitary isomorphism with the model \eqref{pk:eq-fock} that intertwines the
Heisenberg translations and carries vacuum to vacuum; the operators $\partial_w$, being the derivatives
of $\rho$, correspond. Let $c_\rho$ be a constant with
$\rho(w)\rho(y)=e^{\kappa(w,y)}\rho(y)\rho(w)$, $\kappa(w,y)\in i\R$ real bilinear,
$|\kappa(w,y)|\le c_\rho|w||y|$ (in the model \eqref{pk:eq-fock}, $\kappa(w,y)=2i\Iml\operatorname{tr}(y^*w)$ by the
formula for $\psrho_\infty(x)\psrho_\infty(y)$ recalled after it, so $c_\rho=2$ will do).
\begin{enumerate}
\item[\textup{(a)}] For $y\in\C^m$, $\rho(y)$ maps $\mathcal F_s$ into itself, with
$\|\rho(y)\Psi\|_s\le(1+c_\rho|y|)^{2s}\|\Psi\|_s$, and $y\mapsto\rho(y)\Psi$ is continuous into
$\mathcal F_s$ for every $\Psi\in\mathcal F_s$.
\item[\textup{(b)}] If $\|v\|_N<\infty$ then $v=\int v(y)\rho(y)e\,dy$, a Bochner integral converging for
the norm $\|\cdot\|_N$, because $\|\rho(y)e\|_N\le(1+|y|)^N\|e\|_N$. If $\|v\|_{2s}<\infty$, then
$v\in\mathcal F_s$, the same integral converges in $\mathcal F_s$, and $\|v\|_s\le c_1\|v\|_{2s}$ with
$c_1=\max(1,c_\rho)^{2s}\|e\|_s$.
\item[\textup{(c)}] Let $\ell$ be a continuous linear functional on $\mathcal F_s$ with $|\ell(\rho(z)e)|\le
S(1+|z|)^N$ for all $z$. Then $|\ell(v)|\le S\|v\|_N$ for every $v$ with $\|v\|_{\max(N,2s)}<\infty$.
\end{enumerate}
\end{quotedthm}

\begin{quotedthm}{\cite[Proposition~4.10]{P53}}{ (Step (ii))}
Let a square be as in {\cite[\S4.1.3]{P53}}, $\pcm=\pcm_\vartheta$, $y_1$ a form of type
$e_{i,n+1}$ on $U_1$ and $y_2$ a form of type $\overline{e_{i+1,n+1}}$ on $U_2$, $\xi\in\pcMat_{a\times
N}(\pcR_\pcm)$ of rank $a=\dim A$ at every prime of $\pcm$ and $\eta\in\pcMat_{N\times b}(\pcR_\pcm)$ of rank
$b=\dim A'$, and let $\chi_2$ be a unitary character of $\auU(1)$. There are $M_0,C$ such that
for every Fock vector $x$ of $\auHom(W,A)\otimes\R$ and $v$ of $\auHom(A',W)\otimes\R$ with finite
moments of order $M_0$:
\begin{enumerate}
\item[(a)] the function $(g_1,g_2,h)\mapsto y_1(g_1)y_2(g_2)\chi_2(\det g_2)\,\theta^{W,A}(g_1\oplus g_2,h;
[x,\delta_\xi])$ is integrable on $[\auU(U_1)]\times[\auU(U_2)]\times[\auU(A)]$, with integral of absolute
value at most $C\|x\|_{M_0}$;
\item[(b)] the function $(g_1,g_2,h')\mapsto y_1(g_1)y_2(g_2)\chi_2(\det g_2)\,\theta^{A',W}(h',g_1\oplus
g_2;[v,\delta_\eta])$ is integrable on $[\auU(U_1)]\times[\auU(U_2)]\times[\auU(A')]$, with the same bound
in $\|v\|_{M_0}$.
\end{enumerate}
(The characters of the kernels are any; $\|\cdot\|_M$ is the moment of {\cite[Proposition~3.26]{P53}}.)
In particular all iterated integrals in steps (i), (iii) of {\cite[\S2.6]{P53}} converge
absolutely and may be computed in any order.
\end{quotedthm}

\begin{quotedthm}{\cite[Proposition~4.40]{P53}}{ (Reversal)}
For every modulus $\pcm$ there are a character $\psi_\pcm$ of the additive group of $\pcR_\pcm$,
non-trivial on every factor $\auF_w$, a scalar $\zeta_\pcm\in\auU(1)(\auA_{E^+,f})$ equal to
$\zeta_\alpha$ outside the primes of $\pcm$, both independent of $S,T$, and for all $\alpha,\beta$ a
number $c=c(S,T,\pcm,\alpha)$ with $|c|^2=C_{S,T}\,(\#\pcR_\pcm)^{-NN'}$, $C_{S,T}>0$ depending only on
$S,T$, such that for every Fock vector $P$, every $\Psi$ on $\pcMat_{N'\times N}(\pcR_\pcm)$ and all
$g\in\auU(S)(\auA)$, $h\in\auU(T)(\auA)$
\[
\aubar{\theta_{\alpha,\beta}(g,h;P,\Psi)}=c\,\theta_{\beta\lambda_T,\,\alpha\lambda_S}(h,g\zeta_\pcm;
P^\flat,\Psi^\flat),\qquad\Psi^\flat(\varrho)=\sum_{v\in\pcMat_{N'\times N}(\pcR_\pcm)}\aubar{\Psi(v)}\,
\psi_\pcm(\ptvtr(\varrho v)),
\]
where the right side is a kernel of $\auHom(T,S)$, $\varrho\in\pcMat_{N\times N'}(\pcR_\pcm)$, and
$P^\flat=\mathcal A_\infty P$, a Fock vector with the same moments as $P$.
\end{quotedthm}

\begin{quotedthm}{\cite[Lemma~5.37]{P53}}{ (Line reversal when $\alpha$ is a unit at the primes of $\pcm$; the general case is {\cite[Lemma~6.33]{P53}})}
Assume that $\alpha$ is a unit at the primes of $\pcm$. There are an element $\pgupt=\pgupt_\pcm$ of $\auU(1)(\auA_{E^+,f})$ depending only on $\pcm$, a
linear bijection $\tau$ from the polynomials on $\auHom(L,S)\otimes\R$ to those on
$\auHom(S,L')\otimes\R$, independent of $\vartheta$, a primitive additive character $\psi_\pcm$ of
$R$, and $c_{\mathrm{line}}=c_{\mathrm{line}}(\vartheta)\in\C^\times$ with
\[
|c_{\mathrm{line}}|^2=C_0\,(\auNm\pcm)^{-n},\qquad C_0>0\text{ independent of }\vartheta,
\]
such that for every $k\in\auU(L)(\auA_f)$, $g\in\auU(S)(\auA)$, $v$ and $B$
\[
\pguK^{\mathrm c}_{v,B}(k,g)=c_{\mathrm{line}}\,\zeta(k)
\sum_{v'\in\pcMat_{\bullet\times\mathsf N_2}(R)}\psi_\pcm(v'v)\,\pguK^{\mathrm r}_{v',\tau B}(g,k\pgupt)
\]
where $\zeta(k)$ is a number of modulus one depending only on $k$, on the characters and on $\pcm$.
Here $\pgupt$ is the element constructed in the proof (it is $k_0$ outside $S(\pcm)$ and has
$w$-component $\varpi^{-k_w}$ at $v\in S(\pcm)$); it is in general not the element $\pgupt$ of
(G4), and the case of (G4) is treated by {\cite[Lemma~6.33]{P53}}.
\end{quotedthm}

\begin{quotedthm}{\cite[Proposition~6.9]{P53}}{}
\begin{enumerate}
\item[(a)] $\alpha$ is a unit at every prime $w$ of $E$ with $w\nmid q_1$; in particular at every
chosen place, and at every prime of $\pcm(\vartheta)$ if $\auNm Q_\vartheta$ is prime to $q_1$.
\item[(b)] Let $p\mid q_1$ be a prime that ramifies in $\Q(\zeta_{q_1})$ (every odd $p\mid q_1$, and
$p=2$ if $4\mid q_1$). Then $p$ splits in $K$, and every prime $w$ of $E$ above $p$ is split over $E^+$
($w\ne\aubar w$) and ramified over $K$; $\alpha$ is not a unit at the primes $w\mid p$ that divide
$\omega$.
\end{enumerate}
\end{quotedthm}

\begin{quotedthm}{\cite[Lemma~6.30]{P53}}{ (The vacuum at a split place)}
The self-dual lattice $M_v=\alpha\auHom(L_S,L_T)_v$ is $M_w\oplus M_{\aubar w}$ with
$M_w=\alpha_w\pcMat_{N'\times N}(\auO_w)$ and $M_{\aubar w}$ the dual of $M_w$ for the pairing
$(x,a')\mapsto[x,a']$ of $X_w\times X_{\aubar w}$. In the model (W5) the vacuum $\psvac_{M_v}$ is,
up to a scalar, the indicator function $\mathbf 1_{M_w}$.
\end{quotedthm}

\begin{quoteddef}{\cite[Definition~6.31]{P53}}{ (Data at the primes of a modulus)}
For every $w\mid\pcm$ let $a_w$ be the generator of the ideal $\alpha_w\auO_w$ fixed in
\S\ref{ps:ss-weil} ($a_w=\alpha_w$ if $\alpha_w$ is not a unit, $a_w=1$ otherwise). For a function $\Psi$ on
$\pcMat_{N'\times N}(\pcR_\pcm)$ let $\Psi^\natural$ be the function on $\prod_{w\mid\pcm}
\pcMat_{N'\times N}(E_w)$ equal to $\Psi(a^{-1}x\bmod\pcm)$ at $x\in a\,\pcMat_{N'\times N}
(\widehat\auO)$ (that is, $x_w\in a_w\pcMat(\auO_w)$ for all $w\mid\pcm$) and to $0$ elsewhere, and
use at $v\in S(\pcm)$ the $L^2$-norm for the measure giving $M_w=a_w\pcMat(\auO_w)$ the volume $1$.
The datum $[P,\Psi]$, the kernels, the lifts and the transposed lifts are then defined by the
formulas of {\cite[\S2.3]{P53}}.
\end{quoteddef}

\begin{quotedthm}{\cite[Lemma~6.32]{P53}}{ (The normalised coordinates)}
Let $a$ be the idele of \S\ref{ps:ss-weil}, $S,T$ spaces of ranks $N,N'$, $v=w\aubar w$ a place of
$E^+$ split in $E$, and $X_w,X_{\aubar w}$ the components of $\auHom(S,T)_v$, written as matrices $x$
in the frames. Put $\tilde x=a_w^{-1}x$ on $X_w$ and $\tilde x=a_{\aubar w}^{-1}x$ on $X_{\aubar w}$, and
more generally $\tilde x=a^{-1}x$ on $\auHom(S,T)(\auA_f)$.
\begin{enumerate}
\item[(a)] $\alpha_w$ is not a unit if and only if $w$ lies above a prime $p\mid q_1$ and $w\mid\omega$;
then $w\ne\aubar w$, $\aubar w\nmid\omega$ and $\alpha_{\aubar w}=a_{\aubar w}=1$. Every place of $E^+$
above $p\mid q_1$ is $w\aubar w$ with exactly one of $w,\aubar w$ dividing $\omega$. No chosen place
and no prime above $N$ lies above $q_1$. A prime of $Q_\vartheta$ may lie above $q_1$ and divide
$\omega$; such primes are the only primes of a modulus $\pcm_\vartheta$, of $S(\pcm_\vartheta)$, or of
the moduli of the grids at which $a_w\ne1$.
\item[(b)] The self-dual lattice $\alpha\auHom(L_S,L_T)$ is $a\auHom(L_S,L_T)$; at $v$ its components
are $M_w=a_w\pcMat(\auO_w)$ and $M_{\aubar w}=a_{\aubar w}\pcMat(\auO_{\aubar w})$, i.e.\ the integral
matrices in the coordinates $\tilde x$, and they are dual to each other for $[\,,]$. For $x\in X_w$,
$y\in X_{\aubar w}$ one has $[x,y]=\lambda_w(\operatorname{tr}({}^t\tilde y\tilde x))^{\pm1}$ with
$\lambda_w(t)=\psi_w(\jmath(a\aubar a)_wt)$, trivial on $\auO_w$ and not on $\varpi_w^{-1}\auO_w$
($E_{\aubar w}$ identified with $E_w$ by conjugation). In the frame coordinates the dual of
$\pcMat(\auO_w)$ is $\aubar{a_w}\,a_{\aubar w}\pcMat(\auO_{\aubar w})$, which is not
$\pcMat(\auO_{\aubar w})$ when $a_w\ne1$ or $a_{\aubar w}\ne1$.
\item[(c)] On $X_w$ the coordinate $T$ of (N5) and the coordinate $a_w^{-1}x$ of
{\cite[Definition~6.31]{P53}} are both $\tilde x$. Hence for $w\mid\pcm$:
$\Psi^\natural=\Psi(\tilde x\bmod\pcm)\,\mathbf 1[\tilde x\ \text{integral}]$; the datum
$\delta_\iota^\natural$ is $\mathbf 1_{a_w(\iota+\varpi\pcMat(\auO_w))}$ in the frames, which is the
datum of \S\ref{ps:ss-Y} at a level place; and a datum defined as a function of $T$ by (N5) (as in
{\cite[\S\S2.5, 3.2 and~4.1.3]{P54}}) is the same function of $\tilde x$, so it is
of the form $\Psi^\natural$ exactly when, as a function of $T$, it is supported on the integral
matrices and depends only on $T$ modulo $\pcm$.
\item[(d)] For $u\in\auHom(S,T)(\auA_f)$, $u\in\alpha\auHom(L_S,L_T)$ if and only if $\tilde u$ is
integral. If $u=s_0y$ with $s_0=a$, then $\tilde u=y$; for $h\in\auU(T)(\auA_f)$,
$\widetilde{h^{-1}u}=h^{-1}y$.
\item[(e)] Let $u$ be a finite idele of $E$ with $u\aubar u=1$, acting as a central element of
$\auU(T)(\auA_f)$. Then $\omega(1,u)$ carries the vacuum of a self-dual lattice $\Lambda$ of
$\auHom(S,T)(\auA_f)$ to a multiple of modulus one of the vacuum of $u\Lambda$, and for every datum $f$,
all characters $\alpha',\beta'$ and $g\in\auU(S)(\auA)$,
\[
\int_{[\auU(T)]}\theta\bigl(\omega_{\alpha',\beta'}(g,h)\,\omega(1,u)f\bigr)\,dh=\beta'(\det u)^{-1}
\int_{[\auU(T)]}\theta\bigl(\omega_{\alpha',\beta'}(g,h)f\bigr)\,dh
\]
whenever either side converges absolutely.
\item[(f)] The map $w\mapsto w^\dagger$ carries $\alpha\auHom(L_S,L_T)$ onto $\aubar\alpha\auHom(L_T,L_S)=
u\,\alpha\auHom(L_T,L_S)$, $u=\aubar a/a$ (a finite idele with $u\aubar u=1$, central in $\auU(S)$,
equal to $1$ at the places not above $q_1$). At $v$, in the coordinates $\tilde x$ of $\auHom(S,T)$ and
$\tilde X$ of $\auHom(T,S^-)$ (frames of $S$ and $T$), it is $\tilde X=-\nu^{-1}\,{}^t\tilde x$ after
identifying $E_{\aubar w}$ with $E_w$ by conjugation, where $\nu=\aubar{a_{\aubar w}}/a_w$ is the
conjugate of $u_{\aubar w}^{-1}$.
\end{enumerate}
\end{quotedthm}

\begin{quoteddef}{\cite[Definition~6.14]{P54}}{ (The characters attached to the data)}
\begin{enumerate}
\item[(i)] For each chosen place $w$ fix a character $\omega_w$ of $\auF_w^\times$ of order exactly $q_2$ (it exists by
{\cite[Lemma~6.12(d)]{P54}}); for the other finite places put $\omega_w=1$. Write $\omega=(\omega_w)_w$.
\item[(ii)] $\chi_M$ is a character of $\auU(1)$ of weight $0$ with tame part $\omega$: {\cite[Lemma~6.13]{P54}} with $S$ the chosen places,
$c_\tau=0$, $\psi_w=\omega_w$. {\cite[(39)]{P54}} holds factor by factor: $\omega_w(\xi\bmod w)=1$ because $\xi\bmod w$ is a
$q_2$-th power ({\cite[Lemma~6.12(d)]{P54}}) and $\omega_w^{q_2}=1$. We write $\chi_{M,w}=\omega_w$ for its tame parts.
\item[(iii)] $\mu^\flat$ is a character of $\auU(1)$ of weight $-\kappa_0=-(\kappa+\frac\Delta2)$, unramified at all finite places:
{\cite[Lemma~6.13]{P54}} with $S=\emptyset$, $c_\tau=-\kappa(\tau)-r^8$; {\cite[(39)]{P54}} is {\cite[Lemma~6.12(c)]{P54}}.
\item[(iv)] $\mu^0_j=\mu^\flat\,(\upsilon^{-\Delta}\chi_M)^j$ for $1\le j\le m$, $\upsilon$ restricted to $\auU(1)$.
\end{enumerate}
\end{quoteddef}

\begin{quotedthm}{\cite[Proposition~6.15]{P54}}{ (Properties of the tame parts)}
\begin{enumerate}
\item[(T1)] $\mu^0_j$ has weight $-\kappa_j$, tame part $\omega_w^j$ at every chosen place $w$, and is unramified at all other finite
places; so it is a character as required in {\cite[\S14.3]{P46}}. The same characters are obtained, up to a character of weight $0$
unramified everywhere, from {\cite[Lemma~6.13]{P54}} with $c_\tau=-\kappa_j(\tau)$ and $\psi_w=\omega_w^j$: condition {\cite[(39)]{P54}}
holds for every $j$.
\item[(T2)] $\omega_w$ has order $q_2>8m$ at every chosen place; hence $\omega_w^{\,n}\ne1$ for $1\le|n|\le8m$, and in particular
$\omega_w^{\pm j}\ne1$ for $1\le j\le m$.
\item[(T3)] For every twist $\vartheta$, the tame part of $\mu=\mu^0_j\vartheta$ at a chosen place $v$ is $\psi_v=\omega_v^{\,j}$, and
$\psi_v^{q_2}=1$. This is the property used in (F2) of {\cite[Lemma~4.15]{P54}} ({\cite[Remark~4.16]{P54}}).
\item[(T4)] The character $\chi_0$ of {\cite[\S2.7]{P54}} is $\chi_M$, with tame part $\omega$ trivial off $\mathsf M$ and of order
$>8m$ at $\mathsf M$; the tame part of $\mu^0_j$ is $\omega^j$. The characters $\nu_\ast=\mu^0_j\upsilon^{n'}$, $\nu_\pm$,
$\chi'=\chi_M^j\vartheta$, $\nu''_\vartheta$ of {\cite[\S2.7]{P54}}, {\cite[\S3.2.6]{P54}} have the weights and tame parts stated there,
and $\nu^\flat=\nu''_\vartheta/\chi'=\mu^\flat\upsilon^{-j\Delta+j+\lfloor(t+k-1)/2\rfloor}$ is independent of $\vartheta$ with trivial tame part. At every
prime of $\pfqm$ the tame parts of $\nu''_\vartheta$ and $\chi'$ coincide ($\omega^j_w$ at $\mathsf M$, $\psi_w$ at $Q_\vartheta$).
\end{enumerate}
\end{quotedthm}

\begin{thebibliography}{P46}
\bibitem[P46]{P46} D.~A.~Trejo Pizzo, \emph{The sign of the radial intersection form on the idele class group}, preprint, 2026, \url{https://github.com/dtrejopizzo/riemann/tree/main/04-papers/46-row-d-sign-radial-intersection-form}.
\bibitem[P50]{P50} D.~A.~Trejo Pizzo, \emph{Elementary divisors of an edge Hecke operator on a Frobenius cycle}, preprint, 2026, \url{https://github.com/dtrejopizzo/riemann/tree/main/04-papers/50-elementary-divisors-edge-hecke-operator}.
\bibitem[P51]{P51} D.~A.~Trejo Pizzo, \emph{The Spanning Theorem: the second moment of a Poincar\'e map}, preprint, 2026, \url{https://github.com/dtrejopizzo/riemann/tree/main/04-papers/51-spanning-theorem-second-moment-poincare-map}.
\bibitem[P52]{P52} D.~A.~Trejo Pizzo, \emph{The Spanning Theorem: sieve systems at the real and finite places and Siegel's step}, preprint, 2026, \url{https://github.com/dtrejopizzo/riemann/tree/main/04-papers/52-sieve-systems-siegel-step}.
\bibitem[P53]{P53} D.~A.~Trejo Pizzo, \emph{The Covolume Theorem: exchange identities and the determinant of a chain of theta lifts}, preprint, 2026, \url{https://github.com/dtrejopizzo/riemann/tree/main/04-papers/53-covolume-theorem-exchange-identities-chain-determinant}.
\bibitem[P54]{P54} D.~A.~Trejo Pizzo, \emph{The Covolume Theorem: one theta lift, Hecke $L$-values and test forms of least valuation}, preprint, 2026, \url{https://github.com/dtrejopizzo/riemann/tree/main/04-papers/54-covolume-theorem-one-theta-lift-hecke-l-values}.
\bibitem[P55]{P55} D.~A.~Trejo Pizzo, \emph{The Covolume Theorem: regular families of twists, Selmer groups and Ribet's method}, preprint, 2026, \url{https://github.com/dtrejopizzo/riemann/tree/main/04-papers/55-covolume-theorem-regular-families-of-twists}.
\bibitem[P56]{P56} D.~A.~Trejo Pizzo, \emph{Explicit cup exponents for local Galois representations beyond the Fontaine--Laffaille range}, preprint, 2026, \url{https://github.com/dtrejopizzo/riemann/tree/main/04-papers/56-explicit-cup-exponents-local-galois-representations}.
\bibitem{BorelHC1962} A.~Borel and Harish-Chandra, \emph{Arithmetic subgroups of algebraic groups},
Ann.\ of Math.\ (2) \textbf{75} (1962), 485--535.
\bibitem{Banaszczyk} W.~Banaszczyk, \emph{New bounds in some transference theorems in the geometry
of numbers}, Math.\ Ann.\ \textbf{296} (1993), 625--635.
\bibitem{BannaiKobayashi} K.~Bannai and S.~Kobayashi, \emph{Algebraic theta functions and the
$p$-adic interpolation of Eisenstein--Kronecker numbers}, Duke Math.\ J.\ \textbf{153} (2010),
229--295.
\bibitem{Chellali} M.~Chellali, \emph{Congruences entre nombres de Bernoulli--Hurwitz dans le cas
supersingulier}, J.\ Number Theory \textbf{35} (1990), 157--179.
\bibitem{ChowlaSelberg} S.~Chowla and A.~Selberg, \emph{On Epstein's zeta-function}, J.\ reine
angew.\ Math.\ \textbf{227} (1967), 86--110.
\bibitem{GomezZhu} R.~Gomez and C.-B.~Zhu, \emph{Local theta lifting of generalized Whittaker
models associated to nilpotent orbits}, Geom.\ Funct.\ Anal.\ \textbf{24} (2014), 796--853.
\bibitem{GrossAMF} B.~H.~Gross, \emph{Algebraic modular forms}, Israel J.\ Math.\ \textbf{113}
(1999), 61--93.
\bibitem{HoweTheta} R.~Howe, \emph{$\theta$-series and invariant theory}, in: Automorphic Forms,
Representations and $L$-functions, Proc.\ Sympos.\ Pure Math.\ \textbf{33}, Part~1, Amer.\ Math.\
Soc., 1979, 275--285.
\bibitem{KashiwaraVergne} M.~Kashiwara and M.~Vergne, \emph{On the Segal--Shale--Weil
representations and harmonic polynomials}, Invent.\ Math.\ \textbf{44} (1978), 1--47.
\bibitem{KatzDivisibilities} N.~M.~Katz, \emph{Divisibilities, congruences, and Cartier duality},
J.\ Fac.\ Sci.\ Univ.\ Tokyo Sect.\ IA Math.\ \textbf{28} (1981), 667--678.
\bibitem{KudlaSeesaw} S.~S.~Kudla, \emph{Seesaw dual reductive pairs}, in: Automorphic Forms of
Several Variables (Katata, 1983), Progr.\ Math.\ \textbf{46}, Birkh\"auser, 1984, 244--268.
\bibitem{BLR} S.~Bosch, W.~L\"utkebohmert and M.~Raynaud, \emph{N\'eron Models}, Ergebnisse der
Mathematik und ihrer Grenzgebiete (3) \textbf{21}, Springer, 1990.
\bibitem{DEP} C.~De~Concini, D.~Eisenbud and C.~Procesi, \emph{Young diagrams and determinantal
varieties}, Invent.\ Math.\ \textbf{56} (1980), 129--165.
\bibitem{LangEF} S.~Lang, \emph{Elliptic Functions}, 2nd ed., Graduate Texts in Mathematics
\textbf{112}, Springer, 1987.
\bibitem{LubinTate} J.~Lubin and J.~Tate, \emph{Formal complex multiplication in local fields},
Ann.\ of Math.\ (2) \textbf{81} (1965), 380--387.
\bibitem{SilvermanAEC} J.~H.~Silverman, \emph{The Arithmetic of Elliptic Curves}, 2nd ed., Graduate
Texts in Mathematics \textbf{106}, Springer, 2009.
\bibitem{SilvermanAT} J.~H.~Silverman, \emph{Advanced Topics in the Arithmetic of Elliptic Curves},
Graduate Texts in Mathematics \textbf{151}, Springer, 1994.
\bibitem{HKS} M.~Harris, S.~S.~Kudla and W.~J.~Sweet, \emph{Theta dichotomy for unitary groups},
J.\ Amer.\ Math.\ Soc.\ \textbf{9} (1996), 941--1004.
\bibitem{KudlaSplitting} S.~S.~Kudla, \emph{Splitting metaplectic covers of dual reductive pairs},
Israel J.\ Math.\ \textbf{87} (1994), 361--401.
\bibitem{WeilActa} A.~Weil, \emph{Sur certains groupes d'op\'erateurs unitaires}, Acta Math.\
\textbf{111} (1964), 143--211.
\bibitem{FultonYT} W.~Fulton, \emph{Young Tableaux}, London Mathematical Society Student Texts
\textbf{35}, Cambridge University Press, 1997.
\bibitem{GoodmanWallach} R.~Goodman and N.~R.~Wallach, \emph{Symmetry, Representations, and
Invariants}, Graduate Texts in Mathematics \textbf{255}, Springer, 2009.
\bibitem{KubertLang} D.~S.~Kubert and S.~Lang, \emph{Modular Units}, Grundlehren der mathematischen
Wissenschaften \textbf{244}, Springer, 1981.
\bibitem{MumfordAV} D.~Mumford, \emph{Abelian Varieties}, Tata Institute of Fundamental Research
Studies in Mathematics \textbf{5}, Oxford University Press, 1970.
\bibitem{PlatonovRapinchuk} V.~Platonov and A.~Rapinchuk, \emph{Algebraic Groups and Number
Theory}, Pure and Applied Mathematics \textbf{139}, Academic Press, 1994.
\bibitem{TitsCorvallis} J.~Tits, \emph{Reductive groups over local fields}, in: Automorphic Forms,
Representations and $L$-Functions (Corvallis, 1977), Proc.\ Sympos.\ Pure Math.\ \textbf{33},
Part~1, Amer.\ Math.\ Soc., 1979, pp.~29--69.
\bibitem{RallisFunctoriality} S.~Rallis, \emph{Langlands' functoriality and the Weil
representation}, Amer.\ J.\ Math.\ \textbf{104} (1982), 469--515.
\bibitem{CoxPrimes} D.~A.~Cox, \emph{Primes of the Form $x^2+ny^2$: Fermat, Class Field Theory, and
Complex Multiplication}, Wiley, New York, 1989.
\bibitem{SerreLF} J.-P.~Serre, \emph{Local Fields}, Graduate Texts in Mathematics 67, Springer,
New York, 1979.
\bibitem{HarishChandraAF} Harish-Chandra, \emph{Automorphic Forms on Semisimple Lie Groups}, Lecture
Notes in Mathematics 62, Springer, Berlin, 1968.
\bibitem{BorelJacquet} A.~Borel and H.~Jacquet, \emph{Automorphic forms and automorphic
representations}, in: Automorphic Forms, Representations and $L$-functions, Proc.\ Sympos.\ Pure
Math.\ \textbf{33}, Part~1, Amer.\ Math.\ Soc., 1979, 189--202.
\bibitem{HoweJAMS} R.~Howe, \emph{Transcending classical invariant theory}, J.\ Amer.\ Math.\ Soc.\
\textbf{2} (1989), 535--552.
\bibitem{FJ} M.~Flensted-Jensen, Discrete series for semisimple symmetric spaces, Ann. of Math. (2) 111 (1980), 253--311.
\bibitem{Knapp} A.~W.~Knapp, Representation Theory of Semisimple Groups: An Overview Based on Examples, Princeton Univ. Press, 1986.
\bibitem{Parthasarathy72} R.~Parthasarathy, Dirac operator and the discrete series, Ann. of Math. (2) 96 (1972), 1--30.
\bibitem{HuangPandzic} J.-S.~Huang and P.~Pand\v zi\'c, Dirac cohomology, unitary representations and a proof of a conjecture of Vogan, J. Amer. Math. Soc. 15 (2002), 185--202.
\bibitem{Schmid75} W.~Schmid, Some properties of square-integrable representations of semisimple Lie groups, Ann. of Math. (2) 102 (1975), 535--564.
\bibitem{Dixmier} J.~Dixmier, $C^*$-algebras, North-Holland, 1977.
\bibitem{VoganZuckerman} D.~A.~Vogan, Jr. and G.~J.~Zuckerman, Unitary representations with non-zero cohomology, Compositio Math. 53 (1984), 51--90.
\bibitem{KnappBeyond} A.~W.~Knapp, Lie Groups Beyond an Introduction, 2nd ed., Birkh\"auser, 2002.
\bibitem{FollandAHA} G.~B.~Folland, A Course in Abstract Harmonic Analysis, CRC Press, 1995.
\bibitem{BorelFin} A.~Borel, \emph{Some finiteness properties of adele groups over number
fields}, Publ.\ Math.\ Inst.\ Hautes \'Etudes Sci.\ \textbf{16} (1963), 5--30.
\bibitem{Helgason} S.~Helgason, \emph{Groups and Geometric Analysis}, Academic Press, 1984.
\bibitem{WeilSiegel} A.~Weil, \emph{Sur la formule de Siegel dans la th\'eorie des groupes
classiques}, Acta Math.\ \textbf{113} (1965), 1--87.
\bibitem{WallachRRG1} N.~R.~Wallach, \emph{Real Reductive Groups I}, Pure and Applied Mathematics 132, Academic Press, 1988.
\end{thebibliography}

\end{document}
